% !TeX program = pdflatex
% Plik generowany: tekst umowy z psa.tex + odpowiedzi memorandum (niebieskie ramki).
% Pełny tekst odpowiedzi: memorandum.md. Kompilacja: pdflatex memorandum.tex (dwukrotnie).
\documentclass[12pt,a4paper]{report}
\usepackage[T1]{fontenc}
\usepackage[utf8]{inputenc}
\usepackage{basilisk_i18n}
\usepackage{lmodern}
\usepackage{microtype}
\usepackage{setspace}
\usepackage{parskip}
\usepackage{needspace}
\usepackage{tabularx}
\usepackage[normalem]{ulem}

\BusinessPlanCompany{Basilisk Systems}
\BusinessPlanWatermarkText{Basilisk Systems}
\BusinessPlanWatermarkColor{black!6}
\BusinessPlanWatermarkFontSize{38}{46}
\BusinessPlanWatermarkAngle{38}
\BusinessPlanAuthorsText{Basilisk Systems --- memorandum do projektu umowy P.S.A.}
\BusinessPlanConfidentialText{Projekt do weryfikacji prawnej --- nie stanowi części umowy}
\BusinessPlanFooterText{Basilisk Systems --- memorandum do projektu umowy P.S.A.}
\BusinessPlanVersionText{Memorandum 1.2 --- umowa 0.9.4-C z komentarzem}

\setstretch{1.18}
\setlist[itemize]{topsep=0.2em,itemsep=0.15em,parsep=0.1em,leftmargin=2.0em}
\setlist[enumerate]{topsep=0.25em,itemsep=0.18em,parsep=0.1em,leftmargin=2.2em}
\setlength{\parindent}{0pt}
\setlength{\parskip}{0.55em}
\setlength{\emergencystretch}{4em}
\renewcommand{\arraystretch}{1.18}

\newcounter{paragraf}
\newcommand{\paragraf}[1]{%
  \refstepcounter{paragraf}%
  \phantomsection%
  \addcontentsline{toc}{section}{\S~\theparagraf\ #1}%
  \Needspace{6\baselineskip}%
  \subsection*{\S~\theparagraf\ #1}%
}
\newenvironment{ContractClause}[1]{%
  \paragraf{#1}%
}{%
  \par
}

\newcommand{\field}[1]{\textbf{[#1]}}
\newcommand{\CompanyFull}{Basilisk Systems prosta spółka akcyjna}
\newcommand{\CompanyShort}{Basilisk Systems P.S.A.}
\newcommand{\KSH}{Kodeks spółek handlowych}
\newcommand{\BoilerplateStrike}[1]{\sout{#1}}

\hypersetup{
  pdftitle={Basilisk Systems P.S.A. --- memorandum do projektu umowy (umowa z komentarzem)},
  pdfauthor={Basilisk Systems --- memorandum do projektu umowy P.S.A.}
}


% --- memorandum: kolory i makra -------------------------------------------
\definecolor{memoblue}{RGB}{0,70,160}
\newtcolorbox{MemoBox}[1]{breakable,enhanced,colback=memoblue!4,colframe=memoblue!70,coltext=memoblue!85!black,
  boxrule=0.4pt,leftrule=2.4pt,arc=1pt,left=6pt,right=6pt,top=5pt,bottom=5pt,before skip=0.6em,after skip=0.8em,
  fontupper=\small,title={Memorandum --- pytanie #1},fonttitle=\bfseries\small,coltitle=white,colbacktitle=memoblue!80,
  toptitle=1pt,bottomtitle=1pt}
\newcommand{\MemoPyt}[2]{\par\noindent\phantomsection\label{memo:#1}\textbf{#1} \textit{#2}\par\smallskip}
\newcommand{\MemoPytCd}[1]{\par\noindent\textbf{#1 (cd.)}\par\smallskip}
\newcommand{\Wiar}[1]{\textbf{[#1]}}
\newcommand{\Ocena}[1]{\textbf{Ocena:} #1\par}
\newcommand{\Podstawa}[1]{\textbf{Podstawa prawna:} #1\par}
\newcommand{\Uzasadnienie}[1]{\textbf{Uzasadnienie.} #1\par}
\newcommand{\Rekomendacja}[1]{\textbf{Rekomendacja.} #1\par}
\newcommand{\Klasa}[1]{\textbf{Klasyfikacja:} #1\par}
\newcommand{\Kierunek}[1]{\textbf{Kierunek i konsekwencje.} #1\par}
\newenvironment{Brzmienie}{\par\textbf{Proponowane brzmienie:}\begin{quote}\itshape}{\end{quote}}
\newcommand{\Pytanie}[2]{\Needspace{6\baselineskip}\subsubsection*{#1}\textit{#2}\par\medskip}

\begin{document}

\BusinessPlanTitlePage
  {Basilisk Systems}
  {Memorandum 1.2 --- umowa 0.9.4-C z komentarzem}
  {Memorandum}
  {Memorandum do projektu umowy prostej spółki akcyjnej}
  {Tekst umowy z odpowiedziami memorandum przy właściwych postanowieniach}
  {1.2}

\tableofcontents
\newpage
\chapter*{Jak czytać to memorandum}
\addcontentsline{toc}{chapter}{Jak czytać to memorandum}

\textbf{Układ tego dokumentu.} Czarny tekst to projekt umowy w wersji 0.9.4-C, bez zmian (wraz z przekreśleniami i polami do uzupełnienia, jak w \texttt{psa.tex}). \textcolor{memoblue}{Niebieskie ramki to odpowiedzi memorandum}, wstawione przy ustępie, którego dotyczą; każda zaczyna się od numeru i treści pytania z listy szczegółowej, podaje ocenę, klasyfikację uwagi (z numerem K-n dla uwag koniecznych), podstawę prawną z oznaczeniem wiarygodności, zwięzłe uzasadnienie i rekomendację, a przy uwagach koniecznych --- pełne proponowane brzmienie. Pełne odpowiedzi, z całym wywodem, są w pliku \texttt{memorandum.md}. Po tekście umowy następują sekcje tematyczne T1--T9 i załączniki: nr 1 zbiorcza tabela uwag, nr 2 decyzje Założycieli, nr 3 źródła i lista do weryfikacji, nr 4 zmiany dla wariantu B, nr 5 korekty wprowadzone przy weryfikacji, nr 6 indeks pytań z numerami stron.

\textbf{Przedmiot.} Memorandum odpowiada na 79 pytań z listy szczegółowej (\texttt{law.md}, 21 września 2026 r.) do projektu umowy \CompanyFull{} w wersji 0.9.4-C (wariant hybrydowy), w numeracji tej listy, oraz realizuje pozostałe punkty zakresu A--D z zapytania (\texttt{email.md}): zmiany przepisów, wykonalność mechanizmów, ścieżka zawiązania, dokumentacja wkładów, rejestr akcjonariuszy, PKD, oczekiwania inwestorów, podział materii między umowę spółki a umowę akcjonariuszy oraz ryzyka podatkowe. Odesłania do paragrafów (§) odpowiadają numeracji umowy w tym dokumencie i w \texttt{psa.tex}.

\textbf{Status.} Dokument został przygotowany jako materiał roboczy w roli recenzenta prawnego; wersja 1.0 powstała bez dostępu do baz aktów prawnych, wersja 1.1 została zweryfikowana na tekstach aktów krajowych, wersja 1.2 na tekstach prawa Unii i materiałach PFR Ventures. Nie jest opinią prawną i nie zastępuje weryfikacji przez radcę prawnego albo adwokata przed podpisaniem umowy. Każda podstawa prawna niesie oznaczenie wiarygodności:
\begin{itemize}
  \item \Wiar{Z} --- teza potwierdzona w tekście aktu albo w źródle dostępnym w toku prac (\texttt{src/legal/});
  \item \Wiar{W} --- wiedza ogólna o treści przepisu lub praktyce; numer artykułu i brzmienie do sprawdzenia w tekście jednolitym;
  \item \Wiar{?} --- hipoteza; wymaga potwierdzenia przed jakąkolwiek decyzją.
\end{itemize}
Załącznik nr~3 zbiera wszystkie punkty do weryfikacji przez uprawnionego prawnika.

\textbf{Dwie skale.} Lista pytań prosi o ocenę stanu (\emph{zgodne} / \emph{zgodne warunkowo} / \emph{ryzyko}), a zakres prac o priorytet zmiany (\emph{konieczne} / \emph{zalecane} / \emph{opcjonalne}). Stosujemy obie: ocena opisuje projekt, klasyfikacja mówi, co z nim zrobić. Reguła przejścia: „ryzyko” daje uwagę konieczną; „zgodne warunkowo” daje uwagę zalecaną, chyba że warunek spełnia się poza tekstem umowy (wtedy opcjonalna albo bez zmian); „zgodne” nie rodzi uwagi, chyba że proponujemy ulepszenie (opcjonalna). Odstępstwa są uzasadnione przy odpowiedzi. Dla każdej uwagi koniecznej podajemy pełne proponowane brzmienie gotowe do wklejenia do umowy; dla zalecanych i opcjonalnych --- kierunek i konsekwencje, a tam, gdzie zmiana jest prosta, także brzmienie.

\textbf{Stan prawny.} Odpowiedzi odnoszą się do stanu prawnego na 30 września 2026 r. Wersja 1.1 memorandum została zweryfikowana na pełnych tekstach aktów krajowych pobranych do katalogu \texttt{src/legal/} (teksty ujednolicone KSH, KC, KPC, KRO, prawa autorskiego, PWP, PIT, PCC, ustawy o KRS, o kontroli niektórych inwestycji, AML, ustawy koncesyjnej z 2019 r., wykazu WT, PKD 2025, nowelizacji KSH z 23 stycznia 2026 r. i ustawy z 13 marca 2026 r.); Wersja 1.2 została dodatkowo zweryfikowana na tekstach prawa Unii w \texttt{src/legal/eu/}: rozporządzenie 2021/697 (EDF) w wersji polskiej oraz --- wyłącznie w wersji angielskiej, co oznaczono dopiskiem „(tekst EN)” --- rozporządzenia 2021/821 (konsolidacja z 26 maja 2023 r.; załącznik I był później zmieniany), 269/2014 (tekst pierwotny), 2019/947, 2024/1689 (AI Act), 2026/877 (transfer technologii) i 2026/1386 (kontrola inwestycji zagranicznych, uchyla 2019/452 z dniem 17 stycznia 2028 r.). Tezy o TFUE, rozporządzeniach 2019/452, 2024/2847, 2023/1230, 833/2014, 2018/1139 i 2019/945 oraz o orzecznictwie SN pozostają oznaczone \Wiar{W} (lista w \texttt{src/legal/todo.md}). Ustawa z 23 stycznia 2026 r. o zmianie \KSH{} (Dz.U. 2026 poz. 176) wchodzi w życie 18 lutego 2027 r.; jej zmiany dotyczące P.S.A. ograniczają się do rejestru akcjonariuszy i są omówione w odpowiedzi 4.3.3 oraz w sekcji T1.


\textbf{Skróty.} KSH --- Kodeks spółek handlowych; KC --- Kodeks cywilny; KPC --- Kodeks postępowania cywilnego; KRO --- Kodeks rodzinny i opiekuńczy; PWP --- Prawo własności przemysłowej; PIT --- ustawa o podatku dochodowym od osób fizycznych; PCC --- podatek od czynności cywilnoprawnych; KRS --- Krajowy Rejestr Sądowy; PRS --- Portal Rejestrów Sądowych; WZ --- Walne Zgromadzenie; Rada --- Rada Dyrektorów; Kryterium --- Kryterium Bezpieczeństwa EU/NATO z §~\ref{art:dopuszczalny-nabywca}; EDF --- Europejski Fundusz Obronny; DIANA --- akcelerator NATO; ASI --- alternatywna spółka inwestycyjna.


\clearpage
\chapter*{Podsumowanie wykonawcze}
\addcontentsline{toc}{chapter}{Podsumowanie wykonawcze}

\textbf{Ocena ogólna.} Projekt umowy w wersji 0.9.4-C jest konstrukcyjnie spójny i nie zawiera postanowienia, które byłoby wprost sprzeczne z bezwzględnie obowiązującymi przepisami o prostej spółce akcyjnej. Kluczowe mechanizmy (Kryterium Bezpieczeństwa EU/NATO, zgoda Spółki na rozporządzenie akcją, mechanizm zwrotnego zbycia, ograniczenie wstąpienia spadkobierców, serie P i F, tryb rundy) mieszczą się w swobodzie kształtowania umowy P.S.A. Ryzyka koncentrują się nie w dopuszczalności, lecz w \emph{wykonalności}: kilka mechanizmów w obecnym brzmieniu nie zadziała bez podpisu drugiej strony, nie ma oznaczonego dłużnika albo terminu, albo zależy od decyzji jednego organu tam, gdzie ustawa lub uczciwość obrotu wymaga procedury. Na 79 pytań odpowiedzi „ryzyko” padły w kilkunastu przypadkach, a uwag koniecznych jest \textbf{18}; uwag zalecanych \textbf{61}, opcjonalnych \textbf{19} (Załącznik nr~1).

\textbf{Uwagi konieczne} (numeracja K-n używana dalej w załącznikach; pełne brzmienia przy odpowiedziach):
\begin{enumerate}[label=K\arabic*.,leftmargin=2.6em]
  \item (\hyperref[memo:1.2.4]{1.2.4}, §~\ref{art:zbywanie} ust. 3 (wariant B)) W wariancie B dodać wentyl: po 12 miesiącach bezskutecznego poszukiwania nabywcy akcjonariusz może żądać wskazania nabywcy przez Spółkę.
  \item (\hyperref[memo:1.5.1]{1.5.1}, §~\ref{art:zbywanie} ust. 2) Zamknięty katalog odmowy zgody nie pozwala odmówić zbycia akcji objętych obowiązkiem zbycia nabywcy, który nie przystąpił do umowy wykonawczej; dodać lit. f.
  \item (\hyperref[memo:1.5.1]{1.5.1}, §~\ref{art:zwrotne-zbycie} ust. 14) Obowiązek zbycia ma obciążać każdoczesnego posiadacza akcji, a Spółka ma żądać wpisu wzmianek w rejestrze na podstawie art.~300$^{33}$ KSH.
  \item (\hyperref[memo:1.5.3]{1.5.3}, §~\ref{art:zwrotne-zbycie} ust. 13) Rozszerzyć obowiązkową treść umowy wykonawczej (oferta do Spółki lub umowa przedwstępna zredagowana jako podstawa wpisu z art.~300$^{34}$ §~4 KSH, pełnomocnik oznaczony z podpisem poświadczonym wobec art.~300$^{34}$ §~3 KSH od 18.02.2027, wstrzymanie 6-miesięcznego terminu, skutki niezapłacenia ceny, dopłata po zmianie kwalifikacji, zbieg z prawem pierwszeństwa).
  \item (\hyperref[memo:1.6.1]{1.6.1}, §~\ref{art:dziedziczenie} ust. 2) Wskazać dłużnika spłaty (nabywca, solidarnie Spółka), termin wskazania nabywcy i ścieżkę rezerwową przez nabycie akcji własnych lub umorzenie z określeniem niewstąpienia jako przesłanki umorzenia (art.~300$^{45}$ §~1 KSH), pod rygorem bezskuteczności ograniczenia z art.~300$^{41}$ §~1 KSH.
  \item (\hyperref[memo:1.6.1]{1.6.1}, §~\ref{art:dziedziczenie} ust. 2) Akcje serii F objęte za wkład pracy niewniesiony w całości: przesądzić, że wstąpienie spadkobierców podlega §17 zamiast zgody Spółki z art.~300$^{41}$ §~2 KSH, i dodać proporcjonalne pomniejszenie spłaty (art.~300$^{41}$ §~1 zd. 3 KSH).
  \item (\hyperref[memo:1.6.1]{1.6.1}, §~\ref{art:dziedziczenie} ust. 4) Bezwarunkowy termin pierwszej raty spłaty (30 dni od ustalenia Wartości Godziwej) zamiast przesłanki zagrożenia płynności bez wskazania, kto ją ocenia.
  \item (\hyperref[memo:1.6.3]{1.6.3}, §~\ref{art:dziedziczenie} ust. 2a) Dodać tryb wezwania i biegu 90 dni, uprawnienie Spółki do wszczęcia postępowania spadkowego, stwierdzenie niewstąpienia uchwałą Rady oraz umorzenie automatyczne z art.~300$^{46}$ KSH (z zastrzeżeniem środków z art.~300$^{15}$ §~2 KSH) jako ścieżkę wobec spadkobiercy niewspółdziałającego.
  \item (\hyperref[memo:1.7.3]{1.7.3}, §~\ref{art:dopuszczalny-nabywca} ust. 11a) Tryb stwierdzenia statusu Prawnie Niedopuszczalnego Posiadacza, wstrzymania świadczeń i niedopuszczenia do wykonywania praw oraz wyłączenie głosów zawieszonych z mianownika większości umownych, z zastrzeżeniem praw zarządcy przymusowego (art. 6a ustawy z 13.04.2022).
  \item (\hyperref[memo:1.8.2]{1.8.2}, §~\ref{art:wycena} ust. 4) Wybór Niezależnego Eksperta wspólnie przez strony, rezerwowo przez instytucję neutralną zamiast Rady Dyrektorów; procedura kontradyktoryjna z terminem 60 dni.
  \item (\hyperref[memo:1.8.2]{1.8.2}, §~\ref{art:wycena} ust. 7) Dodać wyjątek oczywistej sprzeczności z zasadami wyceny i tryb ponownej wyceny.
  \item (\hyperref[memo:3.1.2]{3.1.2}, §~\ref{art:wklady} ust. 4) Przeznaczyć na kapitał akcyjny także wkłady niepieniężne (art.~300$^{3}$ §~1 KSH); wyłączenie ich z kapitału grozi odmową wpisu i przychodem CIT po stronie Spółki.
  \item (\hyperref[memo:3.1.4]{3.1.4}, Załącznik nr 1) Rozdzielić numery akcji obejmowanych za wkłady niepieniężne i pieniężne (art.~300$^{5}$ §~1 pkt 4 KSH) i dodać kolumnę wartości wkładu z podstawą wyceny.
  \item (\hyperref[memo:3.2.1]{3.2.1}, §~\ref{art:program-motywacyjny} ust. 3) Oprzeć emisje serii P na ustawowych instrumentach: warunkowej emisji akcji (art.~300$^{114}$–300$^{118}$ KSH) jako trybie podstawowym albo uchwale zwykłej z pełnym katalogiem art.~300$^{104}$ KSH; wyłączyć prawo poboru serii P w umowie na podstawie art.~300$^{106}$ §~1 KSH (ust. 5).
  \item (\hyperref[memo:3.3.2]{3.3.2}, §~\ref{art:kwalifikowana-runda} ust. 4) Rozróżnić warunki wymagające zmiany umowy w protokole notarialnym (uprzywilejowanie, prawo wskazania dyrektora, zgoda serii) od kontraktowych; obecne brzmienie sugeruje preferencję bez zmiany umowy.
  \item (\hyperref[memo:4.3.1]{4.3.1}, §~\ref{art:transakcje-akcjonariusze}) Usunąć oznaczenie robocze „[WYMAGA OPINII PRAWNIKA]” przed sporządzeniem aktu notarialnego.
  \item (\hyperref[memo:4.3.1]{4.3.1}, §~\ref{art:final} ust. 1) Usunąć wszystkie przekreślenia (po decyzjach Założycieli) i wypełnić wszystkie pola; akt notarialny nie może zawierać tekstu przekreślonego ani adnotacji roboczych.
  \item (\hyperref[memo:4.3.1]{4.3.1}, §~\ref{art:final} ust. 7) Dodać ustęp o uchwałach założycielskich objętych aktem (pierwsza Rada, podmiot prowadzący rejestr) i o reprezentacji spółki w organizacji (art.~300$^{11}$ §~2 KSH).
\end{enumerate}

\textbf{Siedem rozstrzygnięć, które wynikają z memorandum.}
\begin{enumerate}
  \item \textbf{Wariant A (ratalny) zamiast wariantu B.} Rekomendujemy utrzymanie mechanizmu z §~\ref{art:zbywanie} ust.~3 z poprawkami z odpowiedzi \hyperref[memo:1.2.1]{1.2.1}--\hyperref[memo:1.2.3]{1.2.3} (termin końcowy wyceny, warunki rat, zaświadczenie Spółki dla rejestru, skutek milczenia Rady). Wariant B jest dopuszczalny, ale bez wentyla czasowego grozi zarzutem „akcji uwięzionej” i częściową nieważnością, której skutek byłby dla Spółki gorszy niż słabość wariantu A; jeżeli Założyciele wybiorą B, konieczny jest wentyl z odpowiedzi \hyperref[memo:1.2.4]{1.2.4} (Załącznik nr~4). Dla inwestorów wariant A jest łatwiejszy do obrony (\hyperref[memo:4.4.1]{4.4.1}).
  \item \textbf{Przywrócić wyjątek dla funduszy} (§~\ref{art:dopuszczalny-nabywca} ust.~14) w brzmieniu uszczelnionym: test na poziomie zarządzającego, siedziby i kontroli, wyłączenie wehikułów jednego lub dwóch inwestorów, ujawnienie beneficjentów rzeczywistych w rozumieniu AML i inwestorów powyżej progu, oświadczenie funduszu, a zmiana po stronie pasywnego inwestora funduszu nie jest utratą Kryterium z ust.~15 (odpowiedzi \hyperref[memo:1.4.1]{1.4.1}--\hyperref[memo:1.4.3]{1.4.3}, \hyperref[memo:4.4.2]{4.4.2}). Bez tego każdy fundusz z inwestorami spoza UE/EOG/NATO potrzebuje zgody Walnego Zgromadzenia, a fundusze obronne odrzucą badanie wszystkich inwestorów.
  \item \textbf{Kryterium a kontrola z NATO spoza UE.} Kryterium jest dopuszczalne na tle art.~300$^{39}$ KSH i prawa Unii, ale dopuszcza kontrolę nad Spółką przez podmiot z USA, Wielkiej Brytanii, Turcji lub Kanady, co wyklucza Spółkę z EDF, AGILE i EUDIS. Rekomendujemy nie zawężać Kryterium w umowie spółki, lecz zapisać w umowie akcjonariuszy warunek: dopóki Spółka korzysta z tych programów, akcjonariusze spoza UE/EOG nie uzyskują kontroli ani weta (decyzja D6; odpowiedź \hyperref[memo:1.1.1]{1.1.1}, sekcje T1a i T8).
  \item \textbf{Wszystkie wkłady na kapitał akcyjny.} Obecne §~\ref{art:wklady} ust.~4 (tylko wkłady pieniężne na kapitał) jest sprzeczne z przeważającą wykładnią art.~300$^{3}$ §~1 KSH, nie ma odpowiednika bilansowego i grozi wezwaniem sądu oraz przychodem po stronie Spółki; a umowa P.S.A. według wykładni literalnej ustawy o PCC w ogóle nie podlega temu podatkowi, więc nie ma czego oszczędzać (\hyperref[memo:3.1.2]{3.1.2}, T4, T9). Załącznik nr~1 wymaga przebudowy (\hyperref[memo:3.1.4]{3.1.4}).
  \item \textbf{Wykonalność zwrotnego zbycia, dziedziczenia i wyceny.} Trzy mechanizmy wymagają domknięcia, żeby działały bez podpisu drugiej strony: umowa wykonawcza z ofertą do Spółki albo umową przedwstępną i pełnomocnikiem oznaczonym (\hyperref[memo:1.5.3]{1.5.3}), spłata spadkobierców z oznaczonym dłużnikiem, terminem i ścieżką rezerwową przez umorzenie (\hyperref[memo:1.6.1]{1.6.1}, \hyperref[memo:1.6.3]{1.6.3}), oraz wybór Niezależnego Eksperta wspólnie albo przez instytucję neutralną zamiast Rady (\hyperref[memo:1.8.2]{1.8.2}). Do tego tryb dla Prawnie Niedopuszczalnego Posiadacza z wyłączeniem głosów zawieszonych z mianownika większości (\hyperref[memo:1.7.3]{1.7.3}).
  \item \textbf{Seria P na ustawowych instrumentach emisyjnych.} Weryfikacja obaliła tezę wersji 1.0, że P.S.A. nie zna delegacji emisyjnej: KSH przewiduje upoważnienie Rady do emisji (art.~300$^{110}$--300$^{113}$, do 5 lat i 1/4 akcji) oraz warunkową emisję akcji (art.~300$^{114}$--300$^{118}$), a prawo poboru jest dyspozytywne (art.~300$^{106}$ §~1). Program motywacyjny należy oprzeć na warunkowej emisji akcji albo na uchwale z pełnym katalogiem art.~300$^{104}$, a prawo poboru serii P wyłączyć w umowie; obecna „uchwała ramowa” w §~\ref{art:program-motywacyjny} ust.~3 bez tych elementów jest ryzykiem (\hyperref[memo:3.2.1]{3.2.1}, uwaga konieczna; decyzja D11).
  \item \textbf{Usunąć aparat roboczy przed aktem.} Siedem przekreślonych zdań można usunąć bez straty (\hyperref[memo:3.4.4]{3.4.4}) z poprawką sztywnych odesłań w §~\ref{art:wz}; §~\ref{art:dopuszczalny-nabywca} ust.~14 przywrócić w nowym brzmieniu; znacznik „[WYMAGA OPINII PRAWNIKA]” w §~\ref{art:transakcje-akcjonariusze} usunąć po wdrożeniu uwag z \hyperref[memo:2.3.4]{2.3.4}; wypełnić 58 pól (\hyperref[memo:4.3.1]{4.3.1}, T3).
\end{enumerate}

\textbf{Co nie wymaga zmian.} Firma, siedziba, cel, struktura akcji serii A, program motywacyjny co do zasady (z zastrzeżeniem trybu emisji serii P, \hyperref[memo:3.2.1]{3.2.1}), seria F, prawo pierwszeństwa, tag-along, małżonek, organy w systemie monistycznym, reprezentacja, własność intelektualna (warstwa korporacyjna), poufność, kontrola eksportu jako ramy, finansowanie, zysk, impas, likwidacja. Audyt wersji 0.8 z \texttt{psa\_feedback.tex} Część 2 pozostaje w tych punktach aktualny.

\textbf{Zmiany przepisów.} Nowelizacja KSH z 23 stycznia 2026 r. (w życie 18 lutego 2027 r.), sprawdzona w pełnym tekście, nie zmienia przepisów, na których opierają się mechanizmy umowy (art.~300$^{39}$, 300$^{41}$, 300$^{45}$--300$^{47}$, 300$^{103}$--300$^{107}$); dotyczy rejestru akcjonariuszy: zgłoszenie podmiotu prowadzącego rejestr do KRS, adres e-mail w rejestrze tylko za zgodą akcjonariusza, 7-dniowy termin zgłaszania zmian danych, forma zgody na wpis, PESEL niejawny dla pozostałych akcjonariuszy, rozszerzony katalog nabyć poza wpisem. Nie ma powodu odkładać zawiązania; spółka zawiązana wcześniej ma 2 lata na dostosowanie umowy. Zalecamy klauzulę dostosowawczą, przeredagowanie §~\ref{art:wz} ust.~3 (zgoda na komunikację e-mail) i dostosowanie formy pełnomocnictw w umowie wykonawczej (\hyperref[memo:4.3.3]{4.3.3}, T1). Nowe rozporządzenie o transferze technologii (2026/877) wymaga zmiany klauzuli o ulepszeniach w §~\ref{art:przedsiewziecia} ust.~3 przed pierwszą umową Przedsięwzięcia (\hyperref[memo:2.4.1]{2.4.1}). Ustawa z 13 marca 2026 r. zmienia ustawę koncesyjną z 2019 r. w zakresie omówionym w 4.1.1. \textbf{Nowe rozporządzenie (UE) 2026/1386 z 17 czerwca 2026 r. o kontroli inwestycji zagranicznych} (sprawdzone w tekście angielskim) zastępuje rozporządzenie 2019/452 z dniem 17 stycznia 2028 r., a część przepisów o współpracy stosuje się od 16 lipca 2026 r. Trzy zmiany dotykają Spółki wprost: każde państwo członkowskie musi wymagać uprzedniej zgody na inwestycję zagraniczną w podmiot, który rozwija, produkuje lub komercjalizuje towary podwójnego zastosowania z załącznika I do 2021/821 albo wyroby wojskowe z załącznika do dyrektywy 2009/43/WE (art.~4 ust.~15 lit.~a--b), a inwestycja w uczestnika programu Unii, w tym EDF, podlega dodatkowo opinii Komisji; „inwestor zagraniczny” to osoba bez obywatelstwa państwa członkowskiego albo podmiot z państwa trzeciego, a inwestycja przez spółkę z Unii kontrolowaną przez taki podmiot jest objęta tak samo jak inwestycja bezpośrednia; inwestycje typu greenfield nie wymagają uprzedniej zgody (art.~4 ust.~17). Polski próg podmiotowy (przychód powyżej 10 mln EUR) przestanie chronić Spółkę przed obowiązkiem zgłoszenia, gdy ustawa krajowa zostanie dostosowana; runda z inwestorem spoza Unii (także z NATO) po 17 stycznia 2028 r. powinna zakładać postępowanie kontrolne (\hyperref[memo:1.1.2]{1.1.2}, \hyperref[memo:1.4.1]{1.4.1}, T1).

\textbf{Zastrzeżenie.} Wersja 1.2 została zweryfikowana na pełnych tekstach aktów krajowych (\texttt{src/legal/}) oraz na tekstach prawa Unii i stronach PFR Ventures dostępnych w \texttt{src/legal/eu/} i \texttt{src/legal/pfr/} (część aktów Unii wyłącznie po angielsku, oznaczona „(tekst EN)”); TFUE, pozostałe akty Unii, orzecznictwo SN i tezy o praktyce rynkowej pozostają oznaczone \Wiar{W} i wymagają sprawdzenia według listy w Załączniku nr~3. Memorandum nadal nie zastępuje opinii uprawnionego prawnika przed zmianą umowy. Pełny tekst memorandum (\texttt{memorandum.md}) zawiera 736 oznaczeń \textbf{[Z]}, 184 oznaczeń \textbf{[W]} i 34 oznaczeń \textbf{[?]} (w wersji 1.0: 194, 370 i 36; w wersji 1.1: 607, 163 i 22); pozostałe oznaczenia \textbf{[W]} dotyczą TFUE, aktów Unii bez pobranego tekstu (2019/452, CRA, rozporządzenie maszynowe, 833/2014, 2018/1139, 2019/945), orzecznictwa SN i praktyki rynkowej.

Korekty wprowadzone przy weryfikacji na tekstach aktów (wersje 1.1 i 1.2) wymienia Załącznik nr~5.

\textbf{Decyzje Założycieli} (D1--D16, Załącznik nr~2) warunkują etap wdrożenia: jurysdykcja, wariant A/B, wyjątek dla funduszy, przekreślenia, Kryterium wobec spadkobierców, warstwa EDF, PKD przeważające, osoby objęte mechanizmem zwrotnego zbycia, progi rundy, podmiot prowadzący rejestr, model programu motywacyjnego, podział na umowę wykonawczą i akcjonariuszy, data zawiązania, doradcy, zakres wdrożenia uwag zalecanych i opcjonalnych, dane do Załączników.


\clearpage
\chapter*{Projekt umowy prostej spółki akcyjnej}
\addcontentsline{toc}{chapter}{Projekt umowy prostej spółki akcyjnej}

\begin{center}
{\Large\bfseries UMOWA PROSTEJ SPÓŁKI AKCYJNEJ}\par
\vspace{0.25cm}
{\LARGE\bfseries \CompanyFull}\par
\vspace{0.6cm}
\end{center}

Stawający, wskazani szczegółowo w Załączniku nr 1 do niniejszej umowy, zwani dalej łącznie „Założycielami” albo „Akcjonariuszami Założycielami”, zawiązują prostą spółkę akcyjną na następujących zasadach.

\begin{ContractClause}{Firma, forma prawna i czas trwania}
\label{art:firma}
\begin{enumerate}
  \item Firma Spółki brzmi: „\CompanyFull”.
  \item Spółka może posługiwać się skrótem: „\CompanyShort”.
  \item Spółka może używać wyróżniających oznaczeń handlowych, znaków towarowych, domen i skrótów, w szczególności oznaczenia „Basilisk Systems”.
  \item Spółka zostaje zawiązana na czas nieoznaczony.
\end{enumerate}

\end{ContractClause}

\begin{ContractClause}{Siedziba i obszar działania}
\label{art:siedziba}
\begin{enumerate}
  \item Siedzibą Spółki jest Pruszków.
  \item Spółka może prowadzić działalność w Rzeczypospolitej Polskiej i za granicą, z zastrzeżeniem ograniczeń dotyczących bezpieczeństwa, kontroli eksportu, sankcji oraz dopuszczalnych jurysdykcji określonych w niniejszej umowie.
  \item Spółka może tworzyć oddziały, zakłady, laboratoria, centra badawczo-rozwojowe, przedstawicielstwa i spółki zależne, a także uczestniczyć w innych podmiotach, o ile nie prowadzi to do obejścia §~\ref{art:dopuszczalny-nabywca}.
\end{enumerate}

\end{ContractClause}

\begin{ContractClause}{Cel strategiczny Spółki}
\label{art:cel}
\begin{enumerate}
  \item Głównym celem Spółki jest projektowanie, rozwój i komercjalizacja fizycznej (ucieleśnionej) sztucznej inteligencji oraz systemów autonomicznych, w szczególności oprogramowania autonomii, percepcji, koordynacji zespołowej, dowodzenia i obsługi misji, wraz z platformami sprzętowymi, na których systemy te działają.
  \item Pierwszym i podstawowym polem zastosowania technologii Spółki są bezzałogowe statki powietrzne, w szczególności autonomiczne roje dronów wraz z systemami ich dowodzenia, koordynacji, łączności, percepcji i obsługi misji. Spółka może rozwijać kolejne zastosowania fizycznej sztucznej inteligencji, w tym roboty lądowe, jednostki nawodne i inne maszyny autonomiczne oraz technologie podwójnego zastosowania. Rozwój kolejnych zastosowań mieszczących się w celu określonym w ust.~1 nie stanowi istotnej zmiany profilu działalności w rozumieniu §~\ref{art:walne-zastrzezone}.
  \item Rozwiązania Spółki są przeznaczone w szczególności do zastosowań cywilnych, przemysłowych, rolniczych, ratowniczych, infrastrukturalnych, bezpieczeństwa i obronności, przy zachowaniu zgodności z prawem oraz wymogami bezpieczeństwa państw Unii Europejskiej i NATO.
  \item Działalność wymagająca koncesji, zezwolenia, licencji, certyfikacji, wpisu do rejestru albo poświadczenia bezpieczeństwa może być rozpoczęta dopiero po spełnieniu odpowiednich wymogów. Dotyczy to w szczególności koncesji na wykonywanie działalności gospodarczej w zakresie wytwarzania i obrotu materiałami wybuchowymi, bronią, amunicją oraz wyrobami i technologią o przeznaczeniu wojskowym lub policyjnym, certyfikacji lotniczej bezzałogowych systemów powietrznych oraz zezwoleń wymaganych przepisami o kontroli obrotu towarami o znaczeniu strategicznym.

\begin{MemoBox}{4.1.1}
\MemoPytCd{4.1.1}
\Kierunek{Zwrot „certyfikacja lotnicza bezzałogowych systemów powietrznych” jest zawężający: w~kategorii otwartej i~szczególnej przepisów o~certyfikacji nie stosuje się (art.~156b ust.~3 Prawa lotniczego \Wiar{Z}), a~certyfikacja BSP i~operatora dotyczy kategorii certyfikowanej (art.~3 lit.~c, art.~6 ust.~1--2 rozporządzenia 2019/947 \Wiar{Z} (tekst EN)). Proponujemy zastąpić „certyfikacji lotniczej” zwrotem „oceny zgodności, certyfikacji i~zezwoleń operacyjnych”.}
\end{MemoBox}

\end{enumerate}

\end{ContractClause}


\begin{ContractClause}{Przedmiot działalności i PKD 2025}
\label{art:pkd}
\begin{enumerate}
\item \textbf{Plan działalności Spółki.} Podstawową działalnością Spółki jest projektowanie, badanie i rozwój fizycznej (ucieleśnionej) sztucznej inteligencji oraz systemów autonomicznych --- w szczególności oprogramowania autonomii, percepcji, koordynacji zespołowej, dowodzenia, sterowania, misji, łączności, nawigacji, telemetrii i zarządzania flotą --- oraz produkcja, integracja, wdrażanie, sprzedaż i serwis wykorzystujących te technologie cywilnych bezzałogowych statków powietrznych, w szczególności autonomicznych rojów dronów, wraz z infrastrukturą operatorską. Działalność ta obejmuje zarówno własne platformy sprzętowe, jak i oprogramowanie oraz integrację komponentów dostarczanych przez podmioty trzecie.
\item Spółka prowadzi badania naukowe i prace rozwojowe oraz działalność inżynieryjną i programistyczną dotyczącą robotyki, autonomii zespołowej, sztucznej inteligencji fizycznej i ucieleśnionej, widzenia komputerowego, fuzji sensorów, systemów nawigacyjnych i pomiarowych, działania w środowisku bez dostępu do globalnych systemów nawigacji satelitarnej, symulacji, cyfrowych bliźniaków, generowania danych syntetycznych, testowania i walidacji systemów autonomicznych.
\item Działalność uzupełniająca obejmuje projektowanie, produkcję, integrację, instalowanie, diagnostykę, naprawę, konserwację i modernizację mobilnych robotów lądowych, autonomicznych pojazdów i maszyn dla rolnictwa i leśnictwa, autonomicznych i zdalnie sterowanych jednostek nawodnych, stacji naziemnych, doków, systemów ładowania, komputerów pokładowych, autopilotów, sensorów, aparatury pomiarowej i nawigacyjnej oraz infrastruktury obliczeniowej i przetwarzania danych wykorzystywanej do obsługi takich systemów.
\item Spółka może rozwijać prototypy, demonstratory technologii, platformy badawcze i krótkie serie przedprodukcyjne oraz wdrażać rozwiązania dla zastosowań cywilnych, przemysłowych, infrastrukturalnych, ratowniczych, rolniczych, leśnych, morskich, bezpieczeństwa publicznego i ochrony infrastruktury krytycznej. W dalszym etapie Spółka może rozwijać także warianty przeznaczone do zastosowań obronnych lub wojskowych, wyłącznie po spełnieniu właściwych wymogów prawa.
\item \textbf{Kody PKD 2025 przewidziane do ujawnienia przy pierwszym wpisie do KRS:}
\begin{enumerate}
\item \textbf{72.10.Z — Badania naukowe i prace rozwojowe w dziedzinie nauk przyrodniczych i technicznych} — działalność przeważająca;
\item \textbf{30.31.Z — Produkcja cywilnych statków powietrznych, statków kosmicznych i podobnych maszyn};
\item \textbf{71.12.B — Pozostała działalność w zakresie inżynierii i związane z nią doradztwo techniczne};
\item \textbf{62.10.B — Pozostała działalność w zakresie programowania};
\item \textbf{28.99.Z — Produkcja pozostałych maszyn specjalnego przeznaczenia, gdzie indziej niesklasyfikowana};
\item \textbf{26.51.Z — Produkcja instrumentów i przyrządów pomiarowych, kontrolnych i nawigacyjnych};
\item \textbf{33.20.Z — Instalowanie maszyn przemysłowych, sprzętu i wyposażenia};
\item \textbf{33.16.Z — Naprawa i konserwacja cywilnych statków powietrznych i statków kosmicznych};
\item \textbf{30.11.Z — Produkcja cywilnych statków i konstrukcji pływających};
\item \textbf{28.30.Z — Produkcja maszyn dla rolnictwa i leśnictwa}.
\end{enumerate}
\item \textbf{Pozostałe kody PKD 2025 mieszczące się w planowanym przedmiocie działalności Spółki i przewidziane dla dalszego rozwoju:}
\begin{enumerate}[label=\alph*)]
\item \textbf{33.15.Z — Naprawa i konserwacja cywilnych statków i łodzi};
\item \textbf{33.12.Z — Naprawa i konserwacja maszyn};
\item \textbf{33.13.Z — Naprawa i konserwacja urządzeń elektronicznych i optycznych};
\item \textbf{29.10.E — Produkcja pozostałych pojazdów silnikowych, z wyłączeniem motocykli};
\item \textbf{63.10.D — Pozostała działalność usługowa w zakresie infrastruktury obliczeniowej, przetwarzania danych, zarządzania stronami internetowymi (hosting) i działalności powiązane};
\item \textbf{30.32.Z — Produkcja wojskowych statków powietrznych, statków kosmicznych i podobnych maszyn};
\item \textbf{30.13.Z — Produkcja wojskowych okrętów i jednostek pływających};
\item \textbf{30.40.Z — Produkcja wojskowych pojazdów bojowych};
\item \textbf{33.18.Z — Naprawa i konserwacja wojskowych pojazdów bojowych, okrętów, łodzi, statków powietrznych i statków kosmicznych}.
\end{enumerate}

\begin{MemoBox}{4.1.3}
\MemoPyt{4.1.3}{Czy zestaw kodów PKD 2025 wymaga końcowej kontroli klasyfikacyjnej przed wnioskiem do KRS i czy liczba kodów nie budzi zastrzeżeń sądu?}
\Ocena{zgodne --- kontrola wykonana na tekście rozporządzenia; wszystkie 19 kodów istnieje, nazwy w~umowie są identyczne z~urzędowymi}
\Klasa{bez zmian}
\Podstawa{rozporządzenie Rady Ministrów z~18 grudnia 2024~r. w~sprawie Polskiej Klasyfikacji Działalności (PKD) (Dz.U. 2024 poz.~1936), załącznik \Wiar{Z}; art.~40 pkt~1 ustawy o~KRS (limit dziesięciu pozycji, w~tym jedna przeważająca na poziomie podklasy) \Wiar{Z}; art.~165 i~art.~164 §~3 KSH w~zw. z~art.~300$^{13}$ §~1 KSH \Wiar{Z}.}
\Uzasadnienie{Limit dziesięciu pozycji dotyczy wyłącznie działu~3 rejestru; przedmiot działalności w~umowie nie jest limitowany \Wiar{Z}, więc 19 kodów w~umowie i~10 we wniosku nie budzi zastrzeżeń (praktyka sądów jednolita \Wiar{W}). Kontrolę wykonaliśmy na tekście załącznika do rozporządzenia: wszystkie 19 kodów z~§~\ref{art:pkd} ust.~5--6 istnieje w~PKD 2025, a~nazwy są identyczne z~urzędowymi co do znaku --- ryzyko wezwania z~powodu nazwy kodu nie występuje. Działalność przeważająca 72.10.Z (decyzja D7) jest spójna z~ust.~1 i~profilem B+R; 30.31.Z miałby sens tylko przy produkcji seryjnej od startu. Kod przeważający nie ma znaczenia koncesyjnego.}
\Rekomendacja{Nie zmieniać ust.~5--6; utrzymać 72.10.Z jako przeważający; we wniosku do KRS podać dokładnie dziesięć kodów z~ust.~5 z~nazwami w~brzmieniu umowy; jeżeli rozporządzenie zmieni się przed aktem, powtórzyć porównanie.}
\end{MemoBox}

\item Ujęcie określonego rodzaju działalności w niniejszym paragrafie nie oznacza obowiązku jego natychmiastowego rozpoczęcia. Spółka może rozwijać poszczególne obszary etapowo, odpowiednio do rozwoju produktów, technologii, kompetencji, finansowania i potrzeb rynku.

\begin{MemoBox}{4.1.2}
\MemoPyt{4.1.2}{Czy ujawnienie kodów wojskowych PKD przed uzyskaniem koncesji jest zasadne, czy lepiej dodać je przy zmianie umowy po koncesji?}
\Ocena{zgodne}
\Klasa{opcjonalne}
\Podstawa{art.~300$^{5}$ §~1 pkt~2 KSH (przedmiot działalności jako element umowy) \Wiar{Z}; art.~40 pkt~1 ustawy o~KRS („nie więcej niż dziesięć pozycji, w~tym jeden przedmiot przeważającej działalności na poziomie podklasy”) \Wiar{Z}; art.~5 i~art.~7 ust.~1 ustawy z~13 czerwca 2019~r. (koncesja warunkuje \emph{wykonywanie} działalności) \Wiar{Z}; art.~17 ust.~1 pkt~2 tej ustawy \Wiar{Z}; art.~300$^{98}$ §~2 pkt~1, art.~300$^{100}$ §~2 i~art.~300$^{102}$ §~1--2 KSH (zmiana umowy) \Wiar{Z}; art.~55 ust.~1 ustawy o~kosztach sądowych w~sprawach cywilnych (250~zł od wniosku o~zmianę wpisu) \Wiar{Z}.}
\Uzasadnienie{Projekt rozdziela przedmiot działalności w~umowie (§~\ref{art:pkd} ust.~5--6, wszystkie 19 kodów) i~pozycje ujawniane w~KRS (dziesięć z~ust.~5, bez kodów wojskowych). Prawo nie zakazuje ujęcia w~umowie działalności koncesjonowanej przed koncesją --- koncesja warunkuje jej wykonywanie, a~nie deklarację w~umowie; sąd rejestrowy nie żąda koncesji przy wpisie spółki \Wiar{W}. Zaleta: kod już obecny w~umowie ujawnia się w~KRS przy bramce G3 samym wnioskiem o~zmianę danych (250~zł \Wiar{Z} i~100~zł za ogłoszenie w~MSiG \Wiar{W}), bez protokołu notarialnego i~uchwały 75\%; bez kodu w~umowie potrzebna jest zmiana umowy (trzy czwarte głosów, protokół notarialny, wpis konstytutywny), co przy rundzie może wymagać zgody inwestora i~wydłużyć ścieżkę koncesyjną o~miesiące. Wada: umowa jest jawna w~aktach rejestrowych, więc banki i~kontrahenci w~KYC zobaczą działalność wojskową; część instytucji stosuje ograniczenia polityk sektorowych \Wiar{W}. Dla EDF, EUDIS i~DIANA kody są neutralne albo korzystne \Wiar{W}; dla EIC (Horyzont Europa finansuje wyłącznie zastosowania cywilne --- art.~7 ust.~1 rozporządzenia (UE) 2021/695 \Wiar{W}) co najwyżej neutralne. Otwarte: czy organ koncesyjny w~praktyce oczekuje zgodności zakresu koncesji z~kodami PKD w~KRS; jeżeli tak, ujawnienie 30.32.Z trzeba zaplanować \emph{przed} wnioskiem \Wiar{?}.}
\Rekomendacja{Pozostawić kody wojskowe w~§~\ref{art:pkd} ust.~6 i~nie ujawniać ich w~KRS przy pierwszym wpisie; ujawnić je wnioskiem o~zmianę danych przy bramce G3, w~kolejności ustalonej z~organem koncesyjnym. Opcjonalnie wzmocnić ust.~7 zdaniem, że wskazanie kodu nie stanowi podstawy do rozpoczęcia działalności regulowanej (wartość wyłącznie komunikacyjna wobec banków i~kontrahentów).}
\Kierunek{Bez skutków prawnych; zdanie powtarza treść ust.~8 innymi słowami, więc jego dodanie jest decyzją redakcyjną Założycieli. Proponowane brzmienie zdania trzeciego §~\ref{art:pkd} ust.~7:}
\begin{Brzmienie}
Wskazanie kodu PKD w niniejszym paragrafie ani jego ujawnienie w rejestrze przedsiębiorców nie stanowi podstawy do rozpoczęcia działalności wymagającej koncesji, zezwolenia, licencji, certyfikacji albo wpisu do rejestru działalności regulowanej.
\end{Brzmienie}
\end{MemoBox}

\item Działalność podlegająca koncesji, zezwoleniu, licencji, certyfikacji albo innym szczególnym wymogom może być rozpoczęta dopiero po spełnieniu wymagań wynikających z właściwych przepisów prawa, w szczególności ustawy z dnia 13 czerwca 2019 r. o wykonywaniu działalności gospodarczej w zakresie wytwarzania i obrotu materiałami wybuchowymi, bronią, amunicją oraz wyrobami i technologią o przeznaczeniu wojskowym lub policyjnym, przepisów o kontroli obrotu towarami o znaczeniu strategicznym oraz przepisów prawa lotniczego i regulacji Unii Europejskiej dotyczących bezzałogowych systemów powietrznych.

\begin{MemoBox}{4.1.1}
\MemoPyt{4.1.1}{Czy powołana podstawa (ustawa z 13 czerwca 2019 r. o wykonywaniu działalności gospodarczej w zakresie wytwarzania i obrotu materiałami wybuchowymi, bronią, amunicją oraz wyrobami i technologią o przeznaczeniu wojskowym lub policyjnym) jest aktualna i obejmuje bezzałogowe statki powietrzne oraz ich oprogramowanie; jak ma się do prawa lotniczego i rozporządzeń UE o bezzałogowych systemach powietrznych?}
\Ocena{zgodne warunkowo --- podstawa aktualna i~trafnie powołana; zmiany z~2026~r. sprawdzone w~tekście i~nieistotne dla Spółki; warunkiem jest kwalifikacja wyrobu do wykazu przed bramką G3}
\Klasa{zalecane}
\Podstawa{ustawa z~13 czerwca 2019~r. (Dz.U. 2019 poz.~1214; t.j. Dz.U. 2023 poz.~1743), zmieniona art.~2 ustawy z~13 marca 2026~r. (Dz.U. 2026 poz.~471) \Wiar{Z}; art.~3 ust.~1 pkt~12--13, art.~4, art.~5, art.~7 ust.~1, art.~8 ust.~1--2, art.~9 ust.~1, art.~10 ust.~1 pkt~2--3, art.~11--12, art.~17 i~art.~133 tej ustawy \Wiar{Z}; rozporządzenie Rady Ministrów z~17 września 2019~r. w~sprawie klasyfikacji (Dz.U. 2019 poz.~1888): kategoria WT~V ust.~3, WT~XI ust.~2, WT~XIII ust.~3--4 oraz definicje pkt~12--15 załącznika \Wiar{Z}; wykaz zawiera kategorie WT~I--XIV, nie ma kategorii WT~XXI/XXII \Wiar{Z}; rozporządzenie wykonawcze (UE) 2019/947 (tekst pierwotny): art.~2 pkt~1 i~17, art.~3--6, art.~12, art.~14 ust.~5--6, art.~23 ust.~1--2, załącznik część~A pkt UAS.OPEN.060 ust.~2 lit.~d i~część~B pkt UAS.SPEC.050 ust.~1 lit.~b \Wiar{Z} (tekst EN); rozporządzenie (UE) 2018/1139, art.~2 ust.~3 lit.~a \Wiar{W}; ustawa Prawo lotnicze (t.j. Dz.U. 2025 poz.~1431): art.~1 ust.~3--4, art.~2 pkt~1a, 1b, 2 i~24--26, dział~VIa, art.~156b ust.~3 \Wiar{Z}.}
\Uzasadnienie{\textbf{Aktualność.} Ustawa z~2019~r. obowiązuje i~pozostaje właściwą podstawą; powołanie jej samą nazwą jest prawidłowe. Nowela z~13 marca 2026~r. zmienia w~niej wyłącznie odnośnik do tytułu, art.~44 ust.~2aa oraz odesłania w~art.~44 ust.~4 i~art.~56 ust.~4; definicji, zakresu koncesji i~sankcji nie zmieniono. Zmiany w~ustawie o~obrocie strategicznym obowiązują od 22 kwietnia 2026~r. (ewidencja obrotu --- od 8 października 2026~r.).

\textbf{BSP i~oprogramowanie.} WT~V ust.~3 obejmuje BSP „oraz ich systemy i~urządzenia”, w~tym kierowania i~kontroli lotu, bez wyłączenia cywilnego; kryterium --- „specjalne zaprojektowanie lub zmodyfikowanie do celów wojskowych”. Wykaz nie ma odpowiednika ML21/ML22 (poprzednia wersja memorandum błędnie zakładała WT~XXI/XXII). Oprogramowanie autonomii BSP zaprojektowanego dla celów wojskowych jest składnikiem wyrobu, a~kod, dokumentacja i~dane niezbędne do jego rozwoju --- „technologią”, której udostępnienie (także licencja) jest obrotem wymagającym koncesji; produkt cywilny nie jest wyrobem z~wykazu. Brak wyłączenia dla prototypów i~B+R \Wiar{Z}; już wytworzenie wyrobu uruchamia obowiązek \Wiar{W}. Sankcja: od roku do lat 10 (art.~133).

\textbf{Prawo lotnicze} i~rozporządzenia 2019/945 i~2019/947 to odrębna oś (bezpieczeństwo cywilnych BSP). Operacja autonomiczna (bez możliwości interwencji pilota) mieści się w~kategorii szczególnej, nie otwartej (2019/947 \Wiar{Z} (tekst EN)); prawo lotnicze nie obejmuje lotnictwa państwowego (art.~1 ust.~3--4) \Wiar{Z}. Wymagania operacyjne trzeba sprawdzić na wersji skonsolidowanej 2019/947 \Wiar{W}.}
\Rekomendacja{Utrzymać obie klauzule. Doprecyzować §~\ref{art:pkd} ust.~8 (oprogramowanie i~technologia, w~tym licencjonowanie; przepisy „w~brzmieniu obowiązującym”); w~§~\ref{art:cel} ust.~4 zastąpić „certyfikacji lotniczej” zwrotem „oceny zgodności, certyfikacji i~zezwoleń operacyjnych”. Przed bramką G3: zakwalifikować konkretny wyrób do wykazu WT, wyznaczyć dwóch dyrektorów (albo dyrektora i~prokurenta) spełniających art.~10 ust.~1 pkt~1 (szkolenie z~art.~11, badania z~art.~12) oraz zaplanować pozyskanie zaświadczeń o~niekaralności od akcjonariuszy z~pakietem co najmniej 20\%. Sygnał dla pakietu A1: §~\ref{art:dopuszczalny-nabywca} ust.~15.}
\Kierunek{Zmiana redakcyjna bez skutków dla akcjonariuszy, inwestora, rejestru i~sądu; zwiększa czytelność dla banków, programów finansowania i~partnerów, że reguła dotyczy również oprogramowania i~technologii. Proponowane brzmienie §~\ref{art:pkd} ust.~8 (w~skrócie): działalność podlegająca koncesji, zezwoleniu, licencji, certyfikacji, ocenie zgodności albo innym szczególnym wymogom może być rozpoczęta dopiero po spełnieniu wymagań z~właściwych przepisów „w~brzmieniu obowiązującym w~dniu jej rozpoczęcia”, w~szczególności ustawy z~13 czerwca 2019~r. wraz z~przepisami wykonawczymi o~klasyfikacji wyrobów i~technologii, przepisów o~kontroli obrotu z~zagranicą towarami, technologiami i~usługami o~znaczeniu strategicznym (w~tym rozporządzenia (UE) 2021/821) oraz prawa lotniczego i~rozporządzeń (UE) 2019/945 i~2019/947; nowe zdanie obejmuje także wytwarzanie i~modyfikację wyrobów oraz udostępnianie, w~tym licencjonowanie, oprogramowania i~technologii objętych tymi przepisami.}
\end{MemoBox}

\end{enumerate}
\end{ContractClause}
\begin{ContractClause}{Akcje i kapitał akcyjny}
\label{art:akcje}
\begin{enumerate}
  \item Akcje Spółki nie mają wartości nominalnej, nie stanowią części kapitału akcyjnego i są niepodzielne.
  \item Na dzień zawiązania Spółki emituje się \textbf{1 000 000 (jeden milion)} zwykłych akcji imiennych serii A, oznaczonych numerami od A0000001 do A1000000.
  
\begin{MemoBox}{4.3.3}
\MemoPytCd{4.3.3}
\Kierunek{Opcjonalnie usunąć określenie „imiennych” (także w~§~\ref{art:seria-f}): w~P.S.A. jest zbędne, choć nie błędne --- „akcje nie mają formy dokumentu” (art.~300$^{29}$ §~1 KSH \Wiar{Z}), a~zniesienie podziału na akcje imienne i~na okaziciela w~nowelizacji z~23 stycznia 2026~r. dotyczy spółki akcyjnej. §~\ref{art:akcje} należy do pakietu A4.}
\end{MemoBox}

\item Każda akcja serii A daje prawo do jednego głosu na Walnym Zgromadzeniu i równego udziału w dywidendzie oraz majątku pozostałym po likwidacji, z zastrzeżeniem bezwzględnie obowiązujących przepisów prawa i postanowień niniejszej umowy.
\item Cena emisyjna każdej akcji serii A wynosi \textbf{0,01 zł (jeden grosz)}.
  \item Wysokość kapitału akcyjnego nie jest określana w niniejszej umowie i podlega ujawnieniu oraz zmianom zgodnie z przepisami prawa.
  \item \BoilerplateStrike{Akcje podlegają zarejestrowaniu w rejestrze akcjonariuszy prowadzonym przez podmiot uprawniony zgodnie z prawem.}

\begin{MemoBox}{4.3.2}
\MemoPyt{4.3.2}{Które ograniczenia i obowiązki powinny zostać ujawnione w rejestrze akcjonariuszy i jak dobrać podmiot prowadzący rejestr pod kątem obsługi tych ograniczeń?}
\Ocena{zgodne warunkowo}
\Klasa{opcjonalne --- warunek spełnia się w~umowie z~podmiotem prowadzącym rejestr}
\Podstawa{art.~300$^{31}$ §~1--5 KSH (dom maklerski i~inne podmioty z~art.~4 ust.~1 pkt~1 ustawy o~obrocie instrumentami finansowymi albo notariusz; wybór uchwałą akcjonariuszy, „przy zawiązaniu spółki wyboru dokonują akcjonariusze”) \Wiar{Z}; art.~300$^{32}$ §~1--2 KSH \Wiar{Z}; art.~300$^{33}$ §~1 pkt~10--11 i~§~2 KSH („umowa spółki może zawierać dodatkowe postanowienia dotyczące informacji ujawnianych w~rejestrze akcjonariuszy”) \Wiar{Z}; art.~300$^{34}$ §~1 i~3 KSH \Wiar{Z}; art.~300$^{36}$ §~4 i~art.~300$^{37}$ §~1 KSH \Wiar{Z}; ustawa z~23 stycznia 2026~r. (Dz.U. 2026 poz.~176): od 18 lutego 2027~r. art.~300$^{32}$ §~1$^{1}$--1$^{3}$ i~§~3, art.~300$^{33}$ §~3, art.~300$^{34}$ §~3 zdanie drugie i~§~9, art.~300$^{35}$ §~1$^{1}$ KSH \Wiar{Z}.}
\Uzasadnienie{\textbf{Co ujawniać.} Wpis nabywcy jest konstytutywny (art.~300$^{37}$ §~1), więc tylko rejestr pokazuje ograniczenia nabywcy i~podmiotowi. Ograniczenia rozporządzania (pkt~10): zgoda Spółki, Okres Ograniczenia Zbywania, prawo pierwszeństwa, przyłączenia i~przymusowego współzbycia, Kryterium i~weryfikacja nabywcy, ograniczenia dziedziczenia. Obowiązki związane z~akcją (pkt~11): obowiązek zbycia w~mechanizmie zwrotnego zbycia, zbycie po utracie Kryterium, obowiązki informacyjne, poufność i~--- po zmianie z~4.2.1 --- §~\ref{art:export} ust.~5. Podmioty wpisują z~reguły odesłanie do paragrafu umowy \Wiar{W}; akcje objęte mechanizmem wymagają adnotacji przy numerach, najlepiej jako odesłania do Załącznika nr~2 i~formuły z~§~\ref{art:zwrotne-zbycie}, bo od 18 lutego 2027~r. zmiany z~pkt~10--11 Rada zgłasza w~7 dni pod rygorem grzywny. Bez ujawnienia podmiot wpisze nabywcę na podstawie samej umowy zbycia (art.~300$^{36}$ §~4) --- skuteczność Kryterium wobec rejestru zależy od ujawnienia i~procedury podmiotu. \textbf{Dobór podmiotu} --- kryteria: (1)~wpisy tekstem wolnym przy numerach akcji, (2)~procedura przy zbyciu (zgoda Spółki, zgody z~art.~300$^{34}$ §~3 --- od 2027~r. kwalifikowane, blokady), (3)~adres e-mail, (4)~głosowanie elektroniczne, (5)~koszt i~(6)~czas wpisu, (7)~doświadczenie z~P.S.A. i~funduszami VC, (8)~gotowość do obowiązków z~nowelizacji. Notariusz --- taniej i~elastyczniej; dom maklerski --- narzędzia, lecz często sztywne szablony \Wiar{W}. Zmiana podmiotu wymaga nowej umowy (art.~300$^{32}$ §~2) i~od 2027~r. zgłoszenia do KRS --- to decyzja na lata.}
\Rekomendacja{Sporządzić wykaz ograniczeń i~obowiązków do ujawnienia jako załącznik do umowy z~podmiotem prowadzącym rejestr; wybrać podmiot po zapytaniu ofertowym do 2--3 podmiotów z~kryteriami (1)--(8); wybór ująć w~uchwale założycielskiej objętej aktem (4.3.1). Opcjonalnie przywrócić §~\ref{art:akcje} ust.~6 w~brzmieniu merytorycznym (§~\ref{art:akcje} należy do pakietu A4).}
\Kierunek{Dla Spółki: obowiązek aktualizacji adnotacji przy akcjach objętych mechanizmem (od 2027~r. w~terminie 7 dni); dla akcjonariuszy: pewność, że nabywca widzi ograniczenia; dla inwestora: standard oczekiwany w~due diligence; dla rejestru: jasna dyspozycja oparta na art.~300$^{33}$ §~2 KSH; dla sądu: bez znaczenia. Proponowane brzmienie §~\ref{art:akcje} ust.~6 (w~miejsce zdania przekreślonego):}
\begin{Brzmienie}
Akcje są rejestrowane w rejestrze akcjonariuszy prowadzonym przez podmiot wybrany uchwałą akcjonariuszy; pierwszego wyboru dokonują Założyciele przy zawiązaniu Spółki. Spółka zgłasza do rejestru akcjonariuszy ograniczenia rozporządzania akcjami i obowiązki związane z akcjami wynikające z niniejszej umowy, w tym oznaczenie akcji objętych mechanizmem zwrotnego zbycia, oraz aktualizuje te dane w terminach przewidzianych prawem.
\end{Brzmienie}
\end{MemoBox}

\end{enumerate}

\end{ContractClause}

\begin{ContractClause}{Objęcie akcji przez Założycieli}
\label{art:cap-table}
\begin{enumerate}
  \item Akcje serii A zostają objęte przez Założycieli w następujący sposób:
\end{enumerate}

\begin{center}
\begin{tabularx}{\textwidth}{>{\bfseries}p{2.5cm} >{\raggedright\arraybackslash}X r r}
\toprule
Założyciel & Dane identyfikacyjne & Liczba akcji & Udział w akcjach serii A \\
\midrule
A & \field{imię, nazwisko, PESEL/adres} & 500 000 & 50\% \\
B & \field{imię, nazwisko, PESEL/adres} & 200 000 & 20\% \\
C & \field{imię, nazwisko, PESEL/adres} & 100 000 & 10\% \\
D & \field{imię, nazwisko, PESEL/adres} & 100 000 & 10\% \\
E & \field{imię, nazwisko, PESEL/adres} & 100 000 & 10\% \\
\midrule
Łącznie & & 1 000 000 & 100\% \\
\bottomrule
\end{tabularx}
\end{center}

\begin{enumerate}[resume]
  \item Numery akcji przysługujących poszczególnym Założycielom oraz przyporządkowanie wkładów określa Załącznik nr 1.
  \item Na potrzeby niniejszej umowy „Założycielem Pierwotnym” jest każda osoba zawierająca niniejszą umowę i obejmująca przy zawiązaniu Spółki akcje serii A.
  \item Zachowanie relacji 50/20/10/10/10 dotyczy udziału Założycieli Pierwotnych między sobą na dzień zawiązania Spółki. Każda późniejsza emisja akcji, w tym emisja serii F, realizacja programu motywacyjnego serii P lub runda inwestycyjna, może proporcjonalnie rozwodnić ich udział, chyba że uchwała emisyjna stanowi inaczej.
  \item Wszystkie 1 000 000 akcji emitowanych przy zawiązaniu Spółki stanowią akcje serii A. Akcje dla późniejszych Założycieli Funkcjonalnych mogą być emitowane jako odrębna seria F, a akcje programu motywacyjnego jako odrębna seria P; emisje serii F i P nie zmieniają liczby akcji serii A i są od siebie niezależne.
\end{enumerate}

\end{ContractClause}


\begin{ContractClause}{Wkłady Założycieli}
\label{art:wklady}
\begin{enumerate}
  \item Akcje serii A są obejmowane w zamian za wkłady niepieniężne oraz wkłady pieniężne, w postaci:
  \begin{enumerate}[label=\alph*)]
    \item autorskich praw majątkowych do kodu źródłowego, kodu wynikowego, modeli, bibliotek, dokumentacji technicznej i innych elementów oprogramowania opisanych w Załączniku nr 1;
    \item prawa własności urządzeń, komputerów, komponentów elektronicznych, sensorów, aparatury laboratoryjnej, narzędzi i innego sprzętu opisanego w Załączniku nr 1;
    \item patentów, praw z patentów, praw ochronnych, praw do uzyskania patentu lub innych praw do wynalazków i zgłoszeń opisanych w Załączniku nr 1;
    \item wkładów pieniężnych w kwotach określonych w Załączniku nr 1.
  \end{enumerate}
  \item Praca, świadczenie usług, know-how nieutrwalone w zidentyfikowanej dokumentacji ani inne świadczenia osobiste nie stanowią wkładów Założycieli na pokrycie akcji serii A.
  \item Wartość przypisana każdemu wkładowi niepieniężnemu oraz kwota wkładu pieniężnego każdego Założyciela są określone w Załączniku nr 1. Łączna wartość wkładów przypisanych danemu Założycielowi odpowiada co najmniej łącznej cenie emisyjnej obejmowanych przez niego akcji.
  
\begin{MemoBox}{3.1.1}
\MemoPyt{3.1.1}{Jakie dokumenty i wyceny wkładów niepieniężnych są potrzebne do zawiązania i wpisu (art.~300$^{2}$ KSH), kto ponosi odpowiedzialność za zawyżenie wartości wkładu przeznaczonego na kapitał akcyjny i jak ją ograniczyć, skoro na kapitał idą tylko wkłady pieniężne?}
\Ocena{zgodne warunkowo}
\Klasa{zalecane}
\Podstawa{art.~300$^{2}$ §~1--2, art.~300$^{3}$ §~1--2, art.~300$^{5}$ §~1 pkt~3--5 i §~2 KSH \Wiar{Z}; art.~300$^{9}$ §~1--3, art.~300$^{10}$ §~1--2, art.~300$^{12}$ §~2 pkt~6--7, §~3 pkt~2--3 i §~4 KSH \Wiar{Z}; art.~300$^{123}$, art.~300$^{124}$, art.~587 §~1 KSH \Wiar{Z}; art.~14 §~1--2 KSH \Wiar{Z}; art.~17 ust.~1 pkt~9, art.~19 ust.~1 i~4, art.~22 ust.~1e pkt~3 ustawy o PIT \Wiar{Z}.}
\Uzasadnienie{P.S.A. nie zna sprawozdania założycieli ani badania wkładów przez biegłego rewidenta; wycena jest własną odpowiedzialnością Założycieli i dyrektorów \Wiar{W}. Umowa musi wskazywać przedmiot każdego wkładu niepieniężnego, serie i numery akcji za ten wkład i osobę obejmującą (art.~300$^{5}$ §~1 pkt~4); wartość nie jest jej ustawowym elementem, ale jest potrzebna do oświadczenia Rady o kapitale akcyjnym (art.~300$^{12}$ §~3 pkt~2 \Wiar{Z}), do PIT i do §~\ref{art:wklady} ust.~3 --- a Załącznik nr~1 nie ma kolumny wartości (3.1.4). Do zgłoszenia: umowa z Załącznikiem nr~1, oświadczenia dyrektorów o kapitale i wniesieniu wkładów, lista akcjonariuszy (art.~300$^{12}$ §~3--4 \Wiar{Z}); wniesienie po wpisie stwierdza uchwała Rady (art.~300$^{9}$ §~2 \Wiar{Z}). Za \emph{znaczne} zawyżenie wartości wkładu przeznaczonego na kapitał akcyjny wobec wartości godziwej z dnia objęcia odpowiada wnoszący, a dyrektorzy solidarnie, chyba że nie ponoszą winy, i nie można ich zwolnić (art.~300$^{10}$ §~1--2 \Wiar{Z}) --- surowiej niż art.~175 KSH; ponadto art.~14 §~2, 300$^{123}$ (cywilna, nie karna), 300$^{124}$ i 587 §~1 KSH \Wiar{Z}. Wyłączenie wkładów niepieniężnych z kapitału formalnie wyprowadza je poza art.~300$^{10}$, ale samo jest wątpliwe (3.1.2) i nie powinno służyć ograniczeniu odpowiedzialności. Skuteczne ograniczenie: rzetelna, udokumentowana, ostrożna wycena; zaniżeniu przeciwdziała art.~17 ust.~1 pkt~9 w zw. z art.~19 ust.~1 i~4 PIT \Wiar{Z}, a wysoka wycena praw wytworzonych przez Założyciela to realny podatek w dniu objęcia (koszt: faktyczne wydatki, art.~22 ust.~1e pkt~3 PIT \Wiar{Z}). Jedna obronna wartość dla umowy, ksiąg i PIT.}
\Rekomendacja{Przed dniem podpisania umowy przygotować dla każdego składnika: (a)~kod, modele, dokumentacja --- pisemną umowę przeniesienia autorskich praw majątkowych (3.1.3), wykaz repozytoriów z identyfikatorem wersji, sumy kontrolne, protokół wydania, wykaz twórców z podstawą nabycia praw, wykaz komponentów otwartych z licencjami; (b)~sprzęt --- protokół wydania z numerami seryjnymi, dowody nabycia, oświadczenie o braku obciążeń; (c)~zgłoszenia patentowe i prawa do uzyskania patentu --- cesję pisemną, wnioski o zmianę zgłaszającego do UPRP/EPO, wykaz twórców; (d)~wycenę --- raport rzeczoznawcy albo co najmniej notatkę metodyczną (metoda, założenia, data) podpisaną przez wszystkich Założycieli; (e)~oświadczenia z §~\ref{art:wlasnosc-intelektualna} ust.~5 w odrębnym dokumencie. W §~\ref{art:wklady} ust.~3 dodać wskazanie podstawy wyceny.}
\Kierunek{Zmiana porządkuje dowód wartości i tytułu na potrzeby wpisu, due diligence inwestora i ewentualnej odpowiedzialności z art.~300$^{10}$ KSH; nie zmienia praw akcjonariuszy ani nie wpływa na rejestr akcjonariuszy.}
\begin{Brzmienie}
§~\ref{art:wklady} ust.~3: „Wartość przypisana każdemu wkładowi niepieniężnemu oraz kwota wkładu pieniężnego każdego Założyciela są określone w Załączniku nr~1 wraz ze wskazaniem podstawy wyceny (raportu rzeczoznawcy albo opisanej metody i daty wyceny). Łączna wartość wkładów przypisanych danemu Założycielowi odpowiada co najmniej łącznej cenie emisyjnej obejmowanych przez niego akcji. Założyciele oświadczają, że wartość wkładów niepieniężnych została ustalona z należytą starannością i nie jest zawyżona.”
\end{Brzmienie}
\end{MemoBox}

\item Na kapitał akcyjny przeznacza się wkłady pieniężne w łącznej kwocie określonej w Załączniku nr 1, nie niższej niż 1,00 zł. Przed złożeniem wniosku o wpis Spółki do rejestru Założyciele wniosą w całości wkłady pieniężne. Pozostałe wkłady zostaną wniesione najpóźniej w terminie 14 dni od dnia wpisu Spółki do rejestru przedsiębiorców. Dokumenty rozporządzające prawami własności intelektualnej powinny zostać podpisane najpóźniej w dniu zawarcia niniejszej umowy, ze skutkiem dla Spółki w organizacji albo — jeżeli wymaga tego charakter prawa — ze skutkiem na dzień wpisu Spółki do rejestru.
  
\begin{MemoBox}{3.1.2}
\MemoPyt{3.1.2}{Czy przeznaczenie na kapitał akcyjny wyłącznie wkładów pieniężnych (a wkładów niepieniężnych poza kapitał) jest dopuszczalne i jakie ma skutki podatkowe i bilansowe?}
\Ocena{ryzyko}
\Klasa{konieczne (K12)}
\Podstawa{art.~300$^{3}$ §~1--2 w zw. z art.~14 §~1 KSH \Wiar{Z}; art.~300$^{2}$ §~2, art.~300$^{10}$ §~1, art.~300$^{12}$ §~3 pkt~2, art.~300$^{15}$ §~1--6, art.~300$^{19}$ KSH \Wiar{Z}; art.~17 ust.~1 pkt~9, art.~19 ust.~1 i~3--4, art.~21 ust.~1 pkt~109, art.~22 ust.~1e i~1f ustawy o PIT \Wiar{Z}; art.~12 ust.~1 pkt~2 ustawy o CIT \Wiar{W}; art.~1 ust.~1 pkt~1 lit.~k i ust.~3 pkt~2, art.~1a pkt~2, art.~6 ust.~1 pkt~8 i art.~7 ust.~1 pkt~9 ustawy o PCC \Wiar{Z}; art.~36 ust.~1 i~2aa ustawy o rachunkowości \Wiar{Z}.}
\Uzasadnienie{Art.~300$^{3}$ §~1 KSH: na kapitał akcyjny „przeznacza się wniesione wkłady pieniężne oraz niepieniężne, z uwzględnieniem art.~14 §~1” \Wiar{Z}; art.~14 §~1 wyłącza tylko prawa niezbywalne oraz pracę i usługi \Wiar{Z}. Według przeważającego, naszym zdaniem, odczytania Założyciele nie wybierają, które wkłady zdatne do zaliczenia trafią do kapitału --- swoboda dotyczy wartości, nie alokacji \Wiar{W}; pogląd przeciwny (zwrot „przeznaczonych na kapitał akcyjny” w art.~14 §~1, art.~300$^{10}$ §~1, art.~300$^{12}$ §~3 pkt~2 \Wiar{Z}) jest mniejszościowy i niepotwierdzony w praktyce rejestrowej \Wiar{?}. Akcje nie mają wartości nominalnej (art.~300$^{2}$ §~3 \Wiar{Z}), P.S.A. nie zna agio --- cała wartość wkładu zaliczalnego tworzy kapitał. §~\ref{art:wklady} ust.~4 grozi zakwestionowaniem oświadczenia o kapitale przez sąd rejestrowy i nie ma odpowiednika bilansowego (art.~36 ust.~2aa ustawy o rachunkowości: kapitał akcyjny w wysokości wpisanej w rejestrze \Wiar{Z}); wkład poza kapitałem może być nieodpłatnym przysporzeniem --- przychodem Spółki (art.~12 ust.~1 pkt~2 CIT \Wiar{W}). PIT Założycieli nie zależy od alokacji: przychodem jest wartość wkładu z umowy, a gdy niższa od rynkowej --- rynkowa (art.~17 ust.~1 pkt~9, art.~19 ust.~3--4 PIT \Wiar{Z}); zwolnienie z art.~21 ust.~1 pkt~109 dotyczy tylko (zorganizowanej części) przedsiębiorstwa \Wiar{Z}. Korzyść w PCC jest iluzoryczna: ustawa nie wymienia P.S.A. ani kapitału akcyjnego (art.~1a pkt~2, art.~6 ust.~1 pkt~8 lit.~a \Wiar{Z}), opodatkowanie jest co najmniej sporne, a według dostępnych interpretacji organy przyjmują jego brak \Wiar{W}. Wyższy kapitał nie „zamraża” środków (wypłaty w granicach testów art.~300$^{15}$ §~1--6 \Wiar{Z}).}
\Rekomendacja{Przeznaczyć na kapitał akcyjny wszystkie wkłady z §~\ref{art:wklady} ust.~1 lit.~a--d w wartościach z Załącznika nr~1; kapitał na dzień wpisu obejmuje wkłady wniesione przed zgłoszeniem, a wkłady wnoszone w 14 dni po wpisie zwiększają kapitał na podstawie uchwały Rady bez zmiany umowy. Jeżeli Założyciele chcą utrzymać niski kapitał, drogą jest ostrożna (ale obronna podatkowo) wycena wkładów, nie ich wyłączenie z kapitału. Decyzję poprzedzić potwierdzeniem przez kancelarię wykładni art.~300$^{3}$ §~1 KSH.}
\begin{Brzmienie}
§~\ref{art:wklady} ust.~4: „Na kapitał akcyjny przeznacza się wkłady pieniężne oraz wkłady niepieniężne, o których mowa w ust.~1 lit.~a--c, w wartościach określonych w Załączniku nr~1, z uwzględnieniem art.~14 §~1 KSH; kapitał akcyjny na dzień zgłoszenia Spółki do rejestru wynosi nie mniej niż 1,00~zł i odpowiada wartości wkładów wniesionych do tego dnia. Przed złożeniem wniosku o wpis Spółki do rejestru Założyciele wniosą w całości wkłady pieniężne. Pozostałe wkłady zostaną wniesione najpóźniej w terminie 14 dni od dnia wpisu Spółki do rejestru przedsiębiorców; ich wniesienie stwierdza Rada Dyrektorów uchwałą i ujawnia wynikającą z tego zmianę wysokości kapitału akcyjnego zgodnie z art.~300$^{3}$ §~2 KSH. Dokumenty rozporządzające prawami własności intelektualnej powinny zostać podpisane najpóźniej w dniu zawarcia niniejszej umowy, ze skutkiem dla Spółki w organizacji albo --- jeżeli wymaga tego charakter prawa --- ze skutkiem na dzień wpisu Spółki do rejestru.”
\end{Brzmienie}
\end{MemoBox}

\item Jeżeli skuteczne przeniesienie danego prawa wymaga odrębnej formy, zgody, wpisu, zawiadomienia lub protokołu wydania, właściwy Założyciel jest zobowiązany niezwłocznie dokonać wszystkich wymaganych czynności na własny koszt.
  \item Do czasu skutecznego wniesienia wkładu Założyciel nie może rozporządzać jego przedmiotem, obciążać go ani udzielać praw osobom trzecim, poza niewyłączną licencją konieczną do bieżącej działalności Spółki, udzieloną Spółce nieodpłatnie i bez ograniczeń terytorialnych.
  \item Niezależnie od wkładów Założyciel może udzielić Spółce pożyczki na warunkach rynkowych; do takiej pożyczki stosuje się §~\ref{art:finansowanie} oraz §~\ref{art:transakcje-akcjonariusze}.
\end{enumerate}

\end{ContractClause}

\begin{ContractClause}{Program motywacyjny — pula 15\% po pełnym rozwodnieniu}
\label{art:program-motywacyjny}
\begin{enumerate}
  \item Tworzy się docelową pulę programu motywacyjnego dla kluczowych pracowników, współpracowników, dyrektorów i doradców Spółki, odpowiadającą maksymalnie \textbf{176 471 (stu siedemdziesięciu sześciu tysiącom czterystu siedemdziesięciu jednej)} akcji serii P.
  \item Przy założeniu wyemitowania wszystkich 176 471 akcji serii P i braku innych emisji, łączna liczba akcji po pełnym wykorzystaniu puli wyniesie 1 176 471, a pula programu motywacyjnego będzie odpowiadała 15\% po pełnym rozwodnieniu po zaokrągleniu do jednej dziesiątej punktu procentowego. Różnica od matematycznie dokładnych 15\% wynika wyłącznie z niepodzielności akcji.
  \item Akcje serii P mogą być emitowane jednorazowo albo wielokrotnie, nie później niż do dnia \textbf{31 grudnia 2036 r.}, przy czym każda emisja wymaga uchwały Walnego Zgromadzenia podjętej większością 75\% wszystkich głosów w Spółce i określającej elementy wymagane przez bezwzględnie obowiązujące przepisy prawa. Walne Zgromadzenie może taką większością podjąć jedną uchwałę ramową obejmującą kilka transz emisji akcji serii P, określającą dla każdej transzy maksymalną liczbę akcji, minimalną cenę emisyjną, okres, w którym transza może zostać wykonana, oraz pozostałe elementy wymagane przez bezwzględnie obowiązujące przepisy prawa; w granicach uchwały ramowej wykonanie poszczególnych transz, w tym wybór uczestników i zawarcie umów objęcia akcji, powierza się Radzie Dyrektorów zgodnie z ust.~6 i §~\ref{art:rada-zastrzezone} ust.~1 lit.~j, bez potrzeby odrębnej uchwały Walnego Zgromadzenia dla każdej transzy. Uchwała ramowa stanowi uchwałę o emisji w rozumieniu §~\ref{art:walne-zastrzezone} ust.~1 lit.~b dla wszystkich objętych nią transz.
  
\begin{MemoBox}{3.2.1}
\MemoPyt{3.2.1}{Czy „uchwała ramowa” na wiele transz emisji serii P, wykonywana przez Radę, oraz jedna uchwała 4/5 o pozbawieniu prawa poboru dla wszystkich transz są dopuszczalne na tle art.~300$^{103}$--300$^{107}$ KSH, czy każda transza wymaga własnej uchwały o prawie poboru?}
\Ocena{ryzyko (konstrukcja hybrydowa nieznana ustawie; ustawa daje dwa gotowe instrumenty, z których umowa nie korzysta)}
\Klasa{konieczne (K14)}
\Podstawa{art.~300$^{103}$, art.~300$^{104}$ §~1 pkt~1--7 i §~2, art.~300$^{105}$ §~1--3, art.~300$^{106}$ §~1--2 i~6, art.~300$^{107}$ §~1--3 KSH \Wiar{Z}; art.~300$^{108}$--300$^{113}$ KSH (upoważnienie zarządu albo rady dyrektorów do emisji akcji) \Wiar{Z}; art.~300$^{114}$--300$^{118}$ KSH (warunkowa emisja akcji), art.~300$^{119}$ KSH (warranty subskrypcyjne), art.~300$^{81}$ pkt~4 KSH \Wiar{Z}; art.~300$^{102}$ §~2 KSH (termin sześciu miesięcy na zgłoszenie zmiany umowy) \Wiar{Z}; art.~300$^{37}$ §~2 KSH w brzmieniu od 18 lutego 2027~r. (Dz.U. 2026 poz.~176, art.~1 pkt~6) \Wiar{Z}.}
\Uzasadnienie{\emph{Podstawa w umowie:} art.~300$^{103}$ KSH zwalnia emisję z przepisów o zmianie umowy, jeżeli umowa przewiduje „maksymalną liczbę akcji i termin ich emisji” \Wiar{Z}; §~\ref{art:program-motywacyjny} ust.~1 i~3 to spełniają. \emph{Uchwała o emisji:} wbrew założeniu ust.~3 P.S.A. zna delegację kompetencji emisyjnej: (a)~upoważnienie Rady Dyrektorów (art.~300$^{110}$--300$^{113}$) --- do pięciu lat, łącznie do jednej czwartej akcji istniejących, za wkłady pieniężne; uchwała Rady zastępuje uchwałę WZ, a pozbawienie prawa poboru wymaga uchwały akcjonariuszy 4/5 albo upoważnienia w umowie uchwalonego 4/5 \Wiar{Z}; (b)~warunkowa emisja akcji (art.~300$^{114}$--300$^{118}$) --- ustawowy instrument programów motywacyjnych: akcje obejmują osoby uprawnione z umowy ze spółką zawartej za zgodą WZ 3/4 (art.~300$^{114}$ §~2 pkt~2 i §~3), do dwukrotności akcji istniejących; uchwała sama „skutkuje wyłączeniem prawa poboru” i wymaga 4/5 (art.~300$^{115}$ §~2); uprawnieni obejmują akcje pisemnym oświadczeniem, prawa z akcji powstają z wpisem do rejestru akcjonariuszy na dyspozycję Rady, a do sądu trafia tylko roczny wykaz (art.~300$^{117}$--300$^{118}$) \Wiar{Z}; potwierdza to art.~300$^{37}$ §~2 w brzmieniu od 18 lutego 2027~r. \Wiar{Z}. Uchwała ramowa z ust.~3 nie jest żadnym z nich, lecz zwykłą emisją rozłożoną na transze. Jest dopuszczalna tylko wtedy, gdy sama jest uchwałą o emisji dla wszystkich akcji: liczba, seria i numery, cena i terminy wkładów (albo upoważnienie Rady do ich oznaczenia), dzień dywidendy (art.~300$^{104}$ §~1 pkt~1, 3--5 \Wiar{Z}); wybór uczestników, oferta i umowa objęcia to czynności wykonawcze Rady (art.~300$^{105}$), a skoro przyjęcie oferty nie może być uzależnione od warunku lub terminu (§~2), vesting zostaje w umowie uczestnictwa. Słabości: każda transza wymaga odrębnego zgłoszenia (art.~300$^{107}$ §~2), akcje powstają z wpisem do KRS (§~3), a sześciomiesięczny termin z art.~300$^{102}$ §~2 prawdopodobnie nie obowiązuje, lecz kwestia jest nierozstrzygnięta \Wiar{?}; sąd może traktować każdą transzę jako odrębną emisję wymagającą własnej uchwały. \emph{Prawo poboru:} art.~300$^{106}$ §~1 jest dyspozytywny („jeżeli umowa spółki lub uchwała akcjonariuszy nie stanowią inaczej”) \Wiar{Z}, więc umowa może wyłączyć pobór serii P z góry, przy zawiązaniu jednomyślnie; uchwała 4/5 (§~2) jest potrzebna tylko, gdy umowa milczy. Jedna uchwała 4/5 dla wielu transz jest spójna z jedną uchwałą o emisji, lecz nie wiąże późniejszego inwestora. Podstawowa ścieżka umowy opiera się na nieprawdziwej tezie o braku delegacji i pomija instrument stworzony dla programów motywacyjnych --- wada do poprawy przed pierwszą emisją.}
\Rekomendacja{Przebudować §~\ref{art:program-motywacyjny} ust.~3 i~5: (1)~jako tryb podstawowy przewidzieć warunkową emisję akcji serii P (jedna uchwała 4/5 na całą pulę albo jej część, zgoda na zawieranie umów uczestnictwa w tej samej uchwale, objęcie po nabyciu uprawnień pisemnym oświadczeniem, prawa z akcji z chwilą wpisu do rejestru akcjonariuszy); (2)~zachować tryb zwykły z uchwałą obejmującą kilka transz jako alternatywę, z pełnym katalogiem art.~300$^{104}$ i upoważnieniami dla Rady; (3)~wyłączyć w umowie prawo poboru akcji serii P na podstawie art.~300$^{106}$ §~1, co czyni ust.~5 zbędnym (brzmienie ust.~5 --- ramka przy ust.~5); (4)~potwierdzić z kancelarią, czy termin z art.~300$^{102}$ §~2 dotyczy emisji z art.~300$^{103}$, i do czasu potwierdzenia zgłaszać każdą emisję zwykłą w sześć miesięcy od uchwały. Upoważnienie Rady z art.~300$^{110}$ rozważyć dla drobnych emisji, nie dla puli P. Wybór trybu podstawowego jest decyzją Założycieli; rekomendujemy warunkową emisję.}
\begin{Brzmienie}
§~\ref{art:program-motywacyjny} ust.~3: „Akcje serii P mogą być emitowane jednorazowo albo wielokrotnie, na podstawie uchwał Walnego Zgromadzenia podjętych nie później niż do dnia \textbf{31 grudnia 2036~r.} większością 75\% wszystkich głosów w Spółce. Emisja akcji serii P następuje w trybie warunkowej emisji akcji (art.~300$^{114}$--300$^{118}$ KSH) albo w trybie zwykłym (art.~300$^{103}$--300$^{107}$ KSH). Uchwała o warunkowej emisji akcji serii P określa maksymalną liczbę akcji, cenę emisyjną albo sposób jej ustalenia, cel emisji, termin wykonania prawa do objęcia akcji oraz krąg osób uprawnionych; zgoda Walnego Zgromadzenia na zawarcie umów przyznających prawa do objęcia akcji (art.~300$^{114}$ §~2 pkt~2 i §~3 KSH) może zostać udzielona w tej samej uchwale dla kategorii osób i na warunkach określonych w regulaminie programu motywacyjnego. W granicach takiej uchwały Rada Dyrektorów zawiera umowy uczestnictwa, przyjmuje oświadczenia o objęciu akcji po spełnieniu warunków nabycia uprawnień, wydaje dyspozycje wpisu akcji do rejestru akcjonariuszy i składa zgłoszenia, o których mowa w art.~300$^{118}$ §~2 KSH. Uchwała o emisji w trybie zwykłym określa elementy wymagane przez art.~300$^{104}$ KSH, może obejmować kilka transz akcji serii P z podziałem numerów akcji między transze i okresami ich wykonania oraz może upoważniać Radę Dyrektorów do oznaczenia ceny emisyjnej nie niższej niż minimalna i do wskazania terminów wniesienia wkładów w umowach objęcia akcji; w jej granicach Rada Dyrektorów wybiera uczestników, składa oferty objęcia, zawiera umowy objęcia akcji i zgłasza emisję do rejestru, bez odrębnej uchwały Walnego Zgromadzenia dla każdej transzy. Każda z powyższych uchwał stanowi uchwałę o emisji w rozumieniu §~\ref{art:walne-zastrzezone} ust.~1 lit.~b dla wszystkich objętych nią akcji.”
\end{Brzmienie}
\end{MemoBox}

\item Cena emisyjna akcji serii P nie może być niższa niż 0,01 zł za akcję, chyba że bezwzględnie obowiązujące przepisy prawa wymagają ceny wyższej.
  \item Pozbawienie dotychczasowych akcjonariuszy prawa poboru akcji serii P w całości albo w części wymaga odrębnej uchwały akcjonariuszy podjętej większością co najmniej czterech piątych głosów, z zachowaniem pozostałych wymogów prawa. Uchwała ta może obejmować wszystkie transze objęte uchwałą ramową, o której mowa w ust.~3, w granicach dopuszczalnych przez bezwzględnie obowiązujące przepisy prawa.
  
\begin{MemoBox}{3.2.1}
\MemoPytCd{3.2.1}
Wyłączenie prawa poboru akcji serii P w umowie na podstawie art.~300$^{106}$ §~1 KSH (uzasadnienie --- ramka 3.2.1 przy ust.~3):
\begin{Brzmienie}
§~\ref{art:program-motywacyjny} ust.~5: „Na podstawie art.~300$^{106}$ §~1 KSH dotychczasowym akcjonariuszom nie przysługuje prawo poboru akcji serii P. Jeżeli bezwzględnie obowiązujące przepisy prawa wymagają mimo to uchwały o pozbawieniu prawa poboru, uchwała ta wymaga większości co najmniej czterech piątych głosów i może obejmować wszystkie akcje serii P objęte daną uchwałą o emisji.”
\end{Brzmienie}
\end{MemoBox}

\item Rada Dyrektorów prowadzi program motywacyjny, rekomenduje uczestników i warunki przyznania oraz wykonuje uchwały Walnego Zgromadzenia, lecz nie jest samodzielnie uprawniona do emisji akcji serii P ani do pozbawienia prawa poboru.
  \item Przyznanie praw w programie motywacyjnym nie daje uczestnikowi statusu akcjonariusza do chwili emisji akcji i wpisu uczestnika do rejestru akcjonariuszy.
  \item Akcje serii P mogą podlegać ograniczeniom zbywania, stopniowemu zwalnianiu z obowiązku zbycia, wymogom dotyczącym Dopuszczalnego Nabywcy oraz zasadom nabycia przez wskazanego nabywcę określonym w regulaminie programu i umowie uczestnictwa.
  \item Zwiększenie puli ponad 176 471 akcji wymaga zmiany niniejszej umowy albo odrębnej podstawy emisyjnej przyjętej większością określoną dla Spraw Zastrzeżonych.

\begin{MemoBox}{3.2.2}
\MemoPytCd{3.2.2}
\Kierunek{Ramka główna 3.2.2 --- przy §~\ref{art:seria-f} ust.~4. Uprościć ust.~9: każda emisja poza upoważnieniem z art.~300$^{103}$ KSH jest zmianą umowy (protokół notarialny), więc „odrębna podstawa emisyjna” nie jest odrębną, prostszą ścieżką i może sugerować, że wystarczy zwykła uchwała.}
\begin{Brzmienie}
§~\ref{art:program-motywacyjny} ust.~9: „Zwiększenie puli ponad 176\,471 akcji albo przedłużenie terminu, o którym mowa w ust.~3, wymaga zmiany niniejszej umowy.”
\end{Brzmienie}
\end{MemoBox}

\end{enumerate}

\begin{MemoBox}{3.2.4}
\MemoPyt{3.2.4}{Aspekty podatkowe programu motywacyjnego w P.S.A. (zakres E) --- zakres i koszt interpretacji indywidualnej.}
\Ocena{zgodne (umowa nie przesądza modelu podatkowego; ryzyko leży w prawie podatkowym, nie w tekście umowy)}
\Klasa{bez zmian}
\Podstawa{art.~24 ust.~11--11b i~12a--12b ustawy o PIT (ust.~12b stosuje ust.~11--11b odpowiednio do akcji prostej spółki akcyjnej) \Wiar{Z}; art.~14d §~1, art.~14f §~1--2, art.~14k--14m, art.~14n §~1 pkt~1 Ordynacji podatkowej \Wiar{Z}; art.~300$^{114}$--300$^{118}$ KSH \Wiar{Z}.}
\Uzasadnienie{Preferencja z art.~24 ust.~11 PIT (odroczenie przychodu do odpłatnego zbycia akcji, 19\% z kapitałów pieniężnych, bez składek) wymaga programu utworzonego uchwałą walnego zgromadzenia i stosuje się odpowiednio do akcji P.S.A. (ust.~12b \Wiar{Z}); §~\ref{art:program-motywacyjny} spełnia warunek uchwały, a warunkowa emisja akcji (3.2.1) odpowiada wprost konstrukcji ust.~11b --- otwarte pozostaje, jak organy potraktują objęcie po 0,01~zł i „odpowiednie” stosowanie przy braku wartości nominalnej. Wniosek o interpretację powinien objąć: kwalifikację programu serii P, brak przychodu przy umowie uczestnictwa i objęciu akcji (pracownik, zleceniobiorca, członek Rady), uczestników B2B (poza preferencją), obowiązki płatnika, koszt przy zbyciu, serię F, PCC i składki ZUS; opłata to 40~zł od każdego stanu faktycznego albo zdarzenia przyszłego (art.~14f §~1--2 Ordynacji \Wiar{Z}), termin wydania do trzech miesięcy (art.~14d §~1 \Wiar{Z}), honorarium doradcy rynkowo rzędu kilku do kilkunastu tysięcy złotych \Wiar{W}. Interpretacja chroni tylko wnioskodawcę (art.~14k--14m), więc wnioski powinny złożyć Spółka (przed zawiązaniem --- Założyciele, art.~14n §~1 pkt~1 \Wiar{Z}) i po jednym reprezentatywnym uczestniku z każdej kategorii (pracownik, B2B), przed uchwałą ramową, bo wybór modelu programu zmienia stan faktyczny wniosku.}
\Rekomendacja{Bez zmian w umowie. Zlecić doradcy podatkowemu wnioski w powyższym zakresie po wyborze modelu programu i przed pierwszą umową uczestnictwa; wynik uwzględnić w regulaminie programu.}
\end{MemoBox}

\end{ContractClause}

\begin{ContractClause}{Założyciele Funkcjonalni Serii F — późniejsi współzałożyciele}
\label{art:seria-f}
\begin{enumerate}
  \item „Założycielem Funkcjonalnym Serii F” może zostać wyłącznie osoba fizyczna, która po zawiązaniu Spółki obejmuje trwałą, strategiczną i osobistą odpowiedzialność za istotny obszar rozwoju Spółki, porównywalną co do znaczenia z rolą Założyciela Pierwotnego, oraz zobowiązuje się do długoterminowego, istotnego zaangażowania operacyjnego w rozwój Spółki. Sam status pracownika, współpracownika, dyrektora albo doradcy nie jest wystarczający do uzyskania tego statusu. Posiadanie akcji serii A ani status Założyciela Pierwotnego nie wyłączają możliwości uzyskania tego statusu ani objęcia akcji serii F.
  \item Nadanie statusu Założyciela Funkcjonalnego Serii F wymaga uchwały Walnego Zgromadzenia podjętej większością 75\% Głosów Uprawnionych w Sprawie. Kandydat będący akcjonariuszem jest wyłączony od głosowania w tej sprawie. Uchwała określa co najmniej rolę strategiczną, minimalny zakres zaangażowania, maksymalną liczbę akcji serii F przewidzianych dla danej osoby oraz „Dzień Przyznania Serii F”, od którego liczy się okres stopniowego zwalniania akcji z obowiązku zbycia.
  \item Uchwała o nadaniu statusu Założyciela Funkcjonalnego Serii F ma charakter kwalifikacyjny i nie stanowi sama przez się oferty objęcia akcji ani bezwarunkowego prawa do ich objęcia. Przed złożeniem przez Spółkę wiążącej oferty objęcia akcji Walne Zgromadzenie może większością 75\% wszystkich głosów cofnąć nadany status, jeżeli dana osoba przestała spełniać kryteria określone w ust. 1 lub warunki uchwały kwalifikacyjnej. Po zawarciu umowy objęcia akcji oraz po powstaniu akcji ich los podlega wyłącznie właściwym przepisom prawa, niniejszej umowie i mechanizmowi zwrotnego zbycia; nadanego prawa własności akcji nie cofa sama uchwała o utracie statusu.
  
\begin{MemoBox}{3.2.3}
\MemoPyt{3.2.3}{Czy uchwała kwalifikacyjna serii F (bez oferty i bezwarunkowego prawa do akcji, §~\ref{art:seria-f} ust.~3) nie rodzi po stronie kandydata roszczeń i czy status Założyciela Funkcjonalnego Serii F może być odebrany?}
\Ocena{zgodne}
\Klasa{opcjonalne}
\Podstawa{art.~66 §~1, art.~72 §~2, art.~389 §~1, art.~353$^{1}$ KC \Wiar{Z}; art.~300$^{105}$ §~1--3, art.~300$^{107}$ §~3, art.~300$^{118}$ §~1 KSH \Wiar{Z}; art.~300$^{101}$ KSH w zw. z art.~422 §~2 KSH (legitymacja do zaskarżenia uchwały) \Wiar{Z}.}
\Uzasadnienie{Uchwała kwalifikacyjna jest aktem wewnętrznym Spółki, nie oświadczeniem woli wobec kandydata; ust.~3 odmawia jej charakteru oferty, a akcje powstają z wpisem do rejestru (art.~300$^{107}$ §~3; przy warunkowej emisji prawa z akcji --- z wpisem do rejestru akcjonariuszy, art.~300$^{118}$ §~1 \Wiar{Z}). Nie jest ofertą (nie określa istotnych postanowień umowy objęcia, art.~66 §~1 KC) ani umową przedwstępną (kandydat nie jest jej stroną, art.~389 §~1 KC \Wiar{Z}). Roszczenia kandydata mogą powstać tylko poza uchwałą: (1)~z umowy o pracę, umowy B2B albo term-sheetu założycielskiego, jeżeli przyrzekają akcje; (2)~z winy w kontraktowaniu (art.~72 §~2 KC \Wiar{Z}) --- np. gdy Spółka prowadziła rozmowy o objęciu akcji bez zamiaru ich zakończenia albo cofnęła status arbitralnie po tym, jak kandydat na jej życzenie porzucił inne zatrudnienie --- w granicach ujemnego interesu umownego; (3)~z art.~5 KC przy nadużyciu. Cofnięcie statusu jest możliwe przed wiążącą ofertą, większością 75\% wszystkich głosów, tylko z zamkniętych przesłanek ust.~3; cofnięcie bez przesłanki narusza umowę spółki (między akcjonariuszami). Kandydat niebędący akcjonariuszem nie może zaskarżyć uchwały (art.~422 §~2 KSH w zw. z art.~300$^{101}$ \Wiar{Z}); kandydat-akcjonariusz może ją zaskarżyć jako krzywdzącą, choć od głosowania jest wyłączony. Po zawarciu umowy objęcia cofnięcie statusu nie ma skutku dla akcji. Konstrukcja jest spójna; źródłem sporu, nie wadą prawną, jest to, że okres liczy się od Dnia Przyznania (ust.~8), a kandydat pracuje przez miesiące bez akcji.}
\Rekomendacja{Wymóg jasnego komunikatu przenieść do dokumentów z kandydatem: umowa o pracę albo B2B i term-sheet założycielski powinny stwierdzać, że do zawarcia umowy objęcia kandydatowi nie przysługuje roszczenie o akcje ani o złożenie oferty, że status może zostać cofnięty na warunkach §~\ref{art:seria-f} ust.~3 oraz że wynagrodzenie jest niezależne od akcji. W umowie spółki dodać zdanie o braku roszczenia o złożenie oferty.}
\Kierunek{Zmiana wyjaśniająca; zmniejsza pole do argumentu z art.~72 §~2 KC, nie zmienia większości ani praw akcjonariuszy.}
\begin{Brzmienie}
§~\ref{art:seria-f} ust.~3 zdanie pierwsze: „Uchwała o nadaniu statusu Założyciela Funkcjonalnego Serii F ma charakter kwalifikacyjny i nie stanowi sama przez się oferty objęcia akcji, przyrzeczenia jej złożenia ani bezwarunkowego prawa do objęcia akcji; roszczenie o objęcie akcji serii F powstaje wyłącznie z umowy objęcia akcji.”
\end{Brzmienie}
\end{MemoBox}

\item Na podstawie art. 300$^{103}$ KSH niniejsza umowa przewiduje możliwość wyemitowania, w jednej lub kilku emisjach, nie później niż do dnia \textbf{31 grudnia 2036 r.}, maksymalnie \textbf{250 000 (dwustu pięćdziesięciu tysięcy)} zwykłych akcji imiennych serii F. Łączna liczba wyemitowanych i istniejących akcji serii F nie może ponadto przekroczyć 20\% wszystkich akcji Spółki istniejących bezpośrednio po danej emisji serii F.
  
\begin{MemoBox}{3.2.2}
\MemoPyt{3.2.2}{Czy określenie w umowie maksymalnej liczby akcji i terminu (2036~r.) dla serii P i F spełnia wymogi art.~300$^{103}$ KSH dotyczące upoważnienia do emisji?}
\Ocena{zgodne}
\Klasa{opcjonalne}
\Podstawa{art.~300$^{103}$ KSH \Wiar{Z}; art.~300$^{5}$ §~1 pkt~3, art.~300$^{100}$ §~2, art.~300$^{102}$ §~1--2, art.~300$^{107}$ §~3 KSH \Wiar{Z}; art.~300$^{110}$ §~1--3 KSH (pięcioletnie upoważnienie Rady, porównawczo) \Wiar{Z}; art.~444 §~1 KSH (porównawczo) \Wiar{Z}.}
\Uzasadnienie{Art.~300$^{103}$ KSH wymaga, aby umowa przewidywała „maksymalną liczbę akcji i termin ich emisji” \Wiar{Z}; obie serie to spełniają: seria P --- 176\,471 akcji do 31 grudnia 2036~r. (§~\ref{art:program-motywacyjny} ust.~1 i~3), seria F --- 250\,000 akcji do tej samej daty z dodatkowym limitem 20\% (ust.~4), który zawęża upoważnienie, więc jest dopuszczalny. Ustawa nie ogranicza długości terminu (inaczej niż pięcioletnie upoważnienie Rady z art.~300$^{110}$ §~1 i trzyletni kapitał docelowy w S.A., art.~444 §~1 \Wiar{Z}), więc dziesięcioletni horyzont jest dopuszczalny \Wiar{W}; jest natomiast decyzją biznesową, bo inwestor zażąda zgody na emisje serii F po rundzie. Dwie nieścisłości: (1)~„emitowane [\ldots] nie później niż do dnia” nie przesądza, czy przed terminem ma nastąpić uchwała, umowa objęcia czy wpis; skoro akcje nowej emisji powstają z chwilą wpisu (art.~300$^{107}$ §~3 \Wiar{Z}), sąd może odczytać termin jako termin wpisu --- bezpieczniej odnieść go do dnia uchwały o emisji; (2)~§~\ref{art:program-motywacyjny} ust.~9 dopuszcza zwiększenie puli „odrębną podstawą emisyjną”, choć każda emisja poza upoważnieniem jest zmianą umowy w formie protokołu notarialnego, zgłaszaną w sześć miesięcy od uchwały (art.~300$^{100}$ §~2, art.~300$^{102}$ §~1--2, art.~300$^{103}$ \Wiar{Z}) --- ta alternatywa nie jest prostszą ścieżką i może sugerować, że wystarczy zwykła uchwała. Limity są granicą kompetencji Walnego Zgromadzenia; ich przekroczenie daje podstawę do zaskarżenia uchwały.}
\Rekomendacja{Odnieść termin do dnia uchwały o emisji w obu paragrafach i uprościć §~\ref{art:program-motywacyjny} ust.~9 (brzmienie --- ramka przy tym ustępie). Rozważyć (decyzja Założycieli) skrócenie horyzontu serii F albo zastrzeżenie, że po Kwalifikowanej Rundzie emisje serii F wymagają zgody przewidzianej w dokumentach rundy.}
\Kierunek{Zmiana redakcyjna; usuwa spór o rozumienie terminu i mylącą sugestię o pozanotarialnym zwiększeniu puli; bez skutku dla rejestru i sądu poza jasnością badania.}
\begin{Brzmienie}
§~\ref{art:seria-f} ust.~4 zdanie pierwsze: „Na podstawie art.~300$^{103}$ KSH niniejsza umowa przewiduje możliwość wyemitowania, w jednej lub kilku emisjach, na podstawie uchwał o emisji podjętych nie później niż do dnia \textbf{31 grudnia 2036~r.}, maksymalnie \textbf{250\,000 (dwustu pięćdziesięciu tysięcy)} zwykłych akcji imiennych serii F.”
\end{Brzmienie}
\end{MemoBox}

\item Niezależnie od §~\ref{art:walne-zastrzezone}, każda emisja akcji serii F wymaga odrębnej uchwały Walnego Zgromadzenia podjętej większością 75\% Głosów Uprawnionych w Sprawie i określającej elementy wymagane przez art. 300$^{104}$ KSH. Akcjonariusz obejmujący akcje serii F jest wyłączony od głosowania nad uchwałą dotyczącą jego objęcia. Jeżeli emisja wymaga pozbawienia dotychczasowych akcjonariuszy prawa poboru w całości albo w części, odrębna uchwała w tej sprawie wymaga większości co najmniej czterech piątych głosów i spełnienia pozostałych wymogów art. 300$^{106}$ KSH.
  \item Akcje serii F są zwykłymi akcjami imiennymi. Każda akcja serii F daje prawo do jednego głosu na Walnym Zgromadzeniu oraz, na równi z akcjami serii A, prawo do dywidendy i udziału w majątku pozostałym po likwidacji, z zastrzeżeniem bezwzględnie obowiązujących przepisów prawa i postanowień niniejszej umowy.
  \item \BoilerplateStrike{Cena emisyjna, rodzaj i wartość wkładu oraz pozostałe warunki objęcia akcji serii F są określane w uchwale emisyjnej i umowie objęcia akcji zgodnie z KSH.} Jeżeli wkład obejmuje świadczenie pracy lub usług, uchwała i umowa objęcia określają rodzaj i czas świadczenia zgodnie z art. 300$^{104}$ § 1 pkt 7 KSH.
  \item Okres 48 miesięcy stopniowego zwalniania akcji serii F z obowiązku zbycia liczy się od Dnia Przyznania Serii F wskazanego w uchwale kwalifikacyjnej. Jeżeli akcje serii F powstaną później, przy ich powstaniu uwzględnia się okres nieprzerwanego, zgodnego z uchwałą zaangażowania od Dnia Przyznania Serii F; żadne prawa z akcji ani zwolnienie konkretnych akcji nie powstają jednak przed powstaniem samych akcji.
  \item Od chwili powstania akcji serii F do Założyciela Funkcjonalnego Serii F i jego akcji stosuje się odpowiednio w całości §~\ref{art:zwrotne-zbycie}, w tym zasady Odejścia Usprawiedliwionego, Odejścia Dobrowolnego, Odejścia Zawinionego, przyspieszonego zwolnienia, obowiązkowego zbycia, Wartości Godziwej i odrębnej umowy wykonawczej. Uchwała kwalifikacyjna i umowa wykonawcza określają dane niezbędne do zastosowania tego mechanizmu do konkretnej osoby.
  \item Akcje serii F podlegają na zasadach ogólnych postanowieniom o Dopuszczalnym Nabywcy, Prawnie Niedopuszczalnym Posiadaczu, prawie pierwszeństwa, prawie przyłączenia, prawie przymusowego współzbycia oraz ograniczeniu zbywania określonemu w §~\ref{art:okres-ograniczenia-zbywania}. Dla Założyciela Funkcjonalnego Serii F 12-miesięczny Okres Ograniczenia Zbywania Akcji liczy się od jego Dnia Przyznania Serii F.
  \item Seria F jest całkowicie odrębna od serii P. Emisja akcji serii F nie zmniejsza ani nie konsumuje puli programu motywacyjnego serii P, a przyznanie akcji serii P tej samej osobie wymaga odrębnej podstawy i nie jest zaliczane na limit serii F.
\end{enumerate}
\end{ContractClause}




\begin{ContractClause}{Akcjonariusze Funkcjonalni i mechanizm zwrotnego zbycia akcji}
\label{art:zwrotne-zbycie}
\begin{enumerate}
  \item Na potrzeby niniejszego paragrafu „Akcjonariuszem Funkcjonalnym” jest: (a) Założyciel Pierwotny wskazany w Załączniku nr 2, który zobowiązał się do stałego, osobistego i istotnego zaangażowania operacyjnego w rozwój Spółki, albo (b) Założyciel Funkcjonalny Serii F od chwili powstania przysługujących mu akcji serii F.
  \item Akcje Akcjonariusza Funkcjonalnego podlegają mechanizmowi stopniowego zwalniania akcji z obowiązku zbycia przez okres 48 miesięcy: dla Założyciela Pierwotnego od daty rozpoczęcia wskazanej w Załączniku nr 2, a dla Założyciela Funkcjonalnego Serii F od Dnia Przyznania Serii F określonego zgodnie z §~\ref{art:seria-f}, na następujących zasadach:
  \begin{enumerate}[label=\alph*)]
    \item przez pierwsze 12 miesięcy żadna z akcji objętych mechanizmem nie zostaje zwolniona z obowiązku zbycia;
    \item po nieprzerwanym upływie 12 miesięcy jednorazowo zwalnia się z obowiązku zbycia 25\% akcji objętych mechanizmem;
    \item przez kolejne 36 miesięcy akcje są zwalniane z obowiązku zbycia co miesiąc, w równych częściach, aż do zwolnienia 100\% akcji objętych mechanizmem.
  \end{enumerate}
  \item Do chwili zwolnienia z obowiązku zbycia akcje pozostają akcjami istniejącymi i dają prawa korporacyjne zgodnie z prawem, lecz podlegają obowiązkowi zbycia w przypadkach określonych w niniejszym paragrafie.
  \item „Odejście Usprawiedliwione” oznacza ustanie funkcjonalnego zaangażowania z powodu:
  \begin{enumerate}[label=\alph*)]
    \item śmierci, trwałej niezdolności do pracy albo poważnej choroby uniemożliwiającej dalsze wykonywanie obowiązków;
    \item rozwiązania stosunku prawnego przez Spółkę bez istotnej przyczyny leżącej po stronie Akcjonariusza Funkcjonalnego;
    \item istotnego i trwałego naruszenia przez Spółkę uzgodnionych warunków współpracy, nieusuniętego mimo pisemnego wezwania;
    \item uzgodnionego odejścia zatwierdzonego uchwałą Walnego Zgromadzenia większością 75\% wszystkich głosów;
    \item innego zdarzenia uznanego uchwałą Walnego Zgromadzenia za Odejście Usprawiedliwione.
  \end{enumerate}
  \item „Odejście Zawinione” oznacza wyłącznie ustanie zaangażowania w związku z co najmniej jednym z następujących zdarzeń:
  \begin{enumerate}[label=\alph*)]
    \item umyślnym naruszeniem lub rażącym niedbalstwem wobec Spółki;
    \item oszustwem, przywłaszczeniem, działaniem na szkodę Spółki albo popełnieniem przestępstwa istotnego dla zaufania lub bezpieczeństwa;
    \item naruszeniem poufności, praw własności intelektualnej, zakazu konkurencji, przepisów sankcyjnych, kontroli eksportu albo zasad cyberbezpieczeństwa;
    \item zatajeniem praw osoby trzeciej do technologii, kodu, patentu, danych lub sprzętu wniesionego albo wykorzystywanego przez Spółkę;
    \item porzuceniem obowiązków albo odmową ich wykonywania przez co najmniej 30 dni po bezskutecznym pisemnym wezwaniu; złożenie rezygnacji z zachowaniem co najmniej 30-dniowego okresu przejściowego oraz należyte wykonywanie obowiązków w tym okresie nie stanowi porzucenia obowiązków ani odmowy ich wykonywania w rozumieniu niniejszej litery.
  \end{enumerate}
  
\begin{MemoBox}{1.5.4}
\MemoPyt{1.5.4}{Czy katalog Odejścia Zawinionego (§~\ref{art:zwrotne-zbycie} ust.~5) i tryb jego stwierdzania nie rodzą ryzyka sporu o kwalifikację zdarzenia, i kto powinien je stwierdzać?}
\Ocena{zgodne warunkowo}
\Klasa{zalecane}
\Podstawa{art.~353$^{1}$, art.~58 §~2, art.~5 KC \Wiar{Z}; art.~42 ust.~3 Konstytucji RP (domniemanie niewinności --- pomocniczo) \Wiar{Z}; art.~52 §~1 pkt~2 KP (pomocniczo) \Wiar{Z}; art.~300$^{55}$ §~1 i art.~300$^{58}$ §~2 KSH \Wiar{Z}; art.~1157 pkt~1 KPC \Wiar{Z}}
\Uzasadnienie{Katalog jest zamknięty („wyłącznie”) i w większości obejmuje zdarzenia obiektywne (lit.~b, d, e); nieostre pozostają „rażące niedbalstwo” (lit.~a) i „działanie na szkodę Spółki” (lit.~b). Ryzyko sporu podnoszą trzy elementy: (1)~„popełnienie przestępstwa” bez wskazania, czy wymagane jest prawomocne skazanie --- stwierdzenie przestępstwa przez Radę na własną rękę naraża Spółkę na zarzut arbitralności (art.~42 ust.~3 Konstytucji; kompromis art.~52 §~1 pkt~2 KP: przestępstwo „oczywiste lub stwierdzone prawomocnym wyrokiem”), a czekanie na wyrok trwa lata; (2)~ust.~6 zd.~2 i §~4 ust.~5 umowy wspólników pozwalają „w każdym czasie” przekwalifikować Odejście Dobrowolne na Zawinione --- bez granicy czasowej i bez reguły dla akcji zwolnionych już zbytych za 100\% albo pozostałych przy założycielu; (3)~„ustanie zaangażowania w związku z” zdarzeniem jest luźne. Organ: Rada Dyrektorów z wyłączeniem zainteresowanego i dyrektorów w konflikcie (§~4 ust.~3 umowy wspólników) odpowiada art.~300$^{55}$ §~1 KSH, przy kworum z art.~300$^{58}$ §~2 KSH, i jest praktycznie jedyną możliwością: Walne Zgromadzenie to ci sami założyciele, a ekspert nie rozstrzyga wiążąco, bo klauzula ekspercka nie zamyka drogi sądowej. Spór o kwalifikację ma zdatność arbitrażową (art.~1157 pkt~1 KPC), lecz §~\ref{art:final} ust.~4 wybiera sąd powszechny; zapis na sąd polubowny wymagałby odrębnej decyzji Założycieli. Ochroną jest procedura (pisemne uzasadnienie z literą katalogu, wysłuchanie, sprzeciw) i --- kluczowe --- korekta ex post: po prawomocnym ustaleniu innego rodzaju Odejścia nabywca dopłaca różnicę do ceny właściwej, a Spółka odpowiada solidarnie; spór o akcje staje się sporem o pieniądze.}
\Rekomendacja{Doprecyzować lit.~b (przestępstwo: prawomocne skazanie albo --- do czasu wyroku --- zawieszenie wykonania prawa nabycia akcji zwolnionych, ze zbyciem akcji niezwolnionych po cenie emisyjnej bez zwłoki), ograniczyć przekwalifikowanie do 24 miesięcy od Dnia Odejścia i dodać regułę dopłaty. Stwierdzanie pozostawić Radzie Dyrektorów w trybie umowy wykonawczej.}
\Kierunek{Nowy ust.~6a §~\ref{art:zwrotne-zbycie} (do rozważenia): „Stwierdzenie Odejścia Zawinionego z powodu popełnienia przestępstwa wymaga prawomocnego skazania albo warunkowego umorzenia postępowania; do tego czasu Spółka może wstrzymać wykonanie prawa nabycia akcji zwolnionych, a obowiązek zbycia akcji niezwolnionych wykonuje się po cenie emisyjnej. Zdarzenie z ust.~5 ujawnione po ustaniu zaangażowania może stanowić podstawę stwierdzenia Odejścia Zawinionego nie później niż w terminie 24 miesięcy od dnia ustania zaangażowania. Jeżeli prawomocne orzeczenie sądu albo ugoda ustali inny rodzaj Odejścia niż stwierdzony przez Spółkę, nabywca akcji dopłaca różnicę do ceny właściwej dla ustalonego rodzaju Odejścia w terminie 30 dni, a Spółka odpowiada za tę dopłatę solidarnie.” Skutek: dla Spółki --- przewidywalność; dla założyciela --- gwarancja, że błędna kwalifikacja nie pozbawi go wartości; dla sądu --- wyraźne roszczenie pieniężne zamiast sporu o skuteczność przeniesienia akcji.}
\end{MemoBox}

\item „Odejście Dobrowolne” oznacza ustanie funkcjonalnego zaangażowania wskutek dobrowolnej rezygnacji Akcjonariusza Funkcjonalnego przed upływem 48 miesięcy, które nie stanowi Odejścia Usprawiedliwionego i nie jest związane z żadnym ze zdarzeń wymienionych w definicji Odejścia Zawinionego. Wystąpienie zdarzenia z definicji Odejścia Zawinionego przed ustaniem zaangażowania albo jego ujawnienie po ustaniu zaangażowania wyłącza kwalifikację odejścia jako Dobrowolnego.
  
\begin{MemoBox}{1.5.5}
\MemoPyt{1.5.5}{Czy mechanizm jest neutralny wobec podstawy zaangażowania założyciela (umowa o pracę, B2B, powołanie) i czy zwolnienie akcji nie rodzi skutków w prawie pracy?}
\Ocena{zgodne warunkowo}
\Klasa{zalecane}
\Podstawa{art.~18 §~1, art.~45 §~1, art.~52 §~1, art.~87 §~1, art.~91 §~1, art.~101$^{1}$--101$^{2}$, art.~114--122, art.~300 KP \Wiar{Z}; art.~300$^{2}$ §~2 KSH (wkład w postaci pracy lub usług) \Wiar{Z}; art.~300$^{73}$ §~3 i art.~300$^{74}$ §~1 KSH (odwołanie dyrektora) \Wiar{Z}; art.~12 ust.~1 i art.~17 ust.~1 pkt~6 lit.~a ustawy o PIT \Wiar{Z}}
\Uzasadnienie{Mechanizm jest formalnie neutralny: uruchamia go „ustanie funkcjonalnego zaangażowania” (ust.~4--6) bez względu na jego podstawę. Akcje objęto za wkłady (§~\ref{art:wklady}), więc zwalnianie nie jest przychodem ze stosunku pracy (art.~12 PIT), a zbycie --- przychodem z kapitałów pieniężnych (art.~17 PIT); inaczej może być przy akcjach serii F za wkład w postaci pracy (§~\ref{art:seria-f} ust.~7), co wymaga odrębnej opinii \Wiar{?}. Nieneutralne są punkty styku: (1)~ust.~4 lit.~b i ust.~5 nie pokrywają się z art.~52 i art.~45 KP, więc wyrok sądu pracy czyni kwalifikację sporną --- potrzebne są reguła dopłaty (1.5.4) i zastrzeżenie, że prawomocne ustalenie bezzasadności rozwiązania przez Spółkę oznacza Odejście Usprawiedliwione; (2)~kary umowne za naruszenie obowiązków pracowniczych SN uznaje za niedopuszczalne (art.~114--122, art.~300 KP) \Wiar{W} --- dyskonto 80\% jako cena opcji się broni (1.5.2), ale kary z umowy wykonawczej (§~10 ust.~5, §~12 ust.~4) ograniczyć do okresu po ustaniu stosunku pracy albo naruszeń spoza obowiązków pracowniczych; (3)~ceny za akcje i roszczeń Spółki nie można rozliczać z pensją (art.~87 §~1, art.~91 §~1 KP); (4)~zakaz konkurencji po ustaniu stosunku pracy wymaga odszkodowania co najmniej 25\% wynagrodzenia (art.~101$^{2}$ §~3 KP) --- §~10 ust.~3 umowy wspólników przewiduje to warunkowo, poprawnie; (5)~odwołanie dyrektora bez istotnej przyczyny (art.~300$^{74}$ §~1 KSH) trafnie stanowi Odejście Usprawiedliwione. Art.~18 §~1 KP nie jest naruszony. Zwolnienie akcji nie rodzi skutków w prawie pracy; rodzi je kierunek odwrotny: ocena rozwiązania stosunku pracy rzutuje na kwalifikację Odejścia.}
\Rekomendacja{Dodać do §~\ref{art:zwrotne-zbycie} zdanie o niezależności kwalifikacji od podstawy zaangażowania i o skutku prawomocnego ustalenia bezzasadności rozwiązania; w umowie wykonawczej ograniczyć kary umowne wobec pracowników zgodnie z pkt~(2). Preferować dla założycieli powołanie plus umowa B2B albo kontrakt menedżerski, co Założyciele powinni rozstrzygnąć świadomie.}
\Kierunek{Nowy ustęp w §~\ref{art:zwrotne-zbycie} (po ust.~6): „Rodzaj Odejścia ustala się niezależnie od podstawy prawnej zaangażowania Akcjonariusza Funkcjonalnego (stosunek pracy, umowa cywilnoprawna, powołanie do organu) i od trybu jej ustania. Jeżeli prawomocne orzeczenie ustali, że rozwiązanie stosunku prawnego przez Spółkę było bezzasadne albo nastąpiło z naruszeniem prawa, Odejście uważa się za Odejście Usprawiedliwione, a rozliczenia dokonuje się zgodnie z umową wykonawczą.” Dla Spółki: spójność z prawem pracy bez rezygnacji z mechanizmu; dla założyciela-pracownika: ochrona z KP nie jest omijana; dla inwestora: mniejsze ryzyko podważenia vestingu w due diligence.}
\end{MemoBox}

\item W przypadku Odejścia Usprawiedliwionego:
  \begin{enumerate}[label=\alph*)]
    \item akcje zwolnione z obowiązku zbycia pozostają przy Akcjonariuszu Funkcjonalnym;
    \item akcje nadal podlegające obowiązkowi zbycia podlegają zbyciu na rzecz wskazanego przez Spółkę Dopuszczalnego Nabywcy po cenie równej ich cenie emisyjnej;
    \item Walne Zgromadzenie może przyspieszyć zwolnienie z obowiązku zbycia całości albo części pozostałych akcji.
  \end{enumerate}
  \item W przypadku Odejścia Zawinionego:
  \begin{enumerate}[label=\alph*)]
    \item akcje nadal podlegające obowiązkowi zbycia podlegają zbyciu na rzecz wskazanego przez Spółkę Dopuszczalnego Nabywcy po cenie równej ich cenie emisyjnej;
    \item akcje zwolnione z obowiązku zbycia mogą zostać objęte prawem nabycia przez wskazanego przez Spółkę Dopuszczalnego Nabywcę po cenie odpowiadającej 80\% Wartości Godziwej, nie niższej jednak niż cena emisyjna i z zastrzeżeniem bezwzględnie obowiązujących przepisów prawa;
    \item powyższe nie wyłącza roszczeń Spółki o naprawienie szkody ani innych środków ochrony.
  \end{enumerate}
  
\begin{MemoBox}{1.5.2}
\MemoPyt{1.5.2}{Czy zbycie za 80\,\% Wartości Godziwej przy Odejściu Zawinionym może zostać zakwestionowane (obejście przepisów o umorzeniu lub o nabyciu akcji własnych, kara umowna, nadmierność) i jakie są granice takiej redukcji ceny?}
\Ocena{zgodne warunkowo}
\Klasa{zalecane}
\Podstawa{art.~300$^{45}$ §~1--2 KSH \Wiar{Z}; art.~300$^{47}$ §~9 KSH \Wiar{Z}; art.~353$^{1}$, art.~58 §~2--3, art.~483--484 KC \Wiar{Z}; art.~536 §~1 KC \Wiar{Z}; art.~19 ust.~1 w zw.\ z art.~17 ust.~2 oraz art.~11 ust.~2b ustawy o PIT \Wiar{Z}}
\Uzasadnienie{Zarzut obejścia przepisów o umorzeniu i akcjach własnych jest słaby: akcje nie są unicestwiane, nabywcą nie może być Spółka ani osoba działająca na jej rachunek (ust.~11--12, por.\ art.~300$^{47}$ §~9 KSH), a minimum „spłata, która nie może być niższa od wartości godziwej akcji” (art.~300$^{45}$ §~2 KSH) dotyczy wyłącznie umorzenia przymusowego. Byłby realny tylko przy nabywcy finansowanym przez Spółkę albo działającym na jej zlecenie, czego umowa wykonawcza (§~5 ust.~1) dostatecznie wprost nie wyklucza. Nie jest to też kara umowna: 80\% Wartości Godziwej to cena w opcji kupna, oznaczona przez wskazanie podstaw jej ustalenia (art.~536 §~1 KC), nie suma za niewykonanie zobowiązania niepieniężnego (art.~483 §~1 KC). Sąd może jednak badać „nadmierność” sankcji (art.~58 §~2, art.~353$^{1}$ KC), a po przekwalifikowaniu miarkować karę „rażąco wygórowaną” (art.~484 §~2 KC); 20\% dyskonta z dolną granicą ceny emisyjnej, wyłącznie dla zamkniętego katalogu z ust.~5, jest umiarkowane wobec praktyki rynkowej. Granice: (i)~dyskonto nie może obejmować akcji nabytych poza mechanizmem; (ii)~bez kumulacji z karą umowną za to samo zdarzenie bez wyraźnego zaliczenia; (iii)~przy założycielu-pracowniku sankcja cenowa, nie „potrąceniowa” (zob.~1.5.5); (iv)~stwierdzenie Odejścia Zawinionego musi być weryfikowalne sądownie. Podatkowo: gdy cena „bez uzasadnionej przyczyny, znacznie odbiega od wartości rynkowej”, organ może określić przychód w wysokości wartości rynkowej (art.~19 ust.~1 w zw.\ z art.~17 ust.~2 PIT), a nabywca uzyskuje świadczenie częściowo odpłatne (art.~11 ust.~2b PIT); mechanizm umowny jest „uzasadnioną przyczyną”, co trzeba udokumentować.}
\Rekomendacja{Utrzymać 80\%; w ust.~8 lit.~b dodać zdanie, że cena jest ceną umowną opcji i nie stanowi kary umownej, a różnicę między Wartością Godziwą a ceną zapłaconą zalicza się na ewentualne roszczenia odszkodowawcze Spółki; w umowie wykonawczej wykluczyć jako Nabywcę Wskazanego podmiot zależny od Spółki lub przez nią finansowany.}
\Kierunek{Dopisek w §~\ref{art:zwrotne-zbycie} ust.~8 lit.~b: „Cena określona w niniejszej literze jest ceną umowną prawa nabycia i nie stanowi kary umownej; różnica między Wartością Godziwą a ceną zapłaconą podlega zaliczeniu na poczet roszczeń Spółki o naprawienie szkody wynikłej z tego samego zdarzenia.” Dla Spółki: mniejsze ryzyko miarkowania; dla założyciela: ochrona przed podwójną sankcją; dla inwestora: 80\% jest łagodniejsze od standardu VC i fundusz może żądać zaostrzenia w rundzie (decyzja Założycieli). Podatkowo: przed pierwszym Odejściem uzyskać interpretację indywidualną co do skutków sprzedaży po cenie emisyjnej i po 80\% Wartości Godziwej.}
\end{MemoBox}

\item W przypadku Odejścia Dobrowolnego:
  \begin{enumerate}[label=\alph*)]
    \item akcje nadal podlegające obowiązkowi zbycia podlegają zbyciu na rzecz wskazanego przez Spółkę Dopuszczalnego Nabywcy po cenie równej ich cenie emisyjnej;
    \item akcje zwolnione z obowiązku zbycia pozostają przy Akcjonariuszu Funkcjonalnym; Spółka może jednak, w terminie 6 miesięcy od dnia Odejścia Dobrowolnego, wskazać Dopuszczalnego Nabywcę uprawnionego do nabycia całości albo części tych akcji po cenie równej 100\% Wartości Godziwej; po bezskutecznym upływie tego terminu prawo nabycia wygasa;
    \item Odejście Dobrowolne nie stanowi samo w sobie naruszenia niniejszej umowy i nie uzasadnia sankcji przewidzianych dla Odejścia Zawinionego; nie uchyla natomiast obowiązków trwających po ustaniu zaangażowania, w szczególności wynikających z §~\ref{art:wlasnosc-intelektualna}, §~\ref{art:security}, §~\ref{art:export} i §~\ref{art:konkurencja}.
  \end{enumerate}
  \item W razie zmiany kontroli nad Spółką, a następnie — w ciągu 12 miesięcy — rozwiązania współpracy z Akcjonariuszem Funkcjonalnym przez Spółkę bez istotnej przyczyny, wszystkie pozostałe akcje objęte mechanizmem zostają zwolnione z obowiązku zbycia.
  
\begin{MemoBox}{1.5.6}
\MemoPyt{1.5.6}{Przyspieszenie po zmianie kontroli (§~\ref{art:zwrotne-zbycie} ust.~10) --- czy definicja „istotnej przyczyny” i 12-miesięczne okno są wystarczająco określone?}
\Ocena{zgodne warunkowo}
\Klasa{zalecane}
\Podstawa{art.~353$^{1}$, art.~65 KC \Wiar{Z}; art.~4 §~1 pkt~4 KSH (definicja spółki dominującej --- pomocniczo do pojęcia kontroli) \Wiar{Z}}
\Uzasadnienie{Ust.~10 nie definiuje ani „zmiany kontroli”, ani „istotnej przyczyny”. Pierwsze pojęcie definiuje dopiero umowa wykonawcza (§~9 ust.~3: ponad 50\% głosów albo prawo powoływania większości dyrektorów), co jest odwróceniem hierarchii --- przesłanka skutku korporacyjnego z umowy spółki nie powinna zależeć od umowy, której inwestor nie jest stroną (§~2 ust.~4 umowy wspólników). „Istotna przyczyna” (ust.~4 lit.~b i ust.~10) nie ma definicji, mimo gotowego katalogu zawinionych przyczyn w ust.~5; wykładnia z art.~65 KC prowadzi zapewne do utożsamienia jej z ust.~5, ale nabywca kontroli będzie argumentował szerzej (np.\ „utrata zaufania”, „reorganizacja”). Dwunastomiesięczne okno jest standardem rynkowym i jest określone wystarczająco; brakuje natomiast (a)~objęcia oknem rozwiązania współpracy od ogłoszenia transakcji do jej zamknięcia, (b)~konstruktywnego odejścia po zmianie kontroli (istotne pogorszenie roli lub wynagrodzenia; ust.~4 lit.~c pokrywa to tylko częściowo), (c)~zbycia całości lub zasadniczej części przedsiębiorstwa albo Kluczowej Własności Intelektualnej jako zmiany kontroli. Nieokreśloność nie czyni postanowienia nieważnym, ale przenosi spór na wykładnię, gdzie Spółka po zmianie kontroli jest stroną silniejszą.}
\Rekomendacja{Zdefiniować w §~\ref{art:zwrotne-zbycie} „Zmianę Kontroli” i „Istotną Przyczynę” (jako zdarzenia z ust.~5), rozszerzyć okno o okres od zawarcia umowy zmiany kontroli do jej wykonania i objąć nim rezygnację z powodu istotnego pogorszenia warunków.}
\begin{Brzmienie}
§~\ref{art:zwrotne-zbycie} ust.~10 (nowe brzmienie): „»Zmiana Kontroli« oznacza nabycie przez osobę albo osoby działające w porozumieniu, inne niż Założyciele Pierwotni, bezpośrednio lub pośrednio, akcji dających ponad 50\% głosów w Spółce, prawa do powoływania większości członków Rady Dyrektorów albo całości lub zasadniczej części przedsiębiorstwa Spółki lub Kluczowej Własności Intelektualnej. »Istotną Przyczyną« jest wyłącznie zdarzenie wymienione w ust.~5. Jeżeli w okresie od zawarcia umowy prowadzącej do Zmiany Kontroli do upływu 12 miesięcy od jej dokonania Spółka rozwiąże współpracę z Akcjonariuszem Funkcjonalnym bez Istotnej Przyczyny albo Akcjonariusz Funkcjonalny rozwiąże ją z powodu istotnego pogorszenia jego roli, zakresu odpowiedzialności lub wynagrodzenia, nieusuniętego w terminie 30 dni od pisemnego wezwania, wszystkie pozostałe akcje objęte mechanizmem zostają zwolnione z obowiązku zbycia z dniem ustania zaangażowania.”
\end{Brzmienie}
\end{MemoBox}

\item Mechanizm zwrotnego zbycia akcji nie stanowi umorzenia akcji ani zobowiązania Spółki do nabycia akcji własnych. Nabycie akcji własnych przez Spółkę albo umorzenie akcji może nastąpić wyłącznie w odrębnym trybie i po spełnieniu właściwych wymogów prawa.
  \item Nabywcą akcji w wykonaniu niniejszego mechanizmu nie może być Spółka ani osoba trzecia działająca na rachunek Spółki, chyba że nabycie następuje w odrębnym trybie zgodnym z art. 300$^{47}$ KSH i po spełnieniu wszystkich właściwych wymogów prawa.
  \item Szczegółowy tryb wykonania obowiązku zbycia, w tym zdarzenie uruchamiające procedurę, sposób stwierdzenia Odejścia Usprawiedliwionego, Odejścia Dobrowolnego albo Odejścia Zawinionego, tryb wykonania i rozliczenia prawa nabycia akcji zwolnionych z obowiązku zbycia przy Odejściu Dobrowolnym (w tym sposób liczenia 6-miesięcznego terminu i skutek jego bezskutecznego upływu), wskazanie Dopuszczalnego Nabywcy, ofertę sprzedaży, ograniczone pełnomocnictwo, terminy wykonania, sposób zapłaty i moment przeniesienia akcji, określa odrębna umowa wykonawcza podpisana przez każdego Akcjonariusza Funkcjonalnego najpóźniej w dniu objęcia mechanizmem.
  
\begin{MemoBox}{1.5.3}
\MemoPyt{1.5.3}{Umowa wykonawcza: czy nieodwołalne oferty i pełnomocnictwa (art.~101 §~2 i art.~108 KC) mogą zostać skonstruowane tak, aby mechanizm zadziałał bez dalszego podpisu akcjonariusza, i czy orzeczenie zastępujące oświadczenie woli (art.~64 KC, art.~1047 KPC) jest realną drogą wykonania? Jak wygląda minimalna treść takiej umowy (zdarzenie, stwierdzenie Odejścia, nabywca, cena, terminy)?}
\Ocena{zgodne warunkowo --- konstrukcja poprawna co do zasady, projekt umowy wykonawczej ma braki, z których dwa są istotne}
\Klasa{konieczne (K4)}
\Podstawa{art.~66 §~1--2, art.~66$^{2}$ §~2, art.~89, art.~99 §~1, art.~101 §~1--2, art.~106, art.~108, art.~389--390, art.~393, art.~536 §~1 KC \Wiar{Z}; art.~64 KC, art.~1047 §~1--2 i art.~786 §~1 KPC \Wiar{Z}; art.~300$^{36}$ §~4, art.~300$^{34}$ §~3--5, art.~300$^{37}$ §~1 KSH \Wiar{Z}; SN V CSK 522/18 i II CSKP 593/22 \Wiar{Z} (tezy według \textit{psa\_feedback.tex}; teksty orzeczeń niedostępne); art.~4 pkt~1 i art.~7 ust.~1 pkt~1 lit.~b ustawy o PCC \Wiar{Z}}
\Uzasadnienie{\textbf{Oferta.} Oferty z terminem związania po dojściu do adresata nie da się odwołać (art.~66 §~2 KC a contrario; w obrocie B2B art.~66$^{2}$ §~2 KC), więc „nieodwołalność” jest skuteczna i w istocie deklaratoryjna. Wystarcza forma dokumentowa (art.~300$^{36}$ §~4 KSH), a „oświadczenie akcjonariusza o zobowiązaniu do przeniesienia akcji” jest samodzielną podstawą wpisu (art.~300$^{34}$ §~4 zd.~2 KSH) --- tak warto zredagować ofertę albo umowę przedwstępną. \textbf{Brak nr~1}: oferta z Załącznika~A jest skierowana do osoby nieoznaczonej; bezpieczniej: oferta do Spółki z prawem wskazania osoby trzeciej (art.~393 KC) plus umowa przedwstępna w formie dokumentowej (art.~389 KC) o silniejszym skutku (art.~390 §~2 KC) --- najpewniejsza podstawa roszczenia z art.~64 KC.

\textbf{Pełnomocnictwo.} Zrzeczenie się odwołania i niewygasanie ze śmiercią są dopuszczalne, gdy uzasadnia je stosunek podstawowy (art.~101 §~1--2 KC; SN V CSK 522/18); art.~108 KC jest zachowany, bo pełnomocnik nie jest nabywcą. \textbf{Brak nr~2}: pełnomocnictwo z Załącznika~B dla „osoby wskazanej przez Radę Dyrektorów” jest in blanco i rejestr go nie przyjmie. Pełnomocnikiem powinna być Spółka, z substytucją (art.~106 KC) i zakazem użycia przy nabyciu przez Spółkę (art.~300$^{47}$ KSH). Wystarcza forma dokumentowa (art.~99 §~1 KC), lecz podpis notarialnie poświadczony powinien być obowiązkowy: od 18.02.2027~r.\ zgoda zbywcy na wpis wymaga takiego podpisu, podpisu złożonego przed osobą upoważnioną przez podmiot prowadzący rejestr albo kwalifikowanego podpisu elektronicznego (art.~300$^{34}$ §~3 zd.~2 KSH, Dz.U. 2026 poz.~176).

\textbf{Art.~64 KC.} Droga realna, ale wolna (12--30 miesięcy); art.~64 KC obowiązku nie kreuje (SN II CSKP 593/22) --- oznaczają go §~\ref{art:zwrotne-zbycie} i uchwały o stwierdzeniu Odejścia oraz o wskazaniu nabywcy z ceną. Gdy oświadczenie zależy od zapłaty ceny, skutek powstaje dopiero z prawomocnym nadaniem klauzuli wykonalności (art.~1047 §~2 KPC) po udokumentowaniu zapłaty (art.~786 §~1 KPC) --- depozyt powinien być notarialny albo sądowy. Do wyroku założyciel głosuje; osłoną jest tylko zabezpieczenie. Ścieżką główną musi pozostać pełnomocnictwo, art.~64 KC --- rezerwową.

\textbf{Minimalna treść.} Projekt ma elementy węzłowe; brakuje: (a)~skutku niezapłacenia ceny przez Nabywcę Wskazanego; (b)~wstrzymania terminu z ust.~9 lit.~b na czas sprzeciwu i mediacji; (c)~granicy czasowej stwierdzenia Odejścia Zawinionego; (d)~rozstrzygnięcia zbiegu z §~\ref{art:pierwszenstwo}; (e)~dopłaty po zmianie kwalifikacji (zob.~1.5.4); (f)~PCC 1\% (obowiązek kupującego) i obowiązków płatnika; (g)~depozytu oryginałów oferty i pełnomocnictwa. Ust.~13 powinien je wymieniać jako obowiązkową treść umowy wykonawczej.}
\Rekomendacja{Przebudować Załączniki A--B (oferta do Spółki z prawem wskazania nabywcy plus umowa przedwstępna; pełnomocnictwo dla Spółki z substytucją, z podpisem notarialnie poświadczonym), uzupełnić braki (a)--(g), a w ust.~13 rozszerzyć katalog obowiązkowej treści. Umowę wykonawczą podpisać tego samego dnia co umowę Spółki; wzory uzgodnić z podmiotem prowadzącym rejestr.}
\begin{Brzmienie}
§~\ref{art:zwrotne-zbycie} ust.~13 (nowe brzmienie): „Szczegółowy tryb wykonania obowiązku zbycia określa odrębna umowa wykonawcza podpisana przez każdego Akcjonariusza Funkcjonalnego najpóźniej w dniu objęcia mechanizmem. Umowa wykonawcza określa co najmniej: zdarzenie uruchamiające procedurę i sposób zawiadomienia o nim; organ, tryb i termin stwierdzenia Odejścia Usprawiedliwionego, Odejścia Dobrowolnego albo Odejścia Zawinionego, w tym prawo Akcjonariusza Funkcjonalnego do wysłuchania i sprzeciwu oraz skutki zmiany kwalifikacji Odejścia po przeniesieniu akcji; termin, w którym Odejście Zawinione może zostać stwierdzone po ustaniu zaangażowania; tryb wskazania Dopuszczalnego Nabywcy i jego stosunek do §~\ref{art:pierwszenstwo}; nieodwołalną ofertę sprzedaży albo umowę przedwstępną sprzedaży akcji oraz pełnomocnictwo ograniczone do czynności niezbędnych do wykonania obowiązku zbycia, udzielone osobie oznaczonej, z podpisem notarialnie poświadczonym; tryb wykonania i rozliczenia prawa nabycia akcji zwolnionych z obowiązku zbycia przy Odejściu Dobrowolnym, w tym sposób liczenia 6-miesięcznego terminu, jego wstrzymanie na czas sporu i skutek bezskutecznego upływu; cenę albo sposób jej ustalenia, termin i sposób zapłaty, złożenie ceny do depozytu, skutki niezapłacenia ceny przez wskazanego nabywcę oraz moment przeniesienia akcji.”
\end{Brzmienie}
\end{MemoBox}

\item Obowiązek zbycia akcji określony w niniejszym paragrafie stanowi obowiązek związany z akcją wobec Spółki. Spółka podejmuje czynności niezbędne do ujawnienia właściwych ograniczeń rozporządzania akcją i obowiązków związanych z akcją w rejestrze akcjonariuszy w zakresie przewidzianym przez KSH.
  
\begin{MemoBox}{1.5.1}
\MemoPyt{1.5.1}{Czy obowiązek zbycia jako „obowiązek związany z akcją” jest skuteczny wobec nabywców wtórnych i wobec podmiotu prowadzącego rejestr akcjonariuszy; jak go ujawnić w rejestrze?}
\Ocena{zgodne warunkowo --- skuteczność wobec Spółki i rejestru tak, wobec nabywcy wtórnego z luką w §~\ref{art:zbywanie} ust.~2}
\Klasa{konieczne (K2, K3)}
\Podstawa{art.~300$^{1}$ §~3 KSH \Wiar{Z}; art.~300$^{33}$ §~1 pkt~10--11 i §~2 KSH \Wiar{Z}; art.~300$^{33}$ §~3 i art.~594 §~1 pkt~2$^{1}$ KSH w brzmieniu od 18.02.2027~r.\ (Dz.U. 2026 poz.~176) \Wiar{Z}; art.~300$^{34}$ §~1 i~4--6 KSH \Wiar{Z}; art.~300$^{37}$ §~1 i art.~300$^{38}$ §~1 KSH \Wiar{Z}; art.~300$^{39}$ §~1--2 KSH \Wiar{Z}}
\Uzasadnienie{Obowiązek zbycia jest świadczeniem wobec Spółki „określonym w umowie spółki” (art.~300$^{1}$ §~3 KSH), obowiązkiem korporacyjnym obciążającym każdoczesnego posiadacza akcji; kwalifikacja w ust.~14 jest prawidłowa. Rejestr ujawnia z mocy ustawy ograniczenia rozporządzania i obowiązki związane z akcją wobec spółki (art.~300$^{33}$ §~1 pkt~10--11 KSH); wpis następuje na żądanie Spółki, nie później niż w siedem dni, bez zgody akcjonariusza (art.~300$^{34}$ §~1 i~4 KSH). Od 18.02.2027~r.\ zgłaszanie zmian tych danych będzie obowiązkiem organu zarządzającego pod groźbą grzywny (art.~300$^{33}$ §~3 i art.~594 §~1 pkt~2$^{1}$ KSH w brzmieniu nadanym Dz.U. 2026 poz.~176). Rejestr „uwzględnia ograniczenia co do rozporządzania akcją” (art.~300$^{34}$ §~6 KSH), więc przy konstytutywnym wpisie przejścia akcji (art.~300$^{37}$ §~1 KSH) sprawdzi zgodę Spółki z §~\ref{art:zbywanie}; obowiązek związany z akcją nie jest natomiast ograniczeniem rozporządzania --- rejestr bada treść i formę dokumentów, nie ich zgodność z prawem (art.~300$^{34}$ §~5 KSH), i nie będzie egzekwował obowiązku zbycia ani odmawiał z tego powodu wpisu nabywcy wtórnego.

Wobec nabywcy wtórnego obowiązek przechodzi z akcją, lecz mechanizm jest w istocie osobisty: opiera się na ofercie i pełnomocnictwie z umowy wykonawczej, które nabywcy nie wiążą. Po Okresie Ograniczenia (§~\ref{art:okres-ograniczenia-zbywania}) założyciel może zbyć akcje niezwolnione osobie trzeciej, a zamknięty katalog §~\ref{art:zbywanie} ust.~2 nie pozwala odmówić zgody z powodu nieprzystąpienia nabywcy do obowiązków z §~\ref{art:zwrotne-zbycie}. Wpis luki nie zamyka, bo rejestr musi uwzględnić dopiero odmowę zgody; trzeba ją więc zamknąć w §~\ref{art:zbywanie} ust.~2.}
\Rekomendacja{Dodać w §~\ref{art:zbywanie} ust.~2 podstawę odmowy zgody dla akcji objętych obowiązkiem zbycia, gdy nabywca nie przystąpił do umowy wykonawczej (oferta i pełnomocnictwo), oraz doprecyzować w ust.~14, że obowiązek obciąża każdoczesnego posiadacza akcji, a zdarzenia dotyczące Akcjonariusza Funkcjonalnego uruchamiają go także wobec nabywcy. Przed zawiązaniem uzgodnić z wybranym podmiotem prowadzącym rejestr (notariusz albo dom maklerski) brzmienie wzmianek do wpisu. Brzmienie nowej lit.~f) --- w ramce przy §~\ref{art:zbywanie} ust.~2.}
\begin{Brzmienie}
§~\ref{art:zwrotne-zbycie} ust.~14 (nowe brzmienie): „Obowiązek zbycia akcji określony w niniejszym paragrafie stanowi obowiązek związany z akcją wobec Spółki i obciąża każdoczesnego posiadacza akcji objętych mechanizmem; zdarzenia dotyczące Akcjonariusza Funkcjonalnego, od którego akcje pochodzą, uruchamiają obowiązek zbycia także wobec nabywcy tych akcji. Spółka niezwłocznie po zawiązaniu, a następnie po każdej zmianie, żąda ujawnienia w rejestrze akcjonariuszy wzmianki o ograniczeniu rozporządzania akcjami objętymi mechanizmem oraz o obowiązku zbycia jako obowiązku związanym z akcją, w zakresie przewidzianym przez art.~300$^{33}$ KSH.”
\end{Brzmienie}
\end{MemoBox}

\item W razie niewykonania obowiązku zbycia przez Akcjonariusza Funkcjonalnego Spółka albo uprawniony Dopuszczalny Nabywca może dochodzić wykonania obowiązku zgodnie z prawem, w tym --- jeżeli zachodzą ustawowe przesłanki --- wydania orzeczenia zastępującego wymagane oświadczenie woli zgodnie z art. 64 Kodeksu cywilnego w związku z art. 1047 § 1 Kodeksu postępowania cywilnego. Odrębna umowa wykonawcza powinna określać pełnomocnika w sposób wykluczający niedopuszczalną czynność z samym sobą lub reprezentowanie obu stron bez wymaganej podstawy prawnej.
\end{enumerate}
\end{ContractClause}




\begin{ContractClause}{Dopuszczalny Nabywca, Niedopuszczalny Nabywca i zmiana stanu prawnego po nabyciu}
\label{art:dopuszczalny-nabywca}
\begin{enumerate}
  \item „Dopuszczalnym Nabywcą” jest osoba fizyczna, osoba prawna albo inna jednostka organizacyjna, której nabycie akcji nie jest w chwili planowanego nabycia objęte bezwzględnie obowiązującym zakazem albo ograniczeniem prawnym uniemożliwiającym dokonanie transakcji.
  \item „Niedopuszczalnym Nabywcą” jest osoba lub podmiot, którego nabycie akcji jest w chwili planowanego nabycia objęte bezwzględnie obowiązującym zakazem albo ograniczeniem prawnym uniemożliwiającym dokonanie transakcji, w szczególności wiążącym zakazem dokonania transakcji, zakazem udostępnienia środków lub zasobów gospodarczych, zamrożeniem aktywów albo innym środkiem prawnym wyłączającym możliwość skutecznego nabycia akcji.
  \item „Prawnie Niedopuszczalnym Posiadaczem” jest akcjonariusz, który nabył akcje skutecznie i zgodnie z prawem, lecz po ich nabyciu powstał dotyczący tej osoby albo wykonywania praw z akcji bezwzględnie obowiązujący zakaz lub ograniczenie prawne dotyczące dalszego posiadania akcji, wykonywania praw z akcji, otrzymywania świadczeń z akcji albo rozporządzania akcjami, w szczególności obowiązek zamrożenia aktywów, zawieszenia określonych uprawnień albo zastosowania innego środka prawnego.
  \item Weryfikacja nabywcy służy ustaleniu, czy planowany nabywca nie jest Niedopuszczalnym Nabywcą oraz czy spełnia Kryterium Bezpieczeństwa EU/NATO. W szczególności może obejmować ustalenie tożsamości nabywcy, jego struktury własności i kontroli, beneficjentów rzeczywistych oraz innych informacji niezbędnych do ustalenia, czy transakcja jest prawnie dopuszczalna i czy nabywca spełnia Kryterium Bezpieczeństwa EU/NATO.
  \item Sama narodowość, obywatelstwo, siedziba, miejsce zamieszkania albo jurysdykcja pochodzenia nabywcy nie czynią go Niedopuszczalnym Nabywcą, jeżeli nabycie akcji przez tę osobę lub podmiot jest zgodne z bezwzględnie obowiązującymi przepisami prawa. Niezależnie od tego, nabycie akcji przez osobę lub podmiot niespełniający Kryterium Bezpieczeństwa EU/NATO wymaga zgody Walnego Zgromadzenia zgodnie z ust.~6; w jej braku Spółka odmawia zgody na rozporządzenie akcją na zasadach §~\ref{art:zbywanie} ust.~2 lit.~b, z mechanizmem wskazania innego nabywcy i zapłaty ratalnej określonym w §~\ref{art:zbywanie} ust.~3.
  \item „Kryterium Bezpieczeństwa EU/NATO” spełnia:
  \begin{enumerate}[label=\alph*)]
    \item osoba fizyczna posiadająca obywatelstwo państwa członkowskiego Unii Europejskiej, Europejskiego Obszaru Gospodarczego albo NATO lub mająca stałe miejsce zamieszkania w takim państwie;
    \item osoba prawna albo jednostka organizacyjna, która ma siedzibę i rzeczywisty ośrodek zarządzania w państwie, o którym mowa w lit.~a, nie jest bezpośrednio ani pośrednio kontrolowana przez osobę lub podmiot spoza tych państw, a jej beneficjenci rzeczywiści spełniają kryterium określone w lit.~a.
  \end{enumerate}
  Kryterium Bezpieczeństwa EU/NATO nie spełnia osoba ani podmiot, który nabywa lub posiada akcje na rachunek, w interesie albo na polecenie osoby lub podmiotu niespełniającego tego Kryterium, w szczególności jako powiernik, pełnomocnik albo strona porozumienia dotyczącego wykonywania praw z akcji. Walne Zgromadzenie może, uchwałą podjętą większością 75\% wszystkich głosów, stwierdzić spełnienie Kryterium Bezpieczeństwa EU/NATO w przypadku granicznym albo wyrazić zgodę na nabycie akcji przez osobę lub podmiot niespełniający tego Kryterium, jeżeli nabycie jest zgodne z bezwzględnie obowiązującymi przepisami prawa.
  
\begin{MemoBox}{1.1.1}
\MemoPyt{1.1.1}{Czy ograniczenie kręgu nabywców akcji, spadkobierców, małżonków i dyrektorów według obywatelstwa, siedziby i struktury kontroli jest dopuszczalne na tle KSH (art.~300$^{39}$ i 300$^{41}$), swobód traktatowych UE (przepływ kapitału, swoboda przedsiębiorczości) i zakazu dyskryminacji? Czy odmienne traktowanie państw UE/EOG i pozostałych państw NATO wymaga innej konstrukcji?}
\Ocena{zgodne warunkowo}
\Klasa{zalecane}
\Podstawa{art.~300$^{39}$ §~1--2 KSH \Wiar{Z}; art.~300$^{41}$ §~1 KSH \Wiar{Z}; art.~300$^{1}$ §~3 i art.~2 KSH w zw. z art.~353$^{1}$, art.~57 i art.~58 KC \Wiar{Z}; art.~18, 49, 63, 65 ust.~1 lit.~b i art.~346 ust.~1 lit.~b TFUE \Wiar{W}; art.~2 pkt~6, 7 i~24, art.~5 oraz art.~9 ust.~1--4 i~7 rozporządzenia (UE) 2021/697 (EDF) \Wiar{Z}; art.~2 pkt~1, 5 i~7, art.~4 ust.~15 lit.~a--b i ust.~17, art.~30 ust.~1, art.~31 oraz motywy 14, 15 i~20 rozporządzenia (UE) 2026/1386 (kontrola inwestycji zagranicznych, stosowane od 17 stycznia 2028~r.) \Wiar{Z} (tekst EN).}
\Uzasadnienie{KSH dopuszcza ograniczenie: art.~300$^{39}$ §~1 KSH pozwala rozporządzenie akcją „w inny sposób je ograniczyć”, bez zamkniętego katalogu; obywatelstwo, siedziba i kontrola to klasyczny „inny sposób” \Wiar{W}. Art.~300$^{41}$ §~1 KSH pozwala ograniczyć wstąpienie spadkobierców, wymagając określenia warunków spłaty niewstępujących „pod rygorem bezskuteczności” --- projekt to spełnia (§~\ref{art:dziedziczenie} ust.~2 i~4). Małżonek --- zasady ogólne (§~\ref{art:malzonek} ust.~3), bezpieczniejsze niż wyłączenie wstąpienia; dyrektorzy --- zob.~1.1.5. Granice (art.~353$^{1}$ i~58 KC): ograniczenie nie może trwale wyłączać zbywalności ani obchodzić przepisów o umorzeniu i akcjach własnych; projekt broni się „wentylem” Walnego Zgromadzenia (ust.~6 zd.~ostatnie) i ceną równą Wartości Godziwej.

Swobody z art.~49 i~63 TFUE wiążą przede wszystkim państwa; wobec umowy spółki prywatnej nie działają horyzontalnie jak np.~zakaz dyskryminacji pracowników \Wiar{W}. Kryterium nie różnicuje podmiotów z UE/EOG między sobą (art.~18 TFUE); nawet przy badaniu z art.~63 uzasadnia je bezpieczeństwo publiczne (art.~65 ust.~1 lit.~b) i art.~346 ust.~1 lit.~b TFUE, a proporcjonalność wzmacnia zgoda Walnego Zgromadzenia \Wiar{W}. Ryzyko traktatowe --- niskie; krajowe (art.~58 KC) --- umiarkowane, zależne od wariantu z pkt~1.2.

Odmienne traktowanie UE/EOG i NATO nie jest problemem prawnym (Kryterium traktuje je identycznie), lecz gospodarczym. Art.~9 rozporządzenia 2021/697 (EDF) wymaga siedziby odbiorców i ich zarządczych struktur wykonawczych w Unii lub państwie stowarzyszonym oraz braku kontroli niestowarzyszonego państwa trzeciego lub jego podmiotu (ust.~1--3), z wyjątkiem gwarancji zatwierdzonych przez państwo członkowskie (ust.~4) \Wiar{Z}. Stowarzyszeni są członkowie EFTA w EOG (art.~5), nie USA, Wielka Brytania, Kanada czy Turcja \Wiar{Z}; obecnie faktycznie tylko Norwegia \Wiar{Z} --- listę potwierdzić u Komisji. AGILE i EUDIS przejmują ten test \Wiar{W}. Kryterium dopuszcza więc strukturę wykluczającą Spółkę z EDF albo wymuszającą procedurę gwarancyjną; potrzebna jest nie inna konstrukcja, lecz druga, węższa warstwa dotycząca wyłącznie \emph{kontroli} na czas korzystania z programów.

Druga różnica wynika z prawa publicznego: od 17 stycznia 2028~r. rozporządzenie (UE) 2026/1386 nakazuje państwom objąć uprzednim zezwoleniem inwestycję inwestora spoza Unii (także przez spółkę zależną w Unii), bez wyjątku dla NATO, OECD i EOG, dającą „skuteczny udział w zarządzaniu lub kontroli” podmiotu opracowującego, produkującego lub wprowadzającego do obrotu produkty podwójnego zastosowania albo obronne (art.~2 pkt~1, 5 i~7, art.~4 ust.~15 lit.~a--b, motyw~15; bez nabyć czysto portfelowych, motyw~14; progi głosów może określić ustawa krajowa, motyw~20) \Wiar{Z} (tekst EN); dwa stany prawne --- pkt~1.1.2. Umowa może to tylko uwzględnić jako warunek (§~\ref{art:zbywanie} ust.~2 lit.~e).}
\Rekomendacja{Zachować Kryterium w obecnym kształcie. Dodać w §~\ref{art:dopuszczalny-nabywca} ust.~6 zdanie o dodatkowej przesłance dla nabycia \emph{kontroli} (UE/EOG/państwa stowarzyszone z EDF) obowiązującej w okresie korzystania z programów, z możliwością zgody Walnego Zgromadzenia większością 75\%; alternatywnie zapisać to w umowie akcjonariuszy jako zobowiązanie do niedopuszczenia do takiej kontroli. Wybór miejsca należy do Założycieli: umowa spółki daje skuteczność wobec Spółki i rejestru, umowa akcjonariuszy --- elastyczność.}
\Kierunek{Warstwa „kontroli UE/EOG” bez zmiany zasady Kryterium: zachowanie kwalifikowalności EDF; akcjonariuszy z NATO spoza UE nie dotyczy, dopóki nie dążą do kontroli; inwestor zna zasadę przed rundą; w rejestrze dodatkowa adnotacja o ograniczeniu; dla sądu rejestrowego bez wpływu.
\begin{Brzmienie}
(§~\ref{art:dopuszczalny-nabywca} ust.~6, zdanie dodawane po zdaniu o Walnym Zgromadzeniu) „W okresie, w którym Spółka jest beneficjentem finansowania z Europejskiego Funduszu Obronnego albo innego programu Unii Europejskiej uzależniającego finansowanie od struktury kontroli nad beneficjentem, nabycie akcji albo praw z akcji prowadzące do uzyskania bezpośredniej albo pośredniej kontroli nad Spółką przez osobę lub podmiot spoza państw członkowskich Unii Europejskiej, Europejskiego Obszaru Gospodarczego albo państw stowarzyszonych z danym programem wymaga zgody Walnego Zgromadzenia wyrażonej większością 75\% wszystkich głosów; do odmowy zgody Spółki stosuje się §~\ref{art:zbywanie} ust.~2 lit.~b i ust.~3.”
\end{Brzmienie}}
\end{MemoBox}

\begin{MemoBox}{1.1.3}
\MemoPyt{1.1.3}{Czy wyłączenie nabycia na cudzy rachunek (§~11 ust.~6) jest wystarczająco określone i egzekwowalne --- w szczególności wobec podmiotu prowadzącego rejestr akcjonariuszy --- oraz czy oświadczenie nabywcy z §~11 ust.~8 wystarcza jako narzędzie weryfikacji? Jakie skutki naruszenia (§~12 ust.~8, §~36) są realnie osiągalne?}
\Ocena{zgodne warunkowo}
\Klasa{zalecane}
\Podstawa{art.~300$^{33}$ §~1 pkt~10--11, art.~300$^{34}$ §~4--6, art.~300$^{37}$ §~1 i art.~300$^{38}$ §~1 KSH \Wiar{Z}; art.~83 i art.~58 KC \Wiar{Z}; art.~87 ust.~1 pkt~5--6 ustawy o ofercie publicznej (definicja działania w porozumieniu, jako wzorzec redakcyjny) \Wiar{W}; art.~471 i art.~483 KC \Wiar{Z}.}
\Uzasadnienie{Klauzula jest wystarczająco określona co do istoty (rachunek, interes, polecenie; powiernik, pełnomocnik, strona porozumienia), lecz „w interesie” jest zwrotem ocennym, a „porozumienie dotyczące wykonywania praw z akcji” nie zostało zdefiniowane. Wobec podmiotu prowadzącego rejestr klauzula nie jest egzekwowalna bezpośrednio: bada on treść i formę dokumentów, ale nie ma „obowiązku badania zgodności z prawem oraz prawdziwości dokumentów”, chyba że poweźmie uzasadnione wątpliwości (art.~300$^{34}$ §~5 KSH), a więc nie bada stosunków wewnętrznych między nabywcą a jego mocodawcą \Wiar{Z}; wpisze osobę, która przedstawi umowę zbycia i zgodę Spółki. Bramką jest więc wyłącznie Spółka na etapie zgody (oświadczenie z ust.~8, dokumenty struktury właścicielskiej, publiczne wykazy z ust.~9). Oświadczenie wystarcza na etapie wejścia (fałszywe rodzi odpowiedzialność odszkodowawczą, a przy umyślności w niektórych stanach faktycznych karną z art.~286 KK \Wiar{?}), ale nie na etapie trwania stosunku: brak obowiązku jego odnawiania i zawiadomienia o porozumieniu zawartym po nabyciu, a ust.~15 obejmuje wprawdzie „objęcie akcji na rachunek osoby niespełniającej Kryterium”, ale zależy od lojalnego zawiadomienia. Realnie osiągalne skutki: (a) obowiązek zbycia z ust.~15, wykonalny przez art.~64 KC i art.~1047 KPC, jeżeli umowa wykonawcza obejmie także inwestorów (dziś tylko Założycieli, §~\ref{art:breach} ust.~3); (b) odszkodowanie z art.~471 KC, trudne dowodowo; (c) roszczenie o zaniechanie i usunięcie skutków (§~\ref{art:breach} ust.~1--2). Nieosiągalne bez wyraźnej podstawy: zawieszenie prawa głosu, unieważnienie nabycia wpisanego do rejestru (art.~300$^{37}$ §~1 i art.~300$^{38}$ §~1 KSH), kara umowna (brak zastrzeżenia).}
\Rekomendacja{Zdefiniować „porozumienie” przez odesłanie do pisemnego albo ustnego porozumienia dotyczącego zgodnego głosowania, prowadzenia trwałej polityki wobec Spółki albo nabywania akcji; zastąpić „w interesie” zwrotem „na zlecenie albo na ryzyko ekonomiczne”; dodać obowiązek ponownego złożenia oświadczenia na żądanie Rady (nie częściej niż raz w roku) i obowiązek zawiadomienia o zawarciu porozumienia w 14 dni; rozważyć karę umowną w umowie akcjonariuszy (nie w umowie spółki).}
\Kierunek{Większa przewidywalność testu i narzędzie Rady na etapie trwania stosunku; dla inwestora okresowe oświadczenie to standard KYC; dla rejestru i sądu bez wpływu.
\begin{Brzmienie}
(§~\ref{art:dopuszczalny-nabywca} ust.~6, zdanie zastępujące zdanie o nabyciu na cudzy rachunek) „Kryterium Bezpieczeństwa EU/NATO nie spełnia osoba ani podmiot, który nabywa lub posiada akcje na rachunek, na zlecenie, na ryzyko ekonomiczne albo na polecenie osoby lub podmiotu niespełniającego tego Kryterium, w szczególności jako powiernik, pełnomocnik, zastawnik z prawem głosu albo strona porozumienia --- pisemnego lub ustnego --- dotyczącego zgodnego wykonywania prawa głosu, prowadzenia trwałej polityki wobec Spółki albo nabywania akcji. Akcjonariusz zawiadamia Radę Dyrektorów o zawarciu takiego porozumienia w terminie 14 dni i na żądanie Rady Dyrektorów, nie częściej niż raz w roku obrotowym, ponawia oświadczenie, o którym mowa w ust.~8.”
\end{Brzmienie}}
\end{MemoBox}

\item Postanowienia niniejszej umowy nakładające obowiązki na osobę, która nie jest jeszcze akcjonariuszem ani stroną niniejszej umowy, wiążą tę osobę od chwili skutecznego nabycia akcji albo przystąpienia do umowy. Do tego czasu wykonanie tych obowiązków, w szczególności przekazanie informacji niezbędnych do weryfikacji, zapewnia zbywca, a ich niewykonanie wstrzymuje udzielenie zgody Spółki na rozporządzenie akcją zgodnie z §~\ref{art:zbywanie}.
  \item Każda osoba zamierzająca nabyć prawa dotyczące akcji jest zobowiązana, na uzasadnione żądanie Rady Dyrektorów, przedstawić dokumenty i informacje niezbędne do przeprowadzenia weryfikacji prawnej dopuszczalności nabycia oraz spełnienia Kryterium Bezpieczeństwa EU/NATO, w tym oświadczenie, że nabywa akcje na własny rachunek i we własnym interesie.
  
\begin{MemoBox}{1.3.3}
\MemoPyt{1.3.3}{Obowiązki informacyjne osoby trzeciej przed nabyciem (§~11 ust.~8): skoro nabywca nie jest jeszcze stroną umowy, czy konstrukcja „warunku procedury” (brak informacji = wstrzymanie, nie odmowa) jest wystarczająca, a §~11 ust.~7 skutecznie wiąże nabywcę od chwili nabycia?}

\Ocena{zgodne}
\Klasa{bez zmian}
\Podstawa{art.~300$^{1}$ §~3, art.~300$^{37}$ §~1 i art.~300$^{38}$ §~1 KSH \Wiar{Z}; art.~393 KC (a contrario --- brak obciążenia osoby trzeciej) \Wiar{Z}; art.~300$^{39}$ §~1 KSH \Wiar{Z}.}
\Uzasadnienie{Konstrukcja jest poprawna: umowa spółki nie może nałożyć obowiązku na osobę trzecią, ale może uczynić przedstawienie informacji przesłanką zgody Spółki, co obciąża zbywcę (ust.~7 zd.~2), a ust.~8 ma wobec osoby trzeciej charakter jedynie informacyjny, co nie szkodzi, bo sankcją jest wstrzymanie, a nie roszczenie. Po nabyciu (z chwilą wpisu do rejestru --- art.~300$^{37}$ §~1 KSH) nabywca jest związany całą umową spółki z mocy prawa (art.~300$^{1}$ §~3 KSH), nie z mocy ust.~7, którego zdanie pierwsze jest deklaratoryjne i nieszkodliwe. Oświadczenie nabywcy w umowie zbycia (znajomość umowy spółki, prawdziwość informacji, przystąpienie do umowy wykonawczej w zakresie pełnomocnictw) to materia umowy inwestycyjnej i wzoru umowy zbycia, nie umowy spółki; przy zawiązaniu i pierwszej rundzie zbywcami będą Założyciele, więc praktyczne ryzyko jest niskie.}
\Rekomendacja{Bez zmian w umowie spółki. Przygotować wzór oświadczenia nabywcy (załącznik do regulaminu Rady albo do umowy o prowadzenie rejestru), aby Rada żądała jednolitego zestawu dokumentów.}
\end{MemoBox}

\item Weryfikację prawnej dopuszczalności nabycia przeprowadza Rada Dyrektorów na podstawie dokumentów przedstawionych przez nabywcę, publicznie dostępnych urzędowych wykazów oraz, w razie potrzeby, opinii prawnej. Rada Dyrektorów stwierdza wynik weryfikacji w formie uchwały lub protokołu. Uznanie nabywcy za Niedopuszczalnego Nabywcę wymaga wskazania konkretnej podstawy prawnej uniemożliwiającej planowane nabycie.
  \item Brak wystarczających informacji do zakończenia weryfikacji nie oznacza sam w sobie uznania osoby za Niedopuszczalnego Nabywcę, lecz wstrzymuje zakończenie transakcji do czasu uzyskania informacji koniecznych do ustalenia jej prawnej dopuszczalności.
  \item Uzyskanie statusu Prawnie Niedopuszczalnego Posiadacza nie powoduje samo przez się przymusowego zbycia akcji. Spółka i akcjonariusz stosują środki wymagane przez bezwzględnie obowiązujące przepisy prawa. Jeżeli właściwe przepisy dopuszczają usunięcie stanu prawnej niedopuszczalności przez zbycie akcji, akcjonariusz jest zobowiązany dokonać zbycia na rzecz Dopuszczalnego Nabywcy niezwłocznie po dniu, w którym takie zbycie stało się prawnie dopuszczalne. Do tego czasu wykonywanie praw, otrzymywanie świadczeń, rozporządzanie akcjami oraz inne czynności dotyczące akcji podlegają ograniczeniom wynikającym z właściwych przepisów prawa.
  
\begin{MemoBox}{1.7.3}
\MemoPyt{1.7.3}{Czy konstrukcja Prawnie Niedopuszczalnego Posiadacza (bez automatycznego przymusowego zbycia) jest spójna z sankcjami i kontrolą eksportu, i czy nie potrzebuje wyraźnego trybu zawieszenia wykonywania praw z akcji?}
\Ocena{ryzyko --- konstrukcja spójna, ale brak trybu zawieszenia i korekty progów głosowania może sparaliżować Spółkę}
\Klasa{konieczne (K9)}
\Podstawa{art.~1 lit.~e--g, art.~2 ust.~1--2, art.~4--7 i art.~9 rozporządzenia Rady (UE) nr~269/2014 \Wiar{Z} (tekst EN, pierwotny, bez konsolidacji); art.~1 pkt~2, art.~2, art.~6 ust.~1--2, art.~6a ust.~1, 11, 15--16 i~27, art.~6b ust.~1 i~4 oraz art.~6da ustawy z 13 kwietnia 2022 r.\ o szczególnych rozwiązaniach w zakresie przeciwdziałania wspieraniu agresji na Ukrainę oraz służących ochronie bezpieczeństwa narodowego \Wiar{Z}; art.~1 oraz art.~2 pkt~2 lit.~d i pkt~10 lit.~c rozporządzenia (UE) 2021/821 \Wiar{Z} (tekst EN, wersja skonsolidowana na 26.05.2023); art.~300$^{1}$ §~3, art.~300$^{33}$, art.~300$^{39}$ §~6 KSH \Wiar{Z}; art.~58 §~1 KC \Wiar{Z}}
\Uzasadnienie{Konstrukcja jest trafna: zamrożenie obejmuje akcje (art.~1 lit.~g pkt~iii, art.~2 ust.~1 rozp.~269/2014) \Wiar{Z} (tekst EN, pierwotny; załącznik~I wymaga sprawdzenia w wersji skonsolidowanej), więc automatyczne przymusowe zbycie byłoby zakazane (art.~58 §~1 KC), co ust.~11--12 słusznie wykluczają. Derogacje tekstu pierwotnego (art.~4--6) nie obejmują wprost zbycia zamrożonych akcji \Wiar{Z} (tekst EN); obecny stan trzeba sprawdzić w wersji skonsolidowanej \Wiar{?}. Wobec osób z listy krajowej ustawa z 2022~r.\ przewiduje fakultatywny zarząd przymusowy: zarządca wykonuje prawa z akcji poza zamrożeniem i może je zbyć ze skutkami sprzedaży egzekucyjnej (art.~6a ust.~11, 15--16 i~27, art.~6da) \Wiar{Z}; czy wtedy, jak w egzekucji (art.~300$^{39}$ §~6 KSH), zgoda Spółki jest zbędna, jest niepewne \Wiar{?} --- umowa tego nie wyłączy. Kontrola eksportu (rozp.~2021/821) dotyczy dostępu do technologii, nie posiadania akcji \Wiar{Z} (tekst EN, wersja skonsolidowana na 26.05.2023). Tryb na okres zamrożenia musi stworzyć umowa. Luki: \textbf{(1)}~dywidendy muszą trafić na rachunek zamrożony (art.~2 ust.~2, art.~7 ust.~1) \Wiar{Z} (tekst EN), a głosowanie zmieniające zasoby może naruszać zamrożenie \Wiar{W}; umowa nie pozbawi głosu, ale może ustalić tryb wstrzymania świadczeń i niedopuszczenia do praw „w zakresie wynikającym z właściwych przepisów”; głosy zarządcy trzeba liczyć; \textbf{(2)}~progi „75\% wszystkich głosów” liczą głosy zamrożone --- posiadacz ponad 25\% blokuje każdą Sprawę Zastrzeżoną; należy je wyłączyć z mianownika, lecz tylko przy większościach umownych.}
\Rekomendacja{Dodać w §~\ref{art:dopuszczalny-nabywca} tryb stwierdzenia statusu uchwałą Rady, wstrzymania świadczeń i niedopuszczenia do wykonywania praw w zakresie wynikającym z prawa, oraz regułę wyłączenia głosów zawieszonych z mianownika większości umownych, z zastrzeżeniem wykonywania praw przez zarządcę przymusowego.}
\begin{Brzmienie}
§~\ref{art:dopuszczalny-nabywca} --- nowy ust.~11a: „Status Prawnie Niedopuszczalnego Posiadacza oraz jego ustanie stwierdza Rada Dyrektorów uchwałą wskazującą podstawę prawną i zakres ograniczeń, doręczaną akcjonariuszowi i ujawnianą w rejestrze akcjonariuszy w zakresie przewidzianym przez KSH. Od dnia powstania podstawy prawnej Spółka wstrzymuje wypłatę dywidendy i innych świadczeń z akcji, przekazując je w sposób wymagany przez właściwe przepisy, oraz nie dopuszcza do wykonywania praw z akcji w zakresie, w jakim ich wykonywanie jest zakazane albo wymaga zezwolenia właściwego organu; przewodniczący Walnego Zgromadzenia odnotowuje to w protokole. Głosów z akcji, których wykonywanie jest w ten sposób wstrzymane, nie uwzględnia się przy obliczaniu większości »wszystkich głosów« i Głosów Uprawnionych w Sprawie przewidzianych w niniejszej umowie, w granicach dopuszczalnych przez bezwzględnie obowiązujące przepisy prawa. Postanowienia zdania drugiego i trzeciego nie ograniczają wykonywania praw z akcji przez zarządcę ustanowionego na podstawie przepisów szczególnych, w zakresie jego umocowania, a głosy przez niego oddane uwzględnia się przy obliczaniu większości.”
\end{Brzmienie}
\end{MemoBox}

\item Żaden organ Spółki nie może uznać Niedopuszczalnego Nabywcy za Dopuszczalnego Nabywcę, usunąć statusu Prawnie Niedopuszczalnego Posiadacza bez ustania, uchylenia albo stwierdzenia braku jego podstawy prawnej ani ustanowić wyjątku prowadzącego do czynności zakazanej przez bezwzględnie obowiązujące przepisy prawa. Żadne postanowienie niniejszej umowy nie upoważnia do dokonania zbycia, zapłaty, wypłaty świadczenia ani wykonania prawa z akcji w czasie, w którym dana czynność jest prawnie zakazana.
  \item Zasady określone w niniejszym paragrafie stosuje się odpowiednio do zastawu, użytkowania, powiernictwa, opcji, prawa do głosu, prawa do korzyści ekonomicznych, instrumentów zamiennych oraz innych konstrukcji, jeżeli ich ustanowienie, dalsze utrzymywanie albo wykonywanie na rzecz danej osoby lub podmiotu jest objęte bezwzględnie obowiązującym zakazem lub ograniczeniem prawnym albo prowadziłoby do obejścia Kryterium Bezpieczeństwa EU/NATO.
  \item \BoilerplateStrike{„Kwalifikowanym Inwestorem Finansowym” jest fundusz inwestycyjny, alternatywny fundusz inwestycyjny (w tym alternatywna spółka inwestycyjna), spółka zarządzająca aktywami albo inny podmiot zbiorowego inwestowania, jeżeli łącznie: (a)~jego siedziba, podmiot nim zarządzający (w szczególności komplementariusz albo zarządzający alternatywnym funduszem inwestycyjnym) oraz rzeczywisty ośrodek podejmowania decyzji inwestycyjnych znajdują się w państwie, o którym mowa w ust.~6 lit.~a; (b)~nie jest bezpośrednio ani pośrednio kontrolowany przez Niedopuszczalnego Nabywcę ani przez państwo, organ państwowy lub podmiot kontrolowany przez państwo spoza państw, o których mowa w ust.~6 lit.~a; (c)~żaden z jego inwestorów niespełniających kryterium określonego w ust.~6 lit.~a nie sprawuje kontroli nad decyzjami inwestycyjnymi tego podmiotu. Kwalifikowany Inwestor Finansowy spełnia Kryterium Bezpieczeństwa EU/NATO także wtedy, gdy część jego inwestorów nie spełnia kryterium określonego w ust.~6 lit.~a; w tym zakresie nie stosuje się wymogu z ust.~6 lit.~b dotyczącego beneficjentów rzeczywistych. Spełnienie warunków określonych w niniejszym ustępie stwierdza Rada Dyrektorów w ramach weryfikacji, o której mowa w ust.~9, a w przypadku Kwalifikowanej Rundy Finansowania rozstrzyga o nim uchwała kierunkowa, o której mowa w §~\ref{art:kwalifikowana-runda} ust.~2, przy określaniu dopuszczalnego kręgu inwestorów; odrębna uchwała Walnego Zgromadzenia przewidziana w ust.~6 nie jest wówczas wymagana dla poszczególnej transakcji.}
  
\begin{MemoBox}{1.4.1}
\MemoPyt{1.4.1}{Czy zachowanie wyjątku uprościłoby obsługę zbyć wtórnych i wejścia funduszy (w tym funduszy z inwestorami spoza UE/NATO) bez naruszenia wymogów właścicielskich programów finansowania obronnego i przepisów o kontroli inwestycji?}

\Ocena{zgodne warunkowo}
\Klasa{zalecane}
\Podstawa{art.~2 pkt~6 i~24, art.~5 oraz art.~9 ust.~3--4 i~7 rozporządzenia 2021/697 (EDF: kontrola przez państwo trzecie) \Wiar{Z}; art.~2 pkt~1, 5, 7 i~8, art.~4 ust.~15, art.~5 ust.~1 lit.~a, art.~15 ust.~1 lit.~a--b i ust.~3, art.~19 ust.~2 lit.~e oraz motyw~18 rozporządzenia (UE) 2026/1386 (od 17 stycznia 2028~r.) \Wiar{Z} (tekst EN); warunki NATO DIANA (siedziba w państwie NATO) \Wiar{W}; NATO Innovation Fund (wspierany przez 24 państwa NATO, inwestuje w spółki z siedzibą w jednym z tych państw) \Wiar{Z}; art.~3 ust.~1 pkt~1, art.~12c ust.~6 i art.~12e ust.~4 ustawy o kontroli niektórych inwestycji (podmiot dominujący, nabycie pośrednie, spółki zależne podmiotów z państw trzecich) \Wiar{Z}.}
\Uzasadnienie{Tak. Bez wyjątku każdy fundusz VC z choćby jednym LP spoza UE/EOG/NATO nie spełnia lit.~b Kryterium (beneficjenci rzeczywiści) i potrzebuje uchwały Walnego Zgromadzenia 75\% do każdej transakcji, odbieranej przez fundusze jako sygnał ryzyka \Wiar{W}. Wyjątek nie narusza wymogów programów. EDF bada \emph{kontrolę} (decydujący wpływ, także pośredni) przez niestowarzyszone państwo trzecie albo podmiot z takiego państwa (art.~9 ust.~3, art.~2 pkt~6); fundusz z siedzibą i strukturą zarządczą w Unii nie jest takim podmiotem (art.~2 pkt~24) \Wiar{Z}, a pasywny LP nie ma decydującego wpływu --- tak w praktyce Komisja \Wiar{W}. Gdyby kontrolę uzyskał LP spoza państw stowarzyszonych (EFTA w EOG, art.~5 --- w praktyce Norwegia), odbiorca nie traci automatycznie kwalifikowalności, lecz potrzebuje gwarancji (art.~9 ust.~4) i zgłasza zmianę Komisji (art.~9 ust.~7) \Wiar{Z} --- wyjątek powinien przewidywać obowiązek zawiadomienia. Od 17 stycznia 2028~r. fundusz unijny z pasywnymi LP spoza Unii nie jest spółką zależną inwestora zagranicznego (rozporządzenie 2026/1386, art.~2 pkt~1 i~7, motyw~18), ale państwo może żądać informacji o łańcuchu kontroli (art.~15 ust.~3), a nieprzejrzysta struktura własności jest czynnikiem ryzyka (art.~2 pkt~8, art.~19 ust.~2 lit.~e) \Wiar{Z} (tekst EN). W ustawie o kontroli niektórych inwestycji podmiotem dominującym w funduszu jest zarządzający, nie LP bez uprawnień kontrolnych (art.~3 ust.~1 pkt~1, art.~12c ust.~6) \Wiar{Z}; do progu 10~mln euro przychodów ustawa Spółki nie obejmuje (zob.~1.1.2). Warunek: wyjątek nie obejmie wehikułu jednego inwestora spoza Kryterium (1.4.3), a Kryterium co do kontroli nad Spółką pozostaje nienaruszone.}
\Rekomendacja{Przywrócić ust.~14 w brzmieniu uszczelnionym (1.4.3) wraz z odesłaniem w §~\ref{art:kwalifikowana-runda} ust.~1 i wyłączeniem w §~\ref{art:dopuszczalny-nabywca} ust.~15 zdarzeń po stronie pasywnych inwestorów funduszu. Decyzja należy do Założycieli; jeżeli zdecydują o usunięciu wyjątku, umowa pozostaje poprawna, a koszt przenosi się na każdą rundę (uchwała 75\%).}
\Kierunek{Przywrócenie ustępu zmienia sposób weryfikacji funduszu (test zarządzającego i kontroli zamiast testu każdego LP), nie zmienia bramki kontroli nad Spółką. Dla Spółki: mniejszy koszt każdej rundy; dla Założycieli: brak konieczności uchwały 75\% dla typowego funduszu; dla rejestru i sądu: bez wpływu.}
\end{MemoBox}

\begin{MemoBox}{1.4.2}
\MemoPyt{1.4.2}{Czy odstąpienie od badania beneficjentów rzeczywistych funduszu jest do pogodzenia z przepisami AML i wymogami zamawiających z sektora obronnego?}

\Ocena{zgodne warunkowo}
\Klasa{zalecane}
\Podstawa{art.~2 ust.~2 pkt~1 lit.~a i art.~58 pkt~4a ustawy z 1 marca 2018 r. o przeciwdziałaniu praniu pieniędzy oraz finansowaniu terroryzmu \Wiar{Z}; rozporządzenie (UE) 2024/1624 (pakiet AML, stosowane od 2027 r.) \Wiar{?}; art.~1 lit.~d--g, art.~2 ust.~1--2, art.~3 ust.~1 i art.~9 oraz załącznik~I rozporządzenia (UE) nr~269/2014 (test sankcyjny) \Wiar{Z} (tekst EN, pierwotny, bez konsolidacji); art.~5 ust.~1 lit.~b oraz art.~19 ust.~2 lit.~a pkt~vi i lit.~b rozporządzenia (UE) 2026/1386 \Wiar{Z} (tekst EN); art.~57 ust.~2 pkt~1, art.~58 ust.~2 pkt~2 i art.~64 ust.~2 pkt~2 ustawy z 5 sierpnia 2010 r. o ochronie informacji niejawnych \Wiar{Z}; art.~10 ust.~1 pkt~2 ustawy z 13 czerwca 2019 r. (koncesja) \Wiar{Z}.}
\Uzasadnienie{AML nie nakłada na Spółkę obowiązku badania wszystkich inwestorów funduszu-akcjonariusza: Spółka nie jest instytucją obowiązaną z tytułu samego przyjęcia inwestora \Wiar{W}; zgłasza do CRBR własnych beneficjentów rzeczywistych (art.~58 pkt~4a), tj.~osoby z ponad 25\% akcji lub głosów, także pośrednio, albo sprawujące kontrolę (art.~2 ust.~2 pkt~1 lit.~a) \Wiar{Z}. Przekreślony ust.~14 jest więc \emph{ostrzejszy} niż AML co do kontroli, a łagodniejszy tylko co do pasywnych LP poniżej progów, których AML też nie uznaje za beneficjentów. Zgodność jest zachowana pod warunkiem, że wyjątek nie zwalnia z ustalenia beneficjentów w rozumieniu AML ani z testu sankcyjnego wobec nich; zwolnienie z „wymogu lit.~b dotyczącego beneficjentów” trzeba zawęzić. Test sankcyjny nie może być ograniczony do samego funduszu: rozporządzenie 269/2014 zakazuje udostępniania środków (w tym akcji) „bezpośrednio lub pośrednio” osobom z załącznika~I (tłum. robocze) i świadomego udziału w obchodzeniu tych środków (art.~1 lit.~g, art.~2 ust.~1--2, art.~9) \Wiar{Z} (tekst EN, pierwotny, bez konsolidacji; załącznik~I trzeba sprawdzać w wersji skonsolidowanej). Od 17 stycznia 2028~r. sankcje wobec inwestora, osoby go kontrolującej, beneficjenta rzeczywistego albo spółki zależnej uruchamiają też notyfikację (art.~5 ust.~1 lit.~b rozporządzenia 2026/1386), jeżeli cel jest objęty art.~4 ust.~15 \Wiar{Z} (tekst EN). Postępowania bezpieczeństwa przemysłowego i koncesyjne badają strukturę kapitału, źródła środków i akcjonariuszy \Wiar{Z}; organ może żądać informacji o inwestorach funduszu bezpośrednio, czego umowa nie zablokuje --- wyjątek powinien zawierać zobowiązanie funduszu do ich udzielenia.}
\Rekomendacja{W przywróconym ust.~14 zachować: obowiązek ujawnienia beneficjentów rzeczywistych w rozumieniu przepisów AML, test sankcyjny wobec nich oraz zobowiązanie do współdziałania z organami koncesyjnymi i bezpieczeństwa przemysłowego.}
\Kierunek{Uzupełnienia są warunkiem, pod którym wyjątek nie osłabia zgodności regulacyjnej; brzmienie w 1.4.3 (s.~\pageref{memo:1.4.3}).}
\end{MemoBox}

\begin{MemoBox}{1.4.3}
\MemoPyt{1.4.3}{Jeżeli wyjątek ma zostać, jak go sformułować, aby nie stał się furtką dla podmiotów spoza Kryterium (np.~wehikuły jednego inwestora)?}

\Ocena{zgodne warunkowo}
\Klasa{zalecane}
\Podstawa{ustawa o funduszach inwestycyjnych i zarządzaniu AFI (ZASI, rejestr KNF) \Wiar{W}; dyrektywa 2011/61/UE \Wiar{W}; art.~2 ust.~2 pkt~1 lit.~a ustawy AML (beneficjent rzeczywisty) \Wiar{Z}; art.~2 pkt~1, 7 i~8 oraz art.~19 ust.~2 lit.~e rozporządzenia (UE) 2026/1386 \Wiar{Z} (tekst EN); art.~9 ust.~7 rozporządzenia 2021/697 \Wiar{Z}; strony programów PFR Ventures (Starter, Biznest, Otwarte Innowacje, KOFFI) \Wiar{Z}.}
\Uzasadnienie{Furtki w przekreślonym brzmieniu są trzy: brak wymogu rozproszenia (fundusz z jednym LP spoza Kryterium, zarządzany formalnie z UE, spełniałby test), brak wymogu statusu regulacyjnego zarządzającego oraz zbyt szerokie zwolnienie z lit.~b, obejmujące także beneficjentów kontrolujących. Progi (5 inwestorów, 25\%, 50\%, 10\%) są propozycją do decyzji Założycieli \Wiar{W}. Strony PFR Ventures nie publikują takich progów, tylko strukturę kapitału (minimalny wkład prywatny); w KOFFI „posiadanie jednego LP nie jest formalnie zabronione”, lecz fundusze „powinny mieć zdywersyfikowaną strukturę LPs”, a rejestracja ZASI jest wymagana najpóźniej przy umowie z PFR \Wiar{Z}. Fundusz PFR z zarejestrowanym ZASI i inwestorami spełniającymi Kryterium spełnia je wprost i wyjątku nie potrzebuje; w modelach koinwestycyjnych (OI, Biznest) kapitał funduszu pochodzi w zasadzie od PFR i zarządzającego, więc rozproszenia by nie było, ale koinwestorzy obejmują akcje bezpośrednio i przechodzą test samodzielnie. Strona OI jest wewnętrznie niespójna co do limitu geograficznego (opis: do 40\% portfela w spółki z zagraniczną siedzibą i polskimi founderami; FAQ: do 15\% za granicą) \Wiar{Z} co do brzmienia strony; limit nie dotyczy Spółki z siedzibą w Polsce, a Term Sheetu (zał.~nr~7 do Zasad Naboru) nie ma w materiałach. Wymóg rozproszenia odpowiada definicji spółki zależnej inwestora zagranicznego (art.~2 pkt~1 i~7 rozporządzenia 2026/1386, od 17 stycznia 2028~r.), a ujawnienie i zawiadamianie --- kryterium „nieprzejrzystej struktury własności” (art.~2 pkt~8, art.~19 ust.~2 lit.~e) \Wiar{Z} (tekst EN); zawiadamianie pozwala Spółce informować Komisję o zmianach (art.~9 ust.~7 rozporządzenia 2021/697) \Wiar{Z}.}
\Rekomendacja{Przywrócić ust.~14 w brzmieniu uszczelnionym; przywrócić odesłanie w §~\ref{art:kwalifikowana-runda} ust.~1; dodać w ust.~15 wyłączenie zdarzeń po stronie pasywnych inwestorów.}
\Kierunek{Brzmienie jest gotowe do wklejenia (pełny tekst --- w memorandum.md); jeżeli Założyciele zdecydują o usunięciu wyjątku, żadna inna zmiana nie jest potrzebna. Streszczenie proponowanego ust.~14: „Kwalifikowanym Inwestorem Finansowym” jest fundusz inwestycyjny, alternatywny fundusz inwestycyjny (w tym alternatywna spółka inwestycyjna) albo inny podmiot zbiorowego inwestowania, jeżeli łącznie:
\begin{itemize}
\item[(a)] jego siedziba, rzeczywisty ośrodek decyzji inwestycyjnych i podmiot zarządzający są w państwie z ust.~6 lit.~a, a zarządzający jest zarejestrowany albo licencjonowany przez organ nadzoru nad rynkiem finansowym takiego państwa i sam spełnia Kryterium, także co do beneficjentów rzeczywistych;
\item[(b)] nie jest bezpośrednio ani pośrednio kontrolowany przez Niedopuszczalnego Nabywcę ani przez państwo spoza państw z ust.~6 lit.~a, jego organ lub podmiot przez nie kontrolowany;
\item[(c)] ma co najmniej pięciu inwestorów niepowiązanych; żaden inwestor niespełniający ust.~6 lit.~a nie ma więcej niż 25\% zobowiązań kapitałowych ani uprawnień do współdecydowania o inwestycjach, a łącznie tacy inwestorzy --- nie więcej niż 50\% zobowiązań kapitałowych;
\item[(d)] nie został utworzony w celu nabycia akcji Spółki ani nie jest wehikułem współinwestycyjnym jednego lub dwóch inwestorów;
\item[(e)] ujawnił Radzie beneficjentów rzeczywistych w rozumieniu przepisów AML i każdego inwestora z co najmniej 10\% zobowiązań kapitałowych, a żadna z tych osób nie jest Niedopuszczalnym Nabywcą.
\end{itemize}
Kwalifikowany Inwestor Finansowy spełnia Kryterium także wtedy, gdy część jego inwestorów nie spełnia ust.~6 lit.~a; wymogu z ust.~6 lit.~b co do beneficjentów rzeczywistych nie stosuje się wyłącznie do inwestorów, którzy nie sprawują kontroli ani nie przekraczają progów z lit.~c. Fundusz składa oświadczenie, zawiadamia Radę o zmianach w 14 dni i zobowiązuje się udzielać na żądanie organom właściwym w sprawach koncesji, bezpieczeństwa przemysłowego i kontroli inwestycji informacji o swoich inwestorach; zmiana po stronie inwestora poniżej progów z lit.~c nie jest utratą Kryterium w rozumieniu ust.~15. Spełnienie warunków stwierdza Rada w weryfikacji z ust.~9, a w Kwalifikowanej Rundzie Finansowania --- uchwała kierunkowa (§~\ref{art:kwalifikowana-runda} ust.~2), bez odrębnej uchwały z ust.~6; Walne Zgromadzenie może większością 75\% wszystkich głosów stwierdzić, że podmiot nie spełnia lub przestał spełniać te warunki.}
\end{MemoBox}

\begin{MemoBox}{4.4.2}
\MemoPytCd{4.4.2}
\Kierunek{(I)~§~\ref{art:dopuszczalny-nabywca} ust.~14 --- przywrócić w~brzmieniu uszczelnionym (spójnie z~pytaniem 1.4.3): test na poziomie zarządzającego, siedziby i~kontroli zamiast look-through do LP. „Kwalifikowany Inwestor Finansowy” (fundusz, AFI/ASI, spółka zarządzająca aktywami, inny podmiot zbiorowego inwestowania) --- jeżeli łącznie: (a)~jego siedziba, siedziba i~rzeczywisty ośrodek zarządzania zarządzającego oraz ośrodek decyzji inwestycyjnych są w~państwie z~ust.~6 lit.~a; (b)~ani on, ani zarządzający nie jest bezpośrednio ani pośrednio kontrolowany przez Niedopuszczalnego Nabywcę, osobę spoza kryterium z~ust.~6 lit.~a ani państwo (podmiot państwowy) spoza tych państw; (c)~żaden inwestor spoza tego kryterium nie ma (sam lub z~podmiotami powiązanymi) ponad 25\% zadeklarowanego kapitału ani kontroli nad decyzjami inwestycyjnymi (weto, komitet inwestycyjny, porozumienie z~zarządzającym); (d)~nie jest wehikułem jednego inwestora, chyba że ten spełnia Kryterium; (e)~złożył pisemne oświadczenie z~dokumentami (ust.~8) i~zobowiązał się zawiadamiać Radę o~zmianie w~14 dni. Taki podmiot spełnia Kryterium mimo części inwestorów spoza kryterium --- bez wymogu co do beneficjentów rzeczywistych z~ust.~6 lit.~b; weryfikacja, czy nie jest Niedopuszczalnym Nabywcą (ust.~1--2 i~9), i~obowiązki AML pozostają. Spełnienie warunków stwierdza Rada w~weryfikacji z~ust.~9, a~w~Kwalifikowanej Rundzie --- uchwała kierunkowa z~§~\ref{art:kwalifikowana-runda} ust.~2 (bez odrębnej uchwały z~ust.~6); podmiot tego samego lub powiązanego zarządzającego spełniający lit.~a--e --- bez odrębnej uchwały kierunkowej. Konsekwentnie przywrócić przekreślone odesłanie w~§~\ref{art:kwalifikowana-runda} ust.~1.}
\end{MemoBox}

\item Akcjonariusz, który po nabyciu akcji przestał spełniać Kryterium Bezpieczeństwa EU/NATO, w szczególności wskutek zmiany obywatelstwa lub miejsca zamieszkania, zmiany kontroli nad akcjonariuszem będącym osobą prawną albo jednostką organizacyjną, zmiany jego beneficjentów rzeczywistych albo objęcia akcji na rachunek osoby niespełniającej Kryterium, zawiadamia o tym Radę Dyrektorów w terminie 14 dni od dnia zdarzenia. W terminie 6 miesięcy od dnia zdarzenia akcjonariusz zbywa akcje na rzecz Dopuszczalnego Nabywcy spełniającego Kryterium w trybie §~\ref{art:zbywanie}, chyba że Walne Zgromadzenie większością 75\% wszystkich głosów wyrazi zgodę na dalsze posiadanie akcji. Po bezskutecznym upływie tego terminu Spółka może wskazać Dopuszczalnego Nabywcę spełniającego Kryterium, na rzecz którego akcjonariusz jest zobowiązany zbyć akcje za cenę równą Wartości Godziwej ustalonej zgodnie z §~\ref{art:wycena}, płatną na zasadach określonych w §~\ref{art:zbywanie} ust.~3; obowiązek ten jest obowiązkiem związanym z akcją wobec Spółki, a do jego wykonania stosuje się odpowiednio §~\ref{art:zwrotne-zbycie} ust.~14--15 oraz §~\ref{art:breach}. Niewykonanie obowiązków określonych w niniejszym ustępie stanowi istotne naruszenie niniejszej umowy.

\begin{MemoBox}{1.1.4}
\MemoPyt{1.1.4}{Utrata Kryterium po nabyciu (§~11 ust.~15): czy obowiązek zbycia akcji w 6 miesięcy i prawo Spółki do wskazania nabywcy są dopuszczalne i skuteczne wobec nabywców wtórnych jako obowiązek „związany z akcją”; czy nie grozi zakwalifikowanie tej konstrukcji jako przymusowego wykupu poza trybem ustawowym?}
\Ocena{zgodne warunkowo}
\Klasa{zalecane}
\Podstawa{art.~300$^{33}$ §~1 pkt~11 KSH (obowiązki związane z akcją w rejestrze) \Wiar{Z}; art.~300$^{39}$ §~1 KSH \Wiar{Z}; art.~300$^{45}$ §~2, art.~300$^{47}$ §~9 oraz art.~300$^{49}$--300$^{50}$ KSH \Wiar{Z}; art.~64 KC, art.~1047 KPC \Wiar{Z}; SN II CSKP 593/22 \Wiar{Z}.}
\Uzasadnienie{Obowiązek zbycia na rzecz osoby trzeciej za Wartość Godziwą nie jest umorzeniem (akcje nie są unicestwiane, kapitał akcyjny nie jest obniżany), nabyciem akcji własnych (nabywcą nie jest Spółka ani osoba działająca na jej rachunek --- projekt powinien to powtórzyć wprost, jak w §~\ref{art:zwrotne-zbycie} ust.~11--12) ani przymusowym wykupem: przepisy o P.S.A. nie zawierają odpowiednika art.~418 KSH, a jedynymi ustawowymi trybami przymusowej zmiany składu akcjonariatu są sądowe wyłączenie akcjonariusza (art.~300$^{49}$) i ustąpienie z wykupem „po cenie odpowiadającej wartości godziwej” (art.~300$^{50}$ §~3) \Wiar{Z}. Jest to umowny obowiązek związany z akcją (art.~300$^{1}$ §~3 w zw. z art.~300$^{39}$ §~1 KSH), ujawniany w rejestrze (art.~300$^{33}$ §~1 pkt~11) i skuteczny wobec nabywcy wtórnego, który wstępuje w stosunek spółki ukształtowany umową; ujawnienie wyłącza zarzut nieznajomości. Ryzyko kwalifikacji jako obejścia przepisów o umorzeniu jest niskie, bo cena to pełna Wartość Godziwa bez dyskonta --- art.~300$^{45}$ §~2 KSH nie jest obchodzony nawet przy stosowaniu per analogiam. Cztery luki: (1) brak terminu na wskazanie nabywcy przez Spółkę po 6 miesiącach i skutku braku wskazania --- „istotne naruszenie” trwa bez końca; (2) samodzielne zbycie w trybie §~\ref{art:zbywanie} (zgoda, prawo pierwszeństwa) może zająć większość terminu; (3) u inwestora instytucjonalnego utratę Kryterium może wywołać zmiana jego inwestora (LP), na którą nie ma wpływu \Wiar{W}; (4) pełnomocnictwo z umowy wykonawczej (§~\ref{art:breach} ust.~3) dotyczy Założycieli, więc wobec inwestora egzekucja wymaga procesu z art.~64 KC.}
\Rekomendacja{Uzupełnić ust.~15 o: termin 3 miesięcy dla Spółki na wskazanie nabywcy, a po jego bezskutecznym upływie --- utrzymanie obowiązku z zawieszeniem prawa do wskazania do czasu kolejnego wezwania; wyłączenie z „utraty Kryterium” zdarzeń po stronie inwestorów pasywnych funduszu, jeżeli Założyciele przywrócą ust.~14 (zob.~1.4.3); zdanie, że nabywcą nie może być Spółka ani osoba działająca na jej rachunek poza trybem art.~300$^{47}$ KSH; wymóg przystąpienia każdego inwestora do umowy wykonawczej w zakresie pełnomocnictwa (w umowie inwestycyjnej).}
\Kierunek{Zmiany usuwają stan zawieszenia po 6 miesiącach i czynią mechanizm przewidywalnym dla inwestora; w rejestrze adnotacja o obowiązku związanym z akcją (zob.~1.1.6); dla sądu rejestrowego bez wpływu.
\begin{Brzmienie}
(§~\ref{art:dopuszczalny-nabywca} ust.~15, zdanie trzecie w nowym brzmieniu) „Po bezskutecznym upływie tego terminu Spółka może, w terminie 3 miesięcy, wskazać Dopuszczalnego Nabywcę spełniającego Kryterium, na rzecz którego akcjonariusz jest zobowiązany zbyć akcje za cenę równą Wartości Godziwej ustalonej zgodnie z §~\ref{art:wycena} na dzień zdarzenia, płatną na zasadach określonych w §~\ref{art:zbywanie} ust.~3; nabywcą nie może być Spółka ani osoba działająca na jej rachunek, chyba że nabycie następuje w trybie zgodnym z art.~300$^{47}$ KSH. Jeżeli Spółka nie wskaże nabywcy w tym terminie, obowiązek zbycia trwa nadal, a Spółka może ponowić wskazanie po kolejnym wezwaniu z terminem 3 miesięcy.”
\end{Brzmienie}}
\end{MemoBox}

\begin{MemoBox}{4.1.1}
\MemoPytCd{4.1.1}
\Kierunek{Sygnał dla pakietu A1 (opcjonalne). Warunki niekaralności i~braku postępowania dotyczą także akcjonariuszy posiadających co najmniej 20\% akcji, a~wniosek o~koncesję zawiera ich dane osobowe i~zaświadczenia o~niekaralności (art.~10 ust.~1 pkt~2, art.~17 ustawy z~13 czerwca 2019~r. \Wiar{Z}); inwestor z~takim pakietem będzie badany przez organ koncesyjny. Rozważyć w~ust.~15 obowiązek akcjonariusza dostarczenia dokumentów wymaganych w~postępowaniu koncesyjnym --- dziś ust.~15 nie przewiduje tego wprost.}
\end{MemoBox}

\begin{MemoBox}{4.4.2}
\MemoPytCd{4.4.2}
\Kierunek{(II)~§~\ref{art:dopuszczalny-nabywca} ust.~15 --- zawęzić do zdarzeń po stronie funduszu, nie jego LP. Zawężenie nie zwalnia Spółki z~obowiązku informowania Komisji o~zmianie, „która może budzić wątpliwości w~zakresie spełniania kryteriów kwalifikowalności”, w~trakcie działania EDF (art.~9 ust.~7 rozporządzenia 2021/697 \Wiar{Z}); dlatego zawiadomienie funduszu o~zmianach z~brzmienia (I) musi pozostać. Nowe zdanie drugie (po zdaniu pierwszym, przed „W terminie 6 miesięcy”):}
\begin{Brzmienie}
Zmiana w składzie inwestorów Kwalifikowanego Inwestora Finansowego albo zmiana jego beneficjentów rzeczywistych nie stanowi utraty Kryterium Bezpieczeństwa EU/NATO, dopóki warunki określone w ust.~14 lit.~a--d pozostają spełnione; utratę Kryterium przez Kwalifikowanego Inwestora Finansowego stanowi w szczególności zmiana kontroli nad podmiotem nim zarządzającym, przeniesienie zarządzania na podmiot niespełniający ust.~14 lit.~a--b albo uzyskanie przez inwestora niespełniającego kryterium z ust.~6 lit.~a udziału lub uprawnień, o których mowa w ust.~14 lit.~c.
\end{Brzmienie}
\end{MemoBox}

\end{enumerate}
\end{ContractClause}





\begin{ContractClause}{Zbywanie i obciążanie akcji}
\label{art:zbywanie}
\begin{enumerate}
  \item Zbycie albo obciążenie akcji (rozporządzenie akcją) wymaga zgody Spółki w rozumieniu art.~300$^{39}$ §~1 KSH. Zgody udziela albo odmawia jej Rada Dyrektorów w formie dokumentowej pod rygorem nieważności, w terminie 14 dni od zgłoszenia Spółce zamiaru rozporządzenia. Zgłoszenie zawiera dane planowanego nabywcy, cenę i istotne warunki rozporządzenia oraz informacje niezbędne do weryfikacji.
  
\begin{MemoBox}{1.3.1}
\MemoPyt{1.3.1}{Czy 14-dniowy termin, forma dokumentowa i zamknięty katalog przyczyn odmowy są zgodne z art.~300$^{39}$ §~3 KSH i czy projekt powinien określać skutek milczenia Rady (zgoda dorozumiana albo bieg terminu na wskazanie nabywcy)?}

\Ocena{zgodne warunkowo}
\Klasa{zalecane}
\Podstawa{art.~300$^{39}$ §~3--5 i art.~4 §~2$^{1}$ KSH \Wiar{Z}; art.~300$^{40}$ §~2 KSH (14-dniowy termin ustawowy dla zgody na zbycie akcji nie w pełni pokrytej) \Wiar{Z}; art.~77$^{2}$ KC (forma dokumentowa) \Wiar{Z}; art.~60 KC \Wiar{Z}; art.~300$^{58}$ KSH (uchwały Rady) \Wiar{Z}.}
\Uzasadnienie{Forma dokumentowa pod rygorem nieważności odpowiada art.~300$^{39}$ §~4 KSH (odmowy zgody i wskazania nabywcy dokonuje zarząd w tej formie, „chyba że umowa spółki stanowi inaczej”; w Spółce --- Rada Dyrektorów, art.~4 §~2$^{1}$ KSH) \Wiar{Z}; jej zachowanie w umowie jest wskazane niezależnie od derogacji. Termin 14 dni jest krótszy niż ustawowy miesiąc na wskazanie nabywcy i dotyczy innej czynności (stanowiska w sprawie zgody), więc nie koliduje z ustawą; ten sam termin ustawa przyjmuje w art.~300$^{40}$ §~2 KSH; jest wykonalny, jeżeli Rada obraduje elektronicznie (§~\ref{art:rada} ust.~10). Zamknięty katalog przyczyn odmowy jest dopuszczalny i korzystny: ustawa nie wymaga uzasadnienia odmowy, ale umowa może związać Radę kryteriami obiektywnymi. Luką jest brak skutku milczenia Rady: akcjonariusz nie może wpisać nabywcy (brak zgody), a nie wie, czy biegnie termin na wskazanie nabywcy. Zgoda dorozumiana jest wykluczona --- byłaby nieważna z braku formy dokumentowej. Właściwy skutek to fikcja odmowy z przyczyny lit.~b z upływem 14 dni, uruchamiająca ust.~3 (w wariancie~A termin 3 miesięcy liczy się od zgłoszenia, więc milczenie skraca Radzie jej własny czas na wskazanie nabywcy); jest korzystna dla akcjonariusza, a dla Spółki bezpieczna (nie dopuści przez milczenie nabywcy spoza Kryterium). Fikcja nie działa w czasie wstrzymania terminu (lit.~d i~e).}
\Rekomendacja{Dodać skutek milczenia jako fikcję odmowy z przyczyny lit.~b, z zastrzeżeniem wstrzymania terminu.}
\Kierunek{Zmiana usuwa stan zawieszenia; dyscyplinuje Radę; nie zmienia pozycji rejestru (nadal potrzebna zgoda albo zaświadczenie z 1.2.2).}
\begin{Brzmienie}
(§~\ref{art:zbywanie} ust.~1, zdanie dodawane na końcu) „Jeżeli Rada Dyrektorów nie udzieli zgody ani nie odmówi jej w tym terminie, a bieg terminu nie został wstrzymany na podstawie ust.~2 lit.~d albo~e, uważa się, że Spółka odmówiła zgody z przyczyny określonej w ust.~2 lit.~b z upływem tego terminu; stosuje się wówczas ust.~3.”
\end{Brzmienie}
\end{MemoBox}

\item Rada Dyrektorów może odmówić zgody wyłącznie, jeżeli:
  \begin{enumerate}[label=\alph*)]
    \item planowany nabywca jest Niedopuszczalnym Nabywcą;
    \item planowany nabywca nie spełnia Kryterium Bezpieczeństwa EU/NATO, a Walne Zgromadzenie nie wyraziło zgody zgodnie z §~\ref{art:dopuszczalny-nabywca};
    \item rozporządzenie następowałoby w Okresie Ograniczenia Zbywania Akcji bez zgody wymaganej przez §~\ref{art:okres-ograniczenia-zbywania};
    \item mimo wezwania nie przedstawiono informacji lub oświadczeń niezbędnych do zakończenia weryfikacji --- w tym przypadku bieg terminu z ust.~1 ulega wstrzymaniu do czasu ich przedstawienia, a wstrzymanie nie stanowi odmowy; albo
    \item nabycie wymaga zgody, braku sprzeciwu albo upływu terminu na sprzeciw właściwego organu, w szczególności na podstawie przepisów o kontroli inwestycji albo o kontroli koncentracji, a warunek ten nie został jeszcze spełniony --- w tym przypadku bieg terminu z ust.~1 ulega wstrzymaniu do czasu jego spełnienia, a wstrzymanie nie stanowi odmowy.
  \end{enumerate}
  Odmowa wymaga wskazania konkretnej podstawy z niniejszego ustępu. Rada Dyrektorów nie może odmówić zgody z żadnej innej przyczyny.
  
\begin{MemoBox}{1.1.2}
\MemoPyt{1.1.2}{Jak Kryterium ma się do ustawy o kontroli niektórych inwestycji i innych przepisów o ochronie bezpieczeństwa i obronności --- czy umowne Kryterium może współistnieć z kontrolą publiczną i czy powinno wprost do niej odsyłać?}
\Ocena{zgodne}
\Klasa{opcjonalne}
\Podstawa{art.~12a ust.~1 pkt~1, art.~12c ust.~1 pkt~1, 4 i~5, art.~12d ust.~3 pkt~7 i ust.~4, art.~12f ust.~1 i~5, art.~12h ust.~5 i~8, art.~12k ust.~1 oraz art.~16a ustawy z 24 lipca 2015 r. o kontroli niektórych inwestycji \Wiar{Z}; rozporządzenie (UE) 2019/452 (obowiązuje do 16 stycznia 2028~r.; tekst niedostępny) \Wiar{W}; art.~2 pkt~1, 5 i~7, art.~3 ust.~1--2, art.~4 ust.~2, 4--5, 9 i~15--17, art.~5 ust.~1--2, art.~6, art.~9, art.~11 ust.~3--5, art.~19 ust.~1 lit.~a--b, art.~20 ust.~1 i~4, art.~30 ust.~1--3, art.~31 oraz motywy 14, 15 i~20 rozporządzenia (UE) 2026/1386 \Wiar{Z} (tekst EN); art.~10 ust.~1 pkt~2 ustawy z 13 czerwca 2019 r. (koncesja) \Wiar{Z}; art.~57 ust.~2 pkt~1 ustawy o ochronie informacji niejawnych \Wiar{Z}; §~\ref{art:zbywanie} ust.~2 lit.~e.}
\Uzasadnienie{Kontrola publiczna i Kryterium działają na różnych płaszczyznach. \emph{Do 16 stycznia 2028~r.} art.~4--12 ustawy o kontroli niektórych inwestycji nie dotyczą Spółki (brak jej w wykazie), a art.~12a--12k obejmą ją dopiero po przekroczeniu 10~000~000 euro przychodu w Polsce (art.~12d ust.~3 pkt~7 i ust.~4) i tylko wobec nabywców spoza UE/EOG/OECD nabywających co najmniej 20\% głosów albo dominację (art.~12a ust.~1 pkt~1, art.~12c ust.~1 pkt~1); zawiadomienie składa się przed zawarciem umowy zobowiązującej (art.~12f ust.~1 i~5), a nabycie bez zawiadomienia albo mimo sprzeciwu „jest nieważne” (art.~12k ust.~1; kary --- art.~16a) \Wiar{Z}. Kryterium jest surowsze od ustawy i nie może z nią kolidować. \emph{Od 17 stycznia 2028~r.} rozporządzenie (UE) 2026/1386 (uchyla 2019/452) nakazuje objąć uprzednim zezwoleniem --- bez progu przychodowego (progi głosów może określić ustawa krajowa, motyw~20) --- inwestycje każdego inwestora spoza Unii, także z NATO i OECD, dające „skuteczny udział w zarządzaniu lub kontroli” (art.~2 pkt~1 i~5, motyw~15; zwykle także przy prawie powołania dyrektora), w podmiot opracowujący, produkujący lub wprowadzający do obrotu produkty dual-use albo obronne (art.~4 ust.~15) \Wiar{Z} (tekst EN). Ustawa wymaga zmiany \Wiar{?}; objęcie Spółki rozstrzygnie klasyfikacja produktów (§~\ref{art:export} ust.~2) \Wiar{?}. Od 2028~r. Kryterium jest wobec inwestorów z NATO spoza Unii łagodniejsze od prawa publicznego. Przepisy koncesyjne i o bezpieczeństwie przemysłowym nie zakazują umownych ograniczeń \Wiar{Z}. Odsyłanie wprost do ustawy nie jest wskazane: odesłanie rodzajowe w lit.~e jest odporne na zmiany; uzupełnienia wymaga tylko procedura.}
\Rekomendacja{Bez zmian w definicji Kryterium. Opcjonalnie dopisać w §~\ref{art:zbywanie} ust.~2 lit.~e, że zawiadomienie organu (wniosek o zezwolenie) składa nabywca (art.~12f ust.~1 ustawy; w zakresie współdziałania --- zbywca i Spółka), niezwłocznie po zgłoszeniu Spółce zamiaru rozporządzenia, a w każdym razie przed zawarciem umowy zobowiązującej do nabycia (art.~12f ust.~5; od 17 stycznia 2028~r. --- przed zamknięciem transakcji, art.~4 ust.~9 rozporządzenia 2026/1386), oraz że wstrzymanie terminu ustaje z dniem doręczenia decyzji albo upływu terminu na sprzeciw. W harmonogramach rund z inwestorem spoza Unii od 2028~r. uwzględnić 45-dniowe badanie wstępne i możliwe badanie pogłębione, a przy inwestorze kontrolowanym przez państwo trzecie albo projekcie EDF --- terminy mechanizmu współpracy (art.~5--6, art.~9 i art.~11).}
\Kierunek{Doprecyzowanie procedury nie zmienia zakresu Kryterium; ułatwia Radzie ocenę, kiedy termin biegnie ponownie, i zapobiega zarzutowi ukrytej odmowy (zob.~1.3.2).}
\end{MemoBox}

\begin{MemoBox}{1.3.2}
\MemoPyt{1.3.2}{Czy wstrzymanie biegu terminu (lit.~d i~e) jest dopuszczalne, czy też grozi uznaniem za ukrytą odmowę uruchamiającą mechanizm z ust.~3?}

\Ocena{zgodne warunkowo}
\Klasa{zalecane}
\Podstawa{art.~300$^{39}$ §~2--5 KSH \Wiar{Z}; art.~5 i art.~354 KC \Wiar{Z}; art.~89 KC (warunek) \Wiar{Z}; art.~12f ust.~1 i~5 oraz art.~12h ust.~5 i~8 ustawy o kontroli niektórych inwestycji (uprzednie zawiadomienie nabywcy; 30 dni roboczych i 120 dni) \Wiar{Z}.}
\Uzasadnienie{Wstrzymanie biegu terminu jest modyfikacją trybu ustawowego dopuszczalną na podstawie art.~300$^{39}$ §~2 KSH i gospodarczo uzasadnioną: Rada nie może zająć stanowiska bez informacji, a zgoda na nabycie wymagające decyzji organu byłaby zgodą warunkową. Ryzyko ukrytej odmowy powstaje, gdy wstrzymanie nie ma granicy: Rada mogłaby żądać coraz to nowych dokumentów i nigdy nie odmówić, blokując wpis i bieg terminu na wskazanie nabywcy; sąd oceniłby to jako nadużycie (art.~5 KC) i mógłby przyjąć, że odmowa nastąpiła z chwilą pierwszego kompletnego zgłoszenia --- z nieprzewidywalnym skutkiem dla terminów. Dla lit.~d wystarczy ograniczyć żądania Rady do jednego wezwania i określić maksymalny czas wstrzymania (proponujemy 60 dni), po którym Rada rozstrzyga na podstawie posiadanych informacji, a brak informacji o beneficjentach rzeczywistych oznacza niespełnienie Kryterium (lit.~b). Dla lit.~e wstrzymanie trwa tyle, co postępowanie przed organem (30 dni roboczych postępowania wstępnego i do 120 dni kontrolnego, art.~12h ust.~5 i~8 ustawy o kontroli niektórych inwestycji), na co Spółka nie ma wpływu; zawiadomienie jest z mocy ustawy obowiązkiem nabywcy i musi być uprzednie, tj.~złożone przed zawarciem umowy zobowiązującej (art.~12f ust.~1 i~5). Wystarczy dodać, że składa je nabywca w 30 dni od zgłoszenia, a jeżeli tego nie zrobi, Rada rozstrzyga o zgodzie warunkowej (art.~89 KC: zgoda pod warunkiem zawieszającym uzyskania decyzji organu), co jest dopuszczalne i powszechne w praktyce transakcyjnej \Wiar{W}.}
\Rekomendacja{Ograniczyć wstrzymanie z lit.~d do 60 dni i jednego wezwania; w lit.~e przewidzieć zgodę warunkową, jeżeli pozostałe przesłanki są spełnione.}
\Kierunek{Zmiana zabezpiecza Spółkę przed zarzutem nadużycia, a akcjonariusza przed zawieszeniem bez końca; dla rejestru: zgoda warunkowa jest dokumentem, wpis po spełnieniu warunku (przedstawienie decyzji organu).}
\begin{Brzmienie}
(§~\ref{art:zbywanie} ust.~2 lit.~d w nowym brzmieniu) „mimo wezwania, doręczonego w terminie 7 dni od zgłoszenia i wskazującego wyczerpująco brakujące informacje lub oświadczenia, nie przedstawiono ich w terminie 60 dni od doręczenia wezwania --- do upływu tego terminu bieg terminu z ust.~1 ulega wstrzymaniu, a wstrzymanie nie stanowi odmowy; po upływie tego terminu Rada Dyrektorów rozstrzyga na podstawie posiadanych informacji, a brak informacji niezbędnych do ustalenia spełnienia Kryterium Bezpieczeństwa EU/NATO traktuje się jako jego niespełnienie; albo”

(§~\ref{art:zbywanie} ust.~2 lit.~e, zdanie dodawane) „Jeżeli pozostałe przesłanki zgody są spełnione, Rada Dyrektorów może zamiast wstrzymania udzielić zgody pod warunkiem zawieszającym uzyskania wymaganej zgody, braku sprzeciwu albo upływu terminu na sprzeciw właściwego organu; zawiadomienie do organu składa nabywca przy współdziałaniu zbywcy niezwłocznie, nie później niż w terminie 30 dni od zgłoszenia.”
\end{Brzmienie}
\end{MemoBox}

\begin{MemoBox}{1.5.1}
\MemoPytCd{1.5.1}
\begin{Brzmienie}
§~\ref{art:zbywanie} ust.~2 --- nowa lit.~f): „rozporządzenie dotyczy akcji objętych obowiązkiem zbycia zgodnie z §~\ref{art:zwrotne-zbycie}, a planowany nabywca nie przystąpił do umowy wykonawczej, o której mowa w §~\ref{art:zwrotne-zbycie} ust.~13, w zakresie obowiązków związanych z nabywanymi akcjami, w szczególności nie złożył oferty i nie udzielił pełnomocnictwa przewidzianych w §~\ref{art:breach} ust.~3.”
\end{Brzmienie}
\end{MemoBox}

\item Jeżeli Spółka odmawia zgody z przyczyny określonej w ust.~2 lit.~b, na podstawie art.~300$^{39}$ §~2 KSH strony postanawiają, że Spółka wskazuje innego nabywcę będącego Dopuszczalnym Nabywcą i spełniającego Kryterium Bezpieczeństwa EU/NATO w terminie 3 miesięcy od dnia zgłoszenia Spółce zamiaru rozporządzenia akcją; nabywcą może być także Spółka w granicach określonych w art.~300$^{47}$ KSH. Ceną nabycia jest Wartość Godziwa ustalona zgodnie z §~\ref{art:wycena}, płatna w terminie 30 dni od dnia jej ostatecznego ustalenia, nie później niż w dniu przeniesienia akcji; cena może zostać rozłożona na raty, do czego stosuje się odpowiednio §~\ref{art:dziedziczenie} ust.~4, a pierwsza rata jest wówczas płatna w tym terminie. Jeżeli Spółka nie wskazała nabywcy w terminie albo wskazany nabywca nie uiścił w terminie ceny albo jej pierwszej raty, akcjonariusz może swobodnie zbyć akcję zgodnie z art.~300$^{39}$ §~5 KSH, chyba że nie przyjął oferowanej zapłaty; swoboda zbycia nie obejmuje zbycia na rzecz Niedopuszczalnego Nabywcy w zakresie wynikającym z bezwzględnie obowiązujących przepisów prawa.
  
\begin{MemoBox}{1.2.1}
\MemoPyt{1.2.1}{Czy art.~300$^{39}$ §~2 KSH („chyba że umowa spółki stanowi inaczej”) pozwala zarówno na modyfikację mechanizmu ustawowego (A: termin 3 miesięcy zamiast miesiąca, zapłata ratalna, Spółka jako możliwy nabywca), jak i na całkowite wyłączenie §~3--6 dla jednej przesłanki odmowy (B)? Gdzie leżą granice tej swobody?}
\Ocena{zgodne warunkowo (A: zgodne; B: zgodne warunkowo)}
\Klasa{opcjonalne}
\Podstawa{art.~300$^{39}$ §~1--6 KSH \Wiar{Z}; art.~300$^{1}$ §~3 i art.~4 §~2$^{1}$ KSH \Wiar{Z}; art.~182 §~2--5 KSH (dla porównania: model sp.~z~o.o. z udziałem sądu rejestrowego) \Wiar{Z}; art.~57, art.~353$^{1}$ i art.~58 §~1--2 KC \Wiar{Z}; art.~300$^{47}$ KSH \Wiar{Z}.}
\Uzasadnienie{Konstrukcja art.~300$^{39}$ KSH jest dwustopniowa: §~1 pozwala umowie spółki ograniczyć rozporządzanie „w inny sposób” (bez katalogu), a §~2 („chyba że umowa spółki stanowi inaczej”) czyni ustawowy tryb §~3--6 jedynie modelem domyślnym dla wariantu „zgoda spółki”. Skoro umowa mogłaby w ogóle nie sięgać po zgodę spółki, to tym bardziej może --- w wariancie ze zgodą --- zmodyfikować tryb ustawowy (A) albo wyłączyć go w części (B). Klauzula jest wspólna z sp.~z~o.o. (art.~182 §~2 KSH), ale tam model ustawowy angażuje sąd rejestrowy (art.~182 §~3--4), a w P.S.A. termin do wskazania nabywcy, cenę albo sposób jej określenia i termin zapłaty określa umowa spółki (art.~300$^{39}$ §~3 zd.~2), przy czym w modelu domyślnym termin do wskazania nabywcy nie może przekroczyć miesiąca (§~3 zd.~4) \Wiar{Z}. Wariant~A mieści się w §~2: cena i termin zapłaty to materia oddana umowie, a termin 3 miesięcy jest dopuszczalny tylko jako odstępstwo od §~3 na podstawie §~2 --- projekt czyni to wyraźnie, co należy utrzymać; Spółka jako nabywca --- w granicach art.~300$^{47}$ KSH (upoważnienie Walnego Zgromadzenia i test wypłacalności, zob.~1.3.5). Wariant~B jest dopuszczalny co do zasady (projekt stosuje tę technikę w lock-upie), lecz jego granice wyznacza prawo cywilne.

Granice swobody: (1) §~6 (zbycie w egzekucji) chroni wierzycieli i nie podlega derogacji --- projekt to respektuje; (2) art.~58 §~2 i art.~353$^{1}$ KC: ograniczenie nie może prowadzić do faktycznego, bezterminowego wyłączenia zbywalności akcji (art.~57 §~1 KC) \Wiar{W}; derogacja czasowa (lock-up 12 miesięcy) jest bezpieczna, a bezterminowa, uzależniona od cechy nabywcy --- obronna tylko przy realnej ścieżce wyjścia (rynek nabywców spełniających Kryterium jest szeroki: UE/EOG/NATO); (3) zakaz obejścia przepisów o umorzeniu i akcjach własnych; (4) przy derogacji częściowej (tylko lit.~b) trzeba jasno rozstrzygnąć skutek dla pozostałych przesłanek (zob.~1.2.7), inaczej sąd może zastosować model ustawowy jako domyślny; (5) forma dokumentowa czynności z §~3 (§~4; w Spółce --- Rada Dyrektorów, art.~4 §~2$^{1}$ KSH) jest dyspozytywna, ale nie powinna być uchylana, bo służy pewności obrotu i rejestru.}
\Rekomendacja{Oba warianty są dopuszczalne; decyzja między nimi jest decyzją o rozkładzie ryzyka, nie o legalności (zob.~1.2.5). Niezależnie od wyboru utrzymać wyraźne odesłanie do art.~300$^{39}$ §~2 KSH w każdym miejscu, w którym umowa odstępuje od trybu ustawowego (ust.~3, ust.~4 zd.~1, §~\ref{art:okres-ograniczenia-zbywania} ust.~1), aby sąd rejestrowy i podmiot prowadzący rejestr widzieli świadomą derogację.}
\Kierunek{Technika redakcyjna (jednolite odesłania do §~2); brzmienie --- w 1.2.7 (s.~\pageref{memo:1.2.7}).}
\end{MemoBox}

\begin{MemoBox}{1.2.2}
\MemoPyt{1.2.2}{Jak każdy z wariantów zadziała wobec podmiotu prowadzącego rejestr akcjonariuszy (czy odmówi wpisu nabywcy spoza Kryterium, jeżeli akcjonariusz zbędzie akcje mimo odmowy) oraz wobec sądu rejestrowego przy rejestracji umowy?}
\Ocena{zgodne warunkowo}
\Klasa{zalecane}
\Podstawa{art.~300$^{33}$ §~1 pkt~10, art.~300$^{37}$ §~1 i art.~300$^{38}$ §~1 KSH \Wiar{Z}; art.~300$^{34}$ §~1 i~4--6 KSH --- dokumenty, zakres badania i uwzględnianie ograniczeń przez podmiot prowadzący rejestr \Wiar{Z}; art.~300$^{34}$ §~3 zd.~2 KSH w brzmieniu od 18 lutego 2027 r. (forma zgody na wpis) \Wiar{Z}; art.~23 ust.~1 ustawy o KRS (zakres kognicji sądu rejestrowego) \Wiar{Z}; art.~300$^{5}$ KSH (treść umowy) \Wiar{Z}.}
\Uzasadnienie{Podmiot prowadzący rejestr w obu wariantach nie bada Kryterium, lecz dokumenty: wpis następuje na podstawie „dokumentów uzasadniających dokonanie wpisu” (art.~300$^{34}$ §~4 KSH), podmiot bada ich treść i formę, bez obowiązku badania zgodności z prawem, chyba że poweźmie uzasadnione wątpliwości (§~5), i uwzględnia ograniczenia co do rozporządzania akcją (§~6) \Wiar{Z}. Jeżeli w rejestrze ujawniono wymóg zgody Spółki (art.~300$^{33}$ §~1 pkt~10), brak dokumentu zgody jest brakiem formalnym i podmiot odmówi wpisu; zbycie mimo odmowy nie wywoła skutku nabycia (art.~300$^{37}$ §~1, art.~300$^{38}$ §~1 KSH). Różnica ujawnia się po odmowie: w wariancie~A, po upływie terminu na wskazanie nabywcy albo braku zapłaty, podmiot musi ustalić, czy swoboda zbycia powstała --- bez zaświadczenia Spółki zażąda go albo odmówi wpisu, a spór rozstrzygnie sąd. W wariancie~B taka sytuacja nie powstaje (poza zgodą Walnego Zgromadzenia), ale sfałszowanej albo nieważnej zgody podmiot nie wykryje bez kontaktu ze Spółką. Od 18 lutego 2027 r. zgoda zbywcy na wpis będzie wymagała formy kwalifikowanej (art.~300$^{34}$ §~3 zd.~2), co nie zastąpi weryfikacji zgody Spółki. Sąd rejestrowy bada zgodność dokumentów z przepisami co do formy i treści (art.~23 ust.~1 ustawy o KRS); ograniczenia rozporządzania są w P.S.A. wyraźnie dopuszczone, więc oba warianty przejdą rejestrację, a referendarze nie badają proporcjonalności klauzul \Wiar{W}; ryzyko wezwania do usunięcia braków jest niskie.}
\Rekomendacja{W obu wariantach: (a) ujawnić ograniczenia w rejestrze (1.1.6); (b) w wariancie~A dodać obowiązek Spółki wydania w 7 dni zaświadczenia o powstaniu swobody zbycia (art.~300$^{39}$ §~5 KSH), które stanowi dokument dla rejestru; (c) w umowie o prowadzenie rejestru zastrzec, że wpis nabywcy następuje po przedstawieniu zgody Spółki, potwierdzenia z ust.~7 albo zaświadczenia z lit.~b.}
\Kierunek{Zaświadczenie zamyka lukę dowodową między Spółką a rejestrem; nie zmienia praw akcjonariusza. Odmowa wydania zaświadczenia byłaby naruszeniem umowy przez Spółkę, dochodzonym na drodze sądowej (art.~64 KC).
\begin{Brzmienie}
(§~\ref{art:zbywanie} ust.~3, zdanie dodawane na końcu) „Spółka wydaje akcjonariuszowi, w terminie 7 dni od jego żądania, zaświadczenie w formie dokumentowej stwierdzające powstanie prawa swobodnego zbycia akcji albo wskazujące przyczynę, dla której prawo to nie powstało; zaświadczenie stanowi dokument uzasadniający wpis w rejestrze akcjonariuszy.”
\end{Brzmienie}}
\end{MemoBox}

\begin{MemoBox}{1.2.3}
\MemoPyt{1.2.3}{Wariant~A: czy zapłata ratalna i pierwsza rata w 30 dni odpowiadają wymogowi określenia w umowie ceny (sposobu jej ustalenia) i terminu zapłaty; czy niezapłacenie pierwszej raty prawidłowo uruchamia swobodę zbycia; jak zastrzeżenie „chyba że nie przyjął oferowanej zapłaty” (art.~300$^{39}$ §~5 KSH) działa przy zapłacie ratalnej; czy odesłanie do §~17 ust.~4 (raty przewidziane dla spłaty spadkobierców) jest wystarczająco precyzyjne?}
\Ocena{zgodne warunkowo}
\Klasa{zalecane}
\Podstawa{art.~300$^{39}$ §~2, 3 i 5 KSH \Wiar{Z}; art.~300$^{47}$ KSH \Wiar{Z}; art.~327--329 KC (zastaw na prawach) \Wiar{Z} oraz ustawa o zastawie rejestrowym \Wiar{W}; art.~300$^{33}$ §~1 pkt~6 i art.~300$^{37}$ §~1--2 KSH (ujawnienie zastawu w rejestrze; od 18 lutego 2027 r. zastaw rejestrowy powstaje z wpisem do rejestru zastawów bez wpisu w rejestrze akcjonariuszy) \Wiar{Z}; art.~455 i art.~476 KC \Wiar{Z}.}
\Uzasadnienie{Cena: odesłanie do Wartości Godziwej z §~\ref{art:wycena} jest „sposobem określenia” ceny (art.~300$^{39}$ §~3 zd.~2 KSH) i spełnia wymóg ustawowy. Termin zapłaty („30 dni od ostatecznego ustalenia, nie później niż w dniu przeniesienia”) jest określony, ale otwarty: ustalenie wartości przez Niezależnego Eksperta nie ma w §~\ref{art:wycena} terminu końcowego, więc akcjonariusz może czekać miesiącami bez prawa swobodnego zbycia; sąd mógłby uznać, że umowa nie określa terminu zapłaty dostatecznie, a wówczas §~3 zd.~3 („w braku tych postanowień akcja może być zbyta bez ograniczenia”) otwiera swobodę zbycia od razu \Wiar{Z} co do brzmienia; czy niedostatecznie określony termin jest „brakiem postanowień”, pozostaje kwestią wykładni \Wiar{?}. Raty: KSH ich nie zabrania; uzależnienie swobody zbycia od braku zapłaty pierwszej raty „stanowi inaczej” (§~2) i jest skuteczne. Po zapłacie pierwszej raty ryzyko pozostałych rat spada na zbywcę; ochroną jest wyłącznie zabezpieczenie i wymagalność całości z §~\ref{art:dziedziczenie} ust.~4. Zastrzeżenie „chyba że nie przyjął oferowanej zapłaty” działa wobec pierwszej raty: akcjonariusz, który odmawia jej przyjęcia (np.~kwestionując wiążącą wycenę), nie uzyskuje swobody zbycia. Odesłanie do §~\ref{art:dziedziczenie} ust.~4 jest zbyt ogólne: (1) nie mówi, kto stwierdza zagrożenie płynności --- tu nabywca oceniałby własną płynność; (2) „zabezpieczenie adekwatne” nie jest zdefiniowane, a zbywca nie ma wpływu na jego wybór; (3) nie określa, kiedy następuje przeniesienie akcji, z którym wiąże się termin pierwszej raty.}
\Rekomendacja{Przeredagować ust.~3: (a) termin końcowy: jeżeli Wartość Godziwa nie zostanie ostatecznie ustalona w 90 dni od wskazania nabywcy, nabywca płaci pierwszą ratę od wartości wskazanej przez Radę z wyrównaniem po ustaleniu, a w braku zapłaty powstaje swoboda zbycia; (b) raty tylko za zgodą zbywcy albo, bez jego zgody, wyłącznie gdy pierwsza rata wynosi co najmniej 50\% ceny, liczba rat nie przekracza 12, a zabezpieczeniem jest zastaw rejestrowy na nabywanych akcjach albo gwarancja bankowa; (c) wyraźne określenie dnia przeniesienia (wpis w rejestrze po zapłacie pierwszej raty).}
\Kierunek{Zmiana zachowuje sens wariantu~A (Spółka kupuje czas i rozkłada gotówkę), ale przenosi część ryzyka rat z powrotem na nabywcę; inwestorzy VC odbiorą to jako uczciwsze niż odesłanie do reżimu spadkowego; zastaw na akcjach jest ujawniany w rejestrze. Proponowane nowe brzmienie ust.~3 (pełny tekst --- w memorandum.md) zachowuje wskazanie nabywcy w 3 miesiące (także Spółki w granicach art.~300$^{47}$ KSH) i cenę równą Wartości Godziwej płatną w 30 dni od jej ostatecznego ustalenia, a ponadto: jeżeli Wartość Godziwa nie zostanie ostatecznie ustalona w 90 dni od wskazania nabywcy --- zapłata w kolejnych 30 dniach wartości wskazanej przez Radę Dyrektorów we wskazaniu nabywcy, z wyrównaniem różnicy w 14 dni od ostatecznego ustalenia; raty za zgodą zbywcy na zasadach §~\ref{art:dziedziczenie} ust.~4, a bez niej tylko przy pierwszej racie co najmniej 50\% ceny, najwyżej 12 ratach i zabezpieczeniu zastawem rejestrowym na nabywanych akcjach albo gwarancją bankową (pierwsza rata zawsze w terminie określonym dla ceny); przeniesienie przez wpis w rejestrze akcjonariuszy po zapłacie ceny albo pierwszej raty; skutki braku wskazania nabywcy albo zapłaty --- jak dotąd (art.~300$^{39}$ §~5 KSH, bez zbycia na rzecz Niedopuszczalnego Nabywcy); na końcu zdanie o zaświadczeniu Spółki z 1.2.2.}
\end{MemoBox}

\begin{MemoBox}{1.2.4}
\MemoPyt{1.2.4}{Wariant~B: czy istnieje ryzyko uznania go za nadmierne ograniczenie zbywalności (akcja „uwięziona”, gdy nie ma nabywcy spełniającego Kryterium), a jeżeli tak --- czy wystarcza zgoda Walnego Zgromadzenia z §~11 ust.~6, czy potrzebny jest dodatkowy „wentyl” (np.~obowiązek wskazania nabywcy po dłuższym terminie)?}
\Ocena{ryzyko (w brzmieniu bez wentyla)}
\Klasa{konieczne (jeżeli wybrany zostanie wariant~B) (K1)}
\Podstawa{art.~300$^{39}$ §~1--3 KSH \Wiar{Z}; art.~57 §~1, art.~353$^{1}$, art.~58 §~2--3 KC \Wiar{Z}; art.~5 KC \Wiar{Z}; art.~300$^{45}$ KSH \Wiar{Z}.}
\Uzasadnienie{Ryzyko istnieje i jest realne dla akcjonariusza mniejszościowego. W wariancie~B jedyną drogą zbycia akcji osobie spoza Kryterium jest uchwała Walnego Zgromadzenia większością 75\% wszystkich głosów --- decyzja uznaniowa większości, od której nie ma odwołania i której odmowa nie rodzi żadnego obowiązku po stronie Spółki. Dopóki istnieje rynek nabywców spełniających Kryterium, ograniczenie jest proporcjonalne; jeżeli jednak w konkretnym przypadku nabywcy nie ma, akcja staje się faktycznie niezbywalna na czas nieoznaczony, a większość może to wykorzystać do wymuszenia sprzedaży po zaniżonej cenie. Sąd mógłby wtedy sięgnąć po art.~58 §~2 KC albo art.~5 KC, z niepewnym skutkiem: nieważność częściowa klauzuli (art.~58 §~3 KC) oznaczałaby powrót do modelu ustawowego (§~3--6) z terminem miesiąca i ceną, której umowa dla tej ścieżki nie określa --- czyli, wobec „braku tych postanowień”, swobodę zbycia (§~3 zd.~3) \Wiar{Z}. To ryzyko jest gorsze dla Spółki niż wariant~A. Zgoda Walnego Zgromadzenia nie wystarcza jako wentyl, bo jest prawem większości, nie akcjonariusza. Wentyl powinien być obiektywny: po udokumentowanym, bezskutecznym poszukiwaniu nabywcy przez 12 miesięcy od odmowy akcjonariusz może zażądać, aby Spółka wskazała nabywcę na zasadach wariantu~A; brak wskazania albo zapłaty otwiera swobodę zbycia. „Wariant~B z opóźnionym~A” zachowuje twardość bramki w horyzoncie roku i chroni przed zarzutem uwięzienia.}
\Rekomendacja{Jeżeli Założyciele wybiorą wariant~B, uzupełnić go o wentyl czasowy; bez wentyla nie rekomendujemy wariantu~B.}
\begin{Brzmienie}
(§~\ref{art:zbywanie} ust.~3 w wariancie~B, zdanie dodawane po zdaniu o zgodzie Walnego Zgromadzenia) „Jeżeli w terminie 12 miesięcy od dnia odmowy zgody akcjonariusz, mimo udokumentowanych starań, nie zbył akcji Dopuszczalnemu Nabywcy spełniającemu Kryterium Bezpieczeństwa EU/NATO, a Walne Zgromadzenie nie wyraziło zgody, o której mowa w §~\ref{art:dopuszczalny-nabywca} ust.~6, akcjonariusz może zażądać, aby Spółka wskazała innego nabywcę; do wskazania nabywcy, ceny, terminu zapłaty i skutków ich niedochowania stosuje się wówczas odpowiednio zasady określone dla nabywcy wskazanego przez Spółkę w §~\ref{art:dopuszczalny-nabywca} ust.~15 zdanie trzecie, przy czym termin 3 miesięcy biegnie od dnia doręczenia żądania. Postanowienie to nie ma zastosowania w Okresie Ograniczenia Zbywania Akcji.”
\end{Brzmienie}
\end{MemoBox}

\begin{MemoBox}{1.2.5}
\MemoPyt{1.2.5}{Który wariant jest w Państwa ocenie skuteczniejszy i bezpieczniejszy dla Spółki, a jak oba będą odbierane przez inwestorów VC i partnerów przemysłowych? Prosimy o rekomendację; wyboru dokonają Założyciele.}
\Ocena{zgodne warunkowo (wariant~A z poprawkami z 1.2.3)}
\Klasa{zalecane}
\Podstawa{art.~300$^{39}$ §~2--5 i art.~300$^{47}$ KSH \Wiar{Z}; art.~58 KC \Wiar{Z}; art.~9 ust.~3--4 rozporządzenia 2021/697 \Wiar{Z}; \texttt{vc.md} sekcje~1, 4--6 \Wiar{W}; \texttt{web\_edf\_gowling.txt} \Wiar{Z}.}
\Uzasadnienie{\textit{Skuteczność bramki.} Wariant~B jest szczelniejszy na papierze: akcje nie mogą trafić poza Kryterium bez uchwały 75\%. Wariant~A jest szczelny tylko wtedy, gdy Spółka albo wskazany nabywca ma środki na Wartość Godziwą w 3 miesiące (z ratami: na pierwszą ratę); w przeciwnym razie akcje wychodzą poza Kryterium. \textit{Bezpieczeństwo prawne.} Wariant~A jest niemal wolny od ryzyka podważenia (mieści się w literze §~2--5); wariant~B niesie ryzyko z 1.2.4, którego materializacja jest dla Spółki gorsza niż słabość wariantu~A, bo skutkiem nieważności byłaby swoboda zbycia bez żadnej bramki. \textit{Odbiór rynkowy.} Fundusze VC oczekują zawsze istniejącej ścieżki wyjścia za wartość godziwą; wariant~B bez wentyla będzie kwestionowany w każdym term sheecie, wariant~A jest akceptowalny, a ratalność będzie negocjowana (\texttt{vc.md} sekcja~5 \Wiar{W}). Partnerzy przemysłowi i programy obronne (EDF, DIANA) patrzą na \emph{kontrolę} i strukturę właścicielską w dniu oceny (EDF: odbiorcy „nie podlegają kontroli niestowarzyszonego państwa trzeciego lub podmiotu z niestowarzyszonego państwa trzeciego”, art.~9 ust.~3 rozporządzenia 2021/697 \Wiar{Z}; tak samo \texttt{web\_edf\_gowling.txt} \Wiar{Z}), nie na mechanikę odmowy zgody; dla nich oba warianty są równoważne, a istotna jest warstwa „kontroli UE/EOG” z 1.1.1. \textit{Znaczenie praktyczne.} Kryterium obejmuje 30+ państw; przypadek, w którym jedynym chętnym nabywcą jest podmiot spoza UE/EOG/NATO (zwykle inwestor strategiczny z Azji albo Zatoki), będzie rzadki, a Walne Zgromadzenie i tak zdecyduje wtedy świadomie (w obu wariantach). Różnica sprowadza się do rzadkiego scenariusza, w którym wariant~B daje bramkę kosztem ryzyka prawnego i kosztu negocjacyjnego, a wariant~A daje pewność prawną kosztem gotówki.

Rekomendujemy wariant~A z poprawkami z 1.2.3 (termin końcowy wyceny, raty z zabezpieczeniem, zaświadczenie), z derogacją twardą wyłącznie dla lock-upu (§~\ref{art:okres-ograniczenia-zbywania}). Jeżeli Założyciele przywiązują większą wagę do szczelności niż do pozycji negocjacyjnej wobec funduszy, drugim wyborem jest wariant~B z wentylem z 1.2.4 („B z opóźnionym~A”), odzwierciedlony wtedy także w §~\ref{art:dopuszczalny-nabywca} ust.~5 zd.~2 i ust.~15 (płatność jednorazowa) oraz w Załączniku nr~4 memorandum (wykaz zmian dla wariantu~B).}
\Rekomendacja{Wariant~A z poprawkami; wybór należy do Założycieli. W obu wariantach uzupełnić §~\ref{art:kwalifikowana-runda} albo umowę inwestycyjną o „przeniesienia dozwolone” dla inwestora finansowego (podmioty powiązane, fundusz następca), bo to --- nie mechanika odmowy --- będzie pierwszym punktem sporu z funduszem (\texttt{vc.md} sekcja~6 pkt~2).}
\Kierunek{Utrzymanie wariantu~A nie wymaga zmian poza 1.2.3; wybór wariantu~B wymaga zmian z 1.2.4 oraz spójnych zmian w §~\ref{art:dopuszczalny-nabywca} ust.~5 i~15, §~\ref{art:zbywanie} ust.~4 i §~\ref{art:okres-ograniczenia-zbywania} ust.~1.}
\end{MemoBox}

\item W przypadku odmowy z przyczyny określonej w ust.~2 lit.~a obowiązek wskazania innego nabywcy nie powstaje, ponieważ przeszkoda wynika z bezwzględnie obowiązującego prawa. W przypadku odmowy z przyczyny określonej w ust.~2 lit.~c stosuje się §~\ref{art:okres-ograniczenia-zbywania}, który --- w granicach art.~300$^{39}$ §~2 KSH --- wyłącza na czas Okresu Ograniczenia Zbywania stosowanie mechanizmu wskazania nabywcy z ust.~3.
  
\begin{MemoBox}{1.2.7}
\MemoPyt{1.2.7}{§~12 ust.~4 zdanie pierwsze: przy odmowie z powodu Niedopuszczalnego Nabywcy obowiązek wskazania innego nabywcy nie powstaje, bo przeszkoda wynika z prawa. Czy to ujęcie jest spójne z art.~300$^{39}$ KSH, czy wymaga wyraźnej derogacji jak w ust.~3?}
\Ocena{zgodne warunkowo}
\Klasa{zalecane}
\Podstawa{art.~300$^{39}$ §~2--5 KSH \Wiar{Z}; art.~58 §~1 KC \Wiar{Z}; art.~1 lit.~f--g i art.~2 ust.~1--2 rozporządzenia (UE) nr~269/2014 (zamrożenie środków obejmujące „stocks and shares”; zakaz udostępniania) \Wiar{Z} (tekst EN, pierwotny, bez konsolidacji); rozporządzenie 833/2014 \Wiar{W}; art.~1 pkt~1--2 w zw. z art.~2 ust.~1 ustawy z 13 kwietnia 2022 r. o szczególnych rozwiązaniach w zakresie przeciwdziałania wspieraniu agresji na Ukrainę (lista krajowa; odpowiednie stosowanie środków z art.~2 i art.~9 rozporządzenia 269/2014) \Wiar{Z}.}
\Uzasadnienie{Merytorycznie ujęcie jest trafne: jeżeli nabycie przez planowanego nabywcę jest zakazane przez prawo, umowa zbycia byłaby nieważna (art.~58 §~1 KC), a odmowa zgody Spółki ma charakter deklaratoryjny; nie ma powodu, by Spółka szukała nabywcy zastępczego, skoro zbywca może wskazać innego (ust.~6). Technicznie jednak art.~300$^{39}$ §~3 zd.~1 KSH („jeżeli spółka odmawia zgody na zbycie akcji, powinna wskazać innego nabywcę”) nie różnicuje przyczyn odmowy: każda odmowa uruchamia obowiązek wskazania nabywcy, „chyba że umowa spółki stanowi inaczej” (§~2). Zdanie „obowiązek nie powstaje, ponieważ przeszkoda wynika z prawa” jest uzasadnieniem, nie postanowieniem „stanowiącym inaczej”; ostrożny sąd albo podmiot prowadzący rejestr może uznać, że brak wyraźnej derogacji oznacza stosowanie modelu ustawowego, a wtedy brak wskazania nabywcy w miesiąc dawałby --- paradoksalnie --- swobodę zbycia (choć nadal ograniczoną przepisami sankcyjnymi). Wystarczy formuła derogacyjna analogiczna do ust.~3 i §~\ref{art:okres-ograniczenia-zbywania} ust.~1. To samo dotyczy przesłanek z lit.~d i~e, które nie są odmową, ale warto to potwierdzić w tym samym zdaniu.}
\Rekomendacja{Dodać wyraźną derogację w ust.~4 zd.~1.}
\Kierunek{Zmiana redakcyjna bez skutków gospodarczych; zwiększa pewność wobec rejestru i sądu.
\begin{Brzmienie}
(§~\ref{art:zbywanie} ust.~4, zdanie pierwsze w nowym brzmieniu) „W przypadku odmowy z przyczyny określonej w ust.~2 lit.~a, na podstawie art.~300$^{39}$ §~2 KSH strony postanawiają, że przepisy art.~300$^{39}$ §~3--5 KSH nie znajdują zastosowania, a obowiązek wskazania innego nabywcy nie powstaje, ponieważ przeszkoda wynika z bezwzględnie obowiązującego prawa; zbywca może wskazać innego nabywcę zgodnie z ust.~6. Wstrzymanie biegu terminu na podstawie ust.~2 lit.~d albo~e nie jest odmową zgody i nie uruchamia przepisów art.~300$^{39}$ §~3--5 KSH ani mechanizmu z ust.~3.”
\end{Brzmienie}}
\end{MemoBox}

\item Zbycie albo obciążenie akcji na rzecz Niedopuszczalnego Nabywcy jest niedopuszczalne w zakresie wynikającym z bezwzględnie obowiązujących przepisów prawa. Upływ terminu, brak stanowiska organu Spółki, uchwała akcjonariuszy ani uchwała Rady Dyrektorów nie mogą sanować transakcji zakazanej przez prawo.
  \item Jeżeli planowany nabywca jest Niedopuszczalnym Nabywcą, zbywca może wskazać innego Dopuszczalnego Nabywcę i przeprowadzić transakcję z zachowaniem pozostałych postanowień niniejszej umowy, w szczególności prawa pierwszeństwa, prawa przyłączenia, prawa przymusowego współzbycia oraz ograniczeń obowiązujących Założycieli.
  \item Zgody Spółki nie wymaga zbycie akcji w postępowaniu egzekucyjnym (art.~300$^{39}$ §~6 KSH) ani nabycie akcji w wykonaniu prawa pierwszeństwa, prawa przyłączenia, prawa przymusowego współzbycia, mechanizmu zwrotnego zbycia akcji albo w ramach Kwalifikowanej Rundy Finansowania zatwierdzonej zgodnie z §~\ref{art:kwalifikowana-runda} --- weryfikacja, czy nabywca jest Dopuszczalnym Nabywcą, pozostaje jednak wymagana w każdym przypadku. Wyłączenie wymogu zgody Spółki nie uchyla wymogu spełniania przez nabywcę Kryterium Bezpieczeństwa EU/NATO: nabycie akcji przez osobę lub podmiot niespełniający tego Kryterium --- także w wykonaniu prawa pierwszeństwa, prawa przyłączenia, prawa przymusowego współzbycia, mechanizmu zwrotnego zbycia akcji albo w ramach Kwalifikowanej Rundy Finansowania --- wymaga zgody Walnego Zgromadzenia zgodnie z §~\ref{art:dopuszczalny-nabywca}. Przy wykonaniu prawa przymusowego współzbycia spełnienie Kryterium albo uzyskanie tej zgody weryfikuje się przed doręczeniem zawiadomienia o wykonaniu prawa. Zdania poprzedzające nie dotyczą zbycia akcji w postępowaniu egzekucyjnym.
  
\begin{MemoBox}{1.2.6}
\MemoPytCd{1.2.6}
\Kierunek{Dopisać przeniesienie do podmiotu w całości kontrolowanego (§~\ref{art:okres-ograniczenia-zbywania} ust.~2 lit.~d) do wyłączeń zgody Spółki (i prawa pierwszeństwa, §~\ref{art:pierwszenstwo}), z zachowaniem weryfikacji Dopuszczalnego Nabywcy i Kryterium.
\begin{Brzmienie}
(§~\ref{art:zbywanie} ust.~7, uzupełnienie zdania pierwszego) „[...] albo w ramach Kwalifikowanej Rundy Finansowania zatwierdzonej zgodnie z §~\ref{art:kwalifikowana-runda}, ani przeniesienie akcji na warunkach określonych w §~\ref{art:okres-ograniczenia-zbywania} ust.~2 lit.~d --- weryfikacja, czy nabywca jest Dopuszczalnym Nabywcą, pozostaje jednak wymagana w każdym przypadku.”
\end{Brzmienie}}
\end{MemoBox}

\begin{MemoBox}{1.2.8}
\MemoPyt{1.2.8}{§~12 ust.~7: zgody Spółki nie wymaga zbycie w postępowaniu egzekucyjnym ani nabycie w wykonaniu prawa pierwszeństwa, przyłączenia, przymusowego współzbycia, zwrotnego zbycia i w Kwalifikowanej Rundzie --- przy zachowaniu wymogu Kryterium (poza egzekucją). Czy wymóg Kryterium przy tych nabyciach jest skuteczny, skoro nie ma tam zgody Spółki, a przy drag-along weryfikuje się go przed zawiadomieniem?}
\Ocena{zgodne warunkowo}
\Klasa{zalecane}
\Podstawa{art.~300$^{39}$ §~1 KSH („w inny sposób je ograniczyć”) i §~6 \Wiar{Z}; art.~300$^{33}$ §~1 pkt~10, art.~300$^{34}$ §~6 i art.~300$^{37}$ §~1--2 KSH \Wiar{Z}; art.~300$^{103}$--300$^{105}$ KSH (emisja i objęcie akcji nowej emisji) \Wiar{Z}; art.~910 i art.~911$^{3}$ KPC (egzekucja z praw majątkowych; zajęcie ujawniane w rejestrze akcjonariuszy) \Wiar{Z}.}
\Uzasadnienie{Wymóg Kryterium bez zgody Spółki jest skuteczny jako ograniczenie „w inny sposób” (art.~300$^{39}$ §~1 KSH): zamknięty krąg nabywców jest ograniczeniem niezależnym od mechanizmu zgody i wiąże strony transakcji oraz --- po ujawnieniu (pkt~10) --- podmiot prowadzący rejestr. Słabością jest brak dokumentu: podmiot prowadzący rejestr sam nie oceni Kryterium; wpisze nabywcę na podstawie umowy, a naruszenie ujawni się później. Przy Kwalifikowanej Rundzie problem nie występuje: objęcie akcji nowej emisji (art.~300$^{103}$ KSH) jest wyłączone spod wpisu konstytutywnego (art.~300$^{37}$ §~2 KSH), a Kryterium bada się w uchwale kierunkowej. Przy prawie pierwszeństwa nabywcami są Spółka albo akcjonariusze już zweryfikowani (chyba że utracili Kryterium --- §~\ref{art:dopuszczalny-nabywca} ust.~15). Przy przyłączeniu i współzbyciu weryfikacja osoby trzeciej „przed doręczeniem zawiadomienia” jest prawidłowo umiejscowiona (§~\ref{art:przymusowe-wspolzbycie} ust.~2 lit.~e), ale brak dokumentu dla rejestru; przy zwrotnym zbyciu weryfikacja jest aktem Spółki. Rozwiązaniem jest „potwierdzenie” Rady Dyrektorów: nie zgoda w rozumieniu art.~300$^{39}$ (aby nie uruchamiać §~3--5), lecz stwierdzenie wyniku weryfikacji z §~\ref{art:dopuszczalny-nabywca} ust.~9 w formie dokumentowej, w 14 dni, jako dokument dla rejestru. Egzekucja pozostaje poza Kryterium z mocy §~6 --- słusznie. P.S.A. nie ma odpowiednika art.~185 KSH \Wiar{Z}; umowne prawo wskazania nabywcy nie wiązałoby komornika, więc realną ochroną jest objęcie zgodą Spółki obciążenia akcji oraz --- wobec nabywcy licytacyjnego --- Kryterium jako ograniczenie „w inny sposób” o ograniczonej skuteczności praktycznej \Wiar{W}.}
\Rekomendacja{Dodać w ust.~7 potwierdzenie Rady jako dokument dla rejestru; nie wprowadzać mechanizmu wskazania nabywcy w egzekucji na wzór art.~185 KSH (brak podstawy ustawowej w P.S.A.); ryzyko egzekucyjne ograniczać przez wymóg zgody Spółki na obciążenie akcji.}
\Kierunek{Potwierdzenie nie jest zgodą (nie uruchamia obowiązku wskazania nabywcy), więc nie zmienia pozycji stron; daje rejestrowi dokument, a Spółce moment kontroli. Dla inwestora: neutralne. Dla sądu rejestrowego: bez wpływu.
\begin{Brzmienie}
(§~\ref{art:zbywanie} ust.~7, zdanie dodawane po zdaniu o przymusowym współzbyciu) „W przypadkach, w których zgoda Spółki nie jest wymagana, Rada Dyrektorów stwierdza wynik weryfikacji, o której mowa w §~\ref{art:dopuszczalny-nabywca} ust.~9, w formie dokumentowej w terminie 14 dni od przedstawienia dokumentów; potwierdzenie to nie stanowi zgody w rozumieniu art.~300$^{39}$ KSH, nie uruchamia mechanizmu z ust.~3 i stanowi dokument uzasadniający wpis nabywcy w rejestrze akcjonariuszy.”
\end{Brzmienie}}
\end{MemoBox}

\begin{MemoBox}{4.4.2}
\MemoPytCd{4.4.2}
\Kierunek{(III)~Przeniesienia dozwolone --- nowy ust.~8 w~§~\ref{art:zbywanie}, po ust.~7 (ust.~7 zostaje bez zmian, aby nie naruszać jego struktury). Bez zgody Spółki i~bez prawa pierwszeństwa z~§~\ref{art:pierwszenstwo} Kwalifikowany Inwestor Finansowy przenosi akcje („Przeniesienie Dozwolone”) na rzecz: (a)~podmiotu zarządzanego przez ten sam lub powiązany podmiot zarządzający (fundusz równoległy, kontynuacyjny, następca, spółka celowa), jeżeli jest on Kwalifikowanym Inwestorem Finansowym; (b)~swoich inwestorów przy likwidacji albo wydaniu aktywów w~naturze, proporcjonalnie, jeżeli każdy z~nich jest Dopuszczalnym Nabywcą i~spełnia Kryterium zgodnie z~§~\ref{art:dopuszczalny-nabywca} ust.~6 (tu look-through jest uzasadniony, bo LP staje się akcjonariuszem). Warunki: przystąpienie nabywcy do obowiązków związanych z~akcjami; zawiadomienie Rady co najmniej 14 dni przed przeniesieniem z~dokumentami z~§~\ref{art:dopuszczalny-nabywca} ust.~8; weryfikacja z~§~\ref{art:dopuszczalny-nabywca} ust.~9 pozostaje, a~jej wynik Rada stwierdza przed przeniesieniem; ust.~2 lit.~a, d i~e stosuje się odpowiednio. Przeniesienie Dozwolone nie uruchamia prawa przyłączenia (§~\ref{art:przylaczenie}); gdy nabywca z~lit.~a przestanie spełniać §~\ref{art:dopuszczalny-nabywca} ust.~14 --- stosuje się §~\ref{art:dopuszczalny-nabywca} ust.~15. Rejestr: ograniczenia co do rozporządzania akcją są treścią rejestru (art.~300$^{33}$ §~1 pkt~10 KSH \Wiar{Z}); podmiot prowadzący rejestr będzie oczekiwał dokumentu potwierdzającego spełnienie przesłanek wyjątku --- stąd wymóg protokołu Rady.}
\end{MemoBox}

\item Naruszenie wymogów dotyczących Dopuszczalnego Nabywcy stanowi istotne naruszenie niniejszej umowy. Zbywca i nabywca, w zakresie w jakim są związani niniejszą umową, są zobowiązani niezwłocznie usunąć skutki naruszenia i doprowadzić stan prawny do zgodności z bezwzględnie obowiązującymi przepisami prawa. Postanowienie to nie przesądza sposobu działania podmiotu prowadzącego rejestr akcjonariuszy, który działa zgodnie z przepisami prawa.

\begin{MemoBox}{1.1.6}
\MemoPyt{1.1.6}{Czy dla skuteczności Kryterium wobec osób trzecich potrzebne jest ujawnienie ograniczenia w rejestrze akcjonariuszy, a jeżeli tak --- w jakiej formie i na czyj wniosek?}
\Ocena{zgodne warunkowo}
\Klasa{zalecane}
\Podstawa{art.~300$^{33}$ §~1 pkt~10 („ograniczenia co do rozporządzania akcją”) i pkt~11 („postanowienia umowy spółki o związanych z akcją obowiązkach wobec spółki”) oraz §~2 KSH \Wiar{Z}; art.~300$^{34}$ §~1 i~4--6 KSH (wpis na żądanie spółki albo osoby mającej interes prawny; uwzględnianie ograniczeń) \Wiar{Z}; art.~300$^{37}$ §~1 i art.~300$^{38}$ §~1 KSH (skutek wpisu) \Wiar{Z}; art.~300$^{32}$ §~1 KSH (umowa o prowadzenie rejestru) \Wiar{Z}; art.~300$^{33}$ §~3 i art.~594 §~1 pkt~2$^{1}$ KSH w brzmieniu obowiązującym od 18 lutego 2027 r. (Dz.U. 2026 poz.~176) \Wiar{Z}.}
\Uzasadnienie{Kryterium wiąże każdego akcjonariusza z mocy umowy spółki niezależnie od wpisu; wpis nie jest przesłanką ważności ograniczenia. Ujawnienie jest jednak potrzebne dla skuteczności praktycznej: podmiot prowadzący rejestr „przy dokonywaniu wpisów [...] uwzględnia ograniczenia co do rozporządzania akcją” (art.~300$^{34}$ §~6 KSH), więc znając je zażąda zgody Spółki i bez niej odmówi wpisu nabywcy, a nabycie następuje dopiero z chwilą wpisu (art.~300$^{37}$ §~1 KSH) \Wiar{Z}. Bez adnotacji może wpisać nabywcę na podstawie samej umowy zbycia; ujawnienie ma też walor dowodowy. Forma: rodzajowa adnotacja przy każdej akcji (zgoda Spółki, Kryterium, zgoda Walnego Zgromadzenia, obowiązki zbycia z §~\ref{art:dopuszczalny-nabywca} ust.~15, §~\ref{art:zwrotne-zbycie} i §~\ref{art:dziedziczenie}), a przy akcjach Założycieli --- Okres Ograniczenia Zbywania z datą końcową. Wniosek składa Spółka (Rada Dyrektorów) przy zawarciu umowy o prowadzenie rejestru i przy każdej zmianie; interes prawny ma też każdy akcjonariusz. Projekt przewiduje ujawnienie tylko dla zwrotnego zbycia (§~\ref{art:zwrotne-zbycie} ust.~14) i przez odesłanie w §~\ref{art:dopuszczalny-nabywca} ust.~15; brakuje ogólnego obowiązku Rady. Od 18 lutego 2027 r. lukę częściowo wypełni nowy art.~300$^{33}$ §~3 KSH: zarząd (w Spółce --- Rada Dyrektorów, art.~4 §~2$^{1}$ KSH) zgłasza zmiany danych z §~1 pkt~1--4 i~9--11 „w terminie siedmiu dni”, pod groźbą grzywny (art.~594 §~1 pkt~2$^{1}$ KSH). Umowa powinna przejąć ten termin już teraz (7 dni zamiast 14) i objąć nim pierwsze zgłoszenie, którego nowelizacja wprost nie reguluje; art.~300$^{33}$ §~2 KSH pozwala nadto zapisać w umowie spółki dodatkowe informacje ujawniane w rejestrze.}
\Rekomendacja{Dodać w §~\ref{art:zbywanie} ustęp nakładający na Radę Dyrektorów obowiązek zgłoszenia do rejestru ograniczeń i obowiązków z §~\ref{art:dopuszczalny-nabywca}, §~\ref{art:zbywanie}, §~\ref{art:okres-ograniczenia-zbywania}, §~\ref{art:zwrotne-zbycie} i §~\ref{art:dziedziczenie} w terminie 7 dni od zawarcia umowy o prowadzenie rejestru i od każdego zdarzenia uzasadniającego wpis (zgodnie z art.~300$^{33}$ §~3 KSH w brzmieniu od 18 lutego 2027 r.); zawrzeć w tej umowie zobowiązanie podmiotu prowadzącego rejestr do żądania zgody Spółki albo potwierdzenia z §~\ref{art:zbywanie} ust.~7 (zob.~1.2.8) przed każdym wpisem nabywcy.}
\Kierunek{Zmiana porządkowa; nie wpływa na treść ograniczeń. Dla rejestru: jasna instrukcja; dla inwestora: przewidywalność; dla sądu rejestrowego: bez wpływu.
\begin{Brzmienie}
(nowy ust.~9 w §~\ref{art:zbywanie}) „Rada Dyrektorów zapewnia ujawnienie w rejestrze akcjonariuszy, zgodnie z art.~300$^{33}$ §~1 KSH, ograniczeń rozporządzania akcjami oraz obowiązków związanych z akcją wynikających z §~\ref{art:dopuszczalny-nabywca}, §~\ref{art:zbywanie}, §~\ref{art:okres-ograniczenia-zbywania}, §~\ref{art:zwrotne-zbycie} i §~\ref{art:dziedziczenie}, w terminie 7 dni od zawarcia umowy o prowadzenie rejestru oraz od każdego zdarzenia uzasadniającego zmianę wpisu; postanowienie to stanowi dodatkowe postanowienie dotyczące informacji ujawnianych w rejestrze w rozumieniu art.~300$^{33}$ §~2 KSH. Umowa o prowadzenie rejestru przewiduje, że wpis nabywcy akcji następuje po przedstawieniu zgody Spółki albo potwierdzenia, o którym mowa w ust.~7.”
\end{Brzmienie}}
\end{MemoBox}

\end{enumerate}
\end{ContractClause}
\begin{ContractClause}{Okres Ograniczenia Zbywania Akcji Założycieli}
\label{art:okres-ograniczenia-zbywania}
\begin{enumerate}
  \item Założyciel Pierwotny przez 12 miesięcy od dnia wpisu Spółki do rejestru, a Założyciel Funkcjonalny Serii F przez 12 miesięcy od jego Dnia Przyznania Serii F, nie może zbyć ani obciążyć objętych tym statusem akcji bez uprzedniej zgody Walnego Zgromadzenia podjętej większością 75\% Głosów Uprawnionych w Sprawie, przy czym osoba występująca o zgodę jest w tej sprawie wyłączona od głosowania. Ograniczenie to stanowi ograniczenie rozporządzenia akcją w rozumieniu art.~300$^{39}$ §~1 KSH; na podstawie art.~300$^{39}$ §~2 KSH strony postanawiają, że w Okresie Ograniczenia Zbywania Akcji przepisy art.~300$^{39}$ §~3--6 KSH oraz mechanizm wskazania innego nabywcy z §~\ref{art:zbywanie} ust.~3 nie znajdują zastosowania, a odmowa zgody nie rodzi po stronie Spółki obowiązku wskazania nabywcy.
  
\begin{MemoBox}{1.2.6}
\MemoPyt{1.2.6}{Lock-up (§~13 ust.~1): Założyciel Pierwotny przez 12 miesięcy od wpisu Spółki (Założyciel Serii~F --- od Dnia Przyznania) nie może zbyć ani obciążyć akcji bez zgody Walnego Zgromadzenia (75\% Głosów Uprawnionych w Sprawie); na podstawie art.~300$^{39}$ §~2 KSH w tym okresie wyłączono §~3--6 i mechanizm z §~12 ust.~3. Czy taka derogacja na czas określony jest dopuszczalna i jak ją ujawnić w rejestrze? Czy wyjątki z §~13 ust.~2 (runda, tag-along i drag-along, zwrotne zbycie, przeniesienie do podmiotu w całości kontrolowanego) wymagają jeszcze zgody Spółki z §~12?}
\Ocena{zgodne warunkowo}
\Klasa{zalecane}
\Podstawa{art.~300$^{39}$ §~1--2 i §~6 KSH \Wiar{Z}; art.~57 §~2 KC \Wiar{Z}; art.~300$^{33}$ §~1 pkt~10 i §~2 KSH \Wiar{Z}; art.~300$^{5}$ KSH \Wiar{Z}; art.~910 i art.~911$^{3}$ KPC \Wiar{Z}.}
\Uzasadnienie{Derogacja czasowa jest najbezpieczniejszym zastosowaniem §~2: ograniczenie ma termin końcowy, dotyczy tylko Założycieli (którzy sami je przyjmują przy zawiązaniu), a wyjątki z ust.~2 zachowują drogi wyjścia. Osadzenie lock-upu w art.~300$^{39}$ §~1 KSH (a nie w zobowiązaniu obligacyjnym z art.~57 §~2 KC) czyni go skutecznym wobec Spółki i rejestru. Ujawnienie: adnotacja przy akcjach każdego Założyciela Pierwotnego („ograniczenie rozporządzania: zgoda Walnego Zgromadzenia do dnia [data = wpis + 12 miesięcy], §~\ref{art:okres-ograniczenia-zbywania} umowy”); dla akcji serii~F --- przy wpisie emisji, z datą liczoną od Dnia Przyznania, którą Spółka podaje podmiotowi prowadzącemu rejestr. Niespójność: ust.~1 wyłącza §~3--6; art.~300$^{39}$ §~2 literalnie obejmuje także §~6, ale §~6 (egzekucja) chroni wierzycieli akcjonariusza, a jego umowne uchylenie nie wiązałoby komornika sprzedającego zajęte prawo według KPC (art.~910, art.~911$^{3}$) --- derogację należy ograniczyć do §~3--5. Wyjątki z ust.~2 a zgoda Spółki: to dwie niezależne bramki. Dla lit.~a--c zgodę Spółki wyłącza §~\ref{art:zbywanie} ust.~7; „inna runda” z lit.~a nie jest tam objęta, ale objęcie nowych akcji nie jest rozporządzeniem, więc zgoda nie jest wymagana. Lit.~d nie jest wymieniona w ust.~7, więc wymaga zgody Spółki i podlega prawu pierwszeństwa (§~\ref{art:pierwszenstwo}) --- wynik wadliwy, bo prawo pierwszeństwa niweczy sens przeniesienia do własnej spółki; trzeba ją dopisać do wyłączeń w §~\ref{art:zbywanie} ust.~7 i w §~\ref{art:pierwszenstwo}, z zachowaniem weryfikacji Dopuszczalnego Nabywcy i Kryterium.}
\Rekomendacja{Zmienić „§~3--6” na „§~3--5” w §~\ref{art:okres-ograniczenia-zbywania} ust.~1 (§~\ref{art:zbywanie} ust.~4 odsyła tylko do mechanizmu z ust.~3 i nie wymaga zmiany); dopisać przeniesienie do podmiotu w całości kontrolowanego do wyłączeń zgody i prawa pierwszeństwa; ujawnić lock-up w rejestrze z datą końcową.}
\Kierunek{Zmiana usuwa potencjalny zarzut derogacji §~6 i czyni lit.~d wykonalną; dla inwestora --- standard „permitted transfer” założyciela; dla rejestru --- adnotacja z datą. Uzupełnienie §~\ref{art:zbywanie} ust.~7 --- w ramce przy tym ustępie.
\begin{Brzmienie}
(§~\ref{art:okres-ograniczenia-zbywania} ust.~1, zdanie drugie w nowym brzmieniu) „Ograniczenie to stanowi ograniczenie rozporządzenia akcją w rozumieniu art.~300$^{39}$ §~1 KSH; na podstawie art.~300$^{39}$ §~2 KSH strony postanawiają, że w Okresie Ograniczenia Zbywania Akcji przepisy art.~300$^{39}$ §~3--5 KSH oraz mechanizm wskazania innego nabywcy z §~\ref{art:zbywanie} ust.~3 nie znajdują zastosowania, a odmowa zgody nie rodzi po stronie Spółki obowiązku wskazania nabywcy. Spółka ujawnia Okres Ograniczenia Zbywania Akcji wraz z jego datą końcową w rejestrze akcjonariuszy.”
\end{Brzmienie}}
\end{MemoBox}

\item Ograniczenie nie dotyczy:
  \begin{enumerate}[label=\alph*)]
    \item zbycia albo objęcia akcji w ramach Kwalifikowanej Rundy Finansowania (§~\ref{art:kwalifikowana-runda}) albo innej rundy finansowania zatwierdzonej uchwałą Walnego Zgromadzenia podjętą większością 75\% wszystkich głosów;
    \item wykonania prawa przyłączenia albo prawa przymusowego współzbycia;
    \item obowiązkowego zbycia wynikającego z mechanizmu zwrotnego zbycia akcji;
    \item przeniesienia do podmiotu w całości kontrolowanego przez Założyciela Pierwotnego albo Założyciela Funkcjonalnego Serii F, jeżeli podmiot ten jest Dopuszczalnym Nabywcą, przystąpi do wszystkich obowiązków związanych z przenoszonymi akcjami, a zbywający pozostanie solidarnie odpowiedzialny i zapewni zwrotne przeniesienie akcji przed utratą kontroli.
  \end{enumerate}
  
\begin{MemoBox}{3.3.3}
\MemoPyt{3.3.3}{Jak pogodzić Kwalifikowaną Rundę z prawem poboru (§~\ref{art:kwalifikowana-runda} ust.~5, §~\ref{art:walne-zastrzezone} ust.~2) i z wymogiem Kryterium wobec inwestorów, a także z lock-upem założycieli (§~\ref{art:okres-ograniczenia-zbywania} ust.~2 lit.~a)?}
\Ocena{zgodne warunkowo}
\Klasa{zalecane}
\Podstawa{art.~300$^{106}$ §~1--3 i~6 KSH \Wiar{Z} (§~1: prawo poboru przysługuje, „jeżeli umowa spółki lub uchwała akcjonariuszy nie stanowią inaczej”); art.~300$^{39}$ §~1--2 KSH \Wiar{Z}; §~\ref{art:dopuszczalny-nabywca} ust.~5--6, §~\ref{art:zbywanie} ust.~7, §~\ref{art:pierwszenstwo} ust.~1.}
\Uzasadnienie{\emph{Prawo poboru.} Emisja w Rundzie rodzi prawo poboru wszystkich akcjonariuszy; skierowanie akcji do inwestora wymaga pozbawienia go uchwałą 4/5 w interesie Spółki (art.~300$^{106}$ §~1--2 \Wiar{Z}), co §~\ref{art:kwalifikowana-runda} ust.~5 poprawnie zastrzega; przy strukturze 50/20/10/10/10 koalicje 4/5 pokrywają się z koalicjami 75\% (po emisjach serii F i P układ się zmieni). Pro-rata z §~\ref{art:kwalifikowana-runda} ust.~4 lit.~e jest węższa od ustawowego prawa poboru, więc każda kolejna runda wymaga uchwały 4/5 o pozbawieniu go „w części” (stąd rozszerzenie ust.~3, 3.3.1). Prawo poboru w P.S.A. jest dyspozytywne (art.~300$^{106}$ §~1 \Wiar{Z}): umowa może je wyłączyć dla akcji obejmowanych przez inwestorów wskazanych w uchwale kierunkowej albo zastąpić umownym pro-rata --- po zawiązaniu ostrożnie większością 4/5 (per analogiam art.~300$^{113}$ §~3 \Wiar{W}); obecny model ustawowy to wybór Założycieli, nie wymóg prawa. \emph{Kryterium.} Uchwała kierunkowa zapada tą samą większością co zgoda z §~\ref{art:dopuszczalny-nabywca} ust.~6 i może ją zawierać; zdanie z przekreślonego ust.~14 warto zachować; fundusze z inwestorami spoza UE/EOG/NATO --- pakiet 1.4. \emph{Lock-up.} Objęcie akcji nie jest rozporządzeniem, a „zbycie w ramach Rundy” jest niezdefiniowane (Runda to emisja); chodzi zapewne o transakcję wtórną z inwestorem Rundy --- wyjątek sensowny, jeżeli jest objęta uchwałą kierunkową, wyłączona spod prawa pierwszeństwa (§~\ref{art:pierwszenstwo}; decyzja Założycieli) i nie uchyla zwrotnego zbycia akcji niezwolnionych. Pominięty akcjonariusz może zaskarżyć uchwałę o pozbawieniu prawa poboru, jeżeli nie była w interesie Spółki.}
\Rekomendacja{Doprecyzować §~\ref{art:okres-ograniczenia-zbywania} ust.~2 lit.~a i dodać w §~\ref{art:kwalifikowana-runda} ust.~2, że uchwała kierunkowa może zawierać zgodę z §~\ref{art:dopuszczalny-nabywca} ust.~6 oraz zatwierdzenie transakcji wtórnych z inwestorami Rundy. Rozważyć (decyzja Założycieli) wykorzystanie art.~300$^{106}$ §~1 KSH: zapisanie w umowie, że prawo poboru nie przysługuje w zakresie akcji obejmowanych przez inwestorów wskazanych w uchwale kierunkowej, co usuwa odrębną uchwałę 4/5 z każdej rundy; w zamian uchwała kierunkowa (75\%) staje się jedyną blokadą.}
\Kierunek{Zmiana usuwa lukę definicyjną i skraca ścieżkę zgód; nie zmienia większości; dla rejestru akcjonariuszy oznacza jasną podstawę wpisu inwestora nabywającego akcje wtórnie bez procedury §~\ref{art:pierwszenstwo}. Poniżej brzmienie lit.~a; brzmienie §~\ref{art:kwalifikowana-runda} ust.~2 zdanie drugie --- w ramce przy tym ustępie.}
\begin{Brzmienie}
§~\ref{art:okres-ograniczenia-zbywania} ust.~2 lit.~a: „zbycia akcji na rzecz inwestora Kwalifikowanej Rundy Finansowania (§~\ref{art:kwalifikowana-runda}) albo innej rundy finansowania zatwierdzonej uchwałą Walnego Zgromadzenia podjętą większością 75\% wszystkich głosów, jeżeli zbycie to zostało zatwierdzone w uchwale kierunkowej albo w uchwale zatwierdzającej rundę; do takiego zbycia nie stosuje się §~\ref{art:pierwszenstwo}, a akcje objęte obowiązkiem zbycia zgodnie z §~\ref{art:zwrotne-zbycie} mogą być zbyte tylko w zakresie zwolnionym z tego obowiązku;”
\end{Brzmienie}
\end{MemoBox}

\item Upływ Okresu Ograniczenia Zbywania Akcji nie wpływa na dalsze obowiązywanie mechanizmu zwrotnego zbycia akcji ani obowiązku zbycia akcji, które nie zostały jeszcze zwolnione z tego obowiązku.
\end{enumerate}
\end{ContractClause}

\begin{ContractClause}{Prawo pierwszeństwa nabycia akcji}
\label{art:pierwszenstwo}
\begin{enumerate}
  \item Akcjonariusz zamierzający zbyć akcje osobie trzeciej jest zobowiązany najpierw zaoferować je Spółce, a następnie pozostałym akcjonariuszom, na warunkach nie mniej korzystnych dla nabywcy niż uzgodnione z osobą trzecią. Ofertę składa się po uzyskaniu zgody Spółki na rozporządzenie akcją zgodnie z §~\ref{art:zbywanie} albo równocześnie ze zgłoszeniem zamiaru rozporządzenia; w tym drugim przypadku terminy określone w ust.~3 biegną od dnia udzielenia zgody.
  
\begin{MemoBox}{1.3.4}
\MemoPyt{1.3.4}{Kolejność prawa pierwszeństwa (§~14) i zgody Spółki (§~12): czy dopuszczenie oferty równoczesnej ze zgłoszeniem nie koliduje z wymogiem, by oferta odzwierciedlała warunki uzgodnione z nabywcą, i czy 20-dniowe terminy przyjęcia oferty oraz 60-dniowe okno zbycia są wykonalne?}

\Ocena{zgodne warunkowo}
\Klasa{zalecane}
\Podstawa{art.~300$^{39}$ §~1 KSH \Wiar{Z}; art.~66 §~1--2 KC (oferta i termin związania) \Wiar{Z}; art.~300$^{47}$ KSH \Wiar{Z}.}
\Uzasadnienie{Oferta równoczesna nie koliduje z wymogiem odzwierciedlenia warunków, bo zgłoszenie z §~\ref{art:zbywanie} ust.~1 już zawiera cenę i istotne warunki uzgodnione z nabywcą; oferta jest tym samym zestawem warunków skierowanym do innych adresatów. Kolizja pojawia się gdzie indziej: jeżeli Spółka odmówi zgody z powodu Kryterium i uruchomi wskazanie nabywcy (§~\ref{art:zbywanie} ust.~3), równolegle złożona oferta pierwszeństwa wisi w próżni --- projekt nie mówi, czy wygasa, czy wskazany nabywca ma pierwszeństwo przed akcjonariuszami. Trzeba to rozstrzygnąć: odmowa zgody powoduje wygaśnięcie oferty (terminy z ust.~3 nigdy nie zaczynają biec), a wskazanie nabywcy przez Spółkę zastępuje prawo pierwszeństwa. Terminy: 14 dni (zgoda) + 20 (Spółka) + 20 (akcjonariusze) + 60 (okno zbycia) dają maksymalnie ok.~114 dni, co jest rynkowo akceptowalne przy sprzedaży pakietu w spółce prywatnej \Wiar{W}; 60-dniowe okno powinno biec od bezskutecznego upływu terminów z ust.~3, a nie od nieokreślonej chwili „nieprzyjęcia”. Termin 20 dni dla Spółki jest iluzoryczny bez wcześniejszego upoważnienia Walnego Zgromadzenia do nabycia akcji własnych (zob.~1.3.5). Wreszcie ust.~4 wymaga, aby nabywca był „Dopuszczalnym Nabywcą”, ale nie powtarza wymogu Kryterium --- należy go dodać dla spójności z §~\ref{art:zbywanie} ust.~7.}
\Rekomendacja{Doprecyzować: skutek odmowy zgody dla oferty, początek biegu 60 dni, wymóg Kryterium w ust.~4.}
\Kierunek{Zmiany porządkują sekwencję i usuwają ryzyko sporu o pierwszeństwo między wskazanym nabywcą a akcjonariuszami; dla inwestora --- neutralne; dla rejestru --- bez wpływu. Proponowane brzmienie ust.~4 --- w ramce przy tym ustępie.}
\begin{Brzmienie}
(§~\ref{art:pierwszenstwo} ust.~1, zdanie dodawane na końcu) „Odmowa zgody Spółki powoduje wygaśnięcie oferty; jeżeli Spółka wskazuje nabywcę zgodnie z §~\ref{art:zbywanie} ust.~3, wskazanie to zastępuje prawo pierwszeństwa Spółki i akcjonariuszy.”
\end{Brzmienie}
\end{MemoBox}

\item Oferta zawiera wszystkie istotne warunki transakcji, w tym cenę, świadczenia dodatkowe, sposób i termin zapłaty oraz dane potencjalnego nabywcy.
  \item Spółka może przyjąć ofertę w terminie 20 dni, o ile nabycie akcji własnych jest prawnie dopuszczalne. Po bezskutecznym upływie tego terminu pozostali akcjonariusze mogą przyjąć ofertę w terminie kolejnych 20 dni, proporcjonalnie do liczby posiadanych akcji, z prawem objęcia niewykorzystanej części.
  
\begin{MemoBox}{1.3.5}
\MemoPyt{1.3.5}{§~14 ust.~3: Spółka może przyjąć ofertę, „o ile nabycie akcji własnych jest prawnie dopuszczalne” --- czy odesłanie to wystarcza w świetle art.~300$^{47}$ KSH?}

\Ocena{zgodne warunkowo}
\Klasa{zalecane}
\Podstawa{art.~300$^{47}$ §~1 pkt~2, §~2--3 i §~9 KSH (nabycie na podstawie i w granicach upoważnienia udzielonego w uchwale akcjonariuszy; warunki: pełne pokrycie, próg 25\%, celowy kapitał rezerwowy; osoba trzecia działająca na rachunek spółki) \Wiar{Z}; art.~300$^{15}$ §~2 i §~4--6 KSH (kwota dostępna do wypłaty i test wypłacalności) \Wiar{Z}; art.~362 §~1 pkt~8 KSH (dla porównania: pięcioletni limit upoważnienia w spółce akcyjnej) \Wiar{Z}; §~\ref{art:walne-zastrzezone} ust.~1 lit.~j.}
\Uzasadnienie{Odesłanie jest prawnie wystarczające (nie tworzy pozoru uprawnienia sprzecznego z ustawą), ale praktycznie puste. Dla prawa pierwszeństwa właściwy jest art.~300$^{47}$ §~1 pkt~2 KSH: nabycie „na podstawie i w granicach upoważnienia udzielonego w uchwale akcjonariuszy” \Wiar{Z}, wymagające łącznie (§~2) pełnego pokrycia akcji, nieprzekroczenia 25\% wszystkich akcji oraz łącznej ceny z kosztami nie wyższej od kapitału rezerwowego utworzonego w tym celu z kwoty z art.~300$^{15}$ §~2 (zysk i kapitały rezerwowe z zysku); do zapłaty ceny stosuje się odpowiednio art.~300$^{15}$ §~4--6, w tym test wypłacalności. W umowie nabycie jest ponadto Sprawą Zastrzeżoną (75\% wszystkich głosów). Upoważnienie musi pochodzić z uchwały akcjonariuszy, nie z umowy spółki, a jego wykonanie wymaga celowego kapitału rezerwowego z zysku, którego spółka przed przychodami nie ma --- prawo pierwszeństwa Spółki będzie więc przez pierwsze lata martwe niezależnie od redakcji; w 20 dni od oferty uchwała taka jest osiągalna tylko przy uchwale poza zgromadzeniem (§~\ref{art:wz} ust.~4). Bez upoważnienia „stałego” termin Spółki upływa bezużytecznie, a mechanizm sprowadza się do pierwszeństwa akcjonariuszy. Do rozważenia: (a)~uchwała Walnego Zgromadzenia po zawiązaniu, udzielająca Radzie ramowego upoważnienia do nabywania akcji własnych w wykonaniu prawa pierwszeństwa i §~\ref{art:zbywanie} ust.~3, w granicach kwoty i czasu (art.~300$^{47}$ KSH, inaczej niż art.~362 §~1 pkt~8 KSH, nie określa maksymalnego okresu upoważnienia \Wiar{Z}; rozsądne jest jednak wskazanie okresu i kwoty); (b)~wskazanie przez Spółkę w 20 dni samodzielnego nabywcy, niedziałającego na rachunek Spółki (art.~300$^{47}$ §~9).}
\Rekomendacja{Uzupełnić ust.~3 o prawo Spółki do wskazania samodzielnego nabywcy w terminie 20 dni i o odesłanie do upoważnienia Walnego Zgromadzenia; po zawiązaniu podjąć uchwałę upoważniającą.}
\Kierunek{Zmiana czyni prawo pierwszeństwa Spółki realnym bez naruszenia art.~300$^{47}$; dla inwestora --- standard (spółka może wskazać nabywcę); dla rejestru i sądu rejestrowego --- bez wpływu (upoważnienie w uchwale, nie w umowie).}
\begin{Brzmienie}
(§~\ref{art:pierwszenstwo} ust.~3, zdanie pierwsze w nowym brzmieniu) „Spółka może w terminie 20 dni przyjąć ofertę, o ile nabycie akcji własnych jest prawnie dopuszczalne, w szczególności na podstawie i w granicach upoważnienia udzielonego uchwałą Walnego Zgromadzenia zgodnie z art.~300$^{47}$ §~1 pkt~2 i §~2 KSH, albo w tym samym terminie wskazać Dopuszczalnego Nabywcę spełniającego Kryterium Bezpieczeństwa EU/NATO, niedziałającego na rachunek Spółki, który przyjmie ofertę na jej warunkach.”
\end{Brzmienie}
\end{MemoBox}

\item Jeżeli oferta nie zostanie przyjęta w całości, akcjonariusz może przez 60 dni zbyć akcje nabywcy wskazanemu w zgłoszeniu, będącemu Dopuszczalnym Nabywcą, na warunkach nie korzystniejszych niż przedstawione Spółce i akcjonariuszom.
  
\begin{MemoBox}{1.3.4}
\MemoPytCd{1.3.4}
\begin{Brzmienie}
(§~\ref{art:pierwszenstwo} ust.~4 w nowym brzmieniu) „Jeżeli oferta nie zostanie przyjęta w całości, akcjonariusz może w terminie 60 dni od bezskutecznego upływu terminów określonych w ust.~3 zbyć akcje nabywcy wskazanemu w zgłoszeniu, będącemu Dopuszczalnym Nabywcą i spełniającemu Kryterium Bezpieczeństwa EU/NATO albo dopuszczonemu uchwałą Walnego Zgromadzenia zgodnie z §~\ref{art:dopuszczalny-nabywca}, na warunkach nie korzystniejszych niż przedstawione Spółce i akcjonariuszom.”
\end{Brzmienie}
\end{MemoBox}

\item Istotna zmiana warunków albo upływ terminu wymaga ponowienia procedury.
\end{enumerate}

\end{ContractClause}

\begin{ContractClause}{Prawo przyłączenia do zbycia}
\label{art:przylaczenie}
\begin{enumerate}
  \item Jeżeli jeden lub kilku akcjonariuszy zamierza dokonać transakcji, w wyniku której Dopuszczalny Nabywca lub grupa podmiotów działających w porozumieniu uzyska bezpośrednio albo pośrednio ponad 50\% głosów w Spółce, każdy pozostały akcjonariusz ma prawo zażądać objęcia transakcją wszystkich albo części swoich akcji.
  \item Akcje akcjonariusza korzystającego z prawa przyłączenia są zbywane po tej samej cenie za akcję, w tym samym rodzaju świadczenia i na warunkach nie mniej korzystnych niż akcje zbywających akcjonariuszy.
  \item Zawiadomienie o planowanej transakcji doręcza się co najmniej 30 dni przed jej zamknięciem. Uprawniony akcjonariusz wykonuje prawo w terminie 20 dni od doręczenia kompletnego zawiadomienia.
  \item Zbywający nie może zamknąć transakcji, jeżeli nabywca nie nabędzie również akcji prawidłowo zgłoszonych do przyłączenia.
\end{enumerate}

\end{ContractClause}

\begin{ContractClause}{Prawo przymusowego współzbycia}
\label{art:przymusowe-wspolzbycie}
\begin{enumerate}
  \item Akcjonariusz albo grupa akcjonariuszy posiadająca łącznie co najmniej \textbf{75\% wszystkich głosów} może, w związku z wiążącą ofertą nabycia 100\% akcji Spółki przez niezależnego Dopuszczalnego Nabywcę, zobowiązać pozostałych akcjonariuszy do zbycia wszystkich ich akcji.
  \item Prawo może zostać wykonane wyłącznie, gdy:
  \begin{enumerate}[label=\alph*)]
    \item transakcja jest zawierana w dobrej wierze z niepowiązaną osobą trzecią;
    \item każdy akcjonariusz otrzymuje taką samą cenę za akcję tego samego rodzaju i ten sam rodzaj świadczenia;
    \item odpowiedzialność akcjonariusza mniejszościowego ma charakter indywidualny, nie solidarny, jest ograniczona do otrzymanego wynagrodzenia i dotyczy zasadniczo tytułu prawnego do akcji, zdolności oraz upoważnienia do zawarcia transakcji;
    \item ewentualny depozyt, zatrzymanie ceny lub odroczona płatność obciąża akcjonariuszy proporcjonalnie;
    \item nabywca przeszedł pozytywnie pełną weryfikację na podstawie §~\ref{art:dopuszczalny-nabywca} i §~\ref{art:export}.
  \end{enumerate}
  \item Zawiadomienie o wykonaniu prawa doręcza się co najmniej 30 dni przed planowanym zamknięciem transakcji i dołącza pełną dokumentację warunków sprzedaży.
  \item Akcjonariusz zobowiązany do współzbycia jest obowiązany podpisać dokumenty i wykonać czynności konieczne do zamknięcia transakcji, pod warunkiem spełnienia wymogów niniejszego paragrafu.

\begin{MemoBox}{4.4.2}
\MemoPytCd{4.4.2}
\Kierunek{(IV)~Ochrona serii inwestorskiej przy drag-along: zgoda większości akcji serii z~Kwalifikowanej Rundy albo okres ochronny z~ceną minimalną oraz podział ceny według uprzywilejowania z~§~\ref{art:kwalifikowana-runda} ust.~4 lit.~a. Bezpieczeństwo właścicielskie nie słabnie (ust.~2 lit.~e nadal wymaga pełnej weryfikacji nabywcy), a~Założyciele zyskują symetryczną pewność, że po okresie ochronnym inwestor nie zablokuje sprzedaży. Nowy ust.~5:}
\begin{Brzmienie}
Po Kwalifikowanej Rundzie Finansowania prawo przymusowego współzbycia może zostać wykonane wobec akcjonariuszy posiadających akcje serii wyemitowanej w tej Rundzie wyłącznie, jeżeli: (a)~za wykonaniem prawa opowiedzieli się akcjonariusze posiadający większość akcji tej serii, albo (b)~od dnia wpisu tej serii do rejestru akcjonariuszy upłynęło co najmniej \textbf{[48]} miesięcy, a cena za akcję tej serii nie jest niższa niż cena minimalna określona w uchwale emisyjnej tej Rundy, w jej braku --- niż cena emisyjna. Cenę należną akcjonariuszom dzieli się z uwzględnieniem uprzywilejowania ustanowionego zgodnie z §~\ref{art:kwalifikowana-runda} ust.~4 lit.~a; ust.~2 lit.~b stosuje się w ramach akcji tego samego rodzaju. Uchwała emisyjna Kwalifikowanej Rundy Finansowania może określić okres, o którym mowa w lit.~b, jako krótszy.
\end{Brzmienie}
\end{MemoBox}

\end{enumerate}

\end{ContractClause}

\begin{ContractClause}{Dziedziczenie akcji}
\label{art:dziedziczenie}
\begin{enumerate}
  \item Na podstawie art.~300$^{41}$ §~1 KSH wstąpienie do Spółki spadkobierców na miejsce zmarłego akcjonariusza jest ograniczone w sposób określony w niniejszym paragrafie. Spadkobierca wstępuje do Spółki i może wykonywać prawa z akcji po wykazaniu następstwa prawnego oraz przejściu weryfikacji Dopuszczalnego Nabywcy zgodnie z §~\ref{art:dopuszczalny-nabywca}, a jeżeli nie spełnia Kryterium Bezpieczeństwa EU/NATO --- także po uzyskaniu zgody Walnego Zgromadzenia zgodnie z §~\ref{art:dopuszczalny-nabywca}.
  \item Spadkobierca, który nie jest Dopuszczalnym Nabywcą, który nie spełnia Kryterium Bezpieczeństwa EU/NATO i w terminie 90 dni od wezwania nie uzyskał zgody Walnego Zgromadzenia zgodnie z §~\ref{art:dopuszczalny-nabywca}, albo który nie przedstawi wymaganych informacji w terminie 90 dni od wezwania, nie wstępuje do Spółki. Spadkobiercy niewstępującemu przysługuje spłata w wysokości 100\% Wartości Godziwej akcji przypadających mu ze spadku, ustalonej zgodnie z §~\ref{art:wycena} na dzień otwarcia spadku, płatna na zasadach określonych w ust.~4; warunki spłaty określa się na podstawie art.~300$^{41}$ §~1 zdanie drugie KSH. Akcje spadkobiercy niewstępującego podlegają nabyciu przez wskazanego przez Spółkę Dopuszczalnego Nabywcę za cenę równą spłacie; jeżeli mechanizm ograniczenia wstąpienia okazałby się w danym przypadku bezskuteczny, akcje podlegają obowiązkowemu zbyciu na rzecz wskazanego przez Spółkę Dopuszczalnego Nabywcy za Wartość Godziwą.
  
\begin{MemoBox}{1.6.1}
\MemoPyt{1.6.1}{Czy §~\ref{art:dziedziczenie} prawidłowo wykorzystuje art.~300$^{41}$ §~1 KSH (ograniczenie albo wyłączenie wstąpienia pod warunkiem określenia zasad spłaty), a warunki spłaty (100\,\% Wartości Godziwej, raty z §~\ref{art:dziedziczenie} ust.~4) są określone wystarczająco, aby ograniczenie było skuteczne?}
\Ocena{ryzyko --- konstrukcja poprawna, ale warunki spłaty nie wskazują dłużnika ani ścieżki, gdy nabywcy nie ma}
\Klasa{konieczne (K5, K6, K7)}
\Podstawa{art.~300$^{41}$ §~1--2 KSH \Wiar{Z}; art.~300$^{44}$--300$^{46}$ KSH (umorzenie przymusowe i automatyczne, spłata nie niższa od wartości godziwej) \Wiar{Z}; art.~300$^{15}$ §~2 i~4--6 KSH \Wiar{Z}; art.~300$^{47}$ §~1--2 KSH \Wiar{Z}; art.~922 §~1, art.~924--925, art.~1025 §~2, art.~1027 KC \Wiar{Z}; art.~183 §~1 KSH (analogia; brzmienie identyczne) \Wiar{Z} i orzecznictwo do sp.\ z~o.o.\ \Wiar{W}}
\Uzasadnienie{Art.~300$^{41}$ §~1 KSH pozwala umową ograniczyć lub wyłączyć wstąpienie spadkobierców pod warunkiem określenia warunków spłaty niewstępujących „pod rygorem bezskuteczności ograniczenia lub wyłączenia”. Projekt wyraźnie powołuje przepis (ust.~1), określa przesłanki niewstąpienia i wysokość spłaty (100\% Wartości Godziwej na dzień otwarcia spadku) oraz odsyła do rat z ust.~4. Trafnie nie wyłącza dziedziczenia (akcje przechodzą ex lege, art.~922 KC), lecz tylko wstąpienie --- model ugruntowany na tle art.~183 §~1 KSH o identycznym brzmieniu.

Skuteczności zagrażają braki. \textbf{Po pierwsze}, brak dłużnika spłaty: jest nim przyszły, niewskazany nabywca; gdy Rada nikogo nie wskaże albo nabywca nie zapłaci, spadkobierca nie ma od kogo żądać spłaty, a ścieżka rezerwowa jest tautologiczna. „Warunki spłaty” to co najmniej: kto, ile, kiedy; brak „kto” grozi bezskutecznością ograniczenia i wstąpieniem spadkobiercy (także spoza Kryterium). \textbf{Po drugie}, ust.~4 uzależnia raty od zagrożenia płynności, nie mówiąc, kto to ocenia --- warunek spłaty musi być oznaczony bezwarunkowo. \textbf{Po trzecie}, brak terminu wskazania nabywcy i ścieżki, gdy nabywcy nie ma: nabycie akcji własnych (art.~300$^{47}$ §~1 pkt~1--2 i §~2 KSH) albo umorzenie. Umorzenie przymusowe wymaga przesłanki w umowie już przy zawiązaniu (później --- zgody uprawnionego, art.~300$^{45}$ §~1 KSH), spłata pochodzi tylko ze środków podlegających podziałowi, z testem wypłacalności (art.~300$^{44}$ §~4 w zw.\ z art.~300$^{15}$ §~2 i~4--6 KSH), po wpisie umorzenia do rejestru, a przy braku środków umorzenie automatyczne „nie doszło do skutku” (art.~300$^{46}$ KSH). Umorzenie nie może więc być jedyną ścieżką; solidarna odpowiedzialność Spółki i ubezpieczenie z ust.~5 są konieczne. \textbf{Po czwarte}, akcje serii F za wkład pracy niewniesiony w całości: ustawa uzależnia wstąpienie od zgody Spółki, „chyba że umowa spółki stanowi inaczej”, i nakazuje proporcjonalne pomniejszenie spłaty (art.~300$^{41}$ §~1 zdanie trzecie i §~2 KSH); §~\ref{art:dziedziczenie} tego nie odzwierciedla.

Wysokość 100\% Wartości Godziwej bez dyskonta jest bezpieczna konstytucyjnie (art.~21 ust.~1 i art.~64 Konstytucji) i przewyższa standard art.~300$^{45}$ §~2 KSH. Spłata ceną emisyjną za akcje niezwolnione zmarłego Akcjonariusza Funkcjonalnego (§~\ref{art:zwrotne-zbycie} ust.~7 lit.~b) jest skuteczna, ale surowa; rynkowym standardem jest przyspieszenie przy śmierci, które ust.~7 lit.~c pozostawia uznaniu Walnego Zgromadzenia --- to decyzja Założycieli.}
\Rekomendacja{Uzupełnić ust.~2 i ust.~4 o dłużnika spłaty (nabywca, a solidarnie Spółka), termin wskazania nabywcy, bezwarunkowy termin pierwszej raty i ścieżkę rezerwową przez nabycie akcji własnych albo umorzenie, z wyraźnym określeniem niewstąpienia jako przesłanki umorzenia (art.~300$^{45}$ §~1 KSH); dodać regułę dla akcji za wkład pracy (art.~300$^{41}$ §~1 zdanie trzecie i §~2 KSH). Rozważyć automatyczne zwolnienie akcji niezwolnionych przy śmierci (decyzja Założycieli).}
\begin{Brzmienie}
§~\ref{art:dziedziczenie} ust.~2 zdania trzecie i czwarte (nowe brzmienie): „Akcje spadkobiercy niewstępującego nabywa wskazany przez Spółkę w terminie 3 miesięcy od dnia stwierdzenia niewstąpienia Dopuszczalny Nabywca spełniający Kryterium Bezpieczeństwa EU/NATO, za cenę równą spłacie; dłużnikiem spłaty jest nabywca, a Spółka odpowiada za spłatę solidarnie. Jeżeli Spółka nie wskaże nabywcy w tym terminie albo nabywca nie zapłaci pierwszej raty w terminie, Spółka nabywa akcje spadkobiercy niewstępującego w granicach art.~300$^{47}$ KSH albo --- jeżeli dysponuje środkami, o których mowa w art.~300$^{15}$ §~2 KSH --- umarza je bez zgody spadkobiercy za spłatą równą Wartości Godziwej; stwierdzone niewstąpienie spadkobiercy stanowi przesłankę umorzenia przymusowego w rozumieniu art.~300$^{45}$ §~1 KSH, a spłatę uiszcza się po wpisie umorzenia do rejestru; niezależnie od wybranej ścieżki Spółka odpowiada za spłatę solidarnie. Spadkobierca niewstępujący jest zobowiązany współdziałać przy przeniesieniu akcji, a do wykonania tego obowiązku stosuje się odpowiednio §~\ref{art:zwrotne-zbycie} ust.~14--15 i §~\ref{art:breach}. Wstąpienie spadkobierców akcjonariusza uprawnionego z akcji objętych za wkład w postaci świadczenia pracy lub usług podlega niniejszemu paragrafowi w miejsce zgody Spółki, o której mowa w art.~300$^{41}$ §~2 KSH; spłata za takie akcje, jeżeli wkład nie został w całości wniesiony, ulega pomniejszeniu stosownie do stosunku wartości wkładu wniesionego do wartości wkładu niewniesionego zgodnie z art.~300$^{41}$ §~1 zdanie trzecie KSH.”
\end{Brzmienie}
\end{MemoBox}

\begin{MemoBox}{1.6.2}
\MemoPyt{1.6.2}{Czy uzależnienie wstąpienia od Kryterium i zgody Walnego Zgromadzenia jest dopuszczalne wobec spadkobierców (w tym ustawowych), także z perspektywy konstytucyjnej ochrony dziedziczenia?}
\Ocena{zgodne warunkowo}
\Klasa{zalecane}
\Podstawa{art.~300$^{41}$ §~1 KSH \Wiar{Z}; art.~21 ust.~1, art.~31 ust.~3, art.~64 ust.~1--3 Konstytucji RP \Wiar{Z}; art.~63 TFUE (dziedziczenie jako przepływ kapitału w orzecznictwie TSUE) \Wiar{W}; art.~18 TFUE \Wiar{W}; art.~58 §~2 KC \Wiar{Z}}
\Uzasadnienie{Art.~300$^{41}$ §~1 KSH nie ogranicza przesłanek, od których umowa może uzależnić wstąpienie; Kryterium Bezpieczeństwa EU/NATO i zgoda Walnego Zgromadzenia mieszczą się w tej swobodzie, także wobec spadkobierców ustawowych, bo ustawa nie różnicuje. Ochrona konstytucyjna dziedziczenia (art.~21 ust.~1, art.~64 Konstytucji) obejmuje wartość majątkową spadku, a nie przynależność do korporacji: spadkobierca otrzymuje 100\% Wartości Godziwej bez dyskonta, a ograniczenie ma podstawę ustawową (art.~64 ust.~3) i cel (bezpieczeństwo struktury właścicielskiej spółki obronnej), co spełnia test proporcjonalności z art.~31 ust.~3 --- tym bardziej, że spadkodawca sam zawarł umowę. Nabycie w drodze spadku jest przepływem kapitału (art.~63 TFUE), ale Kryterium obejmuje wszystkie państwa UE/EOG, więc nie różnicuje obywateli Unii (państwa trzecie --- zob.~1.1.1). Warunki oceny pozytywnej: (1)~określona spłata (1.6.1); (2)~zgoda nie może być czysto uznaniowa --- odmowa bez uzasadnienia wobec małżonka i dzieci założyciela spełniających Kryterium byłaby nadużyciem (art.~58 §~2 KC); spadkobierca spełniający Kryterium i niebędący Niedopuszczalnym Nabywcą wstępuje bez zgody (ust.~1 --- prawidłowo); (3)~próg 75\% „wszystkich głosów” liczy też głosy z akcji spadkowych, których nikt nie wykonuje --- przy dużym pakiecie zmarłego (50\% u Założyciela~A) zgoda staje się arytmetycznie niemożliwa. Alternatywa łagodniejsza (\emph{zmiany.md}, ZMIANA~14: spadkobiercy poza Kryterium) jest dopuszczalna, ale sprzeczna z celem programów finansowania obronnego --- rekomendujemy wariant przyjęty; wybór należy do Założycieli.}
\Rekomendacja{Utrzymać Kryterium wobec spadkobierców; w §~\ref{art:dziedziczenie} ust.~2 zastrzec, że próg zgody liczy się z wyłączeniem akcji objętych spadkiem, i dodać wymóg uzasadnienia odmowy.}
\Kierunek{Dopisek w §~\ref{art:dziedziczenie} ust.~2: „Przy obliczaniu większości wymaganej do wyrażenia zgody nie uwzględnia się głosów z akcji wchodzących w skład spadku; uchwała odmawiająca zgody wymaga uzasadnienia.” Dla Spółki: zgoda pozostaje wykonalna; dla spadkobierców: ochrona przed arbitralnością; dla sądu: kryterium kontroli uchwały.}
\end{MemoBox}

\begin{MemoBox}{1.6.3}
\MemoPyt{1.6.3}{Czy 90-dniowe terminy i mechanizm rezerwowy (obowiązkowe zbycie na rzecz wskazanego nabywcy) są wykonalne wobec spadkobiercy, który nie współdziała; kto i jak stwierdza „niewstąpienie”?}
\Ocena{ryzyko}
\Klasa{konieczne (K8)}
\Podstawa{art.~300$^{41}$ §~1, art.~300$^{34}$ §~1 i~4, art.~300$^{37}$ §~2, art.~300$^{38}$ §~1, art.~300$^{44}$--300$^{46}$, art.~300$^{47}$ KSH \Wiar{Z}; art.~1025 §~2, art.~1026, art.~1027, art.~1035--1036 KC \Wiar{Z}; art.~64 KC, art.~1047 KPC \Wiar{Z}; art.~666--667 KPC (kurator spadku) \Wiar{Z}}
\Uzasadnienie{Projekt ma dwie słabości. \textbf{Bieg terminów:} 90 dni liczy się „od wezwania”, ale Spółka często nie zna spadkobierców, a ci nie muszą się ujawnić przed stwierdzeniem nabycia spadku albo poświadczeniem dziedziczenia (art.~1025 §~2, art.~1027 KC), co trwa miesiącami (art.~1026 KC); do tego czasu akcje figurują na zmarłym (art.~300$^{37}$ §~2, art.~300$^{38}$ §~1 KSH), a Spółka nie może ani wezwać, ani stwierdzić niewstąpienia. Potrzebne jest wezwanie znanych spadkobierców, rezerwowo na ostatni adres zmarłego, oraz prawo Spółki do wniosku o stwierdzenie nabycia spadku (art.~1025 §~1 KC) albo o kuratora spadku (art.~666--667 KPC). \textbf{Mechanizm rezerwowy} „obowiązkowego zbycia” wymaga oświadczenia spadkobiercy; bez współdziałania pozostaje art.~64 KC, czyli lata sporu, gdy spadkobierca spoza Kryterium pozostaje właścicielem akcji spółki obronnej (choć bez praw korporacyjnych, jeżeli ograniczenie jest skuteczne). Bez współdziałania działa tylko umorzenie przymusowe albo automatyczne (art.~300$^{44}$--300$^{46}$ KSH): wymaga przesłanki w umowie (stwierdzone niewstąpienie), spłaty nie niższej od wartości godziwej po wpisie umorzenia do rejestru i środków podlegających podziałowi (art.~300$^{44}$ §~4 KSH); w ich braku umorzenie „nie doszło do skutku” i pozostaje nabycie akcji własnych albo art.~64 KC. Domyka to też „warunki spłaty” z 1.6.1. \textbf{Stwierdzenie niewstąpienia:} projekt milczy; powinna to być uchwała Rady po upływie terminu, doręczona spadkobiercy, z prawem sprzeciwu w trybie §~\ref{art:final} ust.~4, stanowiąca z dowodem doręczenia podstawę wpisu wzmianki w rejestrze.}
\Rekomendacja{Dodać do §~\ref{art:dziedziczenie}: tryb wezwania i bieg terminu od wykazania następstwa lub od doręczenia na ostatni adres, uprawnienie Spółki do wszczęcia postępowania spadkowego lub o kuratora, stwierdzenie niewstąpienia uchwałą Rady z doręczeniem, oraz umorzenie jako ścieżkę ostateczną.}
\begin{Brzmienie}
§~\ref{art:dziedziczenie} --- nowy ust.~2a: „Wezwanie, o którym mowa w ust.~2, Spółka kieruje do spadkobierców znanych Spółce, a w razie ich nieznajomości --- na ostatni znany adres zmarłego akcjonariusza; termin 90 dni biegnie od doręczenia wezwania, nie wcześniej jednak niż od dnia, w którym spadkobierca mógł wykazać następstwo prawne postanowieniem o stwierdzeniu nabycia spadku albo aktem poświadczenia dziedziczenia. Spółka może jako zainteresowany wystąpić o stwierdzenie nabycia spadku albo o ustanowienie kuratora spadku. Niewstąpienie spadkobiercy stwierdza Rada Dyrektorów uchwałą, którą doręcza spadkobiercy wraz z informacją o wysokości i terminie spłaty; uchwała stanowi podstawę żądania przez Spółkę wpisu odpowiedniej wzmianki w rejestrze akcjonariuszy. Jeżeli spadkobierca niewstępujący nie współdziała przy przeniesieniu akcji w terminie 30 dni od wezwania, akcje te ulegają umorzeniu bez zgody spadkobiercy i bez uchwały akcjonariuszy (art.~300$^{46}$ KSH) za spłatą równą Wartości Godziwej, nie niższą od wartości godziwej w rozumieniu art.~300$^{45}$ §~2 KSH, płatną po wpisie umorzenia do rejestru; Rada Dyrektorów niezwłocznie podejmuje uchwałę stwierdzającą umorzenie i zgłasza je do rejestru, a jeżeli Spółka nie dysponuje środkami na pełną spłatę --- uchwałę stwierdzającą, że umorzenie nie doszło do skutku, i korzysta z uprawnień z §~\ref{art:zwrotne-zbycie} ust.~15 oraz z nabycia akcji własnych w granicach art.~300$^{47}$ KSH.”
\end{Brzmienie}
\end{MemoBox}

\item Jeżeli akcje przypadają kilku spadkobiercom, są oni zobowiązani wskazać jednego wspólnego przedstawiciela do wykonywania praw korporacyjnych do czasu działu spadku.
  
\begin{MemoBox}{1.6.4}
\MemoPyt{1.6.4}{Wspólny przedstawiciel (§~\ref{art:dziedziczenie} ust.~3) --- spójność z przepisami KSH o współuprawnionych z akcji.}
\Ocena{zgodne}
\Klasa{bez zmian}
\Podstawa{art.~300$^{38}$ §~3--4 KSH \Wiar{Z}; art.~300$^{41}$ §~3 KSH \Wiar{Z}; art.~300$^{33}$ §~1 pkt~5 KSH w brzmieniu od 18.02.2027~r.\ (Dz.U. 2026 poz.~176) \Wiar{Z}; art.~1035--1036 KC \Wiar{Z}}
\Uzasadnienie{Ust.~3 powtarza ustawowy obowiązek wykonywania praw przez wspólnego przedstawiciela (art.~300$^{38}$ §~3--4 KSH) i ogranicza go „do czasu działu spadku”, co odpowiada wspólności majątku spadkowego (art.~1035 KC); kolizji nie ma, a postanowienie ma wartość informacyjną dla spadkobierców. Użyteczne, choć niekonieczne byłoby dopowiedzenie, że do wskazania przedstawiciela prawa z akcji nie są wykonywane, a Spółka może doręczać któremukolwiek ze spadkobierców, oraz odpowiednie stosowanie reguły do innych przypadków współuprawnienia; można też rozważyć wyłączenie albo ograniczenie podziału akcji między spadkobierców (art.~300$^{41}$ §~3 KSH); od 18.02.2027~r.\ rejestr będzie ujawniał wszystkich współwłaścicieli akcji (art.~300$^{33}$ §~1 pkt~5 KSH w nowym brzmieniu), więc wspólny przedstawiciel nie zastąpi wpisu każdego ze spadkobierców.}
\Rekomendacja{Bez zmian; opcjonalnie uzupełnić o zdanie o doręczeniach i o odpowiednim stosowaniu do innych przypadków współuprawnienia.}
\end{MemoBox}

\item Cena może zostać rozłożona na nie więcej niż 36 równych miesięcznych rat, jeżeli jednorazowa zapłata zagrażałaby płynności Spółki lub nabywcy; niespłacona część jest oprocentowana według stopy referencyjnej NBP powiększonej o 3 punkty procentowe. Pierwsza rata jest płatna najpóźniej w dniu przeniesienia akcji, a kolejne w równych odstępach miesięcznych. Nabywca może dokonać wcześniejszej spłaty bez dodatkowej opłaty i ustanawia w umowie zbycia zabezpieczenie pozostałej ceny adekwatne do jej wysokości i ryzyka. Opóźnienie przekraczające 30 dni w zapłacie dwóch rat, nieusunięte w terminie 14 dni od pisemnego wezwania, powoduje wymagalność całej pozostałej ceny.
  
\begin{MemoBox}{1.6.1}
\MemoPytCd{1.6.1}
\begin{Brzmienie}
§~\ref{art:dziedziczenie} ust.~4 zdanie pierwsze (nowe brzmienie): „Spłata jest płatna w terminie 30 dni od dnia ostatecznego ustalenia Wartości Godziwej, nie później niż w dniu przeniesienia akcji; na żądanie nabywcy albo Spółki, zgłoszone przed upływem tego terminu, spłata może zostać rozłożona na nie więcej niż 36 równych miesięcznych rat, przy czym pierwsza rata jest płatna w tym terminie, a niespłacona część jest oprocentowana według stopy referencyjnej NBP powiększonej o 3 punkty procentowe.”
\end{Brzmienie}
\end{MemoBox}

\item Spółka może zawrzeć ubezpieczenie na życie kluczowych Założycieli w celu sfinansowania nabycia akcji od spadkobierców.
\end{enumerate}
\end{ContractClause}

\begin{ContractClause}{Stosunki majątkowe małżeńskie i akcje}
\label{art:malzonek}
\begin{enumerate}
  \item Każdy akcjonariusz będący osobą fizyczną zobowiązuje się, w granicach prawnie dopuszczalnych, niezwłocznie informować Spółkę o zmianie ustroju majątkowego albo innym zdarzeniu mogącym wpływać na przynależność akcji, ich podział lub sposób wykonywania praw z akcji.
  \item \BoilerplateStrike{Nabycie albo ustalenie przynależności akcji w związku ze stosunkami majątkowymi małżeńskimi następuje zgodnie z bezwzględnie obowiązującymi przepisami prawa.} Sam fakt objęcia akcji wspólnością majątkową małżeńską albo ustania tej wspólności nie powoduje sam przez się, wobec Spółki, nabycia przez drugiego małżonka statusu akcjonariusza ani samodzielnego prawa wykonywania praw korporacyjnych. Jeżeli zgodnie z prawem powstanie współuprawnienie z akcji, prawa wobec Spółki wykonuje się zgodnie z właściwymi przepisami KSH, w szczególności art. 300$^{38}$ KSH.
  \item Jeżeli wskutek prawomocnego orzeczenia sądu, ugody, umowy o podział majątku albo innego prawnie skutecznego zdarzenia małżonek albo były małżonek skutecznie nabędzie akcje, stosuje się do niego na zasadach ogólnych postanowienia niniejszej umowy dotyczące Dopuszczalnego Nabywcy i Kryterium Bezpieczeństwa EU/NATO (w tym wymóg zgody Walnego Zgromadzenia zgodnie z §~\ref{art:dopuszczalny-nabywca}, jeżeli Kryterium nie jest spełnione) oraz --- jeżeli po skutecznym nabyciu powstanie odpowiednia podstawa prawna --- Prawnie Niedopuszczalnego Posiadacza. Niniejszy paragraf nie usuwa skutków zakazu wynikającego z bezwzględnie obowiązujących przepisów prawa i nie stanowi samodzielnej podstawy nabycia akcji.
  
\begin{MemoBox}{1.7.1}
\MemoPyt{1.7.1}{Czy rozdzielenie sfery majątkowej (przynależność akcji do majątku wspólnego) od korporacyjnej (wykonywanie praw z akcji przez akcjonariusza wpisanego do rejestru) jest ujęte prawidłowo na tle KRO i KSH, i czy potrzebny jest osobny mechanizm wykupu udziału małżonka?}
\Ocena{zgodne warunkowo}
\Klasa{zalecane}
\Podstawa{art.~31 §~1--2, art.~33 pkt~2, 9 i~10, art.~43, art.~46, art.~47$^{1}$ KRO \Wiar{Z}; art.~300$^{37}$--300$^{38}$ KSH \Wiar{Z}; SN III CZP 109/22, III CSKP 65/21, III CZP 32/16 \Wiar{Z} (tezy według \emph{psa\_feedback.tex}; teksty orzeczeń niedostępne); art.~1035 KC w zw.\ z art.~212 KC \Wiar{Z}}
\Uzasadnienie{Ujęcie jest prawidłowe i odpowiada linii SN cytowanej w \emph{psa\_feedback.tex}: akcje nabyte ze środków wspólnych mogą należeć do majątku wspólnego (III CZP 32/16, analogicznie), lecz akcjonariuszem jest małżonek, który je objął i jest wpisany do rejestru (III CSKP 65/21), a rozdzielność nie daje drugiemu małżonkowi praw korporacyjnych (III CZP 109/22). Tak stanowi ust.~2. Przynależność rozstrzyga się przy objęciu: akcje za wkład z majątku osobistego (prawa twórcy i ich surogaty, art.~33 pkt~9--10 KRO) wchodzą do majątku osobistego, za gotówkę wspólną --- do wspólnego; ponieważ Załącznik nr~1 miesza wkłady, założyciele w ustroju wspólności powinni przy zawiązaniu złożyć oświadczenie o źródle wkładu albo rozważyć umowę majątkową. Osobny mechanizm wykupu udziału małżonka nie jest potrzebny: sąd może przy podziale przyznać akcje akcjonariuszowi ze spłatą (art.~46 KRO w zw.\ z art.~1035 i~212 KC), a ust.~4 zobowiązuje do dążenia do takiego rozwiązania. Brakuje skutku przyznania akcji byłemu małżonkowi spoza Kryterium: nabycie z orzeczenia nie wymaga zgody Spółki, więc nie ma czego odmówić, a §~\ref{art:dopuszczalny-nabywca} ust.~15 dotyczy tylko „utraty” Kryterium po nabyciu. Wobec braku w dziale o P.S.A.\ odpowiednika art.~183$^{1}$ KSH (art.~332$^{1}$ KSH uchylono) ust.~2 opiera się na art.~300$^{38}$ §~1 KSH --- podstawie wystarczającej, ale słabszej od wyraźnego upoważnienia ustawowego; tym bardziej skutek nabycia z mocy orzeczenia trzeba uregulować wprost.}
\Rekomendacja{Dodać w §~\ref{art:malzonek} ust.~3 odesłanie do §~\ref{art:dopuszczalny-nabywca} ust.~15 stosowanego odpowiednio do małżonka, który nabył akcje z mocy orzeczenia lub podziału bez spełnienia Kryterium; przy zawiązaniu odebrać oświadczenia o źródle wkładów.}
\begin{Brzmienie}
§~\ref{art:malzonek} ust.~3 --- nowe zdanie po zdaniu pierwszym: „Jeżeli małżonek albo były małżonek nabył akcje w sposób niewymagający zgody Spółki, a nie spełnia Kryterium Bezpieczeństwa EU/NATO i Walne Zgromadzenie nie wyraziło zgody w terminie 90 dni od wezwania, stosuje się odpowiednio §~\ref{art:dopuszczalny-nabywca} ust.~15, przy czym termin 6 miesięcy biegnie od dnia skutecznego nabycia.”
\end{Brzmienie}
\end{MemoBox}

\item Jeżeli przed skutecznym nabyciem akcji przez małżonka albo byłego małżonka możliwe jest zgodne z prawem rozliczenie pieniężne bez przenoszenia akcji, akcjonariusz zobowiązuje się, w granicach prawnie dopuszczalnych, dążyć w pierwszej kolejności do takiego rozwiązania.
  \item Akcjonariusz jest zobowiązany, na uzasadnione żądanie Spółki i w zakresie dopuszczalnym przez prawo, przedstawić dokumenty niezbędne do ustalenia wpływu stosunków majątkowych małżeńskich na prawa do akcji.

\begin{MemoBox}{1.7.2}
\MemoPyt{1.7.2}{Czy obowiązki informacyjne z §~\ref{art:malzonek} ust.~1 i~5 są egzekwowalne i zgodne z ochroną danych?}
\Ocena{zgodne warunkowo}
\Klasa{opcjonalne}
\Podstawa{art.~300$^{1}$ §~3 KSH \Wiar{Z}; art.~47$^{1}$ KRO \Wiar{Z}; art.~5 ust.~1 lit.~c, art.~6 ust.~1 lit.~c i~f, art.~13--14 RODO \Wiar{W}; art.~300$^{33}$ KSH (zakres danych w rejestrze) \Wiar{Z}}
\Uzasadnienie{Obowiązki są obowiązkami związanymi z akcją (art.~300$^{1}$ §~3 KSH); ich naruszenie uruchamia §~\ref{art:breach}, ale bez sankcji swoistej egzekwowalność ogranicza się do roszczeń o wykonanie i o naprawienie szkody --- to wystarcza, bo celem jest wiedza Spółki, a nie karanie. Spółka ma realny interes: małżonek może powołać się wobec niej na umowę majątkową tylko wtedy, gdy jej zawarcie oraz rodzaj były jej wiadome (art.~47$^{1}$ KRO), a wiedza o postępowaniu o podział majątku pozwala uruchomić ust.~4. Na tle RODO Spółka jest administratorem danych akcjonariusza (art.~6 ust.~1 lit.~c --- rejestr, lit.~f --- uzasadniony interes w ustaleniu uprawnionego z akcji i Kryterium) oraz danych małżonka jako osoby trzeciej (lit.~f; obowiązek informacyjny z art.~14 RODO). Dane o stanie cywilnym i ustroju majątkowym nie są danymi szczególnej kategorii. Warunkiem zgodności jest minimalizacja (art.~5 ust.~1 lit.~c RODO): fakt i data umowy majątkowej, sentencja orzeczenia, wskazanie, czy akcje objęto ze środków wspólnych --- nie pełna treść intercyzy. Ust.~5 jest właściwie zawężony, a klauzule „w granicach prawnie dopuszczalnych” wystarczają. Warunek spełnia się poza umową: polityka przetwarzania danych akcjonariuszy i ich małżonków oraz wzór klauzuli informacyjnej.}
\Rekomendacja{Bez zmian w umowie; w dokumentacji korporacyjnej przygotować klauzulę informacyjną RODO dla akcjonariuszy i ich małżonków oraz wykaz dokumentów, których Spółka może żądać.}
\Kierunek{Ewentualny dopisek w ust.~5: „w zakresie ograniczonym do informacji niezbędnych, z poszanowaniem przepisów o ochronie danych osobowych”; zmiana jest kosmetyczna, istotna jest polityka RODO.}
\end{MemoBox}

\end{enumerate}
\end{ContractClause}



\begin{ContractClause}{Wartość Godziwa i Dzień Wyceny}
\label{art:wycena}
\begin{enumerate}
  \item „Dzień Wyceny” oznacza dzień wystąpienia zdarzenia powodującego powstanie obowiązku zbycia albo prawa nabycia akcji, chyba że właściwe postanowienie niniejszej umowy wyraźnie wskazuje inny dzień.
  \item „Wartość Godziwa” oznacza wartość rynkową akcji ustaloną na Dzień Wyceny przy założeniu kontynuacji działalności Spółki, bez dyskonta z tytułu braku płynności lub pakietu mniejszościowego, chyba że niniejsza umowa wyraźnie stanowi inaczej.
  
\begin{MemoBox}{1.8.1}
\MemoPyt{1.8.1}{Czy umowna definicja może być stosowana we wszystkich mechanizmach projektu (zwrotne zbycie, dziedziczenie, utrata Kryterium, wskazanie nabywcy), czy w niektórych ustawa narzuca własną metodę lub minimalny poziom spłaty?}
\Ocena{zgodne}
\Klasa{opcjonalne}
\Podstawa{art.~353$^{1}$ KC \Wiar{Z}; art.~300$^{39}$ §~2--5 KSH \Wiar{Z}; art.~300$^{41}$ §~1 KSH \Wiar{Z}; art.~300$^{45}$ §~2 KSH \Wiar{Z}; art.~300$^{49}$ §~2 w zw.\ z art.~266 §~3 oraz art.~300$^{50}$ §~3 KSH (wyłączenie i ustąpienie akcjonariusza przez sąd) \Wiar{Z}; art.~28 ust.~6 ustawy o rachunkowości (definicja wartości godziwej) \Wiar{Z}}
\Uzasadnienie{W żadnym z czterech mechanizmów ustawa nie narzuca własnej metody. Zwrotne zbycie i zbycie po utracie Kryterium (§~\ref{art:dopuszczalny-nabywca} ust.~15) to umowne obowiązki sprzedaży --- cena jest przedmiotem swobody umów (art.~353$^{1}$ KC) i może być nawet niższa od rynkowej (cena emisyjna, 80\%). Przy odmowie zgody na zbycie ustawa nie ma reguły cenowej, lecz sankcję za jej brak w umowie --- swobodne zbycie (art.~300$^{39}$ §~3 KSH); ten sam skutek wywołuje niewskazanie nabywcy albo niezapłacenie ceny (art.~300$^{39}$ §~5 KSH). Termin 3 miesięcy w §~\ref{art:zbywanie} ust.~3 zamiast ustawowego miesiąca jest dopuszczalny (art.~300$^{39}$ §~2 KSH), ale warto to zaznaczyć wprost. Przy dziedziczeniu ustawa wymaga „określenia warunków spłaty”, nie poziomu; 100\% Wartości Godziwej to poziom najwyższy, z korektą dla wkładu niewniesionego (art.~300$^{41}$ §~1 zdanie trzecie KSH; seria F). Minimum ustawowe (spłata „nie może być niższa od wartości godziwej akcji”) występuje tylko przy umorzeniu przymusowym (art.~300$^{45}$ §~2 KSH); ustawowa wartość godziwa (por.\ art.~28 ust.~6 ustawy o rachunkowości) dopuszcza dyskonta, więc umowna Wartość Godziwa zawsze spełnia minimum. Rozbieżność dotyczy tylko Dnia Wyceny: przy umorzeniu warto zastrzec datę uchwały, jeżeli jest korzystniejsza dla akcjonariusza. Sąd ustala cenę sam tylko przy wyłączeniu i ustąpieniu akcjonariusza (art.~300$^{49}$ §~2 w zw.\ z art.~266 §~3, art.~300$^{50}$ §~3 KSH); projekt tych trybów nie reguluje i nie musi.}
\Rekomendacja{Bez zmian merytorycznych; opcjonalnie dopisać, że Wartość Godziwa stanowi także wynagrodzenie za umorzenie przymusowe, nie niższe niż wymagane ustawą.}
\Kierunek{Dopisek w §~\ref{art:wycena} ust.~2: „Wartość Godziwa stanowi także wynagrodzenie za umorzenie akcji bez zgody akcjonariusza, jeżeli niniejsza umowa je przewiduje; spłata z tego tytułu w żadnym przypadku nie może być niższa od wartości godziwej akcji w rozumieniu art.~300$^{45}$ §~2 KSH.” Skutek: jedna definicja obsługuje także ścieżkę umorzeniową; sąd rejestrowy widzi wyraźne spełnienie minimum ustawowego.}
\end{MemoBox}

\item W pierwszej kolejności strony mogą uzgodnić wartość w terminie 15 dni od Dnia Wyceny albo od dnia, w którym dowiedziały się o zdarzeniu powodującym konieczność wyceny.
  \item W razie braku uzgodnienia wartość ustala Niezależny Ekspert posiadający doświadczenie w wycenie spółek technologicznych i przemysłowych oraz własności intelektualnej, wybrany przez Radę Dyrektorów spośród podmiotów niepowiązanych ze stronami.
  
\begin{MemoBox}{1.8.2}
\MemoPyt{1.8.2}{Czy wiążący charakter wyceny eksperta (§~\ref{art:wycena} ust.~7) i tryb jego wyboru wytrzymają spór sądowy między wspólnikami?}
\Ocena{ryzyko --- wiążący charakter tak, wybór eksperta przez Radę Dyrektorów nie}
\Klasa{konieczne (K10, K11)}
\Podstawa{art.~353$^{1}$, art.~58 §~2, art.~536 §~1 KC \Wiar{Z}; art.~1157 pkt~1, art.~1161 §~1--2 KPC \Wiar{Z}; art.~278 §~1 KPC \Wiar{Z}; art.~300$^{55}$ §~1 KSH \Wiar{Z}}
\Uzasadnienie{Klauzula ekspercka to dopuszczalna umowa o oznaczenie świadczenia przez osobę trzecią (art.~353$^{1}$, art.~536 §~1 KC), a nie zapis na sąd polubowny: ekspert ustala fakt, więc art.~1157 i~1161 §~1 KPC nie mają zastosowania, a droga sądowa (sąd powszechny, §~\ref{art:final} ust.~4) pozostaje otwarta. Sąd nie zastąpi wyceny własną z powodu samej niezgody strony; bada, czy ustalenie mieści się w granicach umowy i nie jest rażąco niesłuszne, a wadę wykazuje strona kwestionująca. Wyjątki z ust.~7 są dobrze dobrane; warto dodać czwarty --- oczywistą sprzeczność z zasadami z ust.~2 i~5 (np.\ dyskonto). Sąd powoła biegłego (art.~278 §~1 KPC), więc wycena musi być udokumentowana na poziomie opinii biegłego. Słaby punkt to wybór eksperta przez Radę Dyrektorów: Spółka albo pozostali założyciele są ekonomicznie po jednej stronie sporu (Spółka wskazuje nabywcę, założyciele-dyrektorzy są potencjalnymi nabywcami, przy Odejściu Zawinionym Rada je stwierdziła); sąd może uznać, że ustalenie ceny powierzono w istocie jednej stronie, co czyni klauzulę wadliwą (art.~58 §~2 KC; por.\ art.~1161 §~2 KPC) --- to najczęstszy sposób podważania wycen. Ust.~4 powinien być spójny z §~\ref{art:impas} ust.~4 (wybór wspólny), a w braku porozumienia odsyłać do instytucji neutralnej, nie Rady. Procedura powinna obejmować stanowiska liczbowe stron (bez nich ust.~6 jest niestosowalny), projekt wyceny z 14 dniami na uwagi, termin końcowy 60 dni i zaliczkę kosztów przez Spółkę.}
\Rekomendacja{Zmienić tryb wyboru na wspólny, z rezerwowym wskazaniem przez instytucję neutralną (do wyboru Założycieli: prezes właściwego sądu okręgowego, Krajowa Izba Biegłych Rewidentów, sąd arbitrażowy przy izbie gospodarczej); dodać procedurę kontradyktoryjną i czwarty wyjątek w ust.~7.}
\begin{Brzmienie}
§~\ref{art:wycena} ust.~4 (nowe brzmienie): „W razie braku uzgodnienia wartość ustala Niezależny Ekspert posiadający doświadczenie w wycenie spółek technologicznych i przemysłowych oraz własności intelektualnej, niepowiązany ze Spółką, z żadną ze stron ani z Dopuszczalnym Nabywcą, wybrany wspólnie przez strony w terminie 15 dni od bezskutecznego upływu terminu z ust.~3, a w braku porozumienia --- wskazany na wniosek którejkolwiek ze stron przez [instytucja neutralna --- do uzupełnienia]. Każda ze stron przedstawia Niezależnemu Ekspertowi w terminie 14 dni od jego wyboru swoje stanowisko co do wartości wraz z uzasadnieniem i udostępnia mu dane niezbędne do wyceny; Spółka udostępnia dane obu stronom na równych zasadach. Niezależny Ekspert doręcza stronom projekt wyceny, a po rozpatrzeniu uwag zgłoszonych w terminie 14 dni --- wycenę ostateczną, nie później niż w terminie 60 dni od wyboru. Koszty wyceny zalicza Spółka, z prawem rozliczenia zgodnie z ust.~6.”
\end{Brzmienie}
\end{MemoBox}

\item Niezależny Ekspert uwzględnia w szczególności przychody, wynik EBITDA i przepływy pieniężne, portfel zamówień i kontrakty, aktywa produkcyjne, zapasy, certyfikaty i zdolności wytwórcze, portfel własności intelektualnej i poziom gotowości technologicznej, ryzyka regulacyjne, zadłużenie, środki pieniężne, perspektywy komercjalizacji oraz ostatnią rynkową rundę finansowania albo inną transakcję na akcjach Spółki, jeżeli miała miejsce.
  \item Koszt wyceny ponoszą strony po połowie, chyba że różnica między stanowiskiem jednej strony a wyceną przekroczy 20\%; w takim przypadku koszt ponosi ta strona.
  \item Ustalenie Niezależnego Eksperta jest wiążące, z wyjątkiem oczywistego błędu rachunkowego, naruszenia zasad niezależności albo rażącego pominięcia istotnych danych.

\begin{MemoBox}{1.8.2}
\MemoPytCd{1.8.2}
\begin{Brzmienie}
§~\ref{art:wycena} ust.~7 (nowe brzmienie): „Ustalenie Niezależnego Eksperta jest wiążące dla stron i Spółki, z wyjątkiem oczywistego błędu rachunkowego, naruszenia zasad niezależności, rażącego pominięcia istotnych danych udostępnionych ekspertowi albo oczywistej sprzeczności z zasadami określonymi w ust.~2 i~5; w takim przypadku każda ze stron może żądać ponownej wyceny przez innego Niezależnego Eksperta wybranego w trybie ust.~4.”
\end{Brzmienie}
\end{MemoBox}

\end{enumerate}
\end{ContractClause}

\begin{ContractClause}{Organy Spółki}
\label{art:organy}
\begin{enumerate}
  \item Organami Spółki są:
  \begin{enumerate}[label=\alph*)]
    \item Walne Zgromadzenie;
    \item Rada Dyrektorów.
  \end{enumerate}
  \item Spółka przyjmuje system monistyczny. Nie powołuje się zarządu ani rady nadzorczej, chyba że niniejsza umowa zostanie odpowiednio zmieniona.
\end{enumerate}

\end{ContractClause}

\begin{ContractClause}{Rada Dyrektorów}
\label{art:rada}
\begin{enumerate}
  \item Rada Dyrektorów prowadzi sprawy Spółki, reprezentuje Spółkę i sprawuje stały nadzór nad jej działalnością.
  \item Rada Dyrektorów składa się z 3 do 7 dyrektorów. Dyrektorów powołuje Walne Zgromadzenie uchwałą podjętą bezwzględną większością głosów oddanych. Odwołanie dyrektora przed upływem kadencji z ważnych powodów następuje uchwałą podjętą bezwzględną większością głosów oddanych, a bez ważnych powodów --- uchwałą podjętą większością 75\% wszystkich głosów.
  
\begin{MemoBox}{2.3.1}
\MemoPytCd{2.3.1}
\Kierunek{Opcjonalnie dodać, że co najmniej jeden dyrektor niewykonawczy nie jest akcjonariuszem ani Osobą Bliską akcjonariusza (rezerwa na wypadek konfliktu interesów wyłączającego większość Rady, §~\ref{art:konkurencja} ust.~8). Niezależny dyrektor niewykonawczy jest oczekiwaniem inwestorów (vc.md) i wzmacnia bramkę §~\ref{art:rada-zastrzezone} ust.~2; koszt --- wynagrodzenie i D\&O. Decyzja Założycieli: czy Rada ma obejmować co najmniej jednego niezależnego dyrektora niewykonawczego niebędącego akcjonariuszem.}
\end{MemoBox}

\item Kadencja dyrektora trwa trzy lata i jest wspólna, chyba że uchwała o powołaniu stanowi inaczej. Dopuszczalne jest ponowne powołanie. Dyrektorem może być wyłącznie osoba fizyczna spełniająca Kryterium Bezpieczeństwa EU/NATO w rozumieniu §~\ref{art:dopuszczalny-nabywca} i niebędąca osobą, wobec której obowiązuje bezwzględnie obowiązujący zakaz pełnienia funkcji, chyba że bezwzględnie obowiązujące przepisy prawa wymagają rozwiązania odmiennego.
  
\begin{MemoBox}{1.1.5}
\MemoPyt{1.1.5}{Czy wymóg Kryterium wobec dyrektorów (§~21 ust.~3) jest dopuszczalny?}
\Ocena{zgodne warunkowo}
\Klasa{zalecane}
\Podstawa{art.~300$^{52}$ §~1 i art.~300$^{53}$ KSH (dyrektorzy podlegają ograniczeniom ustanowionym w umowie spółki), art.~18 §~1--2 KSH (wymogi ustawowe dla członków organów) \Wiar{Z}; art.~300$^{1}$ §~3 KSH \Wiar{Z}; art.~10 ust.~1 pkt~2 w zw. z pkt~1 lit.~a ustawy z 13 czerwca 2019 r. (obywatelstwo osób kierujących działalnością koncesjonowaną) \Wiar{Z}; art.~11$^{3}$ Kodeksu pracy (dotyczy „dyskryminacji w zatrudnieniu”, nie stosunku organizacyjnego) \Wiar{Z}.}
\Uzasadnienie{KSH określa jedynie minimalne wymogi wobec członków organów (art.~18 §~1--2) i nie zakazuje umowie spółki wymogów dodatkowych (kwalifikacje, wiek, obywatelstwo, szersza niekaralność); dla P.S.A. potwierdza to art.~300$^{53}$ KSH: „wobec spółki członkowie zarządu i dyrektorzy podlegają ograniczeniom ustanowionym w umowie spółki” \Wiar{Z}; praktyka zna takie klauzule w spółkach regulowanych \Wiar{W}. Stosunek organizacyjny dyrektora nie jest stosunkiem pracy, więc zakaz dyskryminacji z Kodeksu pracy nie ma zastosowania; przy równoczesnym zatrudnieniu kryterium dotyczy powołania do organu, nie zatrudnienia. Wymóg jest spójny z reżimem koncesyjnym, który wymaga, aby co najmniej dwie osoby z organu zarządzającego (albo członek organu i prokurent lub pełnomocnik kierujący działalnością koncesjonowaną) miały m.in.~obywatelstwo polskie, innego państwa UE, Szwajcarii albo EFTA-EOG bądź zezwolenie na pobyt stały lub status rezydenta długoterminowego UE (art.~10 ust.~1 pkt~2 w zw. z pkt~1 lit.~a ustawy z 13 czerwca 2019 r.; ustawa z 13 marca 2026 r. go nie zmienia \Wiar{Z}); Kryterium jest szersze (NATO), więc tego wymogu nie zastępuje. Zastrzeżenie o bezwzględnie obowiązujących przepisach jest prawidłowe. Wadą jest sztywność: inwestor rundy ma prawo nominacji dyrektora (§~\ref{art:kwalifikowana-runda} ust.~4 lit.~c), a jego kandydat może nie spełniać Kryterium (np.~partner ze Szwajcarii). Skoro Walne Zgromadzenie może dopuścić nabywcę spoza Kryterium, powinno móc dopuścić także dyrektora, z ograniczeniem dostępu do Kluczowej Własności Intelektualnej (§~\ref{art:security} ust.~6).}
\Rekomendacja{Dodać wyjątek za zgodą Walnego Zgromadzenia (75\% wszystkich głosów) z zastrzeżeniem §~\ref{art:security} ust.~6 i §~\ref{art:export}, a także wymóg, aby w każdym czasie większość dyrektorów, w tym Przewodniczący, spełniała Kryterium.}
\Kierunek{Wyjątek nie osłabia bezpieczeństwa (większość Rady i dostęp do IP pozostają chronione), a usuwa punkt sporny w negocjacjach z inwestorami. Sąd rejestrowy bada jedynie wymogi ustawowe; rejestr akcjonariuszy --- bez wpływu.
\begin{Brzmienie}
(§~\ref{art:rada} ust.~3, zdanie trzecie w nowym brzmieniu) „Dyrektorem może być wyłącznie osoba fizyczna spełniająca Kryterium Bezpieczeństwa EU/NATO w rozumieniu §~\ref{art:dopuszczalny-nabywca} i niebędąca osobą, wobec której obowiązuje bezwzględnie obowiązujący zakaz pełnienia funkcji, chyba że bezwzględnie obowiązujące przepisy prawa wymagają rozwiązania odmiennego. Walne Zgromadzenie może, uchwałą podjętą większością 75\% wszystkich głosów, powołać dyrektora niespełniającego Kryterium, jeżeli po powołaniu większość dyrektorów, w tym Przewodniczący, spełnia Kryterium; dostęp takiego dyrektora do Kluczowej Własności Intelektualnej podlega §~\ref{art:security} ust.~6 i §~\ref{art:export}.”
\end{Brzmienie}}
\end{MemoBox}

\begin{MemoBox}{2.1.1}
\MemoPyt{2.1.1}{Czy większości przy powołaniu i odwołaniu oraz przedłużenie mandatu do 6 miesięcy są zgodne z przepisami KSH o organach P.S.A. i czy przedłużenie mandatu nie koliduje z zasadą wygaśnięcia mandatu?}
\Ocena{zgodne warunkowo (większości --- zgodne; prorogacja mandatu --- wymaga przebudowy redakcyjnej)}
\Klasa{zalecane}
\Podstawa{art.~300$^{73}$ §~3 KSH (powołanie i odwołanie dyrektorów uchwałą akcjonariuszy, „chyba że umowa spółki stanowi inaczej”) \Wiar{Z}; art.~300$^{74}$ §~1--2 KSH (umowa „może zawierać inne postanowienia, w szczególności ograniczać prawo odwołania do ważnych powodów”) \Wiar{Z}; art.~300$^{56}$ §~1--4 KSH (mandat; kadencja liczona w latach obrotowych; wspólna kadencja; wygaśnięcie mandatu --- każdorazowo „chyba że umowa spółki stanowi inaczej”) \Wiar{Z}; art.~300$^{98}$ §~1 KSH (bezwzględna większość głosów jako reguła domyślna) \Wiar{Z}; art.~300$^{96}$ KSH (wyłączenie od głosowania --- nie obejmuje głosowania nad własnym odwołaniem) \Wiar{Z}; art.~39 KC (czynność rzekomego organu, potwierdzenie) \Wiar{Z}.}
\Uzasadnienie{Większości (ust.~2) są poprawne. Skoro art.~300$^{74}$ §~2 KSH pozwala umowie „ograniczać prawo odwołania do ważnych powodów”, tym bardziej dopuszcza większość kwalifikowaną (75\% wszystkich głosów) dla odwołania bez ważnych powodów; bezwzględna większość głosów oddanych jest ustawową regułą domyślną (art.~300$^{98}$ §~1). Uwagi praktyczne: (1) przy podziale 50/20/10/10/10 akcjonariusz z 50\% głosów blokuje własne odwołanie także z ważnych powodów --- art.~300$^{96}$ nie obejmuje odwołania z funkcji, a wyłączenie umowne budzi wątpliwości (pkt~2.2.3); (2) „ważne powody” będą osią sporu --- warto podać przykładowy katalog (naruszenie §~\ref{art:security}, §~\ref{art:export}, utrata Kryterium, Odejście Zawinione).

Prorogacja mandatu. Art.~300$^{56}$ §~2 KSH zastrzega „chyba że umowa spółki stanowi inaczej”, a §~1 dopuszcza nawet powołanie na czas nieoznaczony \Wiar{Z}; ust.~4 mieści się więc w upoważnieniu ustawowym. Wada jest redakcyjna: „mandat dyrektora, którego kadencja upłynęła” odtwarza spór o koncepcję prolongacyjną i redukcyjną (w S.A. na tle art.~369 §~4) i niepewność, czy w okresie „przedłużenia” dyrektor działa z mandatu, czy jako rzekomy organ (art.~39 KC). Umowa nie wskazuje też, czy „trzy lata” to lata obrotowe, i milczy o dyrektorze powołanym w trakcie kadencji wspólnej (lukę wypełnia dyspozytywny art.~300$^{56}$ §~3 --- warto to powtórzyć). Jednoznaczna klauzula ogranicza ryzyko odmowy wpisu czynności z okresu przejściowego i kwestionowania uchwał Rady podjętych po upływie trzech lat.}
\Rekomendacja{Zachować mechanizm ciągłości, ale zbudować go jako element definicji kadencji i mandatu, nie jako „przedłużenie” po wygaśnięciu: kadencja liczona w latach obrotowych, mandat wygasa z dniem powołania następcy, nie później niż 6 miesięcy po upływie kadencji; dodać regułę dla dyrektora dokooptowanego w trakcie kadencji wspólnej i obowiązek Rady zwołania WZ z porządkiem obrad obejmującym wybory nie później niż 3 miesiące przed upływem kadencji. Rozważyć przykładowy katalog ważnych powodów (regulamin Rady albo umowa akcjonariuszy). Decyzja Założycieli: czy bezwzględna większość oddanych ma pozostać przy odwołaniu z ważnych powodów, mimo że akcjonariusz z 50\% może ją zablokować.}
\Kierunek{Zmiana redakcyjna §~\ref{art:rada} ust.~3--4 usuwa spór o naturę „przedłużonego mandatu”, nie zmienia układu sił między akcjonariuszami; dla sądu rejestrowego i kontrahentów daje jednoznaczną datę wygaśnięcia mandatu. Poniżej brzmienie ust.~3; brzmienie ust.~4 --- w ramce przy ust.~4.}
\begin{Brzmienie}
3. Dyrektorów powołuje się na wspólną kadencję trwającą trzy lata obrotowe, liczoną od dnia powołania pierwszego składu Rady Dyrektorów w danej kadencji; kadencja dyrektora powołanego w trakcie kadencji wspólnej kończy się z jej upływem. Dopuszczalne jest ponowne powołanie. Dyrektorem może być wyłącznie osoba fizyczna spełniająca Kryterium Bezpieczeństwa EU/NATO w rozumieniu §~\ref{art:dopuszczalny-nabywca} i niebędąca osobą, wobec której obowiązuje bezwzględnie obowiązujący zakaz pełnienia funkcji, chyba że bezwzględnie obowiązujące przepisy prawa wymagają rozwiązania odmiennego.
\end{Brzmienie}
\end{MemoBox}

\begin{MemoBox}{4.4.2}
\MemoPytCd{4.4.2}
\Kierunek{(V)~Wyjątek dla dyrektora inwestora (A.5; zależnie od decyzji Założycieli o~kręgu inwestorów): za zgodą Walnego Zgromadzenia 75\% (jak przy nabywcy), jedna osoba, z~zachowaniem większości Rady spełniającej Kryterium i~z~dostępem do Kluczowej Własności Intelektualnej wyłącznie w~trybie §~\ref{art:security} ust.~6. Do sprawdzenia, czy przepisy koncesyjne (ustawa z~2019~r. o~wytwarzaniu i~obrocie bronią) albo wymogi zamawiających nie stawiają dyrektorom wymogów obywatelstwa, które wyjątek by naruszał \Wiar{?}. Nowe zdanie ostatnie ust.~3:}
\begin{Brzmienie}
Walne Zgromadzenie może, uchwałą podjętą większością 75\% wszystkich głosów, powołać jednego dyrektora niewykonawczego niespełniającego Kryterium Bezpieczeństwa EU/NATO, wskazanego przez akcjonariusza będącego Kwalifikowanym Inwestorem Finansowym na podstawie §~\ref{art:kwalifikowana-runda} ust.~4 lit.~c, jeżeli osoba ta nie jest Niedopuszczalnym Nabywcą ani obywatelem lub rezydentem państwa objętego bezwzględnie obowiązującymi środkami ograniczającymi, większość dyrektorów, w tym Przewodniczący i każdy dyrektor wykonawczy, spełnia Kryterium, a dostęp tej osoby do Kluczowej Własności Intelektualnej i informacji objętych kontrolą eksportu następuje wyłącznie w trybie §~\ref{art:security} ust.~6 i §~\ref{art:export}; regulamin Rady Dyrektorów określa zakres wyłączenia tej osoby z dostępu do informacji i z głosowania w sprawach objętych tymi ograniczeniami.
\end{Brzmienie}
\end{MemoBox}

\item W celu zapewnienia ciągłości działania Spółki mandat dyrektora, którego kadencja upłynęła, trwa do dnia powołania jego następcy, nie dłużej jednak niż 6 miesięcy od dnia upływu kadencji, chyba że bezwzględnie obowiązujące przepisy prawa stanowią inaczej. Jeżeli liczba dyrektorów spadnie poniżej minimum określonego w ust.~2, każdy akcjonariusz może żądać niezwłocznego zwołania Walnego Zgromadzenia z porządkiem obrad obejmującym uzupełnienie składu Rady Dyrektorów.
  
\begin{MemoBox}{2.1.1}
\MemoPytCd{2.1.1}
\begin{Brzmienie}
4. Na podstawie art.~300$^{56}$ KSH strony postanawiają, że mandat dyrektora wygasa z dniem powołania jego następcy albo z dniem podjęcia uchwały Walnego Zgromadzenia o niepowoływaniu następcy, nie później jednak niż z upływem 6 miesięcy od dnia upływu kadencji. Rada Dyrektorów zwołuje Walne Zgromadzenie z porządkiem obrad obejmującym powołanie dyrektorów na kolejną kadencję nie później niż 3 miesiące przed upływem kadencji. Jeżeli liczba dyrektorów spadnie poniżej minimum określonego w ust.~2, każdy akcjonariusz może żądać niezwłocznego zwołania Walnego Zgromadzenia z porządkiem obrad obejmującym uzupełnienie składu Rady Dyrektorów. Mandat wygasa również wskutek śmierci, rezygnacji albo odwołania dyrektora.
\end{Brzmienie}
\end{MemoBox}

\item Rada wybiera ze swojego grona Przewodniczącego i może wybrać Wiceprzewodniczącego.
  \item Rada wyznacza co najmniej jednego dyrektora wykonawczego, któremu deleguje prowadzenie bieżącej działalności. Jeżeli Rada liczy co najmniej trzy osoby, co najmniej jeden dyrektor powinien być dyrektorem niewykonawczym.
  
\begin{MemoBox}{2.1.3}
\MemoPyt{2.1.3}{Czy podział na dyrektorów wykonawczych i niewykonawczych oraz delegowanie prowadzenia spraw (§~21 ust.~6--8) odpowiada art.~300$^{76}$ KSH i nie ogranicza odpowiedzialności dyrektorów?}
\Ocena{zgodne}
\Klasa{opcjonalne}
\Podstawa{art.~300$^{76}$ KSH (delegacja czynności prowadzenia przedsiębiorstwa umową, regulaminem albo uchwałą Rady, z wyjątkiem art.~300$^{75}$ §~2--3; stały nadzór dyrektorów niewykonawczych) \Wiar{Z}; art.~300$^{75}$ §~2 KSH \Wiar{Z}; art.~300$^{54}$ KSH (staranność zawodowa i lojalność wobec spółki) \Wiar{Z}; art.~300$^{125}$ §~1--2 KSH (odpowiedzialność za szkodę, chyba że członek organu nie ponosi winy; działanie „w granicach uzasadnionego ryzyka gospodarczego”) \Wiar{Z}; art.~300$^{132}$ KSH (odpowiedzialność za zobowiązania spółki przy bezskutecznej egzekucji --- przepis mówi literalnie o „członkach zarządu”, a art.~300$^{133}$ rozciąga go tylko na likwidatorów) \Wiar{Z}; objęcie art.~300$^{132}$ dyrektorów w drodze wykładni \Wiar{?}.}
\Uzasadnienie{§~\ref{art:rada} ust.~6 jest delegacją uchwałą Rady w rozumieniu art.~300$^{76}$ KSH, a ust.~7 powierza dyrektorom niewykonawczym nadzór, który ustawa i tak przypisuje im jako „stały nadzór”. Wyjątek ustawowy --- sprawy z art.~300$^{75}$ §~2 (decyzje strategiczne, plany, struktura) i §~3 (prokura) nie mogą być delegowane --- projekt respektuje przez §~\ref{art:rada-zastrzezone} ust.~1 lit.~a)--b) i §~\ref{art:reprezentacja} ust.~2. Ust.~8 odpowiada stanowi prawnemu: odpowiedzialność każdego dyrektora jest indywidualna i zależna od winy; delegacja zmienia treść obowiązków (niewykonawczy --- nadzór, wykonawczy --- prowadzenie), nie ich istnienie, i nie da się jej umownie wyłączyć wobec Spółki ani wierzycieli (art.~300$^{125}$ §~1; co do art.~300$^{132}$ ustawa mówi o „członkach zarządu”, a stosowanie go do dyrektorów przyjmuje się w drodze wykładni --- bez wpływu na ocenę klauzuli). Uwagi redakcyjne: warunek „jeżeli Rada liczy co najmniej trzy osoby” jest zbędny (Rada zawsze liczy co najmniej 3 osoby), a „powinien” jest miękkie --- przy roli §~\ref{art:rada-zastrzezone} ust.~2 potrzebny jest zawsze co najmniej jeden dyrektor niewykonawczy, więc lepiej „jest”; warto też zapisać, że Rada może zmienić albo cofnąć delegację w każdym czasie i że nie obejmuje ona spraw z §~\ref{art:rada-zastrzezone}. Sąd rejestrowy nie wpisuje podziału na dyrektorów wykonawczych i niewykonawczych, więc działa on wyłącznie wewnętrznie; wobec osób trzecich każdy dyrektor reprezentuje według §~\ref{art:reprezentacja}.}
\Rekomendacja{Bez zmian merytorycznych; opcjonalnie wzmocnić ust.~6 i 8 redakcyjnie.}
\Kierunek{Zmiana porządkowa, bez wpływu na rejestr i akcjonariuszy; dla inwestora czytelniejszy podział ról. Brzmienie ust.~6 (obowiązek zapewnienia co najmniej jednego dyrektora niewykonawczego zamiast miękkiego „powinien”; delegacja nie obejmuje spraw z art.~300$^{75}$ §~2--3 KSH ani z §~\ref{art:rada-zastrzezone} i może być w każdym czasie zmieniona albo cofnięta):}
\begin{Brzmienie}
6. Rada Dyrektorów wyznacza uchwałą co najmniej jednego dyrektora wykonawczego, któremu deleguje prowadzenie bieżącej działalności Spółki zgodnie z art.~300$^{76}$ KSH, oraz zapewnia, aby co najmniej jeden dyrektor był dyrektorem niewykonawczym. Delegacja nie obejmuje spraw, o których mowa w art.~300$^{75}$ §~2--3 KSH ani w §~\ref{art:rada-zastrzezone}, i może być w każdym czasie zmieniona albo cofnięta uchwałą Rady Dyrektorów.
\end{Brzmienie}
\end{MemoBox}

\item Dyrektorzy niewykonawczy sprawują w szczególności nadzór nad finansami, ryzykiem, bezpieczeństwem informacji, kontrolą eksportu, sankcjami, zgodnością, transakcjami z podmiotami powiązanymi oraz ochroną własności intelektualnej.
  \item Podział zadań nie ogranicza ustawowych praw i obowiązków Rady ani odpowiedzialności dyrektorów.
  \item Uchwały Rady zapadają bezwzględną większością głosów przy obecności co najmniej połowy jej składu, chyba że niniejsza umowa wymaga większości kwalifikowanej. W razie równości głosów decyduje głos Przewodniczącego, z wyjątkiem spraw dotyczących jego interesu osobistego.
  
\begin{MemoBox}{2.3.2}
\MemoPyt{2.3.2}{Czy quorum liczone od pełnego składu (§~21 ust.~9) i wyłączenie dyrektorów (§~29 ust.~7) są spójne z przepisami KSH o konflikcie interesów (art.~300$^{55}$ KSH) i o uchwałach organów?}
\Ocena{zgodne warunkowo (spójne co do zasady; niespójna redakcja „składu” i „pełnego składu”)}
\Klasa{zalecane}
\Podstawa{art.~300$^{58}$ §~1--5 KSH (§~2: kworum „co najmniej połowa jej członków”, umowa „może przewidywać surowsze wymagania dotyczące kworum”; §~3: głosujący na piśmie albo na odległość liczeni do kworum, „chyba że umowa spółki lub regulamin organu stanowią inaczej”; §~4 większość i głos rozstrzygający; §~5 protokół) \Wiar{Z}; art.~300$^{55}$ §~1 KSH („powinien ujawnić sprzeczność interesów i wstrzymać się od udziału w rozstrzyganiu takich spraw oraz może żądać zaznaczenia tego w protokole”) \Wiar{Z}; art.~300$^{57}$ §~1 KSH (regulamin organu) \Wiar{Z}.}
\Uzasadnienie{Wyłączenie z §~\ref{art:konkurencja} ust.~7 odtwarza art.~300$^{55}$ §~1 i dodaje, że głosu nie liczy się, a wyłączenie „nie zmniejsza pełnego składu” jako podstawy quorum --- to rozwiązanie surowsze niż ustawa, wprost dopuszczone przez art.~300$^{58}$ §~2 i przyjęte przez audyt. Dwie nieścisłości. (1)~Ust.~9 mówi o „obecności co najmniej połowy jej składu”, a §~\ref{art:konkurencja} ust.~7 o „pełnym składzie”; ustawowe „połowa jego członków” jest w doktrynie rozumiane różnie (skład faktyczny albo umowny). Mianownik trzeba zdefiniować jednolicie: liczba dyrektorów aktualnie powołanych (nie maksymalna z ust.~2, nie tylko obecni). (2)~Umowa nie mówi, czy dyrektor wyłączony, lecz obecny, liczy się do quorum; skoro §~\ref{art:konkurencja} ust.~6 lit.~b)--c) zabrania mu udziału w ocenie i dostępu do dokumentów, spójna jest reguła, że do quorum w danej sprawie liczą się tylko dyrektorzy niewyłączeni, a warunek z §~\ref{art:konkurencja} ust.~7 zd.~3 (co najmniej dwóch dyrektorów niekonfliktowych) jest dolnym progiem. Przykład: Rada trzyosobowa, jeden wyłączony --- quorum 2 z 3, głosują dwaj, przy równości decyduje Przewodniczący, chyba że sprawa dotyczy jego interesu (zgodne z art.~300$^{58}$ §~4); dwóch wyłączonych --- brak zdolności do uchwały, stosuje się §~\ref{art:konkurencja} ust.~8. Protokół Rady powinien wskazywać pełny skład, obecnych, wyłączonych i podstawę wyłączenia, co rozstrzyga dowodowo; art.~300$^{58}$ §~5 wymaga protokołu z nazwiskami głosujących i wynikiem, a wyłączony „może żądać zaznaczenia tego w protokole” (art.~300$^{55}$ §~1). Uchwały pisemne i zdalne są dopuszczalne z mocy art.~300$^{58}$ §~1--3, więc ust.~10 jest tylko potwierdzeniem.}
\Rekomendacja{Ujednolicić definicję „pełnego składu” w §~\ref{art:rada} ust.~9 i §~\ref{art:konkurencja} ust.~7 oraz rozstrzygnąć, że dyrektor wyłączony nie liczy się do quorum w danej sprawie.}
\Kierunek{Zmiana redakcyjna, bez skutków dla rejestru; zapobiega sporowi o ważność uchwały Rady. Brzmienie §~\ref{art:rada} ust.~9 zd.~1:}
\begin{Brzmienie}
9. Uchwały Rady Dyrektorów zapadają bezwzględną większością głosów, jeżeli w posiedzeniu, głosowaniu pisemnym albo głosowaniu przy wykorzystaniu środków bezpośredniego porozumiewania się na odległość uczestniczy co najmniej połowa pełnego składu Rady Dyrektorów, przez który rozumie się liczbę dyrektorów aktualnie powołanych; dyrektora wyłączonego od rozstrzygania danej sprawy na podstawie §~\ref{art:konkurencja} ust.~7 nie uwzględnia się przy ustalaniu quorum w tej sprawie, przy czym pełny skład Rady Dyrektorów nie ulega z tego powodu zmniejszeniu. Niniejsza umowa może wymagać większości kwalifikowanej. W razie równości głosów decyduje głos Przewodniczącego, z wyjątkiem spraw dotyczących jego interesu osobistego.
\end{Brzmienie}
\end{MemoBox}

\item Posiedzenia i uchwały mogą odbywać się przy wykorzystaniu środków komunikacji elektronicznej, jeżeli zapewniono identyfikację uczestników, bezpieczeństwo i możliwość dwustronnej komunikacji.
  \item Regulamin Rady Dyrektorów określa szczegółowy podział kompetencji, limity wydatków, zasady raportowania, komitety i politykę konfliktów interesów.
\end{enumerate}

\end{ContractClause}

\begin{ContractClause}{Reprezentacja Spółki}
\label{art:reprezentacja}
\begin{enumerate}
  \item Do składania oświadczeń w imieniu Spółki wymagane jest współdziałanie dwóch dyrektorów albo jednego dyrektora łącznie z prokurentem.
  \item Powołanie prokurenta wymaga zgody wszystkich dyrektorów. Odwołać prokurę może każdy dyrektor, zgodnie z art. 300$^{75}$ § 3 KSH.
  \item \BoilerplateStrike{Postanowienia ogólne dotyczące reprezentacji Spółki nie naruszają art. 300$^{79}$ KSH.} Jeżeli stroną umowy ze Spółką albo stroną sporu ze Spółką jest dyrektor, Spółkę reprezentuje pełnomocnik powołany uchwałą akcjonariuszy.
  \item W umowie oraz sporze między Spółką a dyrektorem wykonawczym Spółkę może reprezentować także dyrektor niewykonawczy działający na podstawie uchwały Rady Dyrektorów podjętej wyłącznie przez dyrektorów niewykonawczych, jeżeli spełnione są warunki określone w niniejszej umowie oraz art. 300$^{79}$ § 3 KSH.
  \item Jeżeli wszystkie akcje Spółki przysługują jednemu akcjonariuszowi, który jest zarazem dyrektorem, do czynności prawnej między tym akcjonariuszem a reprezentowaną przez niego Spółką stosuje się art. 300$^{79}$ § 4 KSH, w tym wymóg formy aktu notarialnego.
  \item Wewnętrzne limity kompetencji dyrektorów wykonawczych nie wpływają na skuteczność reprezentacji wobec osób trzecich, lecz ich naruszenie może stanowić podstawę odpowiedzialności wobec Spółki.
\end{enumerate}
\end{ContractClause}

\begin{ContractClause}{Sprawy wymagające uchwały Rady Dyrektorów}
\label{art:rada-zastrzezone}
\begin{enumerate}
  \item Uchwały Rady Dyrektorów wymagają w szczególności:
  \begin{enumerate}[label=\alph*)]
    \item przyjęcie i zmiana strategii, rocznego budżetu oraz planu finansowego;
    \item ustalenie struktury organizacyjnej, delegowanie zadań dyrektorom wykonawczym i zatwierdzenie regulaminów;
    \item zawarcie umowy lub dokonanie wydatku poza zatwierdzonym budżetem o wartości przekraczającej 250 000 zł;
    \item zaciągnięcie finansowania, udzielenie zabezpieczenia albo rozporządzenie aktywami o wartości przekraczającej 100 000 zł, z wyjątkiem czynności mieszczących się w Planie Finansowania przyjętym zgodnie z §~\ref{art:finansowanie} ust.~5;
    \item zawarcie transakcji z akcjonariuszem, dyrektorem, osobą bliską lub podmiotem powiązanym;
    \item rozpoczęcie działalności w nowej jurysdykcji, utworzenie spółki zależnej albo stałego zakładu;
    \item udzielenie licencji wyłącznej lub przeniesienie istotnej własności intelektualnej w zwykłym toku działalności, z zastrzeżeniem §~\ref{art:walne-zastrzezone};
    \item klasyfikacja produktu jako podlegającego kontroli eksportu, zatwierdzenie eksportu albo udostępnienia technologii kontrolowanej;
    \item przyjęcie polityk bezpieczeństwa, kontroli eksportu, sankcji, cyberbezpieczeństwa, zarządzania danymi i korzystania z otwartego oprogramowania;
    \item przygotowanie i wykonanie emisji akcji serii P lub F w zakresie powierzonym Radzie Dyrektorów uchwałą Walnego Zgromadzenia, bez kompetencji do samodzielnego podjęcia uchwały o emisji ani do samodzielnego pozbawienia prawa poboru;
    \item utworzenie Przedsięwzięcia Produkcyjnego, przystąpienie do niego, objęcie albo utrata kontroli nad nim oraz zwiększenie zaangażowania Spółki w Przedsięwzięcie Produkcyjne, jeżeli łączne zaangażowanie Spółki w dane Przedsięwzięcie nie przekracza 500 000 zł, a także udzielenie Przedsięwzięciu Produkcyjnemu licencji zgodnie z §~\ref{art:przedsiewziecia} ust.~3.
  \end{enumerate}
  \item Sprawy wymienione w ust. 1 lit. e)--i) oraz k) wymagają pozytywnego głosu co najmniej jednego dyrektora niewykonawczego odpowiedzialnego za bezpieczeństwo lub zgodność, o ile taki dyrektor został powołany.

\begin{MemoBox}{2.1.2}
\MemoPyt{2.1.2}{Czy wymóg pozytywnego głosu określonego dyrektora (§~23 ust.~2) jest dopuszczalny (uprzywilejowanie głosu w Radzie) i jaki jest skutek jego braku dla ważności uchwały?}
\Ocena{zgodne warunkowo}
\Klasa{zalecane}
\Podstawa{art.~300$^{58}$ §~4 KSH („Uchwały organu zapadają bezwzględną większością głosów, chyba że umowa spółki stanowi inaczej”; głos rozstrzygający „prezesa, przewodniczącego lub innego członka organu”), §~2 (surowsze wymagania co do kworum) i §~5 (protokół z imionami i nazwiskami głosujących oraz wynikiem) \Wiar{Z}; art.~300$^{75}$ §~1--2 i art.~300$^{76}$ §~1 KSH \Wiar{Z}; art.~17 §~1--3 KSH (brak uchwały wymaganej ustawą --- nieważność; wymaganej wyłącznie umową --- ważność, odpowiedzialność wobec spółki) \Wiar{Z}; art.~300$^{77}$ §~2 KSH (prawa dyrektora do reprezentowania nie można ograniczyć ze skutkiem wobec osób trzecich) \Wiar{Z}; art.~189 KPC w zw. z art.~58 KC \Wiar{Z} oraz uchwała SN (7) z 18.09.2013~r., III~CZP~13/13 (zaskarżanie uchwał organów menedżerskich powództwem o ustalenie) \Wiar{W}.}
\Uzasadnienie{Dopuszczalność. Art.~300$^{58}$ §~4 KSH pozwala umowie określić inną niż bezwzględna większość i przewiduje głos rozstrzygający „innego członka organu”, więc równość głosów w Radzie P.S.A. nie jest zasadą bezwzględną. Wymóg głosu dyrektora niewykonawczego odpowiedzialnego za bezpieczeństwo lub zgodność w sprawach z ust.~1 lit.~e)--i) oraz k) jest większością kwalifikowaną o charakterze funkcyjnym, a nie personalnym uprzywilejowaniem konkretnej osoby; uzasadnia go art.~300$^{76}$ (nadzór dyrektorów niewykonawczych) i regulacyjny profil Spółki (tak też audyt). Luki: (1) umowa nie mówi, kto i jak wyznacza takiego dyrektora --- §~\ref{art:rada} ust.~7 przypisuje te obszary wszystkim dyrektorom niewykonawczym, więc przy braku wyraźnego wyznaczenia wymóg jest sporny; (2) klauzula „o ile taki dyrektor został powołany” pozwala uniknąć wymogu przez niewyznaczenie nikogo; (3) brak reguły na wypadek wyłączenia (§~\ref{art:konkurencja} ust.~7) albo nieobecności tego dyrektora.

Skutek braku. Uchwała bez wymaganego głosu nie uzyskuje większości umownej, więc nie zostaje podjęta; jej „podjęcie” można kwestionować powództwem o ustalenie (art.~189 KPC). Na zewnątrz skutek jest ograniczony: wymóg uchwały wynika z umowy, nie z ustawy (poza art.~300$^{75}$ §~2), więc czynność pozostaje ważna (art.~17 §~3), a ograniczenia reprezentacji są bezskuteczne wobec osób trzecich (art.~300$^{77}$ §~2; §~\ref{art:reprezentacja} ust.~6). Realną sankcją jest odpowiedzialność dyrektorów wobec Spółki i --- przy naruszeniach eksportowych --- sankcje publicznoprawne. Protokół powinien utrwalać, kto był dyrektorem odpowiedzialnym i jak głosował (art.~300$^{58}$ §~5).}
\Rekomendacja{Doprecyzować: Rada wyznacza uchwałą co najmniej jednego dyrektora niewykonawczego odpowiedzialnego za bezpieczeństwo i zgodność niezwłocznie po ukonstytuowaniu się (obowiązek, nie opcja); do czasu wyznaczenia sprawy z lit.~e)--i) i k) wymagają głosów wszystkich dyrektorów niewykonawczych; przy wyłączeniu albo nieobecności wyznaczonego dyrektora --- głosu innego dyrektora niewykonawczego, a w jego braku sprawa trafia do Walnego Zgromadzenia; dodać wprost skutek: „uchwała nie zostaje podjęta”, z prawem przedłożenia sprawy Walnemu Zgromadzeniu (większość 75\% wszystkich głosów), co odbiera klauzuli charakter weta osobistego. Do decyzji Założycieli: czy wymóg ma mieć taki „wentyl” (przedłożenie sprawy WZ), czy pozostać bezwzględną bramką.}
\Kierunek{Zmiana §~\ref{art:rada-zastrzezone} ust.~2 wzmacnia wykonalność bramki bezpieczeństwa (dla programów EDF/DIANA i kontroli eksportu) i odpowiada na typowe pytanie inwestora o „weto jednej osoby”; nie zmienia skutków wobec osób trzecich.}
\begin{Brzmienie}
2. Sprawy wymienione w ust.~1 lit.~e)--i) oraz k) wymagają, oprócz większości określonej w §~\ref{art:rada} ust.~9, pozytywnego głosu co najmniej jednego dyrektora niewykonawczego wyznaczonego uchwałą Rady Dyrektorów jako odpowiedzialny za bezpieczeństwo lub zgodność. Rada Dyrektorów dokonuje takiego wyznaczenia niezwłocznie po każdej zmianie składu; do czasu wyznaczenia sprawy te wymagają pozytywnego głosu wszystkich dyrektorów niewykonawczych. Jeżeli wyznaczony dyrektor jest wyłączony na podstawie §~\ref{art:konkurencja} ust.~7 albo nie uczestniczy w głosowaniu, wymagany jest pozytywny głos innego dyrektora niewykonawczego. Uchwała podjęta bez tego głosu nie zostaje podjęta; Rada Dyrektorów może przedłożyć sprawę Walnemu Zgromadzeniu, które rozstrzyga większością 75\% wszystkich głosów w Spółce. Postanowienie to nie ogranicza reprezentacji Spółki wobec osób trzecich.
\end{Brzmienie}
\end{MemoBox}

\end{enumerate}

\end{ContractClause}

\begin{ContractClause}{Walne Zgromadzenie i sposób obliczania głosów}
\label{art:wz}
\begin{enumerate}
  \item \BoilerplateStrike{Walne Zgromadzenie może być zwyczajne albo nadzwyczajne i działa zgodnie z przepisami prawa oraz niniejszą umową.}
  \item Walne Zgromadzenia odbywają się w siedzibie Spółki albo w Warszawie. Nie wyłącza to uczestnictwa w Walnym Zgromadzeniu przy wykorzystaniu środków komunikacji elektronicznej na zasadach określonych w ust.~6.
  
\begin{MemoBox}{3.4.4}
\MemoPyt{3.4.4}{Przekreślone zdania i ustęp: czy ich usunięcie jest bezpieczne, czy któreś powinny zostać; spójność definicji po usunięciu.}
\Ocena{zgodne warunkowo}
\Klasa{zalecane}
\Podstawa{art.~2 KSH \Wiar{Z}; art.~300$^{30}$ §~1--2, art.~300$^{31}$ §~5 KSH (rejestr akcjonariuszy, wybór podmiotu przez akcjonariuszy) \Wiar{Z}; art.~300$^{79}$ §~1 KSH \Wiar{Z}; art.~300$^{104}$ §~1 KSH \Wiar{Z}; art.~300$^{38}$ §~1 KSH \Wiar{Z}; §~\ref{art:dopuszczalny-nabywca} ust.~14, §~\ref{art:kwalifikowana-runda} ust.~1.}
\Uzasadnienie{Poza §~\ref{art:dopuszczalny-nabywca} ust.~14 i związanym z nim nawiasem w §~\ref{art:kwalifikowana-runda} ust.~1 przekreślenia to powtórzenia ustawy albo deklaracje, których usunięcie jest bezpieczne: §~\ref{art:akcje} ust.~6 (art.~300$^{31}$ §~5 \Wiar{Z}), §~\ref{art:seria-f} ust.~7 zdanie pierwsze (art.~300$^{104}$), §~\ref{art:malzonek} ust.~1 zdanie pierwsze, §~\ref{art:reprezentacja} ust.~3 i §~\ref{art:konkurencja} ust.~12 zdania pierwsze (art.~300$^{79}$ §~1 stosuje się z mocy ustawy \Wiar{Z}), §~\ref{art:wz} ust.~1 oraz §~\ref{art:final} ust.~1 (art.~2 KSH). Ust.~14 (Kwalifikowany Inwestor Finansowy) --- jedyny fragment merytoryczny: usunięcie formalnie bezpieczne (nawias w §~\ref{art:kwalifikowana-runda} ust.~1 musi zniknąć razem z nim), ale kosztowne --- każdy fundusz z inwestorami spoza UE/EOG/NATO potrzebuje zgody 75\% wszystkich głosów; przywrócenie w wersji uszczelnionej --- pakiet 1.4; zdanie o zgodzie w uchwale kierunkowej zachować w §~\ref{art:kwalifikowana-runda} ust.~2 (3.3.3). Usunięcie całych ustępów zmienia numerację i psuje odesłania wpisane na sztywno: w §~\ref{art:wz} ust.~2 „ust.~6” ma brzmieć „ust.~5”, a w ust.~12 „ust.~11” --- „ust.~10”; usunięcie §~\ref{art:final} ust.~1 przesuwa ust.~4--5, do których odsyła umowa akcjonariuszy i Dokument nr~1; usunięcie ust.~14 przesuwa ust.~15, do którego odwołują się materiały robocze; usunięcie §~\ref{art:akcje} ust.~6 nie narusza odesłań. Definicje pozostają spójne: znika tylko „Kwalifikowany Inwestor Finansowy”; §~\ref{art:okres-ograniczenia-zbywania} ust.~2 lit.~a wymaga korekty z innych powodów (3.3.3).}
\Rekomendacja{Usunąć fragmenty deklaratywne i odesłanie (nawias w §~\ref{art:kwalifikowana-runda} ust.~1 --- razem z §~\ref{art:dopuszczalny-nabywca} ust.~14); o ust.~14 zdecydować łącznie z pakietem 1.4 (rekomendacja: przywrócić uszczelniony). Po usunięciu poprawić odesłania w §~\ref{art:wz} ust.~2 i~12 oraz w umowie akcjonariuszy; przed podpisaniem zweryfikować mechanicznie wszystkie odesłania „ust.~N” do paragrafów, w których usunięto ustęp.}
\Kierunek{Porządkowanie tekstu bez skutku prawnego, ale błędne odesłanie po renumeracji jest realną wadą aktu notarialnego (wykładnia, sprostowanie). Poniżej brzmienie ust.~2; brzmienie ust.~12 (nowego ust.~11) --- w ramce przy tym ustępie.}
\begin{Brzmienie}
§~\ref{art:wz} ust.~2 (po usunięciu dotychczasowego ust.~1): „Walne Zgromadzenia odbywają się w siedzibie Spółki albo w Warszawie. Nie wyłącza to uczestnictwa w Walnym Zgromadzeniu przy wykorzystaniu środków komunikacji elektronicznej na zasadach określonych w ust.~5.”
\end{Brzmienie}
\end{MemoBox}

\item Zawiadomienia o Walnym Zgromadzeniu są przesyłane na adres poczty elektronicznej akcjonariusza wpisany do rejestru akcjonariuszy, z zachowaniem terminów ustawowych. Każdy akcjonariusz jest zobowiązany zapewnić wpisanie swojego adresu poczty elektronicznej do rejestru akcjonariuszy i jego aktualizację; skutki braku aktualnego adresu w rejestrze obciążają akcjonariusza w zakresie dopuszczalnym przez prawo.
  
\begin{MemoBox}{2.2.2}
\MemoPyt{2.2.2}{Czy zawiadomienie e-mailem na adres z rejestru akcjonariuszy spełnia art.~300$^{87}$ KSH i co, gdy akcjonariusz nie wskaże adresu?}
\Ocena{zgodne warunkowo}
\Klasa{zalecane}
\Podstawa{art.~300$^{87}$ §~1 KSH („Walne zgromadzenie zwołuje się pocztą elektroniczną na adres akcjonariusza wpisany do rejestru akcjonariuszy, na adres do doręczeń elektronicznych lub za pomocą listu poleconego lub przesyłki nadanej pocztą kurierską”; wysyłka co najmniej dwa tygodnie przed WZ) \Wiar{Z}; art.~300$^{33}$ §~1 pkt~5 KSH (adres e-mail w rejestrze, „jeżeli akcjonariusz wyraził zgodę na komunikację [\ldots] przy wykorzystaniu poczty elektronicznej”; po nowelizacji Dz.U. 2026 poz.~176, od 18.02.2027~r., pkt~5 zachowuje wymóg zgody i dodaje PESEL albo datę urodzenia) \Wiar{Z}; art.~300$^{34}$ §~1 KSH (wpis „na żądanie spółki lub innej osoby mającej interes prawny”) \Wiar{Z}; art.~300$^{90}$ KSH (uchwały mimo braku formalnego zwołania) \Wiar{Z}; art.~300$^{101}$ w zw. z art.~422 §~2 pkt~4 KSH (legitymacja nieobecnego akcjonariusza przy wadliwym zwołaniu) \Wiar{Z}; art.~238 §~1 KSH (sp.~z~o.o.: e-mail za uprzednią pisemną zgodą wspólnika) \Wiar{Z}.}
\Uzasadnienie{Kanał jest ustawowy: art.~300$^{87}$ §~1 wymienia pocztę elektroniczną „na adres akcjonariusza wpisany do rejestru akcjonariuszy”, bez odrębnej zgody pisemnej znanej z art.~238 §~1 KSH; zgoda akcjonariusza jest jednak warunkiem samego wpisu adresu (art.~300$^{33}$ §~1 pkt~5). §~\ref{art:wz} ust.~3 jest więc zgodny z ustawą, a odesłanie do terminów ustawowych (dwa tygodnie) poprawne; decyduje data wysłania, więc Spółka powinna zachować dowód wysyłki. Luka dotyczy braku adresu: wpisu dokonuje podmiot prowadzący rejestr na żądanie spółki lub osoby mającej interes prawny (art.~300$^{34}$ §~1) --- Spółka może żądać wpisu, ale nie zastąpi zgody akcjonariusza, a obowiązek „zapewnić wpisanie” nie ma sankcji wykonawczej; nowelizacja poz.~176 tego nie zmienia (nowy art.~300$^{33}$ §~3 dotyczy tylko danych z pkt~1--4 i 9--11, nie adresów). Bez adresu Spółka musi użyć innego kanału ustawowego (adres do doręczeń elektronicznych, list polecony, przesyłka kurierska); zwołanie e-mailem na adres znany Spółce, lecz niewpisany do rejestru, jest wadliwe, a nieobecny akcjonariusz ma legitymację do zaskarżenia uchwały. Zdanie o skutkach braku adresu obciążających akcjonariusza nie sanuje wadliwego zwołania, bo sposób zwołania jest normą bezwzględną. Adresy i zgody warto zebrać już przy pierwszym wpisie Założycieli.}
\Rekomendacja{Dodać w ust.~3 kaskadę: (1) e-mail na adres z rejestru; (2) w jego braku przesyłka polecona albo kurierska na adres z rejestru; (3) nakaz podania adresu e-mail przy pierwszym wpisie i zgoda na komunikację elektroniczną w umowie akcjonariuszy oraz w umowie objęcia akcji; (4) przypomnienie, że przy reprezentacji wszystkich akcji uchwały można podejmować bez formalnego zwołania (art.~300$^{90}$).}
\Kierunek{Zmiana zwiększa odporność uchwał na zaskarżenie i usuwa niepewność podmiotu prowadzącego rejestr, kto ma dostarczyć adres; nie wpływa na rejestr sądowy. Proponowane nowe brzmienie ust.~3 (w skrócie): zawiadomienia wysyła się, z zachowaniem terminów ustawowych, pocztą elektroniczną na adres wpisany do rejestru akcjonariuszy, a w braku takiego wpisu --- przesyłką poleconą albo kurierską na adres akcjonariusza wpisany do rejestru; każdy akcjonariusz najpóźniej przy wpisie do rejestru podaje adres e-mail wraz ze zgodą na komunikację ze Spółką i podmiotem prowadzącym rejestr i aktualizuje go, a Spółka może żądać wpisania lub aktualizacji adresu; zawiadomienie uważa się za dokonane z chwilą wysłania na adres wpisany do rejestru; postanowienie nie wyłącza podejmowania uchwał bez formalnego zwołania, jeżeli reprezentowany jest cały kapitał akcyjny i nikt nie sprzeciwia się odbyciu Walnego Zgromadzenia ani wniesieniu spraw do porządku obrad.}
\end{MemoBox}

\begin{MemoBox}{4.3.3}
\MemoPytCd{4.3.3}
\Kierunek{Termin umowny 7 dni dla akcjonariusza, symetryczny do ustawowego terminu Spółki (art.~300$^{33}$ §~3 KSH od 18 lutego 2027~r.), obejmujący także inne dane, w~tym zmianę kontroli nad akcjonariuszem istotną dla Kryterium; ochrona PESEL i~adresów nie obejmuje adresu e-mail. Proponowane brzmienie:}
\begin{Brzmienie}
§~\ref{art:wz} ust.~3 zdanie drugie: Każdy akcjonariusz jest zobowiązany zapewnić wpisanie swojego adresu poczty elektronicznej do rejestru akcjonariuszy oraz zgłosić Spółce i podmiotowi prowadzącemu rejestr każdą zmianę danych podlegających wpisowi, w tym zmianę adresu, nazwiska albo firmy, a także zmianę kontroli nad akcjonariuszem albo jego beneficjentów rzeczywistych istotną dla Kryterium Bezpieczeństwa EU/NATO, w terminie 7 dni od dnia zdarzenia, chyba że bezwzględnie obowiązujące przepisy przewidują termin krótszy; skutki braku aktualnych danych w rejestrze obciążają akcjonariusza w zakresie dopuszczalnym przez prawo.
\end{Brzmienie}
\end{MemoBox}

\item Uchwały akcjonariuszy mogą być podejmowane na Walnym Zgromadzeniu albo poza Walnym Zgromadzeniem, w tym przy wykorzystaniu środków komunikacji elektronicznej.
  \item Głosowanie przy wykorzystaniu środków komunikacji elektronicznej może odbywać się za pośrednictwem systemu teleinformatycznego wskazanego przez Spółkę albo za pomocą innych środków pozwalających na identyfikację akcjonariusza, ustalenie treści oddanego głosu oraz jego utrwalenie.
  \item Umowa Spółki dopuszcza uczestnictwo w Walnym Zgromadzeniu przy wykorzystaniu środków komunikacji elektronicznej zgodnie z art.~300\textsuperscript{92} KSH, w szczególności przez dwustronną komunikację w czasie rzeczywistym oraz wykonywanie prawa głosu na odległość.
  \item Zawiadomienie o Walnym Zgromadzeniu określa sposób identyfikacji akcjonariusza, uzyskania dostępu do transmisji lub komunikacji oraz wykonywania prawa głosu.
  
\begin{MemoBox}{2.2.1}
\MemoPyt{2.2.1}{Czy sformułowania o udziale elektronicznym są zgodne z art.~300$^{88}$ i 300$^{92}$ KSH (miejsce zgromadzenia a uczestnictwo zdalne) i nie zostaną odczytane jako w pełni wirtualne zgromadzenie?}
\Ocena{zgodne warunkowo (konstrukcja hybrydowa --- zgodna; brak regulaminu WZ wymaganego przez art.~300$^{92}$ §~2)}
\Klasa{zalecane}
\Podstawa{art.~300$^{88}$ §~1--2 KSH (WZ „odbywa się w siedzibie spółki, jeżeli umowa spółki nie wskazuje innego miejsca”) \Wiar{Z}; art.~300$^{92}$ §~1--2 KSH (umowa „może dopuszczać” udział elektroniczny; wymogi „jedynie” niezbędne do identyfikacji i bezpieczeństwa; „Szczegółowe zasady dotyczące sposobu uczestnictwa w walnym zgromadzeniu przy wykorzystaniu środków komunikacji elektronicznej określa regulamin walnego zgromadzenia”) \Wiar{Z}; art.~300$^{87}$ §~4 KSH (zawiadomienie zawiera „dokładny opis sposobu uczestnictwa w walnym zgromadzeniu i wykonywania prawa głosu”) \Wiar{Z}; art.~300$^{80}$ §~1--2 KSH (uchwały poza WZ) \Wiar{Z}; art.~300$^{100}$ §~1 KSH (lista akcjonariuszy głosujących elektronicznie) i art.~300$^{101}$ KSH (art.~422--427 stosowane odpowiednio) \Wiar{Z}.}
\Uzasadnienie{Wersja 0.9.4 usunęła sformułowanie kwestionowane przez audyt („WZ może odbywać się z wykorzystaniem środków komunikacji elektronicznej”). §~\ref{art:wz} ust.~2 wskazuje fizyczne miejsce (siedziba albo Warszawa --- oba w Polsce, dopuszczalne na tle art.~300$^{88}$), a ust.~6 mówi wyłącznie o „uczestnictwie” przy wykorzystaniu środków komunikacji elektronicznej zgodnie z art.~300$^{92}$ --- to ustawowa konstrukcja zgromadzenia hybrydowego: przewodniczący, protokolant (przy zmianie umowy --- notariusz) i miejsce obrad pozostają fizyczne, a akcjonariusze mogą być zdalni. Art.~300$^{92}$ jest opt-in, więc wyraźne dopuszczenie w ust.~6 jest konieczne i jest. Ust.~4 dotyczy odrębnego trybu uchwał poza zgromadzeniem (art.~300$^{80}$), który nie jest „wirtualnym WZ”. Doprecyzowania: (1) według art.~300$^{92}$ §~2 zd.~2 szczegółowe zasady „określa regulamin walnego zgromadzenia”; ust.~7 przenosi je do zawiadomienia, które regulaminem nie jest --- zwołujący ustala wymogi jednostronnie, co daje akcjonariuszowi zdalnemu podstawę do zarzutu, że wykraczają poza „niezbędne”, i do zaskarżenia uchwały (art.~300$^{101}$ w zw. z art.~422 §~2 pkt~3 KSH); (2) ust.~8 (problemy techniczne po stronie akcjonariusza) jest standardem rynkowym, ale nie może wyłączyć sądowej kontroli uchwały --- działa jako rozkład ryzyka i domniemanie. Hybrydowe zgromadzenie z notariuszem w miejscu obrad jest przyjęte w praktyce; w sporze o zaliczenie głosów rozstrzyga dowód identyfikacji i utrwalenia głosu (ust.~5) --- do protokołu dołącza się listę głosujących elektronicznie (art.~300$^{100}$ §~1), a Spółka powinna archiwizować logi systemu.}
\Rekomendacja{Konstrukcję hybrydową zachować. W ust.~7 przewidzieć regulamin Walnego Zgromadzenia przyjęty uchwałą akcjonariuszy jako podstawę szczegółowych zasad udziału elektronicznego (art.~300$^{92}$ §~2), z regułą przejściową dla zawiadomienia; w ust.~8 dodać „w zakresie dopuszczalnym przez prawo”. Regulamin przyjąć na pierwszym Walnym Zgromadzeniu po rejestracji.}
\Kierunek{Zmiana wykonuje wymóg art.~300$^{92}$ §~2 KSH i zmniejsza ryzyko zarzutu, że wymogi techniczne ustalone przez zwołującego wykraczają poza niezbędne; bez wpływu na rejestr i układ sił. Brzmienie ust.~7:}
\begin{Brzmienie}
7. Szczegółowe zasady uczestnictwa w Walnym Zgromadzeniu przy wykorzystaniu środków komunikacji elektronicznej określa regulamin Walnego Zgromadzenia przyjęty uchwałą akcjonariuszy; do czasu jego przyjęcia sposób identyfikacji akcjonariusza, uzyskania dostępu do transmisji lub komunikacji oraz wykonywania prawa głosu określa zawiadomienie o Walnym Zgromadzeniu. Regulamin i zawiadomienie nie mogą ustanawiać wymogów ani ograniczeń innych niż niezbędne do identyfikacji akcjonariuszy i zapewnienia bezpieczeństwa komunikacji elektronicznej.
\end{Brzmienie}
\end{MemoBox}

\item Problemy techniczne leżące po stronie akcjonariusza nie wpływają na ważność Walnego Zgromadzenia ani podjętych uchwał, chyba że uniemożliwienie udziału wynikało z okoliczności leżących po stronie Spółki i mogło mieć wpływ na wynik głosowania.
  \item Każda akcja daje jeden głos, chyba że bezwzględnie obowiązujące przepisy albo uchwała emisyjna zgodna z niniejszą umową stanowią inaczej.
  \item O ile ustawa albo niniejsza umowa nie wymaga surowszej większości, uchwały zapadają bezwzględną większością głosów oddanych.
  \item „Głosy Uprawnione w Sprawie” oznaczają głosy z akcji, z których zgodnie z ustawą albo wyraźnym postanowieniem niniejszej umowy prawo głosu może być wykonywane w danej sprawie. Jeżeli akcjonariusz jest wyłączony od głosowania, głosów z jego akcji nie uwzględnia się ani w liczniku, ani w mianowniku większości określonej jako procent Głosów Uprawnionych w Sprawie.
  
\begin{MemoBox}{2.2.3}
\MemoPyt{2.2.3}{Czy dwa sposoby liczenia większości (od głosów uprawnionych w sprawie i od wszystkich głosów) są jasne i zgodne z KSH, zwłaszcza przy wyłączeniu głosu akcjonariusza występującego o zgodę (§~13 ust.~1, §~25 ust.~3)?}
\Ocena{zgodne warunkowo}
\Klasa{zalecane}
\Podstawa{art.~300$^{98}$ §~1--2 KSH (bezwzględna większość jako reguła; 3/4 głosów w sprawach z §~2, „chyba że umowa spółki przewiduje surowsze warunki”) \Wiar{Z}; art.~300$^{98}$ §~3 KSH (zmiana umowy „uszczuplająca prawa indywidualne poszczególnych akcjonariuszy, wymaga zgody wszystkich akcjonariuszy, których dotyczy”) \Wiar{Z}; art.~4 §~1 pkt~9--10 KSH (głosy oddane; bezwzględna większość) \Wiar{Z}; art.~300$^{96}$ KSH (ustawowe wyłączenie od głosowania: odpowiedzialność, absolutorium, zwolnienie z zobowiązania, spór) \Wiar{Z}; art.~300$^{1}$ §~3 KSH (akcjonariusz zobowiązany tylko do świadczeń określonych w umowie) \Wiar{Z}; art.~300$^{39}$ §~1--2 KSH (ograniczenia rozporządzania akcją w umowie) \Wiar{Z}.}
\Uzasadnienie{Konstrukcja jest jasna: §~\ref{art:wz} ust.~11 definiuje „Głosy Uprawnione w Sprawie”, a ust.~12 oddziela progi liczone od „wszystkich głosów w Spółce”. Oba sposoby są dopuszczalne --- próg od wszystkich głosów jest surowszy od ustawowych większości głosów oddanych, a art.~300$^{98}$ pozwala na zaostrzenie. Rachunek działa: przy wniosku Założyciela z 50\% próg 75\% to 38 z 50 głosów (20+10+10 wystarcza), przy wniosku Założyciela z 10\% --- 68 z 90 (50+20). Dwa problemy: (1) wyłączenie „występującego o zgodę” (§~\ref{art:okres-ograniczenia-zbywania}, §~\ref{art:walne-zastrzezone} ust.~3) wynika wyłącznie z umowy, bo art.~300$^{96}$ obejmuje tylko odpowiedzialność, zwolnienie z zobowiązania i spór; katalog ustawowy w S.A. uważa się za zamknięty, a w P.S.A. (art.~300$^{1}$ §~3, art.~300$^{39}$ §~2) wyłączenie ma mocne oparcie, ale nie jest wolne od wątpliwości \Wiar{?} --- bezpieczniejsza jest podwójna konstrukcja: wyłączenie i zobowiązanie do niewykonywania głosu. Wyłączenie z pierwotnej umowy wiąże Założycieli; jego późniejsze dodanie albo rozszerzenie (także proponowane uzupełnienie ust.~11) wymaga zgody wszystkich akcjonariuszy, których dotyczy (art.~300$^{98}$ §~3), a po wejściu inwestora --- także jego zgody. (2) Przy wyłączeniu ustawowym w Sprawie Zastrzeżonej (np. lit.~l --- ugoda w sporze o IP z akcjonariuszem mającym ponad 25\% głosów; lit.~m --- Odejście Usprawiedliwione, jeżeli wiąże się ze zwolnieniem z zobowiązania) próg 75\% wszystkich głosów staje się nieosiągalny, a ust.~12 nie rozstrzyga, jak wtedy liczyć --- potrzebna jest reguła konwersji na Głosy Uprawnione w Sprawie.}
\Rekomendacja{Dodać do §~\ref{art:wz} ust.~12 regułę konwersji przy wyłączeniu ustawowym oraz do ust.~11 zobowiązanie akcjonariusza wyłączonego umownie do niewykonywania głosu. Decyzja Założycieli: czy reguła konwersji ma dotyczyć wszystkich Spraw Zastrzeżonych, czy tylko tych, w których wyłączenie ustawowe jest realne (lit.~l, m).}
\Kierunek{Usuwa ryzyko trwałej blokady uchwał i wzmacnia wyłączenie umowne; dla inwestora --- przewidywalny rachunek głosów. Poniżej uzupełnienie ust.~11; brzmienie ust.~12 --- w ramce przy ust.~12.}
\begin{Brzmienie}
11. [\ldots] Akcjonariusz wyłączony od głosowania na podstawie wyraźnego postanowienia niniejszej umowy zobowiązuje się nie wykonywać prawa głosu w danej sprawie; głos oddany wbrew temu zobowiązaniu nie jest uwzględniany przy ustalaniu wyniku głosowania.
\end{Brzmienie}
\end{MemoBox}

\item Zasada z ust. 11 nie zmienia większości określonych w niniejszej umowie jako procent „wszystkich głosów w Spółce”; takie progi oblicza się od pełnej liczby głosów istniejących w Spółce, chyba że bezwzględnie obowiązujący przepis prawa stanowi inaczej.

\begin{MemoBox}{2.2.3}
\MemoPytCd{2.2.3}
\begin{Brzmienie}
12. Zasada z ust.~11 nie zmienia większości określonych w niniejszej umowie jako procent „wszystkich głosów w Spółce”; takie progi oblicza się od pełnej liczby głosów istniejących w Spółce. Jeżeli jednak w danej sprawie akcjonariusz jest wyłączony od głosowania z mocy bezwzględnie obowiązującego przepisu prawa, próg wyrażony jako procent wszystkich głosów w Spółce oblicza się w tej sprawie od Głosów Uprawnionych w Sprawie.
\end{Brzmienie}
\end{MemoBox}

\begin{MemoBox}{3.4.4}
\MemoPytCd{3.4.4}
\begin{Brzmienie}
§~\ref{art:wz} ust.~12 (nowy ust.~11): „Zasada z ust.~10 nie zmienia większości określonych w niniejszej umowie jako procent „wszystkich głosów w Spółce”; takie progi oblicza się od pełnej liczby głosów istniejących w Spółce, chyba że bezwzględnie obowiązujący przepis prawa stanowi inaczej.”
\end{Brzmienie}
\end{MemoBox}

\end{enumerate}
\end{ContractClause}

\begin{ContractClause}{Sprawy Zastrzeżone dla Walnego Zgromadzenia}
\label{art:walne-zastrzezone}
\begin{enumerate}
  \item Uchwały w następujących sprawach wymagają większości 75\% wszystkich głosów w Spółce:
  \begin{enumerate}[label=\alph*)]
    \item zmiana umowy Spółki, z zastrzeżeniem spraw wymagających większości surowszej na podstawie ustawy albo niniejszej umowy;
    \item emisja akcji, warrantów, opcji, instrumentów zamiennych lub innych praw do akcji, w tym każda emisja akcji serii P oraz każda emisja akcji serii F;
    \item nadanie albo, przed złożeniem wiążącej oferty objęcia akcji, cofnięcie statusu Założyciela Funkcjonalnego Serii F;
    \item zwiększenie, odnowienie albo istotna zmiana programu motywacyjnego;
    \item połączenie, podział, przekształcenie, rozwiązanie lub likwidacja Spółki;
    \item sprzedaż przedsiębiorstwa, zorganizowanej części przedsiębiorstwa albo aktywów stanowiących ponad 25\% wartości aktywów Spółki;
    \item przeniesienie, zbycie albo udzielenie wyłącznej, nieodwołalnej lub bezterminowej licencji na Kluczową Własność Intelektualną poza zwykłym tokiem działalności;
    \item istotna zmiana profilu działalności, rezygnacja z działalności badawczo-rozwojowej lub zmiana zasadniczej strategii produktów podwójnego zastosowania;
    \item zatwierdzenie wynagrodzenia dyrektorów oraz przyjęcie programu ubezpieczenia odpowiedzialności; powołanie i odwołanie dyrektorów następuje większościami określonymi w §~\ref{art:rada} ust.~2;
    \item nabycie akcji własnych, ich umorzenie lub wypłata z kapitału akcyjnego, jeżeli wymaga uchwały;
    \item zaciągnięcie zobowiązania, ustanowienie zabezpieczenia lub dokonanie transakcji poza zatwierdzonym budżetem o wartości przekraczającej 500 000 zł;
    \item zawarcie ugody albo zrzeczenie się roszczenia dotyczącego własności intelektualnej lub naruszenia bezpieczeństwa o wartości przekraczającej 500 000 zł;
    \item zatwierdzenie Odejścia Usprawiedliwionego w przypadkach uznaniowych;
    \item zaangażowanie Spółki w Przedsięwzięcie Produkcyjne przekraczające łącznie 500 000 zł, utrata kontroli nad Przedsięwzięciem Produkcyjnym, do którego Spółka wniosła Kluczową Własność Intelektualną albo aktywa stanowiące ponad 25\% wartości aktywów Spółki, oraz dopuszczenie partnera niespełniającego Kryterium Bezpieczeństwa EU/NATO w przypadkach określonych w §~\ref{art:przedsiewziecia} ust.~4.
  \end{enumerate}
  \item Pozbawienie prawa poboru w całości albo w części wymaga większości co najmniej czterech piątych głosów oraz spełnienia pozostałych wymogów bezwzględnie obowiązujących przepisów prawa.
  \item Zgoda na zbycie akcji Założyciela w Okresie Ograniczenia Zbywania Akcji wymaga większości 75\% Głosów Uprawnionych w Sprawie, zgodnie z §~\ref{art:wz}.
  \item Głos akcjonariusza pozostającego w konflikcie interesów może zostać wyłączony wyłącznie w zakresie dopuszczalnym przez prawo albo wyraźnie przewidzianym w niniejszej umowie.
\end{enumerate}

\begin{MemoBox}{2.2.4}
\MemoPyt{2.2.4}{Czy katalog Spraw Zastrzeżonych i próg 75\% nie blokują czynności, dla których ustawa przewiduje inną większość albo kompetencję Rady?}
\Ocena{zgodne warunkowo}
\Klasa{zalecane}
\Podstawa{art.~300$^{98}$ §~1--2 KSH (większości kwalifikowane 3/4; „chyba że umowa spółki przewiduje surowsze warunki”) \Wiar{Z}; art.~300$^{81}$ KSH (czynności wymagające uchwały akcjonariuszy: pkt~2 zbycie i wydzierżawienie przedsiębiorstwa albo ZCP oraz ustanowienie na nich ograniczonego prawa rzeczowego --- bez opt-out; pkt~3 „nabycie i zbycie nieruchomości, użytkowania wieczystego lub udziału w nieruchomości, chyba że umowa spółki stanowi inaczej”) \Wiar{Z}; art.~300$^{75}$ §~2 KSH (decyzje o strategicznym znaczeniu, plany biznesowe, struktura organizacyjna --- uchwała Rady) \Wiar{Z}; art.~17 §~1--3 KSH \Wiar{Z}; art.~15 §~1 KSH (zgoda WZ na kredyt, pożyczkę, poręczenie lub podobną umowę z członkiem zarządu, rady nadzorczej, komisji rewizyjnej, prokurentem, likwidatorem --- nie wymienia dyrektora P.S.A.) \Wiar{Z}, stosowanie art.~15 do dyrektora w drodze wykładni \Wiar{?}; art.~506 §~1, 541 §~1, 575 i 577 §~1 pkt~1 KSH (każdorazowo „chyba że umowa [\ldots] przewiduje surowsze warunki”) \Wiar{Z}.}
\Uzasadnienie{Zaostrzenie jest dozwolone, więc żadna pozycja katalogu nie jest sprzeczna z ustawową większością: 75\% wszystkich głosów jest surowsze niż 3/4 głosów oddanych (art.~300$^{98}$ §~2) i niż większości przy łączeniu, podziale i przekształceniu, a pozbawienie prawa poboru zachowuje 4/5 (ust.~2). (1) Lit.~h i k nakładają się na „decyzje o strategicznym znaczeniu” zastrzeżone dla Rady (art.~300$^{75}$ §~2); umowa nie może przenieść tej decyzji na WZ, ale może wymagać zgody WZ jako warunku dodatkowego (art.~17 §~3 --- brak zgody nie wpływa na ważność czynności, rodzi odpowiedzialność); trzeba to nazwać, inaczej sąd może uznać, że WZ „prowadzi sprawy Spółki”. (2) Nieruchomości: art.~300$^{81}$ pkt~3 wymaga uchwały akcjonariuszy, „chyba że umowa spółki stanowi inaczej”; katalog milczy, więc wymóg dotyczy każdej nieruchomości, a brak uchwały powoduje nieważność czynności (art.~17 §~1). Opt-out nie jest możliwy dla pkt~2, który lit.~f) obejmuje tylko częściowo (bez dzierżawy i zastawu na przedsiębiorstwie). (3) Art.~15 §~1 KSH nie wymienia dyrektora P.S.A.; objęcie dyrektorów tym przepisem bywa bronione wykładnią funkcjonalną, ale sankcja nieważności (art.~17 §~1) jest na tej podstawie niepewna \Wiar{?}, a §~\ref{art:rada-zastrzezone} ust.~1 lit.~e) przewiduje tylko uchwałę Rady --- umowa powinna sama wymagać zgody WZ. Blokada operacyjna: przy 50/20/10/10/10 każde 26\% blokuje; lit.~i (wynagrodzenie, D\&O) i k to sprawy bieżące, a ich zamrożenie przy sporze uruchamia dopiero §~\ref{art:impas} po dwóch zgromadzeniach.}
\Rekomendacja{Dodać do §~\ref{art:walne-zastrzezone} ustęp wyjaśniający charakter zgody WZ w sprawach należących do Rady, odesłanie do zgód ustawowych (art.~300$^{81}$) oraz własny, umowny wymóg zgody WZ na kredyt, pożyczkę, poręczenie i podobną umowę z dyrektorem albo prokurentem (bo art.~15 KSH nie wymienia dyrektora); zdecydować o nieruchomościach (proponowane: opt-out do progu z lit.~k); rozszerzyć lit.~f) o wydzierżawienie i obciążenie przedsiębiorstwa lub ZCP; rozważyć obniżenie większości dla lit.~i (bezwzględna większość).}
\Kierunek{Zmiana zapobiega nieważności czynności dokonanych bez zgody ustawowej (nieruchomości, przedsiębiorstwo), zamyka lukę art.~15 KSH wobec dyrektorów i chroni podział kompetencji przed zarzutem, że WZ przejęło prowadzenie spraw. Proponowany nowy ust.~5 (w skrócie): (a) uchwały WZ w sprawach, które ustawa zastrzega do decyzji Rady (w szczególności art.~300$^{75}$ §~2 KSH), stanowią zgodę na dokonanie czynności i nie zastępują uchwały Rady; (b) niezależnie od ust.~1 uchwały WZ wymagają czynności wskazane w bezwzględnie obowiązujących przepisach, w szczególności w art.~300$^{81}$ KSH, a także umowa kredytu, pożyczki, poręczenia lub inna podobna umowa z dyrektorem, prokurentem albo likwidatorem albo na ich rzecz, niezależnie od tego, czy zgody wymaga art.~15 KSH; (c) na podstawie art.~300$^{81}$ pkt~3 KSH nabycie i zbycie nieruchomości, użytkowania wieczystego albo udziału w nieruchomości o wartości nieprzekraczającej progu z ust.~1 lit.~k) nie wymaga uchwały WZ i należy do Rady zgodnie z §~\ref{art:rada-zastrzezone} ust.~1 lit.~d).}
\end{MemoBox}

\end{ContractClause}


\begin{ContractClause}{Własność intelektualna i aktywa technologiczne}
\label{art:wlasnosc-intelektualna}
\begin{enumerate}
  \item „Kluczowa Własność Intelektualna” oznacza wszelkie prawa i aktywa niezbędne albo istotne dla produktów i działalności Spółki, w szczególności:
  \begin{enumerate}[label=\alph*)]
    \item kod źródłowy i wynikowy, modele uczenia maszynowego, wagi, architektury, algorytmy, biblioteki, interfejsy i dokumentację;
    \item dane treningowe, testowe, telemetryczne, mapy, zbiory syntetyczne, wyniki eksperymentów i procedury walidacji;
    \item projekty mechaniczne, elektroniczne, schematy, płytki drukowane, pliki CAD, BOM, firmware i dokumentację produkcyjną;
    \item wynalazki, patenty, zgłoszenia patentowe, prawa do uzyskania patentu, wzory, znaki towarowe i tajemnice przedsiębiorstwa;
    \item domeny, repozytoria, konta chmurowe, klucze, certyfikaty, konfiguracje, numery telefoniczne, profile i identyfikatory używane dla działalności Spółki.
  \end{enumerate}
  \item Kluczowa Własność Intelektualna rozwijana dla Spółki lub z wykorzystaniem jej zasobów ma przysługiwać Spółce w najszerszym prawnie dopuszczalnym zakresie.
  
\begin{MemoBox}{3.1.3}
\MemoPyt{3.1.3}{Czy oświadczenia z §~\ref{art:wlasnosc-intelektualna} ust.~5 i obowiązek współdziałania z §~\ref{art:wlasnosc-intelektualna} ust.~7 wystarczą, aby przed rundą wykazać tytuł Spółki do kodu, modeli i zgłoszeń patentowych; jakie odrębne umowy przeniesienia (pola eksploatacji, forma pisemna, wpisy) są konieczne?}
\Ocena{zgodne warunkowo (warunek spełnia się poza umową)}
\Klasa{zalecane}
\Podstawa{art.~9 ust.~1--3, art.~12 ust.~1, art.~14 ust.~1, art.~16, art.~41 ust.~2--4, art.~46, art.~50, art.~53, art.~74 ust.~1, 3 i~4 ustawy o prawie autorskim i prawach pokrewnych \Wiar{Z}; art.~11 ust.~1--3, art.~12 ust.~1--2, art.~20, art.~67 ust.~2--3, art.~72 ust.~1 ustawy Prawo własności przemysłowej \Wiar{Z}; art.~11 ust.~1--2 ustawy o zwalczaniu nieuczciwej konkurencji \Wiar{Z}; ustawa o ochronie baz danych \Wiar{W}; art.~300$^{11}$ §~1--2 KSH \Wiar{Z}.}
\Uzasadnienie{Oświadczenia z ust.~5 (zapewnienia) i obowiązek z ust.~7 nie przenoszą praw (przyznaje to ust.~4), a inwestor bada łańcuch tytułu. Konieczne: (1)~\emph{kod i dokumentacja} --- pisemna pod rygorem nieważności (art.~53) umowa przeniesienia praw majątkowych: znane pola eksploatacji wymienione wprost (art.~41 ust.~2 i~4; dla programów art.~74 ust.~4), prawo zezwalania na wykonywanie praw zależnych (art.~46), zobowiązanie do niewykonywania praw osobistych (art.~16), oznaczenie utworów; umowa o wszystkie przyszłe utwory (także danego rodzaju) twórcy jest nieważna (art.~41 ust.~3 \Wiar{Z}), więc ust.~2 nie zastąpi klauzul w umowach o pracę i B2B; cesję podpisuje każdy współtwórca (art.~9 ust.~1 i~3 \Wiar{Z}); (2)~\emph{łańcuch przed Założycielem} --- tytuł pracodawcy (art.~12 ust.~1, art.~74 ust.~3), a do wynalazków także zamawiającego z umowy B2B (art.~11 ust.~3 PWP) \Wiar{Z}, pierwszeństwo publikacji uczelni (art.~14 ust.~1 \Wiar{Z}); (3)~\emph{modele} --- wagi i dane mogą nie być utworami (ochrona: tajemnica przedsiębiorstwa, art.~11 ust.~2 u.z.n.k. \Wiar{Z}, i prawo sui generis do bazy danych), stąd też nośniki, bazy i poufność; (4)~\emph{zgłoszenia patentowe} --- pisemna pod rygorem nieważności cesja prawa do uzyskania patentu (art.~12 ust.~1--2 PWP \Wiar{Z}), zmiana zgłaszającego w UPRP \Wiar{W}, po udzieleniu --- wpis do rejestru patentowego (art.~67 ust.~2--3 \Wiar{Z}); EPO: Rule~22 EPC \Wiar{W}, PCT: Rule~92bis \Wiar{W}; dla wyłączności --- cesja od każdego współtwórcy (art.~72 ust.~1 \Wiar{Z}); (5)~\emph{open source} --- copyleft wyklucza wyłączność Spółki co do pochodnych. Przeniesienie na Spółkę w organizacji jest dopuszczalne (art.~300$^{11}$ §~1--2 KSH \Wiar{Z}); przy prawach wymagających wpisu --- skutek po wpisie Spółki.}
\Rekomendacja{Zawrzeć w dniu podpisania umowy trzy rodzaje dokumentów przeniesienia (kod i dane; sprzęt; prawa do wynalazków) według powyższych wymogów, a w umowach o pracę i B2B każdego Założyciela i pracownika --- klauzule przeniesienia praw do utworów powstałych w wykonaniu danej umowy, z chwilą ich przyjęcia i z okresowymi protokołami przekazania (granica z art.~41 ust.~3). W §~\ref{art:wlasnosc-intelektualna} dodać obowiązek zawarcia takich umów, aby jego naruszenie było naruszeniem umowy spółki.}
\Kierunek{Zmiana nie tworzy przeniesienia, lecz obowiązek korporacyjny; wzmacnia pozycję Spółki w sporze z Założycielem i odpowiada na pytanie inwestora o podstawę żądania podpisania dokumentów. Proponowane uzupełnienie ust.~2 (w skrócie): każdy Założyciel zobowiązuje się zawrzeć ze Spółką umowę w formie pisemnej pod rygorem nieważności --- co do wkładów najpóźniej w dniu zawarcia umowy spółki, a co do utworów, wynalazków i innych dóbr powstających w wykonaniu jego umowy o pracę albo o współpracę --- w terminie 30 dni od wpisu Spółki do rejestru; umowa przenosi, w zakresie dopuszczonym przez art.~41 ust.~3 ustawy o prawie autorskim i prawach pokrewnych, autorskie prawa majątkowe na wszystkich znanych polach eksploatacji (w tym z art.~74 ust.~4) wraz z prawem zezwalania na wykonywanie praw zależnych, zobowiązaniem do niewykonywania autorskich praw osobistych oraz przeniesieniem praw do uzyskania patentu, praw ze zgłoszeń i praw do baz danych.}
\end{MemoBox}

\item Spółka zapewnia, aby każda osoba uzyskująca dostęp do Kluczowej Własności Intelektualnej, w szczególności pracownik, współpracownik lub podwykonawca, zawarła uprzednio odpowiednią umowę przenoszącą prawa albo zapewniającą Spółce wyłączną, nieodwołalną i przenoszalną podstawę korzystania, obejmującą wszystkie wymagane pola eksploatacji i prawo wykonywania praw zależnych.
  \item Żadne postanowienie niniejszej umowy nie zastępuje odrębnego dokumentu przeniesienia, gdy przepisy wymagają oznaczenia utworu, pól eksploatacji, formy pisemnej, wpisu do rejestru lub innej szczególnej czynności.
  \item Założyciel oświadcza w odniesieniu do wkładów i innych przekazywanych aktywów, że:
  \begin{enumerate}[label=\alph*)]
    \item przysługują mu prawa wystarczające do skutecznego przeniesienia;
    \item aktywa nie naruszają praw osób trzecich i nie są obciążone;
    \item ujawnił wykorzystane komponenty otwartego oprogramowania, ich licencje i obowiązki;
    \item ujawnił wszelkie powiązania z wcześniejszym pracodawcą, uczelnią, grantem, zamówieniem publicznym, projektem finansowanym lub współtwórcą;
    \item nie istnieją nieujawnione ograniczenia eksportowe, licencyjne, publikacyjne ani dotyczące poufności.
  \end{enumerate}
  \item Rada Dyrektorów prowadzi rejestr własności intelektualnej, repozytoriów, komponentów otwartego oprogramowania i materiałów pochodzenia zewnętrznego oraz zapewnia utrzymywanie zestawienia składników oprogramowania.
  \item Założyciele są zobowiązani do podpisywania dokumentów, składania oświadczeń i udzielania pomocy niezbędnej do rejestracji, ochrony, egzekwowania i obrony praw Spółki, także po ustaniu współpracy, za zwrotem uzasadnionych kosztów.
\end{enumerate}

\end{ContractClause}

\begin{ContractClause}{Poufność, bezpieczeństwo informacji i cyberbezpieczeństwo}
\label{art:security}
\begin{enumerate}
  \item Informacje techniczne, handlowe, finansowe, operacyjne i organizacyjne Spółki, w szczególności Kluczowa Własność Intelektualna, są poufne, o ile nie zostały zgodnie z prawem podane do publicznej wiadomości.
  \item Akcjonariusz może korzystać z informacji poufnych wyłącznie w celu wykonywania praw i obowiązków wobec Spółki oraz nie może udostępniać ich osobom trzecim bez zgody Rady Dyrektorów.
  \item Obowiązek poufności trwa przez okres posiadania akcji i przez 10 lat po jego zakończeniu, a w odniesieniu do tajemnicy przedsiębiorstwa — tak długo, jak informacja zachowuje taki charakter.
  
\begin{MemoBox}{4.2.3}
\MemoPyt{4.2.3}{Czy 10-letni okres poufności i ograniczenie dostępu spoza UE/NATO (§ 27 ust. 3 i 6) są egzekwowalne wobec akcjonariuszy i dyrektorów, w tym po ustaniu ich statusu?}
\Ocena{zgodne warunkowo}
\Klasa{zalecane}
\Podstawa{art.~300$^{1}$ §~3 KSH („akcjonariusze są zobowiązani jedynie do świadczeń określonych w~umowie spółki”) \Wiar{Z}; art.~300$^{55}$ §~2 KSH („członek organu nie może ujawniać tajemnic spółki, także po wygaśnięciu mandatu”) \Wiar{Z}; art.~11 ust.~1--2 i~4 oraz art.~18 ust.~1 ustawy o~zwalczaniu nieuczciwej konkurencji \Wiar{Z}; art.~353$^{1}$ KC \Wiar{W}; art.~483--484 KC \Wiar{Z}; art.~2 pkt~2 lit.~d i~pkt~3 lit.~b, art.~3 ust.~1, art.~8 ust.~1--3 oraz załącznik~II sekcje~A--F (EU001--EU006) rozporządzenia (UE) 2021/821 \Wiar{Z} (tekst EN).}
\Uzasadnienie{\textbf{Akcjonariusze.} Poufność jako świadczenie z~umowy spółki wiąże każdego akcjonariusza, także nabywcę wtórnego \Wiar{W}, jest ujawnialna jako obowiązek związany z~akcją i~wzmocniona przez uznk (art.~11 ust.~4, roszczenia z~art.~18 ust.~1). \textbf{Po ustaniu statusu} --- punkt sporny: przeważa pogląd, że obowiązki przetrwałe wiążą byłego akcjonariusza, jeżeli umowa wyraźnie przewiduje ich trwanie (art.~353$^{1}$ KC) \Wiar{W}; pogląd przeciwny widzi w~tym wyjście poza materię umowy spółki \Wiar{?}. 10 lat nie jest nadmierne dla technologii obronnej i~nie wymaga ekwiwalentu. Słabości: brak kary umownej i~brak zdania, że utrata statusu nie uchyla obowiązków; bezpieczne jest podwójne zakotwiczenie (umowa spółki i~dokument indywidualny). \textbf{Dyrektor} niebędący akcjonariuszem nie jest stroną umowy; wiąże go art.~300$^{54}$ i~art.~300$^{55}$ §~2, bez ograniczenia czasowego, ale bez sankcji umownej (poprzednia wersja memorandum błędnie lokowała ten zakaz w~art.~300$^{54}$); ust.~3 nie określa dla niego okresu --- właściwym mechanizmem są umowy Spółki z~dyrektorami. \textbf{Ust.~6} to reguła miejsca dostępu, zgodna z~kontrolą eksportu: zdalny dostęp z~państwa trzeciego do kontrolowanej technologii jest wywozem, a~udostępnienie technologii z~załącznika~I w~innym państwie NATO (np.~w~Turcji) także wymaga zezwolenia \Wiar{Z} (tekst EN), więc zgoda Rady musi być poprzedzona analizą kontroli eksportu. Usterki: brak adresata obowiązku i~pominięcie EOG (Liechtenstein), choć Kryterium z~§~\ref{art:dopuszczalny-nabywca} ust.~6 lit.~a je obejmuje.}
\Rekomendacja{Dodać w~ust.~3 klauzulę przetrwania i~objąć nią dyrektorów przez obowiązek Spółki zawarcia z~nimi umowy; w~ust.~6 wskazać adresata i~dodać EOG. Karę umowną umieścić w~umowie wykonawczej i~umowach objęcia akcji, nie w~umowie spółki.}
\Kierunek{Dla akcjonariuszy: jasność, że obowiązek trwa po zbyciu akcji; dla dyrektorów: obowiązek umowny obok ustawowego z~art.~300$^{55}$ §~2; dla inwestora: standardowe oczekiwanie funduszy (NDA założycieli); dla rejestru: bez zmian; dla sądu: bez znaczenia. Proponowane brzmienie §~\ref{art:security} ust.~3 (ust.~6 --- w~ramce przy ust.~6):}
\begin{Brzmienie}
3. Obowiązek poufności wiąże akcjonariusza przez okres posiadania akcji i przez 10 lat po jego zakończeniu, a w odniesieniu do informacji stanowiących tajemnicę przedsiębiorstwa --- tak długo, jak informacja zachowuje taki charakter. Utrata statusu akcjonariusza nie uchyla obowiązków powstałych w czasie jego trwania. Spółka zapewnia, aby każdy dyrektor niebędący akcjonariuszem zobowiązał się na piśmie, najpóźniej w dniu powołania, do przestrzegania niniejszego paragrafu przez okres pełnienia funkcji i przez 10 lat po jego zakończeniu, oraz aby każda osoba obejmująca albo nabywająca akcje potwierdziła na piśmie związanie niniejszym paragrafem.
\end{Brzmienie}
\end{MemoBox}

\item Spółka wdraża środki bezpieczeństwa adekwatne do ryzyka, w szczególności kontrolę dostępu, zasadę minimalnych uprawnień, rozdzielenie środowisk, szyfrowanie, rejestrowanie zdarzeń, bezpieczny cykl wytwarzania, zarządzanie podatnościami i reagowanie na incydenty.
  \item Każdy akcjonariusz i dyrektor ma obowiązek niezwłocznie zgłosić podejrzenie wycieku, naruszenia, utraty urządzenia, nieuprawnionego dostępu, próby pozyskania technologii lub konfliktu interesów.
  \item Dostęp z jurysdykcji spoza Unii Europejskiej i NATO do Kluczowej Własności Intelektualnej wymaga uprzedniej pisemnej zgody Rady Dyrektorów i analizy kontroli eksportu; zgoda nie może zostać udzielona, jeżeli byłoby to sprzeczne z §~\ref{art:dopuszczalny-nabywca} lub prawem.

\begin{MemoBox}{4.2.3}
\MemoPytCd{4.2.3}
\Kierunek{Ust.~6: wskazanie adresata zakazu, uzupełnienie katalogu terytoriów o~EOG zgodnie z~Kryterium, zgoda Rady poprzedzona analizą kontroli eksportu. Proponowane brzmienie §~\ref{art:security} ust.~6:}
\begin{Brzmienie}
6. Akcjonariusz, dyrektor ani osoba działająca na rzecz Spółki nie może uzyskiwać dostępu do Kluczowej Własności Intelektualnej z terytorium państwa spoza Unii Europejskiej, Europejskiego Obszaru Gospodarczego i NATO ani udzielać takiego dostępu innej osobie bez uprzedniej pisemnej zgody Rady Dyrektorów poprzedzonej analizą kontroli eksportu; zgoda nie może zostać udzielona, jeżeli byłoby to sprzeczne z §~\ref{art:dopuszczalny-nabywca}, §~\ref{art:export} lub prawem.
\end{Brzmienie}
\end{MemoBox}

\end{enumerate}

\end{ContractClause}


\begin{ContractClause}{Kontrola eksportu, sankcje i końcowe zastosowanie}
\label{art:export}
\begin{enumerate}
  \item Spółka, jej organy i akcjonariusze przestrzegają obowiązujących przepisów dotyczących produktów podwójnego zastosowania, uzbrojenia, transferu technologii, pomocy technicznej, pośrednictwa, sankcji, embarg i ograniczeń terytorialnych, w szczególności właściwych regulacji Unii Europejskiej i prawa polskiego.
  
\begin{MemoBox}{4.2.1}
\MemoPyt{4.2.1}{Czy klauzule wewnętrznej zgodności wystarczają jako ramy, czy projekt powinien wprost odsyłać do konkretnych aktów (rozporządzenie UE o produktach podwójnego zastosowania, ustawa o obrocie z zagranicą towarami o znaczeniu strategicznym)?}
\Ocena{zgodne}
\Klasa{zalecane --- odstępstwo od reguły dla oceny „zgodne” uzasadnione tym, że ust.~5 w~obecnym brzmieniu nie jest wykonalny}
\Podstawa{rozporządzenie (UE) 2021/821 (tekst skonsolidowany na 26 maja 2023~r., EN): art.~2 pkt~1--3, 9--10 i~21, art.~4 ust.~1--2, art.~8 ust.~1--3, art.~11 ust.~1 i~9, art.~12 ust.~1 i~4, art.~17 ust.~1, art.~27 ust.~1, 3--4, załącznik~II sekcja~A (EU001) i~sekcja~G (EU007) \Wiar{Z} (tekst EN; załącznik~I zmieniany po 2023~r.); ustawa z~29 listopada 2000~r. o~obrocie z~zagranicą towarami, technologiami i~usługami o~znaczeniu strategicznym (t.j. Dz.U. 2023 poz.~1582, zm. Dz.U. 2026 poz.~471): art.~2, art.~6, art.~9 ust.~2 pkt~11, art.~11, art.~17a, art.~33 ust.~1 i~2a \Wiar{Z}; rozporządzenie (UE) 2024/1689 (AI Act): art.~2 ust.~2, 3, 6 i~8, art.~3 pkt~1, art.~6 ust.~1, art.~108, art.~113, załącznik~I sekcja~B pkt~20 \Wiar{Z} (tekst EN; daty stosowania do sprawdzenia pod kątem późniejszych zmian \Wiar{?}); rozporządzenie (UE) 2021/697, art.~9 ust.~4 \Wiar{Z}; rozporządzenia sankcyjne UE (m.in. 833/2014, 269/2014) \Wiar{W}; art.~300$^{33}$ §~1 pkt~11 KSH \Wiar{Z}; art.~64 KC \Wiar{Z}.}
\Uzasadnienie{Klauzule ramowe są lepsze niż odesłania sztywne: załącznik~I do rozporządzenia 2021/821 jest zmieniany aktami delegowanymi (art.~17 ust.~1 lit.~a) \Wiar{Z} (tekst EN) --- w~praktyce corocznie \Wiar{W} (dostępna wersja skonsolidowana --- z~26 maja 2023~r., później zmieniana w~latach 2023--2025); ustawa krajowa została znowelizowana od 22 kwietnia 2026~r., a~reżim sankcyjny zmienia się co kilka miesięcy. §~\ref{art:export} ust.~1 („obowiązujące przepisy”) jest odesłaniem dynamicznym, a~ust.~2--4 przenoszą materię wykonawczą do programu zgodności Rady. Rozporządzenie potwierdza klauzulę: wywozem jest też udostępnienie oprogramowania lub technologii w~formie elektronicznej poza obszarem celnym Unii (art.~2 pkt~2 lit.~d); ICP jest co do zasady warunkiem zezwolenia globalnego i~EU007; dokumenty transferu wewnątrzunijnego muszą wskazywać, że produkty podlegają kontroli (art.~11 ust.~9) \Wiar{Z} (tekst EN). AI Act: wyłączenie wojskowe jest zastosowaniowe, nie podmiotowe (art.~2 ust.~3); cywilne BSP objęte są reżimem art.~2 ust.~2 (głównie pośrednio, przez akty do 2018/1139); B+R przed wprowadzeniem do obrotu jest wyłączone, poza testami w~warunkach rzeczywistych \Wiar{Z} (tekst EN); daty z~art.~113 --- do sprawdzenia (Digital Omnibus) \Wiar{?}. Słabość: ust.~5 („prawnie dopuszczalne środki prowadzące do zbycia”) nie określa nabywcy, ceny ani terminu, więc nie da się go ujawnić w~rejestrze (art.~300$^{33}$ §~1 pkt~11) ani wykonać w~trybie art.~64 KC; wzorzec daje §~\ref{art:dopuszczalny-nabywca} ust.~15.}
\Rekomendacja{Utrzymać klauzule ramowe. Uzupełnić ust.~1 o~przykładowe wyliczenie aktów i~przebudować ust.~5 w~skonkretyzowany obowiązek zbycia (brzmienie --- ramka przy ust.~5). W~programie zgodności z~ust.~4 wskazać osobę odpowiedzialną za koordynację kontroli obrotu (art.~9 ust.~2 pkt~11 ustawy), ująć oznaczanie dokumentów handlowych i~umów licencyjnych (art.~11 ust.~9 rozporządzenia 2021/821), ewidencję (art.~27) oraz ocenę produktów pod kątem AI Act (art.~2 ust.~2, 3 i~8). Decyzja Założycieli: zbycie po Wartości Godziwej (neutralne dla inwestora) czy z~dyskontem (co fundusze VC zakwestionują --- zob. pakiet 4.4).}
\Kierunek{Ust.~1: brak skutków poza czytelnością; ułatwia due diligence inwestorów i~programów, nie usztywniając klauzuli. Proponowane brzmienie ust.~1 (w~skrócie): po dotychczasowym katalogu materii --- „w~szczególności” rozporządzenie (UE) 2021/821 wraz z~załącznikami, ustawa o~obrocie z~zagranicą towarami, technologiami i~usługami o~znaczeniu strategicznym, ustawa o~wykonywaniu działalności gospodarczej w~zakresie wytwarzania i~obrotu materiałami wybuchowymi, bronią, amunicją oraz wyrobami i~technologią o~przeznaczeniu wojskowym lub policyjnym, akty prawa UE i~prawa polskiego ustanawiające środki ograniczające oraz --- w~zakresie, w~jakim ma zastosowanie do produktów Spółki --- rozporządzenie (UE) 2024/1689 w~sprawie sztucznej inteligencji, „każdorazowo w~brzmieniu obowiązującym w~dniu dokonania czynności”.}
\end{MemoBox}

\item Przed eksportem, transferem, udostępnieniem zdalnym, demonstracją, publikacją, dostarczeniem kodu, danych technicznych, sprzętu lub pomocy technicznej Spółka przeprowadza klasyfikację produktu, odbiorcy, jurysdykcji, końcowego użytkownika i końcowego zastosowania.
  
\begin{MemoBox}{4.2.2}
\MemoPyt{4.2.2}{Czy klasyfikacja produktów Firmy (oprogramowanie autonomii, roje dronów) wymaga odrębnej weryfikacji specjalistycznej i czy mogą Państwo ją wykonać albo wskazać specjalistę?}
\Ocena{zgodne warunkowo --- klauzula wystarcza; warunek (klasyfikacja) spełnia się poza umową}
\Klasa{opcjonalne --- warunek spełnia się poza umową}
\Podstawa{rozporządzenie (UE) 2021/821, załącznik~I (tekst skonsolidowany na 26 maja 2023~r., EN; załącznik zmieniany po tej dacie): uwagi ogólne do technologii (GTN) i~do oprogramowania (GSN), pozycje 9A012.a, 9D001--9D002, 9D004.e, 9E001, 7D002--7D004, 7E004.b, 6A008, 5A002.a, 5D002 oraz załącznik~IV (1C101, 1D103) \Wiar{Z} (tekst EN); art.~3 ust.~1, art.~4 ust.~1--2 i~art.~11 ust.~1 rozporządzenia \Wiar{Z} (tekst EN); Wspólny wykaz uzbrojenia UE, ML10 (BSP) i~ML21/ML22 \Wiar{W}; ustawa z~29 listopada 2000~r.: art.~10 ust.~1 (wiążące wyjaśnienie w~sprawie konieczności uzyskania zezwolenia), art.~17a ust.~1 pkt~2 (catch-all) \Wiar{Z}; §~\ref{art:export} ust.~2 i~§~\ref{art:rada-zastrzezone} ust.~1 lit.~h.}
\Uzasadnienie{Tak, wymaga. Klasyfikacja porównuje parametry konkretnego produktu (czas lotu, odporność na podmuchy, dokładność nawigacji inercyjnej, długość klucza, „specjalne zaprojektowanie” warstwy autonomii) z~progami wykazu i~ocenia wyłączenia oraz catch-all. Kancelaria może zaopiniować reżim, ale przypisanie pozycji wymaga danych inżynierskich i~jest obowiązkiem Spółki; niniejsze memorandum klasyfikacji nie wykonuje. Ścieżka: (1)~karta klasyfikacyjna inżyniera, (2)~weryfikacja przez wyspecjalizowanego doradcę (kancelarie z~praktyką obrotu strategicznego, firmy audytujące i~jednostki certyfikujące WSK) \Wiar{W}, (3)~przy wątpliwościach wiążące wyjaśnienie organu kontroli obrotu w~3 miesiące, przedłużalne do 6 (art.~10 ust.~1) \Wiar{Z}, (4)~uchwała Rady i~wpis do Załącznika nr~3 pkt~4. Wstępna hipoteza (9A012.b.2) wymaga korekty: w~wersji z~2023~r. 9A012.b.1--b.2 „nie stosuje się”. Kandydaci: platforma --- 9A012.a (lot poza zasięgiem wzroku i~czas lotu co najmniej 30~minut przy podmuchach od 46,3~km/h albo co najmniej 1~godziny); oprogramowanie autonomii --- 9D004.e, ale tylko gdy platforma spełnia progi 9A012.a; technologia --- 9E001; niezależnie nawigacja inercyjna i~hybrydowa --- 7D002--7D004, 7E004.b; łączność szyfrowana --- 5A002.a i~5D002; radar SAR --- 6A008.d \Wiar{Z} (tekst EN). Transfer wewnątrzunijny nie wymaga zezwolenia poza 1C101/1D103 (obniżona wykrywalność), lecz wymaga oznaczenia dokumentów \Wiar{Z} (tekst EN; do potwierdzenia na aktualnej wersji załączników). Czy platforma przekracza progi 9A012.a, a~nawigacja poziomy 7A003 --- pozostaje hipotezą \Wiar{?}. Klasyfikację powtarzać przy każdej istotnej zmianie produktu i~załącznika~I.}
\Rekomendacja{Zlecić klasyfikację specjaliście ds. kontroli eksportu przed bramką G2 (pierwsze udostępnienie technologii poza firmę, w~tym due diligence inwestora i~chmura poza UE); przy wątpliwościach wystąpić o~wiążące wyjaśnienie z~art.~10 ustawy, z~wyprzedzeniem wobec terminu do 6 miesięcy; wynik wpisać do Załącznika nr~3 pkt~4. W~umowie nie ma nic do zmiany.}
\Kierunek{Uzupełnienie Załącznika nr~3 pkt~4 o~datę klasyfikacji, wersję wykazu i~osobę odpowiedzialną; bez skutków prawnych dla akcjonariuszy i~sądu, istotne dla inwestora w~due diligence.}
\end{MemoBox}

\item Zakazane jest zawarcie transakcji lub udostępnienie technologii, jeżeli:
  \begin{enumerate}[label=\alph*)]
    \item wymagane zezwolenie nie zostało uzyskane;
    \item odbiorca, pośrednik, bank lub beneficjent jest objęty ograniczeniami;
    \item istnieją uzasadnione przesłanki ryzyka niedozwolonego zastosowania, reeksportu, obejścia sankcji, proliferacji albo wykorzystania sprzecznego z prawem międzynarodowym;
    \item nie można wiarygodnie ustalić końcowego użytkownika lub źródła finansowania.
  \end{enumerate}
  \item Rada Dyrektorów ustanawia wewnętrzny program zgodności eksportowej obejmujący klasyfikację, ewidencję, weryfikację kontrahentów, kontrolę dostępu, szkolenia, zatwierdzanie transakcji i zgłaszanie incydentów.
  \item Naruszenie niniejszego paragrafu przez Akcjonariusza Funkcjonalnego stanowi Odejście Zawinione, a przez innego akcjonariusza — istotne naruszenie umowy i podstawę do dochodzenia pełnego naprawienia szkody oraz zastosowania prawnie dopuszczalnych środków prowadzących do zbycia jego akcji.

\begin{MemoBox}{4.2.1}
\MemoPytCd{4.2.1}
\Kierunek{Ust.~5: dla akcjonariusza --- jasna sankcja zamiast klauzuli otwartej; dla inwestora --- przewidywalność (cena rynkowa); dla rejestru --- obowiązek staje się ujawnialny; dla sądu --- bez znaczenia. Proponowane brzmienie §~\ref{art:export} ust.~5:}
\begin{Brzmienie}
5. Naruszenie niniejszego paragrafu przez Akcjonariusza Funkcjonalnego stanowi Odejście Zawinione. Naruszenie niniejszego paragrafu przez innego akcjonariusza stanowi istotne naruszenie niniejszej umowy; Spółka może wówczas, w terminie 6 miesięcy od dnia stwierdzenia naruszenia uchwałą Rady Dyrektorów, wskazać Dopuszczalnego Nabywcę spełniającego Kryterium Bezpieczeństwa EU/NATO, na rzecz którego akcjonariusz jest zobowiązany zbyć wszystkie akcje za cenę równą Wartości Godziwej ustalonej zgodnie z §~\ref{art:wycena}, płatną na zasadach określonych w §~\ref{art:zbywanie} ust.~3. Obowiązek ten jest obowiązkiem związanym z akcją wobec Spółki, a do jego wykonania stosuje się odpowiednio §~\ref{art:zwrotne-zbycie} ust.~14--15 oraz §~\ref{art:breach}. Powyższe nie wyłącza roszczeń Spółki o naprawienie szkody w pełnej wysokości.
\end{Brzmienie}
\end{MemoBox}

\end{enumerate}

\end{ContractClause}

\begin{ContractClause}{Zakaz konkurencji, konflikt interesów i transakcje z podmiotami powiązanymi}
\label{art:konkurencja}
\begin{enumerate}
  \item Akcjonariusz Funkcjonalny, w okresie zaangażowania w Spółkę, nie może bez uprzedniej zgody Walnego Zgromadzenia prowadzić ani wspierać działalności konkurencyjnej dotyczącej produktów wskazanych w §~\ref{art:pkd}, uczestniczyć w podmiocie konkurencyjnym ani wykorzystywać szans biznesowych należących do Spółki.

  \item Zakaz nie obejmuje pasywnego posiadania nie więcej niż 2\% papierów wartościowych spółki publicznej ani działalności uprzednio ujawnionej i zatwierdzonej w Załączniku nr 3, ani udziału w Przedsięwzięciu Produkcyjnym zatwierdzonego zgodnie z §~\ref{art:przedsiewziecia} ust.~7.

  
\begin{MemoBox}{3.4.1}
\MemoPytCd{3.4.1}
\begin{Brzmienie}
§~\ref{art:konkurencja} ust.~2: „Zakaz nie obejmuje pasywnego posiadania nie więcej niż 2\% papierów wartościowych spółki publicznej ani działalności uprzednio ujawnionej i zatwierdzonej w Załączniku nr~3 albo w późniejszej uchwale Walnego Zgromadzenia, o której mowa w §~\ref{art:final} ust.~6, ani udziału w Przedsięwzięciu Produkcyjnym zatwierdzonego zgodnie z §~\ref{art:przedsiewziecia} ust.~7.”
\end{Brzmienie}
\end{MemoBox}

\item Po ustaniu funkcjonalnego zaangażowania zakaz konkurencji może obowiązywać nie dłużej niż 12 miesięcy, wyłącznie na podstawie odrębnej umowy określającej zakres i --- jeżeli prawo tego wymaga --- odpowiednie świadczenie Spółki.

  \item Na potrzeby niniejszego paragrafu „Osobą Bliską” jest współmałżonek dyrektora oraz jego krewny lub powinowaty do drugiego stopnia włącznie. Stopień pokrewieństwa oraz linię i stopień powinowactwa ustala się zgodnie z art. 61$^{7}$ § 2 oraz art. 61$^{8}$ § 2 Kodeksu rodzinnego i opiekuńczego. Osoba pozostająca z dyrektorem w stosunku pokrewieństwa lub powinowactwa dalszym niż drugi stopień nie jest Osobą Bliską wyłącznie z tego tytułu. Powyższe nie wyłącza zastosowania art. 300$^{55}$ § 1 KSH do osoby, z którą dyrektor jest powiązany osobiście, niezależnie od stopnia albo braku pokrewieństwa lub powinowactwa. „Podmiotem Powiązanym z Dyrektorem” jest podmiot, w którym dyrektor, jego Osoba Bliska albo podmiot przez nich kontrolowany:
  \begin{enumerate}[label=\alph*)]
    \item posiada bezpośrednio albo pośrednio więcej niż 2\% udziałów, akcji, głosów albo praw do udziału w zysku;
    \item pełni funkcję członka organu, prokurenta, wspólnika uprawnionego do prowadzenia spraw albo osoby faktycznie kierującej działalnością;
    \item posiada prawo powołania członka organu; albo
    \item odnosi szczególną korzyść ekonomiczną z zawarcia transakcji ze Spółką.
  \end{enumerate}
  \item Niezależnie od progu 2\% konflikt interesów występuje również wtedy, gdy dyrektor albo osoba z nim powiązana posiada istotny interes osobisty lub ekonomiczny, który może racjonalnie wpływać na bezstronność decyzji dotyczącej Spółki.
  \item Dyrektor pozostający w konflikcie interesów:
  \begin{enumerate}[label=\alph*)]
    \item niezwłocznie ujawnia Radzie Dyrektorów charakter i zakres konfliktu;
    \item nie uczestniczy w przygotowaniu, negocjowaniu, ocenie ani zatwierdzaniu transakcji;
    \item nie otrzymuje dokumentów dotyczących transakcji poza zakresem koniecznym do wykonania obowiązków prawnych;
    \item nie reprezentuje Spółki ani nie podpisuje dokumentów dotyczących tej transakcji.
  \end{enumerate}

  \item Dyrektor pozostający w konflikcie nie uczestniczy w rozstrzyganiu sprawy i jego głosu nie uwzględnia się przy ustalaniu wyniku głosowania. Wyłączenie dyrektora z udziału w rozstrzyganiu sprawy nie zmniejsza pełnego składu Rady Dyrektorów stanowiącego podstawę obliczenia quorum wymaganego przez art. 300$^{58}$ § 2 KSH. Uchwała dotycząca transakcji konfliktowej wymaga ponadto udziału co najmniej dwóch dyrektorów niepozostających w konflikcie, jeżeli Rada Dyrektorów liczy co najmniej trzy osoby.

  \item Jeżeli wskutek wyłączenia dyrektorów pozostających w konflikcie Rada Dyrektorów nie może podjąć wymaganej uchwały, stosuje się następujące zasady:
  \begin{enumerate}[label=\alph*)]
    \item jeżeli wymóg uchwały Rady Dyrektorów wynika wyłącznie z niniejszej umowy, zawarcie transakcji wymaga uprzedniej uchwały Walnego Zgromadzenia zatwierdzającej transakcję;
    \item jeżeli uchwały Rady Dyrektorów wymaga bezwzględnie obowiązujący przepis prawa, w szczególności art.~300$^{75}$ §~2 KSH, uchwała Walnego Zgromadzenia nie zastępuje uchwały Rady; transakcja nie może zostać zawarta do czasu przywrócenia zdolności Rady do podjęcia uchwały, w szczególności przez powołanie przez Walne Zgromadzenie dodatkowego dyrektora niepozostającego w konflikcie interesów; Walne Zgromadzenie może dodatkowo wyrazić właścicielską zgodę na transakcję.
  \end{enumerate}
  Uchwała Walnego Zgromadzenia nie zastępuje szczególnego sposobu reprezentacji Spółki wynikającego z art. 300$^{79}$ KSH.

  
\begin{MemoBox}{2.3.1}
\MemoPyt{2.3.1}{Czy uchwała Walnego Zgromadzenia może zastąpić uchwałę Rady w sprawach, dla których ustawa sama wymaga uchwały Rady (w szczególności art.~300$^{75}$ §~2 KSH), gdy wyłączenia dyrektorów uniemożliwiają jej podjęcie? Jeżeli nie --- jaki tryb rezerwowy jest dopuszczalny?}
\Ocena{zgodne (w brzmieniu 0.9.4)}
\Klasa{opcjonalne}
\Podstawa{art.~300$^{75}$ §~2 KSH \Wiar{Z}; art.~300$^{55}$ §~1 KSH (obowiązek ujawnienia i wstrzymania się od udziału w rozstrzyganiu) \Wiar{Z}; art.~300$^{58}$ §~2 KSH (quorum) \Wiar{Z}; art.~300$^{73}$ §~3 KSH (powoływanie dyrektorów uchwałą akcjonariuszy) \Wiar{Z}; art.~17 §~1--3 KSH \Wiar{Z}; art.~300$^{80}$ §~1--3 KSH (uchwała akcjonariuszy poza WZ; tajne głosowanie z art.~300$^{99}$ §~2 nie ma tu zastosowania) \Wiar{Z}.}
\Uzasadnienie{Nie --- i projekt już to przyjmuje. Art.~300$^{75}$ §~2 przydziela decyzje strategiczne, plany i strukturę Radzie jako organowi; uchwała akcjonariuszy może być zgodą właścicielską, ale nie decyzją organu prowadzącego sprawy. Ust.~8 lit.~b) wyraża to poprawnie, a lit.~a) słusznie ogranicza substytucję do wymogów czysto umownych (tam art.~17 §~3 i tak pozbawia brak uchwały skutku zewnętrznego); wątpliwość audytu do wersji 0.8 jest rozwiązana. Dopuszczalny tryb rezerwowy: (1)~powołanie przez WZ dodatkowego dyrektora niepozostającego w konflikcie --- możliwe w kilka dni, bo uchwała może zapaść poza zgromadzeniem, elektronicznie (art.~300$^{80}$ §~1--2; §~\ref{art:wz} ust.~4--5), a wymóg tajnego głosowania (art.~300$^{99}$ §~2) jej nie dotyczy (art.~300$^{80}$ §~3); (2)~stała „rezerwa” --- co najmniej jeden niezależny dyrektor niewykonawczy niebędący akcjonariuszem, jedyne realne zabezpieczenie przy pięciu Założycielach-akcjonariuszach, bo transakcja z jednym z nich albo z Przedsięwzięciem Produkcyjnym, w którym uczestniczą, może wyłączyć większość Rady; (3)~zawężenie zakresu: większość transakcji z podmiotami powiązanymi nie jest „decyzją o strategicznym znaczeniu”, więc wymóg uchwały Rady wynika z §~\ref{art:rada-zastrzezone} ust.~1 lit.~e), a nie z ustawy, i stosuje się lit.~a) --- Rada powinna w uchwale kwalifikować źródło wymogu. Niedopuszczalne są: udział dyrektora w konflikcie „za zgodą WZ” (art.~300$^{55}$ §~1 jest bezwzględny), obniżenie quorum ad hoc poniżej ustawowego i przekazanie decyzji prokurentowi. Kontrahent może powołać się na art.~17 §~3 tylko przy wymogu umownym; przy ustawowym ryzyko dotyczy odpowiedzialności dyrektorów i zaskarżenia „zastępczej” uchwały WZ.}
\Rekomendacja{Bez zmian normatywnych. Opcjonalnie: w §~\ref{art:konkurencja} ust.~8 --- termin, w którym Rada zwołuje WZ w celu powołania dodatkowego dyrektora (np.~14 dni); w §~\ref{art:rada} ust.~2 --- co najmniej jeden dyrektor niewykonawczy niebędący akcjonariuszem ani Osobą Bliską akcjonariusza (zob. ramkę przy §~\ref{art:rada} ust.~2).}
\Kierunek{Brzmienie uzupełnienia §~\ref{art:konkurencja} ust.~8 lit.~b): „[\ldots] Rada Dyrektorów zwołuje w takim przypadku Walne Zgromadzenie z porządkiem obrad obejmującym powołanie dodatkowego dyrektora w terminie 14 dni od stwierdzenia braku zdolności do podjęcia uchwały.”}
\end{MemoBox}

\item Zawarcie transakcji z Podmiotem Powiązanym z Dyrektorem wymaga uprzedniej uchwały Rady Dyrektorów podjętej wyłącznie przez dyrektorów niepozostających w konflikcie interesów albo, w przypadku określonym w ust. 8, uprzedniej uchwały Walnego Zgromadzenia.

  \item Uchwała zatwierdzająca transakcję, o której mowa w ust. 8 lub 9, powinna potwierdzać, że:
  \begin{enumerate}[label=\alph*)]
    \item transakcja leży w interesie Spółki;
    \item jej warunki odpowiadają warunkom rynkowym;
    \item dokonano porównania dostępnych ofert albo sporządzono uzasadnienie wyboru danego kontrahenta;
    \item zakres świadczeń, cena, odpowiedzialność stron i prawa własności intelektualnej zostały należycie zabezpieczone.
  \end{enumerate}

  \item Umowę i pozostałe oświadczenia dotyczące transakcji podpisują, zgodnie z zasadami reprezentacji Spółki i z zastrzeżeniem §~\ref{art:reprezentacja}, wyłącznie osoby niepozostające w konflikcie interesów.

  \item \BoilerplateStrike{Postanowienia niniejszego paragrafu nie naruszają art. 300$^{79}$ KSH.} Jeżeli stroną umowy ze Spółką albo stroną sporu ze Spółką jest dyrektor, Spółka jest reprezentowana przez pełnomocnika powołanego uchwałą akcjonariuszy, a w przypadku dyrektora wykonawczego — również przez dyrektora niewykonawczego, jeżeli spełnione są warunki określone w §~\ref{art:reprezentacja} i art. 300$^{79}$ § 3 KSH. Jeżeli wszystkie akcje Spółki przysługują jednemu akcjonariuszowi będącemu zarazem dyrektorem, stosuje się art. 300$^{79}$ § 4 KSH. Ogólna zasada działania osób niepozostających w konflikcie interesów nie zastępuje tego szczególnego trybu.
  
\begin{MemoBox}{2.3.3}
\MemoPyt{2.3.3}{Czy przy umowie i sporze z dyrektorem projekt prawidłowo odsyła do art.~300$^{79}$ KSH (§~29 ust.~12), a przekreślone zdanie można bezpiecznie usunąć?}
\Ocena{zgodne}
\Klasa{bez zmian}
\Podstawa{art.~300$^{79}$ §~1, 3 i 4 KSH \Wiar{Z}; art.~300$^{80}$ KSH \Wiar{Z}; art.~15 §~1 KSH (nie wymienia dyrektora P.S.A. --- pkt~2.2.4) \Wiar{Z}.}
\Uzasadnienie{Odesłanie jest prawidłowe: §~\ref{art:reprezentacja} ust.~3--5 i niniejszy ust.~12 odtwarzają trzy ustawowe warianty --- pełnomocnik powołany uchwałą akcjonariuszy (art.~300$^{79}$ §~1), dyrektor niewykonawczy na podstawie uchwały podjętej wyłącznie przez dyrektorów niewykonawczych (tylko wobec dyrektora wykonawczego, §~3; opcja wymagająca upoważnienia w umowie, które zawiera §~\ref{art:reprezentacja} ust.~4) oraz tryb jedynego akcjonariusza-dyrektora z aktem notarialnym (§~4), a ostatnie zdanie ust.~12 trafnie oddziela reżim art.~300$^{79}$ (stroną jest sam dyrektor; sankcją jest nieważność) od ust.~11 (stroną jest Podmiot Powiązany z Dyrektorem; zwykła reprezentacja przez osoby niekonfliktowe). Przekreślone zdania są deklaratywne i można je bezpiecznie usunąć, pod warunkiem pozostawienia zdania drugiego §~\ref{art:reprezentacja} ust.~3 i całego ust.~12. Dla pożyczki, poręczenia i podobnej umowy z dyrektorem potrzebna jest ponadto zgoda WZ, której art.~15 §~1 KSH dla dyrektora P.S.A. literalnie nie przewiduje --- należy ją wprowadzić umową (pkt~2.2.4 i 2.3.4).}
\Rekomendacja{Usunąć przekreślone zdania w §~\ref{art:reprezentacja} ust.~3 i §~\ref{art:konkurencja} ust.~12; resztę pozostawić.}
\end{MemoBox}

\item Jeżeli mimo podjęcia wymaganej uchwały nie można zapewnić prawidłowej reprezentacji Spółki, transakcja nie może zostać zawarta do czasu zapewnienia prawidłowej reprezentacji, w szczególności przez powołanie dodatkowego dyrektora przez Walne Zgromadzenie.
\end{enumerate}
\end{ContractClause}





\begin{ContractClause}{Transakcje z Akcjonariuszami i podmiotami powiązanymi}
\label{art:transakcje-akcjonariusze}
\textbf{[WYMAGA OPINII PRAWNIKA]}
\begin{enumerate}
  \item Akcjonariusz może świadczyć Spółce usługi, wykonywać na jej rzecz prace, dokonywać dostaw, udzielać licencji lub zawierać ze Spółką inne umowy. Sam status akcjonariusza nie daje pierwszeństwa w uzyskaniu zamówienia ani nie stanowi przeszkody do zawarcia umowy ze Spółką.
  \item Transakcja z Akcjonariuszem lub podmiotem z nim powiązanym powinna leżeć w interesie Spółki, a jej warunki powinny być godziwe dla Spółki, z uwzględnieniem wartości świadczenia wzajemnego, warunków rynkowych, kwalifikacji wykonawcy, dostępności, czasu realizacji, własności intelektualnej oraz innych okoliczności istotnych dla Spółki.
  
\begin{MemoBox}{2.3.4}
\MemoPyt{2.3.4}{§~30: jak pojęcie „godziwych warunków” ma się do ustawowych ograniczeń świadczeń na rzecz akcjonariuszy (art.~300$^{15}$ i n.\ KSH); czy brak obowiązku przetargu i progi ustalane w regulaminie albo uchwale są bezpieczne; które większe transakcje powinny wymagać zgody Walnego Zgromadzenia?}
\Ocena{zgodne warunkowo}
\Klasa{zalecane}
\Podstawa{art.~300$^{21}$ KSH (wartość świadczeń spółki na rzecz akcjonariuszy z innego tytułu niż prawa z akcji, a także na rzecz spółek lub spółdzielni z nimi powiązanych, „nie może przekraczać wartości godziwej świadczenia wzajemnego otrzymanego przez spółkę”) \Wiar{Z}; art.~300$^{22}$ §~1--4 KSH (zwrot wypłaty „dokonanej wbrew przepisom prawa lub postanowieniom umowy spółki”; solidarna odpowiedzialność członków organów, „chyba że nie ponoszą winy”; zakaz zwolnienia; przedawnienie 3 lata, chyba że odbiorca wiedział o bezprawności) \Wiar{Z}; art.~15 §~1 KSH (nie wymienia dyrektora P.S.A. --- pkt~2.2.4) \Wiar{Z}; art.~394 §~1 KSH (S.A.: nabycie mienia od założyciela lub akcjonariusza za cenę ponad 1/10 wpłaconego kapitału zakładowego w ciągu dwóch lat od rejestracji --- uchwała WZ 2/3; brak odpowiednika w P.S.A., wzorzec) \Wiar{Z}; art.~23m--23zf ustawy o PIT (ceny transferowe; podmioty powiązane od 25\%, art.~23m ust.~1 pkt~4) \Wiar{Z}; art.~61 rozporządzenia finansowego (UE) 2018/1046 i zasada konkurencyjności w programach dotacyjnych \Wiar{W}.}
\Uzasadnienie{Wypłaty na rzecz akcjonariuszy są dopuszczalne tylko w granicach testów z art.~300$^{15}$ §~2 i 5, a ponadto art.~300$^{21}$ stanowi wprost, że wartość świadczeń Spółki na rzecz akcjonariuszy (i spółek z nimi powiązanych) z innego tytułu niż prawa z akcji nie może przekraczać wartości godziwej świadczenia wzajemnego. Nadwyżka jest wypłatą „wbrew przepisom prawa”, którą odbiorca zwraca; członkowie organów odpowiadają za zwrot solidarnie, chyba że nie ponoszą winy, bez możliwości zwolnienia (art.~300$^{22}$ §~1--3); równolegle organ podatkowy doszacuje dochód według przepisów o cenach transferowych (Założyciel z 50\% jest podmiotem powiązanym od progu 25\%; Założyciele z 10\% --- tylko przy innych powiązaniach). Ust.~2 nie odsyła do art.~300$^{21}$, a „warunki rynkowe” są w nim tylko jedną z okoliczności; kryteria jakościowe są poprawne jako test interesu Spółki, ale nie przesuwają granicy ustawowej. Test trzeba rozciągnąć na podmioty powiązane spoza przepisu (Osoby Bliskie, spółki osobowe, fundacje) --- nadwyżka będzie wypłatą sprzeczną co najmniej z umową. Przetarg: prawo spółek go nie wymaga, a ust.~3 poprawnie zastrzega obowiązki z przepisów szczególnych i warunków finansowania (EDF, Horyzont Europa, NCBR, EUDIS). Progi w regulaminie Rady są bezpieczne, ale ich naruszenie ma skutek tylko wewnętrzny (art.~17 §~3); uchwała WZ ustanawiająca dodatkowe wymogi zgody może być kwestionowana jako instruowanie Rady --- najważniejsze progi powinny znaleźć się w umowie (zob. ramkę przy ust.~5). Wobec zarzutu ukrytych wypłat naruszających art.~20 realną obroną jest dokumentacja z ust.~5.}
\Rekomendacja{Doprecyzować ust.~2 (odesłanie do art.~300$^{21}$ KSH: wartość świadczenia Spółki nie wyższa niż wartość godziwa świadczenia wzajemnego, także wobec podmiotów powiązanych spoza literalnego zakresu przepisu; nadwyżka do zwrotu), przenieść progi z regulaminu do umowy jako nowy ust.~5 z listą transakcji wymagających uchwały WZ, objąć nim umowy kredytu, pożyczki i poręczenia z dyrektorem niezależnie od art.~15 KSH; usunąć oznaczenie „[WYMAGA OPINII PRAWNIKA]” po potwierdzeniu przez radcę prawnego. Decyzja Założycieli: wysokość progu rocznego (500~000~zł) i objęcie mienia nabywanego od Założycieli w okresie 2 lat.}
\Kierunek{Zmiana zabezpiecza przed zwrotem świadczeń i doszacowaniem podatkowym, ogranicza swobodę Rady w dużych transakcjach z Założycielami, nie wpływa na rejestr. Brzmienie ust.~2 (nowy ust.~5 --- zob. ramkę przy ust.~5):}
\begin{Brzmienie}
2. Transakcja z Akcjonariuszem lub podmiotem z nim powiązanym powinna leżeć w interesie Spółki, a jej warunki powinny być godziwe dla Spółki i nie gorsze niż warunki, jakie Spółka uzyskałaby od niezależnego kontrahenta w porównywalnych okolicznościach, z uwzględnieniem wartości świadczenia wzajemnego, kwalifikacji wykonawcy, dostępności, czasu realizacji, własności intelektualnej oraz innych okoliczności istotnych dla Spółki. Wartość świadczenia Spółki nie może przekraczać wartości godziwej świadczenia wzajemnego otrzymanego przez Spółkę (art.~300$^{21}$ KSH), także wtedy, gdy stroną transakcji jest podmiot powiązany z Akcjonariuszem niewymieniony w tym przepisie. Świadczenie Spółki przekraczające wartość godziwą świadczenia wzajemnego stanowi wypłatę na rzecz akcjonariusza sprzeczną z prawem i niniejszą umową i podlega zwrotowi na zasadach art.~300$^{22}$ KSH.
\end{Brzmienie}
\end{MemoBox}

\item Sam fakt, że kontrahentem jest Akcjonariusz lub podmiot z nim powiązany, nie powoduje obowiązku przeprowadzenia przetargu, konkursu ofert ani uzyskania określonej liczby ofert, chyba że obowiązek taki wynika z bezwzględnie obowiązujących przepisów prawa, warunków finansowania, dotacji, grantu, zamówienia albo wiążącej Spółkę umowy.
  \item Istotny konflikt interesów związany z transakcją podlega ujawnieniu. Jeżeli Akcjonariusz jest jednocześnie dyrektorem albo transakcja dotyczy podmiotu powiązanego z dyrektorem, stosuje się zasady konfliktu interesów i reprezentacji określone w §~\ref{art:konkurencja} oraz §~\ref{art:reprezentacja}; zainteresowany dyrektor nie może samodzielnie rozstrzygać o zawarciu własnej transakcji ze Spółką.
  \item Transakcje przekraczające progi wartości lub spełniające inne kryteria określone w niniejszej umowie, regulaminie Rady Dyrektorów albo uchwale Walnego Zgromadzenia wymagają dodatkowej zgody właściwego organu. Rada Dyrektorów może określić zasady dokumentowania godziwości i rynkowości takich transakcji, w tym przypadki, w których wymagane jest porównanie ofert, wycena, benchmark lub pisemne uzasadnienie wyboru kontrahenta.

\begin{MemoBox}{2.3.4}
\MemoPytCd{2.3.4}
\Kierunek{Nowe brzmienie ust.~5 (streszczenie): niezależnie od wymogów §~\ref{art:rada-zastrzezone}, uchwały WZ podjętej większością 75\% Głosów Uprawnionych w Sprawie, przy wyłączeniu od głosowania Akcjonariusza będącego stroną albo beneficjentem transakcji (z zastrzeżeniami z pkt~2.2.3), wymaga:
\begin{itemize}
\item udzielenie Akcjonariuszowi lub podmiotowi z nim powiązanemu pożyczki, kredytu, poręczenia, gwarancji albo innego zabezpieczenia, także gdy Akcjonariusz jest dyrektorem albo prokurentem, niezależnie od art.~15 KSH (dla dyrektora nie jest on pewną podstawą --- pkt~2.2.4);
\item przeniesienie na Akcjonariusza lub podmiot z nim powiązany Kluczowej Własności Intelektualnej albo udzielenie licencji na nią (§~\ref{art:walne-zastrzezone} ust.~1 lit.~g już częściowo);
\item transakcja, po której łączna wartość świadczeń Spółki na rzecz jednego Akcjonariusza i podmiotów z nim powiązanych w roku obrotowym przekroczy 500~000~zł (spójnie z §~\ref{art:walne-zastrzezone} ust.~1 lit.~k);
\item nabycie od Założyciela mienia za cenę przekraczającą 10\% kapitału akcyjnego w ciągu dwóch lat od wpisu Spółki do rejestru (wzorzec art.~394 §~1 KSH: cena ponad 1/10 wpłaconego kapitału zakładowego, dwa lata, większość 2/3; w P.S.A. przepis nie obowiązuje, więc ochrona przed obejściem wyceny wkładów z §~\ref{art:wklady} musi wynikać z umowy).
\end{itemize}
Rada Dyrektorów określa w regulaminie zasady dokumentowania godziwości i rynkowości transakcji (porównanie ofert, wycena, benchmark, pisemne uzasadnienie wyboru kontrahenta).}
\end{MemoBox}

\end{enumerate}

\begin{MemoBox}{4.3.1}
\MemoPytCd{4.3.1}
Uwaga konieczna K16 (ramka główna --- s.~\pageref{memo:4.3.1}): akt notarialny nie może zawierać adnotacji roboczych.
\begin{Brzmienie}
§~\ref{art:transakcje-akcjonariusze} --- usunąć oznaczenie „[WYMAGA OPINII PRAWNIKA]” z nagłówka; treść ustępów bez zmian.
\end{Brzmienie}
\end{MemoBox}

\end{ContractClause}

\begin{ContractClause}{Finansowanie Spółki}
\label{art:finansowanie}
\begin{enumerate}
  \item Spółka może finansować działalność z wkładów, emisji akcji, instrumentów zamiennych, pożyczek, kredytów, leasingu, faktoringu, gwarancji bankowych i ubezpieczeniowych, finansowania eksportu, zaliczek i przedpłat od zamawiających, kredytu kupieckiego, dotacji, grantów, zamówień, programów badawczo-rozwojowych, finansowania projektowego na poziomie Przedsięwzięć Produkcyjnych oraz innych prawnie dopuszczalnych źródeł.
  \item Żaden akcjonariusz nie jest zobowiązany do dodatkowego finansowania Spółki, chyba że wyrazi na to odrębną pisemną zgodę.
  \item Finansowanie nie może prowadzić do powstania bezpośredniej lub pośredniej kontroli, praw do akcji ani praw do Kluczowej Własności Intelektualnej po stronie podmiotu niespełniającego wymogów Dopuszczalnego Nabywcy.
  \item Umowy finansowania zawierające prawo konwersji, warrant, opcję, prawo nominacji dyrektora, prawo weta lub zabezpieczenie na Kluczowej Własności Intelektualnej wymagają uchwały Walnego Zgromadzenia jako Sprawa Zastrzeżona. Zabezpieczenia na aktywach produkcyjnych, zapasach i wierzytelnościach Spółki nie wymagają uchwały Walnego Zgromadzenia, chyba że ich wartość przekracza próg określony w §~\ref{art:walne-zastrzezone} ust.~1 lit.~k.
  
\begin{MemoBox}{3.3.4}
\MemoPyt{3.3.4}{Czy ograniczenia z §~\ref{art:finansowanie} ust.~3--4 (instrumenty zamienne, zabezpieczenia na IP) nie utrudnią finansowania dłużnego i grantowego i czy są egzekwowalne wobec finansujących?}
\Ocena{zgodne warunkowo}
\Klasa{zalecane}
\Podstawa{art.~17 §~1--3 KSH \Wiar{Z}; art.~300$^{77}$ §~2 KSH (nieograniczalność prawa reprezentacji dyrektora wobec osób trzecich), art.~300$^{53}$ KSH \Wiar{Z}; art.~300$^{81}$ pkt~2 i~4, art.~300$^{114}$ §~3 KSH (uchwały akcjonariuszy wymagane ustawą) \Wiar{Z}; art.~9 ust.~3--4, art.~2 pkt~6 i~24, art.~5, art.~20 ust.~3--4 i~9, art.~23 ust.~2--4 rozporządzenia (UE) 2021/697 \Wiar{Z}; art.~2 pkt~1 rozporządzenia (UE) 2026/1386 \Wiar{Z} (tekst EN); §~\ref{art:dopuszczalny-nabywca} ust.~1 i~6.}
\Uzasadnienie{Ust.~3 odsyła do „wymogów Dopuszczalnego Nabywcy”, czyli tylko do testu prawnej dopuszczalności (§~\ref{art:dopuszczalny-nabywca} ust.~1), nie do Kryterium Bezpieczeństwa EU/NATO: bank czy leasingodawca spoza UE/NATO nie jest nim objęty. Jeżeli intencją było Kryterium, trzeba to napisać --- wtedy np. venture debt z USA z warrantem wymagałby zgody Walnego Zgromadzenia z §~\ref{art:dopuszczalny-nabywca} ust.~6. Znaczenie dla EDF: odbiorcy nie mogą podlegać kontroli niestowarzyszonego państwa trzeciego lub podmiotu z takiego państwa (art.~9 ust.~3 rozporządzenia 2021/697 \Wiar{Z}; państwa stowarzyszone: tylko EFTA/EOG, art.~5), chyba że państwo siedziby zatwierdzi gwarancje z art.~9 ust.~4; test Dopuszczalnego Nabywcy tego nie wychwytuje. Pod „prawo weta” z ust.~4 można podciągnąć typowe kowenanty kredytowe, co uczyniłoby każdy kredyt Sprawą Zastrzeżoną; kredyt z zastawem rejestrowym na przedsiębiorstwie, obejmującym z reguły także IP, wymaga uchwały 75\%. Standardowe warunki grantowe (dla EDF z mocy rozporządzenia, art.~20 ust.~3--4, art.~23 ust.~2 i~4 \Wiar{Z}; dla pozostałych programów \Wiar{W}) są wobec partnerów spełniających test prawny dozwolone. \emph{Egzekwowalność:} czynność bez zgody wymaganej wyłącznie umową spółki jest ważna (art.~17 §~3 KSH, art.~300$^{77}$ §~2 \Wiar{Z}), sankcją jest odpowiedzialność dyrektorów; nieważna jest czynność bez uchwały wymaganej ustawą (art.~17 §~1 \Wiar{Z}): zastaw na przedsiębiorstwie, warranty i obligacje zamienne, umowy o prawo objęcia akcji (art.~300$^{81}$ pkt~2 i~4, art.~300$^{114}$ §~3 \Wiar{Z}). Ust.~3--4 są więc normami kompetencyjnymi, nie zabezpieczeniem wobec finansującego. Czy dług z warrantem albo konwersją jest od 17.01.2028~r. „inwestycją zagraniczną” (art.~2 pkt~1 rozporządzenia 2026/1386 \Wiar{Z} (tekst EN)), rozstrzygnie praktyka \Wiar{?}.}
\Rekomendacja{Rozstrzygnąć (decyzja Założycieli), czy ust.~3 ma odsyłać do Kryterium; zawęzić „prawo weta” do prawa sprzeciwu wobec uchwał organów Spółki; dodać wyłączenie dla standardowych postanowień umów o dofinansowanie i kowenantów kredytowych.}
\Kierunek{Zmiana zmniejsza liczbę Spraw Zastrzeżonych przy bieżącym finansowaniu, nie osłabia ochrony IP (zabezpieczenie na Kluczowej Własności Intelektualnej pozostaje przy Walnym Zgromadzeniu) i nie wpływa na skuteczność wobec finansujących, która wynika z art.~17 §~1 i~3 KSH.}
\begin{Brzmienie}
§~\ref{art:finansowanie} ust.~4: „Umowy finansowania zawierające prawo konwersji na akcje, warrant, opcję albo inne prawo do akcji, prawo wskazania dyrektora, prawo sprzeciwu wobec uchwał organów Spółki albo zabezpieczenie na Kluczowej Własności Intelektualnej wymagają uchwały Walnego Zgromadzenia jako Sprawa Zastrzeżona. Nie wymagają takiej uchwały zwyczajowe kowenanty finansowe i informacyjne, zastrzeżenie zgody finansującego na zmianę kontroli nad Spółką ani standardowe postanowienia umów o dofinansowanie ze środków publicznych dotyczące praw dostępu, eksploatacji, publikacji i ograniczeń rozporządzania wynikami, jeżeli nie przewidują przeniesienia ani wyłącznej licencji Kluczowej Własności Intelektualnej na rzecz finansującego. Zabezpieczenia na aktywach produkcyjnych, zapasach i wierzytelnościach Spółki nie wymagają uchwały Walnego Zgromadzenia, chyba że ich wartość przekracza próg określony w §~\ref{art:walne-zastrzezone} ust.~1 lit.~k.”
\end{Brzmienie}
\end{MemoBox}

\item Rada Dyrektorów może przyjąć roczny Plan Finansowania określający limity kredytów, leasingu, faktoringu, gwarancji bankowych i ubezpieczeniowych, w tym gwarancji należytego wykonania i wadialnych, oraz zabezpieczeń na aktywach Spółki innych niż Kluczowa Własność Intelektualna, a także warunki ich zaciągania i ustanawiania. Czynności mieszczące się w Planie Finansowania nie wymagają odrębnej uchwały Rady Dyrektorów na podstawie §~\ref{art:rada-zastrzezone} ust.~1 lit.~d; czynności poza Planem Finansowania podlegają progom określonym w §~\ref{art:rada-zastrzezone} i §~\ref{art:walne-zastrzezone}. Plan Finansowania nie może obejmować zabezpieczenia na Kluczowej Własności Intelektualnej ani instrumentów, o których mowa w ust.~4.
\end{enumerate}

\end{ContractClause}

\begin{ContractClause}{Kwalifikowana Runda Finansowania}
\label{art:kwalifikowana-runda}
\begin{enumerate}
  \item „Kwalifikowana Runda Finansowania” oznacza emisję nowych akcji Spółki skierowaną do jednego lub kilku inwestorów, z których każdy jest Dopuszczalnym Nabywcą i spełnia Kryterium Bezpieczeństwa EU/NATO \BoilerplateStrike{(w tym jako Kwalifikowany Inwestor Finansowy w rozumieniu §~\ref{art:dopuszczalny-nabywca} ust.~14)} albo uzyskał zgodę Walnego Zgromadzenia zgodnie z §~\ref{art:dopuszczalny-nabywca}, jeżeli łączna cena emisyjna nowych akcji wynosi co najmniej \field{2\,000\,000 zł} albo emisja następuje przy wycenie Spółki przed emisją wynoszącej co najmniej \field{kwota}.
  
\begin{MemoBox}{1.4.3}
\MemoPytCd{1.4.3}
\Kierunek{Jeżeli §~\ref{art:dopuszczalny-nabywca} ust.~14 zostanie przywrócony w brzmieniu uszczelnionym, przywrócić przekreślone odesłanie do Kwalifikowanego Inwestora Finansowego w niniejszym ustępie (zob.~też 1.4.1); jeżeli Założyciele zdecydują o usunięciu wyjątku, żadna inna zmiana nie jest potrzebna.}
\end{MemoBox}

\item Rozpoczęcie Kwalifikowanej Rundy Finansowania wymaga uchwały kierunkowej Walnego Zgromadzenia podjętej większością 75\% wszystkich głosów, określającej co najmniej: widełki wyceny lub minimalną cenę emisyjną, maksymalną liczbę nowych akcji, dopuszczalny krąg inwestorów oraz ramowe warunki z ust.~4.
  
\begin{MemoBox}{3.3.3}
\MemoPytCd{3.3.3}
\Kierunek{Uchwała kierunkowa może zawierać zgodę z §~\ref{art:dopuszczalny-nabywca} ust.~6 oraz zatwierdzenie transakcji wtórnych z inwestorami Rundy; uzasadnienie w ramce głównej przy §~\ref{art:okres-ograniczenia-zbywania} ust.~2 (s.~\pageref{memo:3.3.3}).}
\begin{Brzmienie}
§~\ref{art:kwalifikowana-runda} ust.~2 zdanie drugie: „Uchwała kierunkowa może zawierać zgodę, o której mowa w §~\ref{art:dopuszczalny-nabywca} ust.~6, wobec oznaczonych inwestorów albo ich kategorii, a także zatwierdzenie zbycia akcji przez akcjonariuszy na rzecz inwestorów Rundy w oznaczonym zakresie.”
\end{Brzmienie}
\end{MemoBox}

\item Po podjęciu uchwały kierunkowej każdy akcjonariusz zobowiązuje się współdziałać w dobrej wierze w celu przeprowadzenia Kwalifikowanej Rundy Finansowania na warunkach zgodnych z tą uchwałą, w szczególności: uczestniczyć w Walnych Zgromadzeniach zwołanych w związku z Rundą, głosować za uchwałami niezbędnymi do jej przeprowadzenia oraz zawrzeć zwyczajowe dokumenty transakcyjne. Zobowiązanie to nie obejmuje głosowania sprzecznego z bezwzględnie obowiązującymi przepisami prawa, uchwałą kierunkową albo interesem Spółki i nie zobowiązuje żadnego akcjonariusza do objęcia nowych akcji ani wniesienia dodatkowych środków.
  
\begin{MemoBox}{3.3.1}
\MemoPyt{3.3.1}{Czy zobowiązanie do głosowania i współdziałania (§~\ref{art:kwalifikowana-runda} ust.~3) jest skuteczne wobec akcjonariuszy i egzekwowalne, czy pozostaje jedynie zobowiązaniem odszkodowawczym?}
\Ocena{zgodne warunkowo}
\Klasa{zalecane}
\Podstawa{art.~300$^{1}$ §~3 KSH \Wiar{Z}; art.~353$^{1}$, art.~64, art.~101 §~1, art.~471, art.~483 §~1 KC \Wiar{Z}; art.~1047 §~1 KPC \Wiar{Z}; art.~300$^{101}$ KSH w zw. z art.~422 §~1 KSH \Wiar{Z}; SN, wyrok z 9 grudnia 2022~r., II CSKP 593/22 \Wiar{Z}; SN, wyrok z 23 czerwca 2020~r., V CSK 522/18 \Wiar{Z}.}
\Uzasadnienie{Umowy o sposób wykonywania prawa głosu są dopuszczalne w granicach art.~353$^{1}$ KC; zobowiązanie z ust.~3 mieści się w nich, bo jest ograniczone do uchwał zgodnych z uchwałą kierunkową, wyłącza głosowanie sprzeczne z prawem i interesem Spółki i nie nakłada obowiązku wniesienia środków \Wiar{W}. Zapisane w umowie spółki wiąże także każdego późniejszego nabywcę akcji (art.~300$^{1}$ §~3 \Wiar{Z}), a jego naruszenie jest naruszeniem umowy spółki (§~\ref{art:breach}). Granice skuteczności: (1)~głos oddany wbrew zobowiązaniu jest ważny; podstawą uchylenia uchwały jest jej sprzeczność z umową spółki bądź dobrymi obyczajami połączona z godzeniem w interes spółki albo pokrzywdzeniem akcjonariusza (art.~422 §~1 w zw. z art.~300$^{101}$ KSH \Wiar{Z}), a nie samo naruszenie zobowiązania \Wiar{W}; uchwała emisyjna, której zabrakło, nie powstaje; (2)~orzeczenie zastępujące oświadczenie woli (art.~64 KC, art.~1047 §~1 KPC \Wiar{Z}) wymaga dostatecznie oznaczonego obowiązku (II CSKP 593/22) --- czyli konkretnej uchwały kierunkowej i projektu uchwały emisyjnej --- a czas procesu czyni ten środek iluzorycznym w rundzie; (3)~klauzula „interesu Spółki” daje odmawiającemu argument obronny. Realne środki to odszkodowanie (art.~471 KC --- trudny dowód szkody), kara umowna (art.~483 KC) i pełnomocnictwo do głosowania ze zrzeczeniem się odwołania (art.~101 §~1 KC \Wiar{Z}; V CSK 522/18); dwa ostatnie należą do umowy akcjonariuszy, która jest już przewidziana. Luka: ust.~3 nie obejmuje późniejszych uchwał wykonawczych z dokumentów Rundy (emisja antyrozwodnieniowa, pro-rata, uzupełnienie puli P) --- inwestor zażąda rozszerzenia.}
\Rekomendacja{Utrzymać ust.~3 i rozszerzyć go na uchwały wykonawcze przewidziane w dokumentach Rundy oraz na uchwałę o pozbawieniu prawa poboru; w umowie akcjonariuszy dodać karę umowną i nieodwołalne pełnomocnictwo do głosowania w sprawach objętych uchwałą kierunkową, z ograniczeniem do jej treści.}
\Kierunek{Proponowane brzmienie ust.~3 wylicza wprost uchwały, za którymi akcjonariusz głosuje: o emisji, o zmianie umowy Spółki w zakresie określonym w ust.~4 i o pozbawieniu prawa poboru w części niezbędnej do objęcia akcji przez inwestorów, a także uchwały wykonawcze przewidziane w dokumentach transakcyjnych Rundy (w szczególności emisja w wykonaniu ochrony antyrozwodnieniowej albo prawa uczestnictwa w przyszłych emisjach); zachowuje dotychczasowe wyłączenia (prawo, uchwała kierunkowa, interes Spółki, brak obowiązku objęcia akcji i wniesienia środków) i dodaje, że naruszenie zobowiązania stanowi istotne naruszenie umowy. Rozszerzenie wzmacnia pozycję inwestora i Spółki wobec akcjonariusza blokującego; dla mniejszości granicą pozostaje uchwała kierunkowa (75\%), interes Spółki i prawo; sąd rejestrowy i rejestr akcjonariuszy nie są zmianą dotknięci.}
\end{MemoBox}

\item Następujące zwyczajowe warunki inwestorskie mogą zostać przyznane w uchwale emisyjnej i dokumentach transakcyjnych Kwalifikowanej Rundy Finansowania bez konieczności odrębnej zmiany pozostałych postanowień niniejszej umowy, w granicach prawa: (a) uprzywilejowanie likwidacyjne bez prawa udziału (non-participating) do wysokości jednokrotności ceny emisyjnej; (b) ochrona antyrozwodnieniowa metodą szerokiej średniej ważonej (broad-based weighted average); (c) prawo nominacji jednego dyrektora albo obserwatora przy Radzie Dyrektorów, z zachowaniem §~\ref{art:rada} ust.~3; (d) prawa informacyjne; (e) prawo uczestnictwa w przyszłych emisjach proporcjonalnie do udziału (pro-rata). Warunki dalej idące, w szczególności uprzywilejowanie likwidacyjne z prawem udziału, ochrona antyrozwodnieniowa typu full ratchet albo prawo weta wykraczające poza katalog Spraw Zastrzeżonych, wymagają odrębnej uchwały Walnego Zgromadzenia podjętej większością 75\% wszystkich głosów.
  
\begin{MemoBox}{3.3.2}
\MemoPyt{3.3.2}{Czy warunki inwestorskie z §~\ref{art:kwalifikowana-runda} ust.~4 mogą zostać przyznane bez zmiany umowy spółki (uprzywilejowanie akcji, prawa osobiste inwestora), czy ich część wymaga zmiany umowy w formie aktu notarialnego?}
\Ocena{ryzyko (sformułowanie „bez konieczności odrębnej zmiany pozostałych postanowień niniejszej umowy” jest mylące co do formy)}
\Klasa{konieczne (K15)}
\Podstawa{art.~300$^{5}$ §~1 pkt~3, art.~300$^{25}$ §~1--2, art.~300$^{26}$ KSH (akcje uprzywilejowane, akcje założycielskie) \Wiar{Z}; art.~300$^{28}$ §~1--2 KSH (uprawnienia indywidualne akcjonariusza --- odpowiednik art.~354 KSH) \Wiar{Z}; art.~300$^{73}$ §~3 KSH (dyrektorów powołują akcjonariusze uchwałą, chyba że umowa stanowi inaczej) \Wiar{Z}; art.~300$^{98}$ §~4--5, art.~300$^{100}$ §~2, art.~300$^{102}$--300$^{104}$, art.~300$^{106}$ §~1--2, art.~300$^{110}$ §~5 KSH \Wiar{Z}; art.~2 pkt~1, 5 i~7, art.~4 ust.~1, 2, 9 i~15, art.~5 ust.~2 lit.~a, art.~11 ust.~3 i~5, art.~30 ust.~3, art.~31 oraz załącznik~II pkt~14 rozporządzenia (UE) 2026/1386 \Wiar{Z} (tekst EN).}
\Uzasadnienie{(a)~Uprzywilejowanie likwidacyjne musi być „określone w umowie spółki” (art.~300$^{25}$ §~1--2 \Wiar{Z}); przewidująca je uchwała emisyjna (art.~300$^{104}$ §~1 pkt~2 \Wiar{Z}) jest zarazem zmianą umowy w protokole notarialnym (art.~300$^{103}$), co potwierdza zakaz emisji akcji uprzywilejowanych przez Radę z upoważnienia (art.~300$^{110}$ §~5 \Wiar{Z}); preferencja przy exicie jest kontraktowa. (b)~Antyrozwodnienie --- kontraktowe, z zastrzeżeniem 3.3.1. (c)~Prawo wskazania dyrektora skuteczne wobec Spółki to uprawnienie indywidualne wymagające umowy spółki (art.~300$^{28}$ §~1, art.~300$^{73}$ §~3 \Wiar{Z}); inaczej --- tylko zobowiązanie akcjonariuszy do głosowania. (d)~Prawa informacyjne --- kontraktowe. (e)~Pro-rata wynika z prawa poboru, które umowa spółki może ukształtować inaczej (art.~300$^{106}$ §~1; 3.3.3); „super pro-rata” --- kontraktowa. Zgody serii wymagają umowy spółki (art.~300$^{98}$ §~4 \Wiar{Z} daje tylko zaczątek). Runda i tak jest zmianą umowy w protokole notarialnym (brak upoważnienia do emisji, art.~300$^{100}$ §~2); obecny ust.~4 sugeruje preferencję z samej uchwały emisyjnej --- grozi to odmową wpisu albo sporem, czy powstała. Wariant: predefiniowana seria I z upoważnieniem z art.~300$^{103}$; praktyka rejestrowa wymaga potwierdzenia \Wiar{?}. Runda z inwestorem spoza UE (także z EOG lub NATO) zamykana po 17.01.2028~r. (art.~30 ust.~3) wymaga uprzedniego zezwolenia, jeżeli Spółka działa w zakresie załącznika~I rozporządzenia 2021/821 albo załącznika dyrektywy 2009/43/WE, a inwestycja daje efektywny udział w zarządzaniu lub kontroli (art.~2 pkt~1, art.~4 ust.~15, art.~31 rozporządzenia 2026/1386 \Wiar{Z} (tekst EN)); inwestor z prawem wskazania dyrektora najpewniej się w tym mieści \Wiar{W}. Potrzebny warunek zawieszający i harmonogram (ocena wstępna 45~dni, postępowanie pogłębione bez terminu, przy projekcie EDF także mechanizm współpracy) \Wiar{Z} (tekst EN) --- T1.}
\Rekomendacja{Przeredagować ust.~4 tak, aby rozróżniał warunki wprowadzane zmianą umowy Spółki (a, c, zgody serii) od kontraktowych (b, d, e) i utrzymywał zgodę właścicielską na poziomie uchwały kierunkowej. Rozważyć predefiniowaną serię I (z preferencją 1$\times$ bez udziału) jako wariant --- opcjonalny, do decyzji.}
\begin{Brzmienie}
§~\ref{art:kwalifikowana-runda} ust.~4: „Uchwała kierunkowa może dopuścić przyznanie inwestorom następujących zwyczajowych warunków, w granicach prawa: (a)~uprzywilejowanie likwidacyjne bez prawa udziału (non-participating) do wysokości jednokrotności ceny emisyjnej; (b)~ochrona antyrozwodnieniowa metodą szerokiej średniej ważonej (broad-based weighted average); (c)~prawo wskazania jednego dyrektora albo obserwatora przy Radzie Dyrektorów, z zachowaniem §~\ref{art:rada} ust.~3; (d)~prawa informacyjne; (e)~prawo uczestnictwa w przyszłych emisjach proporcjonalnie do udziału (pro-rata). Warunki, które zgodnie z prawem muszą wynikać z umowy Spółki --- w szczególności uprzywilejowanie akcji nowej serii (lit.~a; art.~300$^{25}$ KSH), uprawnienie indywidualne do wskazania dyrektora (lit.~c; art.~300$^{28}$ KSH) oraz wymóg zgody akcjonariuszy danej serii na zmianę ich praw --- wprowadza się uchwałą o zmianie niniejszej umowy podejmowaną łącznie z uchwałą o emisji w trybie §~\ref{art:walne-zastrzezone}, za którą akcjonariusze głosują zgodnie z ust.~3; pozostałe warunki mogą zostać przyznane w uchwale emisyjnej i dokumentach transakcyjnych. Przyznanie warunków w granicach dopuszczonych uchwałą kierunkową nie wymaga odrębnej uchwały Walnego Zgromadzenia poza uchwałami, o których mowa w zdaniu poprzedzającym. Warunki dalej idące, w szczególności uprzywilejowanie likwidacyjne z prawem udziału, ochrona antyrozwodnieniowa typu full ratchet albo prawo weta wykraczające poza katalog Spraw Zastrzeżonych, wymagają odrębnej uchwały Walnego Zgromadzenia podjętej większością 75\% wszystkich głosów.”
\end{Brzmienie}
\end{MemoBox}

\begin{MemoBox}{4.4.2}
\MemoPytCd{4.4.2}
\Kierunek{(VI)~Standardowy, zamknięty katalog spraw wymagających zgody serii inwestorskiej, który może zostać przyznany bez odrębnej uchwały 75\% (A.6; zależnie od decyzji Założycieli o~kręgu inwestorów). Katalog jest węższy niż §~\ref{art:walne-zastrzezone}; wszystko poza nim nadal wymaga osobnej uchwały. Próg 75\% dla samej rundy się nie zmienia, ale znika sygnał „weto to warunek dalej idący”, który fundusze czytają jako niechęć. Nowa lit.~f (po lit.~e, przed zdaniem „Warunki dalej idące”):}
\begin{Brzmienie}
(f)~wymóg zgody akcjonariuszy posiadających większość akcji serii wyemitowanej w Rundzie na: zmianę umowy Spółki naruszającą prawa tej serii, emisję akcji lub instrumentów równych lub uprzywilejowanych wobec tej serii, nabycie akcji własnych poza mechanizmem zwrotnego zbycia, wypłatę dywidendy przed upływem okresu z §~\ref{art:zysk} ust.~2, sprzedaż przedsiębiorstwa lub Kluczowej Własności Intelektualnej w przypadkach z §~\ref{art:walne-zastrzezone} ust.~1 lit.~f--g, zaciągnięcie zobowiązania ponad próg z §~\ref{art:walne-zastrzezone} ust.~1 lit.~k, transakcje z podmiotami powiązanymi ponad próg określony w §~\ref{art:transakcje-akcjonariusze}, zwiększenie puli, o której mowa w §~\ref{art:program-motywacyjny}, oraz emisję akcji serii F po Rundzie; uprawnienie to nie ogranicza większości wymaganych przez §~\ref{art:walne-zastrzezone} i wygasa z chwilą, gdy akcje tej serii przestaną stanowić co najmniej \textbf{[5]}\% wszystkich akcji.
\end{Brzmienie}
\end{MemoBox}

\item Emisja akcji w ramach Kwalifikowanej Rundy Finansowania pozostaje Sprawą Zastrzeżoną i następuje zgodnie z §~\ref{art:walne-zastrzezone}, z zachowaniem zasad prawa poboru i większości wymaganych przez prawo dla jego ograniczenia lub wyłączenia.
  \item Umowę przygotowują na Kwalifikowaną Rundę Finansowania w szczególności: §~\ref{art:akcje} i §~\ref{art:cap-table} (struktura akcji), §~\ref{art:program-motywacyjny} (pula ESOP), §~\ref{art:zwrotne-zbycie} (vesting założycielski), §~\ref{art:dopuszczalny-nabywca} i §~\ref{art:zbywanie} (weryfikacja inwestora), §~\ref{art:pierwszenstwo}, §~\ref{art:przylaczenie} i §~\ref{art:przymusowe-wspolzbycie} (obrót wtórny), §~\ref{art:wycena} (standard wyceny), §~\ref{art:finansowanie} ust.~4 (instrumenty zamienne) oraz §~\ref{art:walne-zastrzezone} (ład emisyjny).

\begin{MemoBox}{4.4.1}
\MemoPyt{4.4.1}{Które postanowienia (Kryterium, mechanizm przy odmowie zgody, zwrotne zbycie, drag-along 75\,\%, Sprawy Zastrzeżone 75\,\%, ograniczenia finansowania) w Państwa doświadczeniu budzą zastrzeżenia funduszy na etapie seed/serii A i co zwykle jest przedmiotem negocjacji?}
\Ocena{zgodne warunkowo --- żadne z ocenianych postanowień nie jest prawnie wadliwe; sześć z nich odbiega od standardu rynkowego w sposób, który wydłuży negocjacje albo zawęzi krąg funduszy}
\Klasa{zalecane}

\Podstawa{art.~300$^{39}$ §~1--6, art.~300$^{47}$, art.~300$^{25}$ §~1--2 i~art.~300$^{28}$ §~1--2 KSH \Wiar{Z}; art.~2 ust.~2 pkt~1 ustawy z~2018~r. o~przeciwdziałaniu praniu pieniędzy oraz finansowaniu terroryzmu (beneficjent rzeczywisty: m.in. „więcej niż 25\% ogólnej liczby udziałów lub akcji” albo głosów) \Wiar{Z}; art.~2 pkt~6 i~24, art.~5 i~art.~9 ust.~1--4 i~7 rozporządzenia (UE) 2021/697 (EDF) \Wiar{Z}; art.~2 pkt~1, 5 i~7, art.~15 ust.~1 lit.~b i~ust.~3, art.~19 ust.~2 lit.~e rozporządzenia (UE) 2026/1386 (kontrola inwestycji zagranicznych, stosowane od 17 stycznia 2028~r., art.~31) \Wiar{Z} (tekst EN); praktyka rynkowa: wzorce NVCA po aktualizacji z~2~października 2025~r., omówienia PFR Startup i~strony programów PFR Ventures \Wiar{Z}, wzorce BVCA i~wzory term sheet PFR Ventures (niepobrane) \Wiar{W}.}
\Uzasadnienie{Fundusze: (G)~generalistyczne; (D)~deep tech i~obronne (NIF, EDF, DIANA; art.~9 rozporządzenia 2021/697 \Wiar{Z}); (Z)~zagraniczne. Oceny praktyki \Wiar{W}.
\begin{enumerate}[label=\arabic*.]
\item \textbf{Kryterium (§~\ref{art:dopuszczalny-nabywca} ust.~6 lit.~b).} Przy wykładni przez ustawę AML \Wiar{Z} fundusze G i~D przeważnie je spełnią; spór budzi badanie LP (ust.~4, 10). G chcą oświadczenia zamiast look-through \Wiar{W} (tak wzorce NVCA z~2025~r. \Wiar{Z}); D akceptują test zarządzającego, siedziby i~kontroli \Wiar{W}, jak od 17 stycznia 2028~r. rozporządzenie 2026/1386 \Wiar{Z} (tekst EN); Z --- tylko przez zgodę z~ust.~6. Negocjuje się definicję testu, nie zasadę.
\item \textbf{Utrata Kryterium (ust.~15)}: zmiana po stronie LP może wymusić zbycie; fundusze zażądają zawężenia do zmiany kontroli nad zarządzającym.
\item \textbf{Odmowa zgody (§~\ref{art:zbywanie} ust.~3).} Wariant~A (3~miesiące zamiast ustawowego miesiąca, art.~300$^{39}$ §~2 KSH \Wiar{Z}; raty) łatwiejszy niż~B (bez wyjścia); fundusze zażądają zapłaty jednorazowej i~krótszego terminu.
\item \textbf{Zwrotne zbycie} --- atut; negocjuje się parametry (cena 80\% przy Odejściu Zawinionym, wybór nabywcy albo podział proporcjonalny).
\item \textbf{Drag 75\%}: inwestora do ok.~25\% można zmusić do sprzedaży; brak zgody serii, ceny minimalnej i~waterfall.
\item \textbf{Sprawy Zastrzeżone 75\%}: próg standardowy \Wiar{Z}, ale inwestor z~20\% nie ma weta; zażąda zgód serii (art.~300$^{25}$ KSH \Wiar{Z}).
\item \textbf{Brak przeniesień dozwolonych} (następca, SPV, wydanie LP).
\item \textbf{Dyrektor} musi spełniać Kryterium (§~\ref{art:rada} ust.~3).
\item \textbf{Finansowanie (§~\ref{art:finansowanie} ust.~3--4)}: ust.~3 to test sankcyjny; Kryterium wraca przy prawach do akcji (§~\ref{art:dopuszczalny-nabywca} ust.~13).
\end{enumerate}
Z~każdym funduszem negocjuje się pkt~1, 2, 5--7; z~Z i~częścią D --- pkt~3 i~8.}
\Rekomendacja{Rozstrzygnąć przed rundą, nie w~jej trakcie, cztery punkty sporne z~każdym funduszem: test Kryterium dla funduszy (§~\ref{art:dopuszczalny-nabywca} ust.~14 i~15), przeniesienia dozwolone (§~\ref{art:zbywanie}), ochrona serii inwestorskiej przy drag (§~\ref{art:przymusowe-wspolzbycie}) i~katalog spraw wymagających zgody serii (§~\ref{art:kwalifikowana-runda} ust.~4). Pozostałe punkty zostawić do term sheetu. Brzmienia --- w~odpowiedzi na pytanie 4.4.2 (s.~\pageref{memo:4.4.2}). Wybór, czy szukać inwestora w~grupie G/D (mniejsze zmiany), czy także Z (zmiany w~§~\ref{art:rada} ust.~3 i~wariant~A w~§~\ref{art:zbywanie} ust.~3), należy do Założycieli.}
\Kierunek{Zmiany z~pytania 4.4.2; tu bez odrębnych brzmień. Dla Spółki: krótsze due diligence i~szerszy krąg funduszy; dla Założycieli: brak utraty kontroli, bo każda zmiana zachowuje próg 75\% i~Kryterium co do kontroli nad Spółką; dla rejestru: każdy wyjątek od zgody Spółki Rada Dyrektorów dokumentuje uchwałą lub protokołem, aby podmiot prowadzący rejestr miał podstawę wpisu.}
\end{MemoBox}

\begin{MemoBox}{4.4.2}
\MemoPyt{4.4.2}{Jakie zmiany ułatwiłyby dostosowanie umowy do rundy bez naruszenia jej założeń (bezpieczeństwo właścicielskie, kontrola eksportu, mechanizm zwrotnego zbycia)?}
\Ocena{zgodne warunkowo --- założenia umowy da się utrzymać w całości; potrzebne są cztery zmiany w umowie spółki przed rundą i przeniesienie reszty do dokumentów rundy}
\Klasa{zalecane --- odstępstwo od reguły „zgodne warunkowo → zalecane” nie jest potrzebne; brzmienia podajemy, bo zmiany są proste i jednoznaczne}

\Podstawa{art.~300$^{25}$ §~1--2, art.~300$^{28}$ §~1--2, art.~300$^{33}$ §~1 pkt~4 i~10, art.~300$^{39}$ §~1--2 („chyba że umowa spółki stanowi inaczej”), art.~300$^{103}$, art.~300$^{104}$ §~1 pkt~2 i~art.~300$^{106}$ §~1--2 KSH \Wiar{Z}; art.~2 pkt~6 i~24, art.~5 i~art.~9 ust.~1--4 i~7 rozporządzenia (UE) 2021/697 \Wiar{Z}; art.~2 pkt~1, 5 i~7, art.~4 ust.~1, 2, 9 i~15 lit.~a--c, art.~5 ust.~1--2, art.~6, art.~11 ust.~3 i~5, art.~20 ust.~4 lit.~a, art.~30--31, motywy 14--15 i~20, zał.~I pkt~1--3 i~zał.~II pkt~14 rozporządzenia (UE) 2026/1386 \Wiar{Z} (tekst EN); rozporządzenie (UE) 2021/821 (do jego zał.~I odsyła art.~4 ust.~15 lit.~a rozporządzenia 2026/1386) \Wiar{Z} (tekst EN, konsolidacja na 26.05.2023~r.); art.~12a ust.~1 pkt~1, art.~12c ust.~1 pkt~1 i~5, ust.~4--5 oraz art.~12d ust.~3 pkt~7 i~ust.~4 ustawy o~kontroli niektórych inwestycji \Wiar{Z}; art.~2 ust.~2 pkt~1 ustawy AML \Wiar{Z}.}
\Uzasadnienie{Założenia --- kontrola nad Spółką z~UE/EOG/NATO (przy decyzji D6: UE/EOG), dostęp spoza UE/NATO do Kluczowej Własności Intelektualnej za zgodą Rady (§~\ref{art:security} ust.~6, §~\ref{art:export}), zwrotne zbycie jako obowiązek związany z~akcją (§~\ref{art:zwrotne-zbycie} ust.~14) --- decydują o~dostępie do EDF (siedziba i~„zarządcze struktury wykonawcze” w~Unii lub państwie stowarzyszonym, brak kontroli niestowarzyszonego państwa trzeciego, art.~9 ust.~1--4 rozporządzenia 2021/697 \Wiar{Z}), EUDIS, DIANA/NIF i~koncesji; zmiany ich nie dotykają.

\textbf{A. Przed rundą, w~umowie spółki:} A.1~uszczelniony §~\ref{art:dopuszczalny-nabywca} ust.~14 (I); A.2~zawężenie ust.~15 (II); A.3~przeniesienia dozwolone (III); A.4~ochrona serii przy drag (IV); A.5~dyrektor inwestora (V); A.6~lit.~f w~§~\ref{art:kwalifikowana-runda} ust.~4 (VI); A.7~opcjonalnie seria~I z~upoważnieniem do emisji (KSH nie wyklucza \Wiar{Z}; praktyka sądów rejestrowych otwarta \Wiar{?}).

\textbf{B. W~dokumentach rundy} (waterfall, prawa informacyjne, lock-up, exit) oraz warunek zawieszający z~kontroli inwestycji: do 16 stycznia 2028~r. --- inwestor spoza UE, EOG i~OECD, co najmniej 20\% głosów albo dominacja, przychód w~Polsce ponad 10~mln euro \Wiar{Z}; od 17 stycznia 2028~r. (inwestycje zamknięte wcześniej --- dotychczasowy reżim), gdy Spółka rozwija, wytwarza lub wprowadza do obrotu produkty z~zał.~I rozporządzenia 2021/821 albo dyrektywy 2009/43/WE --- każdy inwestor spoza Unii (także NATO, EOG, OECD) i~fundusz unijny kontrolowany spoza niej, przy „efektywnym udziale w~zarządzaniu lub kontroli” (tłum. robocze) \Wiar{Z} (tekst EN): lider z~dyrektorem i~zgodami serii --- tak, drobny inwestor pasywny --- zapewne nie \Wiar{W}.}
\Rekomendacja{Wprowadzić zmiany A.1--A.4 przed pierwszym term sheetem (A.1 i~A.3 warunkują w~praktyce rozmowę z~funduszami D i~Z); A.5 i~A.6 zależnie od decyzji Założycieli o~kręgu inwestorów; A.7 po stanowisku kancelarii. Resztę zostawić dokumentom rundy. Brzmienia są sugestią do weryfikacji przez uprawnionego prawnika, w~szczególności co do art.~300$^{39}$ §~2 KSH (wyłączenie zgody Spółki dla kategorii przeniesień) i~ujawnienia w~rejestrze akcjonariuszy. Nie zmieniać: Kryterium co do kontroli nad Spółką i~wymogu zgody 75\% dla nabywcy spoza Kryterium (§~\ref{art:dopuszczalny-nabywca} ust.~5--6), §~\ref{art:security} ust.~6, §~\ref{art:export}, obowiązku związanego z~akcją w~§~\ref{art:zwrotne-zbycie} ust.~14 ani progu 75\% dla emisji i~zmiany umowy --- fundusze D tego oczekują, funduszom G i~Z trzeba to uzasadnić dostępem do programów, a~nie ustępować.}
\Kierunek{Dla Spółki i~Założycieli: pełna kontrola właścicielska przy usunięciu czterech punktów spornych; dla inwestora: przewidywalność wejścia, obrotu w~grupie funduszu i~wyjścia; dla rejestru: dodatkowe kategorie zwolnień od zgody Spółki do udokumentowania protokołem Rady; dla sądu rejestrowego: brak nowych elementów poza uprzywilejowaniem serii przy rundzie; dla sporu między akcjonariuszami: obiektywne przesłanki (oświadczenie, próg 25\%, okres ochronny, cena minimalna) zamiast uznania Rady. Brzmienia (I)--(VI) --- w~ramkach przy §~\ref{art:dopuszczalny-nabywca} ust.~14 i~15, §~\ref{art:zbywanie} ust.~7, §~\ref{art:przymusowe-wspolzbycie} ust.~4, §~\ref{art:rada} ust.~3 i~§~\ref{art:kwalifikowana-runda} ust.~4.}
\end{MemoBox}

\end{enumerate}

\end{ContractClause}

\begin{ContractClause}{Spółki celowe i wspólne przedsięwzięcia produkcyjne}
\label{art:przedsiewziecia}
\begin{enumerate}
  \item „Przedsięwzięciem Produkcyjnym” jest spółka celowa, konsorcjum, wspólne przedsięwzięcie albo inna forma współpracy kapitałowej lub kontraktowej z podmiotami trzecimi, w szczególności partnerami przemysłowymi, integratorami i zamawiającymi, której celem jest produkcja przemysłowa, integracja, certyfikacja, serwis albo dystrybucja produktów wykorzystujących technologie Spółki. Spółka może tworzyć Przedsięwzięcia Produkcyjne, przystępować do nich, obejmować w nich rolę wiodącą, w tym kontrolę, albo uczestniczyć w nich jako wspólnik mniejszościowy. Udział w Przedsięwzięciu Produkcyjnym mieści się w celu określonym w §~\ref{art:cel} i nie stanowi istotnej zmiany profilu działalności w rozumieniu §~\ref{art:walne-zastrzezone}.
  \item Utworzenie Przedsięwzięcia Produkcyjnego, przystąpienie do niego, objęcie albo utrata kontroli nad nim oraz zwiększenie zaangażowania Spółki wymagają uchwały Rady Dyrektorów zgodnie z §~\ref{art:rada-zastrzezone} ust.~1 lit.~k, a w przypadkach określonych w §~\ref{art:walne-zastrzezone} ust.~1 lit.~n --- uchwały Walnego Zgromadzenia. Zaangażowanie obejmuje wkłady, nabycie udziałów lub akcji, pożyczki, poręczenia, gwarancje i inne zabezpieczenia udzielone na rzecz Przedsięwzięcia Produkcyjnego albo jego wierzycieli.
  
\begin{MemoBox}{2.4.4}
\MemoPyt{2.4.4}{Czy progi kwotowe i podział kompetencji Rada / Walne Zgromadzenie są wykonalne przy zmianach zaangażowania w czasie (np.\ dokapitalizowanie spółki celowej)?}
\Ocena{zgodne warunkowo}
\Klasa{zalecane}
\Podstawa{art.~300$^{75}$ §~2 KSH \Wiar{Z}; art.~17 §~1 i 3 KSH \Wiar{Z}; art.~300$^{81}$ pkt~2 KSH („zbycie i wydzierżawienie przedsiębiorstwa albo jego zorganizowanej części oraz ustanowienie na nich ograniczonego prawa rzeczowego” --- uchwała akcjonariuszy z mocy ustawy, bez klauzuli opt-out, gdy wkład do Przedsięwzięcia ma taki charakter) \Wiar{Z}; art.~300$^{98}$ §~2 pkt~2 KSH (zbycie przedsiębiorstwa albo ZCP --- 3/4 głosów) \Wiar{Z}; §~\ref{art:rada-zastrzezone} ust.~1 lit.~d), f), k), §~\ref{art:walne-zastrzezone} ust.~1 lit.~f), n), §~\ref{art:finansowanie} ust.~4--5.}
\Uzasadnienie{Podział jest wykonalny co do zasady: próg jest łączny i obejmuje wkłady, udziały, pożyczki i zabezpieczenia, a utrata kontroli i wniesienie Kluczowej Własności Intelektualnej są odrębnymi wyzwalaczami WZ (§~\ref{art:walne-zastrzezone} ust.~1 lit.~n). (1)~Sposób liczenia: umowa nie mówi, czy zabezpieczenia liczy się według sumy gwarancyjnej, pożyczki według kapitału, wkłady rzeczowe według wartości godziwej, czy spłata pożyczki albo wygaśnięcie gwarancji obniża zaangażowanie ani czy licencja z ust.~3 ma wartość zaliczaną do progu. (2)~Po przekroczeniu progu każde dalsze zwiększenie wymaga kolejnej uchwały 75\%, co przy wezwaniach kapitałowych i prawie poboru w spółce celowej jest nierealne; brak uchwały oznacza rozwodnienie i --- przy utracie kontroli --- zdarzenie z lit.~n, którego nikt nie „dokonał”. (3)~Zdarzenia pasywne: utrata kontroli może wynikać z niewzięcia udziału w emisji, konwersji długu partnera albo zmiany umowy spółki celowej; Rada powinna uprzedzać WZ o decyzjach, których zaniechanie prowadzi do utraty kontroli. (4)~Zbieg progów: zaangażowanie może być zarazem zobowiązaniem z §~\ref{art:walne-zastrzezone} ust.~1 lit.~k (500~000~zł poza budżetem) i zabezpieczeniem z §~\ref{art:finansowanie}; ust.~8 to sygnalizuje, ale bez reguły kolizyjnej. Jeżeli wkładem jest ZCP, uchwała WZ wynika z ustawy (art.~300$^{81}$ pkt~2, większość 3/4) niezależnie od kwoty, a jej brak oznacza nieważność (art.~17 §~1) --- to samo przy wydzierżawieniu ZCP Przedsięwzięciu i obciążeniu jej zastawem za jego długi. (5)~Uchwała WZ jest zgodą; decyzję o strategicznym znaczeniu podejmuje Rada (pkt~2.2.4). W sporze niejasna metoda liczenia to łatwy zarzut przekroczenia kompetencji.}
\Rekomendacja{Dodać w §~\ref{art:przedsiewziecia} ust.~2 zasady liczenia zaangażowania i instytucję Limitu Zaangażowania: WZ zatwierdza dla danego Przedsięwzięcia limit (kwotę i rodzaje instrumentów), w ramach którego Rada decyduje o kolejnych transzach bez odrębnej uchwały; obowiązek informacyjny Rady przed zaniechaniem prowadzącym do utraty kontroli; reguła kolizyjna z §~\ref{art:walne-zastrzezone} ust.~1 lit.~k i §~\ref{art:finansowanie}.}
\Kierunek{Zmiana zachowuje kontrolę właścicielską nad skalą ekspozycji, a Radzie daje operacyjną swobodę w ramach limitu; zmniejsza liczbę uchwał 75\%. Proponowany nowy ust.~2a (streszczenie):
\begin{itemize}
\item łączne zaangażowanie to suma wkładów pieniężnych i wartości godziwej wkładów niepieniężnych, ceny nabycia udziałów lub akcji, kapitału pożyczek pozostałego do spłaty oraz maksymalnej kwoty odpowiedzialności z poręczeń, gwarancji i innych zabezpieczeń pozostających w mocy; licencji z ust.~3 nie zalicza się;
\item WZ może uchwałą 75\% wszystkich głosów ustalić dla Przedsięwzięcia Limit Zaangażowania (kwota i dopuszczalne instrumenty); w jego granicach o zwiększeniu zaangażowania decyduje Rada zgodnie z §~\ref{art:rada-zastrzezone} ust.~1 lit.~k), także powyżej 500~000~zł;
\item Rada informuje akcjonariuszy z wyprzedzeniem umożliwiającym podjęcie uchwały o każdej decyzji albo zaniechaniu, które może prowadzić do utraty kontroli nad Przedsięwzięciem;
\item zaangażowanie zatwierdzone zgodnie z §~\ref{art:przedsiewziecia} nie wymaga odrębnej uchwały na podstawie §~\ref{art:walne-zastrzezone} ust.~1 lit.~k) ani §~\ref{art:finansowanie} ust.~5, chyba że obejmuje zabezpieczenie na Kluczowej Własności Intelektualnej albo zbycie, wydzierżawienie lub obciążenie przedsiębiorstwa lub ZCP (art.~300$^{81}$ pkt~2 KSH).
\end{itemize}
Decyzja Założycieli: czy wprowadzić Limit Zaangażowania zatwierdzany przez WZ dla poszczególnych Przedsięwzięć.}
\end{MemoBox}

\item Kluczowa Własność Intelektualna pozostaje własnością Spółki. Przedsięwzięciu Produkcyjnemu Spółka udziela licencji niewyłącznej, ograniczonej do zakresu Przedsięwzięcia (produkt, pole eksploatacji, terytorium, czas), z prawem wypowiedzenia w razie zmiany kontroli nad Przedsięwzięciem, naruszenia obowiązków odpowiadających §~\ref{art:security} lub §~\ref{art:export} albo objęcia partnera ograniczeniami, o których mowa w §~\ref{art:dopuszczalny-nabywca}. Prawa do ulepszeń i rozwinięć Kluczowej Własności Intelektualnej powstałych w Przedsięwzięciu Produkcyjnym przysługują Spółce, a jeżeli nie jest to prawnie możliwe --- Spółka uzyskuje do nich nieodpłatną, wyłączną i bezterminową licencję z prawem sublicencji. Przeniesienie Kluczowej Własności Intelektualnej albo udzielenie licencji wyłącznej, nieodwołalnej lub bezterminowej podlega §~\ref{art:walne-zastrzezone} ust.~1 lit.~g.
  
\begin{MemoBox}{2.4.1}
\MemoPyt{2.4.1}{Czy licencja niewyłączna ograniczona do zakresu przedsięwzięcia z prawami do ulepszeń dla Spółki jest zgodna z prawem konkurencji (porozumienia o transferze technologii) i nie rodzi obowiązków z kontroli eksportu przy udostępnieniu partnerom?}
\Ocena{zgodne warunkowo}
\Klasa{zalecane}
\Podstawa{art.~101 ust.~1 i 3 TFUE \Wiar{W}; rozporządzenie Komisji (UE) 2026/877 z 16.04.2026~r. w sprawie porozumień o transferze technologii (art.~3: progi 20\% łącznie dla konkurentów i 30\% dla każdej ze stron wśród niekonkurentów; art.~5 ust.~1 lit.~a: wyłączeniem nie jest objęte „bezpośrednie lub pośrednie zobowiązanie licencjobiorcy do udzielenia licencji wyłącznej albo przeniesienia praw, w całości lub w części, na licencjodawcę [\ldots] w odniesieniu do własnych ulepszeń albo własnych nowych zastosowań licencjonowanej technologii” (tłum. robocze); art.~9: rozporządzenia nie stosuje się do licencji w umowach B+R i specjalizacyjnych objętych rozporządzeniami 2023/1066 i 2023/1067; art.~10: okres przejściowy 1.05.2026--30.04.2027~r. tylko dla umów obowiązujących 30.04.2026~r. i spełniających warunki rozporządzenia 316/2014; art.~11: wejście w życie 1.05.2026~r., wygaśnięcie 30.04.2038~r.) \Wiar{Z} (tekst EN); rozporządzenia (UE) 2023/1066 (B+R) i 2023/1067 (specjalizacja) oraz wytyczne horyzontalne 2023 \Wiar{W}; rozporządzenie (UE) 2021/821: art.~2 pkt~2 lit.~d (eksport jako elektroniczne przekazanie oprogramowania lub technologii poza obszar celny Unii) i art.~4 ust.~1 (catch-all) \Wiar{Z} (tekst EN, wersja skonsolidowana 2023), art.~3 ust.~1, art.~8 ust.~1 i 3 lit.~a, art.~11 ust.~1 i załącznik~IV \Wiar{Z} (tekst EN), pozycje załącznika~IV nieweryfikowane, oraz zezwolenia generalne EU001 (bez Turcji) i EU002--EU006 (Turcja tylko dla wąskich kategorii) \Wiar{Z} (tekst EN, wersja skonsolidowana 2023); ustawa z 29.11.2000~r. o obrocie z zagranicą towarami, technologiami i usługami o znaczeniu strategicznym (art.~2, art.~3 pkt~5b i 10), znowelizowana ustawą z 13.03.2026~r. (Dz.U. 2026 poz.~471) \Wiar{Z}; ustawa z 13.06.2019~r. o wytwarzaniu i obrocie wyrobami i technologią o przeznaczeniu wojskowym lub policyjnym (art.~7 ust.~1 pkt~2: koncesja na obrót technologią; art.~3 ust.~1 pkt~12) \Wiar{Z}.}
\Uzasadnienie{Prawo konkurencji. Licencja niewyłączna w zakresie Przedsięwzięcia korzysta z wyłączenia grupowego w granicach progów z art.~3 \Wiar{Z} (tekst EN). Problemem jest grant-back z ust.~3: obowiązek przeniesienia albo wyłącznego licencjonowania ulepszeń na licencjodawcę jest „ograniczeniem wyłączonym” (art.~5 ust.~1 lit.~a) --- nie unieważnia całej umowy, lecz wypada spod wyłączenia grupowego i wymaga oceny indywidualnej (art.~101 ust.~3) \Wiar{Z} (tekst EN). Dotyczy to też własności ulepszeń niewydzielnych, choć jej ocena indywidualna jest z reguły korzystna \Wiar{W}; licencja niewyłączna na ulepszenia pozostaje w wyłączeniu \Wiar{Z} (tekst EN); czy prawo pierwszeństwa jest „pośrednim” zobowiązaniem do przeniesienia, rozporządzenie nie rozstrzyga \Wiar{?}. Wzór umowy partnerskiej nie powinien zawierać bezwzględnego zakazu kwestionowania praw (art.~5 ust.~1 lit.~b). Umowa spółki wyznacza mandat negocjacyjny Rady, obecnie nakazujący klauzulę potencjalnie niedozwoloną. Kontrola eksportu. Elektroniczne udostępnienie technologii poza obszar celny UE to eksport (art.~2 pkt~2 lit.~d rozporządzenia 2021/821); zezwolenia wymagają pozycje z załącznika~I (pozostałe --- po poinformowaniu przez organ), pomoc techniczną reguluje art.~8, a pozycje z załącznika~IV wymagają zezwolenia także przy transferze w UE \Wiar{Z} (tekst EN). Technologia wojskowa: obrót wymaga koncesji, obrót z zagranicą --- zezwolenia ministra właściwego do spraw gospodarki; czy licencja dla spółki celowej to „obrót”, wymaga potwierdzenia \Wiar{?}. Odesłania do §~\ref{art:export} i §~\ref{art:security} są poprawne (EU001 dla UK, USA i Norwegii poza pozycjami z sekcji~I załącznika~II; Turcja --- z reguły zezwolenie indywidualne albo globalne). W sporze grant-back może być podniesiony jako nieważny (art.~58 KC w zw. z art.~101 ust.~2 TFUE) --- warto klauzulę salwatoryjną.}
\Rekomendacja{Zmienić ust.~3 tak, aby cel (Spółka kontroluje ulepszenia Kluczowej Własności Intelektualnej) był realizowany w wariantach o możliwie niskim ryzyku: własność ulepszeń niewydzielnych (także ona wypada spod wyłączenia grupowego --- art.~5 ust.~1 lit.~a rozporządzenia 2026/877 --- i opiera się na ocenie indywidualnej, z reguły korzystnej); dla ulepszeń wydzielnych --- licencja niewyłączna z prawem sublicencji plus prawo pierwszeństwa nabycia; wyłączność tylko w zakresie dozwolonym przez prawo konkurencji, z odpłatnością, gdy jest wymagana; obowiązek Rady dokumentowania oceny konkurencyjnej przy każdej licencji; utrzymać odesłania eksportowe.}
\Kierunek{Zmiana obniża ryzyko nieważności klauzuli ulepszeniowej w umowach z partnerami i zarzutu naruszenia art.~101, przy zachowaniu ochrony IP; dla inwestora --- standard rynkowy. Brzmienie ust.~3 zd.~3:}
\begin{Brzmienie}
Prawa do ulepszeń i rozwinięć Kluczowej Własności Intelektualnej powstałych w Przedsięwzięciu Produkcyjnym, których nie można wykorzystać bez naruszenia Kluczowej Własności Intelektualnej, przysługują Spółce. Do pozostałych ulepszeń Spółka uzyskuje co najmniej nieodpłatną, niewyłączną i bezterminową licencję z prawem sublicencji oraz prawo pierwszeństwa ich nabycia, a licencję wyłączną --- w zakresie i na warunkach dopuszczalnych przez prawo konkurencji, w szczególności przez właściwe rozporządzenie w sprawie porozumień o transferze technologii. Rada Dyrektorów dokumentuje ocenę zgodności każdej umowy licencyjnej z prawem konkurencji przed jej zawarciem.
\end{Brzmienie}
\end{MemoBox}

\item Każdy partner Przedsięwzięcia Produkcyjnego podlega odpowiednio weryfikacji przewidzianej dla Dopuszczalnego Nabywcy w §~\ref{art:dopuszczalny-nabywca}. Partner uzyskujący dostęp do Kluczowej Własności Intelektualnej albo kontrolę nad Przedsięwzięciem Produkcyjnym musi spełniać Kryterium Bezpieczeństwa EU/NATO, chyba że Walne Zgromadzenie wyrazi zgodę większością 75\% wszystkich głosów. Udostępnienie technologii Przedsięwzięciu Produkcyjnemu podlega §~\ref{art:export}.
  
\begin{MemoBox}{2.4.2}
\MemoPyt{2.4.2}{Czy uzależnienie udziału partnera od Kryterium jest dopuszczalne, w tym wobec partnerów z państw NATO spoza UE?}
\Ocena{zgodne warunkowo}
\Klasa{zalecane}
\Podstawa{art.~353$^{1}$ KC (swoboda doboru kontrahenta) \Wiar{Z}; art.~5, 9 i 20 ust.~4 rozporządzenia (UE) 2021/697 (EDF): państwa stowarzyszone to członkowie EFTA należący do EOG (art.~5); odbiorcy i podwykonawcy „nie podlegają kontroli niestowarzyszonego państwa trzeciego lub podmiotu z niestowarzyszonego państwa trzeciego” (art.~9 ust.~3), z odstępstwem przy gwarancjach zatwierdzonych przez państwo siedziby (ust.~4); zmianę w toku działania zgłasza się Komisji (ust.~7); przeniesienie własności wyników albo licencja wyłączna na rzecz podmiotu z niestowarzyszonego państwa trzeciego --- uprzednie powiadomienie Komisji, a przy sprzeczności z interesami bezpieczeństwa i obronności Unii zwrot wsparcia (art.~20 ust.~4) \Wiar{Z}, przy czym które państwa EFTA/EOG faktycznie uczestniczą w Funduszu, nie wynika z tekstu rozporządzenia \Wiar{W}; rozporządzenie (UE) 2026/1386, art.~2 pkt~1, 2, 5 i~7, art.~4 ust.~9 i~15--17, art.~30 ust.~1 i~3, art.~31 (uprzednie zezwolenie przed zamknięciem transakcji; greenfield poza zakresem obowiązkowym; stosowanie od 17.01.2028~r. i z tym dniem uchylenie rozporządzenia 2019/452; nowych przepisów nie stosuje się do inwestycji zakończonych albo kontrolowanych w tym dniu) \Wiar{Z} (tekst EN); rozporządzenie (UE) 2021/821 i zezwolenie generalne EU001 (załącznik~II sekcja~A; część~2: Australia, Kanada, Islandia, Japonia, Nowa Zelandia, Norwegia, Szwajcaria z Liechtensteinem, Zjednoczone Królestwo, USA; bez Turcji, którą obejmują tylko wąskie zezwolenia EU002--EU006) \Wiar{Z} (tekst EN, wersja skonsolidowana 2023); art.~54 ust.~1--2 ustawy o ochronie informacji niejawnych (świadectwo bezpieczeństwa przemysłowego ABW albo SKW) \Wiar{Z}, dostęp podmiotów zagranicznych na podstawie umów dwustronnych \Wiar{W}; art.~4 ust.~1 pkt~7 ustawy o kontroli niektórych inwestycji \Wiar{Z}.}
\Uzasadnienie{Kryterium to wewnętrzny warunek doboru kontrahentów, nie ograniczenie obrotu akcjami --- dopuszczalne, jeśli Spółka nie narusza sankcji, reguł równego traktowania ani warunków programów. Zgoda WZ 75\% to tylko surowszy tryb. Partnerzy z NATO spoza UE (USA, UK, Turcja, Kanada) spełniają Kryterium. (1)~Kontrola partnera z USA/UK/Turcji nad Przedsięwzięciem pozbawia je kwalifikowalności w działaniach EDF (także EUDIS finansowanych z EDF \Wiar{W}), chyba że państwo siedziby zatwierdzi gwarancje z art.~9 ust.~4; państwa te nie są stowarzyszone \Wiar{Z}. Spółka niekontrolowana z takiego państwa kwalifikowalności nie traci, ale przeniesienie na partnera wyników wspieranych z Funduszu albo licencja wyłączna wymaga uprzedniego powiadomienia Komisji (art.~20 ust.~4) \Wiar{Z}; analogiczny wymóg w AGILE jest niepewny \Wiar{?}. (2)~Kryterium jest zbyt wąskie wobec państw zaufanych spoza UE/EOG/NATO (Szwajcaria, Japonia, Korea, Australia, Nowa Zelandia, Ukraina): każdy dostęp do Kluczowej Własności Intelektualnej wymaga uchwały 75\%, co przy 50/20/10/10/10 daje każdemu 26\% prawo blokady. (3)~Partner spoza UE uruchamia kontrolę eksportu (Turcja bez EU001), przy EUCI --- odrębne porozumienia; ITAR/EAR u partnera z USA mogą „skazić” produkt. (4)~Od 17.01.2028~r. nabycie przez partnera spoza UE (bez wyjątku dla EOG i NATO) udziału dającego kontrolę lub wpływ na zarządzanie istniejącym Przedsięwzięciem z zakresu załącznika~I do 2021/821 albo wykazu uzbrojenia wymaga uprzedniego zezwolenia (rozporządzenie 2026/1386); transakcje zakończone albo objęte kontrolą w tym dniu --- według dotychczasowych przepisów, a nowa spółka celowa (greenfield) jest poza zakresem obowiązkowym, chyba że prawo krajowe pójdzie dalej \Wiar{Z} (tekst EN); zob.~T1.}
\Rekomendacja{Uzupełnić ust.~4 o warstwę programową: gdy Przedsięwzięcie korzysta z finansowania UE na obronność albo ma być wykorzystane w takim projekcie, Rada zapewnia zgodność z warunkami programu (w szczególności art.~9 EDF), a kontrola partnera z państwa niestowarzyszonego nad Przedsięwzięciem wymaga uchwały WZ. Równolegle rozszerzyć kompetencję Rady (z głosem dyrektora ds. bezpieczeństwa) na partnerów z państw objętych unijnym zezwoleniem generalnym EU001 albo umową o bezpieczeństwie informacji z UE przy dostępie ograniczonym do zakresu Przedsięwzięcia; kontrola --- nadal WZ. Decyzja Założycieli: lista państw zaufanych i czy Ukraina ma być objęta trybem Rady.}
\Kierunek{Zmiana usuwa kolizję Kryterium z EDF i odblokowuje partnerstwa z państwami zaufanymi bez uchwały 75\%; zwiększa odpowiedzialność Rady i dyrektora ds. bezpieczeństwa. Proponowane brzmienie ust.~4 zd.~2--3 (streszczenie):
\begin{itemize}
\item partner uzyskujący dostęp do Kluczowej Własności Intelektualnej albo kontrolę nad Przedsięwzięciem musi spełniać Kryterium Bezpieczeństwa EU/NATO, chyba że WZ wyrazi zgodę większością 75\% wszystkich głosów;
\item Rada, uchwałą podjętą z zachowaniem §~\ref{art:rada-zastrzezone} ust.~2, może jednak dopuścić dostęp ograniczony do zakresu Przedsięwzięcia dla partnera z siedzibą i rzeczywistym ośrodkiem zarządzania w państwie z części~2 sekcji~A załącznika~II do rozporządzenia 2021/821 (EU001) albo związanym umową o bezpieczeństwie informacji z UE lub NATO, niekontrolowanego z innego państwa;
\item gdy Przedsięwzięcie korzysta z finansowania UE w dziedzinie obronności albo ma być wykorzystane w projekcie tak finansowanym, Rada zapewnia spełnienie warunków programu, w szczególności art.~9 rozporządzenia 2021/697, a uzyskanie kontroli nad Przedsięwzięciem przez podmiot z państwa niestowarzyszonego z programem wymaga uchwały WZ większością 75\% wszystkich głosów;
\item udostępnienie technologii Przedsięwzięciu podlega §~\ref{art:export}.
\end{itemize}}
\end{MemoBox}

\item Umowa Przedsięwzięcia Produkcyjnego zapewnia Spółce co najmniej: prawa informacyjne i kontrolne; obowiązki partnerów w zakresie bezpieczeństwa informacji, kontroli eksportu i sankcji odpowiadające §~\ref{art:security} i §~\ref{art:export}; zakaz przeniesienia udziału w Przedsięwzięciu na Niedopuszczalnego Nabywcę; prawo pierwszeństwa nabycia udziału partnera oraz prawo wyjścia w razie zmiany kontroli nad partnerem; a jeżeli Spółka nie sprawuje kontroli nad Przedsięwzięciem --- prawo sprzeciwu w sprawach dotyczących Kluczowej Własności Intelektualnej, bezpieczeństwa i zmiany przedmiotu Przedsięwzięcia.
  \item Umowy między Spółką a Przedsięwzięciem Produkcyjnym zawiera się na warunkach rynkowych. Jeżeli partnerem Przedsięwzięcia jest Akcjonariusz, dyrektor, Osoba Bliska albo podmiot z nimi powiązany, stosuje się §~\ref{art:konkurencja} ust.~4--13 oraz §~\ref{art:transakcje-akcjonariusze}.
  \item Udział Akcjonariusza Funkcjonalnego albo dyrektora, osobiście albo przez podmiot z nim powiązany, w Przedsięwzięciu Produkcyjnym, w którym uczestniczy Spółka, ujawniony Radzie Dyrektorów i zatwierdzony przez nią z zachowaniem ust.~6, nie stanowi działalności konkurencyjnej ani uczestnictwa w podmiocie konkurencyjnym w rozumieniu §~\ref{art:konkurencja} ust.~1; udział podlega ujawnieniu w Załączniku nr 3.
  
\begin{MemoBox}{2.4.3}
\MemoPyt{2.4.3}{Czy wyłączenie udziału założycieli w przedsięwzięciach spod zakazu konkurencji (§~33 ust.~7) nie osłabia §~29 i jak zabezpieczyć Spółkę przed wyciekiem wartości?}
\Ocena{zgodne warunkowo}
\Klasa{zalecane}
\Podstawa{art.~300$^{54}$ KSH (staranność i lojalność) i art.~300$^{55}$ §~1 KSH (konflikt interesów) \Wiar{Z}; art.~300$^{55}$ §~3 KSH (dyrektor „nie może bez zgody spółki zajmować się interesami konkurencyjnymi ani uczestniczyć w spółce konkurencyjnej”, w tym przy co najmniej 10\% głosów lub udziałów w konkurencyjnej spółce kapitałowej, „chyba że umowa spółki stanowi inaczej”; „Jeżeli umowa spółki nie stanowi inaczej, zgody udziela organ uprawniony do powoływania członka organu”) \Wiar{Z}; art.~300$^{1}$ §~3 KSH (obowiązki akcjonariusza z umowy) \Wiar{Z}; art.~300$^{33}$ §~1 pkt~11 KSH (rejestr ujawnia „postanowienia umowy spółki o związanych z akcją obowiązkach wobec spółki”) \Wiar{Z}; art.~11 ust.~1--2 ustawy o zwalczaniu nieuczciwej konkurencji (tajemnica przedsiębiorstwa) \Wiar{Z}; art.~101 ust.~1 TFUE (klauzule niekonkurowania w JV jako ograniczenia akcesoryjne) \Wiar{W}.}
\Uzasadnienie{Wyjątek jest uzasadniony celowościowo: Założyciel może być potrzebny w spółce celowej jako wspólnik-operator, a bez wyjątku każdy taki udział naruszałby §~\ref{art:konkurencja} ust.~1. Warunki (ujawnienie Radzie, zatwierdzenie z zachowaniem ust.~6, wpis do Załącznika nr~3) są jednak zbyt słabe w czterech miejscach. (1)~Zatwierdza Rada bezwzględną większością dyrektorów niekonfliktowych; przy Radzie złożonej z Założycieli możliwe jest wzajemne zatwierdzanie udziałów („ja tobie, ty mnie”). Ustawowo, jeżeli umowa nie stanowi inaczej, zgody na działalność konkurencyjną dyrektora udziela organ uprawniony do jego powołania, czyli WZ, a zakaz obejmuje udział od 10\% (art.~300$^{55}$ §~3); przeniesienie zgody na Radę jest dopuszczalne, ale odwraca ustawowy kierunek kontroli --- to argument, by powyżej progu 10\% decydowało WZ. (2)~Brak granic wielkości udziału, wynagrodzenia od Przedsięwzięcia i roli Założyciela, a to kanały wycieku wartości: marża Przedsięwzięcia trafia do Założyciela, Założyciel-menedżer negocjuje z samym sobą warunki licencji i dostaw, a po zakończeniu licencji Przedsięwzięcie staje się konkurentem. (3)~Zatwierdzenie nie wygasa: umowa wykonawcza rozciąga wyjątek na 12-miesięczny zakaz po odejściu, a ust.~3 i 5 chronią Spółkę tylko częściowo, bo dotyczą partnera, nie Założyciela-wspólnika. (4)~Obowiązek oferowania Spółce szans biznesowych (§~\ref{art:konkurencja} ust.~1 in fine) nie jest wyłączony, ale warto to wyraźnie potwierdzić. Dla inwestora VC udział founderów w spółkach celowych to sygnał „value leakage”; sąd oceni lojalność dyrektora według art.~300$^{54}$--300$^{55}$ §~1 niezależnie od wyjątku, który nie chroni przed odpowiedzialnością za nielojalność w konkretnej transakcji.}
\Rekomendacja{Ograniczyć wyjątek: udział pośredni przez Spółkę jako reguła; udział bezpośredni do 10\% w Przedsięwzięciu, w którym Spółka ma co najmniej równy udział, zatwierdzany przez Radę z głosem dyrektora ds. zgodności; powyżej tego --- uchwała WZ 75\% Głosów Uprawnionych w Sprawie z wyłączeniem zainteresowanego; wynagrodzenie od Przedsięwzięcia ujawnione i zatwierdzone; zatwierdzenie wygasa z Dniem Odejścia albo z wypowiedzeniem licencji, z obowiązkiem zaoferowania udziału Spółce po Wartości Godziwej; zachowanie obowiązku oferowania Spółce szans biznesowych.}
\Kierunek{Zmiana domyka lukę w zakazie konkurencji, kosztem elastyczności Założycieli; wymaga równoległej zmiany umowy wykonawczej. Proponowane brzmienie ust.~7 (streszczenie): udział Akcjonariusza Funkcjonalnego albo dyrektora, osobiście albo przez podmiot powiązany, w Przedsięwzięciu, w którym uczestniczy Spółka, nie jest działalnością konkurencyjną ani uczestnictwem w podmiocie konkurencyjnym w rozumieniu §~\ref{art:konkurencja} ust.~1, jeżeli łącznie:
\begin{itemize}
\item udział ujawniono Radzie i zatwierdzono uchwałą z zachowaniem ust.~6 oraz §~\ref{art:rada-zastrzezone} ust.~2, a przy udziale ponad 10\% praw udziałowych albo praw do zysku albo w Przedsięwzięciu, w którym Spółka nie ma co najmniej równego udziału --- uchwałą WZ większością 75\% Głosów Uprawnionych w Sprawie, z wyłączeniem zainteresowanego;
\item wynagrodzenie i inne korzyści od Przedsięwzięcia ujawniono i zatwierdzono w tej samej uchwale;
\item udział ujawniono w Załączniku nr~3.
\end{itemize}
Zatwierdzenie wygasa z Dniem Odejścia albo z dniem wypowiedzenia licencji z ust.~3; od tego dnia zainteresowany oferuje Spółce nabycie udziału po Wartości Godziwej. Wyłączenie nie obejmuje obowiązku przedstawienia Spółce szans biznesowych związanych z Przedsięwzięciem. Decyzja Założycieli: próg bezpośredniego udziału (proponowane 10\%) i obowiązek oferowania udziału Spółce po odejściu.}
\end{MemoBox}

\item Finansowanie Przedsięwzięcia Produkcyjnego bez regresu do Spółki jest dopuszczalne. Poręczenia, gwarancje i inne zabezpieczenia udzielone przez Spółkę za zobowiązania Przedsięwzięcia podlegają §~\ref{art:finansowanie} oraz progom określonym w §~\ref{art:rada-zastrzezone} i §~\ref{art:walne-zastrzezone}.
\end{enumerate}

\end{ContractClause}

\begin{ContractClause}{Zysk, wypłaty i rezerwy}
\label{art:zysk}
\begin{enumerate}
  \item O przeznaczeniu zysku decyduje Walne Zgromadzenie, z uwzględnieniem ograniczeń ustawowych, potrzeb finansowania rozwoju, zobowiązań grantowych i bezpieczeństwa płynności.
  \item W pierwszych 36 miesiącach od wpisu Spółki do rejestru rekomendowanym sposobem wykorzystania zysku jest jego reinwestowanie, chyba że Walne Zgromadzenie większością 75\% wszystkich głosów postanowi inaczej.
  \item Wypłata dywidendy, zaliczki, wypłata z kapitału akcyjnego albo inne świadczenie na rzecz akcjonariuszy może nastąpić wyłącznie po przeprowadzeniu wymaganych prawem testów i procedur, zgodnie z warunkami umów finansowania oraz gdy nie zagraża zdolności Spółki do wykonywania zobowiązań.
  \item Spółka tworzy wymagane prawem odpisy na kapitał akcyjny oraz może tworzyć kapitały rezerwowe na badania i rozwój, ochronę własności intelektualnej, cyberbezpieczeństwo, certyfikację i ryzyka eksportowe.
\end{enumerate}

\end{ContractClause}

\begin{ContractClause}{Rozwiązywanie Impasu Decyzyjnego}
\label{art:impas}
\begin{enumerate}
  \item „Impas Decyzyjny” występuje, gdy Sprawa Zastrzeżona nie zostanie rozstrzygnięta na dwóch prawidłowo zwołanych Walnych Zgromadzeniach odbytych w odstępie co najmniej 14 dni, a brak decyzji istotnie zagraża działalności Spółki.
  
\begin{MemoBox}{3.4.3}
\MemoPyt{3.4.3}{Czy tryb impasu (§~\ref{art:impas}) jest wykonalny i czy rozstrzygnięcie Niezależnego Eksperta w sprawach obiektywnych nie wkracza w kompetencje organów?}
\Ocena{zgodne}
\Klasa{opcjonalne}
\Podstawa{art.~353$^{1}$ KC \Wiar{Z}; art.~300$^{73}$--300$^{79}$ KSH (rada dyrektorów) i art.~300$^{80}$--300$^{81}$ KSH (uchwały akcjonariuszy) \Wiar{Z}; art.~17 §~3 KSH \Wiar{Z}; art.~1157 KPC (zdatność arbitrażowa: spory o prawa majątkowe) \Wiar{Z}; §~\ref{art:wycena} ust.~7.}
\Uzasadnienie{Mechanizm to umowa o postępowanie przedsądowe (negocjacje, mediacja) i o ekspertyzę arbitralną --- ustalenie przez osobę trzecią elementu stanu faktycznego, wiążące strony jak uzgodnienie umowne i kontrolowane przez sąd tylko w granicach rażącej nieprawidłowości; to konstrukcja §~\ref{art:wycena} ust.~7, mieszcząca się w art.~353$^{1}$ KC \Wiar{W}. Nie jest to zapis na sąd polubowny (art.~1157 pkt~1 KPC \Wiar{Z}), więc nie koliduje z przepisami o sądownictwie polubownym ani z wyłączną kompetencją sądu do uchylania uchwał. Ust.~4 poprawnie wyklucza zastępowanie uchwał organów: ekspert ustala przesłankę, decyzję podejmuje organ. Granice: (1)~w Sprawach Zastrzeżonych (uznaniowych) ekspert rozwiąże tylko spór o dane wejściowe --- impas właścicielski pozostaje, a brak buy-sell i shotgun jest świadomy; (2)~nie wskazano, kto stwierdza, że brak decyzji „istotnie zagraża działalności”; (3)~wybór eksperta przez Radę przy impasie między Założycielami-dyrektorami może być niemożliwy z powodu konfliktu interesów (§~\ref{art:konkurencja} ust.~6--7); (4)~ograniczenie Rady z ust.~5 działa tylko wewnętrznie (art.~17 §~3, art.~300$^{77}$ §~2 KSH \Wiar{Z}), co jest właściwe; (5)~mediacja z ust.~2 i §~\ref{art:final} ust.~4 dublują się bez szkody. Tryb nie ma znaczenia rejestrowego; przesłanką dopuszczalności powództwa byłby tylko, gdyby strony tak postanowiły --- obecnie nie ma takiego skutku, co jest bezpieczne wobec §~\ref{art:final} ust.~5.}
\Rekomendacja{Doprecyzować, że stan Impasu Decyzyjnego stwierdza Rada Dyrektorów uchwałą (a w razie jej niezdolności --- Przewodniczący), oraz przewidzieć wybór eksperta przez instytucję niezależną (np. organizację rzeczoznawców albo ośrodek mediacyjny), gdy Rada nie może działać; utrzymać brak mechanizmu wyjścia jako świadomą decyzję.}
\Kierunek{Zmiana usuwa dwa punkty sporne (kto stwierdza impas, kto wybiera eksperta przy konflikcie), nie zmienia kompetencji organów ani większości. Poniżej brzmienie ust.~1 zdanie drugie; brzmienie ust.~3 --- w ramce przy ust.~3.}
\begin{Brzmienie}
§~\ref{art:impas} ust.~1 zdanie drugie: „Wystąpienie Impasu Decyzyjnego stwierdza Rada Dyrektorów uchwałą, a jeżeli Rada Dyrektorów nie może podjąć uchwały --- jej Przewodniczący; stwierdzenie to nie przesądza treści przyszłej decyzji organu.”
\end{Brzmienie}
\end{MemoBox}

\item W razie Impasu Decyzyjnego Przewodniczący Rady Dyrektorów organizuje w terminie 10 dni spotkanie negocjacyjne z udziałem stron, a następnie — jeżeli Impas Decyzyjny trwa — mediację prowadzoną przez niezależnego mediatora posiadającego doświadczenie w sporach technologicznych.
  \item Jeżeli po mediacji pozostaje spór dotyczący kwestii możliwej do obiektywnego ustalenia, w szczególności wyceny, rachunkowości, zagadnienia technicznego albo zakresu określonego prawa własności intelektualnej, sprawa może zostać poddana rozstrzygnięciu przez Niezależnego Eksperta wybranego wspólnie przez strony sporu, a w braku porozumienia — przez Radę Dyrektorów spośród osób niezależnych od stron, z zastrzeżeniem zasad dotyczących Wartości Godziwej i konfliktu interesów.
  
\begin{MemoBox}{3.4.3}
\MemoPytCd{3.4.3}
\begin{Brzmienie}
§~\ref{art:impas} ust.~3 fragment końcowy: „[\ldots] a w braku porozumienia --- przez Radę Dyrektorów spośród osób niezależnych od stron, z zastrzeżeniem zasad dotyczących Wartości Godziwej i konfliktu interesów; jeżeli Rada Dyrektorów nie może dokonać wyboru wskutek wyłączenia dyrektorów, eksperta wskazuje niezależna instytucja wskazana w regulaminie Rady Dyrektorów.”
\end{Brzmienie}
\end{MemoBox}

\item Rozstrzygnięcie przez Niezależnego Eksperta nie może zastępować uchwały organu Spółki w sprawie strategicznej, personalnej, właścicielskiej albo innej Sprawie Zastrzeżonej, której rozstrzygnięcie wymaga decyzji uznaniowej lub określonej większości głosów.
  \item Do czasu rozstrzygnięcia Impasu Decyzyjnego obowiązuje ostatni zatwierdzony budżet, proporcjonalnie przedłużony, a Rada Dyrektorów może podejmować wyłącznie czynności konieczne dla ciągłości, bezpieczeństwa, ochrony praw i wykonania obowiązków ustawowych.
  \item Impas Decyzyjny nie uprawnia do obejścia wymogów dotyczących Dopuszczalnego Nabywcy, kontroli eksportu, ochrony własności intelektualnej ani większości kwalifikowanych.
\end{enumerate}
\end{ContractClause}

\begin{ContractClause}{Naruszenia i wykonanie obowiązków akcjonariusza}
\label{art:breach}
\begin{enumerate}
  \item Akcjonariusz, który narusza obowiązki dotyczące zbywania akcji, poufności, własności intelektualnej, sankcji, kontroli eksportu lub bezpieczeństwa, jest zobowiązany niezwłocznie zaniechać naruszenia, usunąć jego skutki i naprawić szkodę.
  \item Spółka może żądać zabezpieczenia roszczeń, wykonania zobowiązania w naturze, wydania korzyści, złożenia oświadczeń i dokumentów oraz zastosować inne środki przewidziane prawem lub odrębnymi umowami.
  \item W zakresie prawnie dopuszczalnym każdy Założyciel udzieli w odrębnej umowie wykonawczej pełnomocnictw i złoży oferty niezbędne do wykonania obowiązkowego zbycia akcji, przy czym pełnomocnictwo nie może być użyte poza ściśle określonym przypadkiem i procedurą.
\end{enumerate}

\end{ContractClause}

\begin{ContractClause}{Rozwiązanie i likwidacja Spółki}
\label{art:likwidacja}
\begin{enumerate}
  \item Rozwiązanie Spółki następuje w przypadkach przewidzianych prawem albo na podstawie uchwały Walnego Zgromadzenia podjętej większością 75\% wszystkich głosów.
  \item Likwidatorami są dyrektorzy, chyba że Walne Zgromadzenie postanowi inaczej.
  \item Przed wydaniem lub zniszczeniem nośników, prototypów i dokumentacji likwidatorzy dokonują klasyfikacji bezpieczeństwa, eksportowej i archiwalnej oraz zapewniają zgodne z prawem zabezpieczenie Kluczowej Własności Intelektualnej.
  \item Majątek pozostały po zaspokojeniu lub zabezpieczeniu wierzycieli dzieli się między akcjonariuszy proporcjonalnie do liczby akcji, z uwzględnieniem praw szczególnych, jeżeli zostaną ustanowione.
\end{enumerate}

\end{ContractClause}


\begin{ContractClause}{Postanowienia końcowe}
\label{art:final}
\begin{enumerate}
  \item \BoilerplateStrike{W sprawach nieuregulowanych stosuje się przepisy \KSH{} oraz inne właściwe przepisy prawa polskiego.}
  
\begin{MemoBox}{4.3.1}
\MemoPytCd{4.3.1}
Uwaga konieczna K17 (ramka główna --- s.~\pageref{memo:4.3.1}): akt notarialny nie może zawierać tekstu przekreślonego; wszystkie przekreślenia usunąć albo przywrócić po decyzji Założycieli, a~pola wypełnić.
\begin{Brzmienie}
§~\ref{art:final} ust.~1 (w miejsce zdania przekreślonego): W sprawach nieuregulowanych niniejszą umową stosuje się przepisy Kodeksu spółek handlowych oraz inne właściwe przepisy prawa polskiego.
\end{Brzmienie}
\end{MemoBox}

\item Jeżeli postanowienie umowy jest nieważne albo niewykonalne, pozostałe postanowienia zachowują moc. Strony zastąpią wadliwe postanowienie rozwiązaniem prawnie dopuszczalnym, możliwie najbliższym jego celowi gospodarczemu i bezpieczeństwa.
  \item Zmiana niniejszej umowy wymaga formy przewidzianej prawem i większości określonej w §~\ref{art:walne-zastrzezone}.
  
\begin{MemoBox}{3.4.2}
\MemoPyt{3.4.2}{Czy brak klauzul zgody indywidualnej (ochrony konkretnego akcjonariusza przed zmianą) jest bezpieczny dla założycieli mniejszościowych i czy inwestor VC będzie ich wymagał?}
\Ocena{zgodne warunkowo}
\Klasa{zalecane}
\Podstawa{art.~300$^{98}$ §~2--5 KSH (większość 3/4 dla zmiany umowy; zgoda wszystkich akcjonariuszy, których dotyczy uchwała zwiększająca świadczenia lub uszczuplająca prawa indywidualne; oddzielne głosowanie w grupie akcji jednorodzajowych) \Wiar{Z}; art.~20 KSH (równe traktowanie) \Wiar{Z}; art.~300$^{101}$ w zw. z art.~422 §~1 KSH (uchylenie uchwały krzywdzącej akcjonariusza) \Wiar{Z}; art.~300$^{25}$, art.~300$^{26}$ (akcje założycielskie), art.~300$^{28}$ KSH \Wiar{Z}; art.~419 §~1 KSH (S.A., porównawczo) \Wiar{Z}.}
\Uzasadnienie{Ustawowe minimum ochrony: zgoda wszystkich akcjonariuszy, których dotyczy zmiana umowy zwiększająca świadczenia lub uszczuplająca prawa indywidualne (art.~300$^{98}$ §~3 \Wiar{Z}); oddzielne głosowanie w grupie akcji jednorodzajowych, których prawa narusza zmiana (§~4--5 \Wiar{Z}); równe traktowanie (art.~20 \Wiar{Z}); uchylenie uchwały krzywdzącej (art.~422 §~1 w zw. z art.~300$^{101}$ \Wiar{Z}). Umowa nie przyznaje uprawnień indywidualnych, a akcje serii A są jednorodzajowe, więc pierwsza ochrona zadziała tylko przy „zwiększeniu świadczeń” (np. zaostrzeniu zwrotnego zbycia; czy obowiązek zbycia akcji jest „świadczeniem”, jest sporne \Wiar{?}), a druga dopiero przy seriach o różnych uprawnieniach. Poza tym 75\% wszystkich głosów (A+B+jeden z C/D/E) może zmienić wszystko; C, D i E nie mają weta, A ma je zawsze, B tylko w koalicji. Mniejszość jest bezpieczna o tyle, o ile ufa Założycielowi A; sąd chroni ex post i powoli. Indywidualne weta pozwoliłyby Założycielowi z 10\% zablokować rundę, czego inwestor nie zaakceptuje. Rozsądny środek to ochrona klasowa: zgoda większości serii na uszczuplenie jej praw i zgoda dotkniętych na zaostrzenie obowiązków związanych z akcją. Inwestor VC zażąda protective provisions w umowie spółki, bo art.~300$^{98}$ §~4 nie obejmuje sprzedaży Spółki, długu ani nowej serii o równych prawach; artykuł na portalu startup.pfr.pl opisuje zgodę funduszu albo bardzo dużą większość (np. 75--80\%) \Wiar{Z}, lecz jako postanowienia umowy inwestycyjnej. Założyciele mogą rozważyć akcje założycielskie chroniące przed rozwodnieniem głosów (art.~300$^{26}$ §~1 \Wiar{Z}) --- inwestor VC zwykle ich nie akceptuje.}
\Rekomendacja{Dodać w §~\ref{art:final} ust.~3 zasadę zgody serii i zgody akcjonariuszy dotkniętych zaostrzeniem obowiązków związanych z akcją; nie wprowadzać indywidualnych wet dla poszczególnych Założycieli (decyzja Założycieli, jeżeli mniejszość oczekuje innej równowagi).}
\Kierunek{Dla mniejszości: ochrona przed zaostrzeniem vestingu, lock-upu i drag bez jej zgody; dla inwestora: gotowa rama protective provisions; dla większości: bez wpływu na zwykłe zmiany umowy; sąd rejestrowy bada dodatkową zgodę jak przy art.~300$^{98}$ §~3--4 KSH.}
\begin{Brzmienie}
§~\ref{art:final} ust.~3: „Zmiana niniejszej umowy wymaga formy przewidzianej prawem i większości określonej w §~\ref{art:walne-zastrzezone}. Niezależnie od art.~300$^{98}$ §~3 i~4 KSH zmiana umowy, która uszczupla prawa akcjonariuszy danej serii akcji w porównaniu z pozostałymi seriami, wymaga ponadto zgody akcjonariuszy posiadających większość akcji tej serii, wyrażonej w odrębnym głosowaniu. Zmiana umowy zwiększająca obowiązki związane z akcją, w szczególności zaostrzająca zasady określone w §~\ref{art:zwrotne-zbycie}, §~\ref{art:okres-ograniczenia-zbywania} albo §~\ref{art:przymusowe-wspolzbycie} wobec akcjonariuszy nimi objętych, wymaga zgody każdego akcjonariusza, którego dotyczy.”
\end{Brzmienie}
\end{MemoBox}

\item Spory związane z niniejszą umową strony w pierwszej kolejności poddają negocjacjom i mediacji. W razie braku ugody właściwy jest sąd powszechny właściwy dla siedziby Spółki, chyba że bezwzględnie obowiązujące przepisy stanowią inaczej.
  \item Obowiązek negocjacji i mediacji nie ogranicza prawa strony do żądania zabezpieczenia, podjęcia czynności koniecznej do przerwania lub zawieszenia biegu przedawnienia ani innej czynności niecierpiącej zwłoki.
  \item Załączniki nr 1--3 stanowią integralną część niniejszej umowy.

\begin{MemoBox}{3.4.1}
\MemoPyt{3.4.1}{Czy aktualizacja Załączników nr~1--3 (np. zmiana danych, wpis nowego Założyciela Serii F, ujawnione aktywności) zawsze wymaga zmiany umowy w formie aktu notarialnego i wpisu do KRS, czy część danych można przenieść do uchwał albo rejestrów?}
\Ocena{zgodne warunkowo}
\Klasa{zalecane}
\Podstawa{art.~300$^{5}$ §~1 pkt~3--4, art.~300$^{100}$ §~2, art.~300$^{102}$ §~1--2 KSH \Wiar{Z}; art.~300$^{1}$ §~3, art.~300$^{9}$ §~2 KSH \Wiar{Z}; art.~300$^{30}$--300$^{33}$ KSH (rejestr akcjonariuszy; art.~300$^{33}$ §~1 pkt~10--11 i §~2) \Wiar{Z}; art.~300$^{33}$ §~3 KSH w brzmieniu od 18 lutego 2027~r. (Dz.U. 2026 poz.~176) \Wiar{Z}; art.~9 ust.~4 ustawy o KRS (tekst jednolity umowy) \Wiar{Z}.}
\Uzasadnienie{Załączniki są częścią umowy (ust.~6), więc każda zmiana ich treści jest zmianą umowy: uchwała 75\% w protokole notarialnym (art.~300$^{100}$ §~2 \Wiar{Z}), zgłoszenie w sześć miesięcy, wpis konstytutywny (art.~300$^{102}$ §~1--2 \Wiar{Z}), tekst jednolity (art.~9 ust.~4 ustawy o KRS \Wiar{Z}). Nie każde zdarzenie wymaga jednak aktualizacji. \emph{Załącznik nr~1} (dane założycielskie) jest historyczny: zbycie akcji, zmiana adresu albo nazwiska są ujawniane w rejestrze akcjonariuszy (art.~300$^{33}$ §~1 pkt~5--6 \Wiar{Z}), a wniesienie wkładu stwierdza uchwała Rady (art.~300$^{9}$ §~2 \Wiar{Z}); zmiany wymaga tylko błąd w danych założycielskich. \emph{Załącznik nr~2}: obowiązek zbycia jest obowiązkiem związanym z akcją, który musi wynikać z umowy (art.~300$^{1}$ §~3), więc krąg osób i data muszą być z niej ustalalne; dla Założycieli Serii F umowa odsyła do uchwały kwalifikacyjnej i umowy wykonawczej (§~\ref{art:seria-f} ust.~2 i~9) --- konstrukcja obronna. Rejestr akcjonariuszy ujawnia ograniczenia rozporządzania i obowiązki związane z akcją (art.~300$^{33}$ §~1 pkt~10--11 \Wiar{Z}), a od 18 lutego 2027~r. Rada zgłasza zmiany w siedem dni (art.~300$^{33}$ §~3 \Wiar{Z}); uchwałę kwalifikacyjną warto przekazywać z dyspozycją wpisu, a na podstawie art.~300$^{33}$ §~2 wskazać, że dane te podlegają ujawnieniu \Wiar{?} (praktyka podmiotów prowadzących rejestr). Zmiana daty albo statusu Założyciela Pierwotnego pozostaje zmianą umowy. \emph{Załącznik nr~3} (ujawnienia, zgody konkurencyjne) może być aktualizowany uchwałą Walnego Zgromadzenia z ewidencją prowadzoną przez Radę --- wymaga to zmiany odesłania w §~\ref{art:konkurencja} ust.~2 i zastrzeżenia w ust.~6.}
\Rekomendacja{Pozostawić Załączniki nr~1--2 w umowie (z zastrzeżeniem, że dane historyczne nie wymagają aktualizacji), a Załącznik nr~3 uczynić dokumentem aktualizowanym uchwałą Walnego Zgromadzenia bez zmiany umowy.}
\Kierunek{Zmniejsza liczbę protokołów notarialnych (każde nowe ujawnienie albo zgoda konkurencyjna), nie osłabia ochrony, bo zgoda i tak wymaga uchwały; sąd rejestrowy przestaje być angażowany w sprawy informacyjne. Poniżej brzmienie ust.~6; brzmienie §~\ref{art:konkurencja} ust.~2 --- w ramce przy tym ustępie.}
\begin{Brzmienie}
§~\ref{art:final} ust.~6: „Załączniki nr~1--3 stanowią integralną część niniejszej umowy. Zmiana danych osobowych lub adresowych Założyciela, zbycie akcji oraz stwierdzenie wniesienia wkładu nie wymagają zmiany Załącznika nr~1 i są ujawniane w rejestrze akcjonariuszy albo uchwałą Rady Dyrektorów. Późniejsze ujawnienia i zgody, o których mowa w Załączniku nr~3, mogą następować uchwałą Walnego Zgromadzenia podjętą większością wymaganą dla danej zgody, bez zmiany niniejszej umowy; Rada Dyrektorów prowadzi ich ewidencję.”
\end{Brzmienie}
\end{MemoBox}

\begin{MemoBox}{4.3.1}
\MemoPyt{4.3.1}{Jaka forma zawiązania i jakie dokumenty do KRS są wymagane (w tym dokumentacja wkładów niepieniężnych) i które postanowienia projektu mogą wywołać zastrzeżenia sądu rejestrowego (ograniczenia zbywalności, Kryterium, obowiązki związane z akcją, uchwały ramowe)?}
\Ocena{zgodne warunkowo --- co do treści; w~obecnej postaci redakcyjnej (przekreślenia, oznaczenia robocze, puste pola) tekst nie nadaje się do aktu i~wywołałby wezwanie sądu, co kwalifikujemy jako ryzyko formalne}
\Klasa{konieczne --- odstępstwo od reguły uzasadnione: bez czynności redakcyjnych tekst nie może stać się aktem notarialnym; zmiany merytoryczne nie są konieczne (K16, K17, K18)}
\Podstawa{art.~300$^{5}$ §~1 KSH (elementy umowy) i~art.~300$^{6}$ KSH (forma aktu notarialnego) \Wiar{Z}; art.~300$^{11}$ §~1--2 KSH (spółka w~organizacji; reprezentacja) i~art.~300$^{12}$ §~1--4 KSH (zgłoszenie, załączniki, lista akcjonariuszy) \Wiar{Z}; art.~300$^{13}$ §~1--2 w~zw. z~art.~164 §~3, art.~165, art.~169 §~1 i~art.~172 KSH \Wiar{Z}; art.~19 ust.~2 i~5--6, art.~19a ust.~5, 5a i~5d, art.~20a ust.~1 i~3, art.~23 ust.~1 ustawy o~KRS \Wiar{Z}; art.~300$^{39}$ §~1--2, art.~300$^{33}$ §~1 pkt~10--11, art.~300$^{98}$ §~2 i~art.~300$^{103}$ KSH \Wiar{Z}; ustawa o~PCC: art.~1 ust.~1 pkt~1 lit.~k, art.~1a pkt~1--2, art.~6 ust.~1 pkt~8 lit.~a, art.~7 ust.~1 pkt~9 \Wiar{Z}.}
\Uzasadnienie{\textbf{Forma.} Umowa nie korzysta ze wzorca, więc wymaga aktu notarialnego (art.~300$^{6}$); pełnomocnik Założyciela --- z~pełnomocnictwem w~formie aktu notarialnego. Załączniki nr~1--3 muszą być objęte aktem (Załącznik nr~1 zawiera elementy obowiązkowe z~art.~300$^{5}$ §~1 pkt~3--5). Spółkę w~organizacji reprezentuje Rada, a~do jej ustanowienia --- pełnomocnik powołany jednomyślną uchwałą akcjonariuszy (art.~300$^{11}$ §~2), dlatego pierwszą Radę należy powołać w~akcie. Opłata 500~zł (art.~52 ust.~1 u.k.s.c.), 100~zł za ogłoszenie w~MSiG \Wiar{W}. PCC: definicje spółek w~art.~1a pkt~1--2 nie obejmują P.S.A., a~podstawą jest kapitał zakładowy, którego P.S.A. nie ma --- literalnie umowa nie podlega PCC, co potwierdzają interpretacje indywidualne \Wiar{W}; wobec braku wyraźnego wyłączenia do potwierdzenia przez doradcę podatkowego (T9).

\textbf{Wkłady.} Kapitał akcyjny pokrywa się przed zgłoszeniem (oświadczenia z~art.~300$^{12}$ §~3 pkt~2--3; §~\ref{art:wklady} ust.~4). Wkłady niepieniężne można wnieść po wpisie, ale dokumenty rozporządzające muszą istnieć (umowy przeniesienia praw autorskich i~praw z~patentów na piśmie pod rygorem nieważności, protokoły wydania). Kierowanie na kapitał akcyjny tylko wkładów pieniężnych ogranicza odpowiedzialność z~art.~300$^{10}$ §~1.

\textbf{Dokumenty do KRS.} Wniosek wyłącznie teleinformatycznie, podpisany przez wszystkich dyrektorów albo pełnomocnika, w~6 miesięcy od zawarcia umowy pod rygorem jej rozwiązania; załączniki: umowa, oświadczenia o~kapitale i~wkładach, dowód ustanowienia organów (chyba że w~akcie), adresy do doręczeń, lista akcjonariuszy z~art.~300$^{12}$ §~4 (osobny dokument), zgody dyrektorów (art.~19a ust.~5) i~pełnomocnik do doręczeń przy adresie poza UE (art.~19a ust.~5a). Po wpisie: umowa z~podmiotem prowadzącym rejestr, CRBR w~14 dni.

\textbf{Ryzyko wezwania} (sąd bada formę i~treść dokumentów; brak usuwalny --- termin pod rygorem odmowy wpisu; nie odmawia z~powodu drobnych uchybień):
\begin{itemize}
\item przekreślenia, oznaczenie „[WYMAGA OPINII PRAWNIKA]”, puste pola --- \textbf{wysokie (formalne)};
\item Kryterium (art.~300$^{39}$ §~1 \Wiar{Z}) --- niskie--umiarkowane; sąd nie bada go z~urzędu;
\item derogacja art.~300$^{39}$ §~3--6 w~lock-upie, obowiązki związane z~akcją, treść §~\ref{art:transakcje-akcjonariusze}, większości 75\% i~4/5 --- niskie;
\item uchwała ramowa serii P (art.~300$^{103}$) --- niskie przy wpisie spółki, umiarkowane przy wpisie emisji;
\item mandat do 6~miesięcy po kadencji --- niskie--umiarkowane.
\end{itemize}
Przy prawidłowej redakcji ryzyko merytoryczne jest niskie.}
\Rekomendacja{Sporządzić wersję „do aktu” bez przekreśleń, oznaczeń i~pól; objąć aktem uchwały założycielskie (powołanie pierwszej Rady, wybór podmiotu prowadzącego rejestr, ewentualnie zatwierdzenie Załączników); przygotować pakiet dokumentów wkładowych na dzień aktu, listę akcjonariuszy z~art.~300$^{12}$ §~4 i~--- jeżeli którykolwiek dyrektor ma adres poza UE --- pełnomocnika do doręczeń. Potwierdzić u~doradcy podatkowego brak PCC. Uprawniony prawnik powinien potwierdzić bieżącą praktykę PRS co do formy załączników.}
Brzmienie do K18 (nowy ust.~7) i~pozostałe czynności redakcyjne objęte K17; K16 i~brzmienie §~\ref{art:final} ust.~1 (K17) --- w~ramkach przy §~\ref{art:transakcje-akcjonariusze} i~§~\ref{art:final} ust.~1.
\begin{Brzmienie}
§~\ref{art:final} --- nowy ust.~7: Pierwszych dyrektorów oraz podmiot prowadzący rejestr akcjonariuszy wskazują Założyciele w uchwałach objętych aktem zawiązania Spółki. Do czasu wpisu Spółki do rejestru przedsiębiorców Spółkę w organizacji reprezentuje Rada Dyrektorów zgodnie z §~\ref{art:reprezentacja}, a zaciąganie zobowiązań innych niż niezbędne do zawiązania i rejestracji Spółki wymaga zgody wszystkich Założycieli.

Pozostałe przekreślone zdania (§~\ref{art:akcje} ust.~6, §~\ref{art:reprezentacja} ust.~3, §~\ref{art:wz} ust.~1, §~\ref{art:konkurencja} ust.~12 i inne) --- usunąć albo przywrócić bez przekreślenia zgodnie z decyzją Założycieli; §~\ref{art:dopuszczalny-nabywca} ust.~14 --- usunąć albo przywrócić w brzmieniu poprawionym (pakiet 1.4); wszystkie pola \field{...} --- wypełnić.
\end{Brzmienie}
\end{MemoBox}

\begin{MemoBox}{4.3.3}
\MemoPyt{4.3.3}{Czy nowelizacja KSH z 23 stycznia 2026 r. (wejście w życie 18 lutego 2027 r.) wpływa na którekolwiek z postanowień projektu i czy warto już dziś przygotować umowę pod nowe przepisy?}
\Ocena{zgodne warunkowo --- tekst ustawy sprawdzony w~całości; żadne postanowienie projektu nie jest z~nią sprzeczne, ale nowelizacja nakłada na Spółkę nowe obowiązki proceduralne i~obowiązek dostosowania umowy w~2 lata}
\Klasa{zalecane}
\Podstawa{ustawa z~23 stycznia 2026~r. o~zmianie ustawy --- Kodeks spółek handlowych oraz niektórych innych ustaw (Dz.U. 2026 poz.~176, ogłoszona 17 lutego 2026~r.), art.~35 (wejście w~życie 18 lutego 2027~r., z~wyjątkiem art.~28 i~33 --- 28 lutego 2026~r.) \Wiar{Z}; art.~1 pkt~2--6 i~30 (art.~300$^{32}$, 300$^{33}$, 300$^{34}$, 300$^{35}$, 300$^{37}$ §~2 i~art.~594 §~1 KSH) \Wiar{Z}; art.~5 (art.~10 ust.~4a pkt~4, art.~25da i~art.~38 pkt~8a lit.~j ustawy o~KRS) \Wiar{Z}; art.~28--34 (przepisy przejściowe) \Wiar{Z}; art.~300$^{29}$ §~1 KSH („akcje nie mają formy dokumentu”) \Wiar{Z}; art.~111 §~2 i~art.~112 KC (obliczanie terminów) \Wiar{Z}.}
\Uzasadnienie{\textbf{Zmiany dla P.S.A.} Rada zgłasza do sądu zawarcie umowy o~prowadzenie rejestru (art.~300$^{32}$ §~1$^{1}$--1$^{3}$), a~KRS ujawnia podmiot (art.~38 pkt~8a lit.~j ustawy o~KRS); zmiany danych z~§~1 pkt~1--4 i~9--11 (m.in. ograniczeń i~obowiązków związanych z~akcją) Rada zgłasza podmiotowi w~7 dni pod rygorem grzywny do 20~000~zł (art.~300$^{33}$ §~3, art.~594 §~1 pkt~2$^{1}$); forma zgody na wpis (art.~300$^{34}$ §~3 zdanie drugie); ochrona PESEL i~adresów (art.~300$^{35}$ §~1$^{1}$). Przejściowo: wniosek (250~zł) o~wpis podmiotu do 18 maja 2027~r. (art.~29; poprzednia wersja podawała 17 maja), dostosowanie umów P.S.A. do 18 lutego 2029~r. (art.~34). Art.~300$^{36}$, 300$^{38}$, 300$^{39}$, przepisy o~organach i~emisjach --- bez zmian; zniesienie akcji imiennych dotyczy spółki akcyjnej; obowiązki „zarządu” wykonuje Rada Dyrektorów \Wiar{W}. Termin 7 dni wiąże Spółkę; termin dla akcjonariusza to materia umowna. \textbf{Wpływ na projekt} --- punkty styku: „akcje imienne” w~§~\ref{art:akcje} ust.~2 i~§~\ref{art:seria-f} (zbędne, nie błędne); brak terminu w~§~\ref{art:wz} ust.~3; zgłaszanie adnotacji o~mechanizmie zwrotnego zbycia (4.3.2); forma zgody przy zwrotnym zbyciu i~zbyciu po utracie Kryterium (pakiet 1.5); art.~34 obejmie Spółkę --- klauzula dostosowawcza usuwa spór o~znaczenie odesłań. Z~zawiązaniem nie ma powodu czekać: koszty obu ścieżek są porównywalne.}
\Rekomendacja{Dodać klauzulę dostosowawczą do §~\ref{art:final} i~termin 7 dni do §~\ref{art:wz} ust.~3. W~umowie wykonawczej (pakiet 1.5) przewidzieć składanie przez Założycieli zgód na wpis w~formie z~art.~300$^{34}$ §~3 zdanie drugie. Zawiązać Spółkę bez oczekiwania na 2027~r.; w~kalendarzu Rady zapisać: zgłoszenie podmiotu prowadzącego rejestr do KRS do 18 maja 2027~r., przegląd umowy pod kątem art.~34 do 18 lutego 2029~r.}
\Kierunek{Klauzula dostosowawcza: dla Spółki i~Rady --- podstawa do wykonywania nowych obowiązków bez zmiany umowy; dla akcjonariuszy --- jasny termin aktualizacji danych; dla inwestora --- neutralne; dla rejestru i~sądu --- zgodność z~przyszłym stanem prawnym. Proponowane brzmienie (§~\ref{art:wz} ust.~3 --- w~ramce przy tym ustępie):}
\begin{Brzmienie}
§~\ref{art:final} --- nowy ust.~8: Odesłania do przepisów prawa zawarte w niniejszej umowie rozumie się jako odesłania do tych przepisów w brzmieniu obowiązującym w dniu dokonania danej czynności, a w razie ich uchylenia albo zmiany numeracji --- do przepisów, które je zastąpiły. Obowiązki Spółki i akcjonariuszy wobec rejestru akcjonariuszy oraz rejestru przedsiębiorców wykonuje się w terminach i w sposób określony przepisami obowiązującymi w dniu ich wykonywania. Rada Dyrektorów przedstawia Walnemu Zgromadzeniu projekt dostosowania umowy w terminie 6 miesięcy od dnia wejścia w życie zmiany przepisów, która wymaga zmiany umowy.
\end{Brzmienie}
\end{MemoBox}

\end{enumerate}
\end{ContractClause}

\newpage
\section*{Załącznik nr 1 — Założyciele, akcje i wkłady}
\addcontentsline{toc}{section}{Załącznik nr 1 — Założyciele, akcje i wkłady}

\begin{longtable}{@{}p{1.1cm} p{3.0cm} p{1.9cm} p{2.1cm} p{6.0cm}@{}}
\toprule
\textbf{Zał.} & \textbf{Dane} & \textbf{Akcje} & \textbf{Cena łączna} & \textbf{Przedmiot wkładu — dokładny opis} \\
\midrule
\endfirsthead
\toprule
\textbf{Zał.} & \textbf{Dane} & \textbf{Akcje} & \textbf{Cena łączna} & \textbf{Przedmiot wkładu — dokładny opis} \\
\midrule
\endhead
A & \field{pełne dane} & A0000001--A0500000 & 5 000 zł & \field{kod / sprzęt / patent; repozytorium, wersja, numery seryjne, numer zgłoszenia, dokument nabycia praw} \\
B & \field{pełne dane} & A0500001--A0700000 & 2 000 zł & \field{kod / sprzęt / patent; dokładny opis} \\
C & \field{pełne dane} & A0700001--A0800000 & 1 000 zł & \field{kod / sprzęt / patent; dokładny opis} \\
D & \field{pełne dane} & A0800001--A0900000 & 1 000 zł & \field{kod / sprzęt / patent; dokładny opis} \\
E & \field{pełne dane} & A0900001--A1000000 & 1 000 zł & \field{kod / sprzęt / patent; dokładny opis} \\
\bottomrule
\end{longtable}

\textbf{Wkłady pieniężne Założycieli (§~\ref{art:wklady} ust.~1 lit.~d i ust.~4):}

\begin{longtable}{@{}p{1.1cm} p{3.8cm} p{3.8cm} p{4.5cm}@{}}
\toprule
\textbf{Zał.} & \textbf{Kwota wkładu pieniężnego} & \textbf{Termin wniesienia} & \textbf{Uwagi} \\
\midrule
A & \field{kwota, np. 25 000 zł} & przed wpisem do rejestru & \field{---} \\
B & \field{kwota} & przed wpisem do rejestru & \field{---} \\
C & \field{kwota} & przed wpisem do rejestru & \field{---} \\
D & \field{kwota} & przed wpisem do rejestru & \field{---} \\
E & \field{kwota} & przed wpisem do rejestru & \field{---} \\
\midrule
\multicolumn{2}{@{}l}{\textbf{Razem (kapitał akcyjny):} \field{suma}} & & \\
\bottomrule
\end{longtable}

Dla każdego składnika należy dołączyć protokół wydania, wykaz plików lub repozytoriów, sumy kontrolne albo oznaczenie wersji, dokumentację nabycia praw, wykaz twórców, wykaz licencji zewnętrznych oraz — w przypadku patentów lub zgłoszeń — numery, terytoria i dokumenty cesji.


\begin{MemoBox}{3.1.4}
\MemoPyt{3.1.4}{Załącznik nr~1 zawiera pola do uzupełnienia --- jakie dane muszą być kompletne w chwili podpisania umowy założycielskiej?}
\Ocena{ryzyko (Załącznik nr~1 nie rozdziela akcji obejmowanych za wkłady niepieniężne od obejmowanych za wkłady pieniężne i nie zawiera wartości wkładów)}
\Klasa{konieczne (K13)}
\Podstawa{art.~300$^{5}$ §~1 pkt~3--5 i §~2 KSH \Wiar{Z}; art.~300$^{6}$ KSH (forma aktu notarialnego) \Wiar{Z}; art.~300$^{9}$ §~3, art.~300$^{12}$ §~2 pkt~7, §~3 i §~4 KSH (zgłoszenie, oświadczenia, lista akcjonariuszy) \Wiar{Z}; art.~92 §~1 ustawy Prawo o notariacie \Wiar{W}; art.~19 ust.~1 ustawy o KRS \Wiar{Z}.}
\Uzasadnienie{Załączniki są integralną częścią umowy (§~\ref{art:final} ust.~6), więc w akcie notarialnym nie może zostać żadne pole otwarte. Obowiązkowe są: (1)~dane identyfikujące każdego Założyciela; (2)~liczba, seria i numery akcji każdego z nich oraz cena emisyjna (art.~300$^{5}$ §~1 pkt~3) --- zapewniają to §~\ref{art:cap-table} i kolumna „Akcje”; (3)~przedmiot każdego wkładu niepieniężnego, \emph{serie i numery akcji obejmowanych za ten wkład} oraz osoba obejmująca (pkt~4) --- tego Załącznik nr~1 nie zapewnia: każdy Założyciel obejmuje akcje częściowo za wkład niepieniężny, częściowo za pieniężny, a tabela przypisuje mu jeden przedział numerów; (4)~praca lub usługi (pkt~5) --- nie dotyczy serii A. Terminy wniesienia wkładów określa §~\ref{art:wklady} ust.~4 (art.~300$^{5}$ §~2 \Wiar{Z}). Przypisanie numerów do wkładu jest odstępstwem od równomiernego zaliczania wkładów na wszystkie akcje, dopuszczalnym, gdy „umowa spółki stanowi inaczej” (art.~300$^{9}$ §~3 \Wiar{Z}), więc musi wynikać wprost z Załącznika. §~\ref{art:wklady} ust.~3 odsyła do wartości z Załącznika nr~1, a ten ma tylko „Cenę łączną” (cenę emisyjną), która nie jest wartością wkładu; po zmianie z 3.1.2 wartość jest też podstawą oświadczenia o kapitale. Załącznik nr~2 wymaga „TAK/NIE” i daty dla każdego Założyciela (§~\ref{art:zwrotne-zbycie} ust.~1), Załącznik nr~3 --- treści albo „brak” w każdym punkcie; „np. 25\,000~zł” to pole robocze. Sąd rejestrowy porówna Załącznik nr~1 z listą akcjonariuszy i oświadczeniami (art.~300$^{12}$ §~3 pkt~2--3 i §~4 \Wiar{Z}); w sporze przypisane numery decydują, których akcji dotyczą art.~300$^{10}$ i art.~14 §~2 KSH.}
\Rekomendacja{Przebudować Załącznik nr~1 na dwie tabele: (A)~akcje za wkłady niepieniężne --- numery akcji, przedmiot, wartość, podstawa wyceny, dokumenty przeniesienia; (B)~akcje za wkłady pieniężne --- numery akcji, kwota, termin. Suma numerów obu tabel dla każdego Założyciela musi odpowiadać tabeli z §~\ref{art:cap-table} ust.~1. Uzupełnić wszystkie pola przed wizytą u notariusza; Załączniki nr~2 i~3 uzupełnić albo wpisać „brak”.}
\textit{Tabele z proponowanego brzmienia (w odpowiedzi: tabele) przedstawiono wierszami: pierwszy punkt --- nagłówki kolumn, kolejne --- wiersze; kolumny rozdziela „\textbar”.}
\begin{Brzmienie}
Załącznik nr~1 --- Założyciele, akcje i wkłady. Część A. Akcje obejmowane za wkłady niepieniężne (§~\ref{art:wklady} ust.~1 lit.~a--c):

\begin{tabularx}{\linewidth}{p{0.9cm} p{2.4cm} X p{1.9cm} p{2.6cm}}
\toprule
Zał. & Numery akcji & Przedmiot wkładu (dokładny opis, repozytorium i wersja, numery seryjne, numer zgłoszenia) & Wartość wkładu & Podstawa wyceny; dokument przeniesienia \\
\midrule
A & \field{A0000001--A0\ldots} & \field{opis} & \field{kwota} & \field{raport / metoda; umowa z dnia} \\
\ldots & & & & \\
\bottomrule
\end{tabularx}

Część B. Akcje obejmowane za wkłady pieniężne (§~\ref{art:wklady} ust.~1 lit.~d i ust.~4):

\begin{tabularx}{\linewidth}{p{0.9cm} p{3.2cm} X p{3.4cm}}
\toprule
Zał. & Numery akcji & Kwota wkładu pieniężnego & Termin wniesienia \\
\midrule
A & \field{A0\ldots--A0500000} & \field{kwota} & przed zgłoszeniem Spółki do rejestru \\
\ldots & & & \\
\bottomrule
\end{tabularx}

„Numery akcji z Części A i Części B przypisane danemu Założycielowi obejmują łącznie wszystkie akcje objęte przez niego zgodnie z §~\ref{art:cap-table} ust.~1. Łączna wartość wkładów z Części A i B stanowi kapitał akcyjny w rozumieniu §~\ref{art:wklady} ust.~4.”
\end{Brzmienie}
\end{MemoBox}

\section*{Załącznik nr 2 — Założyciele Pierwotni objęci mechanizmem czasowym}
\addcontentsline{toc}{section}{Załącznik nr 2 — Założyciele Pierwotni objęci mechanizmem czasowym}

Niniejszy załącznik dotyczy wyłącznie Założycieli Pierwotnych posiadających akcje serii A. Dane Założycieli Funkcjonalnych Serii F, ich role, Dni Przyznania Serii F, liczby akcji i warunki zaangażowania określają właściwe uchwały Walnego Zgromadzenia oraz umowy wykonawcze, bez potrzeby zmiany niniejszego załącznika przy każdej późniejszej emisji serii F.

\begin{longtable}{@{}p{2.2cm} p{3.0cm} p{3.0cm} p{3.0cm} p{2.8cm}@{}}
\toprule
\textbf{Założyciel} & \textbf{Czy funkcjonalny?} & \textbf{Liczba akcji objętych} & \textbf{Data rozpoczęcia} & \textbf{Rola i minimalne zaangażowanie} \\
\midrule
A & \field{TAK/NIE} & \field{liczba} & \field{data} & \field{np. CEO, etat/udział czasu} \\
B & \field{TAK/NIE} & \field{liczba} & \field{data} & \field{np. CTO} \\
C & \field{TAK/NIE} & \field{liczba} & \field{data} & \field{rola} \\
D & \field{TAK/NIE} & \field{liczba} & \field{data} & \field{rola} \\
E & \field{TAK/NIE} & \field{liczba} & \field{data} & \field{rola} \\
\bottomrule
\end{longtable}

\section*{Załącznik nr 3 — Ujawnione aktywności, licencje i ograniczenia}
\addcontentsline{toc}{section}{Załącznik nr 3 — Ujawnione aktywności i ograniczenia}

\begin{enumerate}
  \item Ujawnione wcześniejsze zatrudnienie, działalność akademicka, granty i prawa osób trzecich: \field{uzupełnić}.
  \item Dopuszczone wcześniejsze aktywności konkurencyjne lub inwestycje: \field{uzupełnić albo wpisać „brak”}.
  \item Wykaz otwartego oprogramowania i licencji o podwyższonym ryzyku: \field{uzupełnić}.
  \item Ograniczenia eksportowe, pochodzenie komponentów i klasyfikacje: \field{uzupełnić}.
  \item Udziały i funkcje Założycieli w Przedsięwzięciach Produkcyjnych (§~\ref{art:przedsiewziecia} ust.~7): \field{uzupełnić albo wpisać „brak”}.
\end{enumerate}

\vspace{1cm}
\begin{center}
\begin{tabularx}{\textwidth}{X X}
\field{Założyciel A — podpis} & \field{Założyciel B — podpis} \\
\\[1.3cm]
\field{Założyciel C — podpis} & \field{Założyciel D — podpis} \\
\\[1.3cm]
\field{Założyciel E — podpis} & \\
\end{tabularx}
\end{center}


\clearpage
\chapter*{Część II --- Sekcje tematyczne}
\addcontentsline{toc}{chapter}{Część II --- Sekcje tematyczne}

Sekcje T1--T9 realizują punkty zakresu, które nie mają wprost odpowiadającego pytania na liście szczegółowej: zmiany przepisów (T1), wykonalność mechanizmów (T2), ścieżka zawiązania i dokumenty do KRS (T3), dokumentacja wkładów (T4), rejestr akcjonariuszy (T5), PKD i ryzyka zastrzeżeń sądu (T6), oczekiwania inwestorów (T7), podział między umowę spółki a umowę akcjonariuszy (T8) i ryzyka podatkowe (T9). Tam, gdzie odpowiedź na pytanie w Części I już rozstrzyga zagadnienie, sekcja tylko do niej odsyła.

\section*{T1. Zmiany przepisów i ich wpływ na umowę oraz procedury Spółki}
\addcontentsline{toc}{section}{T1. Zmiany przepisów i ich wpływ na umowę oraz procedury Spółki}

\textbf{Metoda.} Dla każdego aktu podajemy: co się zmienia, które paragrafy umowy są dotknięte, co zmienia się w procedurach Spółki niezależnie od tekstu umowy, oraz czy zmieniać umowę teraz, czy po wejściu w życie. Nowelizacja KSH z 23 stycznia 2026~r. i ustawa z 13 marca 2026~r. zostały zweryfikowane na pełnych tekstach (\texttt{src/legal/}); akty prawa Unii zweryfikowano na tekstach urzędowych: 2021/697 (wersja PL) oraz 2026/877, 2024/1689, 2019/947 i 2026/1386 (wersje EN --- numery przepisów i tłumaczenia robocze z tekstu EN). Ustawę o kontroli niektórych inwestycji sprawdzono na tekście ujednoliconym (\texttt{src/legal/}). Niezweryfikowane pozostają TFUE, rozporządzenie 2019/452, 2023/1230 i 2024/2847 (\Wiar{W}).

\begin{longtable}{@{}p{3.1cm} p{3.7cm} p{3.3cm} p{4.4cm}@{}}
\toprule
\textbf{Akt} & \textbf{Co się zmienia} & \textbf{Dotknięte §} & \textbf{Wpływ na procedury; kiedy zmieniać umowę} \\
\midrule
\endfirsthead
\toprule
\textbf{Akt} & \textbf{Co się zmienia} & \textbf{Dotknięte §} & \textbf{Wpływ na procedury; kiedy zmieniać umowę} \\
\midrule
\endhead
Ustawa z 23.01.2026 r. o zmianie KSH (Dz.U. 2026 poz. 176), w życie 18.02.2027 r. (art.~28 i 33 od 28.02.2026 r.) \Wiar{Z} (pełny tekst w \texttt{src/legal/}) & Zmiany dotyczące P.S.A. ograniczają się do rejestru akcjonariuszy: art.~300$^{32}$ §~1$^{1}$--1$^{3}$ i §~3 (zarząd zgłasza umowę z podmiotem prowadzącym rejestr do sądu rejestrowego; podmiot zawiadamia sąd o jej wygaśnięciu w 7 dni), art.~300$^{33}$ §~1 pkt~5--6 i §~3 (w rejestrze numer PESEL albo data urodzenia, adres do doręczeń elektronicznych, adres e-mail „jeżeli akcjonariusz wyraził zgodę na komunikację przy wykorzystaniu poczty elektronicznej”, dane współwłaścicieli; zarząd zgłasza zmiany danych z pkt~1--4 i 9--11 w 7 dni), art.~300$^{34}$ (termin „tygodnia”; zgoda na wpis w formie pisemnej z podpisem notarialnie poświadczonym, w obecności osoby upoważnionej przez podmiot prowadzący rejestr albo elektronicznie z podpisem kwalifikowanym, zaufanym lub osobistym; powiadomienia elektroniczne), art.~300$^{35}$ §~1$^{1}$ (PESEL i adres zamieszkania nie są udostępniane pozostałym akcjonariuszom), art.~300$^{37}$ §~2 (rozszerzony katalog nabyć z mocy prawa poza wpisem konstytutywnym: spadek, zapis windykacyjny, wniesienie akcji jako wkładu, połączenie, podział, przekształcenie); ustawa o KRS art.~38 (dane podmiotu prowadzącego rejestr w KRS) i art.~25da; przepisy przejściowe art.~29 (3 miesiące na zgłoszenie podmiotu prowadzącego rejestr), art.~31 (dotychczasowe przepisy o akcjach imiennych i na okaziciela do akcji niezarejestrowanych), art.~34 (spółki istniejące dostosowują umowy w 2 lata). Nie zmienia art.~300$^{39}$, 300$^{41}$, 300$^{45}$--300$^{47}$, 300$^{52}$--300$^{98}$ ani 300$^{103}$--300$^{107}$ & §~\ref{art:wz} ust.~3 (zawiadomienia wyłącznie na e-mail z rejestru --- po nowelizacji e-mail figuruje w rejestrze tylko za zgodą akcjonariusza, więc obowiązek z umowy trzeba sformułować jako obowiązek wyrażenia zgody i wpisania adresu), §~\ref{art:zwrotne-zbycie} ust.~13 i umowa wykonawcza (forma zgody na wpis z art.~300$^{34}$ §~3 wyznacza formę pełnomocnictwa i dyspozycji), §~\ref{art:dziedziczenie} (spadkobierca nabywa akcje poza wpisem --- potwierdzone), §~\ref{art:final} & Żadne postanowienie nie jest sprzeczne; mechanizmy obrotu, dziedziczenia i emisji działają w obu stanach prawnych. Procedury: wybór i zmiana podmiotu prowadzącego rejestr zgłaszane do KRS z oświadczeniem zarządu; aktualizacja danych rejestru w 7 dni; dokumenty rejestru w formie z art.~300$^{34}$ §~3. Zalecane już teraz: klauzula dostosowawcza (poniżej), zgoda na komunikację e-mail w §~\ref{art:wz} ust.~3, forma pełnomocnictw w umowie wykonawczej. Zawiązanie bez oczekiwania na 18.02.2027 r.; spółka zawiązana wcześniej ma 2 lata na dostosowanie umowy i 3 miesiące na zgłoszenie podmiotu prowadzącego rejestr. \\
Ustawa z 13.03.2026 r. (Dz.U. 2026 poz. 471), art.~1 w mocy od 22.04.2026 r. (art.~1 pkt~17 lit.~b od 8.10.2026 r.), art.~2 od 8.04.2026 r. \Wiar{Z} (pełny tekst) & W ustawie o obrocie strategicznym: definicje przez odesłanie do rozporządzenia 2021/821, krajowe zezwolenie generalne, zezwolenia na wywóz, transfer, pomoc techniczną i pośrednictwo, osoba odpowiedzialna za koordynację kontroli obrotu (art.~9 ust.~2 pkt~11), elektroniczny rejestr zezwoleń (art.~24a--24d), ważność zezwoleń (uzbrojenie do 1 roku/3 lat, dual-use do 2/4 lat), sankcja za fałszywe dane we wniosku (art.~33 ust.~2a). W ustawie koncesyjnej z 13.06.2019 r. zmienia wyłącznie odnośnik do tytułu, art.~44 ust.~2aa (minimalna głębokość oznakowania broni palnej 0,0762 mm) i odesłanie w art.~56 ust.~4; definicji, zakresu koncesji (art.~7), wymogów wobec organów i akcjonariuszy (art.~10, 17) ani sankcji (art.~133) nie dotyka. & §~\ref{art:cel} ust.~4, §~\ref{art:pkd} ust.~8, §~\ref{art:export} & Umowa odsyła do ustaw ogólnie, bez numerów Dz.U., więc tekstu zmieniać nie trzeba; zmienia się program zgodności eksportowej (§~\ref{art:export} ust.~4) i bramki G2--G5 z \texttt{regulations.md}: wyznaczenie osoby odpowiedzialnej za koordynację kontroli obrotu, wnioski przez rejestr elektroniczny, oświadczenia o oryginałach dokumentów pod rygorem odpowiedzialności karnej. Szczegóły w odpowiedzi 4.1.1. \\
Rozporządzenie RM z 18.12.2024 r. w sprawie PKD 2025 (Dz.U. 2024 poz. 1936) \Wiar{Z} & Nowa klasyfikacja obowiązująca od 1.01.2025 r.; okres przejściowy dla podmiotów już wpisanych & §~\ref{art:pkd} ust.~5--6 & Spółka nowa, więc wpis od razu w PKD 2025; kontrola brzmienia 19 kodów --- sekcja T6. \\
Rozporządzenie Komisji (UE) 2026/877 z 16.04.2026~r. o porozumieniach o transferze technologii: w życie 1.05.2026~r., wygasa 30.04.2038~r. (art.~11); zastępuje 316/2014, które wygasło 30.04.2026~r. \Wiar{Z} (tekst EN) & Konstrukcja jak w 316/2014: wyłączenie (art.~2) dla licencji praw do technologii --- w tym praw autorskich do oprogramowania (art.~1 ust.~1 lit.~b pkt~vii) --- udzielanych między dwoma przedsiębiorstwami „w celu wytwarzania produktów objętych umową” (art.~1 ust.~1 lit.~c pkt~i, tłum. robocze); progi 20\% łącznie (konkurenci) i 30\% każdej ze stron (niekonkurenci) z trzyletnią karencją po ich przekroczeniu (art.~3, art.~8 lit.~e); najcięższe ograniczenia wyłączające całą umowę (art.~4, motyw~17): między konkurentami m.in. ustalanie cen, ograniczanie produkcji, podział rynków i klientów oraz ograniczanie eksploatacji własnej technologii licencjobiorcy i B+R stron (ust.~1), między niekonkurentami m.in. ustalanie cen sprzedaży i ograniczenia sprzedaży biernej (ust.~2); ograniczenia wyłączone (art.~5): wyłączny grant-back albo cesja własnych ulepszeń i nowych zastosowań licencjobiorcy (ust.~1 lit.~a), zakaz kwestionowania ważności praw (ust.~1 lit.~b), a między niekonkurentami --- ograniczenie eksploatacji własnej technologii licencjobiorcy i własnych B+R (ust.~2); okres przejściowy 1.05.2026--30.04.2027~r. tylko dla umów obowiązujących 30.04.2026~r. (art.~10); nie stosuje się do licencji w umowach B+R i specjalizacyjnych z 2023/1066 i 2023/1067 (art.~9) & §~\ref{art:przedsiewziecia} ust.~3 (mandat licencyjny); wzór umowy JV (T8) & Każdy bezpośredni lub pośredni obowiązek licencjobiorcy udzielenia licencji wyłącznej na własne ulepszenia lub nowe zastosowania albo przeniesienia praw do nich (w całości lub w części) wypada spod wyłączenia (tylko ta klauzula, motyw~18) i wymaga oceny indywidualnej z art.~101 ust.~3 TFUE; art.~5 ust.~1 lit.~a nie rozróżnia ulepszeń wydzielnych i niewydzielnych, więc wariant z odpowiedzi \hyperref[memo:2.4.1]{2.4.1} (własność ulepszeń niewydzielnych, licencja niewyłączna z prawem pierwszeństwa dla pozostałych) w części dotyczącej przeniesienia ulepszeń niewydzielnych także nie korzysta z wyłączenia --- stąd konieczna dokumentowana ocena Rady. Okres przejściowy nie dotyczy Spółki (umowy przyszłe). Procedury: przy każdej licencji dla Przedsięwzięcia ocena udziałów rynkowych (art.~8) i kontrola klauzul z art.~4--5, w tym brak bezwzględnej klauzuli no-challenge. Zmiana umowy spółki zalecana przed pierwszą umową Przedsięwzięcia. \\
AI Act --- rozporządzenie (UE) 2024/1689 \Wiar{Z} (tekst EN): stosowane od 2.08.2026~r.; rozdziały I--II (przepisy ogólne, praktyki zakazane) od 2.02.2025~r.; rozdział~III sekcja~4 (organy notyfikujące), rozdział~V (modele ogólnego przeznaczenia), VII (zarządzanie), XII (kary, z wyjątkiem art.~101) i art.~78 od 2.08.2025~r.; art.~6 ust.~1 (systemy wysokiego ryzyka jako element bezpieczeństwa produktów z załącznika~I) od 2.08.2027~r. (art.~113); w pliku tekst pierwotny z 12.07.2024~r. --- późniejsze zmiany terminów dla systemów wysokiego ryzyka \Wiar{?} & Wyłączenie z art.~2 ust.~3 akapit drugi: rozporządzenia nie stosuje się do systemów AI wprowadzanych do obrotu, oddawanych do użytku lub używanych „wyłącznie do celów wojskowych, obronnych lub bezpieczeństwa narodowego” (tłum. robocze) --- słowo „wyłącznie” oznacza, że produkt podwójnego zastosowania oferowany także na rynek cywilny jest objęty; wyłączone są ponadto systemy i modele AI specjalnie rozwijane i oddawane do użytku wyłącznie do celów badań naukowych i rozwojowych (ust.~6) oraz badania, testy i rozwój przed wprowadzeniem do obrotu, z wyjątkiem testów w warunkach rzeczywistych (ust.~8). Oprogramowanie autonomii jest „systemem AI” tylko wtedy, gdy spełnia całą definicję z art.~3 pkt~1: system „zaprojektowany do działania z różnym poziomem autonomii”, który „wnioskuje” z otrzymanych danych wejściowych, jak generować wyniki (predykcje, rekomendacje, decyzje) mogące wpływać na środowisko fizyczne lub wirtualne (tłum. robocze) --- sama autonomia nie wystarcza; moduły uczone (percepcja, planowanie) zwykle spełniają definicję, czysto regułowe sterowanie deterministyczne może jej nie spełniać \Wiar{W}. Cywilny BSP: załącznik~I sekcja~B wymienia rozporządzenie 2018/1139 w zakresie bezzałogowych statków powietrznych, więc autonomia jako element bezpieczeństwa BSP podlegającego ocenie zgodności przez stronę trzecią jest systemem wysokiego ryzyka z art.~6 ust.~1, ale wobec produktów z sekcji~B stosuje się tylko art.~6 ust.~1, 102--109 i 112 (art.~2 ust.~2), a wymogi wejdą przez akty wykonawcze i delegowane do 2018/1139 zmienionego przez art.~108 AI Act; załącznik~III (m.in. biometria, elementy bezpieczeństwa infrastruktury krytycznej i ruchu drogowego, zatrudnienie, egzekwowanie prawa, kontrola graniczna) nie wymienia BSP jako takich --- kwalifikacja z art.~6 ust.~2 zależy od zastosowania, nie od platformy & §~\ref{art:rada-zastrzezone} ust.~1 lit.~i (polityki), §~\ref{art:export}, §~\ref{art:przedsiewziecia} (licencje dla zastosowań cywilnych) & Umowy nie zmieniać. Procedury: kwalifikacja każdego produktu cywilnego (system AI? element bezpieczeństwa BSP certyfikowanego?) w rejestrze obowiązków produktowych regulaminu Rady (bramka G1); testy w warunkach rzeczywistych poza wyłączeniem z art.~2 ust.~8; dokumentowanie „wyłącznie wojskowego” przeznaczenia wersji obronnych. \\
Rozporządzenie maszynowe (2023/1230, stosowane od 2027~r.) i Cyber Resilience Act (2024/2847) \Wiar{W} & Obowiązki produktowe i cyberbezpieczeństwa dla cywilnych systemów autonomicznych & §~\ref{art:rada-zastrzezone} ust.~1 lit.~i & Poza umową: ten sam rejestr obowiązków produktowych (bramka G1). \\
Rozporządzenie wykonawcze (UE) 2019/947 w sprawie przepisów i procedur operacji BSP \Wiar{Z} (tekst EN; w pliku tekst pierwotny, stosowany od 1.07.2020~r. --- art.~23 ust.~1; późniejsze zmiany terminów i scenariuszy standardowych \Wiar{W}) & Trzy kategorie operacji (art.~3): „otwarta” bez zezwolenia (art.~4: masa poniżej 25~kg, VLOS, do 120~m nad ziemią), „szczególna” z zezwoleniem operacyjnym właściwego organu po ocenie ryzyka albo z oświadczeniem w scenariuszu standardowym (art.~5, 11--12), „certyfikowana” z certyfikacją BSP według 2019/945 i operatora (art.~3 lit.~c, art.~6); rejestracja operatora w państwie głównego miejsca prowadzenia działalności, obowiązkowa w kategorii szczególnej dla BSP każdej masy (art.~14 ust.~5--6); definicje z 2018/1139 (art.~2) & brak (działalność operacyjna, nie umowa) & Loty testowe autonomii poza kategorią otwartą (BVLOS, masa od 25~kg, lot autonomiczny --- w kategorii otwartej pilot utrzymuje VLOS i zdolność sterowania BSP, pkt UAS.OPEN.060 ppkt~2 lit.~b i d załącznika) wymagają zezwolenia operacyjnego Prezesa ULC (organ właściwy do zadań z art.~18 rozporządzenia --- art.~21 ust.~2g prawa lotniczego \Wiar{Z}) w kategorii szczególnej, oświadczenia w scenariuszu standardowym albo certyfikatu LUC (art.~5 ust.~4 lit.~b i ust.~6 lit.~a); rejestracja Spółki jako operatora przed pierwszym lotem; wyłączenie lotnictwa państwowego i wojskowego wynika z 2018/1139 \Wiar{W}. Wpis do rejestru obowiązków (bramka G1); zgodność potrzebna także partnerom Przedsięwzięć cywilnych. \\
Rozporządzenie (UE) 2019/452 o monitorowaniu BIZ (tekst niedostępny \Wiar{W}; uchylane z dniem 17.01.2028~r. przez art.~30 rozporządzenia 2026/1386, ze stosowaniem nadal do inwestycji będących w toku kontroli lub zakończonych do tego dnia --- art.~30 ust.~1--2 \Wiar{Z} (tekst EN)) i ustawa o kontroli niektórych inwestycji \Wiar{Z} (tekst ujednolicony na podstawie t.j. Dz.U. 2026 poz.~47 ze zm.) & Stan do 16.01.2028~r.: ramy unijne bez obowiązku posiadania mechanizmu krajowego (mechanizm współpracy) \Wiar{W}; w Polsce dwa reżimy: (a)~podmioty z wykazu określonego rozporządzeniem Rady Ministrów (art.~4) --- kontrola nabycia przez każdego nabywcę; (b)~art.~12a--12k --- kontrola nabycia znaczącego uczestnictwa (m.in. 20\% ogólnej liczby głosów) albo dominacji przez inwestora spoza UE/EOG/OECD (podmiot inny niż osoba fizyczna --- bez siedziby w takim państwie co najmniej od dwóch lat; art.~12a ust.~1, art.~12c ust.~1 pkt~1 i 5) w podmiocie objętym ochroną z art.~12d, m.in. prowadzącym działalność, której przedmiotem jest wytwarzanie i obrót wyrobami i technologią o przeznaczeniu wojskowym lub policyjnym (ust.~3 pkt~7), ale tylko przy przychodzie w Polsce powyżej równowartości 10~mln EUR w jednym z dwóch lat obrotowych przed zgłoszeniem (ust.~4); towarów podwójnego zastosowania ustawa w katalogu nie wymienia & §~\ref{art:zbywanie} ust.~2 lit.~e, §~\ref{art:kwalifikowana-runda} & Umowa prawidłowo przewiduje wstrzymanie terminu do czasu zgody organu; w rundzie z inwestorem spoza UE warunek zawieszający w umowie inwestycyjnej (pytanie \hyperref[memo:1.1.2]{1.1.2}). Od 17.01.2028~r. --- wiersz następny. \\
Rozporządzenie (UE) 2026/1386 z 17.06.2026~r. w sprawie monitorowania inwestycji zagranicznych w Unii \Wiar{Z} (tekst EN): w życie 16.07.2026~r., stosowane od 17.01.2028~r.; już od 16.07.2026~r. art.~3 ust.~2 (notyfikacja krajowych mechanizmów do 17.01.2028~r.), art.~15 ust.~2 (formularz), art.~18 ust.~1--6 i 9--12 (system wymiany, portal, baza) oraz art.~27--29 (akty delegowane, komitet) (art.~31) & Wobec 2019/452: (1)~każde państwo członkowskie musi ustanowić mechanizm kontroli (art.~3 ust.~1); (2)~wspólny minimalny zakres obowiązkowego uprzedniego zezwolenia (art.~4 ust.~15): podmioty, które rozwijają, wytwarzają lub komercjalizują pozycje z załącznika~I do 2021/821 (lit.~a) albo z wykazu wojskowego z dyrektywy 2009/43/WE (lit.~b), technologie półprzewodnikowe, kwantowe i AI z załącznika~I (lit.~c; w AI tylko badania i rozwój modeli AI ogólnego przeznaczenia lub opartych na nich systemów nadających się do rozwoju zastosowań kosmicznych lub obronnych albo modeli z ryzykiem systemowym --- załącznik~I pkt~3), infrastruktura transportowa, energetyczna i cyfrowa uznana za krytyczną, surowce strategiczne, infrastruktura finansowa i wyborcza (lit.~d--g); inwestycje greenfield są wyłączone z tego obowiązkowego zakresu (ust.~17); państwa mogą objąć kontrolą więcej (ust.~16); (3)~„inwestor zagraniczny” to osoba fizyczna bez obywatelstwa państwa członkowskiego albo podmiot utworzony według prawa państwa trzeciego (art.~2 pkt~5) --- bez wyjątku dla EOG, NATO ani OECD; „inwestycja zagraniczna” obejmuje inwestycję dokonaną przez unijną spółkę zależną kontrolowaną przez inwestora zagranicznego, jeżeli umożliwia skuteczny udział w zarządzaniu lub kontroli (art.~2 pkt~1 i~7); progów procentowych rozporządzenie nie ustala (określa je prawo krajowe; państwa publikują wytyczne o progach i przesłankach zgłoszenia --- art.~23 ust.~2); (4)~procedura: zgłoszenie i kontrola przed zamknięciem (art.~4 ust.~9), ocena wstępna 45~dni, potem postępowanie pogłębione (ust.~2), kontrola z urzędu inwestycji nieobjętych obowiązkiem uprzedniego zezwolenia przez co najmniej 15~miesięcy i najwyżej 5~lat po zamknięciu (ust.~4), a inwestycji objętych tym obowiązkiem, lecz niezgłoszonych albo zgłoszonych po zamknięciu --- przez co najmniej 24~miesiące (ust.~5), sankcje (ust.~11), środek zaskarżenia (ust.~7); (5)~mechanizm współpracy: obowiązkowa notyfikacja innym państwom i Komisji inwestycji z zakresu art.~4 ust.~15, gdy inwestor lub jego unijna spółka zależna jest bezpośrednio lub pośrednio kontrolowana przez rząd państwa trzeciego, gdy inwestor lub osoby z nim powiązane są objęte unijnymi środkami ograniczającymi albo gdy uczestniczył w inwestycji wcześniej niezatwierdzonej lub zatwierdzonej z warunkami, których istotnie lub wielokrotnie nie przestrzegano (art.~5 ust.~1), a przy postępowaniu pogłębionym --- gdy podmiot uczestniczy w projekcie lub programie z załącznika~II (m.in. Europejski Fundusz Obronny, EDIP, PESCO, ASAP, EDIRPA, Horyzont Europa) albo on lub jego grupa ma spółki zależne w innym państwie członkowskim (art.~5 ust.~2); terminy 15/45~dni (art.~6); uwagi państw i opinie Komisji (art.~8); w informacji podaje się udział w programach z załącznika~II i granty UE od 750~000~EUR z ostatnich 5~lat (art.~15 ust.~1 lit.~f--g); (6)~czynniki oceny (art.~19 ust.~1) obejmują wpływ na projekty i programy z załącznika~II (lit.~a) oraz na dostępność technologii krytycznych z załącznika~III (lit.~b), gdzie wprost wymieniono „drony i pojazdy (powietrzne, lądowe, nawodne i podwodne)” oraz „systemy wykorzystujące AI” (lit.~h), zaawansowane technologie naprowadzania, nawigacji i sterowania, w tym awionikę (lit.~b), oraz technologie lotnicze (lit.~f) (tłum. robocze) & §~\ref{art:dopuszczalny-nabywca} (Kryterium i wyjątek dla funduszy), §~\ref{art:zbywanie} ust.~2 lit.~e, §~\ref{art:finansowanie}--\ref{art:kwalifikowana-runda}, §~\ref{art:przedsiewziecia} ust.~4; umowa akcjonariuszy & Dwa stany prawne. Do 16.01.2028~r.: kontrola tylko, gdy Spółka jest w wykazie z rozporządzenia Rady Ministrów (art.~4 ustawy o kontroli niektórych inwestycji) albo --- wobec inwestora spoza UE/EOG/OECD --- podmiotem objętym ochroną z art.~12d tej ustawy, co przy działalności wojskowej wymaga także przychodu w Polsce powyżej równowartości 10~mln EUR (art.~12d ust.~3 pkt~7 i ust.~4) \Wiar{Z}; runda na wczesnym etapie z inwestorem spoza UE zwykle więc zgłoszenia nie wymaga. Od 17.01.2028~r.: gdy Spółka rozwija, wytwarza lub komercjalizuje pozycje z załącznika~I do 2021/821 albo z wykazu wojskowego (bramka G3, T1a pkt~4), każda inwestycja spoza UE dająca udział w zarządzaniu lub kontroli --- także funduszu z USA, UK, Szwajcarii czy Norwegii i funduszu unijnego kontrolowanego spoza UE --- wymaga uprzedniego zezwolenia; Kryterium (§~\ref{art:dopuszczalny-nabywca}) jest więc węższe niż zakres kontroli państwowej i jej nie zastępuje. Umowy spółki nie zmieniać: §~\ref{art:zbywanie} ust.~2 lit.~e już wstrzymuje termin do zgody organu; przy przywracaniu wyjątku dla funduszy (pakiet 1.4) i w §~\ref{art:kwalifikowana-runda} ust.~4 rozważyć warunek zawieszający „uzyskanie zezwoleń wymaganych przepisami o kontroli inwestycji” (opcjonalne). Umowa akcjonariuszy/inwestycyjna: oświadczenia inwestora o strukturze własności i beneficjencie rzeczywistym (art.~15 ust.~1 lit.~a--b; organ może żądać tych informacji od inwestora z 15-dniowym terminem, art.~15 ust.~3), obowiązek współdziałania przy zgłoszeniu, long-stop date obejmujący 45~dni oceny wstępnej i postępowanie pogłębione, klauzula na wypadek środków łagodzących. Udział w EDF (załącznik~II pkt~14) jest czynnikiem oceny (art.~19 ust.~1 lit.~a), przesłanką notyfikacji przy postępowaniu pogłębionym (art.~5 ust.~2 lit.~a) i podstawą opinii Komisji (art.~8 ust.~2 lit.~b). Kalendarz: śledzić projekt polskiej ustawy wdrażającej (notyfikacja do 17.01.2028~r.). \\
Ustawa o PIT --- programy motywacyjne (art.~24 ust.~11--12b) \Wiar{W} & Od 2022 r. przepis obejmuje także P.S.A.; brak utrwalonej praktyki interpretacyjnej & §~\ref{art:program-motywacyjny} & Interpretacja indywidualna przed pierwszym grantem (T9, pkt 27 zakresu). \\
\bottomrule
\end{longtable}

\textbf{Klauzula dostosowawcza (zalecane, §~\ref{art:final}).} Chroni umowę przed sporem o to, które z jej postanowień „zmieniła” nowelizacja, bez czekania na akt notarialny.
\begin{Brzmienie}
Jeżeli po zawarciu niniejszej umowy zmiana przepisów prawa spowoduje, że postanowienie niniejszej umowy stanie się sprzeczne z przepisem bezwzględnie obowiązującym, w miejsce tego postanowienia stosuje się przepis ustawy, a pozostałe postanowienia zachowują moc. Rada Dyrektorów przedstawia Walnemu Zgromadzeniu projekt dostosowania umowy na najbliższym Walnym Zgromadzeniu po wejściu zmiany w życie. Jeżeli zmiana przepisów rozszerza swobodę kształtowania umowy, dotychczasowe postanowienia stosuje się nadal do czasu ich zmiany.
\end{Brzmienie}
\Klasa{zalecane}

\textbf{Pytania odłożone przez Założycieli (T1a).} Cztery zagadnienia z listy zadań (\texttt{psa\_todo.md}, sekcja 1) dostają odpowiedź wstępną do rozwinięcia w rundzie uzupełniającej:
\begin{enumerate}
  \item \emph{Koncesja przy licencji dla spółki celowej.} Ustawa z 2019 r. obejmuje koncesją „obrót” wyrobami i technologią o przeznaczeniu wojskowym lub policyjnym; ustawa nie definiuje „wytwarzania” ani „obrotu”, a koncesja obejmuje wytwarzanie i obrót wyrobami oraz obrót technologią (art.~7 ust.~1) \Wiar{Z}; w piśmiennictwie i praktyce organu obrót technologią obejmuje udostępnianie, w tym licencjonowanie \Wiar{W}. Jeżeli oprogramowanie autonomii zostanie zakwalifikowane jako technologia z wykazu WT, licencja dla partnera produkującego wojskowe BSP wymaga koncesji po stronie Spółki niezależnie od koncesji producenta. Do czasu kwalifikacji produktu (bramka G3) §~\ref{art:przedsiewziecia} ust.~3 nie powinien być wykonywany dla zastosowań wojskowych.
  \item \emph{Wyłączenia B+R w ustawie z 2019 r.} Ustawa nie zna ogólnego wyłączenia dla prac badawczo-rozwojowych (art.~4 wymienia tylko wyłączenia podmiotowe i przedmiotowe niezwiązane z B+R) \Wiar{Z}; kluczowe jest, czy fizycznie powstaje wyrób z wykazu (wytwarzanie) albo czy technologia jest udostępniana (obrót). W wykazie WT bezzałogowe statki powietrzne mieszczą się w WT~V ust.~3 „oraz ich systemy i urządzenia, takie jak: startu, lądowania, kierowania i kontroli lotu”, a wyłączenie dla statków dopuszczonych do użytku cywilnego dotyczy tylko WT~V ust.~2 (statki załogowe) \Wiar{Z}; oprogramowanie i technologię obejmują definicje z pkt~12 i 15 załącznika oraz WT~XI ust.~2 i WT~XIII ust.~3--4 \Wiar{Z}. Prace nad algorytmami na platformie cywilnej i symulacje nie są objęte; prototyp „specjalnie zaprojektowany lub zmodyfikowany do celów wojskowych” jest. Ustawa z 13.03.2026 r. niczego tu nie zmieniła (\hyperref[memo:4.1.1]{4.1.1}).
  \item \emph{Kryterium a art.~9 EDF.} Kryterium dopuszcza kontrolę przez podmioty z NATO spoza UE; art.~9 ust.~3 rozporządzenia 2021/697 stanowi, że odbiorcy i podwykonawcy „nie podlegają kontroli niestowarzyszonego państwa trzeciego lub podmiotu z niestowarzyszonego państwa trzeciego”, przy czym „kontrola” to „zdolność do wywierania decydującego wpływu na podmiot prawny, bezpośrednio lub pośrednio” (art.~2 pkt~6), a ust.~4 dopuszcza odstępstwo dla podmiotu z siedzibą w UE lub państwie stowarzyszonym tylko przy gwarancjach zatwierdzonych przez państwo członkowskie, obejmujących m.in. zachowanie przez odbiorcę własności praw własności intelektualnej z działania i brak ich kontroli lub ograniczenia przez to państwo (lit.~c); państwami stowarzyszonymi są wyłącznie członkowie EFTA należący do EOG (art.~5), więc kontrola z USA, UK, Kanady lub Turcji --- dopuszczona przez Kryterium --- wyklucza kwalifikowalność bez gwarancji; zmianę struktury w toku działania zgłasza się Komisji, która ocenia dalszą kwalifikowalność (ust.~7) \Wiar{Z} (tekst PL); czy weto akcjonariusza w sprawach strategicznych jest „decydującym wpływem”, tekst nie przesądza \Wiar{W}. Zalecamy nie zawężać Kryterium w umowie spółki (byłoby to ograniczenie także dla inwestorów, których EDF nie dotyczy), lecz wpisać do umowy akcjonariuszy zobowiązanie, że dopóki Spółka uczestniczy w EDF, AGILE lub EUDIS, akcjonariusze spoza UE/EOG nie uzyskują kontroli ani weta (decyzja D6; sekcja T8).
  \item \emph{Klasyfikacja dual-use produktów cywilnych.} To zadanie techniczno-prawne, nie korporacyjne: klasyfikację według aktualnego załącznika I do rozporządzenia 2021/821 (pozycje 9A012 --- bezzałogowe statki powietrzne, 9D/9E, kategoria~7 --- nawigacja i awionika, 5A002; pozycje te figurują w załączniku~I w wersji skonsolidowanej na 26.05.2023~r. \Wiar{Z} (tekst EN), z zastrzeżeniem późniejszych zmian załącznika) wykonuje specjalista od kontroli eksportu na konkretnej specyfikacji produktu, z zapisem wyniku i wersji wykazu; przy wątpliwości wniosek do ministra właściwego do spraw gospodarki o ustalenie, czy pozycja podlega kontroli \Wiar{W}. Umowa (§~\ref{art:export} ust.~2) wymaga tej klasyfikacji przed pierwszym udostępnieniem. Ta sama klasyfikacja przesądza od 17.01.2028~r., czy Spółka mieści się w obowiązkowym zakresie kontroli inwestycji zagranicznych (art.~4 ust.~15 lit.~a rozporządzenia 2026/1386: podmioty rozwijające, wytwarzające lub komercjalizujące pozycje z załącznika~I do 2021/821; lit.~b --- towary i technologie z wykazu uzbrojenia z dyrektywy 2009/43/WE) \Wiar{Z} (tekst EN) --- wynik klasyfikacji należy więc włączyć do dokumentacji każdej rundy z inwestorem spoza UE.
\end{enumerate}
Pytanie o jurysdykcję i strukturę holdingową (decyzja D1) wykracza poza zakres memorandum; sygnalizujemy tylko, że mechanizmy specyficzne dla P.S.A. (upoważnienie emisyjne z art.~300$^{103}$ KSH, obowiązki związane z akcją w rejestrze, ograniczenie wstąpienia spadkobierców z art.~300$^{41}$) nie mają prostych odpowiedników w sp.~z~o.o. ani w spółkach obcych, a przeniesienie własności intelektualnej do zagranicznej spółki-matki po jej wniesieniu do polskiej spółki rodzi skutki podatkowe i eksportowe.

\section*{T2. Wykonalność mechanizmów wobec rejestru akcjonariuszy, sądu rejestrowego i w sporach}
\addcontentsline{toc}{section}{T2. Wykonalność mechanizmów wobec rejestru, sądu rejestrowego i w sporach}

\textbf{Trzy fora.} Podmiot prowadzący rejestr akcjonariuszy dokonuje wpisu na żądanie spółki albo osoby mającej interes prawny, po sprawdzeniu treści i formy dokumentów uzasadniających wpis; nie rozstrzyga sporów o skuteczność czynności, ale może odmówić wpisu sprzecznego z dokumentami albo z ujawnionym ograniczeniem (art.~300$^{34}$--300$^{35}$ KSH \Wiar{W}). Sąd rejestrowy bada zgodność umowy z prawem przy wpisie spółki i przy każdej zmianie umowy (art.~23 ustawy o KRS \Wiar{W}); nie bada umowy akcjonariuszy ani umów wykonawczych. Spór między akcjonariuszami toczy się przed sądem powszechnym (§~\ref{art:final} ust.~4) i tam decyduje, czy obowiązek jest dostatecznie oznaczony, by zastąpić oświadczenie woli wyrokiem (art.~64 KC, art.~1047 KPC).

\begin{longtable}{@{}p{2.8cm} p{3.8cm} p{3.4cm} p{4.5cm}@{}}
\toprule
\textbf{Mechanizm} & \textbf{Rejestr akcjonariuszy} & \textbf{Sąd rejestrowy} & \textbf{Spór między akcjonariuszami} \\
\midrule
\endfirsthead
\toprule
\textbf{Mechanizm} & \textbf{Rejestr akcjonariuszy} & \textbf{Sąd rejestrowy} & \textbf{Spór między akcjonariuszami} \\
\midrule
\endhead
Zgoda Spółki na rozporządzenie akcją (§~\ref{art:zbywanie}) & Ograniczenie z umowy spółki ujawnia się w rejestrze (art.~300$^{33}$ §~1 pkt~10 \Wiar{W}); przy żądaniu wpisu nabywcy podmiot prowadzący rejestr żąda dowodu zgody Spółki albo dowodu, że zgoda nie jest wymagana; zbycie bez zgody jest bezskuteczne wobec Spółki (art.~300$^{39}$ §~1 w zw. z art.~63 KC \Wiar{W}) & Ograniczenie zbywalności jest wprost dopuszczone ustawą; ryzyko wezwania tylko co do derogacji art.~300$^{39}$ §~3--6 (pytania \hyperref[memo:1.2.1]{1.2.1}, \hyperref[memo:1.2.6]{1.2.6}) & Zbywca kwestionujący odmowę: powództwo o ustalenie bezskuteczności odmowy albo o zobowiązanie do wskazania nabywcy; zamknięty katalog przyczyn odmowy (§~\ref{art:zbywanie} ust.~2) ułatwia Spółce obronę \\
Kryterium EU/NATO (§~\ref{art:dopuszczalny-nabywca}) & Rejestr nie bada obywatelstwa; działa przez wymóg zgody Spółki. Skuteczność wobec nabywcy w drodze egzekucji, dziedziczenia lub podziału majątku ograniczona (pytania \hyperref[memo:1.1.6]{1.1.6}, \hyperref[memo:1.6.1]{1.6.1}, \hyperref[memo:1.7.1]{1.7.1}) & Najwyższe ryzyko pytań sądu (T6); argumentacja: ograniczenie przez zgodę spółki z obiektywnym kryterium, zgodne z art.~300$^{39}$ §~1, uzasadnione art.~65 i 346 TFUE \Wiar{W} oraz wymogami programów obronnych (art.~9 ust.~3 rozporządzenia 2021/697: odbiorcy „nie podlegają kontroli niestowarzyszonego państwa trzeciego” \Wiar{Z}) & Nabywca spoza Kryterium, który nabył akcje bez zgody: bezskuteczność wobec Spółki; roszczenia Spółki z §~\ref{art:zbywanie} ust.~8 i §~\ref{art:breach} wobec zbywcy \\
Lock-up (§~\ref{art:okres-ograniczenia-zbywania}) & Ujawnia się jak ograniczenie zbywalności z terminem & Derogacja art.~300$^{39}$ §~3--6 na 12 miesięcy: umiarkowane ryzyko pytania; obrona: „chyba że umowa spółki stanowi inaczej” w §~2 (pytanie \hyperref[memo:1.2.6]{1.2.6}) & Jak wyżej \\
Zwrotne zbycie (§~\ref{art:zwrotne-zbycie}) & Obowiązek związany z akcją ujawnia się w rejestrze (art.~300$^{33}$ §~1 pkt~11 \Wiar{W}); wpis nabywcy wskazanego przez Spółkę na podstawie umowy zbycia (oferta + przyjęcie) albo wyroku; rejestr nie wykona „automatycznego” przeniesienia bez dokumentu & Sąd nie bada umowy wykonawczej; w umowie spółki sam obowiązek zbycia jest dopuszczalny (art.~300$^{1}$ §~3, art.~300$^{33}$) & Realna droga to pełnomocnictwo (art.~101 §~2 KC) i oferta nieodwołalna; w razie sporu powództwo z art.~64 KC wymaga oznaczonego nabywcy, ceny i akcji --- dlatego umowa wykonawcza (pytanie \hyperref[memo:1.5.3]{1.5.3}) \\
Prawo pierwszeństwa, tag, drag (§~\ref{art:pierwszenstwo}--\ref{art:przymusowe-wspolzbycie}) & Ograniczenia ujawniane; przy drag rejestr wpisuje nabywcę na podstawie umów zbycia albo wyroków wobec opornych & Standardowe, niskie ryzyko & Drag wobec opornego akcjonariusza: art.~64 KC; skuteczniejsze są pełnomocnictwa w umowie akcjonariuszy (T8) \\
Dziedziczenie (§~\ref{art:dziedziczenie}) & Spadkobierca żąda wpisu na podstawie aktu poświadczenia dziedziczenia; rejestr wpisuje go jako akcjonariusza, chyba że umowa ogranicza wstąpienie --- wtedy wpis „z ograniczeniem” albo odmowa do czasu spłaty \Wiar{?}; praktyka podmiotów prowadzących rejestr niejednolita & Ograniczenie z art.~300$^{41}$ §~1 wymaga określenia warunków spłaty; sąd sprawdzi, czy są (pytanie \hyperref[memo:1.6.1]{1.6.1}) & Spadkobierca niewspółdziałający: Spółka wykonuje spłatę do depozytu i żąda wpisu nabywcy wskazanego; spór o Wartość Godziwą rozstrzyga ekspert (§~\ref{art:wycena}) \\
Utrata Kryterium po nabyciu (§~\ref{art:dopuszczalny-nabywca} ust.~15) & Brak ujawnienia zdarzenia; rejestr działa dopiero przy zbyciu & Ryzyko zarzutu „przymusowego wykupu poza ustawą” (pytanie \hyperref[memo:1.1.4]{1.1.4}); obrona: obowiązek związany z akcją przyjęty przy nabyciu, cena równa Wartości Godziwej & Egzekucja jak przy zwrotnym zbyciu; wobec funduszy konieczne wyłączenie zdarzeń po stronie ich inwestorów (T7) \\
\bottomrule
\end{longtable}

\textbf{Wniosek.} Wszystkie mechanizmy są wykonalne pod warunkiem trzech rzeczy: (1) ujawnienia ograniczeń i obowiązków w rejestrze przy pierwszym wpisie akcji (dyspozycja Spółki z wykazem, T5); (2) istnienia w dniu zawiązania umowy wykonawczej z ofertami i pełnomocnictwami (pytanie \hyperref[memo:1.5.3]{1.5.3}); (3) wyboru podmiotu prowadzącego rejestr, który potwierdzi na piśmie, jak obsłuży wpisy warunkowe i dyspozycje pełnomocnika (T5).

\section*{T3. Ścieżka zawiązania poza S24 i lista dokumentów do KRS}
\addcontentsline{toc}{section}{T3. Ścieżka zawiązania poza S24 i lista dokumentów do KRS}

\textbf{Forma.} Umowa P.S.A. zawierana poza wzorcem S24 wymaga formy aktu notarialnego (art.~300$^{6}$ §~1 KSH \Wiar{W}); wkłady niepieniężne wykluczają S24 niezależnie od treści umowy. Z chwilą zawarcia umowy powstaje spółka w organizacji (art.~300$^{11}$ \Wiar{W}). Wniosek o wpis składa się wyłącznie elektronicznie przez Portal Rejestrów Sądowych, podpisany przez wszystkich dyrektorów (podpis kwalifikowany, zaufany albo osobisty), w terminie sześciu miesięcy od zawarcia umowy \Wiar{W}.

\textbf{Kolejność czynności.}
\begin{enumerate}
  \item Uzupełnienie Załączników nr~1--3 i pól umowy (58 pól; T4, decyzje D8, D9); wybór podmiotu prowadzącego rejestr (T5, D10) i podpisanie z nim umowy albo listu intencyjnego.
  \item Umowa wykonawcza zwrotnego zbycia i umowa akcjonariuszy podpisane najpóźniej w dniu aktu (§~\ref{art:zwrotne-zbycie} ust.~13); umowy przeniesienia praw własności intelektualnej podpisane w dniu aktu ze skutkiem dla spółki w organizacji (§~\ref{art:wklady} ust.~4).
  \item Akt notarialny: umowa spółki z Załącznikami; w tym samym akcie albo osobnej uchwale założycieli: powołanie pierwszej Rady Dyrektorów (jeżeli umowa nie wskazuje dyrektorów w treści) \Wiar{W}.
  \item Wniesienie wkładów pieniężnych przed złożeniem wniosku (§~\ref{art:wklady} ust.~4; ustawowo wystarczy 1~zł na kapitał akcyjny, art.~300$^{12}$ \Wiar{W}); wkłady niepieniężne w terminie z umowy, nie później niż 3~lata od wpisu (art.~300$^{9}$ \Wiar{W}) --- umowa skraca do 14~dni.
  \item Wniosek przez PRS z załącznikami (tabela); opłata sądowa i za ogłoszenie w MSiG \Wiar{W}.
  \item Po wpisie: dyspozycja do podmiotu prowadzącego rejestr (wpis akcji serii A, akcjonariuszy, ograniczeń i obowiązków związanych z akcją, adresów e-mail); zgłoszenie beneficjentów rzeczywistych do CRBR w 14~dni \Wiar{W}; NIP-8, VAT-R, zgłoszenie do ZUS; protokoły wydania wkładów niepieniężnych.
\end{enumerate}

\begin{longtable}{@{}p{5.6cm} p{4.2cm} p{4.4cm}@{}}
\toprule
\textbf{Dokument do wniosku o wpis} & \textbf{Podstawa} & \textbf{Uwagi} \\
\midrule
\endfirsthead
\toprule
\textbf{Dokument do wniosku o wpis} & \textbf{Podstawa} & \textbf{Uwagi} \\
\midrule
\endhead
Umowa spółki (wypis aktu notarialnego; notariusz przekazuje elektronicznie do CREWAN, we wniosku podaje się numer) & art.~300$^{12}$ §~3 pkt~1 KSH \Wiar{Z}, art.~19 ustawy o KRS \Wiar{Z} & Załączniki nr~1--3 jako część aktu \\
Oświadczenie wszystkich dyrektorów o wysokości kapitału akcyjnego (suma wartości wkładów przeznaczonych na kapitał) oraz o wniesieniu wkładów w części przewidzianej umową & art.~300$^{12}$ §~3 pkt~2--3 KSH \Wiar{Z} & kapitał akcyjny = suma wkładów pieniężnych z Załącznika nr~1 \\
Lista akcjonariuszy z nazwiskiem i imieniem albo firmą oraz liczbą i serią akcji objętych przez każdego, podpisana przez wszystkich dyrektorów & art.~300$^{12}$ §~4 KSH \Wiar{Z} & zgodna z Załącznikiem nr~1 \\
Dowód powołania dyrektorów (uchwała założycieli albo postanowienie umowy), zgody na powołanie, adresy do doręczeń dyrektorów, oświadczenia o niekaralności w zakresie art.~18 KSH & art.~300$^{12}$ §~3 pkt~4--5 KSH, art.~19a ust.~5 i 5a ustawy o KRS (pełnomocnik do doręczeń, gdy adres dyrektora poza UE), art.~18 KSH \Wiar{Z} & Kryterium wobec dyrektorów (§~\ref{art:rada} ust.~3) sąd nie bada \\
Oświadczenie o adresie Spółki (jeżeli inny niż w umowie) i umowa najmu albo tytuł do lokalu (na żądanie) & praktyka sądów \Wiar{W} & --- \\
Wskazanie podmiotu prowadzącego rejestr akcjonariuszy & do 17.02.2027 r. nie jest załącznikiem do wniosku: umowę o prowadzenie rejestru spółka zawiera niezwłocznie po wpisie (art.~300$^{32}$ §~1 KSH) \Wiar{Z}; od 18.02.2027 r. Rada zgłasza zawarcie umowy do sądu z oświadczeniem o jej zawarciu (art.~300$^{32}$ §~1$^{1}$--1$^{3}$ KSH, art.~38 pkt~8a lit.~j ustawy o KRS) \Wiar{Z} & wybór podmiotu w umowie albo uchwale (art.~300$^{31}$ §~5) \\
Dowody wpłat wkładów pieniężnych & praktyka; na żądanie sądu \Wiar{W} & rachunek spółki w organizacji \\
Pełnomocnictwo, jeżeli wniosek składa pełnomocnik procesowy & KPC \Wiar{W} & --- \\
Opłaty: wpis 500~zł (art.~52 ust.~1 ustawy o kosztach sądowych w sprawach cywilnych) \Wiar{Z}, ogłoszenie w MSiG 100~zł (rozporządzenie) \Wiar{W} & ustawa o kosztach sądowych; zmiana wpisu 250~zł (art.~55 ust.~1) \Wiar{Z} & kwoty do sprawdzenia w dniu wniosku \\
\bottomrule
\end{longtable}

\textbf{Czego nie składa się do KRS.} Umowy wykonawczej, umowy akcjonariuszy, umów przeniesienia praw, regulaminów. Sąd rejestrowy widzi tylko umowę spółki; to argument, by mechanizmy sporne dla sądu (T6) trzymać w umowie spółki tylko w zakresie niezbędnym dla skuteczności wobec osób trzecich, a resztę w umowie akcjonariuszy (T8).

\section*{T4. Dokumentacja wkładów niepieniężnych i przeniesienia praw}
\addcontentsline{toc}{section}{T4. Dokumentacja wkładów niepieniężnych i przeniesienia praw}

\textbf{Zasada.} Umowa spółki określa przedmiot wkładu, osobę wnoszącą i liczbę akcji obejmowanych za wkład (art.~300$^{5}$ §~1 pkt~4--5 KSH \Wiar{W}); w P.S.A. nie ma obowiązkowego badania wkładów przez biegłego rewidenta, a odpowiedzialność za zawyżenie dotyczy wkładów przeznaczonych na kapitał akcyjny (art.~300$^{10}$ \Wiar{W}). Odpowiedź \hyperref[memo:3.1.2]{3.1.2} wykazuje, że obecne §~\ref{art:wklady} ust.~4 (na kapitał akcyjny wyłącznie wkłady pieniężne) jest sprzeczne z przeważającą wykładnią art.~300$^{3}$ §~1 KSH, według której na kapitał akcyjny przeznacza się wszystkie wkłady zdatne do zaliczenia, i rekomenduje przeznaczenie na kapitał także wkładów niepieniężnych. Wtedy odpowiedzialność z art.~300$^{10}$ za znaczne zawyżenie obejmuje wkłady niepieniężne, a ich rzetelny opis i obronna wycena stają się podwójnie ważne: są podstawą tytułu Spółki do technologii wobec inwestorów, urzędów patentowych i organów podatkowych oraz podstawą oświadczenia dyrektorów o wysokości kapitału (pytania \hyperref[memo:3.1.1]{3.1.1}--\hyperref[memo:3.1.2]{3.1.2}).

\textbf{Sama umowa spółki nie przenosi praw} tam, gdzie ustawa wymaga oznaczenia utworu, pól eksploatacji albo formy pisemnej (§~\ref{art:wlasnosc-intelektualna} ust.~4). Dla każdego rodzaju wkładu potrzebny jest osobny dokument rozporządzający, podpisany w dniu aktu ze skutkiem dla spółki w organizacji:

\begin{longtable}{@{}p{2.6cm} p{6.0cm} p{5.6cm}@{}}
\toprule
\textbf{Wkład} & \textbf{Dokument i treść} & \textbf{Załączniki dowodowe} \\
\midrule
\endfirsthead
\toprule
\textbf{Wkład} & \textbf{Dokument i treść} & \textbf{Załączniki dowodowe} \\
\midrule
\endhead
Kod źródłowy, modele, dokumentacja (utwory, w tym programy komputerowe) & Umowa przeniesienia autorskich praw majątkowych: forma pisemna pod rygorem nieważności (art.~53 prawa autorskiego \Wiar{W}); wyraźne wymienienie pól eksploatacji (art.~41 ust.~2; dla programów art.~74 ust.~4 \Wiar{W}); prawo do wykonywania praw zależnych (art.~46); zobowiązanie do niewykonywania praw osobistych w zakresie utrudniającym korzystanie (art.~16); oświadczenie o autorstwie i braku praw osób trzecich; zgoda na dalsze przeniesienie; wynagrodzenie = akcje (wskazać, że obejmuje wszystkie pola) & Wykaz repozytoriów i gałęzi z identyfikatorem commitu, sumy kontrolne archiwum, wykaz plików, wykaz współtwórców (każdy współtwórca musi przenieść swój udział --- art.~9), wykaz komponentów open source z licencjami (copyleft jako ryzyko), protokół wydania, potwierdzenie braku praw pracodawcy lub uczelni (art.~12, 14 prawa autorskiego) \\
Sprzęt & Umowa przeniesienia własności rzeczy ruchomych (forma dowolna; pisemna dla dowodu) z protokołem wydania (art.~155 §~2 KC dla rzeczy oznaczonych co do gatunku \Wiar{W}) & numery seryjne, dowody nabycia, wartość rynkowa (wycena własna albo faktury), oświadczenie o braku obciążeń i zastrzeżeń własności \\
Zgłoszenia patentowe, prawo do uzyskania patentu, patenty & Umowa przeniesienia prawa do uzyskania patentu / prawa z patentu: forma pisemna pod rygorem nieważności (art.~67 ust.~2 i art.~12 PWP \Wiar{W}); wniosek o wpis zmiany uprawnionego do rejestru UPRP (art.~67 ust.~3; skuteczność wobec osób trzecich od wpisu \Wiar{W}); dla zgłoszeń EPO/PCT --- odpowiednio wniosek o zmianę zgłaszającego & numery zgłoszeń, terytoria, terminy priorytetu, wykaz twórców i ich oświadczenia (art.~11 PWP: prawo do patentu przysługuje twórcy, chyba że wynalazek powstał przy wykonywaniu obowiązków z umowy o pracę), dowody wniesienia opłat \\
Know-how utrwalone & Tylko jeżeli jest zidentyfikowane w dokumentacji (§~\ref{art:wklady} ust.~2 wyłącza nieutrwalone); umowa przeniesienia z oznaczeniem nośników i zobowiązaniem poufności & wykaz dokumentów, sumy kontrolne \\
\bottomrule
\end{longtable}

\textbf{Wycena.} Umowa wymaga przypisania wartości każdemu wkładowi (§~\ref{art:wklady} ust.~3). Rekomendujemy wycenę wewnętrzną z uzasadnieniem (koszt odtworzenia albo metoda porównawcza) podpisaną przez wszystkich Założycieli, a przy zgłoszeniach patentowych --- opinię rzecznika patentowego o stanie zgłoszeń. Wartości powinny być spójne z deklaracjami PIT Założycieli (T9). Znaczne zaniżenie wartości nie rodzi odpowiedzialności korporacyjnej, ale zawyżenie wobec organów podatkowych --- tak.

\textbf{Kompletność w dniu podpisania (pytanie \hyperref[memo:3.1.4]{3.1.4}).} Bez wypełnienia pól Załącznika nr~1 (dane, numery akcji, opis i wartość każdego wkładu, kwoty pieniężne i suma) umowa nie spełnia art.~300$^{5}$ KSH; Załącznik nr~2 musi wskazać osoby objęte mechanizmem albo stwierdzić ich brak; Załącznik nr~3 --- ujawnić aktywności albo wpisać „brak”.

\section*{T5. Kryteria wyboru podmiotu prowadzącego rejestr akcjonariuszy}
\addcontentsline{toc}{section}{T5. Kryteria wyboru podmiotu prowadzącego rejestr akcjonariuszy}

Rejestr akcjonariuszy P.S.A. prowadzi podmiot uprawniony do prowadzenia rachunków papierów wartościowych (dom maklerski, bank powierniczy) albo notariusz prowadzący kancelarię w Polsce; wybór należy do Walnego Zgromadzenia, a przy zawiązaniu do założycieli (art.~300$^{31}$ \Wiar{W}). Umowa spółki nie przesądza wyboru (przekreślone zdanie w §~\ref{art:akcje} ust.~6 powtarza ustawę). Rejestr jest jawny dla Spółki i akcjonariuszy, a nabycie akcji następuje z chwilą wpisu (art.~300$^{36}$, 300$^{37}$ \Wiar{W}) --- podmiot prowadzący rejestr jest więc realnym wykonawcą wszystkich mechanizmów obrotu z umowy.

\textbf{Kryteria wyboru (w kolejności wagi).}
\begin{enumerate}
  \item \textbf{Obsługa ograniczeń i obowiązków związanych z akcją.} Czy podmiot wpisze do rejestru: wymóg zgody Spółki, lock-up z terminem, obowiązek zbycia (zwrotne zbycie), ograniczenie wstąpienia spadkobierców; czy przy żądaniu wpisu nabywcy sprawdzi zgodę Spółki i odmówi wpisu bez niej; jak dokumentuje weryfikację. Prosić o pisemny opis procedury.
  \item \textbf{Dyspozycje pełnomocnika i wpisy na podstawie wyroku.} Czy przyjmie dyspozycję zbycia złożoną przez pełnomocnika z umowy wykonawczej (jakiej formy pełnomocnictwa wymaga: pisemna, notarialnie poświadczona) oraz wpis na podstawie prawomocnego wyroku z art.~64 KC.
  \item \textbf{Wpisy warunkowe i depozyt.} Czy obsłuży przeniesienie z warunkiem zapłaty (dyspozycja po wpłacie do depozytu) i spłatę spadkobiercy niewstępującego.
  \item \textbf{Adresy e-mail i zawiadomienia.} Ujawnienie adresu poczty elektronicznej akcjonariusza (art.~300$^{87}$ KSH) i sposób jego aktualizacji, bo umowa opiera na tym zawiadomienia o WZ (§~\ref{art:wz} ust.~3).
  \item \textbf{Głosowanie i WZ elektroniczne.} Czy udostępnia system do identyfikacji i głosowania zdalnego (§~\ref{art:wz} ust.~5--6); czy wystawia zaświadczenia i wyciągi w formie elektronicznej.
  \item \textbf{Obsługa emisji serii P i F.} Wpisy zbiorcze po każdej transzy, ewidencja limitów, świadectwa rejestrowe dla uczestników programu; koszt wpisu per akcjonariusz przy kilkudziesięciu uczestnikach.
  \item \textbf{Weryfikacja nabywcy i AML.} Zakres KYC, jaki podmiot sam prowadzi (dom maklerski jest instytucją obowiązaną), i czy udostępni Spółce dane potrzebne do weryfikacji Kryterium.
  \item \textbf{Koszty i umowa.} Opłata za założenie, roczna, za wpis, za zaświadczenie; kary za opóźnienie; możliwość zmiany podmiotu i format eksportu danych.
  \item \textbf{Ciągłość i bezpieczeństwo.} Dom maklerski podlega KNF i ma procedury ciągłości; notariusz jest tańszy i bliższy, ale obsługa dużych emisji i systemów głosowania jest u niego ograniczona \Wiar{W}.
\end{enumerate}

\textbf{Rekomendacja.} Dla spółki, która planuje program motywacyjny dla kilkudziesięciu osób, rundę i zdalne WZ, dom maklerski z doświadczeniem w P.S.A.; notariusz jako rozwiązanie przejściowe, jeżeli przy zawiązaniu liczy się koszt (decyzja D10). Niezależnie od wyboru --- przed podpisaniem umowy z rejestrem wysłać mu §~\ref{art:zwrotne-zbycie}--\ref{art:dziedziczenie} umowy i uzyskać pisemne potwierdzenie sposobu obsługi (checklista w etapie B).

\section*{T6. Kody PKD i ryzyko zastrzeżeń sądu rejestrowego}
\addcontentsline{toc}{section}{T6. Kody PKD i ryzyko zastrzeżeń sądu rejestrowego}

\textbf{PKD.} Wniosek do KRS obejmuje najwyżej dziesięć pozycji, w tym jedną przeważającą (art.~40 pkt~1 ustawy o KRS) \Wiar{Z}, a przedmiot działalności w umowie może być szerszy; §~\ref{art:pkd} ust.~5--6 stosuje ten podział prawidłowo. Kontrola 19 kodów z §~\ref{art:pkd} na tekście rozporządzenia RM z 18.12.2024 r. w sprawie PKD 2025 (Dz.U. 2024 poz. 1936, \texttt{src/legal/pl\_2024\_1936\_pkd\_2025.txt}) \Wiar{Z}: wszystkie 19 symboli istnieje w PKD 2025, a nazwy w §~\ref{art:pkd} ust.~5--6 są identyczne z urzędowymi co do znaku, w tym podklasy literowe 62.10.B, 62.20.A, 62.20.B, 63.10.D, 71.12.B, 29.10.E oraz rozdzielone w PKD 2025 kody cywilne i wojskowe 30.31.Z / 30.32.Z, 30.13.Z, 30.40.Z, 33.18.Z. Symbol 72.10.Z (badania naukowe i prace rozwojowe w dziedzinie nauk przyrodniczych i technicznych) jest prawidłowy. Zastrzeżenie sądu co do brzmienia kodów jest wykluczone; jedyna kontrola przed wnioskiem to zgodność z formularzem PRS, który podpowiada nazwy z tej samej klasyfikacji.
\textbf{Działalność przeważająca (decyzja D7).} 72.10.Z jest spójne z profilem B+R na starcie i z wnioskami grantowymi (Ścieżka SMART, EIC); 30.31.Z stanie się właściwe, gdy przychody z produkcji przekroczą przychody z prac rozwojowych --- GUS określa działalność przeważającą według udziału w wartości dodanej/przychodach, a zmiana wpisu w KRS nie wymaga zmiany umowy, jeżeli kod jest w §~\ref{art:pkd} \Wiar{W}. Kody wojskowe: ujawnienie w KRS przed koncesją nie jest zakazane, ale banki, kontrahenci i programy cywilne (EIC) czytają wpis dosłownie; rekomendujemy nieujawnianie ich w pierwszym wniosku (jak w §~\ref{art:pkd} ust.~6) i dodanie przy uzyskaniu koncesji (pytanie \hyperref[memo:4.1.2]{4.1.2}). Zdanie z \texttt{psa\_todo.md} („samo wskazanie kodu PKD nie jest podstawą do rozpoczęcia działalności regulowanej”) jest zbędne --- §~\ref{art:pkd} ust.~7--8 mówi to samo.

\textbf{Ryzyko zastrzeżeń sądu.} Sąd rejestrowy bada, czy umowa zawiera elementy obowiązkowe i czy jej postanowienia nie są sprzeczne z prawem; nie ocenia celowości. Katalog postanowień, które w praktyce mogą wywołać wezwanie, z oceną prawdopodobieństwa \Wiar{W} i gotową argumentacją (rozwinięcie w etapie B):

\begin{longtable}{@{}p{4.0cm} p{1.9cm} p{8.3cm}@{}}
\toprule
\textbf{Postanowienie} & \textbf{Ryzyko} & \textbf{Argumentacja / brzmienie awaryjne} \\
\midrule
\endfirsthead
\toprule
\textbf{Postanowienie} & \textbf{Ryzyko} & \textbf{Argumentacja / brzmienie awaryjne} \\
\midrule
\endhead
Kryterium EU/NATO (§~\ref{art:dopuszczalny-nabywca} ust.~5--6, 15) & średnie & Ograniczenie zbywalności przez zgodę spółki z obiektywną przesłanką (art.~300$^{39}$ §~1); przesłanka nie dyskryminuje obywateli UE/EOG; zgoda WZ jako wentyl; uzasadnienie art.~65 ust.~1 lit.~b i art.~346 TFUE \Wiar{W} oraz wymogi programów obronnych (art.~9 ust.~3--4 rozporządzenia 2021/697 \Wiar{Z}; DIANA \Wiar{W}); od 17.01.2028~r. --- jeżeli Spółka mieści się w art.~4 ust.~15 lit.~a lub b rozporządzenia 2026/1386 (pozycje dual-use albo wykaz uzbrojenia) --- także zbieżność z obowiązkową kontrolą inwestycji zagranicznych \Wiar{Z} (tekst EN). Awaryjnie: przenieść ust.~15 (utrata Kryterium) do umowy akcjonariuszy. \\
Derogacja art.~300$^{39}$ §~3--6 (§~\ref{art:zbywanie} ust.~3--4, §~\ref{art:okres-ograniczenia-zbywania} ust.~1) & średnie & Wyraźne upoważnienie ustawowe „chyba że umowa spółki stanowi inaczej”; wariant A modyfikuje, nie wyłącza; lock-up jest czasowy. Awaryjnie: skrócić lock-up albo dodać wskazanie nabywcy po 12 miesiącach. \\
Obowiązek zbycia jako obowiązek związany z akcją (§~\ref{art:zwrotne-zbycie} ust.~14) & niskie & Art.~300$^{33}$ §~1 pkt~11 i art.~300$^{1}$ §~3 przewidują obowiązki akcjonariusza w umowie; sądy znają reverse vesting w sp.~z~o.o. \\
Uchwała ramowa serii P wykonywana przez Radę (§~\ref{art:program-motywacyjny} ust.~3) & niskie przy wpisie, średnie przy zgłoszeniu emisji & Sąd bada przy pierwszym wpisie tylko upoważnienie (liczba, termin); spór pojawi się przy zgłoszeniu transzy (pytanie \hyperref[memo:3.2.1]{3.2.1}). Awaryjnie: odrębna uchwała WZ dla każdej transzy. \\
Mandat dyrektora do 6 miesięcy po kadencji (§~\ref{art:rada} ust.~4) & średnie & Ustawa wiąże wygaśnięcie mandatu z upływem kadencji (pytanie \hyperref[memo:2.1.1]{2.1.1}); awaryjnie: zastąpić prorogację obowiązkiem zwołania WZ i skróconą procedurą uzupełnienia. \\
Głos dyrektora ds. zgodności (§~\ref{art:rada-zastrzezone} ust.~2) & niskie & Umowa może wymagać kwalifikowanej większości albo udziału określonego dyrektora (art.~300$^{58}$ w zw. z art.~300$^{1}$ §~3) \Wiar{W}. \\
§~\ref{art:transakcje-akcjonariusze} (transakcje z akcjonariuszami) & niskie & Nie narusza art.~300$^{15}$ i n.; usunąć znacznik „[WYMAGA OPINII PRAWNIKA]” przed aktem. \\
Ograniczenie wstąpienia spadkobierców (§~\ref{art:dziedziczenie}) & niskie & Wprost z art.~300$^{41}$ §~1, warunki spłaty określone. \\
Liczba i długość postanowień (38 §, 43 strony) & niskie & Nie jest podstawą wezwania; notariusz może sugerować skrócenie ze względu na taksę i czytelność. \\
\bottomrule
\end{longtable}

\section*{T7. Oczekiwania inwestorów VC i zmiany ułatwiające rundę}
\addcontentsline{toc}{section}{T7. Oczekiwania inwestorów VC i zmiany ułatwiające rundę}

Szczegółowa analiza oczekiwań funduszy (standard NVCA/BVCA, praktyka polska, fundusze zagraniczne i obronne) oraz mapowanie na każdy paragraf umowy znajduje się w dokumencie \texttt{vc.md}; odpowiedzi \hyperref[memo:4.4.1]{4.4.1} i \hyperref[memo:4.4.2]{4.4.2} w Części I rozstrzygają, co budzi zastrzeżenia i co zmienić. Tu tylko zbiorcza lista działań:

\textbf{Przed rundą, w umowie spółki (etap B):} (1) przywrócenie uszczelnionego wyjątku dla funduszy w §~\ref{art:dopuszczalny-nabywca} ust.~14; (2) „przeniesienia dozwolone” inwestora finansowego w §~\ref{art:zbywanie} ust.~7; (3) wyłączenie w §~\ref{art:dopuszczalny-nabywca} ust.~15 zdarzeń po stronie inwestorów funduszu; (4) ochrona inwestora przy drag-along w §~\ref{art:przymusowe-wspolzbycie} po Kwalifikowanej Rundzie; (5) wyjątek od Kryterium dla dyrektora wskazanego przez inwestora za zgodą WZ w §~\ref{art:rada} ust.~3; (6) rozważenie zdefiniowania serii inwestorskiej z upoważnieniem emisyjnym; (7) wypełnienie progów §~\ref{art:kwalifikowana-runda} ust.~1.

\textbf{W rundzie, poza umową spółki:} katalog spraw wymagających zgody serii inwestorskiej, prawa informacyjne, lock-up założycieli do wyjścia, zasady wyjścia, oświadczenia i zapewnienia, kary umowne, zobowiązania do głosowania z pełnomocnictwami, arbitraż, waterfall przy sprzedaży spółki. Katalog ten odpowiada temu, co artykuł opublikowany na portalu startup.pfr.pl (autor z kancelarii DLA Piper, 15.03.2023~r.) opisuje jako typową treść umowy inwestycyjnej z funduszem VC, zastrzegając, że nie ma jednego modelowego wzoru: warunki zawieszające z due diligence (w tym umowa przeniesienia praw własności intelektualnej na spółkę), procedura inwestycji, oświadczenia i zapewnienia, ład korporacyjny z listą decyzji wymagających głosu funduszu albo bardzo dużej większości (np.~75--80\%), prawo funduszu do członka rady, non-compete i non-solicitation założycieli (zwykle od 18~miesięcy do 3~lat), pro-rata, anti-dilution, lock-up założycieli, prawo pierwszeństwa, tag i drag along oraz preferencja likwidacyjna --- częściej bez prawa udziału (non-participating) \Wiar{Z} (\texttt{pfr\_startup\_umowa\_inwestycyjna.txt}); według artykułu na portalu PFR, za analizą PwC Raise ponad 200 europejskich term sheetów, klauzula o roli inwestora występuje w 95\%, a antyrozwodnieniowa w 43\% z nich --- przy czym artykuł opisuje tę ostatnią jako prawo do obejmowania nowych udziałów w celu zachowania udziału procentowego (w istocie pro-rata), a nie jako korektę ceny \Wiar{Z} (\nolinkurl{pfr_startup_term_sheet.txt}). Od 17.01.2028~r. dochodzą postanowienia o kontroli inwestycji zagranicznych (T1).

\textbf{Czego nie oddawać:} Kryterium co do kontroli nad Spółką, §~\ref{art:security} ust.~6 i §~\ref{art:export}, obowiązki związane z akcją w rejestrze. Fundusze obronne (NIF; fundusze programu PFR Otwarte Innowacje --- program fund-of-funds dla funduszy inwestujących w MŚP realizujące projekty badawczo-rozwojowe, „w tym w obszarze zaawansowanych technologii obronnych i dual-use”, z ticketem od 5~mln~zł do 15\% budżetu inwestycyjnego funduszu i minimalnym udziałem w spółce 10\% (FAQ) \Wiar{Z} (\texttt{pfr\_pfrv\_otwarte\_innowacje.txt}); strona jest wewnętrznie niespójna co do spółek zagranicznych: opis programu pozwala funduszom inwestować „do 40\% wartości portfela w spółki z zagraniczną siedzibą posiadające polskich founderów, o ile kapitał w całości wspiera rozwój ich działalności w Polsce”, a FAQ oczekuje, że „co najmniej 85\%” wartości portfela trafi do spółek z siedzibą w Polsce, a „do 15\% portfela” do spółek z siedzibą w UE/EFTA/EOG+UK, „które są znacząco obecne w Polsce” --- rozbieżność dotyczy zarówno wysokości limitu (40\% albo 15\%), jak i przesłanki (polscy founderzy i kapitał na działalność w Polsce albo istotna obecność w Polsce i siedziba w UE/EFTA/EOG lub UK) \Wiar{Z} (\texttt{pfr\_pfrv\_otwarte\_innowacje.txt}) co do brzmienia strony; który limit wiąże konkretny fundusz, rozstrzyga jego dokumentacja z PFR (term sheet, zał.~nr~7 do zasad naboru) \Wiar{?}; dla Spółki z siedzibą w Polsce rozbieżność nie ma znaczenia, ma je przy strukturze holdingowej za granicą (D1); nabór do programu jest według strony wstrzymany; fundusze z kapitałem EDF) oczekują tych ograniczeń \Wiar{W}; fundusze generalistyczne przekonuje opis programów, do których dają dostęp. Programy PFR Starter (ticket do 5~mln~zł, przede wszystkim spółki przed pierwszą komercyjną sprzedażą w rozumieniu GBER) i Biznest (koinwestycje funduszu z aniołami biznesu w fazie seed i pre-seed, pakiety mniejszościowe, ticket do 5~mln~zł; nabory do obu programów są według stron wstrzymane) \Wiar{Z} (\texttt{pfr\_pfrv\_starter.txt}, \texttt{pfr\_pfrv\_biznest.txt}) wyznaczają realistyczną wielkość pierwszej rundy krajowej, co ma znaczenie dla progów §~\ref{art:kwalifikowana-runda} ust.~1.

\section*{T8. Co w umowie spółki, co w umowie akcjonariuszy, co w obu}
\addcontentsline{toc}{section}{T8. Co w umowie spółki, co w umowie akcjonariuszy, co w obu}

\textbf{Kryteria podziału.} Do umowy spółki trafia to, co ma działać wobec Spółki, jej organów, rejestru i każdego kolejnego nabywcy akcji (skuteczność erga omnes, art.~300$^{1}$ §~3, art.~300$^{33}$ KSH) oraz to, czego ustawa wymaga w umowie (uprzywilejowanie, ograniczenia zbywalności, ograniczenie wstąpienia spadkobierców, upoważnienia emisyjne). Cena: forma aktu notarialnego przy każdej zmianie, jawność w aktach KRS, kontrola sądu rejestrowego, sztywność wobec inwestora. Do umowy akcjonariuszy trafia to, co wiąże tylko strony (zobowiązania do głosowania, kary umowne, pełnomocnictwa, poufne parametry, zakaz konkurencji po odejściu z odszkodowaniem), co zmienia się często i co nie powinno być jawne. „Oba” oznacza: rdzeń w umowie spółki, wykonanie w umowie akcjonariuszy.

\begin{longtable}{@{}p{4.6cm} p{2.4cm} p{7.2cm}@{}}
\toprule
\textbf{Mechanizm} & \textbf{Gdzie} & \textbf{Uzasadnienie} \\
\midrule
\endfirsthead
\toprule
\textbf{Mechanizm} & \textbf{Gdzie} & \textbf{Uzasadnienie} \\
\midrule
\endhead
Kryterium EU/NATO: definicja, zgoda Spółki, zgoda WZ (§~\ref{art:dopuszczalny-nabywca} ust.~1--13) & umowa spółki & musi wiązać każdego nabywcę i rejestr \\
Utrata Kryterium po nabyciu (§~\ref{art:dopuszczalny-nabywca} ust.~15) & oba & obowiązek w umowie spółki (skuteczność wobec nabywców wtórnych), tryb i wyłączenia dla funduszy w umowie akcjonariuszy/inwestycyjnej \\
Zawężenie kontroli do UE/EOG pod EDF (D6) & umowa akcjonariuszy & warunek czasowy i programowy (art.~5 oraz art.~9 ust.~3--4 i~7 rozporządzenia 2021/697 \Wiar{Z}); nie ograniczać zbywalności w statucie \\
Zgoda Spółki, lock-up, ROFR, tag, drag (§~\ref{art:zbywanie}--\ref{art:przymusowe-wspolzbycie}) & oba & ograniczenia zbywalności tylko w umowie spółki są skuteczne wobec rejestru; procedury, pełnomocnictwa do drag, kary --- w umowie akcjonariuszy \\
Zwrotne zbycie: zasady, harmonogram, ceny (§~\ref{art:zwrotne-zbycie} ust.~1--12, 14) & umowa spółki & obowiązek związany z akcją ujawniany w rejestrze; wiąże nabywców \\
Zwrotne zbycie: oferta, pełnomocnictwo, stwierdzanie Odejścia, depozyt (§~\ref{art:zwrotne-zbycie} ust.~13) & umowa wykonawcza (akcjonariuszy) & dokumenty wykonawcze, dane osobowe, forma pisemna wystarcza \\
Dziedziczenie, małżonek (§~\ref{art:dziedziczenie}--\ref{art:malzonek}) & umowa spółki & wymóg art.~300$^{41}$ §~1; spadkobierca nie jest stroną umowy akcjonariuszy \\
Wartość Godziwa (§~\ref{art:wycena}) & umowa spółki & wspólny standard dla wszystkich mechanizmów, w tym ustawowej spłaty \\
Organy, większości, Sprawy Zastrzeżone (§~\ref{art:organy}--\ref{art:walne-zastrzezone}) & umowa spółki & kompetencje organów działają tylko ze statutu; zgoda serii inwestorskiej w rundzie --- umowa spółki (uprzywilejowanie) plus umowa inwestycyjna \\
Regulamin Rady, bramki G1--G6, karta decyzji & regulamin & §~\ref{art:rada} ust.~11; zmiana uchwałą Rady \\
Własność intelektualna: oświadczenia Założycieli (§~\ref{art:wlasnosc-intelektualna} ust.~5) & oba & oświadczenia w umowie spółki wiążą wobec Spółki; odpowiedzialność, kary i przeniesienia przyszłych utworów w umowie akcjonariuszy i umowach indywidualnych \\
Poufność, kontrola eksportu (§~\ref{art:security}--\ref{art:export}) & oba & obowiązki akcjonariusza w umowie spółki (wiążą nabywców); kary umowne i program zgodności poza nią \\
Zakaz konkurencji w okresie zaangażowania (§~\ref{art:konkurencja} ust.~1--2) & umowa spółki & związany ze statusem Akcjonariusza Funkcjonalnego \\
Zakaz konkurencji po odejściu, odszkodowanie, kary & umowa akcjonariuszy & §~\ref{art:konkurencja} ust.~3 to przewiduje; prawo pracy może wymagać odszkodowania \\
Program motywacyjny: pula, termin, uchwała ramowa (§~\ref{art:program-motywacyjny}) & umowa spółki & upoważnienie z art.~300$^{103}$ musi być w umowie \\
Regulamin programu, umowy uczestnictwa, leaver dla serii P & regulamin i umowy & §~\ref{art:program-motywacyjny} ust.~8 \\
Kwalifikowana Runda: tryb, katalog warunków (§~\ref{art:kwalifikowana-runda} ust.~1--2, 4--5) & umowa spółki & wiąże organy przy emisji \\
Kontrola inwestycji zagranicznych: warunek zawieszający, oświadczenia inwestora o strukturze własności i beneficjencie rzeczywistym, współdziałanie przy zgłoszeniu, long-stop date & umowa inwestycyjna/akcjonariuszy (opcjonalnie ogólny warunek zawieszający w §~\ref{art:kwalifikowana-runda} ust.~4) & od 17.01.2028~r. obowiązkowe uprzednie zezwolenie na inwestycję spoza UE dającą udział w zarządzaniu lub kontroli, gdy Spółka rozwija, wytwarza lub komercjalizuje pozycje z załącznika~I do 2021/821 lub z wykazu uzbrojenia (art.~2 pkt~1, art.~4 ust.~9 i 15, art.~15 ust.~1 i 3 rozporządzenia 2026/1386 \Wiar{Z} (tekst EN)); sekcja T1 \\
Zobowiązanie do współdziałania i głosowania (§~\ref{art:kwalifikowana-runda} ust.~3) & oba & deklaracja w umowie spółki, wykonanie (pełnomocnictwa, kary) w umowie akcjonariuszy (pytanie \hyperref[memo:3.3.1]{3.3.1}) \\
Spółki celowe (§~\ref{art:przedsiewziecia}) & umowa spółki + wzór umowy JV & kompetencje i minimalne prawa Spółki w statucie; klauzule wobec partnera w umowie JV i licencyjnej --- ulepszenia i ewentualny zakaz kwestionowania praw w granicach art.~5 rozporządzenia 2026/877 \Wiar{Z} (tekst EN), a mandat z §~\ref{art:przedsiewziecia} ust.~3 w wersji z odpowiedzi \hyperref[memo:2.4.1]{2.4.1} \\
Impas (§~\ref{art:impas}) & umowa spółki & dotyczy organów \\
Ubezpieczenie na życie, zgoda na badania & umowa akcjonariuszy & dane osobowe, zobowiązanie osobiste \\
Progi kwotowe (100~000, 250~000, 500~000 zł) & umowa spółki, opcjonalnie regulamin & zmiana progów wymaga aktu; rozważyć delegowanie aktualizacji do regulaminu w granicach widełek (opcjonalne) \\
\bottomrule
\end{longtable}

\textbf{Skutek dla etapu C.} Rekomendujemy rozdzielić obecny \texttt{shareholder\_agreement.tex} na umowę wykonawczą zwrotnego zbycia (podpisywaną przez każdego Akcjonariusza Funkcjonalnego, także przyszłych z serii F) i umowę akcjonariuszy (Założyciele, potem inwestorzy przez przystąpienie), bo mają różne strony i różny czas życia (decyzja D12).

\section*{T9. Ryzyka podatkowe widoczne z konstrukcji programu motywacyjnego i wkładów}
\addcontentsline{toc}{section}{T9. Ryzyka podatkowe widoczne z konstrukcji programu motywacyjnego i wkładów}

Sekcja sygnalizuje ryzyka widoczne z konstrukcji umowy; nie zastępuje analizy doradcy podatkowego (decyzja D14). Wszystkie tezy \Wiar{W}.

\begin{enumerate}
  \item \textbf{Program motywacyjny serii P a art.~24 ust.~11--12b PIT.} Preferencja (przychód dopiero przy zbyciu akcji, 19\%) wymaga: programu utworzonego uchwałą WZ spółki akcyjnej albo P.S.A. (od 2022 r.), faktycznego objęcia akcji przez uczestnika i przychodów uczestnika ze stosunku pracy (art.~12) albo działalności wykonywanej osobiście (art.~13). Uchwała ramowa z §~\ref{art:program-motywacyjny} ust.~3 spełnia wymóg „uchwały WZ”, ale wykonywanie transz przez Radę i regulamin przyjmowany przez Radę mogą być kwestionowane jako „program utworzony przez Radę” --- rekomendujemy, aby regulamin programu zatwierdzało WZ razem z uchwałą ramową (zmiana w §~\ref{art:program-motywacyjny} ust.~8: opcjonalne). Ryzyko: objęcie po 0,01~zł przez osobę bez preferencji (współpracownik B2B, dyrektor bez umowy z art.~13, doradca) rodzi przychód w dniu objęcia w wysokości różnicy do wartości rynkowej, z obowiązkami płatnika Spółki przy stosunku pracy. Interpretacja indywidualna przed pierwszym grantem (pkt 27) powinna objąć: kwalifikację P.S.A., uchwałę ramową, moment przychodu dla każdej kategorii uczestników, koszt uzyskania przy zbyciu.
  \item \textbf{Seria F za wkład w postaci pracy lub usług (§~\ref{art:seria-f} ust.~7).} Objęcie akcji za pracę nie zasila kapitału akcyjnego (art.~14 §~1 KSH), a podatkowo może być przychodem z nieodpłatnego świadczenia albo z działalności wykonywanej osobiście w dniu objęcia, według wartości rynkowej akcji; brak przepisu odraczającego. Bezpieczniejszy wariant: wkład pieniężny po niskiej cenie emisyjnej z tym samym pytaniem o wartość rynkową --- też do interpretacji.
  \item \textbf{Aport własności intelektualnej Założycieli (§~\ref{art:wklady}).} Objęcie akcji za wkład niepieniężny to przychód z kapitałów pieniężnych w wysokości wartości wkładu określonej w umowie, nie niższej niż rynkowa (art.~17 ust.~1 pkt~9 PIT), opodatkowany 19\% w roku objęcia; koszt uzyskania to wydatki na wytworzenie wkładu (dla samodzielnie wytworzonego kodu często zerowe). Zwolnienie z art.~21 ust.~1 pkt~109 PIT dotyczy aportu przedsiębiorstwa lub jego zorganizowanej części --- rozważyć, czy zespół kodu, sprzętu i zgłoszeń nie stanowi ZCP (wymaga wyodrębnienia organizacyjnego i finansowego, trudne u osoby fizycznej bez działalności). Wniosek: niska wartość wkładów niepieniężnych w Załączniku nr~1 zmniejsza podatek Założycieli, ale obniża wartość aktywów Spółki w bilansie i „koszt” przy późniejszym zbyciu akcji; wartość musi być obiektywnie uzasadniona.
  \item \textbf{Wkłady pieniężne, kapitał akcyjny, PCC.} Wkłady pieniężne na kapitał akcyjny nie rodzą przychodu po stronie wnoszącego. Ustawa o PCC definiuje spółkę kapitałową jako spółkę z o.o., akcyjną lub europejską (art.~1a pkt~2) i opiera podstawę opodatkowania umowy spółki kapitałowej na kapitale zakładowym (art.~6 ust.~1 pkt~8 lit.~a) \Wiar{Z}; P.S.A. nie ma kapitału zakładowego i nie jest wymieniona, więc według wykładni literalnej umowa P.S.A. i podwyższenie kapitału akcyjnego nie podlegają PCC (stanowisko przeważające w praktyce, do potwierdzenia interpretacją --- pkt 27 zakresu) \Wiar{W}. Tym bardziej wyłączenie wkładów niepieniężnych z kapitału (obecne §~\ref{art:wklady} ust.~4) nie daje żadnej korzyści podatkowej, a --- jak wykazuje odpowiedź \hyperref[memo:3.1.2]{3.1.2} --- jest wątpliwe na tle art.~300$^{3}$ §~1 KSH, nie ma odpowiednika bilansowego (art.~36 ust.~2aa ustawy o rachunkowości odsyła wartość wkładów do kapitału akcyjnego \Wiar{Z}; składnik przyjęty do ksiąg bez źródła w kapitale własnym grozi przychodem CIT z nieodpłatnego przysporzenia \Wiar{W}) i nie zmienia PIT Założycieli. Rekomendacja: wszystkie wkłady na kapitał akcyjny, a wysokość kapitału sterować obronną wyceną.
  \item \textbf{Pożyczki założycielskie (§~\ref{art:wklady} ust.~7).} Pożyczka akcjonariusza dla spółki kapitałowej jest zwolniona z PCC (art.~9 pkt~10 lit.~i), ale ta definicja spółki kapitałowej nie obejmuje P.S.A.; pożyczka dla P.S.A. podlega więc PCC 0,5\% jak pożyczka między osobami prywatnymi, chyba że zwolnienie zostanie potwierdzone interpretacją \Wiar{Z} co do treści przepisów, \Wiar{W} co do wniosku; odsetki rynkowe, ceny transferowe przy powiązaniu, podatek u źródła nie dotyczy rezydentów; odsetki nierynkowe --- przychód z nieodpłatnego świadczenia po stronie Spółki.
  \item \textbf{Zwrotne zbycie po cenie emisyjnej.} Sprzedaż akcji nienabytych po 0,01~zł nabywcy wskazanemu przez Spółkę to dla odchodzącego przychód ze zbycia po cenie umownej; organ może szacować przychód według wartości rynkowej (art.~19 PIT), chyba że przyczyna odbiegania ceny jest uzasadniona --- mechanizm vestingu z umowy spółki jest taką przyczyną, co warto opisać w umowie wykonawczej. Dla nabywcy (inny Założyciel) nabycie poniżej wartości rynkowej może być przychodem z innych źródeł \Wiar{?}.
  \item \textbf{Spółki celowe (§~\ref{art:przedsiewziecia}).} Licencja dla podmiotu powiązanego: ceny transferowe, VAT, wpływ na IP Box; poręczenia Spółki --- nieodpłatne świadczenie po stronie Przedsięwzięcia (\texttt{extra.tex}, Dokument nr~6).
  \item \textbf{Ulgi B+R i IP Box.} Nie wynikają z umowy, ale konstrukcja (72.10.Z, rejestr IP z §~\ref{art:wlasnosc-intelektualna} ust.~6) je wspiera; ewidencja od pierwszego dnia.
\end{enumerate}

% Zestawienie do automatycznego przetworzenia (sekcje tematyczne)

\clearpage
\chapter*{Załączniki}
\addcontentsline{toc}{chapter}{Załączniki}

\section*{Załącznik nr 1 --- Zbiorcza tabela uwag}
\addcontentsline{toc}{section}{Załącznik nr 1 --- Zbiorcza tabela uwag}

Kolumna „Pyt.” wskazuje odpowiedź w Części I, w której znajduje się uzasadnienie i --- dla uwag koniecznych --- pełne proponowane brzmienie; „Miejsce” wskazuje paragraf i ustęp umowy (numeracja ustępów według wersji 0.9.4-C; ustępy oznaczone literą „a” są nowe). Jedna odpowiedź może rodzić kilka uwag.

\subsection*{Uwagi konieczne}
\begin{longtable}{@{}p{1.3cm} p{3.3cm} p{9.6cm}@{}}
\toprule
\textbf{Pyt.} & \textbf{Miejsce} & \textbf{Treść uwagi} \\
\midrule
\endfirsthead
\toprule
\textbf{Pyt.} & \textbf{Miejsce} & \textbf{Treść uwagi} \\
\midrule
\endhead
\hyperref[memo:1.2.4]{1.2.4} & §~\ref{art:zbywanie} ust. 3 (wariant B) & W wariancie B dodać wentyl: po 12 miesiącach bezskutecznego poszukiwania nabywcy akcjonariusz może żądać wskazania nabywcy przez Spółkę. \\
\hyperref[memo:1.5.1]{1.5.1} & §~\ref{art:zbywanie} ust. 2 & Zamknięty katalog odmowy zgody nie pozwala odmówić zbycia akcji objętych obowiązkiem zbycia nabywcy, który nie przystąpił do umowy wykonawczej; dodać lit. f. \\
\hyperref[memo:1.5.1]{1.5.1} & §~\ref{art:zwrotne-zbycie} ust. 14 & Obowiązek zbycia ma obciążać każdoczesnego posiadacza akcji, a Spółka ma żądać wpisu wzmianek w rejestrze na podstawie art.~300$^{33}$ KSH. \\
\hyperref[memo:1.5.3]{1.5.3} & §~\ref{art:zwrotne-zbycie} ust. 13 & Rozszerzyć obowiązkową treść umowy wykonawczej (oferta do Spółki lub umowa przedwstępna zredagowana jako podstawa wpisu z art.~300$^{34}$ §~4 KSH, pełnomocnik oznaczony z podpisem poświadczonym wobec art.~300$^{34}$ §~3 KSH od 18.02.2027, wstrzymanie 6-miesięcznego terminu, skutki niezapłacenia ceny, dopłata po zmianie kwalifikacji, zbieg z prawem pierwszeństwa). \\
\hyperref[memo:1.6.1]{1.6.1} & §~\ref{art:dziedziczenie} ust. 2 & Wskazać dłużnika spłaty (nabywca, solidarnie Spółka), termin wskazania nabywcy i ścieżkę rezerwową przez nabycie akcji własnych lub umorzenie z określeniem niewstąpienia jako przesłanki umorzenia (art.~300$^{45}$ §~1 KSH), pod rygorem bezskuteczności ograniczenia z art.~300$^{41}$ §~1 KSH. \\
\hyperref[memo:1.6.1]{1.6.1} & §~\ref{art:dziedziczenie} ust. 2 & Akcje serii F objęte za wkład pracy niewniesiony w całości: przesądzić, że wstąpienie spadkobierców podlega §17 zamiast zgody Spółki z art.~300$^{41}$ §~2 KSH, i dodać proporcjonalne pomniejszenie spłaty (art.~300$^{41}$ §~1 zd. 3 KSH). \\
\hyperref[memo:1.6.1]{1.6.1} & §~\ref{art:dziedziczenie} ust. 4 & Bezwarunkowy termin pierwszej raty spłaty (30 dni od ustalenia Wartości Godziwej) zamiast przesłanki zagrożenia płynności bez wskazania, kto ją ocenia. \\
\hyperref[memo:1.6.3]{1.6.3} & §~\ref{art:dziedziczenie} ust. 2a & Dodać tryb wezwania i biegu 90 dni, uprawnienie Spółki do wszczęcia postępowania spadkowego, stwierdzenie niewstąpienia uchwałą Rady oraz umorzenie automatyczne z art.~300$^{46}$ KSH (z zastrzeżeniem środków z art.~300$^{15}$ §~2 KSH) jako ścieżkę wobec spadkobiercy niewspółdziałającego. \\
\hyperref[memo:1.7.3]{1.7.3} & §~\ref{art:dopuszczalny-nabywca} ust. 11a & Tryb stwierdzenia statusu Prawnie Niedopuszczalnego Posiadacza, wstrzymania świadczeń i niedopuszczenia do wykonywania praw oraz wyłączenie głosów zawieszonych z mianownika większości umownych, z zastrzeżeniem praw zarządcy przymusowego (art. 6a ustawy z 13.04.2022). \\
\hyperref[memo:1.8.2]{1.8.2} & §~\ref{art:wycena} ust. 4 & Wybór Niezależnego Eksperta wspólnie przez strony, rezerwowo przez instytucję neutralną zamiast Rady Dyrektorów; procedura kontradyktoryjna z terminem 60 dni. \\
\hyperref[memo:1.8.2]{1.8.2} & §~\ref{art:wycena} ust. 7 & Dodać wyjątek oczywistej sprzeczności z zasadami wyceny i tryb ponownej wyceny. \\
\hyperref[memo:3.1.2]{3.1.2} & §~\ref{art:wklady} ust. 4 & Przeznaczyć na kapitał akcyjny także wkłady niepieniężne (art.~300$^{3}$ §~1 KSH); wyłączenie ich z kapitału grozi odmową wpisu i przychodem CIT po stronie Spółki. \\
\hyperref[memo:3.1.4]{3.1.4} & Załącznik nr 1 & Rozdzielić numery akcji obejmowanych za wkłady niepieniężne i pieniężne (art.~300$^{5}$ §~1 pkt 4 KSH) i dodać kolumnę wartości wkładu z podstawą wyceny. \\
\hyperref[memo:3.2.1]{3.2.1} & §~\ref{art:program-motywacyjny} ust. 3 & Oprzeć emisje serii P na ustawowych instrumentach: warunkowej emisji akcji (art.~300$^{114}$–300$^{118}$ KSH) jako trybie podstawowym albo uchwale zwykłej z pełnym katalogiem art.~300$^{104}$ KSH; wyłączyć prawo poboru serii P w umowie na podstawie art.~300$^{106}$ §~1 KSH (ust. 5). \\
\hyperref[memo:3.3.2]{3.3.2} & §~\ref{art:kwalifikowana-runda} ust. 4 & Rozróżnić warunki wymagające zmiany umowy w protokole notarialnym (uprzywilejowanie, prawo wskazania dyrektora, zgoda serii) od kontraktowych; obecne brzmienie sugeruje preferencję bez zmiany umowy. \\
\hyperref[memo:4.3.1]{4.3.1} & §~\ref{art:transakcje-akcjonariusze} & Usunąć oznaczenie robocze „[WYMAGA OPINII PRAWNIKA]” przed sporządzeniem aktu notarialnego. \\
\hyperref[memo:4.3.1]{4.3.1} & §~\ref{art:final} ust. 1 & Usunąć wszystkie przekreślenia (po decyzjach Założycieli) i wypełnić wszystkie pola; akt notarialny nie może zawierać tekstu przekreślonego ani adnotacji roboczych. \\
\hyperref[memo:4.3.1]{4.3.1} & §~\ref{art:final} ust. 7 & Dodać ustęp o uchwałach założycielskich objętych aktem (pierwsza Rada, podmiot prowadzący rejestr) i o reprezentacji spółki w organizacji (art.~300$^{11}$ §~2 KSH). \\
\bottomrule
\end{longtable}

\subsection*{Uwagi zalecane}
\begin{longtable}{@{}p{1.3cm} p{3.3cm} p{9.6cm}@{}}
\toprule
\textbf{Pyt.} & \textbf{Miejsce} & \textbf{Treść uwagi} \\
\midrule
\endfirsthead
\toprule
\textbf{Pyt.} & \textbf{Miejsce} & \textbf{Treść uwagi} \\
\midrule
\endhead
\hyperref[memo:1.1.1]{1.1.1} & §~\ref{art:dopuszczalny-nabywca} ust. 6 & Dodać warstwę „kontroli UE/EOG” obowiązującą w okresie korzystania z EDF/AGILE/EUDIS, z możliwością zgody WZ 75\%. \\
\hyperref[memo:1.1.3]{1.1.3} & §~\ref{art:dopuszczalny-nabywca} ust. 6 & Zdefiniować porozumienie i działanie na cudzy rachunek; dodać obowiązek zawiadomienia o porozumieniu i okresowe ponowienie oświadczenia. \\
\hyperref[memo:1.1.4]{1.1.4} & §~\ref{art:dopuszczalny-nabywca} ust. 15 & Dodać termin 3 miesięcy na wskazanie nabywcy przez Spółkę, wyłączenie Spółki jako nabywcy poza art.~300$^{47}$ i skutek braku wskazania. \\
\hyperref[memo:1.1.5]{1.1.5} & §~\ref{art:rada} ust. 3 & Dopuścić dyrektora spoza Kryterium za zgodą WZ 75\%, jeżeli większość Rady i Przewodniczący spełniają Kryterium, z ograniczeniem dostępu do IP. \\
\hyperref[memo:1.1.6]{1.1.6} & §~\ref{art:zbywanie} ust. 9 (nowy) & Nałożyć na Radę obowiązek ujawnienia ograniczeń i obowiązków związanych z akcją w rejestrze akcjonariuszy w 7 dni (termin z art.~300$^{33}$ §~3 KSH od 18.02.2027 r.). \\
\hyperref[memo:1.2.2]{1.2.2} & §~\ref{art:zbywanie} ust. 3 & Dodać obowiązek Spółki wydania w 7 dni zaświadczenia o powstaniu swobody zbycia jako dokumentu dla rejestru. \\
\hyperref[memo:1.2.3]{1.2.3} & §~\ref{art:zbywanie} ust. 3 & Dodać termin końcowy wyceny (90 dni), ograniczyć raty bez zgody zbywcy (50\% z góry, 12 rat, zastaw rejestrowy albo gwarancja) i określić dzień przeniesienia. \\
\hyperref[memo:1.2.5]{1.2.5} & §~\ref{art:zbywanie} ust. 3 & Rekomendacja wariantu A z poprawkami z \hyperref[memo:1.2.3]{1.2.3}; wariant B tylko z wentylem z \hyperref[memo:1.2.4]{1.2.4}; dodać przeniesienia dozwolone dla inwestora finansowego. \\
\hyperref[memo:1.2.6]{1.2.6} & §~\ref{art:okres-ograniczenia-zbywania} ust. 1--2 & Ograniczyć derogację w §~13 ust. 1 do art.~300$^{39}$ §~3--5 (nie §~6); dopisać przeniesienie do podmiotu kontrolowanego do wyłączeń zgody i prawa pierwszeństwa; ujawnić lock-up z datą końcową. \\
\hyperref[memo:1.2.7]{1.2.7} & §~\ref{art:zbywanie} ust. 4 & Dodać wyraźną derogację art.~300$^{39}$ §~3--5 dla odmowy z powodu Niedopuszczalnego Nabywcy i potwierdzić, że wstrzymanie terminu nie jest odmową. \\
\hyperref[memo:1.2.8]{1.2.8} & §~\ref{art:zbywanie} ust. 7 & Wprowadzić potwierdzenie Rady (niebędące zgodą) jako dokument dla rejestru przy nabyciach zwolnionych ze zgody Spółki; nie wprowadzać mechanizmu na wzór art. 185 KSH. \\
\hyperref[memo:1.3.1]{1.3.1} & §~\ref{art:zbywanie} ust. 1 & Określić skutek milczenia Rady jako fikcję odmowy z przyczyny lit. b uruchamiającą ust. 3. \\
\hyperref[memo:1.3.2]{1.3.2} & §~\ref{art:zbywanie} ust. 2 lit. d--e & Ograniczyć wstrzymanie z lit. d do jednego wezwania i 60 dni; w lit. e dopuścić zgodę warunkową. \\
\hyperref[memo:1.3.4]{1.3.4} & §~\ref{art:pierwszenstwo} ust. 1 i 4 & Określić skutek odmowy zgody dla oferty, początek biegu 60 dni i wymóg Kryterium wobec nabywcy z ust. 4. \\
\hyperref[memo:1.3.5]{1.3.5} & §~\ref{art:pierwszenstwo} ust. 3 & Dodać prawo Spółki do wskazania samodzielnego nabywcy w 20 dni i odesłanie do upoważnienia WZ z art.~300$^{47}$ §~1 pkt 2 KSH; nabycie akcji własnych wymaga celowego kapitału rezerwowego z zysku (art.~300$^{47}$ §~2 pkt 3). \\
\hyperref[memo:1.4.1]{1.4.1} & §~\ref{art:dopuszczalny-nabywca} ust. 14 & Przywrócić wyjątek dla Kwalifikowanego Inwestora Finansowego w brzmieniu uszczelnionym wraz z odesłaniem w §~32 ust. 1. \\
\hyperref[memo:1.4.2]{1.4.2} & §~\ref{art:dopuszczalny-nabywca} ust. 14 & Zachować w wyjątku ujawnienie beneficjentów rzeczywistych w rozumieniu AML, test sankcyjny i współdziałanie z organami koncesyjnymi. \\
\hyperref[memo:1.4.3]{1.4.3} & §~\ref{art:dopuszczalny-nabywca} ust. 14--15 & Uszczelnić wyjątek: regulowany zarządzający, min. 5 inwestorów, progi 25\%/50\%/10\%, wyłączenie wehikułów jednego inwestora, oświadczenie i zawiadomienia. \\
\hyperref[memo:1.5.2]{1.5.2} & §~\ref{art:zwrotne-zbycie} ust. 8 & Zastrzec, że 80\% Wartości Godziwej jest ceną umowną opcji, nie karą umowną, z zaliczeniem różnicy na roszczenia odszkodowawcze; wykluczyć nabywcę finansowanego przez Spółkę. \\
\hyperref[memo:1.5.4]{1.5.4} & §~\ref{art:zwrotne-zbycie} ust. 5 & Doprecyzować przesłankę przestępstwa, ograniczyć przekwalifikowanie Odejścia do 24 miesięcy i dodać regułę dopłaty po prawomocnym ustaleniu innego rodzaju Odejścia. \\
\hyperref[memo:1.5.5]{1.5.5} & §~\ref{art:zwrotne-zbycie} ust. 6 & Dodać niezależność kwalifikacji Odejścia od podstawy zaangażowania oraz skutek prawomocnego ustalenia bezzasadności rozwiązania stosunku prawnego przez Spółkę; ograniczyć kary umowne wobec założycieli-pracowników w umowie wykonawczej. \\
\hyperref[memo:1.5.6]{1.5.6} & §~\ref{art:zwrotne-zbycie} ust. 10 & Zdefiniować w umowie spółki Zmianę Kontroli i Istotną Przyczynę (jako zdarzenia z ust. 5), objąć oknem okres od podpisania do zamknięcia transakcji i konstruktywne odejście. \\
\hyperref[memo:1.6.2]{1.6.2} & §~\ref{art:dziedziczenie} ust. 2 & Próg 75\% dla zgody na wstąpienie spadkobiercy liczyć z wyłączeniem akcji spadkowych; odmowa z uzasadnieniem. \\
\hyperref[memo:1.7.1]{1.7.1} & §~\ref{art:malzonek} ust. 3 & Odesłać do §11 ust. 15 stosowanego odpowiednio do małżonka, który nabył akcje z mocy orzeczenia bez spełnienia Kryterium. \\
\hyperref[memo:2.1.1]{2.1.1} & §~\ref{art:rada} ust. 3--4 & Przebudować prorogację mandatu jako element definicji kadencji i mandatu na podstawie art.~300$^{56}$ KSH (kadencja w latach obrotowych, mandat wygasa z powołaniem następcy, nie później niż 6 miesięcy po kadencji, obowiązek zwołania WZ 3 miesiące przed upływem). \\
\hyperref[memo:2.1.2]{2.1.2} & §~\ref{art:rada-zastrzezone} ust. 2 & Uregulować wyznaczanie dyrektora ds. bezpieczeństwa lub zgodności, tryb przy jego braku, wyłączeniu albo nieobecności oraz skutek braku głosu (uchwała niepodjęta, możliwość przedłożenia sprawy WZ). \\
\hyperref[memo:2.2.1]{2.2.1} & §~\ref{art:wz} ust. 7--8 & Przewidzieć regulamin WZ przyjęty uchwałą akcjonariuszy jako podstawę szczegółowych zasad udziału elektronicznego (art.~300$^{92}$ §~2 KSH: „określa regulamin walnego zgromadzenia”), z regułą przejściową dla zawiadomienia, i dodać „w zakresie dopuszczalnym przez prawo” do klauzuli o problemach technicznych. \\
\hyperref[memo:2.2.2]{2.2.2} & §~\ref{art:wz} ust. 3 & Dodać kaskadę doręczeń (e-mail z rejestru, w jego braku przesyłka polecona/kurierska), obowiązek podania adresu i zgody przy wpisie do rejestru oraz odesłanie do uchwał bez formalnego zwołania. \\
\hyperref[memo:2.2.3]{2.2.3} & §~\ref{art:wz} ust. 11--12 & Dodać zobowiązanie akcjonariusza wyłączonego umownie do niewykonywania głosu oraz regułę konwersji progu „wszystkich głosów” na Głosy Uprawnione w Sprawie przy wyłączeniu ustawowym. \\
\hyperref[memo:2.2.4]{2.2.4} & §~\ref{art:walne-zastrzezone} ust. 5 (nowy) & Określić, że uchwały WZ w sprawach Rady są zgodą, nie decyzją; odesłać do zgód ustawowych (art.~300$^{81}$ KSH), wprowadzić umowny wymóg zgody WZ na kredyt, pożyczkę i poręczenie z dyrektorem lub prokurentem (art. 15 KSH nie wymienia dyrektora P.S.A.) i wyłączyć wymóg uchwały WZ dla nieruchomości do progu z lit. k. \\
\hyperref[memo:2.3.2]{2.3.2} & §~\ref{art:rada} ust. 9 & Ujednolicić pojęcie „pełnego składu” (liczba dyrektorów aktualnie powołanych) i przesądzić, że dyrektor wyłączony nie liczy się do quorum w danej sprawie. \\
\hyperref[memo:2.3.4]{2.3.4} & §~\ref{art:transakcje-akcjonariusze} ust. 2 i 5 & Odesłać do art.~300$^{21}$ KSH (wartość świadczenia Spółki nie wyższa niż wartość godziwa świadczenia wzajemnego, także wobec podmiotów powiązanych spoza przepisu) z sankcją zwrotu nadwyżki (art.~300$^{22}$ KSH); przenieść progi do umowy: pożyczki i zabezpieczenia dla akcjonariusza, w tym dyrektora (niezależnie od art. 15 KSH), Kluczowa IP, 500 000 zł rocznie, mienie od Założyciela w 2 lata od rejestracji — uchwała WZ 75\% Głosów Uprawnionych w Sprawie. \\
\hyperref[memo:2.4.1]{2.4.1} & §~\ref{art:przedsiewziecia} ust. 3 & Zastąpić bezwarunkowy wyłączny grant-back ulepszeń wariantem o niższym ryzyku na tle art. 5 ust. 1 lit. a rozporządzenia 2026/877: własność ulepszeń niewydzielnych (poza wyłączeniem grupowym, do oceny indywidualnej), licencja niewyłączna z sublicencją i prawo pierwszeństwa dla pozostałych, wyłączność tylko w zakresie dozwolonym, z dokumentowaną oceną Rady. \\
\hyperref[memo:2.4.2]{2.4.2} & §~\ref{art:przedsiewziecia} ust. 4 & Dodać warstwę programową (art. 9 EDF: kontrola partnera z państwa niestowarzyszonego wymaga uchwały WZ) i kompetencję Rady do dopuszczenia ograniczonego dostępu dla partnerów z państw wymienionych w EU001 (załącznik II sekcja A do rozporządzenia 2021/821, nie każde unijne zezwolenie generalne) albo związanych umową o bezpieczeństwie informacji z UE lub NATO. \\
\hyperref[memo:2.4.3]{2.4.3} & §~\ref{art:przedsiewziecia} ust. 7 & Ograniczyć wyjątek od zakazu konkurencji: próg 10\% i równy udział Spółki, powyżej — WZ z wyłączeniem zainteresowanego; ujawnienie wynagrodzeń; wygaśnięcie zatwierdzenia z Dniem Odejścia z ofertą sprzedaży udziału Spółce; zachowanie obowiązku oferowania szans biznesowych. \\
\hyperref[memo:2.4.4]{2.4.4} & §~\ref{art:przedsiewziecia} ust. 2a (nowy) & Zdefiniować sposób liczenia łącznego zaangażowania, wprowadzić Limit Zaangażowania zatwierdzany przez WZ, obowiązek informacyjny przed utratą kontroli i regułę kolizyjną z §~25 lit. k i §~31 ust. 5. \\
\hyperref[memo:3.1.1]{3.1.1} & §~\ref{art:wklady} ust. 3 & Wskazać w umowie podstawę wyceny wkładów niepieniężnych i przygotować komplet dokumentów wkładu (protokoły, sumy kontrolne, wykaz twórców, cesje, raport wyceny) przed dniem podpisania. \\
\hyperref[memo:3.1.3]{3.1.3} & §~\ref{art:wlasnosc-intelektualna} ust. 2 & Dodać obowiązek zawarcia pisemnych umów przeniesienia praw (pola eksploatacji, prawa zależne, prawa osobiste, prawa do uzyskania patentu) w oznaczonym terminie. \\
\hyperref[memo:3.3.1]{3.3.1} & §~\ref{art:kwalifikowana-runda} ust. 3 & Rozszerzyć zobowiązanie do głosowania na uchwałę o zmianie umowy, uchwałę 4/5 i uchwały wykonawcze z dokumentów Rundy; karę umowną i pełnomocnictwo umieścić w umowie akcjonariuszy. \\
\hyperref[memo:3.3.3]{3.3.3} & §~\ref{art:okres-ograniczenia-zbywania} ust. 2 lit. a & Zdefiniować wyjątek dla transakcji wtórnych z inwestorem Rundy (zatwierdzenie w uchwale kierunkowej, wyłączenie prawa pierwszeństwa, zachowanie zwrotnego zbycia) i dopuścić zgodę z §~11 ust. 6 w uchwale kierunkowej. \\
\hyperref[memo:3.3.4]{3.3.4} & §~\ref{art:finansowanie} ust. 4 & Zawęzić „prawo weta” do sprzeciwu wobec uchwał organów i wyłączyć standardowe kowenanty kredytowe oraz postanowienia umów o dofinansowanie; rozstrzygnąć, czy ust. 3 ma odsyłać do Kryterium. \\
\hyperref[memo:3.4.1]{3.4.1} & §~\ref{art:final} ust. 6 & Zastrzec, że dane historyczne Załącznika nr 1 nie wymagają aktualizacji, a ujawnienia i zgody z Załącznika nr 3 następują uchwałą WZ bez zmiany umowy. \\
\hyperref[memo:3.4.2]{3.4.2} & §~\ref{art:final} ust. 3 & Dodać zgodę większości serii na uszczuplenie jej praw oraz zgodę dotkniętych akcjonariuszy na zaostrzenie obowiązków związanych z akcją. \\
\hyperref[memo:3.4.4]{3.4.4} & §~\ref{art:wz} ust. 2 & Po usunięciu przekreślonych ustępów poprawić odesłania „ust. 6” i „ust. 11” w §~24 oraz odesłania do §~38 w umowie akcjonariuszy; usunąć nawias z §~32 ust. 1 razem z §~11 ust. 14. \\
\hyperref[memo:4.1.1]{4.1.1} & §~\ref{art:pkd} ust. 8 & Rozszerzyć klauzulę koncesyjną o oprogramowanie i technologię (w tym licencjonowanie) oraz odesłanie dynamiczne do przepisów w brzmieniu obowiązującym; w §~3 ust. 4 zastąpić „certyfikacji lotniczej” szerszym zwrotem. \\
\hyperref[memo:4.2.1]{4.2.1} & §~\ref{art:export} ust. 1 & Dodać przykładowe wyliczenie aktów (2021/821, ustawa o obrocie strategicznym, ustawa koncesyjna, środki ograniczające, AI Act w zakresie zastosowania) z zastrzeżeniem „w brzmieniu obowiązującym”. \\
\hyperref[memo:4.2.1]{4.2.1} & §~\ref{art:export} ust. 5 & Zastąpić nieokreśloną sankcję „środków prowadzących do zbycia akcji” skonkretyzowanym obowiązkiem zbycia na rzecz wskazanego Dopuszczalnego Nabywcy za Wartość Godziwą, jako obowiązek związany z akcją. \\
\hyperref[memo:4.2.3]{4.2.3} & §~\ref{art:security} ust. 3 & Dodać klauzulę przetrwania po utracie statusu akcjonariusza i obowiązek Spółki zawarcia pisemnych zobowiązań z dyrektorami niebędącymi akcjonariuszami oraz nabywcami akcji. \\
\hyperref[memo:4.2.3]{4.2.3} & §~\ref{art:security} ust. 6 & Wskazać adresata zakazu (akcjonariusz, dyrektor, osoba działająca na rzecz Spółki) i uzupełnić katalog terytoriów o EOG zgodnie z Kryterium. \\
\hyperref[memo:4.3.3]{4.3.3} & §~\ref{art:final} ust. 8 & Dodać klauzulę dostosowawczą (odesłania dynamiczne, terminy ustawowe wobec rejestru i KRS, obowiązek Rady przygotowania dostosowania w 6 miesięcy od zmiany prawa; realizuje art. 34 ustawy z 23.01.2026 r.). \\
\hyperref[memo:4.3.3]{4.3.3} & §~\ref{art:wz} ust. 3 & Wpisać 7-dniowy umowny termin zgłaszania przez akcjonariusza zmian danych podlegających wpisowi do rejestru akcjonariuszy, w tym zmian istotnych dla Kryterium (ustawowy termin 7 dni z art.~300$^{33}$ §~3 KSH dotyczy Spółki). \\
\hyperref[memo:4.4.1]{4.4.1} & §~\ref{art:dopuszczalny-nabywca} ust. 6, 14, 15 & Kryterium wobec funduszu bez wyjątku dla funduszy i z utratą Kryterium przy zmianie po stronie LP będzie negocjowane z każdym funduszem; test należy przenieść na poziom zarządzającego i kontroli. \\
\hyperref[memo:4.4.1]{4.4.1} & §~\ref{art:zbywanie} ust. 3 & Dla rundy wariant A (ratalny) jest łatwiejszy do obrony niż wariant B; fundusze zażądają płatności jednorazowej i krótszego terminu dla serii inwestorskiej. \\
\hyperref[memo:4.4.1]{4.4.1} & §~\ref{art:przymusowe-wspolzbycie} ust. 1--2 & Drag 75\% bez zgody serii inwestorskiej, ceny minimalnej i waterfall pozwala zmusić inwestora z udziałem do ok. 25\% do sprzedaży; standard rynkowy wymaga ochrony serii. \\
\hyperref[memo:4.4.1]{4.4.1} & §~\ref{art:walne-zastrzezone} ust. 1 & Sprawy Zastrzeżone jako próg 75\% wszystkich głosów należą do polskiego standardu (obok wymogu głosu funduszu „za”), ale nie dają inwestorowi z udziałem poniżej 25\% żadnego weta; fundusze wymagają wtedy katalogu spraw wymagających zgody serii. \\
\hyperref[memo:4.4.1]{4.4.1} & §~\ref{art:zbywanie} ust. 1 i 7 & Brak przeniesień dozwolonych dla funduszu (fundusz następca, spółka celowa, wydanie LP) oznacza zgodę Spółki, prawo pierwszeństwa i Kryterium przy każdym przeniesieniu w grupie funduszu. \\
\hyperref[memo:4.4.2]{4.4.2} & §~\ref{art:dopuszczalny-nabywca} ust. 14 & Przywrócić wyjątek dla funduszy w brzmieniu uszczelnionym: test zarządzającego, siedziby i kontroli, próg 25\% dla inwestora spoza Kryterium, wyłączenie wehikułów jednego inwestora, oświadczenie funduszu, utrzymanie testu sankcyjnego i AML (brzmienie I). \\
\hyperref[memo:4.4.2]{4.4.2} & §~\ref{art:dopuszczalny-nabywca} ust. 15 & Zawęzić utratę Kryterium przez fundusz do zmiany kontroli nad zarządzającym i uprawnień inwestora spoza Kryterium, z wyłączeniem zwykłych zmian w składzie LP (brzmienie II). \\
\hyperref[memo:4.4.2]{4.4.2} & §~\ref{art:zbywanie} ust. 8 (nowy) & Dodać Przeniesienia Dozwolone dla Kwalifikowanego Inwestora Finansowego bez zgody Spółki i prawa pierwszeństwa, z weryfikacją, przystąpieniem nabywcy i zawiadomieniem Rady (brzmienie III). \\
\hyperref[memo:4.4.2]{4.4.2} & §~\ref{art:przymusowe-wspolzbycie} ust. 5 (nowy) & Po rundzie drag wobec serii inwestorskiej tylko za zgodą większości tej serii albo po okresie ochronnym z ceną minimalną, z podziałem ceny według uprzywilejowania (brzmienie IV). \\
\hyperref[memo:4.4.2]{4.4.2} & §~\ref{art:kwalifikowana-runda} ust. 1 & Przywrócić przekreślone odesłanie do Kwalifikowanego Inwestora Finansowego w definicji Kwalifikowanej Rundy, jeżeli ust. 14 zostaje przywrócony. \\
\bottomrule
\end{longtable}

\subsection*{Uwagi opcjonalne}
\begin{longtable}{@{}p{1.3cm} p{3.3cm} p{9.6cm}@{}}
\toprule
\textbf{Pyt.} & \textbf{Miejsce} & \textbf{Treść uwagi} \\
\midrule
\endfirsthead
\toprule
\textbf{Pyt.} & \textbf{Miejsce} & \textbf{Treść uwagi} \\
\midrule
\endhead
\hyperref[memo:1.1.2]{1.1.2} & §~\ref{art:zbywanie} ust. 2 lit. e & Doprecyzować, kto i kiedy składa zawiadomienie do organu kontroli inwestycji i kiedy ustaje wstrzymanie terminu; od 17.01.2028 r. (rozp. 2026/1386) zezwolenie obejmie każdego inwestora spoza UE, także z NATO/OECD, nabywającego skuteczny udział w zarządzaniu lub kontroli, bez progu przychodowego. \\
\hyperref[memo:1.2.1]{1.2.1} & §~\ref{art:zbywanie} ust. 3--4 & Ujednolicić odesłania do art.~300$^{39}$ §~2 KSH w każdym miejscu derogacji trybu ustawowego. \\
\hyperref[memo:1.7.2]{1.7.2} & §~\ref{art:malzonek} ust. 5 & Dopisek o minimalizacji danych; istotna jest polityka RODO i klauzula informacyjna poza umową. \\
\hyperref[memo:1.8.1]{1.8.1} & §~\ref{art:wycena} ust. 2 & Rozszerzyć Wartość Godziwą na wynagrodzenie za umorzenie z zastrzeżeniem minimum z art.~300$^{45}$ §~2 KSH. \\
\hyperref[memo:2.1.3]{2.1.3} & §~\ref{art:rada} ust. 6 & Zastąpić „powinien być” słowem „jest”, wskazać, że delegacja nie obejmuje art.~300$^{75}$ §~2--3 i §~23 oraz może być cofnięta. \\
\hyperref[memo:2.3.1]{2.3.1} & §~\ref{art:konkurencja} ust. 8 & Dodać termin 14 dni na zwołanie WZ w celu powołania dodatkowego dyrektora i rozważyć niezależnego dyrektora niewykonawczego niebędącego akcjonariuszem. \\
\hyperref[memo:3.2.2]{3.2.2} & §~\ref{art:seria-f} ust. 4 & Odnieść termin 31.12.2036 do dnia uchwały o emisji i uprościć §~8 ust. 9 (zwiększenie puli tylko zmianą umowy). \\
\hyperref[memo:3.2.3]{3.2.3} & §~\ref{art:seria-f} ust. 3 & Dopisać, że uchwała kwalifikacyjna nie jest przyrzeczeniem złożenia oferty, a roszczenie o akcje powstaje wyłącznie z umowy objęcia. \\
\hyperref[memo:3.4.3]{3.4.3} & §~\ref{art:impas} ust. 1 & Wskazać, kto stwierdza Impas Decyzyjny, i przewidzieć wybór Niezależnego Eksperta przez instytucję niezależną, gdy Rada nie może działać. \\
\hyperref[memo:4.1.1]{4.1.1} & §~\ref{art:dopuszczalny-nabywca} ust. 15 & Rozważyć (pakiet A1) obowiązek akcjonariusza z pakietem co najmniej 20\% dostarczenia zaświadczeń o niekaralności i danych wymaganych w postępowaniu koncesyjnym (art. 10 ust. 1 pkt 2 i art. 17 ustawy z 2019 r.). \\
\hyperref[memo:4.1.2]{4.1.2} & §~\ref{art:pkd} ust. 7 & Dodać zdanie, że wskazanie kodu PKD nie jest podstawą do rozpoczęcia działalności regulowanej (wartość wyłącznie komunikacyjna). \\
\hyperref[memo:4.2.2]{4.2.2} & §~\ref{art:export} ust. 2 & Klasyfikacja dual-use przez specjalistę przed bramką G2 (kandydaci: 9A012.a z oprogramowaniem 9D004.e i technologią 9E001, kategoria 7: 7D002--7D004, 7E004.b, 5A002/5D002, 6A008), przy wątpliwościach wiążące wyjaśnienie z art. 10 ustawy o obrocie strategicznym (do 6 miesięcy); wynik do Załącznika nr 3 pkt 4 (bez zmiany tekstu umowy). \\
\hyperref[memo:4.3.2]{4.3.2} & §~\ref{art:akcje} ust. 6 & Przywrócić ustęp w brzmieniu merytorycznym (podstawa: art.~300$^{33}$ §~2 KSH): wybór podmiotu prowadzącego rejestr uchwałą, obowiązek Spółki zgłaszania ograniczeń i obowiązków związanych z akcją oraz oznaczenia akcji objętych mechanizmem zwrotnego zbycia. \\
\hyperref[memo:4.3.3]{4.3.3} & §~\ref{art:akcje} ust. 2 & Usunąć określenie „imiennych” (w P.S.A. zbędne: akcje nie mają formy dokumentu, art.~300$^{29}$ §~1 KSH; zniesienie podziału w nowelizacji dotyczy S.A.) --- §~5 należy do pakietu A4. \\
\hyperref[memo:4.4.1]{4.4.1} & §~\ref{art:zwrotne-zbycie} ust. 7--9 & Fundusz zażąda wpływu na wybór Dopuszczalnego Nabywcy akcji leavera albo podziału proporcjonalnego; parametr do term sheetu, nie zmiana konstrukcji. \\
\hyperref[memo:4.4.1]{4.4.1} & §~\ref{art:rada} ust. 3 & Kryterium wobec dyrektora wyklucza partnera funduszu spoza UE/NATO; istotne tylko dla funduszy zagranicznych i części deep tech. \\
\hyperref[memo:4.4.1]{4.4.1} & §~\ref{art:finansowanie} ust. 3 & Ust. 3 odwołuje się do testu sankcyjnego (Dopuszczalny Nabywca), nie do Kryterium; venture debt spoza UE/NATO nie jest zakazany, Kryterium działa dopiero przy prawach do akcji przez §~11 ust. 13. \\
\hyperref[memo:4.4.2]{4.4.2} & §~\ref{art:rada} ust. 3 & Dopuścić jednego dyrektora niewykonawczego inwestora niespełniającego Kryterium za zgodą WZ 75\%, z większością Rady spełniającą Kryterium i dostępem do IP tylko w trybie §~27 ust. 6 (brzmienie V). \\
\hyperref[memo:4.4.2]{4.4.2} & §~\ref{art:kwalifikowana-runda} ust. 4 lit. f (nowa) & Dodać do katalogu warunków inwestorskich zamknięty katalog spraw wymagających zgody serii inwestorskiej, aby weto nie było „warunkiem dalej idącym” (brzmienie VI). \\
\bottomrule
\end{longtable}


\section*{Załącznik nr 2 --- Decyzje Założycieli}
\addcontentsline{toc}{section}{Załącznik nr 2 --- Decyzje Założycieli}

Decyzje D1--D16 warunkują etap wdrożenia (\texttt{plan\_prac.md}, bramka A~$\rightarrow$~B); stanowisko Założycieli wpisuje się do \texttt{psa\_feedback.tex}, Część 3.

\begin{longtable}{@{}p{1.2cm} p{9.3cm} p{3.7cm}@{}}
\toprule
\textbf{Nr} & \textbf{Decyzja} & \textbf{Warunek dla} \\
\midrule
\endfirsthead
\toprule
\textbf{Nr} & \textbf{Decyzja} & \textbf{Warunek dla} \\
\midrule
\endhead
D1 & jurysdykcja: P.S.A. w Polsce czy inna struktura (pytanie odłożone; \texttt{law\_uzup.md} U.5) & całego wdrożenia \\
D2 & wariant A (ratalny, rekomendowany) czy B (wyłączający z wentylem z \hyperref[memo:1.2.4]{1.2.4}) & 16, 21 \\
D3 & §~\ref{art:dopuszczalny-nabywca} ust.~14 (fundusze): przywrócić w brzmieniu z \hyperref[memo:1.4.3]{1.4.3} (rekomendowane) czy usunąć & 16 \\
D4 & usunięcie siedmiu przekreślonych zdań z poprawką odesłań (\hyperref[memo:3.4.4]{3.4.4}) & 16 \\
D5 & Kryterium wobec spadkobierców i małżonków: utrzymać (rekomendowane, z ubezpieczeniem na życie) czy wrócić do brzmienia 0.9 (\hyperref[memo:1.6.2]{1.6.2}) & 16 \\
D6 & warstwa „kontroli UE/EOG” pod EDF/AGILE/EUDIS: w umowie akcjonariuszy (rekomendowane) czy w umowie spółki (\hyperref[memo:1.1.1]{1.1.1}, T8) & 13, 26 \\
D7 & PKD przeważające 72.10.Z (rekomendowane na start) czy 30.31.Z; kody wojskowe przy wpisie czy po koncesji (rekomendowane po koncesji) (\hyperref[memo:4.1.2]{4.1.2}, T6) & 11, 19 \\
D8 & kto objęty mechanizmem zwrotnego zbycia (Załącznik nr~2 umowy) albo „brak” & 21 \\
D9 & progi Kwalifikowanej Rundy w §~\ref{art:kwalifikowana-runda} ust.~1 & 16 \\
D10 & podmiot prowadzący rejestr: dom maklerski (rekomendowany) czy notariusz (T5) & 19 \\
D11 & model programu motywacyjnego: obietnica i emisja po nabyciu uprawnień (rekomendowany) / emisja z góry / instrument gotówkowy; osoby B2B (\hyperref[memo:3.2.4]{3.2.4}, T9) & 23, 27 \\
D12 & rozdzielenie umowy wykonawczej i umowy akcjonariuszy na dwa dokumenty (rekomendowane, T8) & etap C \\
D13 & data zawiązania i notariusz; bez oczekiwania na 18.02.2027 r. (rekomendowane, \hyperref[memo:4.3.3]{4.3.3}) & 18, 19 \\
D14 & doradca podatkowy do interpretacji indywidualnej (T9) & 27 \\
D15 & specjalista od klasyfikacji dual-use (\hyperref[memo:4.2.2]{4.2.2}, T1a) & program zgodności \\
D16 & które uwagi zalecane i opcjonalne wdrożyć; dane do Załączników nr~1--3 umowy (wkłady, wartości, daty, role) & 16, 19 \\
\bottomrule
\end{longtable}

\subsection*{Punkty decyzyjne wskazane w odpowiedziach}
\begin{itemize}
  \item (\hyperref[memo:1.1.1]{1.1.1}) Czy warstwę „kontroli UE/EOG” dla EDF zapisać w umowie spółki, czy w umowie akcjonariuszy.
  \item (\hyperref[memo:1.1.5]{1.1.5}) Czy dopuścić dyrektora spoza Kryterium za zgodą WZ 75\% z ograniczeniem dostępu do IP.
  \item (\hyperref[memo:1.2.3]{1.2.3}) Czy raty bez zgody zbywcy mają być dopuszczalne i na jakich progach (50\% z góry, 12 rat).
  \item (\hyperref[memo:1.2.5]{1.2.5}) Wybór między wariantem A (ratalny, rekomendowany) a wariantem B z wentylem czasowym.
  \item (\hyperref[memo:1.4.1]{1.4.1}) Czy przywrócić wyjątek dla funduszy (ust. 14) w brzmieniu uszczelnionym, czy pozostać przy uchwale WZ 75\% dla każdego funduszu z LP spoza Kryterium.
  \item (\hyperref[memo:1.4.3]{1.4.3}) Progi uszczelnienia wyjątku dla funduszy: liczba inwestorów, 25\%/50\%/10\%.
  \item (\hyperref[memo:1.5.2]{1.5.2}) Czy utrzymać 80\% Wartości Godziwej dla akcji zwolnionych przy Odejściu Zawinionym, czy przyjąć ostrzejszy standard VC (niższa z ceny objęcia i Wartości Godziwej) przed rundą.
  \item (\hyperref[memo:1.5.5]{1.5.5}) Jaka podstawa zaangażowania założycieli (powołanie, B2B, umowa o pracę) — od tego zależy zakres kar umownych i zakazu konkurencji w umowie wykonawczej.
  \item (\hyperref[memo:1.6.1]{1.6.1}) Czy śmierć Akcjonariusza Funkcjonalnego ma automatycznie zwalniać akcje niezwolnione (obecnie spłata po cenie emisyjnej z fakultatywnym przyspieszeniem przez Walne Zgromadzenie).
  \item (\hyperref[memo:1.6.2]{1.6.2}) Utrzymać Kryterium Bezpieczeństwa EU/NATO wobec spadkobierców (wariant przyjęty) czy przywrócić wyjątek dla rodziny (ZMIANA 14).
  \item (\hyperref[memo:1.8.2]{1.8.2}) Wskazać instytucję neutralną do rezerwowego wyboru Niezależnego Eksperta.
  \item (\hyperref[memo:2.1.1]{2.1.1}) Czy odwołanie z ważnych powodów ma pozostać bezwzględną większością głosów oddanych, skoro akcjonariusz z 50\% może ją zablokować, i czy wprowadzić przykładowy katalog ważnych powodów.
  \item (\hyperref[memo:2.1.2]{2.1.2}) Czy wymóg głosu dyrektora ds. bezpieczeństwa ma mieć „wentyl” w postaci przedłożenia sprawy WZ (75\%), czy pozostać bezwzględną bramką.
  \item (\hyperref[memo:2.2.3]{2.2.3}) Czy reguła konwersji progu „wszystkich głosów” na Głosy Uprawnione w Sprawie ma dotyczyć wszystkich Spraw Zastrzeżonych, czy tylko lit. l i m.
  \item (\hyperref[memo:2.2.4]{2.2.4}) Czy nabycie i zbycie nieruchomości do progu 500 000 zł ma należeć do Rady (opt-out z art.~300$^{81}$ pkt 3 KSH) oraz czy wynagrodzenie dyrektorów (lit. i) ma pozostać przy progu 75\%.
  \item (\hyperref[memo:2.3.1]{2.3.1}) Czy Rada ma obejmować co najmniej jednego niezależnego dyrektora niewykonawczego niebędącego akcjonariuszem.
  \item (\hyperref[memo:2.3.4]{2.3.4}) Wysokość rocznego progu transakcji z jednym Akcjonariuszem (proponowane 500 000 zł) i objęcie nabycia mienia od Założycieli w 2 lata od rejestracji zgodą WZ.
  \item (\hyperref[memo:2.4.2]{2.4.2}) Lista państw zaufanych spoza UE/EOG/NATO, dla których Rada może dopuścić ograniczony dostęp partnera, i czy obejmuje Ukrainę.
  \item (\hyperref[memo:2.4.3]{2.4.3}) Próg bezpośredniego udziału Założyciela w Przedsięwzięciu (proponowane 10\%) i obowiązek oferowania udziału Spółce po odejściu.
  \item (\hyperref[memo:2.4.4]{2.4.4}) Czy wprowadzić Limit Zaangażowania zatwierdzany przez WZ dla poszczególnych Przedsięwzięć.
  \item (\hyperref[memo:3.1.2]{3.1.2}) Czy przeznaczyć wkłady niepieniężne na kapitał akcyjny (rekomendowane), czy utrzymać niski kapitał przez ostrożną wycenę wkładów.
  \item (\hyperref[memo:3.1.4]{3.1.4}) Wartości i podstawa wyceny każdego wkładu oraz podział numerów akcji między wkłady niepieniężne i pieniężne dla każdego Założyciela.
  \item (\hyperref[memo:3.2.1]{3.2.1}) Tryb podstawowy emisji serii P: warunkowa emisja akcji (rekomendowana) czy emisja zwykła w transzach; wyłączenie prawa poboru serii P w umowie na podstawie art.~300$^{106}$ §~1 KSH.
  \item (\hyperref[memo:3.2.2]{3.2.2}) Czy skrócić horyzont serii F albo uzależnić emisje serii F po rundzie od zgody inwestora.
  \item (\hyperref[memo:3.2.4]{3.2.4}) Model programu motywacyjnego (obietnica z emisją po nabyciu uprawnień, emisja z góry, instrument gotówkowy) i traktowanie współpracowników B2B.
  \item (\hyperref[memo:3.3.2]{3.3.2}) Czy zdefiniować już w umowie serię I z preferencją 1x bez udziału i upoważnieniem z art.~300$^{103}$ KSH.
  \item (\hyperref[memo:3.3.3]{3.3.3}) Czy transakcje wtórne z inwestorem Rundy mają być wyłączone spod prawa pierwszeństwa pozostałych Założycieli; czy wyłączyć w umowie prawo poboru wobec inwestorów Kwalifikowanej Rundy (art.~300$^{106}$ §~1 KSH) zamiast uchwały 4/5 przy każdej rundzie.
  \item (\hyperref[memo:3.3.4]{3.3.4}) Czy §~31 ust. 3 ma odsyłać do Kryterium Bezpieczeństwa EU/NATO, czy tylko do testu Dopuszczalnego Nabywcy (przy udziale w EDF kontrola finansującego z niestowarzyszonego państwa trzeciego wyłącza kwalifikowalność Spółki, art. 9 ust. 3 rozporządzenia 2021/697).
  \item (\hyperref[memo:3.4.2]{3.4.2}) Czy Założyciele mniejszościowi oczekują ochrony indywidualnej, czy wystarcza ochrona klasowa i zgoda dotkniętych na zaostrzenie obowiązków związanych z akcją.
  \item (\hyperref[memo:3.4.4]{3.4.4}) Czy przywrócić §~11 ust. 14 (Kwalifikowany Inwestor Finansowy) w wersji uszczelnionej.
  \item (\hyperref[memo:4.1.2]{4.1.2}) Czy pozostawić kody wojskowe w §~4 ust. 6 (rekomendowane) i ujawnić je w KRS dopiero przy bramce G3, w kolejności ustalonej z organem koncesyjnym.
  \item (\hyperref[memo:4.1.3]{4.1.3}) Czy działalnością przeważającą pozostaje 72.10.Z (rekomendowane), czy 30.31.Z (decyzja D7).
  \item (\hyperref[memo:4.2.1]{4.2.1}) Czy sankcją za naruszenie kontroli eksportu przez akcjonariusza niefunkcjonalnego ma być obowiązek zbycia po Wartości Godziwej (rekomendowane, neutralne dla VC), czy z dyskontem.
  \item (\hyperref[memo:4.2.3]{4.2.3}) Czy kara umowna za naruszenie poufności ma znaleźć się w umowie wykonawczej i umowach objęcia akcji (rekomendowane), czy w umowie spółki.
  \item (\hyperref[memo:4.3.1]{4.3.1}) Co zrobić z każdym z dziewięciu przekreślonych fragmentów (usunąć albo przywrócić bez przekreślenia) przed sporządzeniem wersji do aktu; w szczególności §~11 ust. 14.
  \item (\hyperref[memo:4.3.2]{4.3.2}) Wybór podmiotu prowadzącego rejestr (notariusz albo dom maklerski) po zapytaniu ofertowym z kryteriami (1)--(8); wybór ująć w uchwale założycielskiej.
  \item (\hyperref[memo:4.3.3]{4.3.3}) Data zawiązania: rekomendujemy nie czekać na 18 lutego 2027 r.; po wejściu w życie nowelizacji złożyć wniosek o ujawnienie podmiotu prowadzącego rejestr do 18 maja 2027 r. (decyzja D13).
  \item (\hyperref[memo:4.4.1]{4.4.1}) Założyciele rozstrzygają, czy krąg inwestorów ma obejmować fundusze zagraniczne z LP spoza UE/NATO (wtedy potrzebne zmiany w §~11 ust. 14, §~12 i §~21 ust. 3), czy tylko fundusze polskie i obronne (wystarczy §~11 ust. 14 i §~16).
  \item (\hyperref[memo:4.4.1]{4.4.1}) Wybór wariantu A (ratalny) albo B (wyłączający) w §~12 ust. 3 z uwzględnieniem, że wariant B utrudni rundę bez jednoczesnych przeniesień dozwolonych i wyjątku dla funduszy.
  \item (\hyperref[memo:4.4.2]{4.4.2}) Czy katalog spraw wymagających zgody serii inwestorskiej wpisać do §~32 ust. 4 lit. f już teraz, czy pozostawić odrębnej uchwale 75\% w rundzie; oraz jaki okres ochronny i cenę minimalną przyjąć w §~16 ust. 5.
  \item (\hyperref[memo:4.4.2]{4.4.2}) Czy zdefiniować serię I z upoważnieniem do emisji i uprzywilejowaniem przed rundą (A.7), czy zostawić pełną zmianę umowy na moment rundy.
\end{itemize}

\section*{Załącznik nr 3 --- Źródła i lista do weryfikacji}
\addcontentsline{toc}{section}{Załącznik nr 3 --- Źródła i lista do weryfikacji}

\subsection*{Źródła zweryfikowane po pobraniu (wersje 1.1 i 1.2)}
Teksty aktów pobrane skryptem \texttt{src/tools/download.sh} i przepisane do plików tekstowych skryptem \texttt{src/tools/extract\_text.py} (katalog \texttt{src/legal/}, stan na 30 września 2026 r.). Każda podstawa krajowa oznaczona \Wiar{Z} została odszukana w tych tekstach.
\begin{itemize}
  \item \textbf{Kodeksy (teksty ujednolicone ELI):} KSH (tekst ujednolicony ELI i t.j. Dz.U. 2024 poz. 18), KC (Dz.U. 2026 poz. 795), KPC (Dz.U. 2026 poz. 468), KRO (Dz.U. 2026 poz. 236), Kodeks pracy, Konstytucja RP, Ordynacja podatkowa.
  \item \textbf{Ustawy szczególne:} o KRS, o kosztach sądowych, o obrocie instrumentami finansowymi, o kontroli niektórych inwestycji, AML, sankcyjna z 13.04.2022 r., o ochronie informacji niejawnych, o zwalczaniu nieuczciwej konkurencji, o rachunkowości, Prawo lotnicze, o obrocie z zagranicą towarami o znaczeniu strategicznym, prawo autorskie, PWP, PIT, PCC.
  \item \textbf{Obronność:} ustawa z 13.06.2019 r. o wykonywaniu działalności gospodarczej w zakresie wytwarzania i obrotu (tekst ogłoszony i ujednolicony), rozporządzenie RM z 17.09.2019 r. z wykazem WT (tekst ogłoszony i ujednolicony), ustawa z 13.03.2026 r. (Dz.U. 2026 poz. 471).
  \item \textbf{Nowelizacja KSH:} ustawa z 23.01.2026 r. (Dz.U. 2026 poz. 176), pełny tekst z przepisami przejściowymi.
  \item \textbf{PKD 2025:} rozporządzenie RM (Dz.U. 2024 poz. 1936) oraz objaśnienia GUS.
  \item \textbf{Strony rynkowe i informacyjne:} PFR Ventures (term sheet, umowa inwestycyjna, dokumentacja FENG), Startup Poland (term sheet), NVCA (dokumenty modelowe 2025, omówienia), NATO Innovation Fund (strona i raport 2026), omówienia EDF, ustawy z 13.03.2026 r. i koncesji, biznes.gov.pl (P.S.A.), Portal Rejestrów Sądowych.
  \item \textbf{Prawo Unii (wersja 1.2, \texttt{src/legal/eu/}):} rozporządzenie 2021/697 (EDF) PL i EN; w wersji angielskiej: 2021/821 (konsolidacja 26.05.2023), 269/2014 (tekst pierwotny), 2019/947, 2024/1689 (AI Act), 2026/877 (transfer technologii), 2026/1386 (kontrola inwestycji zagranicznych).
  \item \textbf{PFR Ventures (wersja 1.2, \texttt{src/legal/pfr/}):} strony programów PFR Starter, Biznest, Otwarte Innowacje, KOFFI, strona dokumentacji FENG (bez samych wzorów) oraz artykuły startup.pfr.pl o umowie inwestycyjnej i term sheet.
  \item \textbf{Nie udało się pobrać:} orzeczenia SN (pliki \texttt{sn\_*.txt} zawierają stronę błędu 404; tezy orzeczeń pozostają \Wiar{W}), dokumenty modelowe BVCA i część stron PFR Ventures (odpowiedź pusta), część aktów prawa Unii (TFUE, 2019/452, 2024/2847, 2023/1230, 833/2014, 2018/1139, 2019/945) oraz polskie wersje aktów dostępnych po angielsku; wzory term sheet PFR Ventures (adresy w \texttt{src/legal/todo.md}).
\end{itemize}

\subsection*{Pozostaje do sprawdzenia w tekstach}
\begin{enumerate}
  \item Prawo Unii bez pobranego tekstu: art.~49, 63, 65, 346 TFUE; rozporządzenia 2019/452 (obowiązuje do 16.01.2028 r.), 2024/2847 (CRA), 2023/1230 (maszynowe), 2018/1139 i 2019/945 (BSP), 833/2014 (sankcje sektorowe); wytyczne Komisji do 2026/877; aktualna konsolidacja załącznika I do 2021/821 (po 2023 r.) i konsolidacja 269/2014; polskie wersje aktów sprawdzonych po angielsku (cytaty w memorandum są tłumaczeniem roboczym).
  \item Orzecznictwo SN powołane w \texttt{psa\_feedback.tex} (V CSK 522/18, II CSKP 593/22, III CZP 109/22, III CSKP 65/21, III CZP 32/16) oraz orzecznictwo o granicach ograniczeń zbywalności i derogacji art.~300$^{39}$ §~2 KSH (termin dłuższy niż miesiąc, wariant A).
  \item Akty nie objęte pobraniem: ustawa o ofercie publicznej (art.~87), ustawa o funduszach inwestycyjnych i zarządzaniu AFI, ustawa o zastawie rejestrowym, rozporządzenie RM z wykazem podmiotów podlegających ochronie (art.~4 ust.~2 ustawy o kontroli niektórych inwestycji), ustawa o CIT.
  \item Tezy o praktyce: praktyka Komisji przy ocenie kontroli w EDF, praktyka referendarzy przy wpisie P.S.A. z nietypowymi postanowieniami, praktyka podmiotów prowadzących rejestr, stanowisko doktryny o wkładach niepieniężnych poza kapitałem akcyjnym, interpretacje podatkowe do art.~24 ust.~11--12b PIT wobec P.S.A.
\end{enumerate}

\subsection*{Lista do weryfikacji przez uprawnionego prawnika}
\begin{itemize}
  \item (\hyperref[memo:1.1.1]{1.1.1}) Brak skutku horyzontalnego art. 49 i 63 TFUE wobec umowy spółki prywatnej; aktualna lista państw stowarzyszonych z EDF (art. 5 rozp. 2021/697 mówi ogólnie o członkach EFTA/EOG; Norwegia tylko według web\_edf\_gowling.txt) i praktyka Komisji co do pasywnych LP przy teście kontroli.
  \item (\hyperref[memo:1.1.2]{1.1.2}) Czy wykaz podmiotów podlegających ochronie (rozporządzenie RM z art. 4 ust. 2 ustawy o kontroli niektórych inwestycji) nie obejmuje i nie obejmie Spółki; aktualność organu kontroli (minister ds. gospodarki) po kolejnych nowelizacjach; tekst rozp. 2019/452 (niedostępny) i stan prac nad zmianą ustawy wykonującą rozp. 2026/1386 przed 17.01.2028 r. (usunięcie progu 10 mln euro i wyłączenia OECD dla podmiotów z art. 4 ust. 15; czy ustawa wprowadzi próg głosów, motyw 20 rozp. 2026/1386; czy Spółka opracowuje, produkuje lub wprowadza do obrotu produkty z zał. I do rozp. 2021/821 w aktualnym brzmieniu albo z załącznika do dyrektywy 2009/43/WE — klasyfikacja produktów).
  \item (\hyperref[memo:1.2.1]{1.2.1}) Stanowisko doktryny i orzecznictwa co do granic derogacji z art.~300$^{39}$ §~2 KSH, w szczególności czy umowny termin wskazania nabywcy dłuższy niż miesiąc (§~3 zd. 4) jest dopuszczalny na podstawie §~2.
  \item (\hyperref[memo:1.2.4]{1.2.4}) Orzecznictwo o nieważności nadmiernych ograniczeń zbywalności udziałów i akcji (art. 58 §~2 KC) i skutku nieważności częściowej klauzuli.
  \item (\hyperref[memo:1.4.1]{1.4.1}) Warunki uczestnictwa w NATO DIANA (wymóg siedziby w państwie NATO) — brak tekstu źródłowego w materiałach; praktyka organów kontroli inwestycji co do funduszy unijnych z pasywnymi LP spoza UE (motyw 18 rozp. 2026/1386).
  \item (\hyperref[memo:1.4.2]{1.4.2}) Wdrożenie pakietu AML 2024 (rozporządzenie 2024/1624) do ustawy z 2018 r.; czy Spółka nie stanie się instytucją obowiązaną (art. 2 ust. 1 ustawy AML) z tytułu innej działalności; praktyka ABW/SKW co do żądania danych o inwestorach funduszu; aktualna skonsolidowana wersja rozp. 269/2014 (sprawdzono tekst pierwotny) i rozp. 833/2014 (tekst niedostępny).
  \item (\hyperref[memo:1.4.3]{1.4.3}) Term Sheet (zał. nr 7 do Zasad Naboru) programów PFR OI i KOFFI — liczba i struktura LP oraz limit geograficzny (strona OI podaje sprzecznie 40\% w spółki zagraniczne z polskimi founderami albo min. 85\% Polska / do 15\% zagranica).
  \item (\hyperref[memo:1.5.1]{1.5.1}) Praktyka wybranego podmiotu prowadzącego rejestr (notariusz / dom maklerski) co do brzmienia wzmianek z art.~300$^{33}$ §~1 pkt 10-11 KSH (przepis potwierdzony).
  \item (\hyperref[memo:1.5.2]{1.5.2}) Praktyka interpretacyjna KIS co do sprzedaży akcji po cenie emisyjnej i po 80\% Wartości Godziwej (art. 19 ust. 1 w zw. z art. 17 ust. 2 i art. 11 ust. 2b PIT — przepisy potwierdzone) — rozważyć interpretację indywidualną.
  \item (\hyperref[memo:1.5.3]{1.5.3}) Dopuszczalność oferty skierowanej do osoby oznaczalnej w przyszłości, akceptacja przez podmiot prowadzący rejestr pełnomocnictwa z substytucją i dyspozycji składanej przez pełnomocnika; tezy SN V CSK 522/18 i II CSKP 593/22 potwierdzone tylko wg psa\_feedback.tex (pliki sn\_*.txt zawierają stronę 404).
  \item (\hyperref[memo:1.5.5]{1.5.5}) Aktualne orzecznictwo SN o karach umownych wobec pracowników (niedostępne w plikach) i kwalifikacja podatkowa objęcia akcji serii F za wkład w postaci pracy.
  \item (\hyperref[memo:1.7.1]{1.7.1}) Tezy SN III CZP 109/22, III CSKP 65/21, III CZP 32/16 potwierdzone tylko wg psa\_feedback.tex (pliki sn\_*.txt zawierają stronę 404).
  \item (\hyperref[memo:1.7.3]{1.7.3}) Aktualne stanowisko KE i organów krajowych co do wykonywania prawa głosu z zamrożonych akcji; derogacje pozwalające na zbycie zamrożonych akcji i aktualny załącznik I w wersji skonsolidowanej rozp. 269/2014 (tekst pierwotny zna tylko art. 4–6); czy art.~300$^{39}$ §~6 KSH obejmuje zbycie przez zarządcę przymusowego mające skutki sprzedaży egzekucyjnej (art. 6da ustawy z 13.04.2022).
  \item (\hyperref[memo:2.1.1]{2.1.1}) Uchwała SN III CZP 109/22 (plik niedostępny) i dopuszczalność w doktrynie kadencji z mandatem trwającym do powołania następcy na tle art.~300$^{56}$ §~2 KSH (tekst ustawy potwierdzony).
  \item (\hyperref[memo:2.1.2]{2.1.2}) Granice art.~300$^{58}$ §~4 KSH: czy „chyba że umowa spółki stanowi inaczej” obejmuje wymóg głosu określonego (funkcyjnie) dyrektora (doktryna); uchwała SN III CZP 13/13 (niedostępna).
  \item (\hyperref[memo:2.1.3]{2.1.3}) Czy art.~300$^{132}$ KSH (literalnie: „członkowie zarządu”; art.~300$^{133}$ rozciąga tylko na likwidatorów) obejmuje dyrektorów rady dyrektorów.
  \item (\hyperref[memo:2.2.1]{2.2.1}) Czy „regulamin walnego zgromadzenia” z art.~300$^{92}$ §~2 KSH musi być przyjęty uchwałą akcjonariuszy, czy może go przyjąć organ zwołujący (doktryna; tekst ustawy nie wskazuje organu).
  \item (\hyperref[memo:2.2.3]{2.2.3}) Dopuszczalność umownego wyłączenia akcjonariusza od głosowania w P.S.A. poza katalogiem art.~300$^{96}$ KSH (orzecznictwo i doktryna na tle art. 244 i 413 KSH).
  \item (\hyperref[memo:2.2.4]{2.2.4}) Stosowanie art. 15 KSH do dyrektora P.S.A. (doktryna, orzecznictwo) — tekst art. 15 §~1 nie wymienia dyrektora, a przepisy o P.S.A. nie odsyłają do niego.
  \item (\hyperref[memo:2.3.1]{2.3.1}) Czy jakiekolwiek transakcje z podmiotami powiązanymi mogą być „decyzją o strategicznym znaczeniu” z art.~300$^{75}$ §~2 KSH, i orzecznictwo o niedopuszczalności substytucji uchwały rady przez uchwałę akcjonariuszy.
  \item (\hyperref[memo:2.3.2]{2.3.2}) Wykładnia „połowy członków” w art.~300$^{58}$ §~2 KSH (skład faktyczny czy umowny) i czy dyrektor wyłączony na podstawie art.~300$^{55}$ liczy się do quorum.
  \item (\hyperref[memo:2.3.4]{2.3.4}) Zakres podmiotowy art.~300$^{21}$ KSH (osoby fizyczne i podmioty powiązane z akcjonariuszem inne niż spółki i spółdzielnie); progi dokumentacyjne cen transferowych w CIT (tekst niedostępny); zasady konkurencyjnego wyboru w programach EDF/EUDIS/NCBR.
  \item (\hyperref[memo:2.4.1]{2.4.1}) Wytyczne Komisji 2026 w sprawie transferu technologii w części o grant-back i ulepszeniach niewydzielnych (tekst niedostępny; rozporządzenie 2026/877 sprawdzone na tekście EN); art. 6 ustawy o ochronie konkurencji i konsumentów i polskie rozporządzenie wyłączeniowe dla transferu technologii (teksty niedostępne); czy licencja dla spółki celowej jest „obrotem technologią” w rozumieniu art. 7 ust. 1 pkt 2 ustawy koncesyjnej z 2019 r.; czy umowne prawo pierwszeństwa nabycia ulepszeń jest pośrednim zobowiązaniem do przeniesienia praw z art. 5 ust. 1 lit. a rozporządzenia 2026/877.
  \item (\hyperref[memo:2.4.2]{2.4.2}) Aktualne uczestnictwo państw EFTA/EOG w EDF — tekst art. 5 nie wymienia państw (art. 5, 9 i 20 ust. 4 rozporządzenia 2021/697 sprawdzone na tekście PL); warunki kontroli w programie AGILE (akt niedostępny); zmiany załącznika II do 2021/821 po 2023 r. (EU001 sprawdzone na wersji skonsolidowanej 2023); krajowa ustawa wdrażająca rozporządzenie 2026/1386 (zakres greenfield, progi); dostęp podmiotów zagranicznych do informacji niejawnych na podstawie umów dwustronnych.
  \item (\hyperref[memo:2.4.3]{2.4.3}) Skuteczność wygaśnięcia zatwierdzenia i obowiązku zaoferowania udziału Spółce wobec Założyciela po Dniu Odejścia (art.~300$^{1}$ §~3, art.~353$^{1}$ KC) oraz spójność z umową wykonawczą.
  \item (\hyperref[memo:2.4.4]{2.4.4}) Czy wkład Spółki do spółki celowej może stanowić zbycie zorganizowanej części przedsiębiorstwa w rozumieniu art.~300$^{81}$ KSH (uchwała akcjonariuszy z mocy ustawy).
  \item (\hyperref[memo:3.1.2]{3.1.2}) Wykładnia art.~300$^{3}$ §~1 KSH („przeznacza się”): czy wkłady niepieniężne zdatne do zaliczenia muszą w całości tworzyć kapitał akcyjny (tekst potwierdzony, doktryna do sprawdzenia); przychód CIT (art. 12 ust. 1 pkt 2) przy wkładzie poza kapitałem; aktualne interpretacje KIS o PCC od umowy P.S.A. (art. 1a pkt 2 PCC nie wymienia P.S.A.).
  \item (\hyperref[memo:3.1.3]{3.1.3}) Tryb zmiany zgłaszającego w UPRP oraz Rule 22 EPC i Rule 92bis PCT; ustawa o ochronie baz danych (tekst niedostępny).
  \item (\hyperref[memo:3.1.4]{3.1.4}) Wymogi notariusza co do treści Załącznika nr 1 (art. 92 §~1 Prawa o notariacie – tekst niedostępny).
  \item (\hyperref[memo:3.2.1]{3.2.1}) Czy sześciomiesięczny termin z art.~300$^{102}$ §~2 KSH dotyczy emisji na podstawie art.~300$^{103}$ KSH (art.~300$^{107}$ własnego terminu nie zawiera); czy zgoda WZ z art.~300$^{114}$ §~3 KSH może być udzielona zbiorczo w uchwale o warunkowej emisji dla kategorii osób; praktyka rejestrowa wobec uchwały zwykłej obejmującej kilka transz.
  \item (\hyperref[memo:3.2.4]{3.2.4}) Stanowisko organów co do „odpowiedniego” stosowania art. 24 ust. 11–11b PIT do P.S.A. (ust. 12b) przy cenie 0,01 zł i braku wartości nominalnej; art. 18 ustawy o systemie ubezpieczeń społecznych (tekst niedostępny).
  \item (\hyperref[memo:3.3.1]{3.3.1}) Orzecznictwo o skuteczności umów o wykonywanie prawa głosu i o braku wpływu ich naruszenia na ważność uchwały; granice art. 64 KC dla głosowania.
  \item (\hyperref[memo:3.3.2]{3.3.2}) Praktyka rejestrowa wobec serii inwestorskiej zdefiniowanej w umowie przed emisją (art.~300$^{25}$ w zw. z art.~300$^{103}$ KSH); odpowiednie stosowanie art.~300$^{28}$ KSH do rady dyrektorów; czy mniejszościowa inwestycja VC z prawem wskazania dyrektora jest „inwestycją zagraniczną” w rozumieniu art. 2 pkt 1 rozporządzenia 2026/1386 i jakie progi przyjmie polska ustawa wdrażająca; kwalifikacja produktów Spółki według aktualnego zał. I do 2021/821 i załącznika do dyrektywy 2009/43/WE (tekst dyrektywy niedostępny).
  \item (\hyperref[memo:3.3.3]{3.3.3}) Czy wyłączenie prawa poboru w umowie spółki na podstawie art.~300$^{106}$ §~1 KSH wobec kategorii inwestorów wskazanych w uchwale kierunkowej jest akceptowane przez sądy rejestrowe; większość dla późniejszej zmiany umowy w tym przedmiocie (analogia z art.~300$^{113}$ §~3 KSH).
  \item (\hyperref[memo:3.3.4]{3.3.4}) Zakres zastawu rejestrowego na przedsiębiorstwie (ustawa o zastawie rejestrowym – tekst niedostępny) w relacji do art.~300$^{81}$ pkt 2 KSH; czy finansowanie dłużne z warrantem lub konwersją od podmiotu spoza UE jest „inwestycją zagraniczną” (art. 2 pkt 1 rozporządzenia 2026/1386); postanowienia wzoru umowy o dofinansowanie EDF o prawach do wyników.
  \item (\hyperref[memo:3.4.1]{3.4.1}) Praktyka podmiotów prowadzących rejestr akcjonariuszy co do wpisu obowiązków z art.~300$^{33}$ §~1 pkt 11 KSH dla serii F na podstawie umowy i uchwały kwalifikacyjnej.
  \item (\hyperref[memo:3.4.2]{3.4.2}) Czy obowiązek zbycia akcji (zwrotne zbycie, drag-along) jest „świadczeniem akcjonariusza” w rozumieniu art.~300$^{98}$ §~3 KSH.
  \item (\hyperref[memo:3.4.3]{3.4.3}) Charakter ekspertyzy arbitralnej a art. 1157 KPC i kontrola sądowa ustaleń Niezależnego Eksperta.
  \item (\hyperref[memo:3.4.4]{3.4.4}) Mechaniczna kontrola wszystkich odesłań „ust. N” po usunięciu przekreślonych ustępów w §~5, §~11, §~24 i §~38 oraz w umowie akcjonariuszy.
  \item (\hyperref[memo:4.1.1]{4.1.1}) Rozporządzenia (UE) 2018/1139 (art. 2 ust. 3 lit. a --- wyłączenie lotnictwa państwowego) i 2019/945 --- treść niedostępna w toku prac; 2019/947 sprawdzone tylko w tekście pierwotnym (EN) --- potwierdzić na wersji skonsolidowanej zmiany z 2020/639, 2020/746 i późniejsze (data stosowania, scenariusze standardowe, część A załącznika); kwalifikacja konkretnego wyrobu Spółki do WT V ust. 3 (kryterium „specjalnie zaprojektowany lub zmodyfikowany”) przed bramką G3.
  \item (\hyperref[memo:4.1.2]{4.1.2}) Praktyka organu koncesyjnego (MSWiA) co do zgodności zakresu koncesji z kodami PKD w KRS; art. 37 i n. Prawa przedsiębiorców (tekst niedostępny); art. 7 ust. 1 rozporządzenia (UE) 2021/695 (Horyzont Europa --- wyłącznie zastosowania cywilne; tekst niedostępny) i praktyka EDF/DIANA/EIC wobec kodów wojskowych.
  \item (\hyperref[memo:4.2.1]{4.2.1}) Rozporządzenie 2021/821 sprawdzone na tekście skonsolidowanym z 26.05.2023 r. (EN) --- potwierdzić załączniki I, II i IV w wersji aktualnej (zmiany 2023--2025) i polską wersję językową cytatów; rozporządzenia sankcyjne UE (833/2014 niedostępne; 269/2014 tylko w tekście pierwotnym, treść niepowoływana); AI Act: czy daty z art. 113 nie zostały przesunięte (projekt Digital Omnibus z 11.2025), akty wykonawcze i delegowane do 2018/1139 uwzględniające wymogi rozdziału III sekcja 2 (art. 108) oraz czy BSP Spółki podlega ocenie zgodności przez stronę trzecią (art. 6 ust. 1 lit. b); praktyka DIANA co do oceny zgodności z reżimem kontroli eksportu.
  \item (\hyperref[memo:4.2.2]{4.2.2}) Klasyfikacja produktów Spółki do pozycji załącznika I (9A012.a, 9D004.e, 9D001, 9E001, 7D002--7D004, 7E004.b, 5A002/5D002, 6A008) i ewentualnie załącznika IV (1C101/1D103 --- obniżona wykrywalność) na AKTUALNEJ wersji załączników (tekst sprawdzony: konsolidacja 26.05.2023, później zmieniana) przez specjalistę ds. kontroli eksportu; Wspólny wykaz uzbrojenia UE (ML10, ML21, ML22) niedostępny; ewentualny wniosek o wiążące wyjaśnienie z art. 10 ustawy o obrocie strategicznym.
  \item (\hyperref[memo:4.2.3]{4.2.3}) Dopuszczalność obowiązków przetrwałych byłego akcjonariusza w umowie P.S.A. (art.~300$^{1}$ §~3 KSH) --- kwestia doktrynalna, nie tekstowa; dyrektywa (UE) 2016/943 (tekst niedostępny).
  \item (\hyperref[memo:4.3.1]{4.3.1}) Opłata za ogłoszenie w MSiG (100 zł, rozporządzenie wykonawcze); stanowisko organów podatkowych co do braku PCC od umowy P.S.A. (art. 1a PCC); bieżąca praktyka PRS co do formy załączników z art.~300$^{12}$ §~3--4 KSH.
  \item (\hyperref[memo:4.3.3]{4.3.3}) Obliczenie terminu z art. 29 ustawy z 23.01.2026 r. (18 maja 2027 r.) i zakres obowiązku dostosowania z art. 34 wobec spółki zawiązanej przed 18.02.2027 r.; czy przepisy o „zarządzie” w art.~300$^{32}$ §~1$^{1}$ stosuje się wprost do rady dyrektorów.
  \item (\hyperref[memo:4.4.1]{4.4.1}) Uprawniony prawnik: czy wymóg zgody większości akcji serii na oznaczone uchwały mieści się w otwartym katalogu uprzywilejowania z art.~300$^{25}$ §~1–2 KSH (jako prawo związane z akcją), czy wymaga konstrukcji uprawnienia indywidualnego z art.~300$^{28}$ KSH.
  \item (\hyperref[memo:4.4.1]{4.4.1}) Uprawniony prawnik: czy wykładnia „beneficjentów rzeczywistych” w §~11 ust. 6 lit. b może odsyłać do definicji z art. 2 ust. 2 pkt 1 ustawy AML (więcej niż 25\% udziałów lub głosów albo inna kontrola) i czy sam brak definicji nie czyni Kryterium nieostrym wobec rejestru akcjonariuszy.
  \item (\hyperref[memo:4.4.1]{4.4.1}) Uprawniony prawnik: czy fundusz unijny (ASI, SCSp) z LP spoza Unii, którzy mają weto albo miejsce w komitecie inwestycyjnym, może zostać uznany za spółkę zależną inwestora zagranicznego „bezpośrednio lub pośrednio kontrolowaną” (art. 2 pkt 7 rozp. 2026/1386) — co przesądza, czy runda z takim funduszem wymaga uprzedniej zgody.
  \item (\hyperref[memo:4.4.2]{4.4.2}) Uprawniony prawnik: dopuszczalność wyłączenia zgody Spółki i prawa pierwszeństwa dla kategorii przeniesień (Przeniesienia Dozwolone) na podstawie art.~300$^{39}$ §~2 KSH oraz sposób ujawnienia tego wyjątku w rejestrze akcjonariuszy jako „ograniczenia co do rozporządzania akcją” (art.~300$^{33}$ §~1 pkt 10 KSH).
  \item (\hyperref[memo:4.4.2]{4.4.2}) Uprawniony prawnik: czy przepisy koncesyjne (ustawa z 2019 r.) albo wymogi zamawiających obronnych stawiają dyrektorom P.S.A. wymóg obywatelstwa, który wyjątek z §~21 ust. 3 (brzmienie V) by naruszał; oraz praktyka Komisji co do „decydującego wpływu” (art. 2 pkt 6 rozp. 2021/697) LP państwowego spoza UE na fundusz przy teście z art. 9 ust. 3.
  \item (\hyperref[memo:4.4.2]{4.4.2}) Kształt polskiego mechanizmu wykonującego rozp. 2026/1386 (progi głosów dla inwestycji mniejszościowych, rozumienie „efektywnego udziału w zarządzaniu lub kontroli” z art. 2 pkt 1 i motywów 14–15, organ, wyłączenia) po zmianie ustawy o kontroli niektórych inwestycji przed 17.01.2028 r.; czy produkty lub prace B+R Spółki są w zał. I do rozp. 2021/821 (w brzmieniu po zmianach z lat 2023–2025, tekst w repozytorium to konsolidacja na 26.05.2023 r.), w załączniku do dyrektywy 2009/43/WE albo w zał. I pkt 1–3 do rozp. 2026/1386 (półprzewodniki, kwanty, AI).
  \item (\hyperref[memo:4.4.2]{4.4.2}) Uprawniony prawnik: praktyka sądów rejestrowych wobec serii I zdefiniowanej w umowie przed emisją (art.~300$^{25}$ §~1, 300$^{103}$, 300$^{104}$ §~1 pkt 2 KSH) i wobec wyłączenia prawa poboru w samej umowie spółki (art.~300$^{106}$ §~1 KSH) zamiast uchwałą 4/5 (§~2).
  \item (T1) TFUE (art. 65, 101, 346), rozporządzenie 2019/452 (tekst), 2023/1230, 2024/2847, 2018/1139 i 2019/945 — teksty niedostępne; aktualna (zmieniona) wersja 2019/947 co do terminów i scenariuszy standardowych; ewentualne zmiany AI Act po 12.07.2024 r. przesuwające terminy dla systemów wysokiego ryzyka (w pliku tylko tekst pierwotny); czy autonomia Spółki spełnia definicję systemu AI (art. 3 pkt 1 — wnioskowanie); wytyczne Komisji 2026 do rozporządzenia 2026/877; aktualny wykaz podmiotów z rozporządzenia RM wydanego na podstawie art. 4 ustawy o kontroli niektórych inwestycji; projekt polskiej ustawy wdrażającej rozporządzenie 2026/1386 (progi, greenfield, organ).
  \item (T7) Oczekiwania funduszy obronnych (NIF, fundusze z kapitałem EDF) co do ograniczeń bezpieczeństwa — praktyka, poza stronami PFR; który limit inwestycji w spółki zagraniczne (40\% z opisu programu czy 15\% z FAQ) obowiązuje fundusze PFR Otwarte Innowacje — term sheet (zał. nr 7 do zasad naboru) i umowa funduszu z PFR.
\end{itemize}

\section*{Załącznik nr 4 --- Zmiany wymagane przez wariant B}
\addcontentsline{toc}{section}{Załącznik nr 4 --- Zmiany wymagane przez wariant B}

Jeżeli Założyciele wybiorą wariant B (odmowa zgody z powodu Kryterium bez obowiązku wskazania nabywcy), zmiana obejmuje łącznie następujące miejsca (brzmienia w \texttt{law.md} 1.2, \texttt{psa\_feedback.tex} Część 1 pkt 1 oraz w odpowiedziach \hyperref[memo:1.2.4]{1.2.4} i \hyperref[memo:1.2.5]{1.2.5}):
\begin{enumerate}
  \item §~\ref{art:dopuszczalny-nabywca} ust.~5 zdanie drugie: brak Kryterium jako podstawa odmowy zgody na zasadach §~\ref{art:zbywanie} ust.~2 lit.~b i ust.~3, chyba że WZ wyrazi zgodę;
  \item §~\ref{art:zbywanie} ust.~3: wyłączenie art.~300$^{39}$ §~3--6 KSH dla przesłanki z ust.~2 lit.~b, prawo zbycia innemu nabywcy spełniającemu Kryterium oraz \textbf{wentyl} z odpowiedzi \hyperref[memo:1.2.4]{1.2.4}: po 12 miesiącach bezskutecznego poszukiwania nabywcy akcjonariusz może żądać wskazania nabywcy przez Spółkę za Wartość Godziwą (uwaga konieczna);
  \item §~\ref{art:zbywanie} ust.~4: usunięcie odesłania do mechanizmu wskazania nabywcy;
  \item §~\ref{art:okres-ograniczenia-zbywania} ust.~1: usunięcie odesłania do §~\ref{art:zbywanie} ust.~3;
  \item §~\ref{art:dopuszczalny-nabywca} ust.~15: cena płatna jednorazowo w 30 dni zamiast na zasadach §~\ref{art:zbywanie} ust.~3;
  \item dokumenty towarzyszące: \texttt{extra.tex} Dokument nr~1 pkt~3 i Dokument nr~5, \texttt{shareholder\_agreement.tex} (odesłania), \texttt{emisja/inwestor.md}, \texttt{vc.md} sekcja 8.
\end{enumerate}

\clearpage
\section*{Załącznik nr 5 --- Co zmieniła weryfikacja na tekstach aktów}
\addcontentsline{toc}{section}{Załącznik nr 5 --- Co zmieniła weryfikacja na tekstach aktów}

\subsection*{Co zmieniła weryfikacja na tekstach aktów}
Wersja 1.1 powstała po sprawdzeniu każdej powołanej podstawy krajowej w tekście aktu (\texttt{src/legal/}). Poniżej korekty merytoryczne wobec wersji 1.0, w numeracji pytań; pełne uzasadnienia przy odpowiedziach. Oceny i klasy uwag nie zmieniły się poza miejscami wskazanymi w tej liście.
\begin{itemize}
  \item (1.1.2) Było: organ kontroli inwestycji Prezes UOKiK, art. 12a i n. bez numerów; jest: minister ds. gospodarki (art. 12c ust. 1 pkt 4), podmiot objęty ochroną wymaga działalności z art. 12d ust. 3 pkt 7 i przychodu ponad 10 mln euro (art. 12d ust. 4), progi 20\%/40\%/dominacja, nieważność z art. 12k ust. 1, kary z art. 16a; do progu przychodowego Spółka nie podlega ustawie, a ustawa nie zna kategorii dual-use.
  \item (1.1.3) Było: nabycie skuteczne z chwilą wpisu na podstawie art.~300$^{36}$ KSH; jest: art.~300$^{37}$ §~1 (nabycie z chwilą wpisu) i art.~300$^{38}$ §~1 (akcjonariusz wobec spółki), art.~300$^{36}$ dotyczy zbywalności i formy dokumentowej zbycia; zakres badania przez podmiot prowadzący rejestr to art.~300$^{34}$ §~5.
  \item (1.1.4) Było: brak squeeze-out w P.S.A. jako wiedza ogólna; jest: potwierdzone na tekście (brak odpowiednika art. 418; jedyne tryby to art.~300$^{49}$ wyłączenie i art.~300$^{50}$ ustąpienie z wykupem po wartości godziwej).
  \item (1.1.5) Było: art.~300$^{52}$ i n. [numer do sprawdzenia], wymóg koncesyjny „dwóch członków organu”; jest: art.~300$^{53}$ KSH (dyrektorzy podlegają ograniczeniom z umowy spółki) i art. 10 ust. 1 pkt 2 ustawy koncesyjnej (dwie osoby z organu zarządzającego albo członek organu i prokurent/pełnomocnik; poz. 471 tego nie zmienia).
  \item (1.1.6) Było: obowiązek Rady zgłoszenia ograniczeń do rejestru w 14 dni, skutek wpisu z art.~300$^{36}$; jest: termin 7 dni zgodny z nowym art.~300$^{33}$ §~3 KSH (od 18.02.2027 r., sankcja art. 594 §~1 pkt 2$^{1}$), podstawa skutku wpisu art.~300$^{37}$ §~1, dodatkowe informacje w rejestrze na podstawie art.~300$^{33}$ §~2.
  \item (1.2.1) Było: art. 182 KSH jako model bezwzględnie wiążący i „§~4 oddaje umowie termin, cenę i termin zapłaty”; jest: art. 182 §~2 zawiera tę samą klauzulę dyspozytywną, różnica to udział sądu rejestrowego; materię oddaje umowie art.~300$^{39}$ §~3 zd. 2, a 3-miesięczny termin wariantu A przekracza miesięczne maksimum z §~3 zd. 4 i opiera się wyłącznie na derogacji z §~2; forma dokumentowa to §~4 (dyspozytywny).
  \item (1.2.2) Było: art.~300$^{34}$ §~3--4 [do sprawdzenia] i art.~300$^{36}$; jest: art.~300$^{34}$ §~4--6 (dokumenty, badanie treści i formy, uwzględnianie ograniczeń), art.~300$^{37}$ §~1 i 300$^{38}$ §~1; od 18.02.2027 r. forma kwalifikowana zgody na wpis (art.~300$^{34}$ §~3 zd. 2).
  \item (1.2.3) Było: „sposób ustalenia ceny” z art.~300$^{39}$ §~4 i swoboda zbycia z „§~4 zd. ostatnie”; jest: art.~300$^{39}$ §~3 zd. 2 i zd. 3 („w braku tych postanowień akcja może być zbyta bez ograniczenia”); zastaw na akcjach wg art. 327--329 KC, a od 18.02.2027 r. zastaw rejestrowy bez wpisu w rejestrze akcjonariuszy (art.~300$^{37}$ §~2).
  \item (1.2.4) Było: skutek nieważności częściowej jako hipoteza z odesłaniem do §~4; jest: powrót do modelu §~3 (termin miesiąca, zd. 4; brak ceny = swoboda zbycia, zd. 3) potwierdzony na tekście.
  \item (1.2.6) Było: zmiana „§~3--6” na „§~3--5” także w §~12 ust. 4; jest: tylko w §~13 ust. 1 (§~12 ust. 4 nie zawiera tego odesłania); dodano, że §~2 literalnie obejmuje §~6, ale jego derogacja nie wiąże komornika.
  \item (1.2.8) Było: rozważyć mechanizm wskazania nabywcy w egzekucji na wzór art. 185 KSH; jest: brak odpowiednika w przepisach o P.S.A., mechanizm nierekomendowany; objęcie akcji nowej emisji wyłączone spod wpisu konstytutywnego (art.~300$^{37}$ §~2), emisja to zmiana umowy spółki (art.~300$^{103}$).
  \item (1.3.1) Było: forma dokumentowa zgody z art.~300$^{39}$ §~3; jest: §~4 (dyspozytywny, dotyczy czynności z §~3, stosowany odpowiednio do ograniczeń „w inny sposób”; Rada Dyrektorów na podstawie art. 4 §~2$^{1}$ KSH); termin 14 dni ma odpowiednik w art.~300$^{40}$ §~2.
  \item (1.3.2) Było: ogólne odesłanie do terminów postępowania; jest: 30 dni roboczych i 120 dni (art. 12h ust. 5 i 8), zawiadomienie uprzednie składane przez nabywcę przed umową zobowiązującą (art. 12f ust. 1 i 5).
  \item (1.3.5) Było: upoważnienie „walnego zgromadzenia”, test wypłacalności jako wiedza ogólna, okres upoważnienia do sprawdzenia; jest: upoważnienie w uchwale akcjonariuszy (art.~300$^{47}$ §~1 pkt 2), warunki z §~2 (pełne pokrycie, 25\%, celowy kapitał rezerwowy z kwoty art.~300$^{15}$ §~2), art.~300$^{15}$ §~4--6 przez §~3; brak limitu czasowego upoważnienia w P.S.A. (inaczej art. 362 §~1 pkt 8); prawo pierwszeństwa Spółki martwe do czasu utworzenia kapitału rezerwowego z zysku.
  \item (1.4.1) Było: ustawa o kontroli inwestycji przypisuje nabycie podmiotowi dominującemu (wiedza ogólna); jest: art. 12c ust. 6 pkt 1, art. 12e ust. 4 i art. 3 ust. 1 pkt 1 potwierdzone; NIF wspierany przez 24 państwa NATO potwierdzony (web\_nif\_about.txt); jedynym państwem stowarzyszonym z EDF jest Norwegia (web\_edf\_gowling.txt).
  \item (1.4.2) Było: definicja beneficjenta rzeczywistego i CRBR jako wiedza ogólna, ankieta bezpieczeństwa przemysłowego [Z/W]; jest: art. 2 ust. 2 pkt 1 lit. a (więcej niż 25\%, kontrola, wyższe stanowisko kierownicze), art. 58 pkt 4a (P.S.A. zgłasza do CRBR), art. 57 ust. 2 pkt 1, 58 ust. 2 pkt 2 i 64 ust. 2 pkt 2 ustawy o ochronie informacji niejawnych, art. 10 ust. 1 pkt 2 ustawy koncesyjnej.
  \item (1.5.1) Było: wpis wzmianek na żądanie Spółki z art.~300$^{35}$ i forma dokumentowa z art.~300$^{31}$ §~2 KSH; jest: art.~300$^{34}$ §~1 i 4 (wpis na żądanie, w 7 dni, na podstawie dokumentów), art.~300$^{34}$ §~5-6 (rejestr uwzględnia ograniczenia rozporządzania, ale nie egzekwuje obowiązków z pkt 11), art.~300$^{36}$ §~4 (forma dokumentowa); dodano obowiązek zgłaszania zmian z art.~300$^{33}$ §~3 KSH od 18.02.2027 r.
  \item (1.5.2) Było: art. 19 ust. 1 i 4 PIT z zastrzeżeniem numeracji; jest: art. 19 ust. 1 PIT stosowany do zbycia akcji przez art. 17 ust. 2 PIT, art. 11 ust. 2b potwierdzony; dodano, że minimum z art.~300$^{45}$ §~2 dotyczy tylko umorzenia przymusowego dopuszczalnego, gdy umowa je przewiduje (art.~300$^{45}$ §~1).
  \item (1.5.3) Było: art. 66 §~2 KC jako podstawa związania ofertą z terminem, art. 66$^{2}$ „nie ma zastosowania”, skutek z art. 1047 §~2 KPC „po zapłacie”, forma zbycia z art.~300$^{31}$ §~2; jest: art. 66 §~2 a contrario, art. 66$^{2}$ §~2 KC (oferty z terminem nie można odwołać także między przedsiębiorcami), skutek z chwilą prawomocnego nadania klauzuli wykonalności po udokumentowaniu zapłaty (art. 1047 §~2 i art. 786 §~1 KPC), art.~300$^{36}$ §~4 KSH; dodano art.~300$^{34}$ §~4 zd. 2 (oświadczenie o zobowiązaniu do przeniesienia jako podstawa wpisu) i art.~300$^{34}$ §~3 zd. 2 od 18.02.2027 r. (forma zgody na wpis) jako argument za podpisem notarialnie poświadczonym.
  \item (1.5.4) Było: art. 1157 KPC powołany ogólnie; jest: spór o kwalifikację Odejścia ma zdatność arbitrażową (art. 1157 pkt 1 KPC), ale §38 ust. 4 wybiera sąd powszechny; dodano art. 42 ust. 3 Konstytucji i art. 52 §~1 pkt 2 KP jako wzorzec dla przesłanki przestępstwa.
  \item (1.5.5) Było: art.~300$^{62}$ KSH [numer do sprawdzenia] jako podstawa odwołania dyrektora; jest: art.~300$^{74}$ §~1 KSH (odwołanie w każdym czasie, bez utraty roszczeń ze stosunku dotyczącego pełnienia funkcji) i art.~300$^{73}$ §~3 KSH (umowa może wymagać ważnych powodów).
  \item (1.6.1) Było: ścieżka rezerwowa przez umorzenie „za wynagrodzeniem równym spłacie” bez warunków; jest: umorzenie przymusowe wymaga określenia przesłanki w umowie od zawiązania (art.~300$^{45}$ §~1 KSH), spłata tylko ze środków podlegających podziałowi z testem wypłacalności (art.~300$^{44}$ §~4 w zw. z art.~300$^{15}$ §~2 i 4-6) i po wpisie umorzenia (art.~300$^{45}$ §~2 zd. 2), więc solidarna odpowiedzialność Spółki jest konieczna; nowa uwaga: art.~300$^{41}$ §~1 zd. 3 i §~2 KSH (akcje za wkład pracy — zgoda Spółki na wstąpienie i proporcjonalna spłata) nieodzwierciedlone w §17 — dodano do Brzmienia.
  \item (1.6.3) Było: umorzenie automatyczne w P.S.A. jako hipoteza [do sprawdzenia]; jest: potwierdzone art.~300$^{46}$ KSH (bez uchwały akcjonariuszy, uchwała organu zarządzającego stwierdzająca, zastrzeżenie braku środków z art.~300$^{15}$ §~2); art.~300$^{35}$ zastąpiono art.~300$^{34}$ §~1; Brzmienie ust. 2a przebudowane na umorzenie automatyczne z fallbackiem; dodano art. 1026 KC (6 miesięcy do stwierdzenia nabycia spadku) i art. 667 §~1 KPC.
  \item (1.6.4) Było: art.~300$^{38}$ §~1-3 KSH [numeracja do sprawdzenia]; jest: art.~300$^{38}$ §~3-4 KSH; dodano opcję z art.~300$^{41}$ §~3 KSH (ograniczenie podziału akcji między spadkobierców) i wpis współwłaścicieli do rejestru od 18.02.2027 r. (art.~300$^{33}$ §~1 pkt 5).
  \item (1.7.1) Było: art. 47 §~2 KRO i hipoteza o odpowiedniku art.~183$^{1}$ / 332$^{1}$ KSH dla P.S.A.; jest: art. 47$^{1}$ KRO (skuteczność umowy majątkowej wobec osób trzecich); dział o P.S.A. nie zawiera odpowiednika art.~183$^{1}$ KSH, art.~332$^{1}$ KSH uchylony — §18 ust. 2 opiera się wyłącznie na art.~300$^{38}$ §~1 KSH.
  \item (1.7.2) Było: art. 47 §~2 KRO; jest: art. 47$^{1}$ KRO.
  \item (1.7.3) Było: ustawa sankcyjna powołana ogólnie; jest: art. 1 pkt 2 (odpowiednie stosowanie art. 2 i 9 rozp. 269/2014 do listy krajowej), art. 2 (lista), art. 6 ust. 1-2 (kara do 20 mln zł); ustalono, że ustawa nie zawiera trybu zawieszenia praw korporacyjnych — tryb musi stworzyć umowa.
  \item (1.8.1) Było: reguła cenowa w braku postanowień umowy [do sprawdzenia w art.~300$^{39}$ §~4] i numery przepisów o wyłączeniu/ustąpieniu [do sprawdzenia]; jest: w braku postanowień akcja może być zbyta bez ograniczenia (art.~300$^{39}$ §~3 zd. 3, §~5), termin ustawowy miesiąc modyfikowalny przez §~2; wyłączenie: art.~300$^{49}$ §~2 w zw. z art. 266 §~3 (rzeczywista wartość w dniu doręczenia pozwu), ustąpienie: art.~300$^{50}$ §~3 (wartość godziwa ustalona przez sąd na dzień doręczenia pozwu); dodano korektor spłaty z art.~300$^{41}$ §~1 zd. 3.
  \item (2.1.1) Podstawa powołania i odwołania: było art.~300$^{53}$ KSH (przepis dotyczy w istocie ograniczeń wobec spółki), jest art.~300$^{73}$ §~3 i art.~300$^{74}$ §~1–2 KSH; wskazano dyspozytywny art.~300$^{56}$ §~3 dla dyrektora dokooptowanego w kadencji wspólnej oraz tajność głosowania (art.~300$^{99}$ §~2) z wyjątkiem uchwał poza WZ (art.~300$^{80}$ §~3).
  \item (2.1.2) Numeracja: większość i głos rozstrzygający to art.~300$^{58}$ §~4 (nie §~1–2), protokół §~5; bezskuteczność ograniczeń reprezentacji wobec osób trzecich to art.~300$^{77}$ §~2 (nie 300$^{78}$); potwierdzono art.~300$^{125}$ §~1.
  \item (2.1.3) Potwierdzono art.~300$^{54}$ i 300$^{125}$; art.~300$^{132}$ mówi literalnie o członkach zarządu (art.~300$^{133}$ rozciąga tylko na likwidatorów) — objęcie dyrektorów oznaczono jako hipotezę.
  \item (2.2.1) Było: regulamin WZ jako opcja (klasa opcjonalne, ocena zgodne); jest: art.~300$^{92}$ §~2 zd. 2 KSH stanowi, że szczegółowe zasady udziału elektronicznego „określa regulamin walnego zgromadzenia”, więc ocena „zgodne warunkowo”, klasa zalecane, z brzmieniem ust. 7; art.~300$^{100}$–300$^{102}$ doprecyzowano (protokół 300$^{100}$, zaskarżanie 300$^{101}$).
  \item (2.2.2) Było: uchwały bez formalnego zwołania z art.~300$^{89}$; jest: art.~300$^{90}$ KSH; doprecyzowano, że zgoda akcjonariusza jest warunkiem wpisu adresu e-mail do rejestru (art.~300$^{33}$ §~1 pkt 5, także po nowelizacji poz. 176 od 18.02.2027 r.), a wpisu dokonuje podmiot prowadzący rejestr na żądanie spółki (art.~300$^{34}$ §~1); legitymacja nieobecnego z art. 422 §~2 pkt 4 w zw. z art.~300$^{101}$.
  \item (2.2.3) Dodano art.~300$^{98}$ §~3 KSH: późniejsze wprowadzenie albo rozszerzenie umownego wyłączenia od głosowania (w tym proponowany ust. 11) jest zmianą uszczuplającą prawa indywidualne i wymaga zgody wszystkich akcjonariuszy, których dotyczy.
  \item (2.2.4) Było: art. 15 KSH obejmuje dyrektora P.S.A. z mocy ustawy (nieważność bez zgody WZ); jest: art. 15 §~1 nie wymienia dyrektora i przepisy o P.S.A. nie odsyłają do niego, więc wymóg zgody WZ na kredyt, pożyczkę i poręczenie z dyrektorem wpisano do proponowanego ust. 5 jako wymóg umowny; potwierdzono większości art. 506 §~1, 541 §~1, 575, 577 §~1 pkt 1 oraz brak opt-out dla art.~300$^{81}$ pkt 2 (przedsiębiorstwo, ZCP).
  \item (2.3.1) Podstawa powoływania dyrektorów: art.~300$^{73}$ §~3 zamiast 300$^{53}$; dodano art.~300$^{80}$ §~3 (brak wymogu tajności przy uchwale o powołaniu poza WZ).
  \item (2.3.2) Było: art.~300$^{57}$ jako podstawa uchwał pisemnych i zdalnych; jest: art.~300$^{58}$ §~1–3 (300$^{57}$ dotyczy regulaminu i komitetów); numeracja kworum §~2, większości §~4, protokołu §~5.
  \item (2.3.3) Potwierdzono numerację art.~300$^{79}$ (§~1 pełnomocnik, §~3 dyrektor niewykonawczy jako opcja umowna, §~4 jedyny akcjonariusz z art.~300$^{14}$); zgoda WZ na pożyczkę dla dyrektora nie wynika pewnie z art. 15 KSH — odesłano do wymogu umownego.
  \item (2.3.4) Było: doktrynalna koncepcja „ukrytej wypłaty” na tle art.~300$^{15}$ i 300$^{22}$; jest: wprost art.~300$^{21}$ KSH (wartość świadczeń dla akcjonariuszy i spółek powiązanych nie może przekraczać wartości godziwej świadczenia wzajemnego) jako granica ustawowa; ust. 2 i brzmienie przebudowane wokół „wartości godziwej świadczenia wzajemnego” z rozszerzeniem na podmioty powiązane spoza przepisu; art. 394 §~1 (2/3 głosów, 1/10 wpłaconego kapitału, 2 lata) i PIT art. 23m–23zf (25\%) potwierdzone.
  \item (2.4.1) Potwierdzono ustawę koncesyjną z 2019 r. (art. 7 ust. 1 pkt 2, art. 3 ust. 1 pkt 12) i ustawę z 2000 r. po noweli poz. 471 (organ kontroli obrotu: minister właściwy do spraw gospodarki); kwalifikacja licencji dla spółki celowej jako „obrotu technologią” oznaczona jako hipoteza.
  \item (2.4.2) Ustawa o ochronie informacji niejawnych doprecyzowana do art. 54 ust. 1–2 (świadectwo bezpieczeństwa przemysłowego ABW/SKW), ustawa o kontroli niektórych inwestycji do art. 4 ust. 1 pkt 7 (wyroby i technologia wojskowa jako podstawa uznania za podmiot objęty ochroną).
  \item (2.4.3) Dodano art.~300$^{55}$ §~3 KSH: ustawowa zgoda na działalność konkurencyjną dyrektora należy do organu powołującego (WZ), chyba że umowa stanowi inaczej, a zakaz obejmuje udział od 10\% w konkurencyjnej spółce kapitałowej — argument za progiem 10\% i kompetencją WZ; art.~300$^{39}$ (rozporządzanie akcją) zastąpiono art.~300$^{33}$ §~1 pkt 11 (obowiązki związane z akcją w rejestrze).
  \item (2.4.4) Doprecyzowano art.~300$^{81}$ pkt 2 (zbycie i wydzierżawienie przedsiębiorstwa lub ZCP oraz obciążenie ograniczonym prawem rzeczowym — bez opt-out, nieważność z art. 17 §~1, większość 3/4 z art.~300$^{98}$ §~2 pkt 2) i rozszerzono wyjątek w proponowanym ust. 2a.
  \item (3.1.1) Było: oświadczenia przy wpisie z art.~300$^{13}$ §~2, odpowiedzialność dyrektorów za zawyżenie „jeżeli wiedząc zgłosili spółkę” i art.~300$^{123}$ jako odpowiedzialność karna; jest: art.~300$^{12}$ §~3 pkt 2–3 i §~4, odpowiedzialność dyrektorów solidarna „chyba że nie ponoszą winy” (art.~300$^{10}$ §~1–2), art.~300$^{123}$ to odpowiedzialność cywilna wobec wierzycieli, karna – art. 587 KSH; dodano koszt objęcia z art. 22 ust. 1e pkt 3 PIT.
  \item (3.1.2) Było: korzyść z niskiego kapitału akcyjnego to niższy PCC 0,5\%; jest: ustawa o PCC nie wymienia P.S.A. wśród spółek kapitałowych (art. 1a pkt 2) i opiera podstawę na kapitale zakładowym, więc korzyść nie istnieje; dodano art. 36 ust. 2aa ustawy o rachunkowości, art. 21 ust. 1 pkt 109 PIT (brak zwolnienia dla wkładu IP) i pełny test wypłat z art.~300$^{15}$ §~4–6 KSH; ocena i klasa bez zmian.
  \item (3.1.3) Było: rekomendacja klauzul przeniesienia praw do „utworów przyszłych określonych rodzajowo”; jest: art. 41 ust. 3 pr. aut. unieważnia także umowy o wszystkie utwory określonego rodzaju, więc klauzule mają obejmować utwory powstałe w wykonaniu danej umowy; dodano art. 9 pr. aut. (współtwórcy), art. 11 ust. 3 PWP (obejmuje B2B) i art. 72 ust. 1 PWP.
  \item (3.1.4) Było: zgłoszenie i lista akcjonariuszy z art.~300$^{13}$; jest: art.~300$^{12}$ §~2 pkt 7, §~3–4; dodano art.~300$^{9}$ §~3 (przypisanie numerów akcji do wkładu jako odstępstwo od równomiernego zaliczania) i art.~300$^{5}$ §~2.
  \item (3.2.1) Było: „P.S.A. nie zna kapitału docelowego ani delegacji kompetencji emisyjnej”, ocena zgodne warunkowo, klasa zalecane; jest: KSH przewiduje upoważnienie rady do emisji (art.~300$^{110}$–300$^{113}$) i warunkową emisję akcji (art.~300$^{114}$–300$^{118}$), a prawo poboru jest dyspozytywne (art.~300$^{106}$ §~1) – uchwała ramowa wykonywana przez Radę jest konstrukcją hybrydową; ocena ryzyko, klasa konieczne, nowe brzmienie ust. 3 i 5.
  \item (3.2.1) Było: hipoteza sześciomiesięcznego terminu w art.~300$^{107}$; jest: art.~300$^{107}$ terminu nie zawiera, termin z art.~300$^{102}$ §~2 dotyczy zmiany umowy i prawdopodobnie nie obowiązuje przy emisji z art.~300$^{103}$ (do potwierdzenia); art.~300$^{106}$ §~2 nie wymaga zapowiedzi w porządku obrad.
  \item (3.2.4) Było: rozszerzenie art. 24 ust. 11 PIT na P.S.A. „do sprawdzenia”; jest: potwierdzone w art. 24 ust. 12b PIT; warunkowa emisja odpowiada konstrukcji ust. 11b; PCC od emisji – wątpliwe, bo art. 1a pkt 2 PCC nie obejmuje P.S.A.
  \item (3.3.2) Było: brak w P.S.A. odpowiednika art. 354 KSH, powoływanie dyrektorów z art.~300$^{60}$; jest: uprawnienia indywidualne z art.~300$^{28}$ KSH, powoływanie z art.~300$^{73}$ §~3 KSH, zakaz uprzywilejowania przez Radę z art.~300$^{110}$ §~5, zaczątek zgody serii w art.~300$^{98}$ §~4; ocena i klasa bez zmian.
  \item (3.3.3) Było: dyspozytywność prawa poboru jako hipoteza; jest: art.~300$^{106}$ §~1 KSH wprost dopuszcza odmienną regulację w umowie spółki lub uchwale – umowa może wyłączyć prawo poboru wobec inwestorów Rundy albo uregulować pro-rata; dodano wariant do decyzji Założycieli.
  \item (3.3.4) Było: nieograniczalność reprezentacji z art.~300$^{63}$ §~2 i 300$^{76}$; jest: art.~300$^{77}$ §~2 KSH; dodano, że zastaw na przedsiębiorstwie, warranty, obligacje zamienne i umowy o prawa do objęcia akcji wymagają uchwały z ustawy (art.~300$^{81}$ pkt 2 i 4, art.~300$^{114}$ §~3), więc bez niej czynność jest nieważna (art. 17 §~1), nie tylko naruszeniem umowy.
  \item (3.4.1) Było: wątpliwość, czy rejestr akcjonariuszy przyjmie ograniczenia serii F; jest: rejestr zawiera z ustawy ograniczenia rozporządzania i obowiązki związane z akcją (art.~300$^{33}$ §~1 pkt 10–11), umowa może dodać dane (§~2), od 18.02.2027 r. zgłoszenie zmian w 7 dni (§~3, poz. 176); pozostaje pytanie o praktykę.
  \item (3.4.2) Było: odpowiednik art. 419 KSH „do sprawdzenia”, ochrona osobista jako „prawa przyznane osobiście”; jest: art.~300$^{98}$ §~3 (zgoda dotkniętych przy zwiększeniu świadczeń lub uszczupleniu praw indywidualnych) i §~4–5 (oddzielne głosowanie grup akcji) istnieją; dodano akcje założycielskie z art.~300$^{26}$; brzmienie §~38 ust. 3 uzupełnione o odesłanie do art.~300$^{98}$ §~3–4.
  \item (4.1.1) Było: kategorie WT XXI/XXII (oprogramowanie, technologia) w wykazie krajowym; jest: wykaz ma kategorie WT I--XIV, oprogramowanie i technologia ujęte przez definicje pkt 12 i 15 załącznika, WT V ust. 3 (systemy kierowania i kontroli lotu BSP), WT XI ust. 2, WT XIII ust. 3--4 i art. 3 ust. 1 pkt 12 ustawy; tytuł rozporządzenia poprawiony.
  \item (4.1.1) Było: zakres zmian poz. 471 w ustawie koncesyjnej nieznany, data 8 albo 22 kwietnia 2026 r.; jest: art. 2 noweli zmienia tylko odnośnik, art. 44 ust. 2aa (głębokość oznakowania broni palnej) i odesłania, w życie 8.04.2026 r.; zmiany w ustawie o obrocie strategicznym od 22.04.2026 r. (art. 1 pkt 17 lit. b od 8.10.2026 r.).
  \item (4.1.1) Było: definicje wytwarzania i obrotu w art. 3 do sprawdzenia, wymogi wobec organów nieznane; jest: ustawa nie definiuje wytwarzania ani obrotu; koncesja tylko na obrót technologią (art. 7 ust. 1 pkt 2); art. 10 ust. 1 pkt 2 (dwóch członków organu z obywatelstwem UE/EOG/CH, szkoleniem, badaniami; niekaralność także wspólników i akcjonariuszy co najmniej 20\%), art. 17 (dane akcjonariuszy we wniosku), sankcja art. 133 (1--10 lat); brak wyłączenia B+R (art. 4).
  \item (4.1.2) Było: art.~300$^{97}$ KSH jako podstawa zmiany umowy; jest: art.~300$^{98}$ §~2 pkt 1 (3/4), art.~300$^{100}$ §~2 (protokół notarialny), art.~300$^{102}$ §~1--2 (wpis, 6 miesięcy); opłata 250 zł potwierdzona (art. 55 ust. 1 u.k.s.c.); hipoteza o wymogu ujawnienia w KRS: ustawa go nie zawiera (art. 17 ust. 1 pkt 2), praktyka pozostaje do sprawdzenia.
  \item (4.1.3) Było: nazwy kodów niezweryfikowane (63.10.D, 29.10.E i inne pod znakiem zapytania), ocena „zgodne warunkowo”, klasa „zalecane”; jest: wszystkie 19 kodów istnieje w PKD 2025, a nazwy w §~4 ust. 5--6 są identyczne z urzędowymi, ocena „zgodne”, klasa „bez zmian”.
  \item (4.2.1) Uzupełniono o treść ustawy o obrocie strategicznym po nowelizacji: art. 9 ust. 2 pkt 11 (osoba odpowiedzialna za koordynację kontroli obrotu), art. 11 (WSK z certyfikacją), art. 33 ust. 1 i 2a (sankcje); art.~300$^{1}$ §~3 KSH opisany zgodnie z treścią (świadczenia określone w umowie).
  \item (4.2.2) Było: „stanowisko” Departamentu MRiT; jest: wiążące wyjaśnienie organu kontroli obrotu (ministra ds. gospodarki) w sprawie konieczności zezwolenia, art. 10 ust. 1 (3 miesiące, przedłużalne do 6), catch-all art. 17a ust. 1 pkt 2.
  \item (4.2.3) Było: zakaz ujawniania tajemnic po wygaśnięciu mandatu w art.~300$^{54}$ KSH; jest: art.~300$^{55}$ §~2 KSH (art.~300$^{54}$ to staranność i lojalność); dodano art. 11 ust. 4 uznk (naruszenie obowiązku z czynności prawnej) i art. 54 ustawy o ochronie informacji niejawnych.
  \item (4.3.1) Było: PCC 0,5\% od kapitału akcyjnego; jest: P.S.A. nie mieści się w definicjach spółki kapitałowej ani osobowej z art. 1a PCC, a podstawą przy spółce kapitałowej jest kapitał zakładowy --- umowa P.S.A. według wykładni literalnej nie podlega PCC (do potwierdzenia przez doradcę podatkowego).
  \item (4.3.1) Było: lista akcjonariuszy „do sprawdzenia”, braki usuwane „w 7 dni pod rygorem zwrotu”; jest: lista akcjonariuszy wymagana z mocy art.~300$^{12}$ §~4 KSH; braki formalne art. 130 §~1--3 KPC (termin tygodniowy), braki merytoryczne art. 165 KSH (termin pod rygorem odmowy wpisu), art. 164 §~3 (drobne uchybienia); dodano art. 19a ust. 5a ustawy o KRS (pełnomocnik do doręczeń dla dyrektora z adresem poza UE) i art.~300$^{11}$ §~2 (reprezentacja spółki w organizacji).
  \item (4.3.3) Było: ocena na podstawie omówień (KRS „art. 38 do sprawdzenia”, termin do 17.05.2027, zniesienie akcji imiennych, 7-dniowy termin dla akcjonariusza); jest: pełny wykaz zmian z tekstu --- art.~300$^{32}$ §~1$^{1}$--1$^{3}$ i §~3, art.~300$^{33}$ §~1 pkt 5--6 i §~3 (7 dni dla zarządu, obejmuje ograniczenia i obowiązki z pkt 10--11), art.~300$^{34}$ §~1, §~3 zd. 2 (forma zgody), §~9, art.~300$^{35}$ §~1$^{1}$ i §~4, art.~300$^{37}$ §~2, art. 594 §~1 pkt 2$^{1}$ (grzywna do 20 000 zł), KRS art. 38 pkt 8a lit. j, art. 25da, art. 10 ust. 4a pkt 4, Prawo o notariacie art. 90a §~1, u.o.i.f. art. 83a ust. 4ac; przejściowe art. 29 (3 miesiące, do 18.05.2027), art. 30, art. 34 (dostosowanie umów w 2 lata, do 18.02.2029); art. 35 (18.02.2027; art. 28 i 33 od 28.02.2026); zniesienie akcji imiennych dotyczy S.A.
  \item (4.4.1) Było: przepis o akcjach uprzywilejowanych i uprawnieniach indywidualnych jako art.~300$^{26}$ [numer do sprawdzenia]; jest: art.~300$^{25}$ §~1–2 (akcje uprzywilejowane, katalog otwarty) i art.~300$^{28}$ §~1–2 KSH (uprawnienia indywidualne, wygasające z utratą statusu akcjonariusza) — art.~300$^{26}$ dotyczy akcji założycielskich.
  \item (4.4.1) Było: każdy fundusz zażąda katalogu spraw wymagających zgody serii niezależnie od liczby głosów; jest: polska praktyka zna dwa równorzędne sposoby — głos funduszu „za” albo bardzo dużą większość 75–80\% (web\_pfr\_umowa\_inwestycyjna.txt) — więc próg §~25 należy do standardu, a katalog zgody serii jest żądany przy udziale poniżej progu blokującego.
  \item (4.4.1) Było: fundusze D „same podlegają wymogom właścicielskim” (NIF, EDF); jest: NIF stawia spółkom warunek siedziby w jednym z 24 państw NATO będących jego inwestorami, nie test właścicielski (web\_nif\_about.txt); test kontroli państwa trzeciego dotyczy EDF (web\_edf\_gowling.txt).
  \item (4.4.1) Było: oświadczenie własne funduszu zamiast look-through jako praktyka [W]; jest: potwierdzone kierunkiem wzorców NVCA 2025 (oświadczenie inwestora „not a person of concern” w reżimach OISP/DSP, weryfikacja LP po stronie funduszu, web\_nvca\_2025\_foley.txt).
  \item (4.4.2) Było: hipoteza [?], czy uprzywilejowanie serii I można ustalić z góry i czy upoważnienie koliduje z prawem poboru; jest: na tekście KSH — uprzywilejowanie musi być w umowie spółki (300$^{25}$ §~1), uchwała o emisji je wskazuje (300$^{104}$ §~1 pkt 2), prawo poboru może wyłączyć sama umowa spółki (300$^{106}$ §~1), a 4/5 dotyczy tylko pozbawienia uchwałą (§~2); otwarta pozostaje praktyka sądów rejestrowych.
\end{itemize}

\subsection*{Co zmieniła weryfikacja na tekstach prawa Unii i materiałach PFR (wersja 1.2)}
Wersja 1.2 sprawdza tezy o prawie Unii na tekstach pobranych do \texttt{src/legal/eu/} (EDF po polsku; 2021/821, 269/2014, 2019/947, AI Act, 2026/877 i nowe 2026/1386 po angielsku) oraz tezy o praktyce PFR Ventures na stronach z \texttt{src/legal/pfr/}. Korekty merytoryczne wobec wersji 1.1:
\begin{itemize}
  \item (1.1.1) Było: art. 9 EDF „zakazuje kontroli” przez państwo trzecie, jedynym państwem stowarzyszonym Norwegia (za web\_edf\_gowling.txt); jest: na tekście PL — art. 9 ust. 1–3 (siedziba, zarządcze struktury wykonawcze, brak kontroli), definicje kontroli i podmiotu z niestowarzyszonego państwa trzeciego (art. 2 pkt 6, 7, 24), derogacja za gwarancjami państwa członkowskiego (ust. 4) i obowiązek zgłaszania zmian (ust. 7); art. 5 określa państwa stowarzyszone ogólnie jako członków EFTA należących do EOG (stąd brak USA/UK/Kanady/Turcji), Norwegia pozostaje tezą ze źródła internetowego; od 17.01.2028 r. rozp. 2026/1386 obejmuje obowiązkowym zezwoleniem inwestorów z NATO/OECD spoza UE, ale tylko przy nabyciu dającym skuteczny udział w zarządzaniu lub kontroli (art. 2 pkt 1, motywy 14–15; progi głosów może dodać ustawa, motyw 20) i przy celu, który opracowuje, produkuje lub wprowadza do obrotu produkty z wykazów (art. 4 ust. 15 lit. a–b).
  \item (1.1.2) Było: jedno zdanie o rozp. 2019/452 jako ramach współpracy [W]; jest: dwa stany prawne — do 16.01.2028 r. (2019/452 + ustawa z progiem 10 mln euro i wyłączeniem OECD, art. 12a ust. 1 pkt 1) i od 17.01.2028 r. (rozp. 2026/1386: obowiązkowy mechanizm w każdym państwie, art. 3 ust. 1; obowiązkowe zezwolenie dla celów opracowujących, produkujących lub wprowadzających do obrotu produkty dual-use z zał. I do 2021/821 i produkty obronne z załącznika do dyr. 2009/43/WE, art. 4 ust. 15 lit. a–b, bez progu przychodowego, z możliwością krajowego progu głosów (motyw 20) i poza zakresem wewnętrznych restrukturyzacji (art. 1 ust. 5 lit. b); inwestor zagraniczny = osoba bez obywatelstwa UE albo podmiot prawa państwa trzeciego, także przez spółkę zależną w UE, art. 2 pkt 1, 5, 7, a skuteczny udział w zarządzaniu także bez decydującego wpływu, motyw 15; terminy 45 dni (art. 4 ust. 2), kontrola ex post 24 mies. dla niezgłoszonych podlegających zezwoleniu (ust. 5) i 15 mies.–5 lat dla pozostałych (ust. 4) — wcześniej zakresy ust. 4 i 5 opisano nieprecyzyjnie; współpraca art. 5–6, 9, 11 ust. 3–5 (z terminami przy żądaniu informacji i przedłużeniem), 12; EDF w zał. II; drony i AI w zał. III; kryteria z art. 19 ust. 2 poprawione — nie wymieniają kontroli rządu państwa trzeciego (to przesłanka notyfikacji z art. 5 ust. 1 lit. a), lecz cele polityki i zdolności wojskowe państwa trzeciego, odmowy, sankcje, prawo wywiadowcze, nieprzejrzystą strukturę; objęcie Spółki art. 4 ust. 15 zależy od klasyfikacji produktów [?], a nie „wynika” z §~3 i §~28; przepisy stosowane od 16.07.2026 r., art. 31).
  \item (1.2.7) Było: rozporządzenia sankcyjne 269/2014 i 833/2014 jako wiedza ogólna; jest: art. 1 lit. f–g i art. 2 ust. 1–2 rozp. 269/2014 potwierdzone na tekście pierwotnym EN (zamrożenie obejmuje akcje); art. 1 pkt 2 ustawy z 13.04.2022 r. odsyła do art. 2 (zamrożenie, zakaz udostępniania) i art. 9 (zakaz obchodzenia) rozp. 269/2014 — wcześniej art. 9 błędnie opisano jako zamrożenie; 833/2014 nadal [W].
  \item (1.4.1) Było: test EDF cytowany po angielsku ze strony kancelarii, Norwegia jako jedyne państwo stowarzyszone, DIANA i NIF „wymagają siedziby w NATO” [Z]; jest: art. 9 ust. 3 i definicja kontroli (art. 2 pkt 6) z tekstu PL, art. 2 pkt 24 (fundusz z zarządem w UE nie jest podmiotem z państwa trzeciego), art. 9 ust. 4 i 7 (gwarancje, zgłoszenie zmian) oraz test „spółki zależnej inwestora zagranicznego” z art. 2 pkt 7 rozp. 2026/1386 (motyw 18), żądanie informacji o strukturze własności od inwestora i łańcucha kontroli (art. 15 ust. 3 w zw. z ust. 1 lit. b — wcześniej błędnie sam ust. 1, który dotyczy treści notyfikacji między państwami), nieprzejrzysta struktura (art. 2 pkt 8, art. 19 ust. 2 lit. e), notyfikacja przy kontroli rządowej tylko dla celów z art. 4 ust. 15; DIANA obniżona do [W] (brak źródła), NIF — siedziba w jednym z 24 państw-inwestorów (web\_nif\_about.txt).
  \item (1.4.2) Było: test sankcyjny bez podstawy w tekście; jest: art. 2 ust. 1–2 (zamrożenie i zakaz udostępniania „bezpośrednio lub pośrednio”), art. 1 lit. g pkt iii (środki finansowe obejmują akcje), art. 9 (świadome i umyślne obejście) rozp. 269/2014 — uzasadnienie objęcia testem kontrolujących i beneficjentów funduszu; od 2028 r. sankcje wobec inwestora uruchamiają notyfikację z art. 5 ust. 1 lit. b rozp. 2026/1386 (przy celu z art. 4 ust. 15), a kryterium oceny są powody sankcji (art. 19 ust. 2 lit. b) i obchodzenie sankcji (lit. a pkt vi) — wcześniej przywołano wyłącznie lit. a pkt vi.
  \item (1.4.3) Było: progi 5/25\%/50\%/10\% „zbliżone do praktyki funduszy obronnych i programów PFR” [W], a po pierwszej weryfikacji W2 — „wkład prywatny min. 40\% w Biznest, OI i KOFFI, PFR do 60\% w OI i KOFFI” i teza, że fundusz PFR „spełni warunki (a)–(d) bez trudu”; jest: strony PFR Ventures nie podają takich progów; Starter — PFR do 80\%, wkład prywatny min. 20\%; KOFFI i OI w strukturze standardowej — wkład prywatny min. 40\% kapitalizacji, PFR do 60\%; OI w modelu koinwestycyjnym — PFR do 97\% kapitalizacji, 40\% kapitału prywatnego deal by deal w spółkach; Biznest — co najmniej 40\% wartości każdej inwestycji od inwestorów prywatnych, wkłady aniołów poza funduszem (fundusz = PFR + zarządzający min. 2\%); KOFFI: oferty z niewielką liczbą inwestorów niepożądane, ZASI wymagany przy umowie; OI obejmuje wprost technologie obronne i dual-use; fundusz PFR z inwestorami spełniającymi Kryterium spełnia je wprost (w modelach koinwestycyjnych nie spełniłby warunku rozproszenia, a koinwestorzy przechodzą test samodzielnie); dodano rozbieżność strony OI co do limitu geograficznego (opis: do 40\% w spółki zagraniczne z polskimi founderami; FAQ: min. 85\% Polska i do 15\% zagranica, w jednej odpowiedzi UE/EFTA/EOG+UK) z [Z] dla brzmienia strony; warunki (d)–(e) powiązano z art. 2 pkt 8 i art. 19 ust. 2 lit. e rozp. 2026/1386; progi pozostają propozycją [W].
  \item (1.7.3) Było: art. 2 ust. 1–2 rozp. 269/2014 jako wiedza ogólna, rozp. 2021/821 jako [Z] bez przepisu, zbycie po derogacji jako pewna ścieżka i teza, że ustawa z 13.04.2022 odsyła w całości do rozporządzenia; jest: na tekście pierwotnym EN potwierdzono, że zamrożenie obejmuje akcje i dywidendy (art. 1 lit. g pkt iii–iv), definicję zamrożenia (art. 1 lit. f), zakaz udostępniania z dopuszczeniem uznania rachunku zamrożonego (art. 2 ust. 2, art. 7 ust. 1) i zakaz świadomego i celowego obejścia (art. 9); tekst pierwotny zna derogacje tylko z art. 4–6, więc derogacja na zbycie akcji jest do sprawdzenia w wersji skonsolidowanej [?]; dodano zarząd przymusowy z art. 6a, 6b i 6da ustawy z 13.04.2022 (zarządca wykonuje prawa z akcji i może je zbyć ze skutkiem sprzedaży egzekucyjnej, art.~300$^{39}$ §~6 KSH) i zastrzeżenie dla zarządcy w Brzmieniu ust. 11a; zakres 2021/821 oparto na art. 1 i art. 2 pkt 2 lit. d oraz pkt 10 lit. c (tekst EN); ocena prawa głosu pozostaje [W].
  \item (2.4.1) Było: 2026/877 powołane ogólnie (progi, okres przejściowy, grant-back) z wiedzy ogólnej; jest: art. 1 ust. 1 lit. b pkt vii i lit. c pkt i, art. 2, 3, 5 ust. 1 lit. a–b, 8 lit. e, 9, 10, 11 z tekstu EN — ograniczeniem wyłączonym jest także cesja ulepszeń, więc wariant „własność ulepszeń niewydzielnych” również opiera się na ocenie indywidualnej (odpowiednio poprawiono Rekomendację i uwagę, które nazywały go wariantem zgodnym z rozporządzeniem), a prawo pierwszeństwa nabycia ulepszeń oznaczono jako niepewne; okres przejściowy dotyczy tylko umów obowiązujących 30.04.2026 r.; w 2021/821 rozdzielono eksport technologii (art. 2 pkt 2 lit. d, art. 3 ust. 1, art. 4 ust. 1) od pomocy technicznej (art. 8, bez wymogu zezwolenia wobec państw EU001 — art. 8 ust. 3 lit. a) i transferu w UE (art. 11 ust. 1, załącznik IV); dla Turcji wskazano zezwolenie indywidualne albo globalne, bo EU002–EU006 obejmują ją tylko wąsko (tekst EN).
  \item (2.4.2) Było: art. 9 EDF i EU001 na podstawie regulations.md, utrata kwalifikowalności przez Spółkę i w AGILE jako skutek art. 9 ust. 3; jest: art. 5 (państwa stowarzyszone = EFTA/EOG, bez wskazania państw) i art. 9 ust. 3, 4 lit. c, 6, 7 z tekstu PL — kontrola partnera odbiera kwalifikowalność Przedsięwzięciu, Spółce tylko przez art. 20 ust. 4 (powiadomienie Komisji i ewentualny zwrot wsparcia przy przeniesieniu wyników lub licencji wyłącznej), AGILE oznaczono jako niepewne; EU001 z załącznika II sekcja A do 2021/821 (tekst EN 2023), a brzmienie ust. 4 zawężono z „unijnego generalnego zezwolenia” do EU001, bo EU002–EU006 obejmują Turcję; dodano skutek rozporządzenia 2026/1386 dla udziału partnera spoza UE w Przedsięwzięciu (istniejąca spółka: uprzednie zezwolenie od 17.01.2028 r., art. 4 ust. 9 i 15 lit. a–b, transakcje zakończone wcześniej według dotychczasowych przepisów, art. 30 ust. 3; greenfield poza zakresem obowiązkowym, art. 4 ust. 17).
  \item (3.3.2) Dodano skutek rozporządzenia 2026/1386 dla Rundy z inwestorem zagranicznym (także z EOG/NATO) zamykanej po 17.01.2028 r. (art. 30 ust. 3): obowiązkowe uprzednie zezwolenie, gdy Spółka rozwija lub komercjalizuje produkty z zał. I do 2021/821 albo z załącznika do dyrektywy 2009/43/WE (art. 4 ust. 15 lit. a–b), lecz tylko przy inwestycji dającej efektywny udział w zarządzaniu lub kontroli (art. 2 pkt 1, motywy 14–15; bez inwestycji portfelowych, bez progów w rozporządzeniu); warunek zawieszający, 45 dni kalendarzowych oceny wstępnej, bez terminu dla postępowania pogłębionego, a przy udziale w EDF notyfikacja w mechanizmie współpracy i terminy z art. 11; standard 1x non-participating i anti-dilution (średnia ważona) potwierdzony na artykule z portalu startup.pfr.pl (autor zewnętrzny, nie stanowisko PFR).
  \item (3.3.4) Było: art. 9 EDF ogólnie (W); jest (tekst PL): art. 9 ust. 3 – odbiorca EDF nie może podlegać kontroli niestowarzyszonego państwa trzeciego (państwa stowarzyszone to tylko EFTA w EOG, art. 5), co przemawia za odesłaniem §~31 ust. 3 do Kryterium; art. 9 ust. 4 lit. c to wyłącznie treść gwarancji dla podmiotu kontrolowanego, a ograniczenia wyników wynikają z art. 20 ust. 3–4 i art. 23 ust. 2 i 4 (zakaz kontroli z państw niestowarzyszonych, uprzednie powiadomienie Komisji o przeniesieniu własności lub licencji wyłącznej, zwrot wsparcia), eksport poza rozporządzeniem (art. 20 ust. 9, art. 23 ust. 3); dodano definicję inwestycji zagranicznej z art. 2 pkt 1 rozporządzenia 2026/1386 wobec finansowania z konwersją lub warrantem.
  \item (3.4.2) Teza o protective provisions inwestora VC potwierdzona na artykule z portalu startup.pfr.pl o umowie inwestycyjnej (autor z kancelarii, nie stanowisko PFR): głos „za” funduszu albo bardzo duża większość (np. 75–80\%), członek rady z prawem zgody na istotne czynności; artykuł lokuje te mechanizmy w umowie inwestycyjnej — Z z nazwą pliku.
  \item (4.1.1) Było: rozporządzenie 2019/947 powołane z wiedzy ogólnej, a w wersji 1.2 m.in. „kategoria otwarta tylko przy <25 kg, VLOS i 120 m”, „operacja autonomiczna obejmuje warstwę autonomii i roje Spółki” i „stosowanie od 1.07.2020” (Z); jest po weryfikacji sceptycznej na tekście EN (pierwotnym): kategorie otwarta/szczególna/certyfikowana (art. 3--6, w tym art. 6 ust. 2), pełny katalog warunków kategorii otwartej (art. 4 ust. 1 lit. a--f: klasa, <25 kg, zakaz lotu nad zgromadzeniami, VLOS, 120 m, bez towarów niebezpiecznych i zrzutów), zezwolenie na operację (art. 5, 12), oświadczenie w scenariuszu standardowym dopiero po uzupełnieniu dodatku 1 (art. 23 ust. 2); operacja autonomiczna (art. 2 pkt 17) tylko bez możliwości interwencji pilota --- rój kontrolowany przez pilota nią nie jest; poza kategorią otwartą lokuje ją wymóg zdolności utrzymania kontroli (UAS.OPEN.060 ust. 2 lit. d w zw. z art. 5 ust. 1), a UAS.SPEC.050 ust. 1 lit. b przewiduje ją w kategorii szczególnej; rejestracja operatora (art. 14 ust. 5--6, także próg 80 J i wyjątek dla zabawek); data 1.07.2020 wg tekstu pierwotnego (art. 23 ust. 1), przesunięta na 31.12.2020 rozporządzeniem 2020/746 (W); 2019/947 nie zawiera definicji roju ani wyłączenia wojskowego (to 2018/1139, nadal W).
  \item (4.1.2) Było: „dla programów EDF, EUDIS, DIANA i EIC (ścieżka dual-use) kody wojskowe są neutralne albo korzystne” (bez oznaczenia); jest: dla EDF/EUDIS/DIANA --- W (praktyka; rozporządzenie 2021/697 nie odnosi się do PKD), a EIC wyodrębniony: Horyzont Europa finansuje wyłącznie zastosowania cywilne (art. 7 ust. 1 rozporządzenia 2021/695, W), więc kody wojskowe są tam co najwyżej neutralne.
  \item (4.2.1) Było: art. 2, 4, 12 i załączniki 2021/821 z wiedzy ogólnej, „załącznik I zmienia się co roku”; jest: z tekstu EN --- wywóz obejmuje elektroniczne udostępnienie i ustne przekazanie technologii osobom poza obszarem celnym (art. 2 pkt 2 lit. d; data room i chmura poza UE są wywozem tylko w zakresie produktów podwójnego zastosowania), definicja ICP (art. 2 pkt 21) i wymóg ICP przy zezwoleniu globalnym, chyba że organ uzna go za zbędny (art. 12 ust. 4 akapit trzeci), oraz przy EU007 (sekcja G część 3 pkt 3) --- wersja 1.2 twierdziła bezwarunkowo „warunek dostępu do zezwolenia globalnego”; catch-all art. 4 ust. 1 lit. a--c (lit. b opisana zgodnie z tekstem: wbudowanie, użycie urządzeń produkcyjnych/testowych/analitycznych, półprodukty) z obowiązkiem powiadomienia (ust. 2), pomoc techniczna art. 8 ust. 1--3, transfer wewnątrzunijny art. 11 ust. 1 i obowiązek oznaczania dokumentów handlowych (ust. 9), ewidencja 5/3 lata (art. 27), EU001 (państwa, powiadomienie o pierwszym użyciu w 30 dni) i EU007 (tylko spółki zależne/siostrzane w całości kontrolowane, wyłączenie odbiorców wojskowych i służb); aktualizacja załącznika I aktami delegowanymi (art. 17), a nie „co roku” (to praktyka); odesłania art. 6 i 17a ustawy krajowej do numerów artykułów rozporządzenia zgadzają się z tekstem (art. 17a ust. 1 powołuje ich więcej --- dopisano „m.in.”); wtrącenie „EDF, DIANA wymagają wykazania świadomości reżimu” zastąpiono tezą z tekstu (art. 9 ust. 4 akapit drugi lit. c rozporządzenia 2021/697 --- gwarancje zakazu wywozu rezultatów bez zgody państwa siedziby) i W dla praktyki.
  \item (4.2.1) Dodano AI Act (2024/1689), po weryfikacji sceptycznej skorygowany: oprogramowanie autonomii jest systemem AI (art. 3 pkt 1) tylko o ile „wnioskuje” --- systemy oparte wyłącznie na regułach zdefiniowanych przez człowieka są wyłączone (motyw 12); wyłączenie zastosowań wyłącznie wojskowych, obronnych i bezpieczeństwa narodowego jest zastosowaniowe, nie podmiotowe (art. 2 ust. 3), B+R przed wprowadzeniem do obrotu wyłączone poza testami w warunkach rzeczywistych (art. 2 ust. 8), systemy specjalnie opracowane i oddane do użytku wyłącznie do badań naukowych i rozwojowych (ust. 6); stosowanie od 2.08.2026, rozdziały I--II od 2.02.2025, art. 6 ust. 1 od 2.08.2027 (art. 113 --- brzmienie pierwotne, możliwe przesunięcie: ?); wersja 1.2 twierdziła, że obowiązki z art. 6 ust. 1 obejmą cywilne BSP bezpośrednio przez 2018/1139 (załącznik I sekcja B pkt 20) --- jest: do produktów z sekcji B stosuje się wyłącznie art. 6 ust. 1, 102--109 i 112 (art. 2 ust. 2), a wymogi rozdziału III sekcja 2 mają być wprowadzone aktami do 2018/1139 (art. 108), więc AI Act dotyczy cywilnego BSP głównie pośrednio; do proponowanego brzmienia §~28 ust. 1 dodano odesłanie do 2024/1689 (w zakresie, w jakim ma zastosowanie), a do programu zgodności obowiązki z art. 11 ust. 9 i 27 rozporządzenia 2021/821.
  \item (4.2.2) Było: pozycje 9A012, 9A120, kategorie 6/7, 5A002 z wiedzy ogólnej, hipoteza „9A012.b.2 (autonomia) z 9D/9E”; jest: w tekście skonsolidowanym 2021/821 (26.05.2023) 9A012.b.1--b.2 są „Not used”, a 9A120 to zbiorniki paliwa rakietowego (błędne powołanie; właściwa pozycja BSP o zasięgu 300 km to 9A112.a) --- oprogramowanie autonomii jest kontrolowane przez 9D004.e (oprogramowanie do działania BSP z 9A012.a: BVLOS i czas lotu co najmniej 30 min z odpornością na 25 kn albo co najmniej 1 h), 9D001/9D002, technologię 9E001, a niezależnie przez kategorię 7 (7D002, 7D003.a--b, 7D004 lit. a i d, 7E004.b.3--4; 7A003 uwaga 2 wyłącza urządzenia certyfikowane cywilnie); po weryfikacji sceptycznej poprawiono opis 7D003.b (wersja 1.2: „dane inercyjne z radarem dopplerowskim czy nawigacją obrazową”; tekst: ciągłe łączenie danych o kursie z prędkością z radaru dopplerowskiego lub sonaru, z GNSS albo z DBRN --- mapy terenu 3D, grawimetryczne, magnetyczne, gwiazdowe) i 7D004; 9A112.b (autonomia + rozpylacz ponad 20 l) bez znaczenia; GTN/GSN i definicje „wymaganej” technologii, „powszechnie dostępnych” informacji i „podstawowych badań naukowych” z tekstu; teza „żadna pozycja dotycząca Spółki nie jest w załączniku IV” zawężona: załącznik IV obejmuje 1C101 (redukcja wykrywalności, także BSP z 9A012) i 1D103, więc transfer wewnątrzunijny bez zezwolenia (art. 11 ust. 1, z oznaczeniem dokumentów, ust. 9) tylko gdy platforma nie ma takich cech.
  \item (4.2.3) Było: „udostępnienie technologii osobie spoza obszaru celnego UE jest eksportem” (W); jest: potwierdzone art. 2 pkt 2 lit. d i pkt 3 lit. b 2021/821 (tekst EN) oraz dodano, że katalog UE/EOG/NATO w §~27 ust. 6 nie pokrywa się z EU001; po weryfikacji sceptycznej poprawiono: z państw NATO/EOG spoza UE EU001 obejmuje Norwegię, Islandię, Liechtenstein, Kanadę, USA i UK (Szwajcaria nie należy do NATO ani EOG; wersja 1.2 zaliczała ją do tej grupy), pomoc techniczna wymaga zezwolenia tylko w sytuacjach catch-all (art. 8 ust. 1--2), a zwolnienie terytorialne z art. 8 ust. 3 lit. a odsyła do listy EU001 (wersja 1.2: „zwolniona tylko wobec tych państw”); dostęp z innego państwa NATO (np. Turcji) do technologii z załącznika I jest wywozem wymagającym zezwolenia indywidualnego, globalnego albo krajowego generalnego (art. 3 ust. 1, art. 12 ust. 1) --- EU002--EU006 obejmujące Turcję nie dotyczą tej technologii (EU004 wyłącza sekcje D i E); zgoda Rady musi być poprzedzona analizą kontroli eksportu.
  \item (4.4.1) Było: wymogi EDF za stroną kancelarii (web\_edf\_gowling.txt), wzory term sheet PFR Ventures niedostępne [W]; jest: art. 9 ust. 1–4 i 7 oraz art. 2 pkt 6 rozp. 2021/697 z tekstu PL; strony programów PFR Ventures potwierdzają strukturę funduszy G (udział PFR do 80\% w Starterze i do 60\% w KOFFI oraz w OI w strukturze standardowej, wkład prywatny min. 20\%/40\%, rejestracja ZASI i zdywersyfikowani LP według strony KOFFI, OI także dla technologii obronnych i dual-use), ale w modelach koinwestycyjnych (Biznest; OI model 2 z udziałem PFR do 97\% kapitalizacji) prywatne min. 40\% wchodzi bezpośrednio do spółki, więc aniołowie i koinwestorzy podlegają Kryterium osobno; rozp. 2026/1386: inwestorem zagranicznym jest też osoba fizyczna bez obywatelstwa państwa członkowskiego, a spółka zależna w Unii jest kanałem inwestycji zagranicznej (art. 2 pkt 1, 5 i 7), nie samym inwestorem; ujawnienie struktury własności wynika z art. 15 ust. 1 lit. b i ust. 3 oraz art. 19 ust. 2 lit. e.
  \item (4.4.2) Było: jedynym państwem stowarzyszonym z EDF Norwegia (za stroną kancelarii), warunek zawieszający z ustawy o kontroli inwestycji tylko dla inwestora spoza UE; jest: art. 5 rozp. 2021/697 określa państwa stowarzyszone ogólnie (EFTA/EOG), Norwegia pozostaje tezą ze źródła internetowego; dwa stany prawne warunku zawieszającego — do 16.01.2028 r. (inwestor spoza UE, EOG i OECD, co najmniej 20\% głosów albo dominacja, przychód w Polsce ponad 10 mln euro, art. 12a, 12c i 12d ustawy) i od 17.01.2028 r. (rozp. 2026/1386: uprzednia zgoda, gdy Spółka rozwija, wytwarza lub wprowadza do obrotu produkty z zał. I do 2021/821 lub z załącznika do dyr. 2009/43/WE albo rozwija technologie półprzewodnikowe, kwantowe lub AI z zał. I do 2026/1386, art. 4 ust. 15 lit. a–c; wobec każdego inwestora spoza UE i funduszu unijnego kontrolowanego spoza UE, ale tylko przy efektywnym udziale w zarządzaniu lub kontroli, bez inwestycji portfelowych, art. 2 pkt 1 i motywy 14–15; 45 dni; notyfikacja przy inwestorze kontrolowanym przez rząd państwa trzeciego, a przy projekcie EDF tylko po wszczęciu badania pogłębionego, art. 5–6; reżim przejściowy z art. 30); dodano obowiązek zgłaszania Komisji zmian z art. 9 ust. 7 EDF przy zawężeniu ust. 15.
  \item (T1) Było: wiersz 2026/877 ogólny (oznaczenie Z bez tekstu); jest: art. 1–5, 8–11 z tekstu EN, w tym definicja umowy między dwoma przedsiębiorstwami (art. 1 ust. 1 lit. c pkt i), katalog najcięższych ograniczeń z podziałem na konkurentów i niekonkurentów (art. 4 ust. 1–2), wykluczenie każdego bezpośredniego lub pośredniego obowiązku wyłącznego grant-backu albo przeniesienia własnych ulepszeń bez rozróżnienia ulepszeń wydzielnych i niewydzielnych (art. 5 ust. 1 lit. a — także wariant z odpowiedzi 2.4.1 wymaga oceny indywidualnej) i okres przejściowy tylko dla umów obowiązujących 30.04.2026 r.
  \item (T1) Było: AI Act, 2023/1230 i CRA w jednym wierszu (oznaczenie W) z tezą o „wyłączeniu zastosowań wojskowych”; jest: osobny wiersz AI Act (oznaczenie Z, tekst EN): wyłączenie z art. 2 ust. 3 akapit drugi tylko dla użycia „wyłącznie” wojskowego/obronnego/bezpieczeństwa narodowego, wyłączenia B+R z ust. 6 i 8, oprogramowanie autonomii jako system AI tylko przy spełnieniu całej definicji z art. 3 pkt 1 (wnioskowanie, nie sama autonomia), BSP jako produkt z załącznika I sekcja B (2018/1139) pod art. 6 ust. 1 i art. 2 ust. 2, wymogi przez art. 108, pełne terminy z art. 113 (z rozdziałem III sekcja 4, art. 78 i wyjątkiem art. 101) z zastrzeżeniem, że plik zawiera tekst pierwotny.
  \item (T1) Dodano wiersz 2019/947 (kategorie operacji art. 3–6, warunki kategorii otwartej z art. 4 i pkt UAS.OPEN.060 — VLOS i zdolność sterowania, zezwolenie, scenariusz standardowy albo LUC z art. 5, rejestracja operatora art. 14 ust. 5–6, stosowanie od 1.07.2020 r. wg tekstu pierwotnego); organ właściwy — Prezes ULC (art. 21 ust. 2g prawa lotniczego); wpływ tylko na procedury (bramka G1).
  \item (T1) Było: wiersz 2019/452 bez wzmianki o następcy, polski reżim kontroli opisany ogólnie z tezą, że katalog sektorów obejmuje wyroby podwójnego zastosowania (oznaczenie W); jest: uchylenie 2019/452 z dniem 17.01.2028 r. (art. 30 ust. 1–2 rozporządzenia 2026/1386) oraz ustawa o kontroli niektórych inwestycji sprawdzona na tekście (oznaczenie Z): dwa reżimy (wykaz RM z art. 4; art. 12a–12k dla inwestorów spoza UE/EOG/OECD), katalog z art. 12d ust. 3 obejmuje wyroby i technologię o przeznaczeniu wojskowym lub policyjnym, nie wymienia dual-use, i wymaga przychodu w Polsce powyżej 10 mln EUR (art. 12d ust. 4); nowy wiersz 2026/1386 z tekstu EN: obowiązkowy mechanizm w każdym państwie (art. 3), minimalny zakres uprzedniego zezwolenia (art. 4 ust. 15, w AI tylko modele ogólnego przeznaczenia z zał. I pkt 3), definicja inwestora zagranicznego bez wyjątku dla EOG/NATO i objęcie inwestycji przez unijne spółki zależne (art. 2 pkt 1, 5, 7), procedura 45 dni (art. 4 ust. 2), kontrola ex post 15 mies.–5 lat dla inwestycji spoza obowiązku zezwolenia i co najmniej 24 mies. dla niezgłoszonych objętych obowiązkiem (art. 4 ust. 4–5), notyfikacja z art. 5 ust. 1 tylko dla inwestycji z art. 4 ust. 15, mechanizm współpracy (art. 5–8, 15), czynniki oceny z art. 19 ust. 1 i zał. III (drony, systemy AI, nawigacja i awionika), terminy z art. 31; skutek dla §~11, §~12 ust. 2 lit. e, §~31–32 i umowy akcjonariuszy w dwóch stanach prawnych.
  \item (T1a) Pkt 3: art. 9 EDF potwierdzony na tekście PL (ust. 3, 4 lit. c, 7) wraz z art. 5 (państwa stowarzyszone = EFTA/EOG) i definicją kontroli z art. 2 pkt 6; pkt 4: pozycje 9A012 i kategoria 7 potwierdzone w załączniku I do 2021/821 (wersja z 26.05.2023 r.), a klasyfikacja dual-use przesądza też o zakresie obowiązkowej kontroli inwestycji z art. 4 ust. 15 lit. a (i lit. b dla wykazu uzbrojenia) rozporządzenia 2026/1386.
  \item (T2/T6) Argumentacja dla sądu rejestrowego powołuje teraz art. 9 ust. 3–4 rozporządzenia 2021/697 (oznaczenie Z; TFUE nadal W) i — tylko jeżeli Spółka mieści się w art. 4 ust. 15 lit. a lub b rozporządzenia 2026/1386 — zbieżność Kryterium z obowiązkową kontrolą inwestycji od 17.01.2028 r.
  \item (T7) Było: „fundusze PFR Deep Tech” i katalog postanowień rundy bez źródła; jest: program PFR Otwarte Innowacje (B+R, w tym obronność i dual-use; ticket od 5 mln zł do 15\% budżetu; min. 10\% w spółce) z opisem niespójności strony co do spółek zagranicznych (opis programu: do 40\% portfela w spółki zagraniczne z polskimi founderami; FAQ: co najmniej 85\% w Polsce i do 15\% w UE/EFTA/EOG+UK przy znaczącej obecności w Polsce) — brzmienie oznaczone Z, wiążący limit ?; Starter/Biznest z oznaczeniem Z na stronach PFR (nabory wstrzymane); katalog umowy inwestycyjnej z artykułu DLA Piper na startup.pfr.pl (typowy, nie wzorcowy; większość 75–80\%) i dane o term sheetach (95\% rola inwestora, 43\% anti-dilution rozumiana w artykule jako pro-rata, wg PwC Raise); oczekiwania funduszy obronnych pozostają W.
  \item (T8) Dodano wiersz o postanowieniach dotyczących kontroli inwestycji zagranicznych (2026/1386: uprzednie zezwolenie na inwestycję spoza UE dającą udział w zarządzaniu lub kontroli, art. 2 pkt 1, art. 4 ust. 9 i 15, art. 15 ust. 1 i 3) i odesłanie do art. 5 rozporządzenia 2026/877 dla wzoru umowy JV.
\end{itemize}

\clearpage
\section*{Załącznik nr 6 --- Indeks pytań}
\addcontentsline{toc}{section}{Załącznik nr 6 --- Indeks pytań}

Numer pytania z listy szczegółowej (\texttt{law.md}), jego początek, klasyfikacja uwagi i strona, na której odpowiedź memorandum stoi przy postanowieniu umowy.

{\small
\begin{longtable}{@{}p{1.2cm} p{9.6cm} p{2.3cm} p{1.4cm}@{}}
\toprule
\textbf{Pyt.} & \textbf{Pytanie} & \textbf{Klasa} & \textbf{Strona} \\
\midrule
\endfirsthead
\toprule
\textbf{Pyt.} & \textbf{Pytanie} & \textbf{Klasa} & \textbf{Strona} \\
\midrule
\endhead
1.1.1 & Czy ograniczenie kręgu nabywców akcji, spadkobierców, małżonków i dyrektorów według obywatelstwa, siedziby i struktury kontroli jest \ldots & zalecane & s.~\pageref{memo:1.1.1} \\
1.1.2 & Jak Kryterium ma się do ustawy o kontroli niektórych inwestycji i innych przepisów o ochronie bezpieczeństwa \ldots & opcjonalne & s.~\pageref{memo:1.1.2} \\
1.1.3 & Czy wyłączenie nabycia na cudzy rachunek (§~11 ust.~6) jest wystarczająco określone i egzekwowalne --- w szczególności \ldots & zalecane & s.~\pageref{memo:1.1.3} \\
1.1.4 & Utrata Kryterium po nabyciu (§~11 ust.~15): czy obowiązek zbycia akcji w 6 miesięcy i prawo Spółki \ldots & zalecane & s.~\pageref{memo:1.1.4} \\
1.1.5 & Czy wymóg Kryterium wobec dyrektorów (§~21 ust.~3) jest dopuszczalny? & zalecane & s.~\pageref{memo:1.1.5} \\
1.1.6 & Czy dla skuteczności Kryterium wobec osób trzecich potrzebne jest ujawnienie ograniczenia w rejestrze akcjonariuszy, a jeżeli \ldots & zalecane & s.~\pageref{memo:1.1.6} \\
1.2.1 & Czy art.~300$^{39}$ §~2 KSH („chyba że umowa spółki stanowi inaczej”) pozwala zarówno na modyfikację mechanizmu ustawowego \ldots & opcjonalne & s.~\pageref{memo:1.2.1} \\
1.2.2 & Jak każdy z wariantów zadziała wobec podmiotu prowadzącego rejestr akcjonariuszy (czy odmówi wpisu nabywcy spoza Kryterium, \ldots & zalecane & s.~\pageref{memo:1.2.2} \\
1.2.3 & Wariant~A: czy zapłata ratalna i pierwsza rata w 30 dni odpowiadają wymogowi określenia w umowie ceny \ldots & zalecane & s.~\pageref{memo:1.2.3} \\
1.2.4 & Wariant~B: czy istnieje ryzyko uznania go za nadmierne ograniczenie zbywalności (akcja „uwięziona”, gdy nie ma nabywcy \ldots & konieczne & s.~\pageref{memo:1.2.4} \\
1.2.5 & Który wariant jest w Państwa ocenie skuteczniejszy i bezpieczniejszy dla Spółki, a jak oba będą odbierane \ldots & zalecane & s.~\pageref{memo:1.2.5} \\
1.2.6 & Lock-up (§~13 ust.~1): Założyciel Pierwotny przez 12 miesięcy od wpisu Spółki (Założyciel Serii~F --- od Dnia \ldots & zalecane & s.~\pageref{memo:1.2.6} \\
1.2.7 & §~12 ust.~4 zdanie pierwsze: przy odmowie z powodu Niedopuszczalnego Nabywcy obowiązek wskazania innego nabywcy nie powstaje, \ldots & zalecane & s.~\pageref{memo:1.2.7} \\
1.2.8 & §~12 ust.~7: zgody Spółki nie wymaga zbycie w postępowaniu egzekucyjnym ani nabycie w wykonaniu prawa pierwszeństwa, \ldots & zalecane & s.~\pageref{memo:1.2.8} \\
1.3.1 & Czy 14-dniowy termin, forma dokumentowa i zamknięty katalog przyczyn odmowy są zgodne z art.~300$^{39}$ §~3 KSH \ldots & zalecane & s.~\pageref{memo:1.3.1} \\
1.3.2 & Czy wstrzymanie biegu terminu (lit.~d i~e) jest dopuszczalne, czy też grozi uznaniem za ukrytą odmowę uruchamiającą \ldots & zalecane & s.~\pageref{memo:1.3.2} \\
1.3.3 & Obowiązki informacyjne osoby trzeciej przed nabyciem (§~11 ust.~8): skoro nabywca nie jest jeszcze stroną umowy, czy \ldots & bez & s.~\pageref{memo:1.3.3} \\
1.3.4 & Kolejność prawa pierwszeństwa (§~14) i zgody Spółki (§~12): czy dopuszczenie oferty równoczesnej ze zgłoszeniem nie koliduje \ldots & zalecane & s.~\pageref{memo:1.3.4} \\
1.3.5 & §~14 ust.~3: Spółka może przyjąć ofertę, „o ile nabycie akcji własnych jest prawnie dopuszczalne” --- czy \ldots & zalecane & s.~\pageref{memo:1.3.5} \\
1.4.1 & Czy zachowanie wyjątku uprościłoby obsługę zbyć wtórnych i wejścia funduszy (w tym funduszy z inwestorami spoza \ldots & zalecane & s.~\pageref{memo:1.4.1} \\
1.4.2 & Czy odstąpienie od badania beneficjentów rzeczywistych funduszu jest do pogodzenia z przepisami AML i wymogami zamawiających \ldots & zalecane & s.~\pageref{memo:1.4.2} \\
1.4.3 & Jeżeli wyjątek ma zostać, jak go sformułować, aby nie stał się furtką dla podmiotów spoza Kryterium \ldots & zalecane & s.~\pageref{memo:1.4.3} \\
1.5.1 & Czy obowiązek zbycia jako „obowiązek związany z akcją” jest skuteczny wobec nabywców wtórnych i wobec podmiotu \ldots & konieczne & s.~\pageref{memo:1.5.1} \\
1.5.2 & Czy zbycie za 80\,\% Wartości Godziwej przy Odejściu Zawinionym może zostać zakwestionowane (obejście przepisów o umorzeniu \ldots & zalecane & s.~\pageref{memo:1.5.2} \\
1.5.3 & Umowa wykonawcza: czy nieodwołalne oferty i pełnomocnictwa (art.~101 §~2 i art.~108 KC) mogą zostać skonstruowane tak, \ldots & konieczne & s.~\pageref{memo:1.5.3} \\
1.5.4 & Czy katalog Odejścia Zawinionego (§~\ref{art:zwrotne-zbycie} ust.~5) i tryb jego stwierdzania nie rodzą ryzyka sporu o kwalifikację \ldots & zalecane & s.~\pageref{memo:1.5.4} \\
1.5.5 & Czy mechanizm jest neutralny wobec podstawy zaangażowania założyciela (umowa o pracę, B2B, powołanie) i czy zwolnienie \ldots & zalecane & s.~\pageref{memo:1.5.5} \\
1.5.6 & Przyspieszenie po zmianie kontroli (§~\ref{art:zwrotne-zbycie} ust.~10) --- czy definicja „istotnej przyczyny” i 12-miesięczne okno są wystarczająco \ldots & zalecane & s.~\pageref{memo:1.5.6} \\
1.6.1 & Czy §~\ref{art:dziedziczenie} prawidłowo wykorzystuje art.~300$^{41}$ §~1 KSH (ograniczenie albo wyłączenie wstąpienia pod warunkiem określenia zasad spłaty), \ldots & konieczne & s.~\pageref{memo:1.6.1} \\
1.6.2 & Czy uzależnienie wstąpienia od Kryterium i zgody Walnego Zgromadzenia jest dopuszczalne wobec spadkobierców (w tym ustawowych), \ldots & zalecane & s.~\pageref{memo:1.6.2} \\
1.6.3 & Czy 90-dniowe terminy i mechanizm rezerwowy (obowiązkowe zbycie na rzecz wskazanego nabywcy) są wykonalne wobec spadkobiercy, \ldots & konieczne & s.~\pageref{memo:1.6.3} \\
1.6.4 & Wspólny przedstawiciel (§~\ref{art:dziedziczenie} ust.~3) --- spójność z przepisami KSH o współuprawnionych z akcji. & bez & s.~\pageref{memo:1.6.4} \\
1.7.1 & Czy rozdzielenie sfery majątkowej (przynależność akcji do majątku wspólnego) od korporacyjnej (wykonywanie praw z akcji przez \ldots & zalecane & s.~\pageref{memo:1.7.1} \\
1.7.2 & Czy obowiązki informacyjne z §~\ref{art:malzonek} ust.~1 i~5 są egzekwowalne i zgodne z ochroną danych? & opcjonalne & s.~\pageref{memo:1.7.2} \\
1.7.3 & Czy konstrukcja Prawnie Niedopuszczalnego Posiadacza (bez automatycznego przymusowego zbycia) jest spójna z sankcjami i kontrolą eksportu, \ldots & konieczne & s.~\pageref{memo:1.7.3} \\
1.8.1 & Czy umowna definicja może być stosowana we wszystkich mechanizmach projektu (zwrotne zbycie, dziedziczenie, utrata Kryterium, wskazanie \ldots & opcjonalne & s.~\pageref{memo:1.8.1} \\
1.8.2 & Czy wiążący charakter wyceny eksperta (§~\ref{art:wycena} ust.~7) i tryb jego wyboru wytrzymają spór sądowy między wspólnikami? & konieczne & s.~\pageref{memo:1.8.2} \\
2.1.1 & Czy większości przy powołaniu i odwołaniu oraz przedłużenie mandatu do 6 miesięcy są zgodne z przepisami \ldots & zalecane & s.~\pageref{memo:2.1.1} \\
2.1.2 & Czy wymóg pozytywnego głosu określonego dyrektora (§~23 ust.~2) jest dopuszczalny (uprzywilejowanie głosu w Radzie) i jaki \ldots & zalecane & s.~\pageref{memo:2.1.2} \\
2.1.3 & Czy podział na dyrektorów wykonawczych i niewykonawczych oraz delegowanie prowadzenia spraw (§~21 ust.~6--8) odpowiada art.~300$^{76}$ KSH \ldots & opcjonalne & s.~\pageref{memo:2.1.3} \\
2.2.1 & Czy sformułowania o udziale elektronicznym są zgodne z art.~300$^{88}$ i 300$^{92}$ KSH (miejsce zgromadzenia a uczestnictwo \ldots & zalecane & s.~\pageref{memo:2.2.1} \\
2.2.2 & Czy zawiadomienie e-mailem na adres z rejestru akcjonariuszy spełnia art.~300$^{87}$ KSH i co, gdy akcjonariusz nie \ldots & zalecane & s.~\pageref{memo:2.2.2} \\
2.2.3 & Czy dwa sposoby liczenia większości (od głosów uprawnionych w sprawie i od wszystkich głosów) są jasne \ldots & zalecane & s.~\pageref{memo:2.2.3} \\
2.2.4 & Czy katalog Spraw Zastrzeżonych i próg 75\% nie blokują czynności, dla których ustawa przewiduje inną większość \ldots & zalecane & s.~\pageref{memo:2.2.4} \\
2.3.1 & Czy uchwała Walnego Zgromadzenia może zastąpić uchwałę Rady w sprawach, dla których ustawa sama wymaga uchwały \ldots & opcjonalne & s.~\pageref{memo:2.3.1} \\
2.3.2 & Czy quorum liczone od pełnego składu (§~21 ust.~9) i wyłączenie dyrektorów (§~29 ust.~7) są spójne z \ldots & zalecane & s.~\pageref{memo:2.3.2} \\
2.3.3 & Czy przy umowie i sporze z dyrektorem projekt prawidłowo odsyła do art.~300$^{79}$ KSH (§~29 ust.~12), a \ldots & bez & s.~\pageref{memo:2.3.3} \\
2.3.4 & §~30: jak pojęcie „godziwych warunków” ma się do ustawowych ograniczeń świadczeń na rzecz akcjonariuszy (art.~300$^{15}$ i \ldots & zalecane & s.~\pageref{memo:2.3.4} \\
2.4.1 & Czy licencja niewyłączna ograniczona do zakresu przedsięwzięcia z prawami do ulepszeń dla Spółki jest zgodna z \ldots & zalecane & s.~\pageref{memo:2.4.1} \\
2.4.2 & Czy uzależnienie udziału partnera od Kryterium jest dopuszczalne, w tym wobec partnerów z państw NATO spoza \ldots & zalecane & s.~\pageref{memo:2.4.2} \\
2.4.3 & Czy wyłączenie udziału założycieli w przedsięwzięciach spod zakazu konkurencji (§~33 ust.~7) nie osłabia §~29 i jak \ldots & zalecane & s.~\pageref{memo:2.4.3} \\
2.4.4 & Czy progi kwotowe i podział kompetencji Rada / Walne Zgromadzenie są wykonalne przy zmianach zaangażowania w \ldots & zalecane & s.~\pageref{memo:2.4.4} \\
3.1.1 & Jakie dokumenty i wyceny wkładów niepieniężnych są potrzebne do zawiązania i wpisu (art.~300$^{2}$ KSH), kto ponosi \ldots & zalecane & s.~\pageref{memo:3.1.1} \\
3.1.2 & Czy przeznaczenie na kapitał akcyjny wyłącznie wkładów pieniężnych (a wkładów niepieniężnych poza kapitał) jest dopuszczalne i \ldots & konieczne & s.~\pageref{memo:3.1.2} \\
3.1.3 & Czy oświadczenia z §~\ref{art:wlasnosc-intelektualna} ust.~5 i obowiązek współdziałania z §~\ref{art:wlasnosc-intelektualna} ust.~7 wystarczą, aby przed rundą wykazać \ldots & zalecane & s.~\pageref{memo:3.1.3} \\
3.1.4 & Załącznik nr~1 zawiera pola do uzupełnienia --- jakie dane muszą być kompletne w chwili podpisania umowy \ldots & konieczne & s.~\pageref{memo:3.1.4} \\
3.2.1 & Czy „uchwała ramowa” na wiele transz emisji serii P, wykonywana przez Radę, oraz jedna uchwała 4/5 \ldots & konieczne & s.~\pageref{memo:3.2.1} \\
3.2.2 & Czy określenie w umowie maksymalnej liczby akcji i terminu (2036~r.) dla serii P i F spełnia \ldots & opcjonalne & s.~\pageref{memo:3.2.2} \\
3.2.3 & Czy uchwała kwalifikacyjna serii F (bez oferty i bezwarunkowego prawa do akcji, §~\ref{art:seria-f} ust.~3) nie rodzi \ldots & opcjonalne & s.~\pageref{memo:3.2.3} \\
3.2.4 & Aspekty podatkowe programu motywacyjnego w P.S.A. (zakres E) --- zakres i koszt interpretacji indywidualnej. & bez & s.~\pageref{memo:3.2.4} \\
3.3.1 & Czy zobowiązanie do głosowania i współdziałania (§~\ref{art:kwalifikowana-runda} ust.~3) jest skuteczne wobec akcjonariuszy i egzekwowalne, czy pozostaje \ldots & zalecane & s.~\pageref{memo:3.3.1} \\
3.3.2 & Czy warunki inwestorskie z §~\ref{art:kwalifikowana-runda} ust.~4 mogą zostać przyznane bez zmiany umowy spółki (uprzywilejowanie akcji, prawa \ldots & konieczne & s.~\pageref{memo:3.3.2} \\
3.3.3 & Jak pogodzić Kwalifikowaną Rundę z prawem poboru (§~\ref{art:kwalifikowana-runda} ust.~5, §~\ref{art:walne-zastrzezone} ust.~2) i z wymogiem Kryterium wobec \ldots & zalecane & s.~\pageref{memo:3.3.3} \\
3.3.4 & Czy ograniczenia z §~\ref{art:finansowanie} ust.~3--4 (instrumenty zamienne, zabezpieczenia na IP) nie utrudnią finansowania dłużnego i grantowego \ldots & zalecane & s.~\pageref{memo:3.3.4} \\
3.4.1 & Czy aktualizacja Załączników nr~1--3 (np. zmiana danych, wpis nowego Założyciela Serii F, ujawnione aktywności) zawsze wymaga \ldots & zalecane & s.~\pageref{memo:3.4.1} \\
3.4.2 & Czy brak klauzul zgody indywidualnej (ochrony konkretnego akcjonariusza przed zmianą) jest bezpieczny dla założycieli mniejszościowych i \ldots & zalecane & s.~\pageref{memo:3.4.2} \\
3.4.3 & Czy tryb impasu (§~\ref{art:impas}) jest wykonalny i czy rozstrzygnięcie Niezależnego Eksperta w sprawach obiektywnych nie wkracza \ldots & opcjonalne & s.~\pageref{memo:3.4.3} \\
3.4.4 & Przekreślone zdania i ustęp: czy ich usunięcie jest bezpieczne, czy któreś powinny zostać; spójność definicji po \ldots & zalecane & s.~\pageref{memo:3.4.4} \\
4.1.1 & Czy powołana podstawa (ustawa z 13 czerwca 2019 r. o wykonywaniu działalności gospodarczej w zakresie wytwarzania \ldots & zalecane & s.~\pageref{memo:4.1.1} \\
4.1.2 & Czy ujawnienie kodów wojskowych PKD przed uzyskaniem koncesji jest zasadne, czy lepiej dodać je przy zmianie \ldots & opcjonalne & s.~\pageref{memo:4.1.2} \\
4.1.3 & Czy zestaw kodów PKD 2025 wymaga końcowej kontroli klasyfikacyjnej przed wnioskiem do KRS i czy liczba \ldots & bez & s.~\pageref{memo:4.1.3} \\
4.2.1 & Czy klauzule wewnętrznej zgodności wystarczają jako ramy, czy projekt powinien wprost odsyłać do konkretnych aktów (rozporządzenie \ldots & zalecane & s.~\pageref{memo:4.2.1} \\
4.2.2 & Czy klasyfikacja produktów Firmy (oprogramowanie autonomii, roje dronów) wymaga odrębnej weryfikacji specjalistycznej i czy mogą Państwo \ldots & opcjonalne & s.~\pageref{memo:4.2.2} \\
4.2.3 & Czy 10-letni okres poufności i ograniczenie dostępu spoza UE/NATO (§ 27 ust. 3 i 6) są \ldots & zalecane & s.~\pageref{memo:4.2.3} \\
4.3.1 & Jaka forma zawiązania i jakie dokumenty do KRS są wymagane (w tym dokumentacja wkładów niepieniężnych) i \ldots & konieczne & s.~\pageref{memo:4.3.1} \\
4.3.2 & Które ograniczenia i obowiązki powinny zostać ujawnione w rejestrze akcjonariuszy i jak dobrać podmiot prowadzący rejestr \ldots & opcjonalne & s.~\pageref{memo:4.3.2} \\
4.3.3 & Czy nowelizacja KSH z 23 stycznia 2026 r. (wejście w życie 18 lutego 2027 r.) wpływa \ldots & zalecane & s.~\pageref{memo:4.3.3} \\
4.4.1 & Które postanowienia (Kryterium, mechanizm przy odmowie zgody, zwrotne zbycie, drag-along 75\,\%, Sprawy Zastrzeżone 75\,\%, ograniczenia finansowania) \ldots & zalecane & s.~\pageref{memo:4.4.1} \\
4.4.2 & Jakie zmiany ułatwiłyby dostosowanie umowy do rundy bez naruszenia jej założeń (bezpieczeństwo właścicielskie, kontrola eksportu, mechanizm \ldots & zalecane & s.~\pageref{memo:4.4.2} \\
\bottomrule
\end{longtable}
}

\end{document}
